%%%%

\documentclass[11pt,a4paper,twoside]{article}

\usepackage[utf8]{inputenc}

\usepackage[T1,T2A]{fontenc}

\usepackage[english.ancient,russian.ancient]{babel}

\usepackage{graphicx}

\usepackage{array}

\usepackage[doi]{samgtu-bib}

\usepackage{samgtu}

\usepackage{samgtu-name}

\usepackage{mathtools}

\mathtoolsset{showonlyrefs}

\usepackage{desclist}

\usepackage{euscript}

\usepackage{mathrsfs} 

\usepackage{multicol}

\usepackage{mathrsfs}

\usepackage{cite}

\usepackage{cmap}

\usepackage{qrcode}

\allowdisplaybreaks[4]

\usepackage[nodayofweek]{datetime}

\usepackage[thinlines]{easybmat}



\renewcommand{\JournalYear}{2018}

\renewcommand{\JournalVolume}{22}

\renewcommand{\JournalNumber}{4}



\newdate{JournalDateSign}{29}{12}{2018}% Дата подписания Вестника в печать

\renewcommand{\JournalDateSign}{\displaydate{JournalDateSign}}

\newdate{JournalDatePrint}{16}{01}{2019}% Дата выхода Вестника из печати

\renewcommand{\JournalDatePrint}{\displaydate{JournalDatePrint}}

%\renewcommand{\JournalUPL}{16.56} % Усл. печ. л.

%\renewcommand{\JournalUIL}{16.53} % Уч.-изд. л.

\renewcommand{\JournalUPL}{15.24} % Усл. печ. л.

\renewcommand{\JournalUIL}{15.21} % Уч.-изд. л.

\renewcommand{\JournalRegNumber}{220/18} % Регистрационный номер

\renewcommand{\JournalOrderNumber}{2} % Номер заказа



%\renewcommand{\JournalRMonth}{Март} % Месяц выхода Вестника

%\renewcommand{\JournalEMonth}{March} % Месяц выхода Вестника

%\renewcommand{\JournalRMonth}{Июнь} % Месяц выхода Вестника

%\renewcommand{\JournalEMonth}{June} % Месяц выхода Вестника

%\renewcommand{\JournalRMonth}{Сентябрь} % Месяц выхода Вестника

%\renewcommand{\JournalEMonth}{September} % Месяц выхода Вестника

\renewcommand{\JournalRMonth}{Декабрь} % Месяц выхода Вестника

\renewcommand{\JournalEMonth}{December} % Месяц выхода Вестника



\newcommand{\totop}[1]{\raisebox{6pt}[1pt][2pt]{#1}}

\newcommand{\tottop}[1]{\raisebox{12pt}[1pt][2pt]{#1}} 

\newcommand{\totttop}[1]{\raisebox{18pt}[1pt][2pt]{#1}} 

\newcommand{\tottttop}[1]{\raisebox{24pt}[1pt][2pt]{#1}} 

\newcommand{\totttttop}[1]{\raisebox{30pt}[1pt][2pt]{#1}} 

\newcommand{\tobot}[1]{\raisebox{-6pt}[1pt][2pt]{#1}} 

\newcommand{\tobbot}[1]{\raisebox{-12pt}[1pt][2pt]{#1}} 

\newcommand{\tobbbot}[1]{\raisebox{-18pt}[1pt][2pt]{#1}}



\graphicspath{{./00PIC/}}





\usepackage{rotating_pr}

\usepackage{desclist}

\usepackage{cmap}

\usepackage{textcomp}

\usepackage{nccboxes}

\usepackage{euscript}

\usepackage{mathrsfs} 

\usepackage{multicol}

\usepackage{nccfoots}

\usepackage{mathrsfs}

\usepackage{calrsfs}

\allowdisplaybreaks[4]

\usepackage{qrcode}

\usepackage{afterpage}

\usepackage{wasysym}

\usepackage{wrapfig}

\usepackage{slashbox}



\usepackage{multirow}

\newcommand{\specialcell}[2][c]{%

  \begin{tabular}[#1]{@{}c@{}}#2\end{tabular}}





\usepackage{pgf,pgfplots,tikz}

\usetikzlibrary{arrows,snakes,backgrounds,patterns,shapes.geometric,calc}



\nofiles



\oek % для генерации pdf, отдаваемого в oek.rsl.ru

%\samgtupubl % для генерации pdf, отдаваемого в типографию СамГТУ

%\pdfcrop % обрезка pdf для размещения на math-net.ru

%\draft % На корректуре





%\DeclareMathOperator{\tr}{tr}

%\DeclareMathOperator{\ink}{Ink}

%\DeclareMathOperator{\erf}{erf}







%\makeatletter

%\newcommand{\PersonRadRef}{\@langrustrue\def\refname{{\bf\small Избранные  труды \rupid{19071}{О.~А.~Репина} за последние 5 лет}}}

%\makeatother



\DeclareMathOperator{\tr}{tr}

\DeclareMathOperator{\arcsh}{arcsh}









\usepackage{listings}

\usepackage{color}

\definecolor{dkgreen}{rgb}{0,0.6,0}

\definecolor{gray}{rgb}{0.5,0.5,0.5}

\definecolor{mauve}{rgb}{0.58,0,0.82}



\renewcommand*{\lstlistingname}{Листинг}

\usepackage{caption}

%\DeclareCaptionFont{white}{\color{white}} 

\DeclareCaptionFormat{listing}{\parbox{\textwidth}{\centering \small  #1.~#3}}

\captionsetup[lstlisting]{format=listing}



\lstset{ %

  inputencoding=utf8,

  language=R,                     % the language of the code

  basicstyle=\footnotesize,       % the size of the fonts that are used for the code

  numbers=left,                   % where to put the line-numbers

  numberstyle=\tiny\color{gray},  % the style that is used for the line-numbers

  stepnumber=1,                   % the step between two line-numbers. If it's 1, each line

                                  % will be numbered

  numbersep=5pt,                  % how far the line-numbers are from the code

  backgroundcolor=\color{white},  % choose the background color. You must add \usepackage{color}

  showspaces=false,               % show spaces adding particular underscores

  showstringspaces=false,         % underline spaces within strings

  showtabs=false,                 % show tabs within strings adding particular underscores

  frame=single,                   % adds a frame around the code

  rulecolor=\color{black},        % if not set, the frame-color may be changed on line-breaks within not-black text (e.g. commens (green here))

  tabsize=2,                      % sets default tabsize to 2 spaces

  captionpos=t,                   % sets the caption-position to bottom

  breaklines=true,                % sets automatic line breaking

  breakatwhitespace=false,        % sets if automatic breaks should only happen at whitespace

  %title=\lstname,                 % show the filename of files included with \lstinputlisting;

                                  % also try caption instead of title

  keywordstyle=\color{blue},      % keyword style

  commentstyle=\color{dkgreen},   % comment style

  stringstyle=\color{mauve},      % string literal style

  %escapeinside={\%*}{*)},         % if you want to add a comment within your code

  morekeywords={*,in,+,-,=}            % if you want to add more keywords to the set

} 







\begin{document}







\setcounter{page}{601}

%\renewcommand{\bfdefault}{bx}
\thispagestyle{empty}
\begin{table}[t!]
\centering
\begin{tabular}{p{7cm}p{6cm}}
{\fontsize{22pt}{12pt}\selectfont\vspace*{0ex}
\bf~~Вестник  \smallskip \par
\bf~~Самарского \smallskip \par
\bf~~государственного \smallskip \par
\bf~~технического\phantom{p} \smallskip \par
\bf~~университета
\par\bigskip} &
{\fontsize{12pt}{8pt}\selectfont\vspace*{0ex}
\phantom{1}\medskip\par
НАУЧНЫЙ ЖУРНАЛ \medskip\par
Издаётся с 1996 г. \medskip\par
Выходит 4 раза в год \bigskip\par
\phantom{1}\medskip\par
\JournalRMonth\ --- \JournalYear}
\\
\multicolumn{2}{p{12cm}}{\fontsize{12pt}{12pt}\selectfont\vspace*{0ex} \bf~~\;Серия}\\
\multicolumn{2}{c}{\fontsize{11pt}{9pt}\selectfont\vspace*{0ex} 
\bf <<ФИЗИКО-МАТЕМАТИЧЕСКИЕ НАУКИ>> (\RISSUE)}\smallskip\par  \\
\hline
\multicolumn{2}{p{13cm}}{
\fontsize{9pt}{8pt}\selectfont\vspace*{0ex}
\par \textbf{Главный редактор} В.\,П.~Радченко  --- д.ф.-м.н., проф. (Самара, Россия)
\par \textbf{Заместитель главного редактора} А.\,И.~Жданов  --- д.ф.-м.н., проф. (Самара, Россия)
\par \textbf{Отв. секретарь} М.\,Н. Саушкин --- к.ф.-м.н., доц. (Самара, Россия)
\par \textbf{Cекретарь} Е.\,А. Козлова --- к.ф.-м.н. (Самара, Россия)
\smallskip\par
\textbf{Редакционный совет:}\par
$~~\bullet$ С.\,А.~Авдонин --- д.ф.-м.н., проф. (Фэрбенкс, Аляска, США)  \par
$~~\bullet$ Б.\,Д.~Аннин --- акад. РАН, д.ф.-м.н., проф. (Новосибирск, Россия)\par
$~~\bullet$ П.\,Б.~Бабаджанов  --- акад. АН РТ, д.ф.-м.н. проф. (Душанбе, Таджикистан) \par
$~~\bullet$ А.\,А.~Буренин --- чл. кор. РАН, д.ф.-м.н., проф. (Комсомольск-на-Амуре, Россия)\par
$~~\bullet$ Р.~Вельмураган  --- доктор наук (Ченнай, Индия) \par
$~~\bullet$ С.\,П.~Верёвкин  --- д.х.н., проф. (Росток, Германия) \par
$~~\bullet$ Н.\,А.~Вирченко --- д.ф.-м.н., проф. (Киев, Украина) \par
$~~\bullet$ Ш.~Иматани --- доктор наук (Киото, Япония) \par
$~~\bullet$ О.\,И.~Маричев  д.ф.-м.н., проф. (Ларами, Вайоминг, США) \par
$~~\bullet$ В.\,П.~Матвеенко --- акад. РАН, д.т.н., проф. (Пермь, Россия) \par
$~~\bullet$ П.\,В.~Севастьянов  --- д.т.н., проф. (Ченстохова, Польша) \par
$~~\bullet$ З.\,Д.~Усманов  --- акад. АН РТ, д.ф.-м.н., проф. (Душанбе, Таджикистан) \par
\smallskip \par
\textbf{Редакционная коллегия:}\par 
$~~\bullet$ В.\,Н.~Акопян --- д.ф.-м.н., проф. (Ереван, Армения) \par
$~~\bullet$ А.\,П.~Амосов --- д.ф.-м.н., проф. (Самара, Россия) \par 
$~~\bullet$ В.\,А.~Ковалёв --- д.ф.-м.н., проф. (Москва, Россия) \par
$~~\bullet$ А.\,И.~Кожанов --- д.ф.-м.н., проф. (Новосибирск, Россия) \par 
$~~\bullet$ В.\,А.~Кудинов --- д.ф.-м.н., проф. (Самара, Россия) \par 
%$~~\bullet$ \fbox{А.\,В.~Манжиров} --- д.ф.-м.н., проф. (Москва, Россия) \par
$~~\bullet$ А.\,М.~Нахушев --- д.ф.-м.н., проф. (Нальчик, Россия) \par
$~~\bullet$ А.\,А. Пимерзин  --- д.х.н., проф. (Самара, Россия) \par
$~~\bullet$ Л.\,С.~Пулькина --- д.ф.-м.н., проф. (Самара, Россия) \par 
$~~\bullet$ Ю.\,Н.~Радаев --- д.ф.-м.н., проф. ((Москва, Россия) \par 
$~~\bullet$ Е.\,В.~Радкевич --- д.ф.-м.н., проф. (Москва, Россия) \par
%$~~\bullet$ \fbox{О.\,А.~Репин} --- д.ф.-м.н., проф. (Самара, Россия)\par  
$~~\bullet$  А.\,В.~Саакян --- д.ф.-м.н., проф. (Ереван, Армения) \par
$~~\bullet$ К.\,Б.~Сабитов --- д.ф.-м.н., проф. (Стерлитамак, Россия)\par 
%$~~\bullet$ Ю.\,Э.~Сеницкий --- д.т.н., проф. (Самара, Россия) \par 
$~~\bullet$ А.\,П.~Солдатов --- д.ф.-м.н., проф. (Белгород, Россия) \par
$~~\bullet$ В.\,В.~Стружанов --- д.ф.-м.н., проф. (Екатеринбург, Россия) \par 
$~~\bullet$ А.\,И.~Хромов --- д.ф.-м.н., проф. (Комсомольск-на-Амуре, Россия)\par
\smallskip \par
}  \\
\end{tabular}
\end{table}

\clearpage
\thispagestyle{empty}


\begin{table}[h!]
\centering
\begin{tabular}{p{6.5cm}p{6.5cm}}
\multicolumn{2}{p{12cm}}{ 
\centering 
\footnotesize 
НАУЧНОЕ ИЗДАНИЕ \par \smallskip
Вестник Самарского государственного технического университета. \par 
Серия <<Физико-математические науки>>~(\RISSUE) \par \bigskip
Учредитель --- ФГБОУ\,ВО <<Самарский государственный технический университет>> 
443100, г.~Самара, ул. Молодогвардейская, 244, Главный корпус\par\bigskip
Редактор Е.~С.~З\,а\,х\,а\,р\,о\,в\,а \par \smallskip
\par Выпускающий редактор Е.~В.~А\,б\,р\,а\,м\,о\,в\,а \par \smallskip
Компьютерная вёрстка М.~Н.~С\,а\,у\,ш\,к\,и\,н, О.~С.~А\,ф\,а\,н\,а\,с\,ь\,е\,в\,а, Е.~Ю.~А\,р\,л\,а\,н\,о\,в\,а} 
\\&\\
%\textit{Сотрудники редакции\/\rm:} \par\small
%Е.\,Ю. Арланова --- секретарь (к.ф.-м.н.)\par 
%О.\,С. Афанасьева --- секретарь  (к.т.н.)\par 
%А.\,А.~Андреев (к.ф.-м.н., доц.) \par
%А.\,Ф.~Заусаев (д.ф.-м.н., проф.) \par 
%В.\,Е.~Зотеев (д.т.н., доц.)  \par
%Э.\,Я.~Рапопорт (д.т.н., проф.) \par \bigskip 
\textit{Адрес редакции и издателя\/\rm:} \par \small
ФГБОУ\,ВО <<СамГТУ>>,\par 
443100, г. Самара,\par 
ул. Молодогвардейская, 244 \par\bigskip
{\it Телефон}\/: +7 (846) 337 04 43 \par 
{\it Факс}\/: +7 (846) 278 44 00 \par  
{\it E-mail}\/: \href{mailto:vsgtu@samgtu.ru}{{\tt vsgtu@samgtu.ru}} \par 
{\it URL}\/: \url{http://www.mathnet.ru/vsgtu}
\par\medskip
\textit{Адрес типографии\/\rm:} \par\small 
443100, г. Самара,\par 
ул. Молодогвардейская, 244.\par 
Корпус 8
&
\textit{Свидетельство о регистрации} \par\smallskip \small
\href{https://rkn.gov.ru/mass-communications/reestr/media/?id=83010}{ПИ \No~ФС 77--66685 от 27.07.2016}. 
\par \bigskip \small
Подписано в печать \JournalDateSign\par 
Дата выхода в свет\, \JournalDatePrint\par
Формат $70 \times  108~{\raise0.5ex\hbox{$\scriptstyle 1\!$} \kern-0.1em/\kern-0.15em \lower0.25ex\hbox{$\scriptstyle\, 16$}}$.\par 
Усл. печ. л. \JournalUPL. 
Уч.-изд. л. \JournalUIL. \par 
Тираж 500 экз. Рег. \No~\JournalRegNumber. \par 
Заказ \No~{\tt \JournalOrderNumber}.\par\medskip
{\fontsize{10pt}{6pt}\selectfont\vspace*{0ex}
Отпечатано в типографии
\par Самарского государственного
\par технического университета\phantom{C}
}
\\&\\
\multicolumn{2}{p{13cm}}{\footnotesize
Журнал индексируется в Web of Science Core Collection (\href{http://mjl.clarivate.com/cgi-bin/jrnlst/jlresults.cgi?PC=MASTER\&ISSN=1991-8615}{Emerging Sources Citation Index}), в \href{http://wokinfo.com/products_tools/multidisciplinary/rsci/}{Russian Science Citation Index} и~входит в ядро
\href{http://elibrary.ru/projects/citation/cit_index.asp}{Российского индекса научного цитирования}.
\par \smallskip
Журнал включен в  \href{http://vak.ed.gov.ru/87}{Перечень рецензируемых научных изданий}, в которых должны быть опубликованы основные научные результаты диссертаций на соискание ученой степени кандидата наук, на соискание ученой степени доктора наук по следующим группам научных специальностей:
01.01.00  – математика; 01.02.00  – механика; 05.13.00  – информатика,  вычислительная техника и  управление.
\par \smallskip
Журнал включен в реферативные базы данных ВИНИТИ (\href{http://www.viniti.ru}{$\tt http{:}{/\!/}www.viniti.ru$}).%,
%Zentralblatt MATH ($\tt http{:}{/\!/}www.zentralblatt{-}math.org{/}portal{/}en{/}$) и
%ULRICH'S Pe\-ri\-o\-di\-cals Directory ($\tt http{:}{/\!/}www.ulrichsweb.com$).
\par \smallskip
Полнотекстовый доступ к статьям журнала осуществляется на Общероссийском математическом портале (\href{http://goo.gl/EeAvfJ}{$\tt http{:}{/\!/}www.mathnet.ru$}), на сайтах научных электронных библиотек eLibrary.ru
(\href{http://elibrary.ru/title_about.asp?id=5784}{$\tt http{:}{/\!/}elibrary.ru$}) и КиберЛенинка (\href{http://cyberleninka.ru/journal/n/vestnik-samarskogo-gosudarstvennogo-tehnicheskogo-universiteta-seriya-fiziko-matematicheskie-nauki}{$\tt http{:}{/\!/}cyberleninka.ru$}).
\par \smallskip
Статьи журнала индексируются в Scholar.Google.com (\href{http://scholar.google.com/scholar?q=Vestn.+Samar.+Gos.+Tekhn.+Univ.+Ser.+Fiz.-Mat.+Nauki}{$\tt http{:}{/\!/}scholar.google.com$}).
\par \bigskip 
\copyright\ Самарский государственный технический университет, \JournalYear\ (составление)\par\smallskip
\OAlogo~\ccby\  Контент публикуется на условиях  лицензии \href{https://creativecommons.org/licenses/by/4.0/deed.ru}{Creative Commons Attribution 4.0 International} (\url{https://creativecommons.org/licenses/by/4.0/deed.ru})
\par \medskip \bf
Подписной индекс в Каталоге Роспечать 18108 \par 
ISSN 1991--8615 (print), 2310--7081 (online)
}\\
\end{tabular}
\par \medskip \small
\begin{tabular}{|c|p{5.5cm}|c}
\cline{1-2}
ФЗ & Издание не подлежит маркировке &  \\
№ 436-ФЗ & в соответствии с п.~1 ч.~2 ст.~1  & \multicolumn{1}{c}{~~~Цена свободная}  \\
\cline{1-2}
\end{tabular}
\end{table}


\clearpage
\thispagestyle{empty}

\begin{table}[t!]
\centering
\begin{tabular}{p{7.4cm}p{5.6cm}}
{\fontsize{22pt}{12pt}
\selectfont\vspace*{0ex}
\bf~Journal \phantom{p} \par\smallskip 
\bf~of Samara State  \phantom{p} \par \smallskip 
\bf~Technical University
}&{
\fontsize{12pt}{8pt}\selectfont\vspace*{0ex}
SCIENTIFIC JOURNAL \medskip\par
Published since 1996 \medskip\par
4 issues per year  \par \medskip
\JournalEMonth{} --- \JournalYear
}
\\
\multicolumn{2}{p{12cm}}{
\fontsize{12pt}{11pt}\selectfont\vspace*{0ex}
\mbox{\bf~~Ser. Physical \& Mathematical Sciences, \EISSUE}}
\smallskip\par\\
\hline
\multicolumn{2}{p{13cm}}{\vspace{-1mm}\centering  \fontsize{9pt}{7pt}\selectfont Founder --- Samara State Technical University, Samara, 443100, Russian Federation} \par\bigskip
\\
\multicolumn{2}{p{13cm}}{
\fontsize{9pt}{8pt}\selectfont\vspace*{0ex}
\textbf{Editor-in-Chief} V.\,P.~Radchenko (Samara, Russian Federation) \par 
\textbf{Deputy Editor-in-Chief} A.\,I.~Zhdanov (Samara, Russian Federation) \par 
\textbf{Executive Secretary} M.\,N.~Saushkin (Samara, Russian Federation)\par
\textbf{Secretary} E.\,A.~Kozlova (Samara, Russian Federation)\par
\bigskip
\textbf{Editorial Council:} \par \smallskip
$~~\bullet$ B.\,D.~Annin (Novosibirsk, Russian Federation) \par
$~~\bullet$ S.\,A.~Avdonin (Fairbanks, Alaska, USA) \par
$~~\bullet$ P.\,B.~Babadzhanov (Dushanbe, Tajikistan) \par
$~~\bullet$ A.\,A.~Burenin (Komsomolsk-on-Amur, Russian Federation) \par
$~~\bullet$  Shõji Imatani  (Kyoto, Japan) \par 
$~~\bullet$ O.\,I.~Marichev  (Laramie, Wyoming, USA) \par
$~~\bullet$ V.\,P.~Matveenko (Perm, Russian Federation) \par
$~~\bullet$ P.\,V.~Sevastiyanov (Cz\c{e}stochowa, Poland)  \par
$~~\bullet$ Z.\,D.~Usmanov  (Dushanbe, Tajikistan) \par
$~~\bullet$ Ramachandran Velmurugan  (Chennai, India) \par
$~~\bullet$ S.\,P.~Verevkin  (Rostock, Germany) \par
$~~\bullet$ N.\,A.~Virchenko (Kiev, Ukraine) \par\bigskip
\textbf{Editorial Board:} \par \smallskip
$~~\bullet$ A.\,P.~Amosov (Samara, Russian Federation) \par 
$~~\bullet$ V.\,N.~Hakobyan (Yerevan, Armenia)\par 
$~~\bullet$ A.\,I.~Khromov (Komsomolsk-on-Amur, Russian Federation)\par
$~~\bullet$ V.\,A.~Kovalev (Moscow, Russian Federation)\par
$~~\bullet$ A.\,I.~Kozhanov (Novosibirsk, Russian Federation) \par 
$~~\bullet$ V.\,A.~Kudinov (Samara, Russian Federation) \par 
%$~~\bullet$ \fbox{A.\,V.~Manzhirov} (Moscow, Russian Federation)\par
$~~\bullet$ A.\,M.~Nakhushev (Nalchik, Russian Federation)\par 
$~~\bullet$ A.\,A.~Pimerzin (Samara, Russian Federation) \par
$~~\bullet$ L.\,S.~Pulkina (Samara, Russian Federation)\par
$~~\bullet$ Yu.\,N.~Radayev  (Moscow, Russian Federation)\par
$~~\bullet$ E.\,V.~Radkevich (Moscow, Russian Federation)\par
%$~~\bullet$ \fbox{O.\,A.~Repin} (Samara, Russian Federation) \par 
$~~\bullet$ K.\,B.~Sabitov  (Sterlitamak, Russian Federation)\par
$~~\bullet$ A.\,V.~Sahakyan (Yerevan, Armenia)\par 
%$~~\bullet$ Yu.\,\'E.~Senitskii (Samara, Russian Federation)\par 
$~~\bullet$ A.\,P.~Soldatov (Belgorod, Russian Federation)\par
$~~\bullet$ V.\,V.~Struzhanov (Ekaterinburg, Russian Federation)\par 
}\\
\end{tabular}
\end{table}


\clearpage
\thispagestyle{empty}

$\;$
\vfill

\begin{minipage}{12cm}
\footnotesize
{\centering \par
Edited by E.~S.~Z\,a\,k\,h\,a\,r\,o\,v\,a, E.~V.~A\,b\,r\,a\,m\,o\,v\,a \par \smallskip
Compiled and typeset by M.~N.~S\,a\,u\,s\,h\,k\,i\,n, O.~S.~A\,f\,a\,n\,a\,s'\,e\,v\,a, E.~Yu.~A\,r\,l\,a\,n\,o\,v\,a \par \bigskip}



 

%{\small \textit{The Editorial Staff:}\/}
%\par \smallskip
%O.\,S.~Afanas'eva (Secretary, Cand. Techn. Sci.; Samara, Russian Federation)\par
%E.\,Yu.~Arlanova    (Secretary, Cand. Phys. \& Math.; Samara, Russian Federation)\par
%A.\,A.~Andreev (Ph.\,D. Phys. \& Math., Assoc. Prof.; Samara, Russian Federation) \par 
%\'E.\,Ya.~Rapoport (Dr. Techn. Sci., Prof.; Samara, Russian Federation) \par 
%A.\,F.~Zausaev (Dr. Phys. \& Math. Sci., Prof.; Samara, Russian Federation) \par 
%V.\,E.~Zoteev (Dr. Techn. Sci.; Samara, Russian Federation) \par
%\bigskip
{\small \textit{The Editorial Board Address}\/:}\par \smallskip
Dept. of Applied Mathematics and Computer Science, \par 
\ensamgtu \par\medskip
\textit{Phone}\/: +7 (846) 337 04 43 \par 
\textit{Phax}\/:~\,\,\,+7 (846) 278 44 00 \par 
\textit{E-mail}\/: \href{mailto:vsgtu@samgtu.ru}{vsgtu@samgtu.ru}\par 
\textit{URL}\/: \mbox{\href{http://www.mathnet.ru/eng/vsgtu}{$\tt http{:}{/\!/}www.mathnet.ru{/}eng{/}vsgtu$}}
\par\medskip
Printed at the Samara State Technical University Press.
\par \bigskip
\fontsize{9pt}{8pt}\selectfont
The journal covered in Web of Science Core Collection (\href{http://mjl.clarivate.com/cgi-bin/jrnlst/jlresults.cgi?PC=MASTER\&ISSN=1991-8615}{Emerging Sources Citation Index}) and  \href{http://wokinfo.com/products_tools/multidisciplinary/rsci/}{Russian Science Citation Index}.
\par \smallskip
The journal is included in VINITI (\href{http://www.viniti.ru}{$\tt http{:}{/\!/}www.viniti.ru$}) abstracts databases.\par \smallskip
The full-text electronic version of journal is hosted by all-Russian mathematical portal Math-Net.Ru (\href{http://www.mathnet.ru}{$\tt http{:}{/\!/}www.mathnet.ru$}) and by the web-site of scientific electronic library eLibrary.ru (\href{http://elibrary.ru/title_about.asp?id=5784}{$\tt http{:}{/\!/}elibrary.ru$}). \par\smallskip
Journal articles indexed in Scholar.Google.com (\href{http://scholar.google.com/scholar?q=Vestn.+Samar.+Gos.+Tekhn.+Univ.+Ser.+Fiz.-Mat.+Nauki}{$\tt http{:}{/\!/}scholar.google.com$}).\par\medskip
\par \bigskip 
\copyright\ Samara State Technical University, \JournalYear\ (Compilation)\par\smallskip
\OAlogo~\ccby\ The content is published under the terms of the \href{http://creativecommons.org/licenses/by/4.0/}{Creative Commons \mbox{Attribution} 4.0 International} License (\url{http://creativecommons.org/licenses/by/4.0/})
\par \medskip \bf
The subscription index in Rospechat catalogue 18108\par \smallskip
ISSN 1991--8615 (print), 2310--7081 (online)  
\end{minipage}

\clearpage
\normalfont

\vsgtu{}
\chead[\fancyplain{}{\selectlanguage{russian}\theruauthor}]{\fancyplain{}{\selectlanguage{russian} \therutitle}}
\rhead[]{\fancyplain{\RuJournalCont}{}}
\lhead[\fancyplain{\RuJournalCont}{}]{}
\cfoot[]{}
\lfoot[\thepage]{}
\rfoot[]{\thepage}
\renewcommand{\plainheadrulewidth}{0.7pt}
\renewcommand{\headrulewidth}{0.7pt}
\pagestyle{fancyplain}
\thispagestyle{plain}
\normalfont
\makecontents
\renewcommand{\baselinestretch}{0.9}

\newpage
\normalfont

\vsgtu{}
\chead[\fancyplain{}{\selectlanguage{english}\theenauthor}]{\fancyplain{}{\selectlanguage{english} \theentitle}}
\rhead[]{\fancyplain{\EnJournalCont}{}}
\lhead[\fancyplain{\EnJournalCont}{}]{}
\cfoot[]{}
\lfoot[\thepage]{}
\rfoot[]{\thepage}
\renewcommand{\plainheadrulewidth}{0.7pt}
\renewcommand{\headrulewidth}{0.7pt}
\pagestyle{fancyplain}
\thispagestyle{plain}
\normalfont
\makeengcontents
\renewcommand{\baselinestretch}{0.9}

\newpage
\normalfont
 







%%%%

\input{preamble}





\usepackage{rotating_pr}

\usepackage{desclist}

\usepackage{cmap}

\usepackage{textcomp}

\usepackage{nccboxes}

\usepackage{euscript}

\usepackage{mathrsfs} 

\usepackage{multicol}

\usepackage{nccfoots}

\usepackage{mathrsfs}

\usepackage{calrsfs}

\allowdisplaybreaks[4]

\usepackage{qrcode}

\usepackage{afterpage}

\usepackage{wasysym}

\usepackage{wrapfig}

\usepackage{slashbox}



\usepackage{multirow}

\newcommand{\specialcell}[2][c]{%

  \begin{tabular}[#1]{@{}c@{}}#2\end{tabular}}





\usepackage{pgf,pgfplots,tikz}

\usetikzlibrary{arrows,snakes,backgrounds,patterns,shapes.geometric,calc}



\nofiles



\oek % для генерации pdf, отдаваемого в oek.rsl.ru

%\samgtupubl % для генерации pdf, отдаваемого в типографию СамГТУ

%\pdfcrop % обрезка pdf для размещения на math-net.ru

%\draft % На корректуре





%\DeclareMathOperator{\tr}{tr}

%\DeclareMathOperator{\ink}{Ink}

%\DeclareMathOperator{\erf}{erf}







%\makeatletter

%\newcommand{\PersonRadRef}{\@langrustrue\def\refname{{\bf\small Избранные  труды \rupid{19071}{О.~А.~Репина} за последние 5 лет}}}

%\makeatother



\DeclareMathOperator{\tr}{tr}

\DeclareMathOperator{\arcsh}{arcsh}









\usepackage{listings}

\usepackage{color}

\definecolor{dkgreen}{rgb}{0,0.6,0}

\definecolor{gray}{rgb}{0.5,0.5,0.5}

\definecolor{mauve}{rgb}{0.58,0,0.82}



\renewcommand*{\lstlistingname}{Листинг}

\usepackage{caption}

%\DeclareCaptionFont{white}{\color{white}} 

\DeclareCaptionFormat{listing}{\parbox{\textwidth}{\centering \small  #1.~#3}}

\captionsetup[lstlisting]{format=listing}



\lstset{ %

  inputencoding=utf8,

  language=R,                     % the language of the code

  basicstyle=\footnotesize,       % the size of the fonts that are used for the code

  numbers=left,                   % where to put the line-numbers

  numberstyle=\tiny\color{gray},  % the style that is used for the line-numbers

  stepnumber=1,                   % the step between two line-numbers. If it's 1, each line

                                  % will be numbered

  numbersep=5pt,                  % how far the line-numbers are from the code

  backgroundcolor=\color{white},  % choose the background color. You must add \usepackage{color}

  showspaces=false,               % show spaces adding particular underscores

  showstringspaces=false,         % underline spaces within strings

  showtabs=false,                 % show tabs within strings adding particular underscores

  frame=single,                   % adds a frame around the code

  rulecolor=\color{black},        % if not set, the frame-color may be changed on line-breaks within not-black text (e.g. commens (green here))

  tabsize=2,                      % sets default tabsize to 2 spaces

  captionpos=t,                   % sets the caption-position to bottom

  breaklines=true,                % sets automatic line breaking

  breakatwhitespace=false,        % sets if automatic breaks should only happen at whitespace

  %title=\lstname,                 % show the filename of files included with \lstinputlisting;

                                  % also try caption instead of title

  keywordstyle=\color{blue},      % keyword style

  commentstyle=\color{dkgreen},   % comment style

  stringstyle=\color{mauve},      % string literal style

  %escapeinside={\%*}{*)},         % if you want to add a comment within your code

  morekeywords={*,in,+,-,=}            % if you want to add more keywords to the set

} 







\begin{document}







\setcounter{page}{601}

\input{specialpages.tex} 







\input{./01du/partition.tex}

\input{./02mdtt/partition.tex}

\input{./03matmod/partition.tex}

\input{./04short/partition.tex}



%\input{./polieces.tex}



\end{document}



\newpage



$\;$





\chead[\fancyplain{}{}]{\fancyplain{}{}}

  \rhead[]{\fancyplain{}{}}

  \lhead[\fancyplain{}{}]{}

  \cfoot[]{}

  \lfoot[]{}

  \rfoot[]{}

\renewcommand{\plainheadrulewidth}{0pt}

\renewcommand{\headrulewidth}{0pt}





  \pagestyle{fancyplain}



  \thispagestyle{plain}





\newpage



$\;$





\chead[\fancyplain{}{}]{\fancyplain{}{}}

  \rhead[]{\fancyplain{}{}}

  \lhead[\fancyplain{}{}]{}

  \cfoot[]{}

  \lfoot[]{}

  \rfoot[]{}

\renewcommand{\plainheadrulewidth}{0pt}

\renewcommand{\headrulewidth}{0pt}





  \pagestyle{fancyplain}



  \thispagestyle{plain}





%\newpage

%

%$\;$

%

%

%\chead[\fancyplain{}{}]{\fancyplain{}{}}

%  \rhead[]{\fancyplain{}{}}

%  \lhead[\fancyplain{}{}]{}

%  \cfoot[]{}

%  \lfoot[]{}

%  \rfoot[]{}

%\renewcommand{\plainheadrulewidth}{0pt}

%\renewcommand{\headrulewidth}{0pt}

%

%

%  \pagestyle{fancyplain}

%

%  \thispagestyle{plain}







\end{document}





\Russian



\addtocontents{toc}{\bigskip\par\hspace{-7mm}}

\addtocontents{etoc}{\bigskip\par\hspace{-7mm}}



$\;$



\hfill \qrcode[hyperlink,height=.8in]{http://doi.org/10.14498/vsgtu1651}

\vspace{-1.5in}





%$$\quad $$ \vskip.6in \par \hskip-7mm  {\huge  Краткие сообщения}

$$\quad $$ \vskip.6in \par \hskip-7mm  {\huge  Short Communications}



%$$\quad $$ \vskip.6in \par \hskip-7mm  {\huge  Mathematical Modeling,\\ Numerical Methods \\[1mm] and Software Complexes}



\addtocontents{toc}{{\large\bf Краткие сообщения}\bigskip}

\addtocontents{etoc}{{\large\bf Short Communications}\bigskip}



\vspace{-5mm}



%\Partition{Математическое моделирование}{Mathematical Modeling}



%\input{preamble}

%\usepackage{samgtu-name}

%

%

%

\renewcommand{\RuCitation}{\noindent\textbf{\CYRO\cyrb\cyrr\cyra\cyrz\cyre\cyrc\ \cyrd\cyrl\cyrya\ \cyrc\cyri\cyrt\cyri\cyrr\cyro\cyrv\cyra\cyrn\cyri\cyrya\\}\selectlanguage{english}{\EnAuthorInContents} {\EnTitleInContents}, \textit{\EnJournalName}\/ [\ParallelJournalName],  \JournalYear, vol.~\JournalVolume, no.~\JournalNumber,  pp.~\thefirstpage--\thelastpage. \doiurl. \selectlanguage{russian}\smallskip}



\renewcommand{\EnCitation}{\noindent\textbf{Please cite this article in press as:\\} {\EnAuthorInContents} {\EnTitleInContents}, \textit{\EnJournalName}\/ [\ParallelJournalName],  \JournalYear, vol.~\JournalVolume, no.~\JournalNumber,  pp.~\thefirstpage--\thelastpage.  \midoiurl. }

%

\vsgtu{1651} %registration number

\FirstPage{735}

\LastPage{749}

%

%

%

%\pdfcrop

%\draft

%

\ShortCommunication

%

%\begin{document}

%

%

%

%$\;$

%\vspace{1mm}

%

%\hfill \qrcode[hyperlink,height=.8in]{http://doi.org/10.14498/vsgtu1651}

%\vspace{-.8in}





\renewcommand{\baselinestretch}{0.98}



\udc{532.51, 517.958:531.3-324}

\msc{76F02, 76F45, 76M45, 76R05, 76U05}



\rutitle{Динамические равновесия неизотермической жидкости}

\entitle{Dynamic equilibria of a nonisothermal fluid}



\ruauthor{Е.~Ю.~Просвиряков\affil{1,2}}{П\,р\,о\,с\,в\,и\,р\,я\,к\,о\,в~Е.~Ю.}

\enauthor{E.~Yu.~Prosviryakov\affil{1,2}}{P\,r\,o\,s\,v\,i\,r\,y\,a\,k\,o\,v~E.~Yu.}





\ruaffil{\ruaffid{2665}{Институт машиноведения УрО РАН,\\

Россия, 620049, Екатеринбург, ул. Комсомольская, 34}.

\and

\ruaffid{7370}{Уральский федеральный университет \\ им.~первого Президента России Б.~Н.~Ельцина, \\ 

Россия, 620002, Екатеринбург, ул. Мира, 19}.}



\enaffil{\enaffid{2665}{Institute of Engineering Science, Urals Branch, Russian Academy of Sciences,\\

34, Komsomolskaya st., Ekaterinburg, 620049, Russian Federation}.

\and

\enaffid{7370}{Ural Federal University named after the First President of Russia B.~N.~Yeltsin, \\ 19, Mira st., Ekaterinburg, 620002, Russian Federation}.

}



\enauthorsdetails{1}{

\fullauthor{\enpid{41426}{Evgeny Yu. Prosviryakov}}

\CAuthor 

\orcid{0000-0002-2349-7801} 

\regalia{Dr. Phys. \& Math. Sci., Professor} 

\position{Head of Sector} 

\dept{Sector of Nonlinear Vortex Hydrodynamics}

\email[br]{evgen_pros@mail.ru}

}



\ruauthorsdetails{1}{

\fullauthor{\rupid{41426}{Евгений Юрьевич Просвиряков}} 

\CAuthor 

\orcid{0000-0002-2349-7801}

\regalia{доктор физико-математических наук} 

\position{заведующий сектором} 

\dept{сектор нелинейной вихревой гидродинамики}

\email{evgen_pros@mail.ru}

}





\enkeywords{dynamic equilibrium,  rotating fluid flow, exact solution, Boussinesq approximation, counterflows, increased velocities}



\rukeywords{динамическое  равновесие, вращающийся поток жидкости, точное решение, аппроксимация Буссинеска, противотечение, увеличенные скорости}



\ruabstract{В рамках точности приближения Буссинеска рассмотрены стационарные динамические равновесия вращающейся массы неизотермической жидкости. Показано, что в этом случае в жидкости наблюдается конечное число противотечений и усиление скоростей по сравнению с заданными на границе значениями, а также формирование зон положительного и отрицательного давления и температуры. }



\enabstract{In this paper, stationary dynamic equilibria of the rotating mass of a nonisothermal fluid are discussed within the accuracy limits of the Boussinesq approximation. It is demonstrated that, in this case, a fluid exhibits a finite number of counterflows, higher values of velocities than those specified on the boundary and the formation of zones of positive and negative pressures and temperatures.}



\newdate{privalovaa}{07}{10}{2018}

\date{\displaydate{privalovaa}}

\newdate{privalovab}{10}{11}{2018}

\revisiondate{\displaydate{privalovab}}

\newdate{privalovac}{12}{11}{2018}

\accepteddate{\displaydate{privalovac}}



\newdate{privalovae}{13}{12}{2018}

\renewcommand{\JournalDataOnline}{\displaydate{privalovae}}





%\ForPeerReview

\makeentitle





\Section[n]{Introduction}

The study of the motion of a rotating self-gravitating fluid was started as long ago as by Isaak Newton~[\citen{ulg:bib:Newtoni},~\citen{ulg:bib:Turnbull}], and it is currently far from being completed. Newton proved to be the first to endeavour to obtain, from the motion equations, a solution simulating the shape of the Earth. Newton’s pioneering study gave rise to numerous investigations in this field of mechanics and physics and a great impetus to the development of the mathematics used practically in all the sections of mathematics. The development of the potential theory, the creation of the theory of special functions and the formation of the rapidly developing theory of integrable systems should be mentioned first.



Practically all the most prominent scientists, such as Appel, Basset, Betti, Gagen, Helmholtz, Dedekind, Dirichlet, Kirchhoff, Lipschitz, Lame, McLoren, Poincare, Riemann, Chandrasekhar, Jacoby, took part in the construction of the theory~[\citen{ulg:bib:Poincare}--\citen{ulg:bib:Dirichlet}]. This list of discoverers of exact solutions can be extended~[\citen{ulg:bib:Poincare}--\citen{ulg:bib:Dirichlet}]. The history of discovering the solutions can be found in the monographs and surveys~[\citen{ulg:bib:Poincare}--\citen{ulg:bib:Dirichlet}]. It should be noted that the scientific breakthrough in the simulation of rotating equilibrium figures was made by the Russian scientists Steklov, Lyapunov, Ovsyannikov and Zel'dovich.



Note that the development of the theory of motion stability described by systems of ordinary differential equations was initiated by studying the stability of equilibrium figures and the motion of an elastic body partially filled with a fluid~[\citen{ulg:bib:Lyapunov1},~\citen{ulg:bib:Lyapunov2}]. It is the work in these fields that offered a formulation to the well-known Lyapunov theorems of stability and instability, among other things, from the first approximation~[\citen{ulg:bib:Lyapunov1},~\citen{ulg:bib:Lyapunov2}]. Original research has currently been done on the stability of Dirichlet and Jacobi ellipsoids~\cite{ulg:bib:Borisov}.



Ovsyannikov’s exact gas-hydrodynamic solution became the first solution    si\-mulating gas cloud motion in the class of linear velocities~\cite{ulg:bib:Borisov}. Afterwards models of Dyson and Fujimoto were proposed~\cite{ulg:bib:Borisov}.



There is another substantial model, which influenced the development of the theory of rotating fluids. It is the generalized solid proposed for research by Arnold~\cite{ulg:bib:Arnold}. In the survey made by Dolzhansky~\cite{ulg:bib:Dolzhansky} there is a detailed analysis of this approach with some generalizations for different classes of fluids.



As a rule, dynamic fluid equilibria are simulated with the application of two mathematical models. One model is Lagrangian fluid simulation, i.e. the use of the ordinary differential equations of Lagrange, Hamilton and Jacoby to describe the relative equilibrium (solid-state) of the continuum. The other model consists in using the Euler equations with potential forces to analyze the structures of a relative fluid equilibrium. Thus, in the simulation of dynamic fluid equilibria no account is taken of dissipation and temperature effect. In other words, an isothermal perfect incompressible fluid is dealt with. 

	

It is apparent that this deficiency in studying fluid flows can be compensated if this problem is analyzed with the application of the Oberbeck--Boussinesq system of equations. Note that it is not only the Boussinesq approximation that has not been used before to describe dynamic fluid equilibria, but also the characteristic scale of solutions. This statement needs to be clarified. Despite Chandrasekhar~\cite{ulg:bib:Chandrasekhar} used the Boussinesq approximation to simulate astrophysical objects, he used the virial method of physical field expansion for solving his problems. The application of this body of mathematics enabled one to reduce the Oberbeck--Boussinesq system to a system of ordinary differential equations and to use the classical problem statements when interpreting results. Similar reasoning holds for~[\citen{ulg:bib:Arnold},~\citen{ulg:bib:Dolzhansky}].

	

There are studies dealing with large-scale fluid or gas motions~\cite{ulg:bib:Chandrasekhar}, when it is necessary to take into account the inhomogeneity of the gravitational field, and with scales for which surface tension forces are principally important, e. g. the Hadamard-Rybczynski problem~[\citen{ulg:bib:Hadamard},~\citen{ulg:bib:Rybczynski}]. In both extreme cases, the principal equilibrium shape is spherical. The consideration of rotation was always viewed as some disturbance of the principal flow, however it was not viewed as a forming factor for equilibrium structures.



Note that the very possibility of the existence of equilibrium figures in the region of intermediate scales, when neither the inhomogeneity of the gravitational field nor the surface tension forces are essential, does not seem to have been studied. At any rate, the authors are unfamiliar with works viewing the problem this way.



Thus, we intend to ascertain the possible existence of equilibrium structures in the region of the so-called intermediate scales, when the Boussinesq approximation is obviously applicable to the simulation of a heat-conducting viscous incompres\-sible fluid.





\smallskip











\Section{Problem statement}

In the Boussinesq approximation~[\citen{ulg:bib:Chandrasekhar},~\citen{ulg:bib:Aristov1}] in the Cartesian coordinates (the $Oz$  axis is directed upwards), the system of equations describing the heat convection of a viscous incompressible fluid are written as

\begin{gather}

\frac{\partial \bf{V}}{\partial t}+\left(\bf{V}\cdot\nabla\right)\bf{V}=-\nabla \textit{P} + \nu \Delta \bf{V}+g \beta \textit{T} \bold{k},

\\    

      \label{eq:Aristov_Prosviryakov:NSE}

\frac{\partial T}{\partial t}+\bf{V}\cdot\nabla \textit{T}=\chi \Delta \textit{T},

\\     

\nabla \cdot \bf{V}=0.

 \end{gather}

The system of equations~\eqref{eq:Aristov_Prosviryakov:NSE} involves the following symbols: $\bf{V}=$ $\left(V_x , V_y , V_z\right)$  is the flow velocity vector;  $P$  is the pressure deviation from hydrostatic, which is related to the constant average fluid density  $\rho;$  $T$ is the deviation from the average temperature;   $\nu$  and $\chi$ are the coefficients of kinematic viscosity and thermal diffusivity of the fluid, respectively; $\beta$ is the temperature coefficient of fluid volumetric expansion; $\bold{k}$  is the unit vector directed vertically upwards, $\nabla$  is the Hamilton operator,  $\Delta$ is the Laplace operator (Laplacian).



The solution of system~\eqref{eq:Aristov_Prosviryakov:NSE} is sought in the approximation to the large scale. For large-scale flows, the characteristic horizontal scale greatly exceeds in magnitude the characteristic dimension during the entire fluid motion. Free surface curvature can in this scale be neglected, the vertical velocity $V_z$  being assumed zero. Thus, system~\eqref{eq:Aristov_Prosviryakov:NSE} is transformed to the following system:

\begin{gather}

\frac{\partial V_x}{\partial t}+V_x \frac{\partial V_x}{\partial x}+V_y \frac{\partial V_x}{\partial y}=-\frac{\partial P}{\partial x} + \nu \Bigl(\frac{\partial^2 V_x}{\partial x^2}+\frac{\partial^2 V_x}{\partial y^2}+\frac{\partial^2 V_x}{\partial z^2}\Bigr),

\\       

\frac{\partial V_y}{\partial t}+V_x \frac{\partial V_y}{\partial x}+V_y \frac{\partial V_y}{\partial y}=-\frac{\partial P}{\partial y} + \nu \Bigl(\frac{\partial^2 V_y}{\partial x^2}+\frac{\partial^2 V_y}{\partial y^2}+\frac{\partial^2 V_y}{\partial z^2}\Bigr),

\\

       \label{eq:Aristov_Prosviryakov:system}

\frac{\partial P}{\partial z}=g\beta T,

\\

\frac{\partial T}{\partial t}+V_x \frac{\partial T}{\partial x}+V_y \frac{\partial T}{\partial y}=\chi \Bigl(\frac{\partial^2 T}{\partial x^2}+\frac{\partial^2 T}{\partial y^2}+\frac{\partial^2 V_x}{\partial z^2}\Bigr),

\\

\frac{\partial V_x}{\partial x}+\frac{\partial V_y}{\partial y}=0.

 \end{gather}

 

Obviously, system \eqref{eq:Aristov_Prosviryakov:system} is nonlinear and overdetermined, four unknown functions of velocity and temperature being required from five equations. The hydrodynamic and temperature fields are computed as~\cite{ulg:bib:Aristov1}

\begin{gather}

V_x=U+xu_1+y u_2,

\quad 

V_y=V+xv_1+y v_2,

\\

       \label{eq:Aristov_Prosviryakov:solution}

P=P_0+xP_1+yP_2+\frac{x^2}{2}P_{11}+\frac{y^2}{2}P_{22}+xyP_{12},

\\

T=T_0+xT_1+yT_2+\frac{x^2}{2}T_{11}+\frac{y^2}{2}T_{22}+xyT_{12}.     

 \end{gather}

 

For horizontal coordinates, all the functions in the formulae of system \eqref{eq:Aristov_Prosviryakov:solution}  depend on the variable $z$  ant time $t.$ Substituting relation \eqref{eq:Aristov_Prosviryakov:solution} into system \eqref{eq:Aristov_Prosviryakov:system}, we arrive at the system

\begin{gather}

\hat{L}U+Uu_1+Vu_2+P_1=0,

\\

\hat{L}u_1+u^{2}_1+v_1u_2+P_{11}=0,

\\

\hat{L}u_2+u_1u_2+u_2v_2+P_{12}=0,

\\

\hat{L}V+Uv_1+Vv_2+P_{2}=0,

\\

\hat{L}v_1+u_1v_1+v_1v_2+P_{12}=0,

\\

\hat{L}v_2+u_2v_1+v^{2}_2+P_{22}=0,

\\

       \label{eq:Aristov_Prosviryakov:main system }

\hat{M}T_0+UT_1+VT_2-\chi\left(T_{11}+T_{22}\right)=0,

\\

\hat{M}T_1+UT_{11}+VT_{12}+u_1T_1+v_1T_2=0,

\\

\hat{M}T_2+UT_{12}+VT_{22}+u_2T_1+v_2T_2=0,

\\

\hat{M}T_{11}+2u_1T_{11}+2v_1T_{12}=0,

\\

\hat{M}T_{22}+2u_2T_{12}+2v_2T_{22}=0,

\\

\hat{M}T_{12}+u_1T_{12}+u_2T_{11}+v_1T_{22}+v_2T_{12}=0,

\\

u_1+v_2=0,

\\

\frac{\partial P_0}{\partial z}=g\beta T_0,

\quad 

\frac{\partial P_1}{\partial z}=g\beta T_1,

\quad 

\frac{\partial P_2}{\partial z}=g\beta T_2,

\\

\frac{\partial P_{11}}{\partial z}=g\beta T_{11},

\quad 

\frac{\partial P_{12}}{\partial z}=g\beta T_{12},

\quad 

\frac{\partial P_{22}}{\partial z}=g\beta T_{22}.

 \end{gather}

Here, the following differential parabolic operators in partial derivatives are intro\-duced:

$$\hat{L}=\frac{\partial }{\partial t}-\nu\frac{\partial^2 }{\partial z^2},\quad 

\hat{M}=\frac{\partial }{\partial t}-\chi\frac{\partial^2 }{\partial z^2}. 

$$





\smallskip





\Section{Analyzing the solvability of system}

In what follows we do not discuss the solvability of system \eqref{eq:Aristov_Prosviryakov:main system } as a whole, but restrict our analysis to one special case. We will seek a solution to the system of equations when the boundaries of the fluid layer rotate. In~\cite{ulg:bib:Aristov2} the validity of equalities for the velocity gradients in the simulation of axisymmetric flows within class \eqref{eq:Aristov_Prosviryakov:solution} was demonstrated,

\begin{gather}

       \label{eq:Aristov_Prosviryakov:sol1}

u_1=0,

\quad 

v_2=0,       

\quad 

u_2=-v_1.       

 \end{gather}

Next the lower boundary of the layer  $z=0$ is assumed to be immobile. The upper boundary describable by the equation $z=h$  rotates with constant angular velocity  $\Omega$. Using equalities \eqref{eq:Aristov_Prosviryakov:sol1}, we obtain the following solutions to system~\eqref{eq:Aristov_Prosviryakov:main system }, which determine a number of functions in class \eqref{eq:Aristov_Prosviryakov:solution} for the coordinates  $x$ and  $y$:

\begin{gather}

u_2=\frac{\Omega z}{h}, 

\quad 

v_1=-\frac{\Omega z}{h},

\quad 

P_{12}=0, 

\\

\label{eq:Aristov_Prosviryakov:sol2}

 P_{11}=P_{22}=-u_2v_1=\frac{\Omega^2 z^2}{h^2},

\quad 

T_{12}=0, 

\\ g\beta T_{11}=g\beta T_{22}=\frac{2\Omega^2 z}{h^2}.

\end{gather}



For the homogeneous summands in expressions \eqref{eq:Aristov_Prosviryakov:solution}, due to solution \eqref{eq:Aristov_Prosviryakov:sol2}, system \eqref{eq:Aristov_Prosviryakov:main system } yields two closed subsystems. First system has of the form

\begin{gather}

\hat{L}U+\frac{\Omega z}{h}V+P_1=0, 

\\

\hat{L}V-\frac{\Omega z}{h}U+P_{2}=0,

\\

\label{eq:Aristov_Prosviryakov:sys1}

\hat{M}T_1+\frac{2\Omega^2 z}{g\beta h^2}U-\frac{\Omega z}{h}T_2=0, 

\\ \hat{M}T_2+\frac{2\Omega^2 z}{g\beta h^2}V+\frac{\Omega z}{h}T_1=0, 

\\

\frac{\partial P_1}{\partial z}=g\beta T_1, 

\quad  

\frac{\partial P_2}{\partial z}=g\beta T_2.

\end{gather}



Second system has of the form

\begin{gather}

\label{eq:Aristov_Prosviryakov:sys2}

\hat{M}T_0+UT_1+VT_2=\frac{4\chi\Omega^2 z}{g\beta h^2},

\\

\frac{\partial P_0}{\partial z}=g\beta T_0.

\end{gather}



Note that the integration of system \eqref{eq:Aristov_Prosviryakov:sys2} is possible only after the integration of system \eqref{eq:Aristov_Prosviryakov:sys1}. In order to find a solution to system \eqref{eq:Aristov_Prosviryakov:sys1}, we introduce the following complex functions:

$$

W=U+iV,\quad 

S=P_1+iP_2,\quad 

\Theta=g\beta\left(T_1+iT_2\right).$$

Here  $i$ is an imaginary unit. Thus, system \eqref{eq:Aristov_Prosviryakov:sys1} is in this case transformed to the following form:

 \begin{gather}

(\hat{L}-i\Omega z) W+S=0,

\\

       \label{eq:Aristov_Prosviryakov:sys1complex}

(\hat{M}+i\Omega z) \Theta+2\Omega^2 z W=0,       

\\

\frac{d S}{dz}=\Theta.      

 \end{gather}

 

To reduce the order of the system, the following transformation is made:

 \begin{gather}

       \label{eq:Aristov_Prosviryakov:trans}

S=S_1+i\Omega z W,

\quad 

\Theta=\Theta_1+2i\Omega W.    

\end{gather}

Substituting \eqref{eq:Aristov_Prosviryakov:trans} into system \eqref{eq:Aristov_Prosviryakov:sys1complex}, we obtain the equalities

 \begin{gather}

\nu\frac{d^2 W}{dz^2}+S_1=0,

\\

       \label{eq:Aristov_Prosviryakov:trans2}

\chi\frac{d^2 \Theta_1}{d z^2}+i\Omega \Bigl(z \Theta_1- 2\frac{\chi}{\nu}S_1\Bigr)=0,       

\\

\frac{d S_1}{dz}=\Theta_1-i\Omega z^2 \frac{d}{dz}\Bigl(\frac{W}{z}\Bigr). 

 \end{gather}

 

For further transformations, we use the differential identity

$$

\frac{d^2 W}{dz^2}=\frac{1}{z}\frac{d}{dz}\left(z^2\frac{d}{dz}\Bigl(\frac{W}{z}\Bigr)\right).

$$

This identity holds for any twice-differentiable function. Consequently, by differentiating the last equation in system \eqref{eq:Aristov_Prosviryakov:trans2}, we arrive at a system for determining the functions  $\Theta_1$  and  $S_1,$

 \begin{gather}

       \label{eq:Aristov_Prosviryakov:trans3}

\chi\frac{d^2 \Theta_1}{d z^2}+i\Omega \left(z \Theta_1- 2\frac{\chi}{\nu}S_1\right)=0,

\quad 

\nu\frac{d^2 S_1}{dz^2}=i\Omega z S_1 +\nu \frac{d \Theta_1}{dz}.          

 \end{gather}

 

Taking into account that equations \eqref{eq:Aristov_Prosviryakov:trans2} offer an explicit expression for  $W$, we obtain an integral of the initial system with automatic order reduction, system \eqref{eq:Aristov_Prosviryakov:trans3} being transformed to the operator equation

$$

\left[\Bigl(\nu\frac{d^2}{dz^2}-i\Omega z\Bigr)\Bigl(\chi\frac{d^2}{dz^2}+i\Omega z\Bigr)-2i\Omega\chi\frac{d}{dz}\right]\Theta_1=0.

$$

Simplifying the latter equality, we derive the following fourth-order ordinary differential equation with complex coefficients:

 \begin{gather}

       \label{eq:Aristov_Prosviryakov:trans4}

\Bigl[\nu\chi\frac{d^4}{dz^4}+i\Omega\left(\nu-\chi\right)\frac{d^2}{dz^2}z+\Omega^2 z^2\Bigr]\Theta_1=0.    

 \end{gather}



\smallskip







\Section{Solving equation when the dissipative coefficients coincide}

 We assume in equation \eqref{eq:Aristov_Prosviryakov:trans4} that $\nu=\chi$  (the Prandtl number equals one). In this case, equation \eqref{eq:Aristov_Prosviryakov:trans4} acquires the form

  \begin{gather}

       \label{eq:Aristov_Prosviryakov:Pr1}

\frac{d^4 \Theta_1}{dz^4}+\left(36\sigma^3 z\right)^2 \Theta_1=0,    

 \end{gather}

where  $\sigma=\frac{1}{6}\sqrt[3]{\frac{\Omega}{6\nu h}}$ is the dispersion relation. For the convenience of computations, in order to illustrate the physical meaning of the solutions and to make the notation concise, we introduce a new variable  $Z=\frac{z}{h \sigma}.$ The general solution of equation \eqref{eq:Aristov_Prosviryakov:Pr1} is written as

\begin{multline}

\label{eq:Aristov_Prosviryakov:soltemp}

\Theta_1={^{3}_{0}F} \Bigl[\Bigl\{\frac{1}{2}, \frac{2}{3}, \frac{5}{6}\Bigr\},-Z^6\Bigr]C_1 +

{^{3}_{0}F} 

\Bigl[\Bigl\{\frac{2}{3}, \frac{5}{6}, \frac{7}{6}\Bigr\},-Z^6\Bigr] C_2 Z +

\\

+

{^{3}_{0}F} \Bigl[\Bigl\{\frac{5}{6}, \frac{7}{6}, \frac{4}{3}\Bigr\},-Z^6\Bigr] C_3 Z^2  +

{^{3}_{0}F} \Bigl[\Bigl\{\frac{7}{6}, \frac{4}{3},\frac{3}{2}\Bigr\},-Z^6\Bigr] C_4 Z^3  . 

\end{multline}

Here,  ${^{3}_{0}F}=F\left(a; b; Z\right)=F\left(a_1, a_2, \ldots, a_p; b_1, b_2, \ldots, b_q; Z\right)$ is a generalized hypergeometric function~\cite{ulg:bib:Polyanin}. The choice of using the generalized hypergeometric function for writing the solutions is dictated by notation conciseness. Note that there is another form of writing the solution of equation \eqref{eq:Aristov_Prosviryakov:Pr1}, which is based on using the Mittag--Leffler functions~\cite{ulg:bib:Polyanin}. Note that each particular solution \eqref{eq:Aristov_Prosviryakov:soltemp} is a function of a real variable and that the general solution is considered in the complex plane. Thus, at least one integration constant is a complex number.



We now turn to finding the generalized pressure $S_1$  from the differential equation \eqref{eq:Aristov_Prosviryakov:trans2} 

 \begin{gather}

       \label{eq:Aristov_Prosviryakov:eqpr}

S_1=-\frac{i}{2\sigma}\frac{d^2 \Theta_1}{d z^2}+\frac{z}{2}\Theta_1.

\end{gather}



The substitution of solution \eqref{eq:Aristov_Prosviryakov:soltemp} into equation \eqref{eq:Aristov_Prosviryakov:eqpr} gives

\begin{multline}

\label{eq:Aristov_Prosviryakov:solstress}

S_1= {^{3}_{0}F} 

\Bigl[\Bigl\{\frac{1}{2}, \frac{2}{3}, \frac{5}{6}\Bigr\},-Z^6\Bigr] \frac{C_1}{2}Z+

{^{3}_{0}F} 

\Bigl[\Bigl\{\frac{2}{3}, \frac{5}{6}, \frac{7}{6}\Bigr\},-Z^6\Bigr] \frac{C_2}{2}Z^2+

\\

+{^{3}_{0}F} 

\Bigl[\Bigl\{\frac{5}{6}, \frac{7}{6}, \frac{4}{3},\Bigr\},Z\Bigr] \sqrt[3]{36\sigma^2} \frac{2\cdot 5C_3}{6!}Z^3+

{^{3}_{0}F} 

\Bigl[\Bigl\{\frac{7}{6}, \frac{4}{3}, \frac{3}{2}\Bigr\},Z\Bigr] \sigma \frac{2\cdot 5C_4}{6!}Z^4-

\\

-\frac{1}{72\sigma}i\left\{ {^{3}_{0}F} 

\Bigl[\Bigl\{\frac{5}{6}, \frac{7}{6}, \frac{4}{3}\Bigr\},Z\Bigr] \sqrt[3]{676\sigma^2} C_3 Z^3 

-{^{3}_{0}F} 

\Bigl[\Bigl\{\frac{7}{6}, \frac{4}{3}, \frac{3}{2}\Bigr\},Z\Bigr] 6 \sigma C_4 Z+

\right.

\\

+{^{3}_{0}F} 

\Bigl[\Bigl\{\frac{5}{6}, \frac{7}{6}, \frac{4}{3}\Bigr\},Z\Bigr] 3 \sigma^2 C_1 Z^4 +

\\

+{^{3}_{0}F} 

\Bigl[\Bigl\{\frac{5}{3}, \frac{11}{6}, \frac{13}{6},\Bigr\},Z\Bigr] \frac{2^4\cdot 3^3\sigma^2}{7!}\sqrt[3]{\frac{3\sigma}{4}} C_2 Z^5

+

\\

+{^{3}_{0}F} 

\Bigl[\Bigl\{\frac{5}{3}, \frac{11}{6}, \frac{13}{6},\Bigr\},Z\Bigr] \frac{2^4\cdot 3^4\sigma^2}{7!}\sqrt[3]{6\sigma} C_2 Z^5+

\\

+{^{3}_{0}F} 

\Bigl[\Bigl\{\frac{11}{6}, \frac{13}{6}, \frac{7}{3} \Bigr\},Z\Bigr] \frac{2^2\cdot 3^4\sigma^2}{7!}\sqrt[3]{\frac{9\sigma^2}{2}} C_3 Z^6+

\\

+{^{3}_{0}F} \Bigl[\Bigl\{\frac{13}{6}, \frac{7}{3}, \frac{5}{2} \Bigr\},Z\Bigr] \frac{2\cdot 5\cdot 11\sigma^3}{7!} C_4 Z^7-

\\

-{^{3}_{0}F} \Bigl[\Bigl\{\frac{5}{2}, \frac{8}{3}, \frac{17}{6} \Bigr\},Z\Bigr] \frac{2^7\cdot 3^3\cdot 7\sigma^4}{11!} C_1 Z^{10}-

\\

-{^{3}_{0}F} \Bigl[\Bigl\{\frac{8}{3}, \frac{17}{6}, \frac{19}{6} \Bigr\},Z\Bigr] \frac{2^9\cdot 3^6\sigma^4}{13!}\sqrt[3]{\frac{3\sigma}{4}} C_2 Z^{11}-

\\

-{^{3}_{0}F} \Bigl[\Bigl\{\frac{17}{6}, \frac{19}{6}, \frac{10}{3} \Bigr\},Z\Bigr] \frac{2^8\cdot 3^5\cdot 5\sigma^4}{14!}\sqrt[3]{\frac{9\sigma^2}{2}} C_3 Z^{12}-

\\

\left.

-{^{3}_{0}F} \Bigl[\Bigl\{\frac{19}{6}, \frac{10}{3}, \frac{7}{2} \Bigr\},Z\Bigr] \frac{2^7\cdot 3^2\cdot 5\cdot 11\sigma^5}{14!} C_4 Z^{13}

\right\}.

\end{multline}





Complex velocity $W$  can be expressed in two ways. Firstly, it can be obtained by integrating the first equation in system \eqref{eq:Aristov_Prosviryakov:sys1complex}, but it is a second-order equation. Hence, we obtain a solution depending on six integration constants, whereas system \eqref{eq:Aristov_Prosviryakov:trans2} is a fifth-order equation system. Thus our problem is complicated by additional analysis of overdetermination. A correct expression for velocity $W$  can be obtained by the integration of the third equation in system \eqref{eq:Aristov_Prosviryakov:sys1complex},

\begin{multline}

\label{eq:Aristov_Prosviryakov:solvelo}

\Omega W=C_5 Z+ \left({^{3}_{0}F} \Bigl[\Bigl\{\frac{2}{3}, \frac{5}{6}, \frac{7}{6} \Bigr\},Z\Bigr]-1\right) \frac{i}{2} \sqrt[3]{\frac{\sigma}{36}} C_2 Z-

\\

-{^{4}_{1}F} \Bigl[\Bigl\{-\frac{1}{6} \Bigr\}, \Bigl\{\frac{1}{2},\frac{2}{3}, \frac{5}{6}, \frac{5}{6} \Bigr\},Z\Bigr]\frac{2i \Gamma\left(-\frac{1}{6}\right)}{4!\Gamma\left(\frac{5}{6}\right)} C_1 Z^2+

\\

+{^{4}_{1}F} \Bigl[\Bigl\{-\frac{1}{6} \Bigr\},\Bigl\{\frac{5}{6},\frac{7}{6}, \frac{4}{3}, \frac{3}{2} \Bigr\},Z\Bigr]\frac{2\cdot 5\Gamma\left(-\frac{1}{6}\right)}{6!\Gamma\left(\frac{5}{6}\right)} C_4+

\\

+{^{4}_{1}F} \Bigl[\Bigl\{\frac{1}{6} \Bigr\},\Bigl\{\frac{5}{6},\frac{7}{6}, \frac{7}{6}, \frac{4}{3} \Bigr\},Z\Bigr]\frac{2\cdot 5i \Gamma\left(\frac{1}{6}\right)}{6!\Gamma\left(\frac{7}{6}\right)} \sqrt[3]{\frac{\sigma^2}{6}} C_3 Z^2+

\\

+{^{4}_{1}F} \Bigl[\Bigl\{\frac{1}{3} \Bigr\},\Bigl\{\frac{7}{6},\frac{4}{3}, \frac{4}{3}, \frac{3}{2} \Bigr\},Z\Bigr]\frac{2^4\cdot 3\cdot 5\cdot 7i \sigma \Gamma\left(\frac{1}{3}\right)}{9!\Gamma\left(\frac{4}{3}\right)} C_4 Z^3-

\\

-{^{4}_{1}F} \Bigl[\Bigl\{\frac{1}{3} \Bigr\},\Bigl\{\frac{4}{3},\frac{3}{2}, \frac{5}{3}, \frac{11}{6} \Bigr\},Z\Bigr]\frac{2^2\cdot 5\sigma \Gamma\left(\frac{1}{3}\right)}{6!\Gamma\left(\frac{4}{3}\right)} C_1 Z^3-

\\

-{^{4}_{1}F} \Bigl[\Bigl\{\frac{1}{2} \Bigr\},\Bigl\{\frac{3}{2},\frac{5}{3}, \frac{11}{6}, \frac{13}{6} \Bigr\},Z\Bigr]\frac{\sigma }{4!}\sqrt[3]{\frac{\sigma}{36}} C_2 Z^4-

\\

-{^{4}_{1}F} \Bigl[\Bigl\{\frac{2}{3} \Bigr\},\Bigl\{\frac{5}{3},\frac{11}{6}, \frac{13}{6}, \frac{7}{3} \Bigr\},Z\Bigr] \frac{2^4\cdot 3^2\cdot 7\sigma \Gamma\left(\frac{2}{3}\right)}{9!\Gamma\left(\frac{5}{3}\right)}\sqrt[3]{\frac{\sigma}{6}} C_3 Z^5-

\\

-{^{4}_{1}F} \Bigl[\Bigl\{\frac{5}{6} \Bigr\},\Bigl\{\frac{3}{2},\frac{5}{3}, \frac{11}{6}, \frac{11}{6} \Bigr\},Z\Bigr] \frac{2^3\cdot 3^2\cdot 7i\sigma^2 \Gamma\left(\frac{5}{6}\right)}{9!\Gamma\left(\frac{11}{6}\right)} C_1 Z^6-

\\

-{^{4}_{1}F} \Bigl[\Bigl\{\frac{5}{6} \Bigr\},\Bigl\{\frac{11}{6},\frac{13}{6}, \frac{7}{3}, \frac{5}{2} \Bigr\},Z\Bigr] \frac{2^3\cdot 5^2\cdot 11\cdot 13\cdot 83\sigma^2 \Gamma\left(\frac{5}{6}\right)}{13! \Gamma\left(\frac{11}{6}\right)} C_4 Z^6-

\\ 

-{^{4}_{1}F} \Bigl[\Bigl\{\frac{7}{6} \Bigr\},\Bigl\{\frac{11}{6},\frac{13}{6}, \frac{13}{6}, \frac{7}{3} \Bigr\},Z\Bigr] \frac{2\cdot 3^2i\sigma^2 \Gamma\left(\frac{7}{6}\right)}{9! \Gamma\left(\frac{13}{6}\right)}\sqrt[3]{\frac{\sigma^2}{6}} C_3 Z^8-

\\

-{^{4}_{1}F} \Bigl[\Bigl\{\frac{4}{3} \Bigr\},\Bigl\{\frac{13}{6},\frac{7}{3}, \frac{7}{3}, \frac{5}{2} \Bigr\},Z\Bigr] \frac{2^3\cdot 5^2\cdot 11\cdot 13 i\sigma^3 \Gamma\left(\frac{4}{3}\right)}{13!\Gamma\left(\frac{7}{3}\right)} C_4 Z^9+

\\

+{^{4}_{1}F} \Bigl[\Bigl\{\frac{4}{3} \Bigr\},\Bigl\{\frac{7}{3},\frac{5}{2}, \frac{8}{3}, \frac{17}{6} \Bigr\},Z\Bigr] \frac{2^3\cdot 3\cdot 5\cdot 7\sigma^3 \Gamma\left(\frac{4}{3}\right)}{11! \Gamma\left(\frac{7}{3}\right)} C_1 Z^9+

\\

+{^{4}_{1}F} \Bigl[\Bigl\{\frac{3}{2} \Bigr\},\Bigl\{\frac{5}{2},\frac{8}{3}, \frac{17}{6}, \frac{19}{6} \Bigr\},Z\Bigr] \frac{2^7\cdot 3^5\sigma^3}{13!} \sqrt[3]{\frac{\sigma}{36}} C_2 Z^{10}+

\\

+{^{4}_{1}F} \Bigl[\Bigl\{\frac{5}{3} \Bigr\},\Bigl\{\frac{8}{3},\frac{17}{6}, \frac{19}{6}, \frac{10}{3} \Bigr\},Z\Bigr] \frac{2^3\cdot 3^4\cdot 5\sigma^3 \Gamma\left(\frac{5}{3}\right)} {13! \Gamma\left(\frac{8}{3}\right)} \sqrt[3]{\frac{\sigma^2}{6}} C_3 Z^{11}+

\\

+{^{4}_{1}F} \Bigl[\Bigl\{\frac{11}{6} \Bigr\},\Bigl\{\frac{17}{6},\frac{19}{6}, \frac{10}{3}, \frac{7}{2} \Bigr\},Z\Bigr] \frac{2^6\cdot 5\cdot 11\sigma^4 \Gamma\left(\frac{11}{6}\right)} {14! \Gamma\left(\frac{17}{6}\right)} C_4 Z^{12}-

\\

-{^{4}_{1}F} \Bigl[\Bigl\{\frac{7}{3} \Bigr\},\Bigl\{\frac{10}{3},\frac{7}{2}, \frac{11}{3}, \frac{23}{6} \Bigr\},Z\Bigr] \frac{2^6\cdot 3\cdot 7^2\cdot 13\sigma^5 \Gamma\left(\frac{7}{3}\right)} {17! \Gamma\left(\frac{10}{3}\right)} C_1 Z^{15}-

\\

-{^{4}_{1}F} \Bigl[\Bigl\{\frac{5}{2} \Bigr\},\Bigl\{\frac{7}{2},\frac{11}{3}, \frac{23}{6}, \frac{25}{6} \Bigr\},Z\Bigr] \frac{2^8\cdot 3^7\cdot 7\sigma^5 } {19!} \sqrt[3]{\frac{\sigma}{36}} C_2 Z^{16}-

\\

-{^{4}_{1}F} \Bigl[\Bigl\{\frac{8}{3} \Bigr\},\Bigl\{\frac{11}{3},\frac{23}{6}, \frac{25}{6}, \frac{13}{3} \Bigr\},Z\Bigr] \frac{2^7\cdot 3^6\cdot 5\sigma^5 \Gamma\left(\frac{8}{3}\right)C_3 Z^{17}} {19! \Gamma\left(\frac{11}{3}\right)} \sqrt[3]{\frac{\sigma^2}{6}} -

\\

-{^{4}_{1}F} \Bigl[\Bigl\{\frac{8}{3} \Bigr\},\Bigl\{\frac{23}{6},\frac{25}{6}, \frac{13}{3}, \frac{9}{2} \Bigr\},Z\Bigr] \frac{2^8\cdot 3^2\cdot 5^2\cdot 11\cdot 17\sigma^6 \Gamma\left(\frac{17}{6}\right)} {21! \Gamma\left(\frac{23}{6}\right)} \sqrt[3]{\frac{\sigma^2}{6}} C_4 Z^{18}.

\end{multline}

Here $\Gamma\left(\lambda\right)= \displaystyle\int_{0}^{+\infty}r^{\lambda-1}\exp\left(-r\right)dr$ is the Euler gamma function~\cite{ulg:bib:Polyanin}. Solutions~\eqref{eq:Aristov_Prosviryakov:soltemp}, \eqref{eq:Aristov_Prosviryakov:solstress},~\eqref{eq:Aristov_Prosviryakov:solvelo} are written as a linear  combination of generalized hypergeometric functions in the form of infinite power series. Consequently, examination of solution convergence is an urgent problem. Using the properties of convergence of these functions, we can state that the convergence radius for each solution equals infinity~\cite{ulg:bib:Polyanin}. It means that the solution of our problem always converges.



Let us now formulate the boundary conditions to determine the integration constants of equation system \eqref{eq:Aristov_Prosviryakov:sys1}. Without any claim to solution generality and without studying the shapes of the fluid, we give a simple example of such solutions. Let the simplest and most physically meaningful solution be a flow between two infinite disks rotating with free velocities. This is possible, since the solution can always be shifted along the $Oz$-axis, an appropriate distance between the disks being chosen. The disks must be nonuniformly heated to a proper temperature (the temperature must be related to the radius by a well-defined law). Assume that the lower boundary is isothermal (the temperature of the lower disk is zero); the adhesion condition holds on both boundaries. Thus, the following equalities are valid:

$$\text{when} \ \ z=0, \ \ U=V=0, \ \ T_1=T_2=0;$$

$$\text{when} \ \ z=h, \ \ U=A\cos\alpha, \ \ V=A\sin\alpha, \ \ T_1=B, \ \ T_2=C, \ \ P_1=\tau_1, \ \ P_2=\tau_2.$$



The conditions for determining the integration constants, written on the boundary, do not determine the integration constants uniquely. The system of linear algebraic equations for determining the integration constants is underdetermined, since there are few boundary conditions. To close the system of equations, we specify the fluid flow rate conditions as

$$\int _{0}^{h}Udz =Q_1,\quad  \int_{0}^{h}Vdz =Q_2.$$



When the flow rate is specified, it is generally assumed that  $Q_1=Q_2=0$~\cite{ulg:bib:Aristov2}. It is a physically justified convention~[\citen{ulg:bib:Aristov1},~\citen{ulg:bib:Aristov2}]. However, it is possible to obtain a solution by specifying a nonzero flow rate. Next, not limiting the generality of the reasoning, we will study the solutions simulating fluid motions with the zero flow rate. 



For the above-introduced complex velocities, pressure and temperature ~\eqref{eq:Aristov_Prosviryakov:soltemp}, \eqref{eq:Aristov_Prosviryakov:solstress},~\eqref{eq:Aristov_Prosviryakov:solvelo}, in view of transformations \eqref{eq:Aristov_Prosviryakov:trans}, the following boundary conditions are valid:

\begin{gather}

\label{eq:Aristov_Prosviryakov:bond}

\text{when} \ \ z=0: \ \ W=0, \ \ \Theta_1=0;

\\

\text{when} \ \ z=h: \ \ W=A\exp\left(i\alpha\right), \ \ \Theta_1=B+2A\Omega\sin\alpha+i\left(C-2A\Omega\cos\alpha\right),

\\

\int_{0}^{h}Wdz =0.

\end{gather}



With the boundary conditions \eqref{eq:Aristov_Prosviryakov:bond}, it follows from equality \eqref{eq:Aristov_Prosviryakov:soltemp} that  $C_1=0$, and from solution \eqref{eq:Aristov_Prosviryakov:solvelo} that  $C_4=0$. Thus, the expressions for the determination of hydrodynamic fields become much simpler. Remember that, in the computation of the integration constants, the variable $z=h\sigma Z$  is to be replaced.











\Section{Analysis of the solutions}

 Using conditions \eqref{eq:Aristov_Prosviryakov:bond}, it is easy, but rather cumbersome, to compute the complex  $C_2,$ $C_3,$ and $C_5$  which are related to the dispersion relation   $\sigma$ (kinematic viscosity) by the fractional linear law. Analyzing conditions \eqref{eq:Aristov_Prosviryakov:bond}, we find that the boundary value problem under study is meaningful at any positive value of  $\sigma.$ Nevertheless, the structure and behavior of the solutions depend considerably on the value of  $\sigma.$



It is well known that, to analyze the polynomial solutions of system \eqref{eq:Aristov_Prosviryakov:sys1}, it would suffice to localize the roots of the linear combination of the generalized hypergeometric functions~\cite{ulg:bib:Polyanin}. Any hypergeometric function is known to be monotonically increasing with a zero value at the origin. Therefore, the appearance of other zero values is possible only when there is a combination of functions with at least one negative weight coefficient~\cite{ulg:bib:Polyanin}. The appearance of negative coefficients depends on $\sigma$  and the function domain, which also depends on the dispersion relation $Z\in\left[0, \frac{1}{\sigma}\right].$  Temperature is characterized by three types of solutions. 

When  $\sigma\geq\frac{3}{2},$ the function describing temperature is a strictly monotonically increasing or decreasing function (fig.~\ref{fig_prosv1}). 

If the evaluation  $\frac{9}{10}<\sigma<\frac{3}{2}$  is valid for the dispersion relation, there is exactly one zero value of temperature inside the fluid layer (fig.~\ref{fig_prosv2}). With the further decrease, temperature gets a finite number of zeros (fig.~\ref{fig_prosv3}). Consequently, the fluid layer is divided into the regions with positive and negative temperature, which are characterized by alternating signs. The analysis of pressure is similar to that of temperature. 





Velocities are characterized by only two qualitative behaviors of the functions (figs.~\ref{fig_prosv2} and ~\ref{fig_prosv3}). Thus, at any  $\sigma,$ there appear fluid counterflows. It was shown in~[\citen{ulg:bib:Aristov3}--\citen{ulg:bib:Gorshkov}] that the appearance of counterflows is a purely nonlinear effect, i.\,e., the viscous forces are comparable to the convective acceleration, and even without Coriolis forces. However, when  $\sigma$ tends to zero, counterflow zones are formed, whose thickness decreases with  $\sigma.$ Other values of the dispersion relation are characteristic of velocity, and they are easy to compute. However, the values obtained for the temperature and pressure fields are a good reference point for refining the characteristic dispersion numbers of velocities.





When analyzing the solutions (fig.~\ref{fig_prosv3}), we can see their ``fluctuation''. In~\cite{ulg:bib:Aristov7} localized convective flows in terms were studied in the axisymmetric formulation. Solution ``damping'' was observed there. We use a wider class of solutions. There may be not only ``damping'' (when $\sigma$  is close to one), but also ``resonance'' initiation. Yet, the solution does not become singular, and the Boussinesq approxi\-mation is not violated in the simulation of the convective flow. This difference is attributed to the account of the quadratic summands for temperature. If tempera\-ture is considered to be a linear function with respect to the horizontal coordinates, the solutions obtained in this study coincide with those discussed in~\cite{ulg:bib:Aristov4}.



\begin{figure}[h!]



\begin{tabular}{p{.45\textwidth}p{.45\textwidth}}

\centerline{\includegraphics[width=0.45\textwidth]{fig_prosv1}} &

\centerline{\includegraphics[width=0.45\textwidth]{fig_prosv2}} \\

\caption{A qualitative form of the solutions  $T_1,$ $T_2,$ $P_1,$ $P_2$ of system (\ref{eq:Aristov_Prosviryakov:sys1}) when $\sigma\geq\frac{3}{2}$ ($1-\tau_1>0$,  $\tau_2>0$, $A>0$, $B>0$, $C>0$; $2-\tau_1<0$, $\tau_2<0$, $A<0$, $B<0$, $C<0$) \label{fig_prosv1}} &

\caption{A qualitative form of the solutions  $T_1,$ $T_2,$ $P_1,$ $P_2,$ $U,$ $V$ of system (\ref{eq:Aristov_Prosviryakov:sys1}) when $\frac{9}{10}<\sigma<\frac{3}{2}$ ($1-\tau_1>0$,  $\tau_2>0$, $A>0$, $B>0$, $C>0$; $2-\tau_1<0$, $\tau_2<0$, $A<0$, $B<0$, $C<0$) \label{fig_prosv2} }

\\

\centerline{\includegraphics[width=0.45\textwidth]{fig_prosv3}} &



\vspace{-2cm}

\caption{A qualitative form of the solutions  $T_1,$ $T_2,$ $P_1,$ $P_2,$ $U,$ $V$ of system (\ref{eq:Aristov_Prosviryakov:sys1}) when $\sigma\leq\frac{9}{10}$ ($1-\tau_1>0$,  $\tau_2>0$, $B>0$, $C>0$; $2-\tau_1<0$, $\tau_2<0$, $B<0$, $C<0$) \label{fig_prosv3} } \\





\end{tabular}



\vspace{-3cm}



%%\noindent

%%\begin{minipage}[h]{0.48\textwidth}

%%

%%\smallskip \footnotesize

%%Fig.~\ref{fig1}.

%%A qualitative form of the solutions  $T_1,$ $T_2,$ $P_1,$ $P_2$ of system (\ref{eq:Aristov_Prosviryakov:sys1}) when $\sigma\geq\frac{3}{2}$ ($1-\tau_1>0$,  $\tau_2>0$, $A>0$, $B>0$, $C>0$; $2-\tau_1<0$, $\tau_2<0$, $A<0$, $B<0$, $C<0$) \label{fig1} 

%%\end{minipage} \hfill

%%\begin{minipage}[h]{0.5\textwidth}

%%\centering

%%\includegraphics[width=0.98\textwidth]{fig1}

%%\end{minipage}

%\end{figure}

%

%

%\begin{figure}[h!]

%\noindent

%\begin{minipage}[h]{0.48\textwidth}

%\caption{A qualitative form of the solutions  $T_1,$ $T_2,$ $P_1,$ $P_2,$ $U,$ $V$ of system (\ref{eq:Aristov_Prosviryakov:sys1}) when $\frac{9}{10}<\sigma<frac{3}{2}$ ($1-\tau_1>0$,  $\tau_2>0$, $A>0$, $B>0$, $C>0$; $2-\tau_1<0$, $\tau_2<0$, $A<0$, $B<0$, $C<0$) 

%\label{fig2}

% }

%\smallskip \footnotesize

%Fig.~\ref{fig2}.

%A qualitative form of the solutions  $T_1,$ $T_2,$ $P_1,$ $P_2,$ $U,$ $V$ of system (\ref{eq:Aristov_Prosviryakov:sys1}) when $\frac{9}{10}<\sigma<\frac{3}{2}$ ($1-\tau_1>0$,  $\tau_2>0$, $A>0$, $B>0$, $C>0$; \centerline{$2-\tau_1<0$, $\tau_2<0$, $A<0$, $B<0$, $C<0$)} \label{fig2} 

%\end{minipage} \hfill

%\begin{minipage}[h]{0.5\textwidth}

%\centering

%\includegraphics[width=0.98\textwidth]{fig2}

%\end{minipage}

%\end{figure}

%

%

%

%\begin{figure}[h!]

%\noindent

%\begin{minipage}[h]{0.48\textwidth}

%\caption{A qualitative form of the solutions  $T_1,$ $T_2,$ $P_1,$ $P_2,$ $U,$ $V$ of system (\ref{eq:Aristov_Prosviryakov:sys1}) when $\sigma\leq\frac{9}{10}$ ($1-\tau_1>0$,  $\tau_2>0$, $B>0$, $C>0$; $2-\tau_1<0$, $\tau_2<0$, $B<0$, $C<0$) 

%\label{fig3}

% }

%\smallskip \footnotesize

%Fig.~\ref{fig3}.

%A qualitative form of the solutions  $T_1,$ $T_2,$ $P_1,$ $P_2,$ $U,$ $V$ of system (\ref{eq:Aristov_Prosviryakov:sys1}) when $\sigma\leq\frac{9}{10}$ ($1-\tau_1>0$,  $\tau_2>0$, $B>0$, $C>0$; 

%\centerline{$2-\tau_1<0$, $\tau_2<0$, $B<0$, $C<0$)}

% \label{fig3} 

%\end{minipage} \hfill

%\begin{minipage}[h]{0.5\textwidth}

%\centering

%\includegraphics[width=0.98\textwidth]{fig3}

%\end{minipage}

\end{figure}









\newpage















\Section{Solution of equation  with the Prandtl number different from one}

An important particular case of equation \eqref{eq:Aristov_Prosviryakov:trans4} was studied above; namely, when the values of kinematic viscosity and thermal diffusivity of the fluid coincide. 

Let us now write the solution of equation \eqref{eq:Aristov_Prosviryakov:trans4} in the general case:

\begin{multline}

\label{eq:Aristov_Prosviryakov:soltempPr}

\Theta_1=\text{\sf Ai}\Bigl(\xi+\frac{4\Omega z}{\sqrt{\nu\chi}}\sqrt[3]{\frac{\nu^2\chi^2} {16\Omega^2}}\Bigr)

\left(C_1\text{\sf Ai}\left(\xi\right)+C_2\text{\sf Bi}\left(\xi\right)\right)+

\\

+\text{\sf Bi}\Bigl(\xi+\frac{4\Omega z}{\sqrt{\nu\chi}}\sqrt[3]{\frac{\nu^2\chi^2}{16\Omega^2}}\Bigr)\left(C_3\text{\sf Ai}\left(\xi\right)+C_4\text{\sf Bi}\left(\xi\right)\right).

\end{multline}

Here $$\xi=\sqrt[3]{\nu^2\chi^2}\displaystyle\frac{ i\left(\nu-\chi\right)+2\sqrt{\nu\chi}\Omega z}{\nu\chi\sqrt[3]{16\Omega^2}}$$ is an auxiliary variable introduced in order to make the solution notation concise;  $\text{\sf Ai} \left(\xi\right)$   and $\text{\sf Bi}\left(\xi\right)$  are the Airy function and the associated function~\cite{ulg:bib:Polyanin}. Solution \eqref{eq:Aristov_Prosviryakov:soltempPr} is complex-valued, since it is written for arbitrary differences between $\nu$  and $\chi.$  This is done for the convenience of the solution analysis. Solution \eqref{eq:Aristov_Prosviryakov:soltempPr} being known, pressure  $S_1$ and velocity $W$ are computed similarly to the previous case.	



When the dissipative coefficients  $\nu$ and  $\chi$ of the nonisothermal fluid are equal, solution \eqref{eq:Aristov_Prosviryakov:soltempPr}  is reduced to formula \eqref{eq:Aristov_Prosviryakov:soltemp}. Studying the convergence of series \eqref{eq:Aristov_Prosviryakov:soltempPr}, we find that the solution expressed in terms of this series converges on the solution domain, i.\,e. layer thickness. 



The consideration of equation \eqref{eq:Aristov_Prosviryakov:trans4} with an arbitrary relation between $\nu$  and~$\chi$ offers solutions to system \eqref{eq:Aristov_Prosviryakov:sys1}, \eqref{eq:Aristov_Prosviryakov:sys2} in qualitative terms, see figs.~\ref{fig_prosv1} to~\ref{fig_prosv3}. In the simulation of the hydrodynamic fields by the linear combination of the Airy functions the behavior of the solutions is predictable from the very beginning, since these special functions satisfy the simplest differential equation having a point where the fluctuation of the solution is replaced by its exponential growth. It was shown above that the solution does not go off to infinity.



\smallskip



\Section{Conclusion}

The problem of simulating rotating fluid masses has a long history. To describe fluid motion (relative equilibrium), ordinary differential equa\-tions are used that enable one to obtain solutions simulating fluid flows by the theory of motion stability with the disturbance of the background (principal) flow. This paper discusses the exact solution to an overdetermined boundary value problem, which describes stationary flow at any scales where the Boussinesq approximation is valid. In this case, stationary dynamic fluid equilibria are chara\-cterized by the formation of counterflows, the sign of velocity being able to alternate several times, depending on the boundary conditions and the geometric dimensions of the fluid layer. A similar situation is true for temperature and pressure, namely, the formation of zones with the negative and positive values of temperature and pressure with respect to the reference value. 



\smallskip 

 



\AdditionalContent{Competing interests}{I declare that I have no competing interests.} 

\AdditionalContent{Author's Responsibilities}{I take full responsibility for submitting the final manuscript in print. I approved the final version of the manuscript.} 





\AdditionalContent{Funding}{The work was done within the state assignment from FASO Russia, theme No. AAAA-A18-118020790140-5.}









\begin{thebibliography}{28}



\Bibitem{ulg:bib:Newtoni}

\by Newton~I.

\book Opera quae exstant omnia

\bookinfo Faksimile-Neudruck der Ausgabe von Samuel Horsley, London 1779--1785 in f\"unf B\"anden. Band 1,~4

\lang In German

\zmath{0129.25102}

\publaddr Stuttgart-Bad Cannstatt

\publ Friedrich Frommann Verlag (G\"unther Holzboog)

\totalpages xx+592

\yr 1964



\Bibitem{ulg:bib:Turnbull}

\book The Correspondence of Isaac Newton

\bookvol 1. 1661-1675

\ed H.~W.~Turnbull

\publ Cambridge Univ. Press

\publaddr New York

\yr 1959

\totalpages xxxviii+468 





\Bibitem{ulg:bib:Poincare}

\by Poincar\'e~H.

\book Cin\'ematique et m\'ecanismes. Potentiel et m\'ecanique des fluides

\bookinfo Cours profess\'e \`a la Sorbonne

\lang In French

\zmath{30.0608.01}

\publaddr Paris

\publ G.~Carr\'e et C.~Naud

\totalpages iv+385 

\yr 1899



\Bibitem{ulg:bib:Lyapunov1}

\by Lyapunov~A.~M.

\paper On the stability of ellipsoidal equilibrium forms of rotating fluid

\inbook Collected works. \rm  Vol. III

\yr 1959

\publ Akad. Nauk SSSR

\publaddr Moscow

\pages 5--113

\zmath{0192.31203}

\lang In Russian



\Bibitem{ulg:bib:Lyapunov2}

\by Lyapunov~A.~M.

\paper On the equilibrium figures of rotating homogeneous liquid mass slightly different from ellipsoids

\inbook Collected works. \rm  Vol. IV

\yr 1959

\publ Akad. Nauk SSSR

\publaddr Moscow

\pages 5--645

\zmath{0192.31204}

\lang In Russian



\Bibitem{ulg:bib:Chandrasekhar}

\by Chandrasekhar~S.

\book Ellipsoidal figures of equilibrium

\zmath{0213.52304}

\publaddr New Haven, London

\publ Yale University Press 

\totalpages ix+252 

\yr 1969



\Bibitem{ulg:bib:Hagihara}

\by Hagihara~Y.

\book Theories of equilibrium figures of a rotating homogeneous fluid mass

\yr 1970

\serial NASA Special Publication 

\vol 186 

\publaddr Washington

\publ US Government Printing Office

\adsnasa{1970NASSP.186.....H}



\Bibitem{ulg:bib:Borisov}

\by Borisov~A.~V., Kilin~A.~A., Mamaev~I.~S.

\paper The Hamiltonian Dynamics of Self-gravitating Liquid and Gas Ellipsoid

\jour Regul. Chaotic Dyn. 

\yr 2009

\vol 14

\issue 2

\pages 179--217

\zmath{https://zbmath.org/?q=an:1229.70049}

\crossref{10.1134/S1560354709020014}



\Bibitem{ulg:bib:Appell}

\by Appell~P.

\book Trait\'e de m\'ecanique rationnelle

\bookinfo Tome IV, 1: Figures d'\'equilibre d'une masse liquide homog\`ene en rotation

\lang In French

\zmath{58.1300.02}

\totalpages viii+342 

\publaddr  Paris

\publ Gauthier-Villars 

\yr 1932







\Bibitem{ulg:bib:Betti}

\by Betti~E.

\paper Sopra i moti che conservano la figura ellissoidale a una massa fluida eterogenea

\lang In Italian

\zmath{13.0718.01}

\jour Brioschi Ann

\issueinfo  (2)~X

\pages 173--187

\yr 1879







\Bibitem{ulg:bib:Volterra}

\by Volterra~V.

\paper Sur la stratification d\'une masse fluide en \'equilibre

\jour Acta Math.

\yr 1903

\vol 27

\issue 1

\pages 105--124



\Bibitem{ulg:bib:Liouville}

\by Liouville~J.

\paper Sur la figure d\'une masse fluide homog\'ene, en \'equilibre et dou\'ee d\'un mouvement de rotation

\jour J. de l'\'Ecole Polytech.

\yr 1834

\vol 14

\pages 289–296



\Bibitem{ulg:bib:Dirichlet}

\by Dirichlet~G.~L.

\paper Untersuchungen \"uber ein Problem der Hydrodynamik (Aus dessen Nachlass hergestellt von Herrn R. Dedekind zu Z\"urich)

\jour J. Reine Angew. Math. (Crelle's Journal)

\yr 1861

\vol 58

\pages181--216



\Bibitem{ulg:bib:Arnold}

\by Arnold~V. I., Khesin~B.~A.

\book Topological Methods in Hydrodynamics

\yr 1999

\publ Springer

\publaddr New York

\totalpages 392



\Bibitem{ulg:bib:Dolzhansky}

\by Dolzhansky~F.~V.

\paper On the mechanical prototypes of fundamental hydrodynamic invariants and slow manifolds

\jour Physics–Uspekhi

\yr 2005

\vol 48

\issue 12

\pages 1205--1234

\crossref{10.1070/PU2005v048n12ABEH002375}



\Bibitem{ulg:bib:Hadamard}

\by Hadamard~J.

\paper Mouvement permanent lent d'une sph\`ere liquide et visqueuse dans un liquide visqueux

\zmath{42.0813.01}

\jour C.~R.~Acad. Sci., Paris 

\vol 152

\issue 25

\pages 1735--1738 

\yr 1911





\Bibitem{ulg:bib:Rybczynski}

\by Rybczi\'nski, W.

\paper \"Uber die fortschreitende Bewegung einer fl\"ussigen Kugel in einem z\"ahen Medium

\zmath{42.0811.02}

\jour Bull. Int. Acad. Sci. Cracovie. Ser. A

\vol 1

\pages 40--46 

\yr 1911



\Bibitem{ulg:bib:Aristov1}

\by Aristov~S.~N., Prosviryakov~E.~Yu. 

\paper A New Class of Exact Solutions for Three Dimensional Thermal Diffusion Equations

\jour Theor. Found. Chem. Eng.

\yr 2016

\vol 50

\issue 3

\pages 286--293

\crossref{10.1134/S0040579516030027}



\Bibitem{ulg:bib:Aristov2}

\by Aristov~S.~N., Knyazev~D.~V., Polyanin~A.~D.

\paper Exact solutions of the Navier-Stokes equations with the linear dependence of velocity components on two space variables

\jour Theor. Found. Chem. Eng.

\crossref{10.1134/S0040579509050066}

\yr 2009

\vol 43

\issue 5

\pages 642--662



\Bibitem{ulg:bib:Polyanin}

\by Polyanin~A.~D., Zaitsev~V.~F.

\book Handbook of ordinary differential equations. Exact solutions, methods, and problems

\zmath{06776009}

\publaddr Boca Raton, FL

\publ CRC Press 

\totalpages xxix+1456 

\yr 2018



\Bibitem{ulg:bib:Aristov3}

\by Aristov~S.~N., Prosviryakov~E.~Yu.

\paper Unsteady Layered Vortical Fluid Flows

\jour Fluid Dyn.

\crossref{10.1134/S0015462816020034}

\yr 2016

\vol 51

\issue 2

\pages 148--154



\Bibitem{ulg:bib:Aristov4}

\by Aristov~S.~N., Prosviryakov~E.~Yu., Spevak~L.~F.

\paper Unsteady-State B\'enard--Marangoni Convection  in Layered Viscous Incompressible Flows

\jour Theor. Found. Chem. Eng.

\crossref{10.1134/S0040579516020019}

\yr 2016

\vol 50

\issue 2

\pages 132--141



\Bibitem{ulg:bib:Aristov5}

\by Aristov~S.~N., Prosviryakov~E.~Yu.

\paper On one class of analytic solutions of the stationary axisymmetric convection B\'enard--Marangoni viscous incompressible fluid

\jour Vestn. Samar. Gos. Tekhn. Univ., Ser. Fiz.-Mat. Nauki \rm [J. Samara State Tech. Univ., Ser. Phys. Math. Sci.]

\yr 2013

\vol 3(32)

\pages 110--118

\mathnet{http://mi.mathnet.ru/vsgtu1205}

\crossref{10.14498/vsgtu1205}

\zmath{https://zbmath.org/?q=an:06968790}

\elib{http://elibrary.ru/item.asp?id=20710187}

\lang In Russian



\Bibitem{ulg:bib:Aristov6}

\by Aristov~S.~N., Privalova~V.~V., Prosviryakov~E.~Yu.

\paper Stationary nonisothermal Couette flow. Quadratic heating of the upper boundary of the fluid layer

\jour Nelin. Dinam.

\yr 2016

\vol 12

\issue 2

\pages 167--178

\mathnet{http://mi.mathnet.ru/nd519}

\elib{http://elibrary.ru/item.asp?id=26193555}

\crossref{10.20537/nd1602001}

\lang In Russian



\Bibitem{ulg:bib:Burmasheva1}

\by Burmasheva~N.~V., Prosviryakov~E.~Yu.

\paper A large-scale layered stationary convection of an incompressible viscous fluid under the action of shear stresses at the upper boundary. Velocity field investigation

\jour Vestn. Samar. Gos. Tekhn. Univ., Ser. Fiz.-Mat. Nauki \rm [J. Samara State Tech. Univ., Ser. Phys. Math. Sci.]

\yr 2017

\vol 21

\issue 1

\pages 180--196

\mathnet{http://mi.mathnet.ru/vsgtu1527}

\crossref{10.14498/vsgtu1527}

\elib{http://elibrary.ru/item.asp?id=29245104}

\lang In Russian



\Bibitem{ulg:bib:Burmasheva2}

\by Burmasheva~N.~V., Prosviryakov~E.~Yu.

\paper A large-scale layered stationary convection of a incompressible viscous fluid under the action of shear stresses at the upper boundary. Temperature and presure field investigation

\jour Vestn. Samar. Gos. Tekhn. Univ., Ser. Fiz.-Mat. Nauki \rm [J. Samara State Tech. Univ., Ser. Phys. Math. Sci.]

\yr 2017

\vol 21

\issue 4

\pages 736--751

\mathnet{http://mi.mathnet.ru/vsgtu1568}

\crossref{10.14498/vsgtu1568}

\zmath{https://zbmath.org/?q=an:06964885}

\elib{http://elibrary.ru/item.asp?id=32712835}

\lang In Russian



\Bibitem{ulg:bib:Gorshkov}

\by Gorshkov~A.~V., Prosviryakov~E.~Yu.

\paper Ekman Convective Layer Flow of a Viscous Incompressible Fluid

\jour Izv. Atmos. Ocean. Phys.

\yr 2018

\vol 54

\issue 2

\pages 89–195

\crossref{10.1134/S0001433818020081}





\Bibitem{ulg:bib:Aristov7}

\by Aristov~S.~N., Knyazev~D.~V.

\paper Localized convective flows in a nonuniformly heated liquid layer

\jour Fluid Dyn.

\yr 2014

\vol 49

\issue 5

\pages 565--575

\crossref{10.1134/S0015462814050020}











\end{thebibliography}





\renewcommand{\baselinestretch}{0.95}



\newpage



\selectlanguage{russian}



$\;$







\makerutitle



\vspace{-7mm}



\AdditionalContent{Конкурирующие интересы}{Конкурирующих интересов не имею.}



\AdditionalContent{Авторский вклад и ответственность}{Я несу полную ответственность за предоставление окончательной версии рукописи в печать. Окончательная версия 

рукописи мною одобрена.}

\AdditionalContent{Финансирование}{Работа выполнена в рамках государственного задания ФАНО, тема № АААА-А18-118020790140-5.}





\renewcommand{\baselinestretch}{0.9}







\renewcommand{\RuCitation}{\noindent\textbf{\CYRO\cyrb\cyrr\cyra\cyrz\cyre\cyrc\ \cyrd\cyrl\cyrya\ \cyrc\cyri\cyrt\cyri\cyrr\cyro\cyrv\cyra\cyrn\cyri\cyrya\\} {\RuAuthorInContents} {\RuTitleInContents}~/\!/ \textit{\RuJournalName}, \JournalYear.  \CYRT.~\JournalVolume, \textnumero~\JournalNumber.  \CYRS.~\thefirstpage--\thelastpage. \doiurl.}



\renewcommand{\EnCitation}{\noindent\textbf{Please cite this article in press as:\\} {\EnAuthorInContents} {\EnTitleInContents}, \textit{\EnJournalName}\/ [\ParallelJournalName],  \JournalYear, vol.~\JournalVolume, no.~\JournalNumber,  pp.~\thefirstpage--\thelastpage.  \midoiurl\ (In Russian).}









%\end{document}

 \newpage

%\input{preamble}

%

%\usepackage{samgtu-name}

%

%\usepackage{slashbox}



\vsgtu{1653} %registration number

\FirstPage{750}

\LastPage{761}

%

%\pdfcrop

%\draft



\ShortCommunication



%\begin{document}





\hfill \qrcode[hyperlink,height=.8in]{http://doi.org/10.14498/vsgtu1653}

\vspace{-.8in}





\udc{517.958:539.3(1)}

\msc{74F05, 74K20}



\rutitle{Динамическая устойчивость геометрически нерегулярной нагретой пологой цилиндрической оболочки в~сверхзвуковом потоке газа}

\rutitleshort{Динамическая устойчивость геометрически нерегулярной\dots}



\entitle{Dynamic stability of heated geometrically irregular cylindrical shell in supersonic gas flow}

\entitleshort{Dynamic stability of heated geometrically irregular cylindrical shell in supersonic gas flow}





\ruauthor{Г.~Н.~Белосточный, О.~А.~Мыльцина}{Б\,е\,л\,о\,с\,т\,о\,ч\,н\,ы\,й~Г.~Н.,\ М\,ы\,л\,ь\,ц\,и\,н\,а~О.~А.}

\enauthor{G.~N.~Belostochnyi, O.~A.~Myltcina}{B\,e\,l\,o\,s\,t\,o\,c\,h\,n\,y\,i~G.~N.,\ M\,y\,l\,t\,c\,i\,n\,a~O.~A.}



\ruaffil{\saratovuni.}



\enaffil{\engsaratovuni.}



\ruauthorsdetails{3}{% Количество авторов



\fullauthor{\rupid{98285}{Григорий Николаевич Белосточный}}

%\CAuthor % Автор-корреспондент

\orcid{0000-0003-4471-6599} % ORCID ID автора 

\regalia{доктор технических наук, профессор} 

\position{профессор} 

\dept{каф. математической теории упругости и биомеханики} 

\email{belostochny@mail.ru}

\fullauthor{\rupid{98283}{Ольга Анатольевна Мыльцина}}

\CAuthor % Автор-корреспондент

\orcid{0000-0003-4718-2772} % ORCID ID автора 

\regalia{кандидат физико-математических наук} 

\position{ассистент} 

\dept{каф. теории функций и стохастического анализа} 

\email{omyltcina@yandex.ru}

}



%Olga Myltcina https://orcid.org/0000-0003-4718-2772 

%Grigory Belostochnyihttps://orcid.org/0000-0001-6949-4174

%Acel Polienko https://orcid.org/0000-0003-4471-6599







\enauthorsdetails{3}{% Количество авторов

\fullauthor{\rupid{98285}{Grigory N. Belostochnyi}}

%\CAuthor % Автор-корреспондент

\orcid{0000-0003-4471-6599} % ORCID ID автора 

\regalia{Dr. Techn. Sci.} 

\position{Professor} 

\dept{Dept. of Mathematic Theory of Elasticity \&  Biomechanics} 

\email[n]{belostochny@mail.ru} 

\fullauthor{\rupid{98283}{Olga A. Myltcina}}

\CAuthor % Автор-корреспондент

\orcid{0000-0003-4718-2772} % ORCID ID автора 

\regalia{Cand. Phys. \& Math. Sci.} 

\position{Assistant} 

\dept{Dept. of Functions \& Approxmation Theory} 

\email[n]{omyltcina@yandex.ru}

}





\rukeywords{динамическая устойчивость, температура, пологие оболочки, сверхзвук, континуальность, обобщенные функции, поршневая теория, аэродинамика, критические скорости, ребра жесткости, демпфирование, кривизна, критерий Гурвица, изотропия}

\enkeywords{dynamic stability, temperature, flat shells, supersonic, continuity, generalized functions, piston theory, aerodynamics, critical velocities, ribs, damping, curvature, Routh–Hurwitz stability criterion, isotropy}



\ruabstract{На базе модели типа Лява рассматривается нагретая до постоянной температуры геометрически нерегулярная пологая цилиндрическая оболочка, обдуваемая сверхзвуковым потоком газа со стороны одной из ее основных поверхностей. За основу взята континуальная модель термоупругой системы в~виде тонкостенной оболочки, подкрепленной ребрами вдоль набегающего газового потока. Сингулярная система уравнений динамической термоустойчивости геометрически нерегулярной оболочки содержит слагаемые, учитывающие «растяжение-сжатие» и~сдвиг подкрепляющих элементов в~тангенциальной плоскости, тангенциальные усилия, вызванные нагревом оболочки, и~поперечную нагрузку, стандартным образом записанную по «поршневой теории».



Тангенциальные усилия предварительно определяются как решение сингулярных дифференциальных уравнений безмоментной термоупругости геометрически нерегулярной оболочки с~учетом краевых усилий.



Решение системы динамических уравнений термоупругости оболочки разыскивается в~виде суммы двойного тригонометрического ряда (для функции прогиба) с~переменными по временной координате коэффициентами. На основании метода Галеркина получена однородная система для коэффициентов аппроксимирующего ряда, которая сведена к~одному дифференциальному уравнению четвертого порядка. Решение приводится во втором приближении, что соответствует двум полуволнам в~направлении потока и~одной полуволне в~перпендикулярном направлении. На основании стандартных методов анализа динамической устойчивости тонкостенных конструкций определяются критические значения скоростей газового потока.



Количественные результаты приводятся в виде таблиц, иллюстрирующих влияние геометрических параметров термоупругой системы «обо\-лоч\-ка-ребра», температуры  на устойчивость геометрически нерегулярной цилиндрической оболочки в~сверхзвуковом потоке газа с~учетом демпфирования.}



\enabstract{On the basis of the Love model, a geometrically irregular heated cylindrical shell blown by a supersonic gas flow from one of its main surfaces is considered. The continuum model of a thermoelastic system in the form of a thin-walled shell supported by ribs along the incoming gas flow is taken as a basis. The singular system of equations for the dynamic thermal stability of a geometrically irregular shell contains terms that take into account the tension-compression and the shift of the reinforcing elements in the tangential plane, the tangential forces caused by the heating of the shell and the transverse load, as standard recorded by the piston theory. The solution of a singular system of differential equations in displacements, in the second approximation for the deflection function, is sought in the form of a double trigonometric series with time coordinate variables. 





Tangential forces are predefined as the solution of singular differential equations of non-moment thermoelasticity of a geometrically irregular shell taking into account boundary forces. 



The solution of the system of dynamic equations of thermoelasticity of the shell is sought in the form of the sum of the double trigonometric series (for the deflection function) with time coordinate variable coefficients. On the basis of the Galerkin method, a homogeneous system for the coefficients of the approximating series is obtained, which is reduced to one fourth-order differential equation. The solution is given in the second approximation, which corresponds to two half-waves in the direction of flow and one half-wave in the perpendicular direction. On the basis of standard methods of analysis of dynamic stability of thin-walled structures are determined critical values of the gas flow rate. 



The quantitative results are presented in the form of tables illustrating the influence of the geometrical parameters of the thermoelastic shell-edge system, temperature and damping on the stability of a geometrically irregular cylindrical shell in a supersonic gas flow.}



\newdate{myla}{11}{10}{2018}

\date{\displaydate{myla}}

\newdate{mylb}{10}{11}{2018}

\revisiondate{\displaydate{mylb}}

\newdate{mylc}{12}{11}{2018}

\accepteddate{\displaydate{mylc}}



\newdate{myle}{21}{12}{2018}

\renewcommand{\JournalDataOnline}{\displaydate{myle}}



\makerutitle







Геометрически нерегулярные тонкостенные упругие системы, обширный класс которых составляют ребристые оболочки и~пластинки, широко используются в~различных областях современной техники. Условия эксплуатации таких систем предусматривают совместное воздействие температурных полей и~высокоскоростных газовых потоков. Исследованию упругого поведения гладких пластин и~оболочек на основе атермической

теории посвящено большое число работ, полный перечень которых содержал бы десятки

наименований. Ограничимся некоторыми из них \cite{myl-bib:1, myl-bib:2, myl-bib:3,myl-bib:4,myl-bib:5,myl-bib:5a}.

Значительно меньше работ

посвящено  исследованиям совместного воздействия температурных факторов и сверхзвукового потока

на тонкостенные системы. Важные для практики результаты в этой области содержатся в работах \cite{myl-bib:6, myl-bib:7, myl-bib:8}.



%Работы, в которых анализируется влияние подкрепляющих  оболочку ребер под действием сверхзвукового потока на базе термической теории в открытой научной литературе отсутствуют. 

Работы, в которых анализируется влияние подкрепляющих  ребер на поведение оболочки, находящейся под действием сверхзвукового потока, на базе термической теории, в~открытой научной литературе отсутствуют. 

Это связано не с маловажностью проблемы, а прежде всего с чрезвычайной математической сложностью таких задач, решаемых на основе дискретной модели «оболочка–ребра». Использование континуальной модели типа «конструктивная анизотропия» мало пригодно для практических целей, так как количественные результаты, полученные на ее основе, далеки от физической реальности.



Обращение к  континуальным моделям на основе теории обобщенных функций, основные положения которых содержится в работах \cite{myl-bib:9, myl-bib:10, myl-bib:11, myl-bib:12, myl-bib:15,myl-bib:16,myl-bib:18}, позволило сводить решения задач статической и динамической термоупругости ребристых  оболочек к интегрированию систем сингулярных дифференциальных уравнений точными и~приближенными методами высшего анализа \cite{myl-bib:19, myl-bib:21}, что немаловажно для инженерной практики.



\smallskip



{\bf 1.} Система сингулярных дифференциальных уравнений динамической термоупругости пологой геометрически нерегулярной цилиндрической оболочки, стандартным образом отнесенной к декартовым координатам на основе континуальной модели, в которых учет тангенциальных усилий записывается в форме Рейснера \cite{myl-bib:22a}, в компонентах поля перемещений примет вид

\begin{gather}

u_{\bigcom 11}+\frac{1-\nu}{2}u_{\bigcom 22}+\frac{1+\nu}{2} v_{\bigcom 12}-k_{11}w_{\bigcom 1}+

\varepsilon_1\frac{1-\nu}{2} \sum_{i=1}^n \frac{h_i}{h}(u_{\bigcom 2}+v_{\bigcom 1} )_{\bigcom 2}\delta(x-x_i)=0,

\\

\frac{1+\nu}{2}u_{\bigcom 21}+ v_{\bigcom 22}+\frac{1-\nu}{2}v_{\bigcom 11}-\nu k_{11}w_{\bigcom 2}+\varepsilon_2\sum_{i=1}^n\frac{h_i}{h}a_i ( \nu u_{\bigcom 1}+v_{\bigcom 2})_{\bigcom 2} \delta(x-x_i)=0,

\end{gather}



\vspace{-12mm}



\begin{multline}\label{bel_eq1}

\nabla^2\nabla^2 w+\frac{B}{D}k^2_{11}w-k_{11}\frac{B}{D} ( u_{\bigcom 1}+\nu v_{\bigcom 2})- 

\\

- \frac{1}{D} \bigl[ 

( T^{11}_0 w_{\bigcom 1})_{\bigcom 1}+

( T^{12}_0 w_{\bigcom 1})_{\bigcom 2}+  

(T^{12}_0 w_{\bigcom 2})_{\bigcom 1}+ 

(T^{22}_0 w_{\bigcom 2})_{\bigcom 2}\bigr]-

\\

-

\frac{1}{D}\sum_{i=1}^n\frac{h_i}{h}a_i

(T^{22}_0 w_{\bigcom 2})\delta(x-x_i)+ 

\\

+\sum_{i=1}^n \Bigl( \frac{h_i}{h}\Bigr)^3 \Phi_{3i}a_iw_{\bigcom 2222}\delta(x-x_i)+

\\

+

2(1-\nu)\sum_{i=1}^n\Bigl( \frac{h_i}{h}\Bigr)^3 \Phi_{3i}a_iw_{\bigcom 122}\bigr|_{x_i}\frac{d\delta(x-x_i)}{dx}= \\

=\frac{\gamma h}{g D}w_{\bigcom tt}\biggl[ 1+\sum_{i=1}^n \frac{h_i}{h}a_i\delta(x-x_i)\biggr] 

-\frac{\mu} {h D}w_{\bigcom t}-p_0\frac{\varkappa \mu}{D}w_{\bigcom 2}, 

\end{multline}

где $p_0\frac{\varkappa \mu}{D}w_{\bigcom 2}$ "---  относительная интенсивность поперечной нагрузки, вызванная прогибом пластинки, стандартным образом записана согласно поршневой теории \cite{myl-bib:1,myl-bib:22, myl-bib:30}; $M={v_y}/{c_0}$ "--- число Маха; 

$v_y$ "--- невозмущенная скорость набегающего газового потока; 

$c_0$ "--- скорость звука на бесконечности; $u$, $v$, $w$ "---  компоненты поля (векторы) перемещений; 

${h_i}/{h}$ "--- относительная высота $i$-того ребра, число которых $n$; 

$a_i$ "--- ширина $i$-того ребра; 

$k_{11}$ "--- кривизна цилиндрической оболочки; 

$D=\frac{E h^3}{12(1-\nu)}$; $\nu$ "--- коэффициент Пуассона; 

$E$ "--- модуль Юнга, 

$\gamma$ "--- удельный вес; 

$g$ "--- интенсивность поля тяжести; 

$$\Phi_{3i}=1+3\frac{h_i}{h}+\Bigl( \frac{h_i}{h}\Bigr)^3 ;$$ 

$\mu$ "--- коэффициент демпфирования; 

$\delta(x-x_i)$ "--- обобщенная дельта"=функция Дирака~\cite{myl-bib:25};  

$\varepsilon_j$ $(j=0,1)$ "--- знаковые числа, равные 1 или 0; 

$ T^{ij}_0 $ "--- тангенциальные усилия, вызванные нагревом пластинки до постоянной температуры $\Theta_0$, которые определяются на основании решений системы сингулярных однородных дифференциальных уравнений, описывающих безмоментное состояние термоупругой системы в компонентах поля перемещений $\bar{u}^0(u^0,v^0,w^0)$:

 \begin{multline}

 u^0_{\bigcom 11}+\frac{1-\nu}{2}u^0_{\bigcom 22}+\frac{1+\nu}{2}v^0_{\bigcom 12}-k_{11}w^0_{\bigcom 1}+

\\

 +\frac{1-\nu}{2}\sum_{i=1}^{n}\frac{h_i}{h}a_i ( u^0_{\bigcom 2}+v^0_{\bigcom 1})_{\bigcom 2}\delta(x-x_i) =0,

 \end{multline}

 \begin{multline}

 \frac{1+\nu}{2}u^0_{\bigcom 12}+v^0,_{22}+\frac{1-\nu}{2}v^0_{\bigcom 11}-\nu k_{11}w^0_{\bigcom 2}+

\\

+ 

 \sum_{i=1}^{n}\frac{h_i}{h}a_i( v^0_{\bigcom 2}+\nu u^0_{\bigcom 1})_{\bigcom 2}\delta(x-x_i)=0,

 \end{multline}

\begin{equation}\label{eq:bel:momentnot}

k_{11}^2 w^0+k_{11} (u^0_{\bigcom 1}+\nu v^0_{\bigcom 2} ) =\alpha(1+\nu)k_{11}\Theta_0.

\end{equation}





При неоднородных краевых условиях

\begin{equation*}

\begin{array}{ll}

\text{$x=0$, $x=a$:}&

u^0_{\bigcom 2}+v^0_{\bigcom 1}=0, \quad u^0_{\bigcom 1}+\nu v^0_{\bigcom 2}-k_{11}w^0=\alpha(1+\nu \Theta_0), \quad w^0=0;

\\[2mm]

\text{$y=0$, $y=b$:}&

v^0=0, \quad u^0_{\bigcom 2}+ v^0_{\bigcom 1}=0, \quad w^0=0

\end{array}

\end{equation*}

решения сингулярной системы \eqref{eq:bel:momentnot} запишутся в элементарных функциях:

\begin{equation}\label{bel_gr1}

v^0=0, \quad u^0=\alpha(1+\nu)\Theta_0 x, \quad w^0=0,

\end{equation}

а тангенциальные усилия примут вид

\begin{equation*}

T^{11}_0=T^{12}_0=0, \quad T^{22}_0=-E h\alpha\Theta_0.

\end{equation*}



Решение системы \eqref{bel_eq1} (в~которой отсутствуют инерционные слагаемые в~тангенциальной плоскости \cite{myl-bib:1,myl-bib:6, myl-bib:7}), тождественно удовлетворяющее условиям шарнирного закрепления всех сторон пологой цилиндрической оболочки

\begin{equation*}

\begin{array}{ll}

\text{при $x=0$, $x=a$:} &

T^{12}=0, \quad T^{11}=0, \quad w=0 \quad M^{11}=0;

\\[2mm]

\text{при $y=0$, $y=b$:} &

v=0, \quad T^{12}=0, \quad w=0 \quad M^{22}=0,

\end{array}

\end{equation*}

зададим в виде

\begin{gather}\label{bel_appr1}

u(x,y,t)=\tilde{u}(t)u_1 ({x}/{a} )u_2 ({y}/{b} ), \quad 

v(x,y,t)=\tilde{v}(t)v_1({x}/{a} )v_2 ({y}/{b} ),

\\

w(x,y,t)=w_{11}(t)\sin \frac{\pi x}{a} \sin\frac{\pi y}{b} +w_{12}(t)\sin \frac{\pi x}{a} \sin\frac{2\pi y}{b},

\end{gather}

где 

\begin{gather}

u_1({x}/{a}) =( {x}/{a})^4-2( {x}/{a})^3+( {x}/{a})^2,\quad

u_2({y}/{b}) =-({2}/{3})( {y}/{b})^3+( {y/}{b})^2,\\ 

v_1({x}/{a}) =({x}/{a})^4-2( {x}/{a})^3+( {x}/{a})^2,

\quad 

v_2({y}/{b}) =({y}/{b})( {y}/{b}-1).

\end{gather}



Первые два уравнения системы \eqref{bel_eq1} на основании подстановок \eqref{bel_appr1} с последующим применением процедуры Галеркина \cite{myl-bib:27,  myl-bib:29}  перепишутся в виде 

\begin{gather}

e_{11}\tilde{u}(t)+e_{12}\tilde{v}(t)=k_{11}a\Bigl[ \pi \frac{b}{a}\mathit{I}_4 w_{11}(t)+\pi\frac{b}{a}\mathit{I}_5 w_{12}(t) \Bigr],

\\

\label{bel_eq7}

e_{21}\tilde{u}(t)+e_{22}\tilde{v}(t)=\nu k_{11} a \bigl[ \pi \mathit{I}_9 w_{11}(t)+2 \pi \mathit{I}_{10} w_{12}(t)\bigr].

\end{gather}

Здесь введены следующие обозначения:

\begin{gather}

e_{11}=\frac{b}{a}\mathit{I}_1+\frac{1-\nu}{2}\frac{a}{b}\mathit{I}_2, \quad 

e_{12}=\frac{1+\nu}{2}\mathit{I}_3,\\

e_{21}=\frac{1+\nu}{2}\mathit{I}_6,\quad 

e_{22}=\frac{a}{b}\mathit{I}_7+\frac{1-\nu}{2}\frac{b}{a}\mathit{I}_8;

\\

\mathit{I}_1=\int_0^1 \int_0^1 \tilde{u}_{1 \bigcom 11} \tilde{u}_1 \tilde{u}_2^2dX dY,\quad 

\mathit{I}_2=\int_0^1 \int_0^1 \tilde{u}_{2  \bigcom 22}\tilde{u}_2\tilde{u}_1^2dX dY,

\\

\mathit{I}_3=\int_0^1 \int_0^1 \tilde{v}_{1\bigcom 1}\tilde{u}_1\tilde{u}_2 \tilde{v}_2,_{2}dX dY, \dots,

\\

\mathit{I}_{10}=\int_0^1 \int_0^1 \tilde{v}_1 \tilde{v}_2 \sin (\pi X )\cos (2\pi Y )dX dY,\quad  X={x}/{a}, \quad  Y={y}/{b}.

\end{gather}



Разрешая неоднородную алгебраическую систему относительно коэффициентов $\tilde{u}(t)$ и $\tilde{v}(t)$, получим

\begin{equation}\label{bel_eq9}

\tilde{u}=k_{11}a\bigl(L_1 w_{11}(t)+L_2 w_{12}(t) \bigr),\quad 

\tilde{v}=k_{11}a\bigl(L_3 w_{11}(t)+L_4 w_{12}(t) \bigr),

\end{equation}

где $L_j$ "--- безразмерные величины:

\begin{gather}

L_1=\frac{\pi \frac{b}{a}\mathit{I}_4e_{22}-\nu \pi \mathit{I}_9 e_{12} }{e_{11} e_{22}-e_{12}e_{21}},\quad L_2=\frac{\pi \frac{b}{a}\mathit{I}_5e_{22}-\nu 2 \pi \mathit{I}_{10} e_{12} }{e_{11} e_{22}-e_{12}e_{21}}, 

\\

L_3=\frac{-\pi \frac{b}{a}\mathit{I}_4e_{21}+\nu \pi \mathit{I}_9 e_{11} }{e_{11} e_{22}-e_{12}e_{21}},\quad L_4=\frac{-\pi \frac{b}{a}\mathit{I}_5e_{21}+\nu 2 \pi \mathit{I}_{10} e_{11} }{e_{11} e_{22}-e_{12}e_{21}}. 

\end{gather}

Окончательно  функции $u(x,y,t)$ и $v(x,y,t)$ примут следующий вид:

\begin{gather}\label{bel_eq11}

u(x,y,t)=k_{11} a \bigl(L_1 w_{11}(t)+L_2 w_{12}(t) \bigr) u_1({x}/{a} )u_2({y}/{b}),

\\

v(x,y,t)=k_{11} a \bigl(L_3 w_{11}(t)+L_4 w_{12}(t) \bigr) v_1({x}/{a} )v_2({y}/{b} ).

\end{gather}



Из третьего уравнения системы \eqref{bel_eq1}  на основании аналогичных преобразований получим систему обыкновенных однородных дифференциальных уравнений относительно коэффициентов $w_{11}(t)$ и $w_{12}(t)$:

\begin{equation}\label{bel_eq12}

\xi_{11}\left[w_{11} \right]+\xi_{12} w_{12}=0, \quad  \xi_{21}w_{11}+\xi_{22} \left[ w_{12}\right] =0,

\end{equation}

где $\xi_{kl}$ "--- дифференциальные операторы:

\begin{gather}

\xi_{11}=\frac{\gamma h a^4}{g D}\biggl( 1+2 \sum\limits_{i=1}^n \tilde{\beta}_i^s\biggr)

\frac{d^2}{dt^2}+\tilde{\mu} \frac{d}{dt}+ f_{11}\Bigl( k_{11}, \Theta_0, \frac{h_i}{h}\Bigr), 

\\

\xi_{12}=\frac{p_0 \varkappa M a^3}{D} 4\pi \frac{a}{b} \mathit{I}_{13}-\tilde{f}_{12}, \quad \xi_{21}=\frac{p_0 \varkappa M a^3}{D} 2\pi \frac{a}{b} \mathit{I}_{16}-\tilde{f}_{11},

\\

\xi_{22}=\frac{\gamma h a^4}{g D}

\biggl( 1+2 \sum\limits_{i=1}^n \tilde{\beta}_i^s\biggr)\frac{d^2}{dt^2}+\tilde{\mu} \frac{d}{dt}+ f_{12}

\Bigl( k_{11},\Theta_0,\frac{h_i}{h}\Bigr), 

\end{gather}

в которых

\begin{multline}

f_{11}=\Bigl(\pi^2+\Bigl( \frac{\pi a}{b}\Bigr) ^2 \Bigr) ^2+

12

\Bigl( \frac{h_i}{h}\Bigr)^2(k_{11}a)^2-

48

\Bigl( \frac{a}{h}\Bigr)^2 (k_{11}a)^2\Bigl(L_1\mathit{I}_{11}+\nu \frac{a}{b}L_3 \mathit{I}_{12} \Bigr) -

\\

-12(1-\nu^2)\Bigl( \frac{\pi a}{b}\Bigr)^2 \Bigl(\frac{a}{h} \Bigr)^2

\biggl(1+\sum\limits_{i=1}^n \tilde{\beta}_i^s \biggr) \alpha \Theta_0+

\Bigl( \frac{\pi a}{b}\Bigr)^4 2\sum\limits_{i=1}^n \beta_i^s+

\\

+4(1-\nu)\pi^2\Bigl(\frac{\pi a}{b} \Bigr)^2  \sum\limits_{i=1}^n \beta_i^c,

\end{multline}

\begin{equation*}

\tilde{f}_{12}=48 \Bigl(\frac{a}{h} \Bigr)^2(k_{11}a)^2\Bigl(L_2 \mathit{I}_{11}+\nu\frac{a}{b}L_4\mathit{I}_{12}\Bigr),

\end{equation*}

\begin{equation*}

\tilde{f}_{11}=48 \Bigl(\frac{a}{h} \Bigr)^2(k_{11}a)^2\Bigl(L_1 \mathit{I}_{14}+\nu \frac{a}{b}L_3\mathit{I}_{15}\Bigr),

\end{equation*}

\begin{multline}

f_{12}=\Bigl(\pi^2+\Bigl( \frac{2\pi a}{b}\Bigr) ^2 \Bigr)^2-

12\Bigl( \frac{a}{h}\Bigr)^2(k_{11}a)^2-

48\Bigl( \frac{a}{h}\Bigr)^2 (k_{11}a)^2 \Bigl(L_2\mathit{I}_{14}+\nu \frac{a}{b}L_4 \mathit{I}_{15} \Bigr) -

\\

-12(1-\nu^2)\Bigl( \frac{2\pi a}{b}\Bigr)^2

\Bigl(\frac{a}{h} \Bigr)^2

\biggl(1+\sum\limits_{i=1}^n \tilde{\beta}_i^s \biggr) \alpha \Theta_0+\Bigl( \frac{2\pi a}{b}\Bigr)^4 2\sum\limits_{i=1}^n \beta_i^s+

\\

+4(1-\nu)\pi^2\Bigl(\frac{2\pi a}{b} \Bigr)^2  \sum\limits_{i=1}^n \beta_i^c, 

\end{multline}

\begin{equation*}

\beta_i^c=\Bigl( \frac{h_i}{h}\Bigr) ^3\frac{a_i}{a}\varPhi_{3i}\cos^2 \frac{\pi x_i}{a},

\quad 

\beta_i^s=\Bigl( \frac{h_i}{h}\Bigr) ^3\frac{a_i}{a}\varPhi_{3i}\sin^2 \frac{\pi x_i}{a},

\quad 

\tilde{\beta}_i^s= \frac{h_i}{h}\frac{a_i}{a}\sin^2 \frac{\pi x_i}{a}.

\end{equation*}



Система дифференциальных уравнений \eqref{bel_eq12} подстановкой 

\begin{equation*}

w_{11}=-\xi_{12}\varPhi(t), \quad w_{12}=\xi_{11}\left[\varPhi(t) \right] 

\end{equation*}

сводится к одному дифференциальному уравнению четвертого порядка, которое в случае отсутствия  демпфирования ($\mu=0$) содержит только четные производные от функции $\varPhi(t)$:

\begin{multline}\label{bel_eq15}

\biggl[ 

\frac{\gamma h a^4}{g D}\biggl( 1+2\sum\limits_{i=1}^n \tilde{\beta}_i^s\biggr) \biggr]^2

\frac{d^4\varPhi}{dt^4} +

\biggl[ \frac{\gamma h a^4}{g D}\biggl( 1+2\sum\limits_{i=1}^n \tilde{\beta}_i^s\biggr)

( f_{11}+f_{12})  \biggr]\frac{d^2\varPhi}{dt^2}+

\\

+\Bigl[ f_{11}f_{12}-8\pi \Bigl( \frac{a}{b}\Bigr)^2\mathit{I}_{13}\mathit{I}_{16}\tilde{M}+\tilde{M} 

\Bigl(4 \pi \frac{a}{b}\mathit{I}_{13}\tilde{f}_{11}+2 \pi \frac{a}{b}\mathit{I}_{16}\tilde{f}_{12} \Bigr) - \tilde{f}_{11}\tilde{f}_{12}\Bigr] \varPhi=0,

\end{multline}

где $\tilde{M}=\frac{p_0 \varkappa M a^3}{D}$.



Если $\mu\neq0$, то дифференциальное уравнение для функции $\varPhi(t)$ запишется в~виде

\begin{equation}\label{bel_eq16}

\frac{d^4\varPhi}{dt^4}+2 e_1\frac{d^3\varPhi}{dt^3}+e_3\frac{d^2\varPhi}{dt^2}+e_4\frac{d\varPhi}{dt}+e_5\varPhi=0,

\end{equation}

где

\begin{equation*}

e_1=\frac{\tilde{\mu} }{ \frac{\gamma h a^4}{g D}\left(1+2\sum_{i=1}^n\tilde{\beta}_i^s \right) }, \quad e_3=\frac{\frac{\gamma h a^4}{g D}\left(1+2\sum_{i=1}^n\tilde{\beta}_i^s \right)\left( f_{11}+f_{12}\right)+ \tilde{\mu}}{\left( \frac{\gamma h a^4}{g D}\left(1+2\sum_{i=1}^n\tilde{\beta}_i^s \right)\right)^2},

\end{equation*}

\begin{equation*}

e_4=\frac{\tilde{\mu}\left(f_{11}+f_{12} \right)  }{\left(  \frac{\gamma h a^4}{g D}\left(1+2\sum_{i=1}^n\tilde{\beta}_i^s \right)\right)^2},\quad e_5=\frac{f_{11} f_{12}-\xi_{21}\xi_{12}}{\left( \frac{\gamma h a^4}{g D}\left(1+2\sum_{i=1}^n\tilde{\beta}_i^s \right)\right)^2}, \quad \tilde{\mu}=\frac{\mu a^4 h}{D}.

\end{equation*}



Из условия, при котором хотя бы один корень из четырех  характеристического уравнения для дифференциального уравнения \eqref{bel_eq15} будет положительным, определяется наименьшее значение  относительной скорости потока газа ${v_y}/{c_0}$, при котором прогиб термоупругой системы растет. При учете демпфирования ($\mu\neq0$) возникает вопрос об устойчивости системы \eqref{bel_eq12}, которая сведена к одному дифференциальному уравнению \eqref{bel_eq16}. 

В~этом случае это значение скорости потока определяется на основании критерия Гурвица~\cite{myl-bib:32}:

\begin{equation*}

\Delta_1=2e_1>0, \quad \Delta_2=2e_1e_3-e_4>0, \quad \Delta_3=2e_1e_3e_4-e_4^2-4e_1^2e_5>0.

\end{equation*}



Кривизна пологой цилиндрической оболочки задавалась в виде $k_{11}\hm=-{4 \tilde{\delta}}/{a^2}$ \cite{myl-bib:33, myl-bib:34}, где $\tilde{\delta}$ "--- наибольший подъем оболочки над ее планом. Выражение для относительной кривизны $\tilde{k}_{11}$ принимает вид

\begin{equation*}

\tilde{k}_{11}=k_{11} a=-4\frac{\tilde{\delta}}{h}\frac{h}{a},

\end{equation*}

где ${\tilde{\delta}}/{h}\in[0, 5]$ \cite{myl-bib:33, myl-bib:34}.



\smallskip



{\bf 2.} Численные результаты приведены в~таблице. 

Вычисления проводились при следующих значениях параметров: 

${h}/{a}=0.01$, 

${a_i}/{a}=0.01$, 

$\mu=0.001$, 

${a}/{b}=1$.





\begin{table}[b!]

	

\small \centering



Значения параметра ${v_y}/{c_0}$  в зависимости от других параметров



[The values of the parameter ${v_y}/{c_0}$ depending on other parameters]

	

\smallskip

	\renewcommand{\tabcolsep}{0.35cm}

\begin{tabular}{c|c|c|c|c|c|c|c}

	\hline 

	\multicolumn{2}{c|} {\rule{0mm}{12pt}}&  \multicolumn{3}{c|} {${h_i}/{h}=2.5$}    &  \multicolumn{3}{c} {${h_i}/{h}=5$}   \\ 

	\hline 

 \backslashbox{$\tilde{k}_{11}$}{$n$} & 0 & 1 & 3 & 5 & 1 & 3 & 5 \\ 

	\hline 

		\multicolumn{8}{c} {\rule{0mm}{12pt}for $\Theta_0=0$,  $\varepsilon_1=\varepsilon_2=0$} \\

	\hline 	

	0.10 & 0.24 & 1.50 & 3.16 & 3.50 & 6.47 & 13.35 & 14.79 \\ 

	0.15 & 0.54 & 1.40 & 3.06 & 3.40 & 6.37 & 13.25 & 14.69 \\

	0.20 & 0.97 & 1.26 & 2.92 & 3.26 & 6.23 & 13.11 & 14.55 \\  

	\hline 

		\multicolumn{8}{c} {\rule{0mm}{12pt}for $\Theta_0=2$,  $\varepsilon_1=\varepsilon_2=0$} \\

	\hline 

	0.10 & 1.08 & 0.61 & 2.22 & 2.54 & 5.55 & 12.34 & 13.73 \\ 

	0.15 & 1.39 & 0.51 & 2.12 & 2.44 & 5.45 & 12.24 & 13.63 \\ 

	0.20 & 1.81 & 0.37 & 1.98 & 2.30 & 5.31 & 12.10 & 13.49 \\ 

	\hline 

		\multicolumn{8}{c} {\rule{0mm}{12pt}for $\Theta_0=4$,  $\varepsilon_1=\varepsilon_2=0$} \\

	\hline 

		0.10 & 1.93 & 0.27 & 1.28 & 1.58 & 4.63 & 11.32 & 12.67 \\ 

		0.15 & 2.23 & 0.37 & 1.18 & 1.48 & 4.53 & 11.22 & 12.57 \\ 

		0.20 & 2.66 & 0.51 & 1.04 & 1.34 & 4.39 & 11.08 & 12.43 \\ 

	\hline 

		\multicolumn{8}{c} {\rule{0mm}{12pt}for $\Theta_0=0$,  $\varepsilon_1=\varepsilon_2=1$} \\

	\hline 		

		0.10 & 0.24 & 1.50 & 3.16 & 3.51 & 6.48 & 13.37 & 14.80 \\ 

		0.15 & 0.54 & 1.40 & 3.07 & 3.42 & 6.39 & 13.28 & 14.72 \\ 

		0.20 & 0.97 & 1.27 & 2.94 & 3.29 & 6.26 & 13.16 & 14.60 \\ 

	\hline 

		\multicolumn{8}{c} {\rule{0mm}{12pt}for $\Theta_0=2$,  $\varepsilon_1=\varepsilon_2=1$} \\

	\hline 	

		0.10 & 1.08 & 0.61 & 2.23 & 2.55 & 5.56 & 12.35 & 13.74 \\ 

		0.15 & 1.39 & 0.52 & 2.13 & 2.46 & 5.47 & 12.26 & 13.66 \\ 

		0.20 & 1.81 & 0.38 & 2.00 & 2.33 & 5.34 & 12.14 & 13.54 \\ 

	\hline 

		\multicolumn{8}{c} {\rule{0mm}{12pt}for $\Theta_0=4$,  $\varepsilon_1=\varepsilon_2=1$} \\

	\hline 	

		0.10 & 1.93 & 0.26 & 1.29 & 1.58 & 4.64 & 11.33 & 12.68 \\ 

		0.15 & 2.23 & 0.36 & 1.19 & 1.49 & 4.55 & 11.25 & 12.60 \\ 

		0.20 & 2.66 & 0.49 & 1.06 & 1.36 & 4.42 & 11.13 & 12.48 \\ 

		\hline 

	\end{tabular}  

\end{table} 



Количественный анализ выявил закономерности влияния геометрических параметров  и температуры  на величины относительных скоростей потока,  начиная с которых прогибы термоупругой системы увеличиваются во времени, и 

сделать нижеследующий выводы.

 

\begin{enumerate}

	\item[1.] Учет в дифференциальных уравнениях системы \eqref{bel_eq1} слагаемых, учитывающих <<растяжение"=сжатие>> ребер и их <<сдвиг>> в~тангенциальной плоскости (слагаемых, содержащих знаковые числа), практически не влияет на величину относительной скорости ${v_y}/{c_0}$ газового потока. По этой причине удерживать эти слагаемые в уравнениях нет необходимости.

	\item[2.] Увеличение температуры ведет к уменьшению предельной скорости потока при прочих равных условиях.

	\item[3.] Существенно повышают динамическую устойчивость геометрически нерегулярной оболочки параметры ${h_i}/{h}$ и число подкрепляющих элементов~$n$. 

	\item[4.] Величина относительной предельной скорости потока мало чувствительна к изменению относительной кривизны пологой оболочки.

\end{enumerate} 





Перечисленные закономерности выявлены для случая двух полуволн в~направлении потока и одной полуволны в перпендикулярном направлении.  

Следует отметить, что при $k_{11}=0$ полученные решения аналитически преобразуются к решениям, приведенным в работах \cite{myl-bib:31, myl-bib:35}. В случае <<холодной>> ($\Theta_0=0$) гладкой пластины   ($k_{11}=0$, ${a_i}/{a}=0$) решения принимают вид, приведенный в~\cite{myl-bib:30}.



\AdditionalContent{Конкурирующие интересы}{Заявляем, что в отношении авторства и публикации этой статьи конфликта интересов не имеем.} 



\AdditionalContent{Авторский вклад и ответственность}{Все авторы принимали участие в разработке концепции статьи и в написании рукописи. Авторы несут ответственность за предоставление окончательной рукописи в печать. Окончательная версия рукописи была одобрена всеми авторами.} 



\AdditionalContent{Финансирование}{Результаты получены в рамках выполнения государственного задания Минобрнауки России № 9.8570.2017/8.9.}



\begin{thebibliography}{10}



\RBibitem{myl-bib:1}

\by Вольмир~А.~С.

\book Оболочки в потоке жидкости и газа

\yr 1979

\publ Наука

\publaddr М.

\totalpages 320



\RBibitem{myl-bib:2}

\by Амбарцумян~C.~А., Багдасарян~Ж.~Е.

\paper Об устойчивости ортотропных пластинок, обтекаемых сверхзвуковым потоком газа

\jour Изв. АН СССР, ОТН, Механика и машиностроение

\yr 1961

\issue 4

\pages 91--96



\RBibitem{myl-bib:3}

\by Болотин~В.~В., Новичков~Ю.~Н.

\paper Выпучивание и установившийся флайтер термически сжатых панелей, находящихся в сверхзвуковом потоке

\jour Инж. журн.

\yr 1961

\issue 2

\pages 82--96







\RBibitem{myl-bib:4}

\by Мовчан~А.~А.

\paper О колебаниях пластинки, движущейся в газе

\jour ПММ

\yr 1956

\vol 20

\issue 2

\pages 211--222





\RBibitem{myl-bib:5}

\by Дун~Мин-дэ,

\paper Об устойчивости упругой пластинки при сверхзвуковом обтекании

\jour Докл. АН СССР

\yr 1958

\vol 120

\issue 4

\pages 726--729

\mathnet{http://mi.mathnet.ru/dan23108}



\RBibitem{myl-bib:5a}

\by Веденеев~В.~В. 

\paper Высокочастотный флаттер прямоугольной пластины

\jour Изв. РАН. МЖГ

\yr 2006

\issue 4

\pages 173--181









\RBibitem{myl-bib:6}

\by Огибалов~П.~М., Грибанов~В.~Ф.

\book Термоустойчивость пластин и оболочек

\yr 1968

\publ МГУ

\publaddr М.

\totalpages 520







\RBibitem{myl-bib:7}

\by Огибалов~П.~М. 

\book Вопросы динамики и устойчивости оболочек

\yr 1963

\publ МГУ

\publaddr М.

\totalpages 417

\zmath{0115.41504}





\RBibitem{myl-bib:8}

\by Болотин~В.~В.

\paper Температурное выпучивание пластин и пологих оболочек в сверхзвуковом потоке газа 

\inbook Расчеты на прочность

\yr 1960

\volinfo Вып.~6

\publaddr М.

\publ Машгиз

\pages 190--216



\RBibitem{myl-bib:9}

\by Жилин~П.~А.

\paper Линейная теория ребристых оболочек

\jour Изв. АН СССР. МТТ

\yr 1970

\issue 4

\pages 150--166



 



\RBibitem{myl-bib:10}

\by Белосточный~Г.~Н., Ульянова~О.~И.

\paper Континуальная модель композиции из оболочек вращения с термочувствительной толщиной

\jour Изв. РАН. МТТ

\yr 2011

\issue 2

\pages 184--191

\elib{https://elibrary.ru/item.asp?id=15785027}



 

 \RBibitem{myl-bib:11}

 \by Белосточный~Г.~Н., Рассудов~В.~М.

 \paper Континуальная модель термочувствительной ортотропной системы «оболочка–ребра» с учетом влияния больших прогибов

 \inbook Механика деформируемых сред

 \yr 1983

\volinfo Вып.~8

\publaddr Саратов

\publ  Сарат. политехн. ин-т

 \pages 10--22



  

\RBibitem{myl-bib:12}

\by Белосточный~Г.~Н., Рассудов~В.~М.

\paper Континуальный подход к термоустойчивости упругих систем «пластинка–ребра» 

\inbook Прикладная теория упругости

\yr 1980

\publaddr Саратов

\publ  Сарат. политехн. ин-т

\pages 94--99





 









\RBibitem{myl-bib:15}

\by Жилин~П.~А.

\paper Общая теория ребристых оболочек. Прочность гидротурбин

\inbook Тр. ЦКТИ

\yr 1968

\volinfo Вып.~8

\publaddr Л.

\pages 46--70







\RBibitem{myl-bib:16}

\by Карпов~В.~В., Сальников~А.~Ю.

\paper Вариационный метод вывода нелинейных уравнений движения пологих ребристых оболочек 

\jour Вестн. гражд. инженеров

\yr 2008

\issue 4(17)

\pages 121--124

\elib{https://elibrary.ru/item.asp?id=11675045}







\RBibitem{myl-bib:18}

\by  Белосточный~Г.~Н., Мыльцина~О.~А.

\paper Уравнения термоупругости композиций из оболочек вращения 

\jour Вестник СГТУ

\yr 2011

\issue 4 (59)

\issueinfo Вып.~1

\pages 56--64



   

\RBibitem{myl-bib:19}

\by  Онанов~Г.~Г.

\paper Уравнения с сингулярными коэффициентами типа дельта-функции и~ее производных

\jour Докл. АН СССР

\yr 1970

\vol 191

\issue 5

\pages 997--1000

\mathnet{http://mi.mathnet.ru/dan35339}

\mathscinet{http://www.ams.org/mathscinet-getitem?mr=0470298}

\zmath{https://zbmath.org/?q=an:0215.20101}









\RBibitem{myl-bib:21}

\by  Белосточный~Г.~Н.

\paper  Аналитические методы определения замкнутых интегралов сингулярных дифференциальных уравнений термоупругости геометрически нерегулярных оболочек

\jour Докл. Акад. воен. наук

\yr 1999

\issue 1

\pages 14--26





\Bibitem{myl-bib:22a}

\by Geckeler~J.~W.

\paper Elastostatik

\lang In German

\zmath{54.0844.01}

\inbook Handbuch der Physik 

\bookvol 6

\pages 141--308

\yr 1928





\RBibitem{myl-bib:22}

\by Ильюшин~А.~А.

\paper  Закон плоских сечений в аэродинамики больших сверхзвуковых скоростей 

\jour ПММ

\yr 1956

\issue 6

\pages 733--755







 

\RBibitem{myl-bib:30}

\by Вольмир~А.~С. 

\book  Устойчивость деформируемых систем

\yr 1967

\publ Наука

\publaddr М.

\totalpages 984



 

 



 

\RBibitem{myl-bib:25}

\by Владимиров~В.~С. 

\book Обобщенные функции в математической физике

\yr 1979

\publ Наука

\publaddr М.

\totalpages 319

\zmath{0515.46033}









\RBibitem{myl-bib:27}

\by Канторович~Л.~В., Крылов~В.~И.

\book  Приближенные методы высшего анализа

\yr 1962

\publ Физматлит

\publaddr М.

\totalpages 708

\zmath{0100.33102}









\Bibitem{myl-bib:29}

\by Rektorys~K.

\book Variational methods in mathematics, science and engineering

\zmath{0481.49002}

\publaddr Dordrecht, Boston,  London

\publ D.~Reidel Publ. 

\totalpages 571 

\yr 1980

\moreref

\crossref{10.1007/978-94-011-6450-4}. \nofrills





\RBibitem{myl-bib:31}

\by Белосточный~Г.~Н., Рассудов~В.~М. 

\paper  Термоупругие системы типа «пластинка–ребра» в сверхзвуковом потоке газа  

\inbook Прикладная теория упругости

\publaddr Саратов

\volinfo Вып.~8

\publ Сарат. политехн. ин-т

\yr 1983

\pages 114--121







\RBibitem{myl-bib:32}

\by Егоров~К.~В.  

\book  Основы теории автоматического регулирования

\yr 1967

\publ Энергия

\publaddr М.

\totalpages 648





\RBibitem{myl-bib:33}

\by Рассудов~В.~М., Красюков~В.~П., Панкратов~Н.~Д.  

\book  Некоторые задачи термоупругости пластинок и пологих оболочек 

\yr 1973

\publ Сарат. ун-т

\publaddr Саратов

\totalpages 155





\RBibitem{myl-bib:34}

\by Назаров~А.~А.

\book Основы теории и методы расчета пологих оболочек

\yr 1966

\publ  Стройиздат

\publaddr Л., М.

\totalpages 304





\RBibitem{myl-bib:35}

\by Мыльцина~О.~А., Белосточный~Г.~Н. 

\paper  Устойчивость нагретой ортотропной геометрически нерегулярной пластинки в сверхзвуковом потоке газа 

\jour Вестник Пермского национального исследовательского политехнического университета. Механика

\yr 2017

\issue 4

\pages 109--120

\crossref{10.15593/perm.mech/2017.4.08}















\end{thebibliography}





\newpage



\makeentitle



\vspace{-5mm}





\AdditionalContent{Competing interests}{We declare that we have no conflicts of interest in the authorship and publication of this article.} 



\AdditionalContent{Authors' contributions and responsibilities}{Each author has participated in the article concept development and in the manuscript writing. The authors are absolutely responsible for submitting the final manuscript in print. Each author has approved the final version of manuscript.} 





\AdditionalContent{Funding}{The results have been obtained within the State Assignment of the Ministry of Education

and Science of the Russian Federation no. 9.8570.2017/8.9.}









\begin{thebibliography}{10}



\Bibitem{myl-bib:1:en}

\lang In Russian

\by Vol'mir~A.~S.

\book Obolochki v potoke zhidkosti i gaza \rm [The shells in flow liquid and gas] 

\publaddr Moscow

\publ Nauka

\yr 1979

\totalpages 320



\Bibitem{myl-bib:2:en}

\lang In Russian

\by Ambartsumyan~S.~A. Bagdasaryan~Zh.~E.

\paper Of the stability of orthotropic plates streamlined by a supersonic gas flow

\jour Izv. Akad. Nauk SSSR, OTN, Mekhanika i Mashinostroenie

\yr 1961

\issue 4

\pages 91--96



\Bibitem{myl-bib:3:en}

\lang In Russian

\by Bolotin~V.~V., Novichkov~Yu.~N. 

\paper Buckling and steady-state flayter thermally compressed panels in supersonic flow

\jour Inzhenernyi Zhurnal 

\yr 1961

\issue 2

\pages 82--96







\Bibitem{myl-bib:4:en}

\lang In Russian

\by Movchan~A.~A. 

\paper Oscillations of the plate moving in the gas

\jour Prikladnaia Matematika i Mekhanika

\yr 1956

\vol 20

\issue 2

\pages 211--222





\Bibitem{myl-bib:5:en}

\lang In Russian

\by Tung~Ming-teh,

\paper The stability of an elastic plate in a supersonic flow

\jour Dokl. Akad. Nauk SSSR

\yr 1958

\vol 120

\issue 4

\pages 726--729



\Bibitem{myl-bib:5a}

\by Vedeneev~V.~V.

\jour Fluid Dyn.

\crossref{10.1007/s10697-006-0083-2}

\yr 2006

\vol 41

\issue 4

\pages 641–648 

\paper High-frequency flutter of a rectangular plate









\Bibitem{myl-bib:6en}

\lang In Russian

\by Ogibalov~P.~M., Gribanov~V.~F. 

\book Termoustoichivost' plastin i obolochek \rm  [Thermostability of plates and shells]

\yr 1968

\publaddr Moscow

\publ Moscow State Univ.

\totalpages 520







\Bibitem{myl-bib:7en}

\lang In Russian

\by Ogibalov~P.~M.

\book Voprosy dinamiki i ustoichivosti obolochek\rm  [The problems of dynamics and stability of shells]

\publaddr Moscow

\publ Moscow State Univ.

\yr 1963

\totalpages 417







\Bibitem{myl-bib:8:en}

\lang In Russian

\by Bolotin~V.~V.

\paper Thermal buckling of plates and shallow shells in a supersonic gas flow

\inbook Raschety na prochnost'

\publaddr Moscow

\publ Mashgiz

\yr 1960

\volinfo Issue~6

\pages 190--216



\Bibitem{myl-bib:9:en}

\lang In Russian

\by Zhilin~P.~A.

\paper Linear theory of ribbed shells

\jour Izv. Akad. Nauk SSSR, Mekhanika Tverdogo Tela

\yr 1970

\issue 4

\pages 150--166



 



\Bibitem{myl-bib:10:en}

\paper Continuum model for a composition of shells of revolution with thermosensitive thickness

\by Belostochnyi~G.~N., Ul'yanova~O.~I.

\jour Mechanics of Solids

\yr 2011

\vol 46

\issue 2

\pages 184--191

\elib{https://elibrary.ru/item.asp?id=16982711}

\crossref{10.3103/S0025654411020051}



 

\Bibitem{myl-bib:11:en}

\lang In Russian

\by Belostochnyi~G.~N., Rassudov~V.~M.

\paper Continuum model of orthotropic heat-sensitive ``shell-ribs'' system taking into account the influence of large deflection

 \yr 1983

\volinfo Issue~8

\inbook Prikladnaia teoriia uprugosti

\publ Saratov Polytechnic Inst.

\publaddr Saratov 

 \pages 10--22



  

\Bibitem{myl-bib:12:en}

\lang In Russian

\by Belostochnyi~G.~N., Rassudov~V.~M.

\paper Continuum approach to the thermal stability of elastic ``plate-ribs'' systems

\inbook Prikladnaia teoriia uprugosti

\publ Saratov Polytechnic Inst.

\publaddr Saratov 

\yr 1980

\pages 94--99







\Bibitem{myl-bib:15:en}

\lang In Russian

\by Zhilin~P.~A.

\paper General theory of ribbed shells. The strength of turbines

\inbook  Proc. of the Central Turbine-Boiler Institute

\yr 1968

\volinfo Issue~8

\publaddr Leningrad

\pages 46--70







\Bibitem{myl-bib:16:en}

\lang In Russian

\by Karpov~V.~V., Sal'nikov~A.~Yu.

\paper Variational method output nonlinear equations of motion of shallow ribbed shells

\jour Vestnik Grazhdanskikh Inzhenerov

\yr 2008

\issue 4(17)

\pages 121--124





\Bibitem{myl-bib:18:en}

\lang In Russian

\by Belostochnyi~G.~N., Myltcina~O.~A.

\paper Thermoelasticity equations of shells compositions

\jour Vestnik Saratovskogo Tekhnicheskogo Universiteta

\yr 2011

\issue 4 (59)

\issueinfo Issue~1

\pages 56--64



   

\Bibitem{myl-bib:19:en}

\lang In Russian

\by Onanov~G.~G.

\paper Equations with singular coefficients of the type of the delta-function and its derivatives

\jour Dokl. Akad. Nauk SSSR

\yr 1970

\vol 191

\issue 5

\pages 997--1000









\Bibitem{myl-bib:21:en}

\lang In Russian

\by Belostochnyi~G.~N.

\paper Analytical methods for the determination of closed integrals of singular differential equations of thermoelasticity, geometrically irregular shells

\jour Doklady Akademii Voennykh Nauk

\yr 1999

\issue 1

\pages 14--26





\Bibitem{myl-bib:22a:en}

\by Geckeler~J.~W.

\paper Elastostatik

\lang In German

\zmath{54.0844.01}

\inbook Handbuch der Physik 

\bookvol 6

\pages 141--308

\yr 1928





\Bibitem{myl-bib:22:en}

\lang In Russian

\by Il'yushin~A.~A.

\paper Law of plane sections in aerodynamics of high supersonic velocity

\jour Prikladnaia Matematika i Mekhanika

\yr 1956

\issue 6

\pages 733--755







 

\Bibitem{myl-bib:30:en}

\lang In Russian

\by Vol'mir~A.~S.

\book  Ustoichivost' deformiruemykh sistem \rm [Stability of deformable systems]

\publaddr Moscow

\publ Nauka

\yr 1967

\totalpages 984



 

 



 

\Bibitem{myl-bib:25:en}

\by Vladimirov~V.~S.

\book Generalized functions in mathematical physics

\zmath{0515.46034}

\publaddr Moscow

\publ Mir Publ.

\totalpages  362 

\yr 1979









\Bibitem{myl-bib:27:en}

\by Kantorovich~L.~V., Krylov~V.~I.

\book Approximate methods of higher analysis

\zmath{0083.35301}

\publaddr Groningen

\publ P.~Noordhoff Ltd.

\totalpages xii+681

\yr 1958











\Bibitem{myl-bib:29:en}

\by Rektorys~K.

\book Variational methods in mathematics, science and engineering

\zmath{0481.49002}

\publaddr Dordrecht, Boston,  London

\publ D.~Reidel Publ. 

\totalpages 571 

\yr 1980

\moreref

\crossref{10.1007/978-94-011-6450-4}. \nofrills





\Bibitem{myl-bib:31:en}

\lang In Russian

\by Belostochnyi~G.~N., Rassudov~V.~M.

\paper Thermoelastic system of ``plate-ribs'' type  in a supersonic gas flow

\inbook Prikladnaia teoriia uprugosti

\publ Saratov Polytechnic Inst.

\publaddr Saratov 

\volinfo Issue~8

\yr 1983

\pages 114--121







\Bibitem{myl-bib:32:en}

\lang In Russian

\by Egorov~K.~V.

\book Osnovy teorii avtomaticheskogo regulirovaniia \rm [Fundamentals of the theory of automatic control]

\publaddr Moscow

\publ Energiia

\yr 1967

\totalpages 648





\Bibitem{myl-bib:33:en}

\lang In Russian

\by Rassudov~V.~M., Krasiukov~V.~P., Pankratov~N.~D.  

\book  Nekotorye zadachi termouprugosti plastinok i pologikh obolochek  \rm [Some problems of thermoelasticity of plates and flat covers]

\yr 1973

\publ Saratov Univ.

\publaddr Saratov

\totalpages 155

 





\Bibitem{myl-bib:34:en}

\lang In Russian

\by Nazarov~A.~A.

\book Osnovy teorii i metody rascheta pologikh obolochek \rm [Fundamentals of the Theory and Methods for Designing Shallow Shells]

\yr 1966

\publ  Stroiizdat

\publaddr Leningrad, Moscow

\totalpages 304





\Bibitem{myl-bib:35:en}

\lang In Russian

\by Myltcina~O.~A., Belostochnyi~G.~N.

\paper Stability of heated orthotropic geometrically irregular plate in a supersonic gas flow

\jour PNRPU Mechanics Bulletin

\yr 2017

\issue 4

\pages 109--120

\crossref{10.15593/perm.mech/2017.4.08}





\end{thebibliography}







%\end{document}

 \newpage

%\input{./preamble}
%
%\usepackage{samgtu-name}
%
\vsgtu{1624} %registration number
\FirstPage{762}
\LastPage{773}
%
%\pdfcrop
%\draft  
%
\ShortCommunication
%
%\begin{document} 




\hfill \qrcode[hyperlink,height=.8in]{http://doi.org/10.14498/vsgtu1624}
\vspace{-.9in}

\udc{539.3}
\msc{74C10}

\rutitle{Об одном подходе к определению предельной несущей~способности
	механических систем с~разупрочняющимися элементами}
\rutitleshort{Об одном подходе к определению предельной несущей способности
	механических систем\dots}

\entitle{One approach to determination of the ultimate load-bearing capacity of mechanical systems with softening elements}
\entitleshort{One approach to determination of the ultimate load-bearing capacity of mechanical systems\dots}

\ruauthor{В.~В.~Стружанов\affil{1}, А.~В.~Коркин\affil{2}, А.~Е.~Чайкин\affil{2}}
{С\,т\,р\,у\,ж\,а\,н\,о\,в~В.~В., К\,о\,р\,к\,и\,н~А.~В., Ч\,а\,й\,к\,и\,н~А.~Е.}
\enauthor{V.\,V.~Struzhanov\affil{1}, A.\,V.~Korkin\affil{2}, A.\,E.~Chaykin\affil{2}} {S\,t\,r\,u\,z\,h\,a\,n\,o\,v~V.~V., K\,o\,r\,k\,i\,n~A.~V., C\,h\,a\,y\,k\,i\,n~A.~E.}

\ruaffil{\ruimash.
\and 
\ruurfuien.}

\enaffil{\enimash.
\and 
\enurfuien.}

\ruauthorsdetails{3}{% Количество авторов
\fullauthor{\rupid{39278}{Валерий Владимирович Стружанов}}
\CAuthor % Автор-корреспондент
\orcid{0000-0002-3669-2032} % ORCID ID автора 
\regalia{доктор физико-математических наук, профессор} 
\position{главный научный сотрудник} 
\dept[p]{~лаб. микромеханики материалов} 
\email{stru@imach.uran.ru}
\fullauthor{\rupid{73471}{Александр Владимирович Коркин}}
\orcid{0000-0003-3533-4257} % ORCID ID автора 
\regalia{аспирант} 
\email{alexkorkin@list.ru}
\fullauthor{\rupid{143469}{Алексей Евгеньевич Чайкин}}
\orcid{0000-0001-5582-2384} % ORCID ID автора 
\regalia{аспирант}  
\email{chaykin.ae@yandex.ru}\vspace{-3mm}
}

\enauthorsdetails{3}{
\fullauthor{\enpid{39278}{Valery V. Struzhanov}}
\CAuthor % Сorresponding Author
\orcid{0000-0002-3669-2032} % ORCID ID avtora 
\regalia{Dr. Phys. \& Math. Sci., Professor} 
\position{Chief Reseacher} 
\dept{Lab. of Matherial Micromechanics}
\email[br]{stru@imach.uran.ru}
\fullauthor{\enpid{73471}{Aleksandr V. Korkin}}
\orcid{0000-0003-3533-4257} % ORCID ID avtora 
\regalia{Postgraduate Student} 
\email{alexkorkin@list.ru}
\fullauthor{\enpid{143469}{Aleksey E.  Chaykin}}
\orcid{0000-0001-5582-2384} % ORCID ID avtora 
\regalia{Postgraduate Student} 
\email{chaykin.ae@yandex.ru}
}

\rukeywords{потенциальная функция, вырожденные критические точки, единый потенциал, матрица Гессе, тонкостенный резервуар, предельное давление}
\enkeywords{potential function, degenerate critical points, unified potential, Hesse matrix, thin walled reservoir, limiting pressure}

\ruabstract{Изложены основные положения теории расчета предельных нагрузок, действующих на дискретные механические системы с~разупрочняющимися элементами. Методика опирается на численное определение вырожденных критических точек потенциальной функции системы, где происходит переход от устойчивости процесса нагружения к~неустойчивости (катастрофа, разрушение), и~позволяет избежать решения большого числа нелинейных уравнений равновесия. 
В~качестве примера применения методики решена задача об определении предельного внутреннего давления в~тонкостенном цилиндрическом резервуаре. При построении потенциальной функции системы использован специально построенный единый потенциал для плоского квадратного элемента материала в~условиях двухосного растяжения, описывающий все стадии деформирования, включая и~разупрочнение.}

\enabstract{The fundamental provisions of the limiting load calculation theory are presented for a discrete mechanical system with softening elements. The method is based on the numerical determination of degenerate critical points for the potential function of the system. At these points there is a transition from the stability of the loading process to instability such as a catastrophe or a failure. This approach helps to avoid solving a large number of nonlinear equilibrium equations. The problem of determining the limiting internal pressure in a thin walled cylindrical tank is solved as an example. A unified potential specially defined for a flat square element of material in biaxial tension is used in developing a potential function of the system. It describes all stages of deformation including the softening stage.}

\newdate{struzha}{11}{05}{2018}
\date{\displaydate{struzha}}
\newdate{struzhb}{11}{10}{2018}
\revisiondate{\displaydate{struzhb}}
\newdate{struzhc}{12}{11}{2018}
\accepteddate{\displaydate{struzhc}}

\newdate{struzhe}{11}{12}{2018}
\renewcommand{\JournalDataOnline}{\displaydate{struzhe}}
%\ForPeerReview
\makerutitle

\newpage

\Section[n]{Введение}
Вычисление предельной несущей способности механических систем основывается на расчете параметров равновесных состояний при постепенно возрастающей нагрузке. Затем эти параметры тем или иным способом сводят к некоторому скаляру и сравнивают его с предельной величиной, полученной в эксперименте. Например, на каждом шаге догружения определяют напряженное состояние элемента конструкции, находят соответствующий инвариант тензора напряжений, который сравнивают с его предельной величиной, взятой из эксперимента.

Возможен и другой подход, основанный на исследовании устойчивости положений равновесия, когда разрушение связывают с невозможностью равновесия, то есть с моментом потери устойчивости [\citen{bib1, bib2}]. Реализация данного подхода возможна, если учитывать физически неустойчивые состояния материала (разупрочнение) [\citen{bib2}--\citen{bib6}]. Но и в этом случае возникает необходимость в решении не удовлетворяющих условиям корректности по Адамару [\citen{bib8}]  больших систем нелинейных уравнений равновесия для определения критических точек потенциальной функции механической системы, что является нетривиальной задачей [\citen{bib9}]. Кроме того, выделение из них вырожденных критических точек, где происходит смена устойчивости на неустойчивость, требует применения нетрадиционного аппарата теории катастроф [\citen{bib10,bib11}]. Однако информация о параметрах равновесия (критических точках потенциальной функции системы) при возрастающих нагрузках, вообще говоря, является излишней. При определении предельной несущей способности нужно знать только вырожденные критические точки. В данной работе представлена методика, которая позволяет избежать аналитического решения нелинейных уравнений. 
Она основана на специальной числовой процедуре приближенного выбора вырожденных критических точек. В качестве примера рассмотрена задача о~разрушении цилиндрического резервуара под действием внутреннего давления. Для построения потенциальной функции разработана методика построения единого потенциала для элемента материала, находящегося в плоском напряженном состоянии и описывающего его свойства как при упрочнении, так и при разупрочнении.

\smallskip

\Section{Общие положения} Положение механической системы в пространстве определяется конечным набором обобщенных координат (параметров состояния)  $ q_i$, $i = 1, 2, \ldots, N $. 
Величины параметров состояния зависят от набора задаваемых параметров управления системой $ s_j$, 
$j = 1, 2, \ldots, M $ 
(внешние силы, перемещения некоторых точек системы, физико"=механические свойства элементов и т.~п.). 
Причем параметры состояния должны принимать такие значения, при которых механическая система находится в состоянии равновесия. 
Разрушение системы есть явление того же порядка, что и явление невозможности равновесия [\citen{bib1}]. 
Когда равновесие становится неустойчивым, то реализуется или скачкообразный переход в ближайшее устойчивое равновесие, или глобальная катастрофа (разрушение). Поэтому под предельными значениями параметров управления следует понимать такие их величины, при реализации которых нарушается устойчивость системы.

В тех случаях, когда механическая система является градиентной [\citen{bib11}], для определения предельных значений параметров управления можно воспользоваться следующей методикой. Градиентная система при квазистатическом нагружении с постоянной температурой описывается потенциальной функцией (лагранжианом) $ \Pi(q_i ,s_j) $, связывающей параметры состояния и~управления. 
Эта функция есть сумма потенциальных функций элементов системы, которые могут быть как вогнутыми, так и~выпукло"=вогнутыми. 
Выпуклые вниз (вогнутые) участки отвечают  устойчивым состояниям элемента, выпуклые (выпуклые вверх) "--- неустойчивым.  
Точки перехода от выпуклости к~вогнутости "--- пограничные точки, являющиеся в общем случае седловыми точками соответствующих потенциальных функций. 

Когда механическая система находится в~равновесии (устойчивом или неустойчивом), для связи параметров состояния и~управления в~этом равновесии можно использовать уравнение Лагранжа второго рода [\citen{bib12}]. 
Тогда уравнение равновесия имеет вид
\begin{equation}\label{S_eq1}
\frac{\partial\Pi(q_i,s_j)}{\partial q_i} = 0, \quad  i=1, 2, \ldots,N.
\end{equation}
Если функция   является выпукло"=вогнутой, то при заданных управляющих параметрах система \eqref{S_eq1} может иметь одно или более одного решения, как устойчивых, так и~неустойчивых. 
Следовательно, в~данном случае не выполняются условия Адамара [\citen{bib8}].

Все решения системы \eqref{S_eq1} для всевозможных параметров управления образуют совокупность критических точек потенциальной функции $ \Pi $.  
Известно [\citen{bib11}], что смена типа равновесия (устойчивого или неустойчивого) происходит в~вырожденных критических точках, которые определяются из совместного решения уравнений \eqref{S_eq1} и~уравнения, получающегося приравниванием к~нулю детерминанта матрицы устойчивости (матрицы Гессе потенциальной функции):
\begin{equation}\label{S_eq2}
\mathop{\rm det}\Bigl| \frac{\partial^{2}\Pi(q_i,s_j)}{\partial q_m \partial s_n}\Bigr|  = 0,  \quad m,n = 1, 2, \ldots, N.
\end{equation}
Очевидно, что уравнение \eqref{S_eq2} выделяет из множества критических точек вырожденные критические точки. 

Таким образом, для отыскания предельных значений параметров управления, при достижении которых система теряет устойчивость, сначала, как правило, вычисляются все критические точки (решения уравнений \eqref{S_eq1}), затем выделяются вырожденные точки и после подстановки их в уравнение \eqref{S_eq1} определяются соответствующие параметры управления. Те параметры управления, при которых в первый раз достигается вырожденная критическая точка, то есть происходит смена устойчивого состояния механической системы на неустойчивое, и будут предельными значениями внешних нагрузок.

При реализации данной системы нахождения предельных значений параметров управления необходимо искать решения большого числа нелинейных уравнений, что представляет достаточно сложную задачу. 
Однако задача о~вычислении предельных значений параметров управления не требует, вообще говоря, знания всех критических точек функции $ \Pi $ (параметров всех положений равновесия механической системы). Необходимы только вырожденные точки, удовлетворяющие уравнению \eqref{S_eq2}. Кроме того, для технических приложений можно использовать их приближенные значения, вычисленные с достаточной степенью точности. 

Вычисление приближенных значений вырожденных критических точек можно осуществить, применяя следующую численную процедуру. Рассмотрим евклидово пространство $ \mathbb R^M \times \mathbb R^N $ "--- декартово произведение евклидовых пространств размерности $ M $ и $ N $. 
Элементами этого пространства являются числа $ ( q_i,s_j) $. 
Задавая физически обоснованные пределы изменения параметров состояния и управления, получаем $ (M{+}N) $"=мерный куб в пространстве $\mathbb  R^M \times \mathbb R^N $, 
который разбиваем сеткой узлов на множество «кубиков» с~заданным шагом. 
Затем в~каждом узле сетки разбиения вычисляем значения компонент матрицы Гессе и~величину ее определителя. 
Таким образом, каждому узлу ставится в~соответствие значение детерминанта матрицы Гессе. 

После этого выделяем те узлы, в которых детерминант близок к нулю с достаточной степенью точности. 
Полученное множество $B$ точек пространства $ \mathbb R^M \times \mathbb R^N $ содержит в~себе все вырожденные критические точки потенциальной функции~$ \Pi $. 
Чтобы выделить из множества $B$ критические точки функции $ \Pi $, следует подставить элементы множества $B$ в~уравнения равновесия \eqref{S_eq1}.  
Те~из них, которые удовлетворяют этим уравнениям с~заданной степенью точности, будут вырожденными критическими точками. 
Параметры управления,  соответствующие ближайшей (по евкледовой метрике) к~началу координат вырожденной точке, образуют совокупность искомых предельных значений управлений. 
Если существует группа вырожденных точек, находящихся на одинаковом наименьшем расстоянии от начала координат, то получим различные сочетания значений параметров управления.

Таким образом, минуя стадию расчета параметров равновесных состояний, можно найти критические нагрузки, что достаточно для оценки прочности механических систем.

\smallskip

\Section{Модельный пример} Применим изложенную выше методику к определению предельных параметров управления (нагрузок) при растяжении модельной стержневой системы, изображенной на  рис.~\ref{Sfig1}. 
Стержень {\sl 2} "--- абсолютно упругий с жесткостью при растяжении, равной $c$. 
Стержень {\sl 1} изготовлен из разупрочняющегося материала.

\begin{figure}[b!]
\centering
		\includegraphics[width=0.4\textwidth]{fig_stru1}
		\caption{Растягиваемая стержневая система с разупрочняющимся элементом} \label{Sfig1}
		\smallskip \footnotesize
		[Figure~\ref{Sfig1}. 		A rod system with a softening element under tension]

\end{figure}

\smallskip

{\bf 2.1.}
При мягком нагружении растягивающая сила $ p $ "--- параметр управления, 
а~удлинение $ x $  стержня {\sl 1}  и перемещение $ u $ правого свободного конца   стержня~{\sl 2}  "--- параметры состояния (см.  рис.~\ref{Sfig1}). 
Сопротивление растяжению стержня {\sl 1} задано функцией $ g(x) $ 
(зависимость растягивающей силы от удлинения), обладающей как восходящей ветвью (упрочнение), так и~ветвью, падающей до нуля (стадия разупрочнения) [\citen{bib2, bib5}].  
В~этом случае его потенциальная функция, равная $\displaystyle \int ^{x}_{0} g(x)\, dx $, является выпукло"=вогнутой, а лагранжиан системы~[\citen{bib5}] 
$$
Q^p = \int^{x}_{0} g(x)  dx + \frac{c}{2} (u-x)^2 - \int ^{u}_{0} p  du,
$$	
где второе слагаемое "--- энергия упругих деформаций стержня {\sl 2}, последнее "--- работа внешней силы, взятая со знаком минус.

Уравнения равновесия имеют вид
\begin{equation}\label{S_eq3}
Q^{p}_{\bigcom x} = g(x) - c(u-x) = 0, \quad  Q^{p}_{\bigcom u} = c(u-x) - p = 0.
\end{equation}
Здесь запятой обозначены частные производные по соответствующим аргументам.
Приравнивая к нулю детерминант матрицы Гессе, получим уравнение 
\begin{equation}\label{S_eq4}
g_{\bigcom x} = 0.
\end{equation}

Отметим, что это уравнение определяет максимум функции  $ g(x) $. 
Его можно разрешить, используя следующую процедуру. 

Проведем разбиение оси $X $ узлами с~достаточно малым шагом. 
Подставляя координаты этих узлов в~уравнение \eqref{S_eq4}, находим среди них значение $ x' $, приближенно ему удовлетворяющее. 
Затем из уравнений \eqref{S_eq3} для данного $ x' $ определяем $ p' $ "--- предельную величину растягивающей силы.

\smallskip

{\bf 2.2.}
В случае жесткого нагружения, когда растяжение осуществляется заданным перемещением $ u $ (параметр управления), а $ x $ является единственным параметром состояния, лагранжиан системы есть функция $ Q^u $, равная двум первым слагаемым в выражении для $ Q^p $. 

Тогда имеем лишь одно уравнение равновесия (первое уравнение в~системе \eqref{S_eq3}). 
В~матрице Гессе только один элемент. 
Тогда равенство детерминанта нулю дает уравнение $ g'_{x} + c = 0 $. 
Его приближенное решение ищем так же, как и~решение уравнения \eqref{S_eq4}. 

В зависимости от величины параметра $ c $ возможны три варианта. 

В первом "--- решения не существует. Следовательно, деформирование системы идет равновесно вплоть до разделения стержня {\sl 1} на части при значении $ x = x^z $, когда падающая ветвь функции $ g(x) $ достигает оси $ X $. Предельное значение параметра управления получаем из уравнения равновесия после подстановки $ x = x^z $: 
$$ u^z = g(x^z)/c+ x^z.$$ 

Во втором варианте имеется только одно приближенное решение $ x =  x_1 $. 
После достижения параметром $ x $ значения этой величины система переходит в~неустойчивое состояние. 
Поэтому предельный параметр управления $$ u_1 = g(x_1)/c + x_1 .$$ 

И, наконец, в третьем варианте имеем два решения $ x_1$, $x_2$  $(x_1 > x_2) $. 
Тогда предельное значение  $ u = u_1 $.

Отметим, что такие же результаты были получены и в работах [\citen{bib5, bib9}] только после решения нелинейных уравнений равновесия при постепенном возрастании параметров управления и построения соответствующих многообразий катастроф и сепаратрис, состоящих из вырожденных критических точек.

\smallskip

\Section{Вогнуто"=выпуклый потенциал при двухосном растяжении тонкой пластины} 
Для реализации положений теории, изложенной выше, необходимо знать вогнуто"=выпуклый потенциал элементов системы. 
В~механике сплошной среды эта функция должна быть определена для элемента материала вплоть до разупрочнения и~разделения на фрагменты.

Рассмотрим плоский квадратный элемент материала с~единичными размерами, находящийся в плоско"=напряженном состоянии, причем касательные напряжения и~деформации отсутствуют. 
Деформирование осуществляется заданием нормальных деформаций $ \varepsilon_1$, $\varepsilon_2 $ (двуосное растяжение). Если материал находится в состоянии упругости, то его  потенциальная энергия 
\begin{equation}\label{S_eq5}
V^e(\varepsilon_1, \varepsilon_2) = \frac{1}{2}(\sigma_1\varepsilon_1 + \sigma_2\varepsilon_2) = \frac{E}{2(1-\nu^2)} (\varepsilon^2_1 + 2\nu \varepsilon_1 \varepsilon_2 + \varepsilon^2_2 ) ,
\end{equation}
где $ E $ "--- модуль Юнга, $ \nu $ "--- коэффициент Пуассона. 
Здесь использован закон Гука, связывающий нормальные напряжения $ \sigma_1$, $\sigma_2 $ и~нормальные деформации $ \varepsilon_1$, $\varepsilon_2 $ и записанный в векторно"=матричной форме:
\begin{equation}\label{S_eq6}
\begin{pmatrix}
\sigma_1 \\ \sigma_2
\end{pmatrix}
 = C 
\begin{pmatrix} \varepsilon_1 \\ \varepsilon_2
\end{pmatrix}.
\end{equation}
В этом выражении 
$$
C = \frac{E}{(1-\nu^2)} \begin{pmatrix} 1 & \nu \\ \nu & 1 \end{pmatrix}
$$
можно трактовать как матричный коэффициент пропорциональности.

Формула \eqref{S_eq5} определяет кривые второго порядка на плоскости $\varepsilon_1 \varepsilon_2 $ 
(линии уровня потенциала $ V^e $), а~именно центральные эллипсы, главные оси которых наклонены к~декартовым осям $ \varepsilon_1$, $\varepsilon_2 $ под углом $ \pi/4 $. 
Большая ось располагается во втором и~четвертом квадрантах, малая "--- в~первом и~третьем. 
Для полуосей  выполняется отношение 
$$ \dfrac{a^{2}}{\rho^2} = \frac{(1-\nu)}{(1+\nu)}, $$ 
где $ a $ "--- малая полуось, $ \rho $ "--- большая полуось. 
Точки, расположенные на малой оси, отвечают равномерному растяжению $ (\varepsilon_1 = \varepsilon_2) $ или сжатию $ (-\varepsilon_1 = -\varepsilon_2) $, а~точки на большой оси определяют чистый сдвиг   
$ (\varepsilon_1 = -\varepsilon_2$, $ -\varepsilon_1 = \varepsilon_2 ) $. 
Произвольная точка  $ (\varepsilon_1,\varepsilon_2) $ лежит на линии уровня с большой полуосью, равной
\begin{equation}\label{S_eq7}
\rho = \frac{1}{\sqrt{2}}\sqrt{{(\varepsilon_1 + \varepsilon_2)}^2 \frac{1+\nu}{1-\nu} + {(\varepsilon_1 - \varepsilon_2)}^2} .
\end{equation}

При чистом сдвиге $ (\varepsilon_1 = \varepsilon$, $\varepsilon_2 = - \varepsilon)$  $\rho = \gamma /\sqrt{2}$, где    $ \gamma = \varepsilon_1 - \varepsilon_2 = 2 \varepsilon $ "--- сдвиг.
После исчерпания упругости для описания свойств материала можно было бы использовать теорию малых упруго-пластических деформаций [\citen{bib13, bib14}], в~которой предполагается упругость объемной деформации и~пропорциональность девиаторов тензоров напряжений и деформаций. 
Однако после выхода на стадию разупрочнения первая гипотеза не выполняется, так как материал уже существенно поврежден множеством микротрещин и объемный модуль изменяется. 

Чтобы остаться в рамках деформационной теории, следует отказаться от первой гипотезы, а~вместо второй сохранить пропорциональность напряжений и~деформаций с некоторым уже переменным, зависящим от деформаций матричным коэффициентом. 
Например, в рассматриваемой задаче %после перехода материала  при активном нагружении 
этот матричный коэффициент в законе \eqref{S_eq6} можно принять в~виде $ C \psi(\varepsilon_1,\varepsilon_2)$, 
где $ \psi $ "--- некоторая скалярная функция, требующая специального определения.

Следствием принятого предположения является тот факт, что в~области неупругости линии уровня функции энергии деформаций $ V $ (потенциальной энергии) подобны линиям уровня потенциала $ V^{e} $, то есть представляют собой эллипсы с~тем же самым отношением главных осей. 
Необходимо только знать распределение значений потенциала по  этим линиям уровня.

Воспользуемся полной диаграммой для чистого сдвига $ \tau(\gamma) $ ($ \tau $ "--- касательное напряжение, $ \gamma $ "--- деформация сдвига), которая обладает восходящей и нисходящей до нуля ветвями. 
Такую диаграмму материала можно получить, используя опытную диаграмму кручения стержней  круглого поперечного сечения с~пересчетом на диаграмму материала [\citen{bib15}]. 
Тогда энергия, затраченная на деформирование, есть   
$$ V = \int ^{\gamma}_{0}\tau(\gamma)  d\gamma .$$

После взятия интеграла и замены $ \gamma $ на $ \sqrt{2} \rho $  с~учетом выражения \eqref{S_eq7} получаем искомый единый потенциал $ V(\varepsilon_1,\varepsilon_2) $. 

Пусть $ \tau = 0.1G   \gamma(1-5 \gamma)$, если $ 0 \le \gamma \le 0.2 $, где  
$ G = \frac{E}{2(1+\nu)} $ "---  модуль сдвига в упругости. 
Тогда
\begin{equation}\label{S_eq8}
V(\varepsilon_1,\varepsilon_2) = 
\left\lbrace  
\begin{array}{cl}
\frac{V^{e}}{2}\Bigl( 0.1 - \frac{1}{3}\sqrt{\frac{V^{e}}{G}}\Bigr), & (\varepsilon_1,\varepsilon_2) \in \Omega , \\
0.002 G/3 , & (\varepsilon_1,\varepsilon_2) \notin \Omega.
\end{array}
\right. 
\end{equation}


\begin{figure}[b!]
\begin{tabular}{p{0.35\textwidth}p{0.6\textwidth}}

\includegraphics[width=0.3\textwidth]{fig_stru2} & 

\vspace{-5.5cm}

~~~\,Связь между напряжениями и деформациями (определяющие соотношения) в данном случае будет задана законом
$$
\begin{pmatrix} \sigma_1  \\ 
\sigma_2 
\end{pmatrix}
=
\begin{pmatrix} \frac{\partial V}{\partial \varepsilon_1} \\ 
\frac{\partial V}{\partial \varepsilon_2} \end{pmatrix}
=
C \psi (\varepsilon_1,\varepsilon_2) 
\begin{pmatrix} \varepsilon_1 \\ 
\varepsilon_2 \end{pmatrix},
$$
где 
$$ 
\psi (\varepsilon_1,\varepsilon_2) = \frac{1}{2} \Bigl(0.1 - \frac{1}{2} \sqrt{\frac{V^e}{G}} \Bigr).  
$$


~~~\,Отметим, что напряжения являются компонентами градиента функции~$ V $.

\\

\vspace{-.8cm}
\caption{Качественный вид поверхности потенциала $ V(\varepsilon_1,\varepsilon_2) $ } \label{Sfig2}
		\smallskip \footnotesize
		[Figure~\ref{Sfig2}. Qualitative form of the \centerline{potential surface $ V(\varepsilon_1,\varepsilon_2) $]} & \\
\end{tabular}
\end{figure}

%\begin{figure}[b!]
%	\begin{center}
%		\includegraphics[width=0.3\textwidth]{fig_stru2}
%		\caption{Качественный вид поверхности потенциала $ V(\varepsilon_1,\varepsilon_2) $ } \label{Sfig2}
%		\smallskip \footnotesize
%		[Figure~\ref{Sfig2}. Qualitative form of the potential surface $ V(\varepsilon_1,\varepsilon_2) $]
%	\end{center}
%\end{figure}



Очевидно, что область $ \Omega $ ограничена эллипсом с~полуосями 
$ \rho = 0.141$; $a = \rho\sqrt{(1+\nu)/(1-\nu)} $.  
Качественный вид поверхности $ V(\varepsilon_1,\varepsilon_2) $ изображен на рис.~\ref{Sfig2}. 
Функция $ V $ описывает все состояния материала (выпуклость вниз "--- устойчивость, упрочнение; 
выпуклость вверх "--- неустойчивость, разупрочнение).





\smallskip


\Section{Разрушение цилиндрического резервуара} Рассмотрим тонкостенный резервуар в~виде длинной трубы с~заделанными абсолютно жесткими крышками торцами. 
Радиус трубы равен $ b $, длина $ l $, толщина стенки $ t$  $(t\ll b)$. 
Резервуар находится под действием  монотонно и квазистатически возрастающего при постоянной температуре внутреннего давления с~интенсивностью~$ p $ [\citen{bib16, bib17}]. 
Материал обладает свойством деформационного разупрочнения. 
Применим изложенную выше методику для расчета предельного значения давления. 

Каждый элемент трубы находится в~плоском напряженном состоянии при всестороннем растяжении: 
$ \varepsilon_1 = \varepsilon_z $ "--- продольная деформация, 
$ \varepsilon_2 = \varepsilon_{\theta} $ "--- тангенциальная деформация. Его потенциальная функция $ V $ задана выражением \eqref{S_eq8}. 
%где  $ \varepsilon_1 = \varepsilon_z , \varepsilon_1 = \varepsilon_{\theta} $. 

Механическая система (резервуар) обладает двумя параметрами состояния (обобщенными координатами) $ \varepsilon_z$, $\varepsilon_{\theta} $ (однородное напряженно"=деформированное состояние) и~одним параметром управления $ p $. 

Лагранжиан (потенциальная функция) этой системы
$$
W = V(\varepsilon_z , \varepsilon_{\theta})2\pi b l t - p \pi b^2 l \varepsilon_z - 2\pi p b^2 l \varepsilon_{\theta}.
$$
Здесь первое слагаемое "--- полная энергия деформаций стенок трубы, второе и третье "--- взятые со знаком минус соответственно работа сил давления на перемещении жестких крышек и~перемещении, равном увеличению радиуса. 
Тогда уравнения равновесия, определяющие критические точки функции $ W $, имеют вид
\begin{equation}\label{S_eq9}
\frac{\partial W}{\partial \varepsilon_z} = \frac{\partial V}{\partial \varepsilon_z} 2\pi b l t  - p \pi b^2 l = 0;
\quad  
\frac{\partial W}{\partial \varepsilon_{\theta}} = \frac{\partial V}{\partial \varepsilon_{\theta}} 2\pi b l t  - 2p \pi b^2 l = 0.
\end{equation}
Уравнения \eqref{S_eq9} сводятся к одному:
$$
2 \frac{\partial V}{\partial \varepsilon_z} = \frac{\partial V}{\partial \varepsilon_{\theta}}.
$$
Производя необходимые вычисления, находим, что
\begin{equation}\label{S_eq10}
\varepsilon_{\theta} = \frac{2-\nu}{1-2\nu} \varepsilon_{z},
\end{equation}
то есть в пространстве деформаций путь деформирования элемента материала является прямой линией, уравнение которой задается равенством \eqref{S_eq10}. 
Кроме этого, из уравнений \eqref{S_eq9} следует, что  $ \sigma_z = pb/2t$, $\sigma_{\theta} = pb/t $. 
Отсюда путь нагружения в~пространстве напряжений также является прямой линией. 
Таким образом, имеем ситуацию с пропорциональным нагружением элемента материала.

Определим, наконец, предельное значение внутреннего давления. 
Как следует из приведенных выше теоретических положений, для этого необходимо гессиан функции $ W $ приравнять к~нулю, найти все решения полученного уравнения, по крайней мере, приближенно, и выделить из них вырожденные критические точки потенциала $ W $, которые принадлежат множеству критических точек (решений уравнений \eqref{S_eq9}) потенциала. 
Имеем 
$$
\Bigl\lbrace 
\frac{\partial^{2} V}{\partial \varepsilon^{2}_{z}} \cdot \frac{\partial^{2} V}{\partial \varepsilon^{2}_{\theta}} - 
\Bigl( \frac{\partial^{2} V}{\partial \varepsilon_{z}\partial \varepsilon_{\theta}}\Bigr) 
\Bigr\rbrace 
4 \pi^{2} b^2 l^2 t = 0.
$$
После несложных преобразований получаем
\begin{multline}\label{S_eq11}
\bigl(0.05 - \alpha \beta - \alpha \beta^{-1}(\varepsilon_{z} + \nu \varepsilon_{\theta})^2\bigr) 
\bigl(0.05 - \alpha \beta - \alpha \beta^{-1}(\varepsilon_{\theta} + \nu \varepsilon_{z})^2\bigr)-
\\
 - \nu^2 \bigl( 0.05 - \alpha \beta - \alpha \beta^{-1}(\varepsilon_{z} + \nu \varepsilon_{\theta}) (\varepsilon_{\theta} + \nu \varepsilon_{z})\bigr)^2 = 0
,                                          
\end{multline}
где 
$ \alpha=0.25 (1-\nu)^{ - 1/2}$, 
$ \beta = (\varepsilon^{2}_{z} + 2\nu \varepsilon_{z} \varepsilon_{\theta} + \varepsilon^{2}_{\theta})^{1/2}$;  
$ \varepsilon_{z}$, $\varepsilon_{\theta} \ge 0$, $(\varepsilon_{z}, \varepsilon_{\theta}) \in \Omega $.

Для вычисления вырожденных критических точек функции $ W $ применим следующую численную процедуру. 




Построим на плоскости  $ \varepsilon_{z} \varepsilon_{\theta} $ квадрат $ \Omega_1 $  такой, что в нем $ \varepsilon_{z}$, $\varepsilon_{\theta} > 0$,     $(\varepsilon_{z}, \varepsilon_{\theta}) \in \Omega \subset  \Omega_1 $ (рис.~\ref{Sfig3}). 
Для определения стороны этого квадрата найдем на осях $ \varepsilon_{\theta}$ и $\varepsilon_{z} $ равные отрезки, внутри которых детерминант матрицы Гессе меняет знак. 
Получаем отрезки длиной $0.1$. 
Эту величину и берем за сторону квадрата $\Omega_1$. 
Отметим, что данный прием может быть использован при построении $M \times N$-мерного куба в общей теории. 



\begin{figure}[h!]
\centering
	\includegraphics[width=0.4\textwidth]{fig_stru3}
	\caption{Вид кривой, в точках которой гессиан равен нулю} \label{Sfig3} 
	\smallskip \footnotesize
		[Figure~\ref{Sfig3}. The curve where the Hessian is equal to zero at all points]
\end{figure}

Разобьем этот квадрат прямоугольной сеткой с~достаточно малым шагом (на рис. \ref{Sfig3} показана укрупненная сетка). В~каждом узле сетки по формуле \eqref{S_eq11} при  $ \nu = 0.3 $ вычисляем детерминант (гессиан) матрицы Гессе. 
Выделяем те из них, где детерминант мало отличен от нуля, что определяется заданной точностью вычислений.

Производя расчеты, получаем множество точек, расположенных на кривой {\sl 1} (рис. \ref{Sfig3}). 
Затем координаты точек кривой {\sl 1} подставляем в уравнения \eqref{S_eq9} и~выделяем те из них, которые удовлетворяют уравнениям с~заданной точностью. 
В~результате находим только одну точку (точка $A$, рис. \ref{Sfig3}). Точка $ A $ лежит на прямой \eqref{S_eq10} и~имеет координаты $ \varepsilon_{z}  = 0.018$, $\varepsilon_{\theta}  = 0.076 $ (с~точностью до округления). 
Это и~есть искомая вырожденная критическая точка потенциальной функции $ W $. 
Подставляя ее координаты в~уравнения равновесия, находим критическое внутреннее давление, при котором происходит потеря устойчивости системы (катастрофа, разрушение). 
Если  $ E = 2 \cdot 10^5 $~МПа, $ b = 100 $ мм,  $ t = 2 $ мм, то критическое (предельное) значение давления $ p = 18 $~МПа. Заметим, что величина предельного давления увеличивается пропорционально толщине стенки.


\pagebreak



\AdditionalContent{Конкурирующие интересы}{Мы не имеем конкурирующих интересов.} 
%\AdditionalContent{Конкурирующие интересы}{Конкурирующих интересов не имею.} 
\AdditionalContent{Авторский вклад и ответственность}{Все авторы принимали участие в разработке концепции статьи и в написании рукописи. Все авторы несут полную ответственность за представление окончательной рукописи в печать. Окончательная версия рукописи была одобрена всеми авторами.} 
%\AdditionalContent{Авторская ответственность}{Я несу полную ответственность за предоставление окончательной версии рукописи в печать. Окончательная версия рукописи мною одобрена.} 
%\AdditionalContent{Авторский вклад и ответственность}{Все  авторы  принимали  участие  в  разработке  концепции  \mbox{статьи}  и  в  написании  рукописи. Авторы несут  полную  ответственность  за  предоставление  окончательной рукописи в печать. Окончательная  версия рукописи была одобрена всеми авторами.} 
\AdditionalContent{Финансирование}{Исследование не имело финансирования.}
%\AdditionalContent{Финансирование}{Исследование выполнялось без финансирования.}
%\AdditionalContent{Финансирование}{Работа выполнена при поддержке РФФИ (проект № xx–xx–xxxxx_a).}
%\AdditionalContent{Благодарность}{Выразите свою благодарность.}
%\AdditionalContent{Благодарность}{Авторы благодарны рецензенту за тщательное прочтение статьи и~ценные предложения и комментарии.}

\begin{thebibliography}{10}

\RBibitem{bib1}
\by Седов~Л.~И.
\book Механика сплошной среды
\bookvol 1
\yr 1970
\publ Наука
\publaddr М.
\totalpages 492
%\misctransl
%\by Sedov ~L.~I.
%\book Mekhanika sploshnoy sredy \rm [Continuum Mechanics] vol. 1 
%\yr 1970
%\publ Nauka
%\publaddr Moscow
%\lang In Russian
%\totalpages 492

\RBibitem{bib2}
\by Стружанов~В.~В., Миронов~В.~И.
\book Деформационное разупрочнение материала в элементах конструкций
\yr 1995
\publ УрО РАН 
\publaddr Екатеринбург
\totalpages 190
%\misctransl
%\by Struzhanov ~V.~V., Mironov ~V.~I.
%\book Deformatsionnoye razuprochneniye materiala v elementakh konstruktsiy \rm [Deformational Weakening of the Material in Structural Elements]
%\yr 1995
%\publ UrO RAN
%\publaddr Ekaterinburg
%\lang In Russian
%\totalpages 190

\RBibitem{bib3}
\by Вильдеман~В.~Э., Чаусов~Н.~Г.
\paper Условия деформационного разупрочнения материала при растяжении образца специальной конфигурации
\jour Заводская лаборатория. Диагностика материалов
\yr 2007
\vol 73
\issue 10
\pages 55--59
\elib{https://elibrary.ru/item.asp?id=9907854}
%\misctransl
%\lang In Russian
%\by Vildeman ~V.~E., Chausov ~N.~G.
%\paper Usloviya deformatsionnogo razuprochneniya materiala pri rastyazhenii obraztsa spetsial'noy konfiguratsii \rm [Conditions of Deformational Weakening of the Material with the Specimen of a Special Configuration under Tension]
%\jour Zavodskaya laboratoriya. Diagnostika materialov.
%\yr 2007
%\vol 73
%\issue 10
%\pages 55--59

\RBibitem{bib4}
\by Вильдеман~В.~Э., Третьяков~М.~П.
\paper Испытания материалов с построением полных диаграмм деформирования
\jour Проблемы машиностроения и надежности машин
\yr 2013
\vol 
\issue 2
\pages 93--98
\elib{https://elibrary.ru/item.asp?id=19548776}
%\misctransl
%\lang In Russian
%\by Vildeman ~V.~E., Tretyakov ~M.~P.
%\paper Ispytaniya materialov s postroyeniyem polnykh diagramm deformirovaniya \rm [Material Tests and Development of Complete Deformation Diagrams]
%\jour Problemy mashinostroyeniya i nadezhnosti mashin
%\yr 2013
%\vol 
%\issue 2
%\pages 93--98

\RBibitem{bib5}
\by Стружанов~В.~В., Коркин~А.~В.
\paper Об устойчивости процесса растяжения одной стержневой системы с~разупрочняющимися элементами
\jour Вестник Уральского государственного университета путей сообщения
\yr 2016
\issue 3(31)
\pages 4--17
\elib{https://elibrary.ru/item.asp?id=27249016}
\crossref{10.20291/2079-0392-2016-3-4-17}
%\misctransl
%\lang In Russian
%\by Struzhanov ~V.~V., Korkin ~A.~V.
%\paper Ob ustoychivosti protsessa rastyazheniya odnoy sterzhnevoy sistemy s razuprochnyayushchimisya elementami \rm [On the Stability of a Single Rod System with Weakening Elements under Tension]
%\jour Vestnik Ural'skogo gos. un–ta putey soobshcheniya
%\yr 2016
%\vol 
%\issue 3(31)
%\pages 4--17

\RBibitem{bib6}
\by Андреева~Е.~А.
\paper Решение одномерных задач пластичности для разупрочняющегося материала
\jour Вестн. Сам. гос. техн. ун-та. Сер. Физ.-мат. науки
\yr 2008
\issue 2(17)
\pages 152--160
\mathnet{http://mi.mathnet.ru/vsgtu642}
\crossref{10.14498/vsgtu642}


%\RBibitem{bib7}
%\by Горбунов~С.~В., Радченко~В.~П.
%\paper Вариант решения краевой задачи неупругого деформирования разупрочняющейся среды и его экспериментальная проверка
%\jour Материалы Х Всероссийской конференции по механике деформируемого твёрдого тела (18–22 сентября 2017 г., Самара, Россия)
%\yr 2017
%\vol 1
%\issue
%\pages 179--181
%\misctransl
%\lang In Russian
%\by Gorbunov ~S.~V., Radchenko ~V.~P.
%\paper Variant resheniya krayevoy zadachi neuprugogo deformirovaniya razuprochnyayushcheysya sredy i yego eksperimental'naya proverka \rm [A Solution Case of the Boundary-Value Problem for Inelastic Deformation of a Weakening Medium and its Experimental Verification]
%\jour Materialy X  Vserossiyskoy konferentsii po mekhanike deformiruyemogo tvordogo tela.  
%\yr 2017
%\vol 1
%\issue
%\pages 179--181

\RBibitem{bib8}
\by Арсенин~В.~Я.
\book Методы математической физики и специальные функции
\yr 1974
\publ Наука
\publaddr М.
\totalpages 431
%\misctransl
%\by Arsenin ~V. ~Ya.
%\book Metody matematicheskoy fiziki i spetsial'nyye funktsii. \rm [Methods of Mathematical Physics and Special Functions]
%\yr 1974
%\publ Nauka
%\publaddr Moscow
%\lang In Russian
%\totalpages 431

\RBibitem{bib9}
\by Стружанов~В.~В., Коркин~А.~В.
\paper О решении нелинейных уравнений равновесия  одной растягиваемой стержневой системы с разупрочняющимися элементами методом простых итераций
\jour Вестник Уральского государственного университета путей сообщения
\yr 2017
\issue 2(34)
\pages 4--16
\crossref{10.20291/2079-0392-2017-2-4-16 }
\elib{https://elibrary.ru/item.asp?id=29409142}
%
%\misctransl
%\lang In Russian
%\by Struzhanov ~V.~V., Korkin ~A.~V.
%\paper O reshenii nelineynykh uravneniy ravnovesiya odnoy rastyagivayemoy sterzhnevoy sistemy s razuprochnyayushchimisya elementami metodom prostykh iteratsiy \rm [The solution of Nonlinear Equilibrium Equations for a Single Rod System with Weakening Elements under Tension Using the Simple Iterations Method]
%\jour Vestnik Ural'skogo gos. un-ta putey soobshcheniya
%\yr 2017
%\vol 
%\issue 2(34)
%\pages 4--16

\Bibitem{bib10}
\by Poston~T., Stewart~I.
\book Catastrophe Theory and Its Applications
\yr 1978
\publ Pitman
\serial Surveys and Reference Works in Mathematics
\vol 2
\zmath{0382.58006}
\publaddr London, San Francisco, Melbourne
\totalpages xviii+491 

\Bibitem{bib11}
\by Gilmore~R.
\book Catastrophe theory for scientists and engineers
\zmath{0497.58001}
\serial A Wiley-Interscience Publication
\publaddr New York 
\publ John Wiley \& Sons
\totalpages xvii+666 
\yr 1981

\Bibitem{bib12}
\by Pars~L.~A.
\book A treatise on analytical dynamics
\zmath{0125.12004}
\publaddr London
\publ Heinemann Educational Books 
\totalpages xxi+641 
\yr 1965
 
\RBibitem{bib13}
\by Ильюшин~А.~А.
\paper Об основах общей математической теории пластичности
\jour Вестник Московского университета. Серия 1. Математика. Механика
\yr 1961
\issue 3
\pages 31--36
\zmath{0100.20304}
%\misctransl
%\lang In Russian
%\by Ilyushin ~A.~A.
%\paper Ob osnovakh obshchey matematicheskoy teorii plastichnosti \rm [Fundamentals of the General Mathematical Theory of Plasticity]
%\jour Voprosy teorii plastichnosti. Izd-vo Akademii nauk SSSR
%\yr 1961
%\vol 
%\issue 
%\pages 3--29

\RBibitem{bib14}
\by Качанов~Л.~М.
\book Основы теории пластичности
\yr 1969
\publ Наука
\publaddr М.
\totalpages 420
%\misctransl
%\by Kachanov ~L.~M.
%\book Osnovy teorii plastichnosti \rm [Fundamentals of the Theory of Plasticity]
%\yr 1969
%\publ Nauka
%\publaddr Moscow
%\lang In Russian
%\totalpages 420

\RBibitem{bib15}
\by Стружанов~В.~В.
\paper Определение диаграммы деформирования материала с~падающей ветвью по диаграмме кручения цилиндрического образца
\jour Сиб. журн. индустр. матем.
\yr 2012
\vol 15
\issue 1
\pages 138--144
\mathnet{http://mi.mathnet.ru/sjim718}

\RBibitem{bib16}
\by Радченко~В.~П.
\book Введение в механику деформируемых систем
\yr 2009
\publ СамГТУ
\publaddr Самара
\totalpages 241
%\misctransl
%\by Radchenko ~V.~P.
%\book Vvedeniye v mekhaniku deformiruyemykh system \rm [Introduction to Mechanics of De-formation Systems]
%\yr 2009
%\publ Samarskiy gos. tekhn. un-t.
%\publaddr Samara
%\lang In Russian
%\totalpages 241

\Bibitem{bib17}
\by Southwell~R.~V.
\book An introduction to the theory of elasticity. For engineers and physicists
\zmath{62.1529.04}
\totalpages ix+510 
\publaddr Oxford
\publ Clarendon Press
\serial Oxford Engineering Science Series
\yr 1936


\end{thebibliography}


\newpage

\makeentitle

\vspace{-5mm}

\AdditionalContent{Competing interests}{We have no competing interests.} 
%\AdditionalContent{Competing interests}{I have ... (or) no competing interests.} 
%\AdditionalContent{Competing interests}{We have ... (or) no competing interests.} 
\AdditionalContent{Authors' contributions and responsibilities}{Each authors has participated in the article concept development  and in the manuscript writing. The authors are absolutely responsible for submitting the final manuscript in print. Each authors has approved the final version of manuscript.}
%\AdditionalContent{Author's Responsibilities}{I take full responsibility for submitting the final manuscript in print. I approved the final version of the manuscript.} 
%\AdditionalContent{Authors' contributions and responsibilities}{Each author has participated in the article concept development and in the manuscript writing. The authors are absolutely responsible for submitting the final manuscript in print. Each author has approved the final version of manuscript.} 

\AdditionalContent{Funding}{This research received no specific grant from any funding agency in the public, commercial, or not-for-profit sectors.}




\begin{thebibliography}{10}

\Bibitem{bib1:1}
\lang In Russian
\by Sedov~L.~I.
\book Mekhanika sploshnoi sredy \rm [Continuum Mechanics] 
\bookvol 1
\yr 1970
\publ Nauka
\publaddr Moscow
\totalpages 492

\Bibitem{bib2:1}
\lang In Russian
\by Struzhanov~V.~V., Mironov~V.~I.
\book Deformatsionnoe razuprochnenie materiala v elementakh konstruktsii \rm [Deformational Softening of Material in Structural Elements]
\yr 1995
\publ UrO RAN 
\publaddr Ekaterinburg
\totalpages 190

\Bibitem{bib3}
\lang In Russian 
\by Vil'deman~V.~E., Chausov~N.~G.
\paper Conditions of strain softening upon stretching of the specimen of special configuration
\jour Zavodskaia laboratoriia. Diagnostika materialov
\yr 2007
\vol 73
\issue 10
\pages 55--59


\Bibitem{bib4:en}
\paper Material testing  by plotting total deformation curves
\by Vil'deman~V.~E., Tretyakov~M.~P.
\jour J. Mach. Manuf. Reliab.
\yr  2013
\vol 42
\issue 2
\pages 166--170
\elib{https://elibrary.ru/item.asp?id=20435521}
\crossref{10.3103/S1052618813010159}



\Bibitem{bib5:1}
\lang In Russian
\by Struzhanov~V.~V., Korkin~A.~V.
\paper Regarding stretching process stability of one bar system with softening elements
\jour Vestnik Ural'skogo gosudarstvennogo universiteta putei soobshcheniia
\yr 2016
\issue 3(31)
\pages 4--17
\crossref{10.20291/2079-0392-2016-3-4-17}


\Bibitem{bib6}
\lang In Russian
\by Andreeva~E.~A.
\paper Solution of one-dimensional softening materials plasticity problems
\jour Vestn. Samar. Gos. Tekhn. Univ., Ser. Fiz.-Mat. Nauki \rm [J. Samara State Tech. Univ., Ser. Phys. Math. Sci.]
\yr 2008
\issue 2(17)
\pages 152--160
\crossref{10.14498/vsgtu642}



\Bibitem{bib8:1}
\lang In Russian
\by Arsenin~V.~Ya.
\book Metody matematicheskoi fiziki i spetsial'nye funktsii \rm [Methods of mathematical physics and special functions]
\yr 1974
\publ Nauka
\publaddr Moscow
\totalpages 431


\Bibitem{bib9:1}
\lang In Russian
\by Struzhanov~V.~V., Korkin~A.~V.
\paper On the solution of non-linear equations of trim of a single expandable frame structure with softening elements by the method of simple iterations
\jour Vestnik Ural'skogo gosudarstvennogo universiteta putei soobshcheniia
\yr 2017
\issue 2(34)
\pages 4--16
\crossref{10.20291/2079-0392-2017-2-4-16}


\Bibitem{bib10:1}
\by Poston~T., Stewart~I.
\book Catastrophe Theory and Its Applications
\yr 1978
\publ Pitman
\serial Surveys and Reference Works in Mathematics
\vol 2
\zmath{0382.58006}
\publaddr London, San Francisco, Melbourne
\totalpages xviii+491 

\Bibitem{bib11:1}
\by Gilmore~R.
\book Catastrophe theory for scientists and engineers
\zmath{0497.58001}
\serial A Wiley-Interscience Publication
\publaddr New York 
\publ John Wiley \& Sons
\totalpages xvii+666 
\yr 1981

\Bibitem{bib12:1}
\by Pars~L.~A.
\book A treatise on analytical dynamics
\zmath{0125.12004}
\publaddr London
\publ Heinemann Educational Books 
\totalpages xxi+641 
\yr 1965
 
\Bibitem{bib13:1}
\lang In Russian
\by Ilyushin~A.~A.
\paper On the foundations of the general mathematical theory of plasticity
\jour Vestn. Mosk. Univ., Ser.~I 
\yr 1961
\issue 3
\pages 31--36



\Bibitem{bib14:1}
\by Kachanov~L.~M.
\book Foundations of the theory of plasticity
\zmath{0231.73015}
\serial North-Holland Series in Applied Mathematics and Mechanics
\vol 12
\publaddr Amsterdam
\publ North-Holland Publ.
\totalpages xiii+482 
\yr 1971


\Bibitem{bib15:!}
\by Struzhanov~V.~V.
\paper The determination of the deformation diagram of a~material with a~falling branch using the torsion diagram of a~cylindrical sample
\jour Sib. Zh. Ind. Mat.
\yr 2012
\vol 15
\issue 1
\pages 138--144

\Bibitem{bib16:1}
\lang In Russian
\by Radchenko~V.~P.
\book Vvedenie v mekhaniku deformiruemykh sistem \rm [Introduction to the mechanics of deformable systems]
\yr 2009
\publ Samara State Technical Univ.
\publaddr Samara
\totalpages 241


\Bibitem{bib17:1}
\by Southwell~R.~V.
\book An introduction to the theory of elasticity. For engineers and physicists
\zmath{62.1529.04}
\totalpages ix+510 
\publaddr Oxford
\publ Clarendon Press
\serial Oxford Engineering Science Series
\yr 1936


\end{thebibliography}


%\end{document}
 \newpage

%\input{preamble}

%

\vsgtu{1639} %registration number

\FirstPage{774}

\LastPage{784}

%

%\usepackage{samgtu-name}

%

%\pdfcrop

%\draft

%

\ShortCommunication

%

%\begin{document}





$\;$

\vspace{1mm}



\hfill \qrcode[hyperlink,height=.8in]{https://doi.org/10.14498/vsgtu1639}

\vspace{-.8in}



\udc{517.955, 517.968.7}

\msc{26A33, 35K15, 35R11}



\rutitle{К проблеме единственности решения задачи Коши для~уравнения дробной диффузии с оператором Бесселя}

\rutitleshort{К проблеме единственности решения задачи Коши\dots}

\rutitlehack{К проблеме единственности решения \\ задачи Коши для~уравнения дробной диффузии с~оператором Бесселя}



\entitle{On the uniqueness of the solution of the Cauchy problem for the equation of fractional diffusion with~Bessel~operator}

\entitleshort{On the uniqueness of the solution of the Cauchy problem\dots}



\ruauthor{Ф.~Г.~Хуштова}{Х\,у\,ш\,т\,о\,в\,а~Ф.~Г.}

\enauthor{F.~G.~Khushtova}{K\,h\,u\,s\,h\,t\,o\,v\,a~F.~G.}



\ruaffil{\runiipma.}



\enaffil{\enniipma.}



\ruauthorsdetails{1}{% Количество авторов

\fullauthor{\rupid{53181}{Фатима Гидовна Хуштова}}

\CAuthor % Автор-корреспондент

\orcid{0000-0003-4088-3621} % ORCID ID автора 

%\regalia{без ученой степени} 

\position[n]{научный сотрудник} 

\dept{отдел дробного исчисления} 

\email{khushtova@yandex.ru}

}



\enauthorsdetails{1}{

\fullauthor{\enpid{53181}{Fatima G. Khushtova}}

\CAuthor % Сorresponding Author

\orcid{0000-0003-4088-3621} % ORCID ID

%\regalia{Cand. Phys. \& Math. Sci., Associate Professor} 

\position[n]{Researcher} 

\dept{Dept. of Fractional Calculus}

\email{khushtova@yandex.ru}

}



\rukeywords{уравнение дробной диффузии, оператор дробного дифференцирования, оператор Бесселя, задача Коши, единственность решения, условие Тихонова, функция Райта}

\enkeywords{fractional diffusion equation, fractional differentiation operator, Bessel operator, Cauchy problem, solution uniqueness, Tikhonov condition,	 Wright function}



\ruabstract{Рассматривается уравнение дробной диффузии с сингулярным оператором Бесселя, действующим по пространственной переменной, и оператором дробного дифференцирования Римана--Лиувилля, действующим по временной переменной. 

Когда порядок дробной производной равен единице, а особенность у оператора Бесселя отсутствует, рассматриваемое уравнение совпадает с~классическим уравнением теплопроводности. 

Ранее для уравнения дробной диффузии с оператором Бесселя было построено решение задачи Коши и доказана теорема единственности решения в классе функций экспоненциального роста. 



Построен пример, показывающий, что увеличение показателя степени в условии, гарантирующем единственность решения задачи Коши, влечет за собой неединственность решения. 

С помощью известных свойств функции Райта получены оценки для построенной функции.

Показывается, что она, будучи не равной тождественно нулю, удовлетворяет однородному уравнению и однородному условию Коши. 

}





\enabstract{In this paper, we consider fractional diffusion equation involving the Bessel operator acting with respect to a spatial variable and the Riemann-Liouville fractional differentiation operator acting with respect to a time variable. When the order of the fractional derivative is unity, and the singularity of the Bessel operator is absent, this equation coincides with the classical heat equation. Earlier, a solution of the Cauchy problem has been considered for the considered equation and a uniqueness theorem has been proved for a class of functions satisfying the analog of the Tikhonov condition. 



In this paper, we have constructed an example to show that the exponent (power) at the condition of the uniqueness of the solution to the Cauchy problem cannot be raised under. Its increase leads to a non-uniqueness of the solution. Using the well-known properties of the Wright function, we have obtained estimates for constructed function, which satisfies the homogeneous equation and the zero Cauchy condition.

}



\newdate{khusha}{28}{08}{2018}

\date{\displaydate{khusha}}

\newdate{khushb}{25}{10}{2018}

\revisiondate{\displaydate{khushb}}

\newdate{khushc}{12}{11}{2018}

\accepteddate{\displaydate{khushc}}



\newdate{khushe}{28}{11}{2018}

\renewcommand{\JournalDataOnline}{\displaydate{khushe}}



%\ForPeerReview

\makerutitle



\newpage



\Section[n]{Введение}

В области $\Omega=\{(x,y): -\infty<x<\infty,\ 0<y<T \}$ рассмотрим уравнение

\begin{equation}\label{GL3_1}

B_xu(x,y)-D^{\alpha}_{0y}u(x,y)=0, 

\end{equation}

где $D_{0y}^{\alpha}$ "--- оператор дробного дифференцирования в смысле  

Римана"--~Лиувилля  порядка $\alpha,$ определяемый равенствами~\cite{Drob, Pskhu}   

$$D_{0y}^{\alpha} \, g(y)= \frac{dg}{dy},$$ если $\alpha=1$, и

$$ 

D_{0y}^{\alpha}\,g(y)=\frac{1}{\Gamma(1-\alpha)}\frac{d}{dy}

\int _{0}^{y}\frac{g(t)}{(y-t)^{\alpha}}dt,

$$ 

если $0<\alpha<1;$

$$

B_x=x^{-b} \frac{d}{dx}\left(x^b \frac{d}{dx}\right)=\frac{d^2}{d x^2}+\frac{b}{x} \frac{d}{d x}

$$

"--- оператор Бесселя, $|b|<1.$



Уравнение \eqref{GL3_1} при $\alpha=1,$ то есть уравнение

\begin{equation}\label{UrAlpha-1}

u_{xx}(x,y)+\frac{b}{x}u_{x}(x,y)-u_{y}(x,y)=0,

\end{equation}

совпадает с %~названным И.~А.~Киприяновым 

$B$-параболическим уравнением.

Краевые задачи в ограниченной и~неограниченной областях для этого уравнения при различных значениях параметра $b$ рассматривали многие авторы. 

Например, в работе V.~Alexia\-des~\cite{Alexiades} для него решены первая, вторая и~третья краевые задачи в~области с~подвижной границей, в~работе D.~Calton\cite{Calton} для него  исследована задача Коши.

 

 Уравнение \eqref{UrAlpha-1} заменой $z=x^2/(4n)$ сводится к~уравнению

\begin{equation} \label{UrKep2}

u_{zz}(z,y)+\frac{m+1}{z}\,u_{z}(z,y)-\frac{n}{z} u_{y}(z,y)=0,\quad m=(b-1)/2,

\end{equation}

для которого S.~K\c{e}pi\'nski  \cite{Kepinski_2} исследовал краевую задачу в~полуполосе $x>0,$ $0<y<y_0.$

Уравнение \eqref{UrAlpha-1} было также объектом исследования  монографий С.~А.~Терсенова~\cite{Tersenov} и~М.~И.~Матийчука~\cite{Matiychuk}.



Дифференциальные уравнения, содержащие оператор Бесселя, наиболее подробно исследованы в~работах И.~А.~Киприянова и~его учеников \cite{Kipriyanov-Monographiya,  Kipriyanov}.

Отметим здесь, что очень важные результаты по применению операторов преобразования в~теории дифференциальных уравнений с~особенностями в~коэффициентах, в~частности с~операторами Бесселя, получены В.~В.~Катраховым и~С.~М.~Ситником~\cite{KatrakhovSitnik, Sitnik}.





В случае, когда $b=0,$ $0<\alpha<2,$ уравнение \eqref{GL3_1} совпадает с диффузионно"=волновым уравнением. 

Различные краевые задачи для него, а также для многомерных его обобщений подробно исследованы в работах многих авторов.  Приведем некоторые из них. 



А.~Н.~Кочубей \cite{Kochubei_2} исследовал задачу Коши для уравнения диффузии дробного порядка с~регуляризованной дробной производной (производной Капуто) и~эллиптическим оператором с~непрерывными ограниченными вещественными коэффициентами

$$

L=\sum\limits_{i,j=1}^{n}a_{i,j}(x)\frac{\partial^2}{\partial x_i \partial x_j}+

\sum\limits_{j=1}^{n}b_{j}(x)\frac{\partial}{\partial x_j}+c(x) u.

$$

В случае, когда $L=\Delta$ "--- оператор Лапласа, фундаментальное решение построено в терминах $H$-функции Фокса. 

Исследованию задачи Коши для диффузионного и диффузионно"=волнового уравнений дробного порядка с~производной Капуто  посвящены также  работы А.~Н.~Кочубея и С.~Д.~Эйдельмана \cite{Kochubei_1, Kochubei_3, Kochubei_4, Kochubei_5}.



А.~В.~Псху~\cite{Pskhu2} построил фундаментальное решение и исследовал задачу Коши  для многомерного диффузионно"=волнового уравнения с~оператором дробного дифференцирования Джрбашяна"--~Нерсесяна, который  включает в~себя как частный случай операторы дробного дифференцирования Римана"--~Лиувилля и Капуто.



В~работах  \cite{Pskhu-SibirJournal, Pskhu3} также построено фундаментальное решение и исследована задача Коши для уравнения дробной диффузии соответственно с~оператором дискретно распределенного дифференцирования и оператором Джрбашяна"--~Нерсесяна, действующим по многим переменным времени.

В работе \cite{Pskhu4} решена первая краевая задача в~нецилиндрической области для диффузионно"=волнового уравнения с~оператором Джрбашяна"--~Нерсесяна. 

Доказана теорема существования и~единственности исследуемой задачи, конструктивно построено ее решение. 



А.~А.~Ворошилов и А.~А.~Килбас \cite{VorKilbas} методом интегральных преобразований построили решение задачи Коши для уравнения

\begin{equation} \label{MnogPerem}

D^{\alpha}_{0t}\,u(x,t)=\lambda^2 \Delta_x u(x,t), \quad  x\in{\mathbb R}^m,\quad  t>0, 

\end{equation}

где $\Delta_x=\sum_{j=1}^{m}\partial^2/\partial x_{j}^{2},$ $n-1<\alpha <n,$ $n\in\mathbb{N}.$ 

При $0<\alpha\leqslant 1$ и $1<\alpha<2$ решения  выписаны в терминах $H$-функции Фокса.

В работе~\cite{VorKilbas2} было также  построено решение задачи Коши в~случае, когда в уравнении

\eqref{MnogPerem} вместо оператора Римана"--~Лиувилля содержится оператор Капуто. 



Задача Коши для уравнения \eqref{MnogPerem} при $\lambda=1,$ 

$0<\alpha\leqslant 1$  была ранее рассмотрена С.~Х.~Геккиевой~\cite{Gekk}, а~в~работе \cite{Gekkieva1} ею исследована первая краевая задача в~первом квадранте для уравнения \eqref{MnogPerem} при $m=\lambda=1,$ 

$0<\alpha\leqslant 1.$  







Уравнение диффузии дробного порядка и некоторые его обобщения рассматривали 

F.~Mainardi, G.~Pagnini, Yu.~Luchko, P.~Paradisi~\cite{Mainardi-1, Mainardi-3, Mainardi-2, Pagnini_1,   Pagnini_Paradisi}  и~многие другие. 

Более подробную библиографию по данному вопросу можно найти в~монографиях \cite{Pskhu, MathaiSaxenaHaubold, Podlubny}.

 

%Чтобы подготовить Вашу статью для рецензента, т.~е. обезличить её, воспользуйтесь командой \verb"\ForPeerReview" перед командой 

%\verb"\makerutitle".



\smallskip





\Section{Постановка задачи Коши и теорема единственности решения}

Пусть $\Omega^{+}=\Omega\cap\{x>0 \},$ $\Omega^{-}=\Omega\cap\{x<0\},$ 

$\bar{\Omega}$ "--- замыкание области $\Omega.$ 



\smallskip



{\sc\small Определение.}

{\it Регулярным  решением} уравнения \eqref{GL3_1} в области $\Omega$ 

будем называть функцию

$u=u(x,y),$ удовлетворяющую уравнению \eqref{GL3_1} в области

$\Omega^{+}\cup\Omega^{-},$ и такую, что $y^{\,1-\alpha}u\in{C(\bar{\Omega})},$ 

$|x|^b u_{x}\in{C(\Omega)},$ $u_{xx},$ $D_{\,0y}^{\,\alpha} u\in{C(\Omega^{+}\cup\Omega^{-})}.$



\smallskip



{\sc\small Задача Коши.}

{\it Найти регулярное в области $\Omega$ решение

уравнения} \eqref{GL3_1}, {\it удовлетворяющее условию

\begin{equation}\label{UslSMCoshi}

\lim\limits_{y \rightarrow 0} D_{0y}^{\alpha-1}u(x,y)=\varphi(x), \quad  -\infty<x<\infty,

\end{equation}

где $\varphi(x)$ "--- заданная функция}.



\smallskip



В работе \cite{KhushtovaVestn2} доказана следующая теорема единственности.



\smallskip



\begin{newthm}{Теорема}  

Существует не более одного регулярного решения задачи Коши {\rm \eqref{GL3_1}, \eqref{UslSMCoshi} } в классе функций$,$ удовлетворяющих для некоторой положительной постоянной $k$ условию

\begin{equation} \label{AnalogTikhonovaInfty}

\lim\limits_{|x|\to \infty}y^{1-\alpha}\,u(x,y)\exp{\big(-k\,|x|^{\,\frac{2}{2-\alpha}} \big)}=0.

\end{equation}  

\end{newthm}



\smallskip





Отметим, что условие \eqref{AnalogTikhonovaInfty} имеет место и в случае задачи Коши для уравнения дробной диффузии.  

Теоремы единственности решения задачи Коши в~классе ограниченных функций  и~в~классе функций экспоненциального роста для многомерного уравнения дробной диффузии с~оператором Капуто были доказаны в работе А.~Н.~Кочубея \cite{Kochubei_2},

теорема единственности решения задачи Коши в классе функций экспоненциального роста для многомерного диффузионно"=волнового уравнения с~оператором Джрбашяна"--~Нерсесяна  "--- в~работе А.~В.~Псху~\cite{Pskhu2}. 

В этих же работах построены примеры, показывающие, что увеличение показателя $2/(2-\alpha)$ ведет к~утрате единственности решения соответствующих задач. 

При $\alpha=1$ условие \eqref{AnalogTikhonovaInfty} совпадает с~хорошо известным условием Тихонова 

$$

\lim\limits_{|x|\to \infty}u(x,y)\exp{\big(-k\,|x|^{2} \big)}=0.

$$

Как известно, пример неединственности решения задачи Коши для уравнения теплопроводности был впервые построен А.~Н.~Тихоновым в~основополагающей работе~\cite{tikhonovUsl-1} (см.~также \cite{tikhonovUsl-2}).



В данной работе эти идеи развиваются применительно к~задаче \eqref{GL3_1}, \eqref{UslSMCoshi}.



\smallskip



\Section{Неулучшаемость показателя степени в условии единственности решения задачи Коши}

Рассмотрим функцию

\begin{equation} \label{T(z,zeta)}  

T_{\mu}(z,\zeta)=\sum\limits_{n=0}^{\infty}\frac{z^{\,\beta+n}}{\,n!\,\Gamma(1+\beta+n)}\,\phi\left(-\rho,\mu-\alpha\beta-\alpha n; -\zeta\right),

\end{equation}    

где  $\phi\,(-\rho,\delta;z)$ "--- {\it функция Райта}, определяемая степенным рядом \cite{Pskhu, GorenfloLuchkoMainardi}

$$

\phi\,(-\rho,\delta;z)=\sum\limits_{n=0}^{\infty}\frac{z^{n}}{\,n!\,\Gamma(\delta-\rho\, n)}.

$$



Далее будем считать, что $\rho\in(0;1).$ 

Для любых $\nu$, $\delta\in \mathbb{R}$ справедливы формулы  \cite[с.~26, 44]{Pskhu}

\begin{gather}\label{diff_x}

\frac{d}{dz}\phi(-\rho,\delta;-z)=-\phi(-\rho,\delta-\rho;-z),

\\

\label{diff_Right}

D_{ay}^{\nu} \,y^{\delta-1}\phi(-\rho,\delta;-c\,y^{-\rho})=y^{\delta-\nu-1}\phi(-\rho,\delta-\nu;-c\,y^{-\rho}).

\end{gather}



\smallskip



\begin{lemma}[1]Если $\delta\geqslant 1,$ то для любого положительного $z$ справедливо нера\-вен\-ство~{\rm \cite[с.~47]{Pskhu} }

\begin{equation}\label{Right_delta>1}

\phi\,(-\rho,\delta;-z)\leqslant \frac{1}{\Gamma\left( \delta \right)}

\exp\left[-(1-\rho)\,\rho^{\frac{\rho}{1-\rho}}\, z^{\frac{1}{1-\rho}} \right].

\end{equation}

\end{lemma}



\smallskip



\begin{lemma}[2]Если $\delta<1,$ то для любых положительных $x$ и $y,$ 

$a\in(0;x),$ $\omega\in(1/2; \omega_0),$ $\omega_0=\min\{1,1/(2\rho)\},$

справедливы неравенства~{\rm \cite[с.~49]{Pskhu}}

\begin{gather}\label{Right-0}

\Bigl|\,y^{\,\delta-1}\phi\Bigl(-\rho,\delta;-\frac{x}{y^{\rho}}\Bigr)\Bigr|

\leqslant 

\frac{\Gamma\left( (1-\delta)/\rho\right)}{\,\rho\, \pi \left[aC_{\rho}(\rho,\omega) \right]^{(1-\delta)/\rho}}\,

\phi\Bigl(-\rho,1;-\frac{x-a}{y^{\rho}}\Bigr),

\\

\label{Right-2}

\Bigl|y^{\,\delta-1}\phi\Bigl(-\rho,\delta;-\frac{x}{y^{\rho}}\Bigr)\Bigr|

\leqslant \frac{\Gamma\left( 1-\delta\right)}{\pi \left[yC_{\rho}(1,\omega) \right]^{1-\delta}},

\end{gather}

где 

$C_{\rho}(\rho,\omega)=\frac1{\rho\pi}{\cos \rho\omega\pi},$ $C_{\rho}(1,\omega)=-\frac1{\pi}{\cos \omega\pi}.

$

\end{lemma}



\smallskip



Как следует из \eqref{Right-2}, ряд \eqref{T(z,zeta)} сходится абсолютно для любых  $z\in\mathbb{C},$ $\zeta>0,$ $\rho, \alpha, \beta\in(0;1),$ $\mu\in\mathbb{R}.$



\smallskip



Докажем следующую лемму.



\smallskip

\begin{lemma}[3]Если $0<\alpha<\rho<1,$ $\mu\geqslant 0,$ то при любых $x\in\mathbb{R},$ $y>0$ справедливо неравенство

\begin{equation} \label{T_muOtsenka}  

\Bigl| y^{\mu-1} T_{\mu}\Bigl(\frac{x^2}{4y^{\alpha}},\frac{1}{y^{\rho}}\Bigr)\Bigr|

\leqslant 

\frac{C

|x|^{\frac{2[1-\rho(1-\beta)]}{2\rho-\alpha}} y^{\mu}}{\Gamma(\mu+1)}\exp\left[ k_1  |x|^{\frac{2\rho}{2\rho-\alpha}}-k_2   y^{-\frac{\rho}{1-\rho}}\right],

\end{equation}    

где $C,$ $k_1,$ $k_2$ "--- положительные постоянные$,$ зависящие только от 

$\alpha,$ $\beta$ и $\rho.$

\end{lemma}



\smallskip



\begin{proof}

Из \eqref{Right-0} следует оценка

$$

\Bigl| y^{-\alpha\beta-\alpha n-1} \phi\Bigl(-\rho,-\alpha\beta-\alpha n; -\frac{1}{y^{\rho}}\Bigr)\Bigr|

\leqslant

 \frac{\Gamma\left( 1/\rho+\varepsilon\beta+\varepsilon n\right)}{ \rho \pi\left(aC_{\rho} \right)^{1/\rho+\varepsilon\beta+\varepsilon n}}

\phi\Bigl(-\rho,1; -\frac{1-a}{y^{\rho}}\Bigr),

$$

где $\varepsilon=\alpha/\rho,$ $a\in(0;1),$ $C_{\rho}=C_{\rho}(\rho,\omega).$



С помощью последней оценки получим

\begin{multline}

\Bigl| y^{\,\mu-1} T_{\mu}\Bigl(\frac{x^2}{4y^{\alpha}},\frac{1}{y^{\rho}}\Bigr)\Bigr|=

\\

=\left|

 D_{0y}^{-\mu}

  \sum\limits_{n=0}^{\infty}\frac{|x|^{2\beta+2n}}{ 2^{2\beta+2n}n! \Gamma(1+\beta+n)}

y^{-\alpha\beta-\alpha n-1}\phi \Bigl(-\rho,-\alpha\beta-\alpha n; -\frac{1}{y^{\rho}}\Bigr)

 \right|\leqslant

\\

 \label{WrightRyadZ}  

\leqslant

\frac{ |x|^{2\beta}y^{\mu}\phi\bigl(-\rho,\mu+1; -\frac{1-a}{y^{\rho}}\bigr)}{2^{2\beta}\rho \pi \left(aC_{\rho} \right)^{1/\rho+\varepsilon\beta}}

\sum\limits_{n=0}^{\infty}

\frac{\Gamma\left(1/\rho+\varepsilon\beta+\varepsilon n\right)}{n!

\Gamma(1+\beta+n)}

\Bigl( \frac{x^2}{4\left(aC_{\rho} \right)^{\varepsilon}}\Bigr)^{n}.

\end{multline}



Степенной ряд, стоящий справа от неравенства \eqref{WrightRyadZ}, представляет собой 

обобщенную функцию Райта

$$

{_1\Psi_{1}}\left[ \,z \,\bigg|

\begin{array}{ll}

 \big(1/\rho+\varepsilon\beta, \,\varepsilon\,\big)\\

\big(\,1+\beta, 1\,\big)\\

\end{array}

\right]=

\sum\limits_{n=0}^{\infty}

\frac{\Gamma\left(1/\rho+\varepsilon\beta+\varepsilon n\right)}{n!

\,\Gamma(1+\beta+n)} z^{n}, \; \; z=\frac{x^{\,2}}{4(aC_{\rho})^{\varepsilon}}.

$$

Из ее асимптотических свойств при $z\to \infty$ следует~\cite{Wright}

$$

{_1\Psi_{1}}\left[ \,z\,\bigg|

\begin{array}{ll}

 \big(1/\rho+\varepsilon\beta, \,\varepsilon\,\big)\\

\big(\,1+\beta, 1\,\big)\\

\end{array}

\right]=Z^{\,1/\rho-\beta(1-\varepsilon)-1}\,e^{Z}

\left[ \sum\limits_{m=0}^{M-1}A_m Z^{-m}+O\left(Z^{-M}\right) \right],

$$

где

$Z=(2-\varepsilon)\,\varepsilon^{\,\frac{\varepsilon}{2-\varepsilon}}z^{\,\frac{1}{2-\varepsilon}},$  $\varepsilon<2,$ а коэффициенты $A_m,$ $m=0, 1, 2, ...,$ зависят только от $\rho,$ $\varepsilon$ и $\beta.$



Учитывая, что $a$ "--- произвольное число из интервала $(0;1),$ из последней оценки и~соотношений \eqref{WrightRyadZ}, \eqref{Right_delta>1} получаем \eqref{T_muOtsenka}. \hfill \end{proof}



\smallskip



Покажем, что функция $y^{\mu-1} T_{\mu}\bigl(\frac{x^2}{4y^{\alpha}},\frac{1}{y^{\rho}}\bigr)$ является решением уравнения \eqref{GL3_1}. 

Из \eqref{T(z,zeta)}, \eqref{diff_x}  и \eqref{diff_Right} имеем

\begin{multline}

B_x T_{\mu}\Bigl(\frac{x^2}{4y^{\alpha}},\frac{1}{y^{\rho}}\Bigr)=

\\

=y^{-\alpha}\sum\limits_{n=1}^{\infty}\frac{1}{(n-1)! \Gamma(\beta+n)}

\Bigl(\frac{x^2}{4y^{\alpha}}\Bigr)^{\beta+n-1}

\phi\Bigl(-\rho,\mu-\alpha\beta-\alpha n; -\frac{1}{y^{\rho}}\Bigr)=

\\

=y^{-\alpha}\sum\limits_{k=0}^{\infty}\frac{1}{k! \Gamma(1+\beta+k)}

\Bigl(\frac{x^2}{4y^{\alpha}}\Bigr)^{\beta+k}

\phi\Bigl(-\rho,\mu-\alpha-\alpha\beta-\alpha k; -\frac{1}{y^{\rho}}\Bigr)=

\\ 

=y^{-\alpha} T_{\mu-\alpha}\Bigl(\frac{x^2}{\,4y^{\alpha}},\frac{1}{y^{\rho}}\Bigr),

\end{multline}   

\begin{equation}

\label{DT(z,zeta)}  

D_{0y}^{\nu} y^{\mu-1} T_{\mu}\Bigl(\frac{x^2}{4y^{\alpha}},\frac{1}{y^{\rho}}\Bigr)=

y^{\mu-\nu-1} T_{\mu-\nu}\Bigl(\frac{x^2}{4y^{\alpha}},\frac{1}{y^{\rho}}\Bigr).

\end{equation}



Из последних двух равенств следует

\begin{equation} \label{BDyT}  

\left(B_x-D_{0y}^{\alpha}\right)y^{\mu-1} T_{\mu}\Bigl(\frac{x^2}{4y^{\alpha}}, \frac{1}{y^{\rho}}\Bigr)=0.

\end{equation}   





Из равенства \eqref{DT(z,zeta)} и оценки \eqref{T_muOtsenka} также имеем

\begin{equation} \label{D(alpha-1)T}  

\lim\limits_{y \rightarrow 0} D_{0y}^{\alpha-1} y^{\mu-1} T_{\mu}\Bigl(\frac{x^2}{4y^{\alpha}}, \frac{1}{y^{\rho}}\Bigr)=0.

\end{equation}   



Таким образом, из соотношений  \eqref{T_muOtsenka}, \eqref{BDyT} и \eqref{D(alpha-1)T} следует, что нетривиальная функция 

$$

u(x,y)=y^{\mu-1} T_{\mu}\Bigl(\frac{x^2}{4y^{\alpha}},\frac{1}{y^{\rho}}\Bigr)

$$

является решением однородного уравнения \eqref{GL3_1}, удовлетворяет однородному условию \eqref{UslSMCoshi} $(\varphi(x)\equiv 0)$ и для любого положительного сколь угодно малого $\delta,$ 

некоторых положительных постоянных $k$ и $C,$ а также параметра 

$\rho,$  выбранного из условия 

$\rho=\frac{\alpha}{2}\frac{2+\delta(2-\alpha)}{\alpha+\delta(2-\alpha)},$

имеет место оценка

$$

\left| u(x,y) \right|\leqslant C \exp\left(k |x|^{\delta+\frac{2}{2-\alpha}} \right).

$$



Последняя оценка показывает, что увеличение показателя  $2/(2-\alpha)$ в~условии \eqref{AnalogTikhonovaInfty} ведет к~неединственности решения задачи  \eqref{GL3_1},  \eqref{UslSMCoshi} и~в~этом смысле условие  \eqref{AnalogTikhonovaInfty} является необходимым.





\smallskip









\Section[N]{Заключение}

В данной работе показывается, что показатель степени в~условии, обеспечивающем единственность решения задачи Коши для уравнения дробной диффузии с оператором Бесселя, является предельным, и его увеличение приводит к нарушению единственности. 



\smallskip





\AdditionalContent{Конкурирующие интересы}{Конкурирующих интересов не имею.} 



\AdditionalContent{Авторская ответственность}{Я несу полную ответственность за предоставление окончательной версии рукописи в печать. Окончательная версия рукописи мною одобрена.} 





\AdditionalContent{Финансирование}{Исследование выполнялось без финансирования.}



\smallskip





\begin{thebibliography}{99}



\RBibitem{Drob}

\by Нахушев~А.~М.

\book Дробное исчисление и его применение

\yr 2003

\publ Физматлит

\publaddr М.

\totalpages 271

%\sortdate 1/1/2003

\zmath{https://zbmath.org/?q=an:1066.26005}

%\misctransl

%\by Nakhushev~A.~M.

%\book Drobnoe ischislenie i ego primenenie \rm [Fractional calculus and its applications]

%\yr 2003

%\publ Fizmatlit

%\publaddr Moscow

%\lang In Russian

%\totalpages 271

%%\sortdate 1/1/2003



\RBibitem{Pskhu}

\by Псху~А.~В.

\book Уравнения в частных производных дробного порядка

\yr 2005

\publ Наука

\publaddr М.

\totalpages 199

%\sortdate 1/1/2005

\zmath{https://zbmath.org/?q=an:1193.35245}

%\misctransl

%\by Pskhu~A.~V.

%\book Uravneniia v chastnykh proizvodnykh drobnogo poriadka \rm [Partial differential equations of fractional order]

%\yr 2005

%\publ Nauka

%\publaddr Moscow

%\lang In Russian

%\totalpages 199

%%\sortdate 1/1/2005



\Bibitem{Alexiades}

\by Alexiades~V.

\paper Generalized axially symmetric heat potentials and singular parabolic initial boundary value problems

\jour Arch. Rational Mech. Anal.

\yr 1982

\vol 79

\issue 4

\pages 325--350

%\citnumscopus 3

%\sortdate 1/1/1982

\crossref{10.1007/BF00250797}

\mathscinet{http://www.ams.org/mathscinet-getitem?mr=656798}

\zmath{https://zbmath.org/?q=an:0511.35050}

\scopus{http://www.scopus.com/record/display.url?origin=inward&eid=2-s2.0-0020402046}



\Bibitem{Calton}

\by Calton~D.

\paper Cauchy's problem for a singular parabolic partial differential equation

\jour J.~Diff. Equations

\yr 1970

\vol 8

\issue 2

\pages 250--257

%\citnumscopus 8

%\sortdate 1/1/1970

\crossref{10.1016/0022-0396(70)90004-5}

\mathscinet{http://www.ams.org/mathscinet-getitem?mr=271547}

\scopus{http://www.scopus.com/record/display.url?origin=inward&eid=2-s2.0-49849107701}



\Bibitem{Kepinski_2}

\by K\c{e}pi\'nski~S.

\paper \"Uber die Differentialgleichung $\frac{\partial^2 z }{\partial x^2}+\frac{m+1}{x}\frac{\partial z }{\partial x} -\frac{n}{x}\frac{\partial z }{\partial t}=0$

\jour Math. Ann.

\yr 1905

\vol 61

\issue 3

\pages 397--405

\lang In German

%\sortdate 1/1/1905

\mathscinet{http://www.ams.org/mathscinet-getitem?mr=1511354}

\zmath{https://zbmath.org/?q=an:36.0416.01}



\RBibitem{Tersenov}

\by Терсенов~С.~А.

\book Параболические уравнения с~меняющимся направлением времени

\yr 1985

\publ Наука

\publaddr М.

\totalpages 105

%\sortdate 1/1/1985

%\misctransl

%\by Tersenov~S.~A.

%\book Parabolicheskie uravneniia s~meniaiushchimsia napravleniem vremeni \rm [Parabolic Equations with a Changing Time Direction]

%\yr 1985

%\publ Nauka

%\publaddr Moscow

%\totalpages 105

%%\sortdate 1/1/1985



\RBibitem{Matiychuk}

\by Матiйчук~M.~I.

\book Параболiчнi сингулярнi крайовi задачi

\yr 1999

\publ Iн-т математики НАН Укра\"\iни

\publaddr Ки\"\iв

\totalpages 176

%%\sortdate 1/1/1999

%\misctransl

%\by Matiichuk~M.~I.

%\book Parabolichni singuliarni kraiovi zadachi \rm [Parabolic Singular Boundary-Value Problems]

%\yr 1999

%\publ Institute of Mathematics of the Ukrainian National Academy of Sciences

%\publaddr Kiev

%\lang In Ukrainian

%\totalpages 176

%%\sortdate 1/1/1999



\RBibitem{Kipriyanov-Monographiya}

\by Киприянов~И.~А.

\book Сингулярные эллиптические краевые задачи

\yr 1997

\publ Наука

\publaddr М.

\totalpages 208

%\sortdate 1/1/1997

%\misctransl

%\by Kipriyanov~I.~A.

%\book Singuliarnye ellipticheskie kraevye zadachi \rm [Singular Elliptic Boundary-Value Problems]

%\yr 1997

%\publ Nauka

%\publaddr Moscow

%\lang In Russian

%\totalpages 208

%%\sortdate 1/1/1997



\RBibitem{Kipriyanov}

\by Киприянов~И.~А., Катрахов~В.~В., Ляпин~В.~М.

\paper О краевых задачах в областях общего вида для сингулярных параболических систем уравнений

\jour Докл. АН СССР

\yr 1976

\vol 230

\issue 6

\pages 1271--1274

%\sortdate 1/1/1976

\mathnet{http://mi.mathnet.ru/rus/dan40691}

\mathscinet{http://www.ams.org/mathscinet-getitem?mr=0425366}

\zmath{https://zbmath.org/?q=an:0354.35056}

%\misctransl

%\by Kipriyanov~I.~A., Katrakhov~V.~V., Lyapin~V.~M.

%\paper On boundary value problems in domains of general type for singular parabolic systems of equations

%\jour Sov. Math., Dokl.

%\yr 1976

%\vol 17

%\pages 1461--1464

%%\sortdate 1/1/1976

%\mathscinet{http://www.ams.org/mathscinet-getitem?mr=0425366}

%\zmath{https://zbmath.org/?q=an:0354.35056}



\RBibitem{Sitnik}

\by Ситник~С.~М.

\yr 2016

\publaddr Воронеж

\thesis Применение операторов преобразования Бушмана--Эрдейи и их обобщений в~теории дифференциальных уравнений с~особенностями в коэффициентах

\thesisinfo Дис.~\dots\ д-ра физ.-мат. наук: 01.01.02

\totalpages 307

%\sortdate 1/1/2016

%\misctransl

%\by Sitnik~S.~M.

%\yr 2016

%\publaddr Voronezh

%\thesis Application of the Bushman--Erd\'elyi transformation operators and their generalizations in the theory of differential equations with singularities in the coefficients

%\thesisinfo Dr. Phys. Math. Sci. Thesis

%\lang In Russian

%\totalpages 307

%%\sortdate 1/1/2016



\RBibitem{KatrakhovSitnik}

\by Катрахов~В.~В., Ситник~С.~М.

\paper Метод операторов преобразования и краевые задачи для сингулярных эллиптических уравнений

\inbook Сингулярные дифференциальные уравнения

\serial СМФН

\yr 2018

\vol 64

\issue 2

\pages 211--426

\publ Российский университет дружбы народов

\publaddr М.

%\sortdate 1/1/2018

\mathnet{http://mi.mathnet.ru/rus/cmfd355}

\crossref{10.22363/2413-3639-2018-64-2-211-426}

%\misctransl

%\by Katrakhov~V.~V., Sitnik~S.~M.

%\paper The transmutation method and boundary-value problems for singular elliptic equations

%\inbook Singular differential equations

%\serial CMFD

%\yr 2018

%\vol 64

%\issue 2

%\pages 211--426

%\publ Peoples' Friendship University of Russia

%\publaddr Moscow

%\lang In Russian

%%\sortdate 1/1/2018

%\crossref{https://doi.org/10.22363/2413-3639-2018-64-2-211-426}



\RBibitem{Kochubei_2}

\by Кочубей~А.~Н.

\paper Диффузия дробного порядка

\jour Дифференц. уравнения

\yr 1990

\vol 26

\issue 4

\pages 660--670

%\sortdate 1/1/1990

\mathnet{http://mi.mathnet.ru/rus/de7144}

\mathscinet{http://www.ams.org/mathscinet-getitem?mr=1061448}

\zmath{https://zbmath.org/?q=an:0712.35049}

%\transl

%\by Kochubei~A.~N.

%\paper Diffusion of fractional order

%\jour Differ. Equ.

%\yr 1990

%\vol 26

%\issue 4

%\pages 485--492

%%\sortdate 1/1/1990

%\mathscinet{http://www.ams.org/mathscinet-getitem?mr=1061448}

%\zmath{https://zbmath.org/?q=an:0729.35064}



\RBibitem{Kochubei_1}

\by Кочубей~А.~Н.

\paper Задача Коши для эволюционных уравнений дробного порядка

\jour Дифференц. уравнения

\yr 1989

\vol 25

\issue 8

\pages 1359--1368

%\sortdate 1/1/1989

\mathnet{http://mi.mathnet.ru/rus/de6938}

\mathscinet{http://www.ams.org/mathscinet-getitem?mr=1014153}

%\transl

%\by Kochubei~A.~N.

%\paper The Cauchy problem for evolution equations of fractional order

%\jour Differ. Equ.

%\yr 1989

%\vol 25

%\issue 8

%\pages 967--974

%%\sortdate 1/1/1989

%\mathscinet{http://www.ams.org/mathscinet-getitem?mr=1014153}



\RBibitem{Kochubei_3}

\by Кочубей~А.~Н., Эйдельман~С.~Д.

\paper Задача Коши для эволюционных уравнений дробного порядка

\jour Докл. РАН

\yr 2004

\vol 394

\issue 2

\pages 159--161

%\sortdate 1/1/2004

\mathnet{http://mi.mathnet.ru/rus/dan1624}

\mathscinet{http://www.ams.org/mathscinet-getitem?mr=2089744}

\zmath{https://zbmath.org/?q=an:1129.35427}

%\misctransl

%\by Kochubeii~A.~N., \'Eidel'man~S.~D.

%\paper The Cauchy problem for evolution equations of fractional order

%\jour Dokl. Akad. Nauk

%\yr 2004

%\vol 394

%\issue 2

%\pages 159--161

%\lang In Russian

%%\sortdate 1/1/2004



\Bibitem{Kochubei_4}

\by Kochubei~A.~N.

\paper Asymptotic Properties of Solutions of the Fractional Diffusion-Wave Equation

\jour Fract. Calc. Appl. Anal.

\yr 2014

\vol 17

\issue 3

\pages 881--896

%\citnumscopus 9

%\sortdate 1/1/2014

\crossref{10.2478/s13540-014-0203-3}

\mathscinet{http://www.ams.org/mathscinet-getitem?mr=3260311}

\zmath{https://zbmath.org/?q=an:1323.35197}

\elib{http://elibrary.ru/item.asp?id=24541812}

\scopus{http://www.scopus.com/record/display.url?origin=inward&eid=2-s2.0-84904312827}



\Bibitem{Kochubei_5}

\by Kochubei~A.~N.

\paper Cauchy problem for fractional diffusion-wave equations with variable coefficients

\jour Applicable Analysis

\yr 2014

\vol 93

\issue 10

\pages 2211--2242

%\citnumscopus 7

%\sortdate 1/1/2014

\crossref{10.1080/00036811.2013.875162}

\mathscinet{http://www.ams.org/mathscinet-getitem?mr=3240385}

\zmath{https://zbmath.org/?q=an:1297.35274}

\elib{http://elibrary.ru/item.asp?id=24642244}

\scopus{http://www.scopus.com/record/display.url?origin=inward&eid=2-s2.0-84905387477}



\RBibitem{Pskhu2}

\by Псху~А.~В.

\paper Фундаментальное решение диффузионно-волнового уравнения дробного порядка

\jour Изв. РАН. Сер. матем.

\yr 2009

\vol 73

\issue 2

\pages 141--182

%\sortdate 1/1/2009

\mathnet{http://mi.mathnet.ru/rus/im2429}

\crossref{10.4213/im2429}

\mathscinet{http://www.ams.org/mathscinet-getitem?mr=2532450}

\zmath{https://zbmath.org/?q=an:1172.26001}

\adsnasa{http://adsabs.harvard.edu/cgi-bin/bib_query?2009IzMat..73..351P}

\elib{http://elibrary.ru/item.asp?id=20425204}

%\transl

%\by Pskhu~A.~V.

%\paper The fundamental solution of a diffusion-wave equation of fractional order

%\jour Izv. Math.

%\yr 2009

%\vol 73

%\issue 2

%\pages 351--392

%%\sortdate 1/1/2009

%\crossref{https://doi.org/10.1070/IM2009v073n02ABEH002450}

%\mathscinet{http://www.ams.org/mathscinet-getitem?mr=2532450}

%\zmath{https://zbmath.org/?q=an:1172.26001}

%\isi{http://gateway.isiknowledge.com/gateway/Gateway.cgi?GWVersion=2&SrcApp=PARTNER_APP&SrcAuth=LinksAMR&DestLinkType=FullRecord&DestApp=ALL_WOS&KeyUT=000266177900006}

%\elib{http://elibrary.ru/item.asp?id=13597908}

%\scopus{http://www.scopus.com/record/display.url?origin=inward&eid=2-s2.0-65349154473}



\RBibitem{Pskhu-SibirJournal}

\by Псху~А.~В.

\paper Уравнение дробной диффузии с оператором дискретно распределенного дифференцирования

\jour Сиб. электрон. матем. изв.

\yr 2016

\vol 13

\pages 1078--1098

%\sortdate 1/1/2016

\mathnet{http://mi.mathnet.ru/rus/semr736}

\crossref{10.17377/semi.2016.13.086}

\mathscinet{http://www.ams.org/mathscinet-getitem?mr=3580054}

\zmath{https://zbmath.org/?q=an:1370.35268}

\elib{http://elibrary.ru/item.asp?id=28127183}

%\misctransl

%\by Pskhu~A.~V.

%\paper Fractional diffusion equation with discretely distributed differentiation operator

%\jour Sib. \`Elektron. Mat. Izv.

%\yr 2016

%\vol 13

%\pages 1078--1098

%\lang In Russian

%%\sortdate 1/1/2016

%\crossref{https://doi.org/10.17377/semi.2016.13.086}



\Bibitem{Pskhu3}

\by Pskhu~A.~V.

\paper Multi-time fractional diffusion equation

\jour Eur. Phys. J.~Spec. Top.

\yr 2013

\vol 222

\issue 8

\pages 1939--1950

%\citnumscopus 2

%\sortdate 1/1/2013

\crossref{10.1140/epjst/e2013-01975-y}

\scopus{http://www.scopus.com/record/display.url?origin=inward&eid=2-s2.0-84885034281}



\RBibitem{Pskhu4}

\by Псху~А.~В.

\paper Первая краевая задача для дробного диффузионно-волнового уравнения в~нецилиндрической области

\jour Изв. РАН. Сер. матем.

\yr 2017

\vol 81

\issue 6

\pages 158--179

%\sortdate 1/1/2017

\mathnet{http://mi.mathnet.ru/rus/im8520}

\crossref{10.4213/im8520}

\mathscinet{http://www.ams.org/mathscinet-getitem?mr=3733319}

\zmath{https://zbmath.org/?q=an:06844352}

\adsnasa{http://adsabs.harvard.edu/cgi-bin/bib_query?2017IzMat..81.1212P}

\elib{http://elibrary.ru/item.asp?id=30737847}

%\transl

%\by Pskhu~A.~V.

%\paper The first boundary-value problem for a fractional diffusion-wave equation in a non-cylindrical domain

%\jour Izv. Math.

%\yr 2017

%\vol 81

%\issue 6

%\pages 1212--1233

%%\sortdate 1/1/2017

%\crossref{https://doi.org/10.1070/IM8520}

%\mathscinet{http://www.ams.org/mathscinet-getitem?mr=3733319}

%\zmath{https://zbmath.org/?q=an:06844352}

%\isi{http://gateway.isiknowledge.com/gateway/Gateway.cgi?GWVersion=2&SrcApp=PARTNER_APP&SrcAuth=LinksAMR&DestLinkType=FullRecord&DestApp=ALL_WOS&KeyUT=000418891300007}

%\scopus{http://www.scopus.com/record/display.url?origin=inward&eid=2-s2.0-85040983678}



\RBibitem{VorKilbas}

\by Ворошилов~А.~А., Килбас~А.~А.

\paper Задача типа Коши для диффузионно-волнового уравнения с частной производной Римана--Лиувилля

\jour Докл. РАН

\yr 2006

\vol 406

\issue 1

\pages 12--16

%\sortdate 1/1/2006

\mathnet{http://mi.mathnet.ru/rus/dan1045}

\mathscinet{http://www.ams.org/mathscinet-getitem?mr=2256848}

\zmath{https://zbmath.org/?q=an:1155.35396}

\elib{http://elibrary.ru/item.asp?id=9176502}

%\misctransl

%\by Voroshilov~A.~A., Kilbas~A.~A.

%\paper A Cauchy-type problem for a diffusion-wave equation with a Riemann-Liouville partial derivative

%\jour Dokl. Akad. Nauk

%\yr 2006

%\vol 406

%\issue 1

%\pages 12--16

%\lang In Russian

%%\sortdate 1/1/2006



\RBibitem{VorKilbas2}

\by Ворошилов~А.~А., Килбас~А.~А.

\paper Задача Коши для диффузионно-волнового уравнения с~частной производной Капуто

\jour Дифференц. уравнения

\yr 2006

\vol 42

\issue 5

\pages 599--609

%\sortdate 1/1/2006

\mathnet{http://mi.mathnet.ru/rus/de11489}

\mathscinet{http://www.ams.org/mathscinet-getitem?mr=2292158}

\zmath{https://zbmath.org/?q=an:1123.35302}

\elib{http://elibrary.ru/item.asp?id=9245260}

%\transl

%\by Voroshilov~A.~A., Kilbas~A.~A.

%\paper The Cauchy problem for the diffusion-wave equation with the Caputo partial derivative

%\jour Differ. Equ.

%\yr 2006

%\vol 42

%\issue 5

%\pages 638--649

%%\citnumscopus 11

%%\sortdate 1/1/2006

%\crossref{https://doi.org/10.1134/S0012266106050041}

%\mathscinet{http://www.ams.org/mathscinet-getitem?mr=2292158}

%\zmath{https://zbmath.org/?q=an:1123.35302}

%\elib{http://elibrary.ru/item.asp?id=27786841}

%\scopus{http://www.scopus.com/record/display.url?origin=inward&eid=2-s2.0-33745304762}



\RBibitem{Gekk}

\by Геккиева~С.~Х.

\paper Задача Коши для обобщенного уравнения переноса с дробной по времени производной

\jour Доклады Адыгской (Черкесской) Международной академии наук

\yr 2000

\vol 5

\issue 1

\pages 16--19

%\sortdate 1/1/2000

%\misctransl

%\by Gekkieva~S.~Kh.

%\paper The Cauchy problem for the generalized transport equation with time-fractional derivative

%\jour Dokl. Adyg. (Cherkess.) Mezdunar. Akad. Nauk

%\yr 2000

%\vol 5

%\issue 1

%\pages 16--19

%\lang In Russian

%%\sortdate 1/1/2000



\RBibitem{Gekkieva1}

\by Геккиева~С.~Х.

\paper Краевая задача для обобщенного уравнения переноса с дробной производной в полубесконечной области

\jour Изв. КБНЦ РАН

\yr 2002

\issue 1(8)

\pages 6--8

%\sortdate 1/1/2002

%\misctransl

%\by Gekkieva~S.~Kh.

%\paper A boundary value problem for the generalized transfer equation with a fractional derivative in a semi-infinite domain

%\jour Izv. KBNTs RAN

%\yr 2002

%\issue 1(8)

%\pages 6--8

%\lang In Russian

%%\sortdate 1/1/2002



\Bibitem{Mainardi-1}

\by Mainardi~F.

\paper The fundamental solutions for the fractional diffusion-wave equation

\jour Appl. Math. Letters

\yr 1996

\vol 6

\issue 1

\pages 23--28

%\citnumscopus 343

%\sortdate 1/1/1996

\crossref{10.1016/0893-9659(96)00089-4}

\mathscinet{http://www.ams.org/mathscinet-getitem?mr=1419811}

\zmath{https://zbmath.org/?q=an:0879.35036}

\scopus{http://www.scopus.com/record/display.url?origin=inward&eid=2-s2.0-30244460855}



\Bibitem{Mainardi-3}

\by Mainardi~F.

\paper The time fractional diffusion-wave equation

\jour Radiophys. Quantum Electron.

\yr 1995

\vol 38

\issue 1-2

\pages 13--24

%\citnumscopus 39

%\sortdate 1/1/1995

\crossref{10.1007/BF01051854}

\mathscinet{http://www.ams.org/mathscinet-getitem?mr=1427165}

\scopus{http://www.scopus.com/record/display.url?origin=inward&eid=2-s2.0-0029427754}



\Bibitem{Mainardi-2}

\by Mainardi~F., Luchko~Yu., Pagnini~G.

\paper The fundamental solution of the space-time fractional diffusion equation

\jour Fract. Calc. Appl. Anal.

\yr 2001

\vol 4

\issue 2

\pages 153--192

\arxiv \href{https://arxiv.org/abs/cond-mat/0702419}{cond-mat/0702419} [cond-mat.stat-mech]

%\sortdate 1/1/2001

\mathscinet{http://www.ams.org/mathscinet-getitem?mr=1829592}

\zmath{https://zbmath.org/?q=an:1054.35156}



\Bibitem{Pagnini_1}

\by Pagnini~G.

\paper The M-Wright function as a generalization of the Gaussian density for fractional diffusion processes

\jour Fract. Calc. Appl. Anal.

\yr 2013

\vol 16

\issue 2

\pages 436--453

%\citnumscopus 22

%\sortdate 1/1/2013

\crossref{10.2478/s13540-013-0027-6}

\mathscinet{http://www.ams.org/mathscinet-getitem?mr=3033698}

\zmath{https://zbmath.org/?q=an:1312.33061}

\elib{http://elibrary.ru/item.asp?id=28682912}

\scopus{http://www.scopus.com/record/display.url?origin=inward&eid=2-s2.0-84879620364}



\Bibitem{Pagnini_Paradisi}

\by Pagnini~G., Paradisi~P.

\paper A stochastic solution with Gaussian stationary increments of the symmetric space-time fractional diffusion equation

\jour Fract. Calc. Appl. Anal.

\yr 2016

\vol 19

\issue 2

\pages 408--440

%\citnumscopus 9

%\sortdate 1/1/2016

\crossref{10.1515/fca-2016-0022}

\mathscinet{http://www.ams.org/mathscinet-getitem?mr=3513003}

\zmath{https://zbmath.org/?q=an:1341.60073}

\elib{http://elibrary.ru/item.asp?id=27104131}

\scopus{http://www.scopus.com/record/display.url?origin=inward&eid=2-s2.0-84969800118}



\Bibitem{Podlubny}

\by Podlubny~I.

\book Fractional differential equations. An introduction to fractional derivatives, fractional differential equations, to methods of their solution and some of their applications

\serial Mathematics in Science and Engineering

\yr 1999

\vol 198

\publ Academic Press

\publaddr San Diego, CA

\totalpages xxiv+340

%\sortdate 1/1/1999

\mathscinet{http://www.ams.org/mathscinet-getitem?mr=1658022}

\zmath{https://zbmath.org/?q=an:0924.34008}



\Bibitem{MathaiSaxenaHaubold}

\by Mathai~A.~M., Saxena~R.~K., Haubold~H.~J.

\book The $H$-function. Theory and Applications

\yr 2010

\publ Springer

\publaddr Dordrecht

\totalpages xiv+268~pp \nofrills

%\citnumscopus 411

%\sortdate 1/1/2010

\crossref{10.1007/978-1-4419-0916-9}

\mathscinet{http://www.ams.org/mathscinet-getitem?mr=2562766}

\zmath{https://zbmath.org/?q=an:1181.33001}

\scopus{http://www.scopus.com/record/display.url?origin=inward&eid=2-s2.0-84891408876}



\RBibitem{KhushtovaVestn2}

\by Хуштова~Ф.~Г.

\paper Задача Коши для уравнения параболического типа~с~оператором Бесселя и~частной производной~Римана--Лиувилля

\jour Вестн. Сам. гос. техн. ун-та. Сер. Физ.-мат. науки

\yr 2016

\vol 20

\issue 1

\pages 74--84

%\sortdate 1/1/2016

\mathnet{http://mi.mathnet.ru/rus/vsgtu1455}

\crossref{10.14498/vsgtu1455}

\zmath{https://zbmath.org/?q=an:06964474}

\elib{http://elibrary.ru/item.asp?id=26898117}

%\misctransl

%\by Khushtova~F.~G.

%\paper Cauchy problem for a parabolic equation with Bessel operator and Riemann--Liouville partial derivative

%\jour Vestn. Samar. Gos. Tekhn. Univ., Ser. Fiz.-Mat. Nauki \rm [J. Samara State Tech. Univ., Ser. Phys. Math. Sci.]

%\yr 2016

%\vol 20

%\issue 1

%\pages 74--84

%%\sortdate 1/1/2016



\Bibitem{GorenfloLuchkoMainardi}

\by Gorenflo~R., Luchko~Y., Mainardi~F.

\paper Analytical properties and applications of the Wright function

\jour Fract. Calc. Appl. Anal.

\yr 1999

\vol 2

\issue 4

\pages 383--414

\arxiv \href{https://arxiv.org/abs/math-ph/0701069}{math-ph/0701069}

%\sortdate 1/1/1999

\mathscinet{http://www.ams.org/mathscinet-getitem?mr=1752379}

\zmath{https://zbmath.org/?q=an:1027.33006}



\Bibitem{tikhonovUsl-1}

\by Tychonoff~A.

\paper Th\'eor\`emes d'unicit\'e pour l'\'equation de la chaleur

\jour Mat. Sb.

\yr 1935

\vol 42

\issue 2

\pages 199--216

\lang In French

%\sortdate 1/1/1935

\mathnet{http://mi.mathnet.ru/rus/sm6410}

\zmath{https://zbmath.org/?q=an:0012.35501|61.1203.05}



\RBibitem{tikhonovUsl-2}

\by Тихонов~А.~Н.

\paper Теоремы единственности для уравнения теплопроводности

\inbook Собрание научных трудов: в десяти томах

\yr 2009

\pages 371--382

\publ Наука

\publaddr М.

\volinfo Т.~II. Математика. Ч.~2. Вычислительная математика 1956--1979. Математическая физика 1933--1948. Ред.-сост. Т.~А.~Сушкевич, А.~В.~Гулин

%\sortdate 1/1/2009

\zmath{https://zbmath.org/?q=an:1200.01055}

%\misctransl

%\by Tikhonov~A.~N.

%\paper Uniqueness theorems for the heat equation

%\inbook Collection of scientific works. In ten volumes

%\yr 2009

%\pages 371--382

%\publ Nauka

%\publaddr Moscow

%\lang In Russian and French

%\volinfo Vol. 2: Mathematics. Part 2: Numerical mathematics 1956–1979. Mathematical physics 1933–1948. Edited by T.~A.~Sushkevich and A.~V.~Gulin

%%\sortdate 1/1/2009

%\zmath{https://zbmath.org/?q=an:1200.01055}



\Bibitem{Wright}

\by Wright~E.~M.

\paper The asymptotic expansion of the generalized hypergeometric function

\jour Proc. Lond. Math. Soc., II. Ser.

\yr 1940

\vol 46

\issue 1

\pages 389--408

%\citnumscopus 136

%\sortdate 1/1/1940

\crossref{10.1112/plms/s2-46.1.389}

\mathscinet{http://www.ams.org/mathscinet-getitem?mr=3876}

\zmath{https://zbmath.org/?q=an:0025.40402|66.1243.01}

\scopus{http://www.scopus.com/record/display.url?origin=inward&eid=2-s2.0-84963078619}



\end{thebibliography}



	





\newpage



\makeentitle



\vspace{-6mm}



%\AdditionalContent{Competing interests}{Specify what conflicts of interest are available.} 

\AdditionalContent{Competing interests}{I have no competing interests.} 

%\AdditionalContent{Competing interests}{We have ... (or) no competing interests.} 

%\AdditionalContent{Authors' contributions and responsibilities}{What is the author's responsibility for each author?}

\AdditionalContent{Author's Responsibilities}{I take full responsibility for submitting the final ma\-nuscript in print. I approved the final version of the manuscript.} 

%\AdditionalContent{Authors' contributions and responsibilities}{Each author has participated in the article concept development and in the manuscript writing. The authors are absolutely responsible for submitting the final manuscript in print. Each author has approved the final version of manuscript.} 



%\AdditionalContent{Funding}{Indicate the source(s) of your funding.}



\AdditionalContent{Funding}{This research received no specific grant from any funding agency in the public,

commercial, or not-for-profit sectors.}



\newpage



\begin{thebibliography}{10}



\Bibitem{Drob:2}

\lang In Russian

\by Nakhushev~A.~M.

\book Drobnoe ischislenie i ego primenenie \rm [Fractional calculus and its applications]

\publaddr Moscow

\publ Fizmatlit 

\totalpages 271

\yr 2003





\Bibitem{Pskhu:w}

\lang In Russian

\by Pskhu~A.~V.

\book Uravneniia v chastnykh proizvodnykh drobnogo poriadka \rm [Partial differential equations of fractional order]

\publaddr Moscow

\publ Nauka 

\totalpages 199 

\yr 2005







\Bibitem{Alexiades:w}

\by Alexiades~V.

\paper Generalized axially symmetric heat potentials and singular parabolic initial boundary value problems

\jour  Arch. Rational Mech. Anal. 

\yr 1982

\vol 79

\issue 4

\pages 325--350

\crossref{10.1007/BF00250797}







\Bibitem{Calton:w}

\by Calton~D.

\paper Cauchy's problem for a singular parabolic partial differential equation

\jour J.~Diff. Equations

\yr 1970

\vol 8

\issue 2

\pages 250--257

\crossref{10.1016/0022-0396(70)90004-5}







\Bibitem{Kepinski_2:w}

\by K\c{e}pi\'nski~S.

\paper \"Uber die Differentialgleichung $\frac{\partial^2 z }{\partial x^2}+\frac{m+1}{x}\frac{\partial z }{\partial x} -\frac{n}{x}\frac{\partial z }{\partial t}=0$

\jour   Math. Ann.

\yr 1905

\vol 61

\issue 3

\pages 397--405

\zmath{36.0416.01}

\lang In German



\Bibitem{Tersenov:en}

\lang In Russian

\by Tersenov~S.~A.

\book Parabolicheskie uravneniia s~meniaiushchimsia napravleniem vremeni \rm [Parabolic Equations with a Changing Time Direction]

\yr 1985

\publ Nauka

\publaddr Moscow

\totalpages 105



\Bibitem{Matiychuk:en}

\by Matiichuk~M.~I.

\book  Parabolichni singuliarni kraiovi zadachi \rm [Parabolic Singular Boundary-Value Problems]

\yr 1999

\publ Institute of Mathematics of the Ukrainian National  Academy of Sciences

\publaddr Kiev

\totalpages 176

\lang In Ukrainian





\Bibitem{Kipriyanov-Monographiya:en}

\lang In Russian

\by Kipriyanov~I.~A.

\book  Singuliarnye ellipticheskie kraevye zadachi \rm [Singular Elliptic Boundary-Value Problems]

\yr 1997

\publ Nauka

\publaddr Moscow

\totalpages 208



\Bibitem{Kipriyanov:en}

\by Kipriyanov~I.~A., Katrakhov~V.~V., Lyapin~V.~M.

\paper On boundary value problems in domains of general type for singular parabolic systems of equations

\jour Sov. Math., Dokl. 

\vol 17

\yr 1976

\pages 1461--1464 

\mathscinet{http://www.ams.org/mathscinet-getitem?mr=0425366}

\zmath{https://zbmath.org/?q=an:0354.35056}





\Bibitem{Sitnik:en}

\lang In Russian

\by Sitnik~S.~M. 

\thesis Application of the Bushman--Erd\'elyi transformation operators and their generalizations in the theory of differential equations with singularities in the coefficients

\thesisinfo Dr. Phys. Math. Sci. Thesis

\yr 2016

\publaddr Voronezh

\totalpages 307







\Bibitem{KatrakhovSitnik:en}

\lang In Russian

\by Katrakhov~V.~V., Sitnik~S.~M.

\paper The transmutation method and boundary-value problems for singular elliptic equations

\inbook Singular differential equations

\serial CMFD

\yr 2018

\vol 64

\issue 2

\pages 211--426

\publ Peoples' Friendship University of Russia

\publaddr Moscow

\crossref{10.22363/2413-3639-2018-64-2-211-426}







\Bibitem{Kochubei_2:en}

\by Kochubei~A.~N.

\paper Diffusion of fractional order

\jour Differ. Equ.

\yr 1990

\vol 26

\issue 4

\pages 485--492







\Bibitem{Kochubei_1:w}

\by Kochubei~A.~N.

\paper The Cauchy problem for evolution equations of fractional order

\jour Differ. Equ.

\yr 1989

\vol 25

\issue 8

\pages 967--974



\Bibitem{Kochubei_3:e}

\lang In Russian

\by Kochubeii~A.~N., \'Eidel'man~S.~D.

\paper The Cauchy problem for evolution equations of fractional order

\jour Dokl. Akad. Nauk

\yr 2004

\vol 394

\issue 2

\pages 159--161







\Bibitem{Kochubei_4:e}

\by Kochubei~A.~N.

\paper Asymptotic Properties of Solutions of the Fractional Diffusion-Wave Equation

\jour Fract. Calc. Appl. Anal.

\yr 2014

\vol 17

\issue 3

\pages 881--896

\crossref{10.2478/s13540-014-0203-3}





\Bibitem{Kochubei_5:e}

\by Kochubei~A.~N.

\paper Cauchy problem for fractional diffusion-wave equations with variable coefficients

\jour  Applicable Analysis

\yr 2014

\vol 93

\issue 10

\pages 2211--2242

\crossref{10.1080/00036811.2013.875162}







\Bibitem{Pskhu2:en}

\paper The fundamental solution of a diffusion-wave equation of fractional order

\by Pskhu~A.~V.

\jour Izv. Math.

\yr 2009

\vol 73

\issue 2

\pages 351--392

\crossref{10.1070/IM2009v073n02ABEH002450}

\isi{http://gateway.isiknowledge.com/gateway/Gateway.cgi?GWVersion=2&SrcApp=PARTNER_APP&SrcAuth=LinksAMR&DestLinkType=FullRecord&DestApp=ALL_WOS&KeyUT=000266177900006}

\elib{http://elibrary.ru/item.asp?id=13597908}

\scopus{http://www.scopus.com/record/display.url?origin=inward&eid=2-s2.0-65349154473}





\Bibitem{Pskhu-SibirJournal:en}

\lang In Russian

\by Pskhu~A.~V.

\paper Fractional diffusion equation with discretely distributed differentiation operator

\jour Sib. \`Elektron. Mat. Izv.

\yr 2016

\vol 13

\pages 1078--1098

\crossref{10.17377/semi.2016.13.086}







\Bibitem{Pskhu3:en}

\by Pskhu~A.~V.

\paper Multi-time fractional diffusion equation

\jour Eur. Phys. J.~Spec. Top.

\yr 2013

\vol 222

\issue 8

\pages 1939--1950

\crossref{10.1140/epjst/e2013-01975-y}





\Bibitem{Pskhu4:en}

\paper The first boundary-value problem for a fractional diffusion-wave equation in a non-cylindrical domain

\by Pskhu~A.~V.

\jour Izv. Math.

\yr 2017

\vol 81

\issue 6

\pages 1212--1233

\crossref{10.1070/IM8520}

\isi{http://gateway.isiknowledge.com/gateway/Gateway.cgi?GWVersion=2&SrcApp=PARTNER_APP&SrcAuth=LinksAMR&DestLinkType=FullRecord&DestApp=ALL_WOS&KeyUT=000418891300007}

\scopus{http://www.scopus.com/record/display.url?origin=inward&eid=2-s2.0-85040983678}







\Bibitem{VorKilbas:a} 

\lang In Russian

\by Voroshilov~A.~A., Kilbas~A.~A.

\paper A Cauchy-type problem for a diffusion-wave equation with a

Riemann-Liouville partial derivative

\jour Dokl. Akad. Nauk

\yr 2006

\vol 406

\issue 1

\pages 12--16



\Bibitem{VorKilbas2:en} 

\by Voroshilov~A.~A., Kilbas~A.~A.

\paper The Cauchy problem for the diffusion-wave equation with the Caputo partial derivative

\mathscinet{http://www.ams.org/mathscinet-getitem?mr=2292158}

\jour Differ. Equ.

\yr 2006

\vol 42

\issue 5

\pages 638--649

\crossref{10.1134/S0012266106050041}



\Bibitem{Gekk:en} 

\lang In Russian

\by Gekkieva~S.~Kh.

\paper The Cauchy problem for the generalized transport equation with time-fractional derivative

\jour Dokl. Adyg. (Cherkess.) Mezdunar. Akad. Nauk

\yr 2000

\vol 5

\issue 1

\pages 16--19







\Bibitem{Gekkieva1:en}

\lang In Russian

\by Gekkieva~S.~Kh.

\paper A boundary value problem for the generalized transfer equation with a fractional derivative in a semi-infinite domain

\jour Izv. KBNTs RAN

\yr 2002

\issue 1(8)

\pages 6--8











  





\Bibitem{Mainardi-1:en}

\by Mainardi~F.

\paper The fundamental solutions for the fractional diffusion-wave equation

\jour Appl.  Math. Letters

\yr 1996

\vol 6

\issue 1

\pages 23--28

\crossref{10.1016/0893-9659(96)00089-4}



\Bibitem{Mainardi-3}

\by Mainardi~F.

\paper  The time fractional diffusion-wave equation

\jour Radiophys. Quantum Electron.

\yr 1995

\vol 38

\issue 1-2

\pages 13--24

\crossref{10.1007/BF01051854}



\Bibitem{Mainardi-2:en}

\by Mainardi~F., Luchko~Yu., Pagnini~G.

\paper The fundamental solution of the space-time fractional diffusion equation

\jour Fract. Calc. Appl. Anal.

\yr 2001

\vol 4

\issue 2

\pages 153--192

\arxiv  	\href{https://arxiv.org/abs/cond-mat/0702419}{cond-mat/0702419} [cond-mat.stat-mech]















\Bibitem{Pagnini_1:en}

\by Pagnini~G.

\paper The M-Wright function as a generalization of the Gaussian density for fractional diffusion processes

\jour Fract. Calc. Appl. Anal.

\yr 2013

\vol 16

\issue 2

\pages 436--453

\crossref{10.2478/s13540-013-0027-6}





\Bibitem{Pagnini_Paradisi:en}

\by Pagnini~G., Paradisi~P.

\paper A stochastic solution with Gaussian stationary increments of the symmetric space-time fractional diffusion equation

\jour Fract. Calc. Appl. Anal.

\yr 2016

\vol 19

\issue 2

\pages 408--440

\crossref{10.1515/fca-2016-0022}





\Bibitem{Podlubny:en}

\by Podlubny~I.

\book Fractional differential equations. An introduction to fractional derivatives, fractional differential equations, to methods of their solution and some of their applications

\yr 1999

\publ Academic Press

\publaddr  San Diego, CA

\zmath{0924.34008}

\serial Mathematics in Science and Engineering

\vol 198

\totalpages xxiv+340





\Bibitem{MathaiSaxenaHaubold:en}

\by Mathai~A.~M., Saxena~R.~K., Haubold~H.~J.

\book The $H$-function. Theory and Applications

\yr 2010

\publ Springer

\publaddr  Dordrecht

\totalpages xiv+268~pp \nofrills

\crossref{10.1007/978-1-4419-0916-9}

\zmath{1181.33001}









\Bibitem{KhushtovaVestn2}

\by Khushtova~F.~G.

\paper Cauchy problem for a parabolic equation with Bessel operator and Riemann--Liouville partial derivative

\jour Vestn. Samar. Gos. Tekhn. Univ., Ser. Fiz.-Mat. Nauki \rm [J. Samara State Tech. Univ., Ser. Phys. Math. Sci.]

\yr 2016

\vol 20

\issue 1

\pages 74--84

\crossref{10.14498/vsgtu1455}





\Bibitem{GorenfloLuchkoMainardi:en}

\by Gorenflo~R., Luchko~Y., Mainardi~F.

\paper Analytical properties and applications of the Wright  function

\jour Fract. Calc. Appl. Anal.

\yr 1999

\vol 2

\issue 4

\pages 383--414

\arxiv \href{https://arxiv.org/abs/math-ph/0701069}{math-ph/0701069}







\Bibitem{tikhonovUsl-1:en}

\by Tychonoff~A.

\paper Th\'eor\`emes d'unicit\'e pour l'\'equation de la chaleur

\jour Mat. Sb.

\yr 1935

\vol 42

\issue 2

\pages 199--216

\lang In French

\mathnet{http://mi.mathnet.ru/msb6410}

\zmath{https://zbmath.org/?q=an:0012.35501|61.1203.05}









\Bibitem{tikhonovUsl-2:en}

\by Tikhonov~A.~N.

\paper Uniqueness theorems for the heat equation

\inbook Collection of scientific works. In ten volumes

\volinfo Vol. 2: Mathematics. Part 2: Numerical mathematics 1956–1979. Mathematical physics 1933–1948. Edited by T.~A.~Sushkevich and  A.~V.~Gulin

\yr 2009

\publ Nauka

\publaddr Moscow

\pages 371--382

\zmath{1200.01055}

\lang In Russian and French





\Bibitem{Wright:en}

\by Wright~E.~M.

\paper The asymptotic expansion of the generalized hypergeometric function

\jour Proc. Lond. Math. Soc., II. Ser. 

\yr 1940

\vol 46

\issue 1

\pages 389--408

\crossref{10.1112/plms/s2-46.1.389}

\zmath{0025.40402|66.1243.01}











\end{thebibliography}



%\end{document}

 \newpage

































%%%%

\input{preamble}





\usepackage{rotating_pr}

\usepackage{desclist}

\usepackage{cmap}

\usepackage{textcomp}

\usepackage{nccboxes}

\usepackage{euscript}

\usepackage{mathrsfs} 

\usepackage{multicol}

\usepackage{nccfoots}

\usepackage{mathrsfs}

\usepackage{calrsfs}

\allowdisplaybreaks[4]

\usepackage{qrcode}

\usepackage{afterpage}

\usepackage{wasysym}

\usepackage{wrapfig}

\usepackage{slashbox}



\usepackage{multirow}

\newcommand{\specialcell}[2][c]{%

  \begin{tabular}[#1]{@{}c@{}}#2\end{tabular}}





\usepackage{pgf,pgfplots,tikz}

\usetikzlibrary{arrows,snakes,backgrounds,patterns,shapes.geometric,calc}



\nofiles



\oek % для генерации pdf, отдаваемого в oek.rsl.ru

%\samgtupubl % для генерации pdf, отдаваемого в типографию СамГТУ

%\pdfcrop % обрезка pdf для размещения на math-net.ru

%\draft % На корректуре





%\DeclareMathOperator{\tr}{tr}

%\DeclareMathOperator{\ink}{Ink}

%\DeclareMathOperator{\erf}{erf}







%\makeatletter

%\newcommand{\PersonRadRef}{\@langrustrue\def\refname{{\bf\small Избранные  труды \rupid{19071}{О.~А.~Репина} за последние 5 лет}}}

%\makeatother



\DeclareMathOperator{\tr}{tr}

\DeclareMathOperator{\arcsh}{arcsh}









\usepackage{listings}

\usepackage{color}

\definecolor{dkgreen}{rgb}{0,0.6,0}

\definecolor{gray}{rgb}{0.5,0.5,0.5}

\definecolor{mauve}{rgb}{0.58,0,0.82}



\renewcommand*{\lstlistingname}{Листинг}

\usepackage{caption}

%\DeclareCaptionFont{white}{\color{white}} 

\DeclareCaptionFormat{listing}{\parbox{\textwidth}{\centering \small  #1.~#3}}

\captionsetup[lstlisting]{format=listing}



\lstset{ %

  inputencoding=utf8,

  language=R,                     % the language of the code

  basicstyle=\footnotesize,       % the size of the fonts that are used for the code

  numbers=left,                   % where to put the line-numbers

  numberstyle=\tiny\color{gray},  % the style that is used for the line-numbers

  stepnumber=1,                   % the step between two line-numbers. If it's 1, each line

                                  % will be numbered

  numbersep=5pt,                  % how far the line-numbers are from the code

  backgroundcolor=\color{white},  % choose the background color. You must add \usepackage{color}

  showspaces=false,               % show spaces adding particular underscores

  showstringspaces=false,         % underline spaces within strings

  showtabs=false,                 % show tabs within strings adding particular underscores

  frame=single,                   % adds a frame around the code

  rulecolor=\color{black},        % if not set, the frame-color may be changed on line-breaks within not-black text (e.g. commens (green here))

  tabsize=2,                      % sets default tabsize to 2 spaces

  captionpos=t,                   % sets the caption-position to bottom

  breaklines=true,                % sets automatic line breaking

  breakatwhitespace=false,        % sets if automatic breaks should only happen at whitespace

  %title=\lstname,                 % show the filename of files included with \lstinputlisting;

                                  % also try caption instead of title

  keywordstyle=\color{blue},      % keyword style

  commentstyle=\color{dkgreen},   % comment style

  stringstyle=\color{mauve},      % string literal style

  %escapeinside={\%*}{*)},         % if you want to add a comment within your code

  morekeywords={*,in,+,-,=}            % if you want to add more keywords to the set

} 







\begin{document}







\setcounter{page}{601}

\input{specialpages.tex} 







\input{./01du/partition.tex}

\input{./02mdtt/partition.tex}

\input{./03matmod/partition.tex}

\input{./04short/partition.tex}



%\input{./polieces.tex}



\end{document}



\newpage



$\;$





\chead[\fancyplain{}{}]{\fancyplain{}{}}

  \rhead[]{\fancyplain{}{}}

  \lhead[\fancyplain{}{}]{}

  \cfoot[]{}

  \lfoot[]{}

  \rfoot[]{}

\renewcommand{\plainheadrulewidth}{0pt}

\renewcommand{\headrulewidth}{0pt}





  \pagestyle{fancyplain}



  \thispagestyle{plain}





\newpage



$\;$





\chead[\fancyplain{}{}]{\fancyplain{}{}}

  \rhead[]{\fancyplain{}{}}

  \lhead[\fancyplain{}{}]{}

  \cfoot[]{}

  \lfoot[]{}

  \rfoot[]{}

\renewcommand{\plainheadrulewidth}{0pt}

\renewcommand{\headrulewidth}{0pt}





  \pagestyle{fancyplain}



  \thispagestyle{plain}





%\newpage

%

%$\;$

%

%

%\chead[\fancyplain{}{}]{\fancyplain{}{}}

%  \rhead[]{\fancyplain{}{}}

%  \lhead[\fancyplain{}{}]{}

%  \cfoot[]{}

%  \lfoot[]{}

%  \rfoot[]{}

%\renewcommand{\plainheadrulewidth}{0pt}

%\renewcommand{\headrulewidth}{0pt}

%

%

%  \pagestyle{fancyplain}

%

%  \thispagestyle{plain}







\end{document}





\Russian



\addtocontents{toc}{\bigskip\par\hspace{-7mm}}

\addtocontents{etoc}{\bigskip\par\hspace{-7mm}}



$\;$



\hfill \qrcode[hyperlink,height=.8in]{http://doi.org/10.14498/vsgtu1651}

\vspace{-1.5in}





%$$\quad $$ \vskip.6in \par \hskip-7mm  {\huge  Краткие сообщения}

$$\quad $$ \vskip.6in \par \hskip-7mm  {\huge  Short Communications}



%$$\quad $$ \vskip.6in \par \hskip-7mm  {\huge  Mathematical Modeling,\\ Numerical Methods \\[1mm] and Software Complexes}



\addtocontents{toc}{{\large\bf Краткие сообщения}\bigskip}

\addtocontents{etoc}{{\large\bf Short Communications}\bigskip}



\vspace{-5mm}



%\Partition{Математическое моделирование}{Mathematical Modeling}



%\input{preamble}

%\usepackage{samgtu-name}

%

%

%

\renewcommand{\RuCitation}{\noindent\textbf{\CYRO\cyrb\cyrr\cyra\cyrz\cyre\cyrc\ \cyrd\cyrl\cyrya\ \cyrc\cyri\cyrt\cyri\cyrr\cyro\cyrv\cyra\cyrn\cyri\cyrya\\}\selectlanguage{english}{\EnAuthorInContents} {\EnTitleInContents}, \textit{\EnJournalName}\/ [\ParallelJournalName],  \JournalYear, vol.~\JournalVolume, no.~\JournalNumber,  pp.~\thefirstpage--\thelastpage. \doiurl. \selectlanguage{russian}\smallskip}



\renewcommand{\EnCitation}{\noindent\textbf{Please cite this article in press as:\\} {\EnAuthorInContents} {\EnTitleInContents}, \textit{\EnJournalName}\/ [\ParallelJournalName],  \JournalYear, vol.~\JournalVolume, no.~\JournalNumber,  pp.~\thefirstpage--\thelastpage.  \midoiurl. }

%

\vsgtu{1651} %registration number

\FirstPage{735}

\LastPage{749}

%

%

%

%\pdfcrop

%\draft

%

\ShortCommunication

%

%\begin{document}

%

%

%

%$\;$

%\vspace{1mm}

%

%\hfill \qrcode[hyperlink,height=.8in]{http://doi.org/10.14498/vsgtu1651}

%\vspace{-.8in}





\renewcommand{\baselinestretch}{0.98}



\udc{532.51, 517.958:531.3-324}

\msc{76F02, 76F45, 76M45, 76R05, 76U05}



\rutitle{Динамические равновесия неизотермической жидкости}

\entitle{Dynamic equilibria of a nonisothermal fluid}



\ruauthor{Е.~Ю.~Просвиряков\affil{1,2}}{П\,р\,о\,с\,в\,и\,р\,я\,к\,о\,в~Е.~Ю.}

\enauthor{E.~Yu.~Prosviryakov\affil{1,2}}{P\,r\,o\,s\,v\,i\,r\,y\,a\,k\,o\,v~E.~Yu.}





\ruaffil{\ruaffid{2665}{Институт машиноведения УрО РАН,\\

Россия, 620049, Екатеринбург, ул. Комсомольская, 34}.

\and

\ruaffid{7370}{Уральский федеральный университет \\ им.~первого Президента России Б.~Н.~Ельцина, \\ 

Россия, 620002, Екатеринбург, ул. Мира, 19}.}



\enaffil{\enaffid{2665}{Institute of Engineering Science, Urals Branch, Russian Academy of Sciences,\\

34, Komsomolskaya st., Ekaterinburg, 620049, Russian Federation}.

\and

\enaffid{7370}{Ural Federal University named after the First President of Russia B.~N.~Yeltsin, \\ 19, Mira st., Ekaterinburg, 620002, Russian Federation}.

}



\enauthorsdetails{1}{

\fullauthor{\enpid{41426}{Evgeny Yu. Prosviryakov}}

\CAuthor 

\orcid{0000-0002-2349-7801} 

\regalia{Dr. Phys. \& Math. Sci., Professor} 

\position{Head of Sector} 

\dept{Sector of Nonlinear Vortex Hydrodynamics}

\email[br]{evgen_pros@mail.ru}

}



\ruauthorsdetails{1}{

\fullauthor{\rupid{41426}{Евгений Юрьевич Просвиряков}} 

\CAuthor 

\orcid{0000-0002-2349-7801}

\regalia{доктор физико-математических наук} 

\position{заведующий сектором} 

\dept{сектор нелинейной вихревой гидродинамики}

\email{evgen_pros@mail.ru}

}





\enkeywords{dynamic equilibrium,  rotating fluid flow, exact solution, Boussinesq approximation, counterflows, increased velocities}



\rukeywords{динамическое  равновесие, вращающийся поток жидкости, точное решение, аппроксимация Буссинеска, противотечение, увеличенные скорости}



\ruabstract{В рамках точности приближения Буссинеска рассмотрены стационарные динамические равновесия вращающейся массы неизотермической жидкости. Показано, что в этом случае в жидкости наблюдается конечное число противотечений и усиление скоростей по сравнению с заданными на границе значениями, а также формирование зон положительного и отрицательного давления и температуры. }



\enabstract{In this paper, stationary dynamic equilibria of the rotating mass of a nonisothermal fluid are discussed within the accuracy limits of the Boussinesq approximation. It is demonstrated that, in this case, a fluid exhibits a finite number of counterflows, higher values of velocities than those specified on the boundary and the formation of zones of positive and negative pressures and temperatures.}



\newdate{privalovaa}{07}{10}{2018}

\date{\displaydate{privalovaa}}

\newdate{privalovab}{10}{11}{2018}

\revisiondate{\displaydate{privalovab}}

\newdate{privalovac}{12}{11}{2018}

\accepteddate{\displaydate{privalovac}}



\newdate{privalovae}{13}{12}{2018}

\renewcommand{\JournalDataOnline}{\displaydate{privalovae}}





%\ForPeerReview

\makeentitle





\Section[n]{Introduction}

The study of the motion of a rotating self-gravitating fluid was started as long ago as by Isaak Newton~[\citen{ulg:bib:Newtoni},~\citen{ulg:bib:Turnbull}], and it is currently far from being completed. Newton proved to be the first to endeavour to obtain, from the motion equations, a solution simulating the shape of the Earth. Newton’s pioneering study gave rise to numerous investigations in this field of mechanics and physics and a great impetus to the development of the mathematics used practically in all the sections of mathematics. The development of the potential theory, the creation of the theory of special functions and the formation of the rapidly developing theory of integrable systems should be mentioned first.



Practically all the most prominent scientists, such as Appel, Basset, Betti, Gagen, Helmholtz, Dedekind, Dirichlet, Kirchhoff, Lipschitz, Lame, McLoren, Poincare, Riemann, Chandrasekhar, Jacoby, took part in the construction of the theory~[\citen{ulg:bib:Poincare}--\citen{ulg:bib:Dirichlet}]. This list of discoverers of exact solutions can be extended~[\citen{ulg:bib:Poincare}--\citen{ulg:bib:Dirichlet}]. The history of discovering the solutions can be found in the monographs and surveys~[\citen{ulg:bib:Poincare}--\citen{ulg:bib:Dirichlet}]. It should be noted that the scientific breakthrough in the simulation of rotating equilibrium figures was made by the Russian scientists Steklov, Lyapunov, Ovsyannikov and Zel'dovich.



Note that the development of the theory of motion stability described by systems of ordinary differential equations was initiated by studying the stability of equilibrium figures and the motion of an elastic body partially filled with a fluid~[\citen{ulg:bib:Lyapunov1},~\citen{ulg:bib:Lyapunov2}]. It is the work in these fields that offered a formulation to the well-known Lyapunov theorems of stability and instability, among other things, from the first approximation~[\citen{ulg:bib:Lyapunov1},~\citen{ulg:bib:Lyapunov2}]. Original research has currently been done on the stability of Dirichlet and Jacobi ellipsoids~\cite{ulg:bib:Borisov}.



Ovsyannikov’s exact gas-hydrodynamic solution became the first solution    si\-mulating gas cloud motion in the class of linear velocities~\cite{ulg:bib:Borisov}. Afterwards models of Dyson and Fujimoto were proposed~\cite{ulg:bib:Borisov}.



There is another substantial model, which influenced the development of the theory of rotating fluids. It is the generalized solid proposed for research by Arnold~\cite{ulg:bib:Arnold}. In the survey made by Dolzhansky~\cite{ulg:bib:Dolzhansky} there is a detailed analysis of this approach with some generalizations for different classes of fluids.



As a rule, dynamic fluid equilibria are simulated with the application of two mathematical models. One model is Lagrangian fluid simulation, i.e. the use of the ordinary differential equations of Lagrange, Hamilton and Jacoby to describe the relative equilibrium (solid-state) of the continuum. The other model consists in using the Euler equations with potential forces to analyze the structures of a relative fluid equilibrium. Thus, in the simulation of dynamic fluid equilibria no account is taken of dissipation and temperature effect. In other words, an isothermal perfect incompressible fluid is dealt with. 

	

It is apparent that this deficiency in studying fluid flows can be compensated if this problem is analyzed with the application of the Oberbeck--Boussinesq system of equations. Note that it is not only the Boussinesq approximation that has not been used before to describe dynamic fluid equilibria, but also the characteristic scale of solutions. This statement needs to be clarified. Despite Chandrasekhar~\cite{ulg:bib:Chandrasekhar} used the Boussinesq approximation to simulate astrophysical objects, he used the virial method of physical field expansion for solving his problems. The application of this body of mathematics enabled one to reduce the Oberbeck--Boussinesq system to a system of ordinary differential equations and to use the classical problem statements when interpreting results. Similar reasoning holds for~[\citen{ulg:bib:Arnold},~\citen{ulg:bib:Dolzhansky}].

	

There are studies dealing with large-scale fluid or gas motions~\cite{ulg:bib:Chandrasekhar}, when it is necessary to take into account the inhomogeneity of the gravitational field, and with scales for which surface tension forces are principally important, e. g. the Hadamard-Rybczynski problem~[\citen{ulg:bib:Hadamard},~\citen{ulg:bib:Rybczynski}]. In both extreme cases, the principal equilibrium shape is spherical. The consideration of rotation was always viewed as some disturbance of the principal flow, however it was not viewed as a forming factor for equilibrium structures.



Note that the very possibility of the existence of equilibrium figures in the region of intermediate scales, when neither the inhomogeneity of the gravitational field nor the surface tension forces are essential, does not seem to have been studied. At any rate, the authors are unfamiliar with works viewing the problem this way.



Thus, we intend to ascertain the possible existence of equilibrium structures in the region of the so-called intermediate scales, when the Boussinesq approximation is obviously applicable to the simulation of a heat-conducting viscous incompres\-sible fluid.





\smallskip











\Section{Problem statement}

In the Boussinesq approximation~[\citen{ulg:bib:Chandrasekhar},~\citen{ulg:bib:Aristov1}] in the Cartesian coordinates (the $Oz$  axis is directed upwards), the system of equations describing the heat convection of a viscous incompressible fluid are written as

\begin{gather}

\frac{\partial \bf{V}}{\partial t}+\left(\bf{V}\cdot\nabla\right)\bf{V}=-\nabla \textit{P} + \nu \Delta \bf{V}+g \beta \textit{T} \bold{k},

\\    

      \label{eq:Aristov_Prosviryakov:NSE}

\frac{\partial T}{\partial t}+\bf{V}\cdot\nabla \textit{T}=\chi \Delta \textit{T},

\\     

\nabla \cdot \bf{V}=0.

 \end{gather}

The system of equations~\eqref{eq:Aristov_Prosviryakov:NSE} involves the following symbols: $\bf{V}=$ $\left(V_x , V_y , V_z\right)$  is the flow velocity vector;  $P$  is the pressure deviation from hydrostatic, which is related to the constant average fluid density  $\rho;$  $T$ is the deviation from the average temperature;   $\nu$  and $\chi$ are the coefficients of kinematic viscosity and thermal diffusivity of the fluid, respectively; $\beta$ is the temperature coefficient of fluid volumetric expansion; $\bold{k}$  is the unit vector directed vertically upwards, $\nabla$  is the Hamilton operator,  $\Delta$ is the Laplace operator (Laplacian).



The solution of system~\eqref{eq:Aristov_Prosviryakov:NSE} is sought in the approximation to the large scale. For large-scale flows, the characteristic horizontal scale greatly exceeds in magnitude the characteristic dimension during the entire fluid motion. Free surface curvature can in this scale be neglected, the vertical velocity $V_z$  being assumed zero. Thus, system~\eqref{eq:Aristov_Prosviryakov:NSE} is transformed to the following system:

\begin{gather}

\frac{\partial V_x}{\partial t}+V_x \frac{\partial V_x}{\partial x}+V_y \frac{\partial V_x}{\partial y}=-\frac{\partial P}{\partial x} + \nu \Bigl(\frac{\partial^2 V_x}{\partial x^2}+\frac{\partial^2 V_x}{\partial y^2}+\frac{\partial^2 V_x}{\partial z^2}\Bigr),

\\       

\frac{\partial V_y}{\partial t}+V_x \frac{\partial V_y}{\partial x}+V_y \frac{\partial V_y}{\partial y}=-\frac{\partial P}{\partial y} + \nu \Bigl(\frac{\partial^2 V_y}{\partial x^2}+\frac{\partial^2 V_y}{\partial y^2}+\frac{\partial^2 V_y}{\partial z^2}\Bigr),

\\

       \label{eq:Aristov_Prosviryakov:system}

\frac{\partial P}{\partial z}=g\beta T,

\\

\frac{\partial T}{\partial t}+V_x \frac{\partial T}{\partial x}+V_y \frac{\partial T}{\partial y}=\chi \Bigl(\frac{\partial^2 T}{\partial x^2}+\frac{\partial^2 T}{\partial y^2}+\frac{\partial^2 V_x}{\partial z^2}\Bigr),

\\

\frac{\partial V_x}{\partial x}+\frac{\partial V_y}{\partial y}=0.

 \end{gather}

 

Obviously, system \eqref{eq:Aristov_Prosviryakov:system} is nonlinear and overdetermined, four unknown functions of velocity and temperature being required from five equations. The hydrodynamic and temperature fields are computed as~\cite{ulg:bib:Aristov1}

\begin{gather}

V_x=U+xu_1+y u_2,

\quad 

V_y=V+xv_1+y v_2,

\\

       \label{eq:Aristov_Prosviryakov:solution}

P=P_0+xP_1+yP_2+\frac{x^2}{2}P_{11}+\frac{y^2}{2}P_{22}+xyP_{12},

\\

T=T_0+xT_1+yT_2+\frac{x^2}{2}T_{11}+\frac{y^2}{2}T_{22}+xyT_{12}.     

 \end{gather}

 

For horizontal coordinates, all the functions in the formulae of system \eqref{eq:Aristov_Prosviryakov:solution}  depend on the variable $z$  ant time $t.$ Substituting relation \eqref{eq:Aristov_Prosviryakov:solution} into system \eqref{eq:Aristov_Prosviryakov:system}, we arrive at the system

\begin{gather}

\hat{L}U+Uu_1+Vu_2+P_1=0,

\\

\hat{L}u_1+u^{2}_1+v_1u_2+P_{11}=0,

\\

\hat{L}u_2+u_1u_2+u_2v_2+P_{12}=0,

\\

\hat{L}V+Uv_1+Vv_2+P_{2}=0,

\\

\hat{L}v_1+u_1v_1+v_1v_2+P_{12}=0,

\\

\hat{L}v_2+u_2v_1+v^{2}_2+P_{22}=0,

\\

       \label{eq:Aristov_Prosviryakov:main system }

\hat{M}T_0+UT_1+VT_2-\chi\left(T_{11}+T_{22}\right)=0,

\\

\hat{M}T_1+UT_{11}+VT_{12}+u_1T_1+v_1T_2=0,

\\

\hat{M}T_2+UT_{12}+VT_{22}+u_2T_1+v_2T_2=0,

\\

\hat{M}T_{11}+2u_1T_{11}+2v_1T_{12}=0,

\\

\hat{M}T_{22}+2u_2T_{12}+2v_2T_{22}=0,

\\

\hat{M}T_{12}+u_1T_{12}+u_2T_{11}+v_1T_{22}+v_2T_{12}=0,

\\

u_1+v_2=0,

\\

\frac{\partial P_0}{\partial z}=g\beta T_0,

\quad 

\frac{\partial P_1}{\partial z}=g\beta T_1,

\quad 

\frac{\partial P_2}{\partial z}=g\beta T_2,

\\

\frac{\partial P_{11}}{\partial z}=g\beta T_{11},

\quad 

\frac{\partial P_{12}}{\partial z}=g\beta T_{12},

\quad 

\frac{\partial P_{22}}{\partial z}=g\beta T_{22}.

 \end{gather}

Here, the following differential parabolic operators in partial derivatives are intro\-duced:

$$\hat{L}=\frac{\partial }{\partial t}-\nu\frac{\partial^2 }{\partial z^2},\quad 

\hat{M}=\frac{\partial }{\partial t}-\chi\frac{\partial^2 }{\partial z^2}. 

$$





\smallskip





\Section{Analyzing the solvability of system}

In what follows we do not discuss the solvability of system \eqref{eq:Aristov_Prosviryakov:main system } as a whole, but restrict our analysis to one special case. We will seek a solution to the system of equations when the boundaries of the fluid layer rotate. In~\cite{ulg:bib:Aristov2} the validity of equalities for the velocity gradients in the simulation of axisymmetric flows within class \eqref{eq:Aristov_Prosviryakov:solution} was demonstrated,

\begin{gather}

       \label{eq:Aristov_Prosviryakov:sol1}

u_1=0,

\quad 

v_2=0,       

\quad 

u_2=-v_1.       

 \end{gather}

Next the lower boundary of the layer  $z=0$ is assumed to be immobile. The upper boundary describable by the equation $z=h$  rotates with constant angular velocity  $\Omega$. Using equalities \eqref{eq:Aristov_Prosviryakov:sol1}, we obtain the following solutions to system~\eqref{eq:Aristov_Prosviryakov:main system }, which determine a number of functions in class \eqref{eq:Aristov_Prosviryakov:solution} for the coordinates  $x$ and  $y$:

\begin{gather}

u_2=\frac{\Omega z}{h}, 

\quad 

v_1=-\frac{\Omega z}{h},

\quad 

P_{12}=0, 

\\

\label{eq:Aristov_Prosviryakov:sol2}

 P_{11}=P_{22}=-u_2v_1=\frac{\Omega^2 z^2}{h^2},

\quad 

T_{12}=0, 

\\ g\beta T_{11}=g\beta T_{22}=\frac{2\Omega^2 z}{h^2}.

\end{gather}



For the homogeneous summands in expressions \eqref{eq:Aristov_Prosviryakov:solution}, due to solution \eqref{eq:Aristov_Prosviryakov:sol2}, system \eqref{eq:Aristov_Prosviryakov:main system } yields two closed subsystems. First system has of the form

\begin{gather}

\hat{L}U+\frac{\Omega z}{h}V+P_1=0, 

\\

\hat{L}V-\frac{\Omega z}{h}U+P_{2}=0,

\\

\label{eq:Aristov_Prosviryakov:sys1}

\hat{M}T_1+\frac{2\Omega^2 z}{g\beta h^2}U-\frac{\Omega z}{h}T_2=0, 

\\ \hat{M}T_2+\frac{2\Omega^2 z}{g\beta h^2}V+\frac{\Omega z}{h}T_1=0, 

\\

\frac{\partial P_1}{\partial z}=g\beta T_1, 

\quad  

\frac{\partial P_2}{\partial z}=g\beta T_2.

\end{gather}



Second system has of the form

\begin{gather}

\label{eq:Aristov_Prosviryakov:sys2}

\hat{M}T_0+UT_1+VT_2=\frac{4\chi\Omega^2 z}{g\beta h^2},

\\

\frac{\partial P_0}{\partial z}=g\beta T_0.

\end{gather}



Note that the integration of system \eqref{eq:Aristov_Prosviryakov:sys2} is possible only after the integration of system \eqref{eq:Aristov_Prosviryakov:sys1}. In order to find a solution to system \eqref{eq:Aristov_Prosviryakov:sys1}, we introduce the following complex functions:

$$

W=U+iV,\quad 

S=P_1+iP_2,\quad 

\Theta=g\beta\left(T_1+iT_2\right).$$

Here  $i$ is an imaginary unit. Thus, system \eqref{eq:Aristov_Prosviryakov:sys1} is in this case transformed to the following form:

 \begin{gather}

(\hat{L}-i\Omega z) W+S=0,

\\

       \label{eq:Aristov_Prosviryakov:sys1complex}

(\hat{M}+i\Omega z) \Theta+2\Omega^2 z W=0,       

\\

\frac{d S}{dz}=\Theta.      

 \end{gather}

 

To reduce the order of the system, the following transformation is made:

 \begin{gather}

       \label{eq:Aristov_Prosviryakov:trans}

S=S_1+i\Omega z W,

\quad 

\Theta=\Theta_1+2i\Omega W.    

\end{gather}

Substituting \eqref{eq:Aristov_Prosviryakov:trans} into system \eqref{eq:Aristov_Prosviryakov:sys1complex}, we obtain the equalities

 \begin{gather}

\nu\frac{d^2 W}{dz^2}+S_1=0,

\\

       \label{eq:Aristov_Prosviryakov:trans2}

\chi\frac{d^2 \Theta_1}{d z^2}+i\Omega \Bigl(z \Theta_1- 2\frac{\chi}{\nu}S_1\Bigr)=0,       

\\

\frac{d S_1}{dz}=\Theta_1-i\Omega z^2 \frac{d}{dz}\Bigl(\frac{W}{z}\Bigr). 

 \end{gather}

 

For further transformations, we use the differential identity

$$

\frac{d^2 W}{dz^2}=\frac{1}{z}\frac{d}{dz}\left(z^2\frac{d}{dz}\Bigl(\frac{W}{z}\Bigr)\right).

$$

This identity holds for any twice-differentiable function. Consequently, by differentiating the last equation in system \eqref{eq:Aristov_Prosviryakov:trans2}, we arrive at a system for determining the functions  $\Theta_1$  and  $S_1,$

 \begin{gather}

       \label{eq:Aristov_Prosviryakov:trans3}

\chi\frac{d^2 \Theta_1}{d z^2}+i\Omega \left(z \Theta_1- 2\frac{\chi}{\nu}S_1\right)=0,

\quad 

\nu\frac{d^2 S_1}{dz^2}=i\Omega z S_1 +\nu \frac{d \Theta_1}{dz}.          

 \end{gather}

 

Taking into account that equations \eqref{eq:Aristov_Prosviryakov:trans2} offer an explicit expression for  $W$, we obtain an integral of the initial system with automatic order reduction, system \eqref{eq:Aristov_Prosviryakov:trans3} being transformed to the operator equation

$$

\left[\Bigl(\nu\frac{d^2}{dz^2}-i\Omega z\Bigr)\Bigl(\chi\frac{d^2}{dz^2}+i\Omega z\Bigr)-2i\Omega\chi\frac{d}{dz}\right]\Theta_1=0.

$$

Simplifying the latter equality, we derive the following fourth-order ordinary differential equation with complex coefficients:

 \begin{gather}

       \label{eq:Aristov_Prosviryakov:trans4}

\Bigl[\nu\chi\frac{d^4}{dz^4}+i\Omega\left(\nu-\chi\right)\frac{d^2}{dz^2}z+\Omega^2 z^2\Bigr]\Theta_1=0.    

 \end{gather}



\smallskip







\Section{Solving equation when the dissipative coefficients coincide}

 We assume in equation \eqref{eq:Aristov_Prosviryakov:trans4} that $\nu=\chi$  (the Prandtl number equals one). In this case, equation \eqref{eq:Aristov_Prosviryakov:trans4} acquires the form

  \begin{gather}

       \label{eq:Aristov_Prosviryakov:Pr1}

\frac{d^4 \Theta_1}{dz^4}+\left(36\sigma^3 z\right)^2 \Theta_1=0,    

 \end{gather}

where  $\sigma=\frac{1}{6}\sqrt[3]{\frac{\Omega}{6\nu h}}$ is the dispersion relation. For the convenience of computations, in order to illustrate the physical meaning of the solutions and to make the notation concise, we introduce a new variable  $Z=\frac{z}{h \sigma}.$ The general solution of equation \eqref{eq:Aristov_Prosviryakov:Pr1} is written as

\begin{multline}

\label{eq:Aristov_Prosviryakov:soltemp}

\Theta_1={^{3}_{0}F} \Bigl[\Bigl\{\frac{1}{2}, \frac{2}{3}, \frac{5}{6}\Bigr\},-Z^6\Bigr]C_1 +

{^{3}_{0}F} 

\Bigl[\Bigl\{\frac{2}{3}, \frac{5}{6}, \frac{7}{6}\Bigr\},-Z^6\Bigr] C_2 Z +

\\

+

{^{3}_{0}F} \Bigl[\Bigl\{\frac{5}{6}, \frac{7}{6}, \frac{4}{3}\Bigr\},-Z^6\Bigr] C_3 Z^2  +

{^{3}_{0}F} \Bigl[\Bigl\{\frac{7}{6}, \frac{4}{3},\frac{3}{2}\Bigr\},-Z^6\Bigr] C_4 Z^3  . 

\end{multline}

Here,  ${^{3}_{0}F}=F\left(a; b; Z\right)=F\left(a_1, a_2, \ldots, a_p; b_1, b_2, \ldots, b_q; Z\right)$ is a generalized hypergeometric function~\cite{ulg:bib:Polyanin}. The choice of using the generalized hypergeometric function for writing the solutions is dictated by notation conciseness. Note that there is another form of writing the solution of equation \eqref{eq:Aristov_Prosviryakov:Pr1}, which is based on using the Mittag--Leffler functions~\cite{ulg:bib:Polyanin}. Note that each particular solution \eqref{eq:Aristov_Prosviryakov:soltemp} is a function of a real variable and that the general solution is considered in the complex plane. Thus, at least one integration constant is a complex number.



We now turn to finding the generalized pressure $S_1$  from the differential equation \eqref{eq:Aristov_Prosviryakov:trans2} 

 \begin{gather}

       \label{eq:Aristov_Prosviryakov:eqpr}

S_1=-\frac{i}{2\sigma}\frac{d^2 \Theta_1}{d z^2}+\frac{z}{2}\Theta_1.

\end{gather}



The substitution of solution \eqref{eq:Aristov_Prosviryakov:soltemp} into equation \eqref{eq:Aristov_Prosviryakov:eqpr} gives

\begin{multline}

\label{eq:Aristov_Prosviryakov:solstress}

S_1= {^{3}_{0}F} 

\Bigl[\Bigl\{\frac{1}{2}, \frac{2}{3}, \frac{5}{6}\Bigr\},-Z^6\Bigr] \frac{C_1}{2}Z+

{^{3}_{0}F} 

\Bigl[\Bigl\{\frac{2}{3}, \frac{5}{6}, \frac{7}{6}\Bigr\},-Z^6\Bigr] \frac{C_2}{2}Z^2+

\\

+{^{3}_{0}F} 

\Bigl[\Bigl\{\frac{5}{6}, \frac{7}{6}, \frac{4}{3},\Bigr\},Z\Bigr] \sqrt[3]{36\sigma^2} \frac{2\cdot 5C_3}{6!}Z^3+

{^{3}_{0}F} 

\Bigl[\Bigl\{\frac{7}{6}, \frac{4}{3}, \frac{3}{2}\Bigr\},Z\Bigr] \sigma \frac{2\cdot 5C_4}{6!}Z^4-

\\

-\frac{1}{72\sigma}i\left\{ {^{3}_{0}F} 

\Bigl[\Bigl\{\frac{5}{6}, \frac{7}{6}, \frac{4}{3}\Bigr\},Z\Bigr] \sqrt[3]{676\sigma^2} C_3 Z^3 

-{^{3}_{0}F} 

\Bigl[\Bigl\{\frac{7}{6}, \frac{4}{3}, \frac{3}{2}\Bigr\},Z\Bigr] 6 \sigma C_4 Z+

\right.

\\

+{^{3}_{0}F} 

\Bigl[\Bigl\{\frac{5}{6}, \frac{7}{6}, \frac{4}{3}\Bigr\},Z\Bigr] 3 \sigma^2 C_1 Z^4 +

\\

+{^{3}_{0}F} 

\Bigl[\Bigl\{\frac{5}{3}, \frac{11}{6}, \frac{13}{6},\Bigr\},Z\Bigr] \frac{2^4\cdot 3^3\sigma^2}{7!}\sqrt[3]{\frac{3\sigma}{4}} C_2 Z^5

+

\\

+{^{3}_{0}F} 

\Bigl[\Bigl\{\frac{5}{3}, \frac{11}{6}, \frac{13}{6},\Bigr\},Z\Bigr] \frac{2^4\cdot 3^4\sigma^2}{7!}\sqrt[3]{6\sigma} C_2 Z^5+

\\

+{^{3}_{0}F} 

\Bigl[\Bigl\{\frac{11}{6}, \frac{13}{6}, \frac{7}{3} \Bigr\},Z\Bigr] \frac{2^2\cdot 3^4\sigma^2}{7!}\sqrt[3]{\frac{9\sigma^2}{2}} C_3 Z^6+

\\

+{^{3}_{0}F} \Bigl[\Bigl\{\frac{13}{6}, \frac{7}{3}, \frac{5}{2} \Bigr\},Z\Bigr] \frac{2\cdot 5\cdot 11\sigma^3}{7!} C_4 Z^7-

\\

-{^{3}_{0}F} \Bigl[\Bigl\{\frac{5}{2}, \frac{8}{3}, \frac{17}{6} \Bigr\},Z\Bigr] \frac{2^7\cdot 3^3\cdot 7\sigma^4}{11!} C_1 Z^{10}-

\\

-{^{3}_{0}F} \Bigl[\Bigl\{\frac{8}{3}, \frac{17}{6}, \frac{19}{6} \Bigr\},Z\Bigr] \frac{2^9\cdot 3^6\sigma^4}{13!}\sqrt[3]{\frac{3\sigma}{4}} C_2 Z^{11}-

\\

-{^{3}_{0}F} \Bigl[\Bigl\{\frac{17}{6}, \frac{19}{6}, \frac{10}{3} \Bigr\},Z\Bigr] \frac{2^8\cdot 3^5\cdot 5\sigma^4}{14!}\sqrt[3]{\frac{9\sigma^2}{2}} C_3 Z^{12}-

\\

\left.

-{^{3}_{0}F} \Bigl[\Bigl\{\frac{19}{6}, \frac{10}{3}, \frac{7}{2} \Bigr\},Z\Bigr] \frac{2^7\cdot 3^2\cdot 5\cdot 11\sigma^5}{14!} C_4 Z^{13}

\right\}.

\end{multline}





Complex velocity $W$  can be expressed in two ways. Firstly, it can be obtained by integrating the first equation in system \eqref{eq:Aristov_Prosviryakov:sys1complex}, but it is a second-order equation. Hence, we obtain a solution depending on six integration constants, whereas system \eqref{eq:Aristov_Prosviryakov:trans2} is a fifth-order equation system. Thus our problem is complicated by additional analysis of overdetermination. A correct expression for velocity $W$  can be obtained by the integration of the third equation in system \eqref{eq:Aristov_Prosviryakov:sys1complex},

\begin{multline}

\label{eq:Aristov_Prosviryakov:solvelo}

\Omega W=C_5 Z+ \left({^{3}_{0}F} \Bigl[\Bigl\{\frac{2}{3}, \frac{5}{6}, \frac{7}{6} \Bigr\},Z\Bigr]-1\right) \frac{i}{2} \sqrt[3]{\frac{\sigma}{36}} C_2 Z-

\\

-{^{4}_{1}F} \Bigl[\Bigl\{-\frac{1}{6} \Bigr\}, \Bigl\{\frac{1}{2},\frac{2}{3}, \frac{5}{6}, \frac{5}{6} \Bigr\},Z\Bigr]\frac{2i \Gamma\left(-\frac{1}{6}\right)}{4!\Gamma\left(\frac{5}{6}\right)} C_1 Z^2+

\\

+{^{4}_{1}F} \Bigl[\Bigl\{-\frac{1}{6} \Bigr\},\Bigl\{\frac{5}{6},\frac{7}{6}, \frac{4}{3}, \frac{3}{2} \Bigr\},Z\Bigr]\frac{2\cdot 5\Gamma\left(-\frac{1}{6}\right)}{6!\Gamma\left(\frac{5}{6}\right)} C_4+

\\

+{^{4}_{1}F} \Bigl[\Bigl\{\frac{1}{6} \Bigr\},\Bigl\{\frac{5}{6},\frac{7}{6}, \frac{7}{6}, \frac{4}{3} \Bigr\},Z\Bigr]\frac{2\cdot 5i \Gamma\left(\frac{1}{6}\right)}{6!\Gamma\left(\frac{7}{6}\right)} \sqrt[3]{\frac{\sigma^2}{6}} C_3 Z^2+

\\

+{^{4}_{1}F} \Bigl[\Bigl\{\frac{1}{3} \Bigr\},\Bigl\{\frac{7}{6},\frac{4}{3}, \frac{4}{3}, \frac{3}{2} \Bigr\},Z\Bigr]\frac{2^4\cdot 3\cdot 5\cdot 7i \sigma \Gamma\left(\frac{1}{3}\right)}{9!\Gamma\left(\frac{4}{3}\right)} C_4 Z^3-

\\

-{^{4}_{1}F} \Bigl[\Bigl\{\frac{1}{3} \Bigr\},\Bigl\{\frac{4}{3},\frac{3}{2}, \frac{5}{3}, \frac{11}{6} \Bigr\},Z\Bigr]\frac{2^2\cdot 5\sigma \Gamma\left(\frac{1}{3}\right)}{6!\Gamma\left(\frac{4}{3}\right)} C_1 Z^3-

\\

-{^{4}_{1}F} \Bigl[\Bigl\{\frac{1}{2} \Bigr\},\Bigl\{\frac{3}{2},\frac{5}{3}, \frac{11}{6}, \frac{13}{6} \Bigr\},Z\Bigr]\frac{\sigma }{4!}\sqrt[3]{\frac{\sigma}{36}} C_2 Z^4-

\\

-{^{4}_{1}F} \Bigl[\Bigl\{\frac{2}{3} \Bigr\},\Bigl\{\frac{5}{3},\frac{11}{6}, \frac{13}{6}, \frac{7}{3} \Bigr\},Z\Bigr] \frac{2^4\cdot 3^2\cdot 7\sigma \Gamma\left(\frac{2}{3}\right)}{9!\Gamma\left(\frac{5}{3}\right)}\sqrt[3]{\frac{\sigma}{6}} C_3 Z^5-

\\

-{^{4}_{1}F} \Bigl[\Bigl\{\frac{5}{6} \Bigr\},\Bigl\{\frac{3}{2},\frac{5}{3}, \frac{11}{6}, \frac{11}{6} \Bigr\},Z\Bigr] \frac{2^3\cdot 3^2\cdot 7i\sigma^2 \Gamma\left(\frac{5}{6}\right)}{9!\Gamma\left(\frac{11}{6}\right)} C_1 Z^6-

\\

-{^{4}_{1}F} \Bigl[\Bigl\{\frac{5}{6} \Bigr\},\Bigl\{\frac{11}{6},\frac{13}{6}, \frac{7}{3}, \frac{5}{2} \Bigr\},Z\Bigr] \frac{2^3\cdot 5^2\cdot 11\cdot 13\cdot 83\sigma^2 \Gamma\left(\frac{5}{6}\right)}{13! \Gamma\left(\frac{11}{6}\right)} C_4 Z^6-

\\ 

-{^{4}_{1}F} \Bigl[\Bigl\{\frac{7}{6} \Bigr\},\Bigl\{\frac{11}{6},\frac{13}{6}, \frac{13}{6}, \frac{7}{3} \Bigr\},Z\Bigr] \frac{2\cdot 3^2i\sigma^2 \Gamma\left(\frac{7}{6}\right)}{9! \Gamma\left(\frac{13}{6}\right)}\sqrt[3]{\frac{\sigma^2}{6}} C_3 Z^8-

\\

-{^{4}_{1}F} \Bigl[\Bigl\{\frac{4}{3} \Bigr\},\Bigl\{\frac{13}{6},\frac{7}{3}, \frac{7}{3}, \frac{5}{2} \Bigr\},Z\Bigr] \frac{2^3\cdot 5^2\cdot 11\cdot 13 i\sigma^3 \Gamma\left(\frac{4}{3}\right)}{13!\Gamma\left(\frac{7}{3}\right)} C_4 Z^9+

\\

+{^{4}_{1}F} \Bigl[\Bigl\{\frac{4}{3} \Bigr\},\Bigl\{\frac{7}{3},\frac{5}{2}, \frac{8}{3}, \frac{17}{6} \Bigr\},Z\Bigr] \frac{2^3\cdot 3\cdot 5\cdot 7\sigma^3 \Gamma\left(\frac{4}{3}\right)}{11! \Gamma\left(\frac{7}{3}\right)} C_1 Z^9+

\\

+{^{4}_{1}F} \Bigl[\Bigl\{\frac{3}{2} \Bigr\},\Bigl\{\frac{5}{2},\frac{8}{3}, \frac{17}{6}, \frac{19}{6} \Bigr\},Z\Bigr] \frac{2^7\cdot 3^5\sigma^3}{13!} \sqrt[3]{\frac{\sigma}{36}} C_2 Z^{10}+

\\

+{^{4}_{1}F} \Bigl[\Bigl\{\frac{5}{3} \Bigr\},\Bigl\{\frac{8}{3},\frac{17}{6}, \frac{19}{6}, \frac{10}{3} \Bigr\},Z\Bigr] \frac{2^3\cdot 3^4\cdot 5\sigma^3 \Gamma\left(\frac{5}{3}\right)} {13! \Gamma\left(\frac{8}{3}\right)} \sqrt[3]{\frac{\sigma^2}{6}} C_3 Z^{11}+

\\

+{^{4}_{1}F} \Bigl[\Bigl\{\frac{11}{6} \Bigr\},\Bigl\{\frac{17}{6},\frac{19}{6}, \frac{10}{3}, \frac{7}{2} \Bigr\},Z\Bigr] \frac{2^6\cdot 5\cdot 11\sigma^4 \Gamma\left(\frac{11}{6}\right)} {14! \Gamma\left(\frac{17}{6}\right)} C_4 Z^{12}-

\\

-{^{4}_{1}F} \Bigl[\Bigl\{\frac{7}{3} \Bigr\},\Bigl\{\frac{10}{3},\frac{7}{2}, \frac{11}{3}, \frac{23}{6} \Bigr\},Z\Bigr] \frac{2^6\cdot 3\cdot 7^2\cdot 13\sigma^5 \Gamma\left(\frac{7}{3}\right)} {17! \Gamma\left(\frac{10}{3}\right)} C_1 Z^{15}-

\\

-{^{4}_{1}F} \Bigl[\Bigl\{\frac{5}{2} \Bigr\},\Bigl\{\frac{7}{2},\frac{11}{3}, \frac{23}{6}, \frac{25}{6} \Bigr\},Z\Bigr] \frac{2^8\cdot 3^7\cdot 7\sigma^5 } {19!} \sqrt[3]{\frac{\sigma}{36}} C_2 Z^{16}-

\\

-{^{4}_{1}F} \Bigl[\Bigl\{\frac{8}{3} \Bigr\},\Bigl\{\frac{11}{3},\frac{23}{6}, \frac{25}{6}, \frac{13}{3} \Bigr\},Z\Bigr] \frac{2^7\cdot 3^6\cdot 5\sigma^5 \Gamma\left(\frac{8}{3}\right)C_3 Z^{17}} {19! \Gamma\left(\frac{11}{3}\right)} \sqrt[3]{\frac{\sigma^2}{6}} -

\\

-{^{4}_{1}F} \Bigl[\Bigl\{\frac{8}{3} \Bigr\},\Bigl\{\frac{23}{6},\frac{25}{6}, \frac{13}{3}, \frac{9}{2} \Bigr\},Z\Bigr] \frac{2^8\cdot 3^2\cdot 5^2\cdot 11\cdot 17\sigma^6 \Gamma\left(\frac{17}{6}\right)} {21! \Gamma\left(\frac{23}{6}\right)} \sqrt[3]{\frac{\sigma^2}{6}} C_4 Z^{18}.

\end{multline}

Here $\Gamma\left(\lambda\right)= \displaystyle\int_{0}^{+\infty}r^{\lambda-1}\exp\left(-r\right)dr$ is the Euler gamma function~\cite{ulg:bib:Polyanin}. Solutions~\eqref{eq:Aristov_Prosviryakov:soltemp}, \eqref{eq:Aristov_Prosviryakov:solstress},~\eqref{eq:Aristov_Prosviryakov:solvelo} are written as a linear  combination of generalized hypergeometric functions in the form of infinite power series. Consequently, examination of solution convergence is an urgent problem. Using the properties of convergence of these functions, we can state that the convergence radius for each solution equals infinity~\cite{ulg:bib:Polyanin}. It means that the solution of our problem always converges.



Let us now formulate the boundary conditions to determine the integration constants of equation system \eqref{eq:Aristov_Prosviryakov:sys1}. Without any claim to solution generality and without studying the shapes of the fluid, we give a simple example of such solutions. Let the simplest and most physically meaningful solution be a flow between two infinite disks rotating with free velocities. This is possible, since the solution can always be shifted along the $Oz$-axis, an appropriate distance between the disks being chosen. The disks must be nonuniformly heated to a proper temperature (the temperature must be related to the radius by a well-defined law). Assume that the lower boundary is isothermal (the temperature of the lower disk is zero); the adhesion condition holds on both boundaries. Thus, the following equalities are valid:

$$\text{when} \ \ z=0, \ \ U=V=0, \ \ T_1=T_2=0;$$

$$\text{when} \ \ z=h, \ \ U=A\cos\alpha, \ \ V=A\sin\alpha, \ \ T_1=B, \ \ T_2=C, \ \ P_1=\tau_1, \ \ P_2=\tau_2.$$



The conditions for determining the integration constants, written on the boundary, do not determine the integration constants uniquely. The system of linear algebraic equations for determining the integration constants is underdetermined, since there are few boundary conditions. To close the system of equations, we specify the fluid flow rate conditions as

$$\int _{0}^{h}Udz =Q_1,\quad  \int_{0}^{h}Vdz =Q_2.$$



When the flow rate is specified, it is generally assumed that  $Q_1=Q_2=0$~\cite{ulg:bib:Aristov2}. It is a physically justified convention~[\citen{ulg:bib:Aristov1},~\citen{ulg:bib:Aristov2}]. However, it is possible to obtain a solution by specifying a nonzero flow rate. Next, not limiting the generality of the reasoning, we will study the solutions simulating fluid motions with the zero flow rate. 



For the above-introduced complex velocities, pressure and temperature ~\eqref{eq:Aristov_Prosviryakov:soltemp}, \eqref{eq:Aristov_Prosviryakov:solstress},~\eqref{eq:Aristov_Prosviryakov:solvelo}, in view of transformations \eqref{eq:Aristov_Prosviryakov:trans}, the following boundary conditions are valid:

\begin{gather}

\label{eq:Aristov_Prosviryakov:bond}

\text{when} \ \ z=0: \ \ W=0, \ \ \Theta_1=0;

\\

\text{when} \ \ z=h: \ \ W=A\exp\left(i\alpha\right), \ \ \Theta_1=B+2A\Omega\sin\alpha+i\left(C-2A\Omega\cos\alpha\right),

\\

\int_{0}^{h}Wdz =0.

\end{gather}



With the boundary conditions \eqref{eq:Aristov_Prosviryakov:bond}, it follows from equality \eqref{eq:Aristov_Prosviryakov:soltemp} that  $C_1=0$, and from solution \eqref{eq:Aristov_Prosviryakov:solvelo} that  $C_4=0$. Thus, the expressions for the determination of hydrodynamic fields become much simpler. Remember that, in the computation of the integration constants, the variable $z=h\sigma Z$  is to be replaced.











\Section{Analysis of the solutions}

 Using conditions \eqref{eq:Aristov_Prosviryakov:bond}, it is easy, but rather cumbersome, to compute the complex  $C_2,$ $C_3,$ and $C_5$  which are related to the dispersion relation   $\sigma$ (kinematic viscosity) by the fractional linear law. Analyzing conditions \eqref{eq:Aristov_Prosviryakov:bond}, we find that the boundary value problem under study is meaningful at any positive value of  $\sigma.$ Nevertheless, the structure and behavior of the solutions depend considerably on the value of  $\sigma.$



It is well known that, to analyze the polynomial solutions of system \eqref{eq:Aristov_Prosviryakov:sys1}, it would suffice to localize the roots of the linear combination of the generalized hypergeometric functions~\cite{ulg:bib:Polyanin}. Any hypergeometric function is known to be monotonically increasing with a zero value at the origin. Therefore, the appearance of other zero values is possible only when there is a combination of functions with at least one negative weight coefficient~\cite{ulg:bib:Polyanin}. The appearance of negative coefficients depends on $\sigma$  and the function domain, which also depends on the dispersion relation $Z\in\left[0, \frac{1}{\sigma}\right].$  Temperature is characterized by three types of solutions. 

When  $\sigma\geq\frac{3}{2},$ the function describing temperature is a strictly monotonically increasing or decreasing function (fig.~\ref{fig_prosv1}). 

If the evaluation  $\frac{9}{10}<\sigma<\frac{3}{2}$  is valid for the dispersion relation, there is exactly one zero value of temperature inside the fluid layer (fig.~\ref{fig_prosv2}). With the further decrease, temperature gets a finite number of zeros (fig.~\ref{fig_prosv3}). Consequently, the fluid layer is divided into the regions with positive and negative temperature, which are characterized by alternating signs. The analysis of pressure is similar to that of temperature. 





Velocities are characterized by only two qualitative behaviors of the functions (figs.~\ref{fig_prosv2} and ~\ref{fig_prosv3}). Thus, at any  $\sigma,$ there appear fluid counterflows. It was shown in~[\citen{ulg:bib:Aristov3}--\citen{ulg:bib:Gorshkov}] that the appearance of counterflows is a purely nonlinear effect, i.\,e., the viscous forces are comparable to the convective acceleration, and even without Coriolis forces. However, when  $\sigma$ tends to zero, counterflow zones are formed, whose thickness decreases with  $\sigma.$ Other values of the dispersion relation are characteristic of velocity, and they are easy to compute. However, the values obtained for the temperature and pressure fields are a good reference point for refining the characteristic dispersion numbers of velocities.





When analyzing the solutions (fig.~\ref{fig_prosv3}), we can see their ``fluctuation''. In~\cite{ulg:bib:Aristov7} localized convective flows in terms were studied in the axisymmetric formulation. Solution ``damping'' was observed there. We use a wider class of solutions. There may be not only ``damping'' (when $\sigma$  is close to one), but also ``resonance'' initiation. Yet, the solution does not become singular, and the Boussinesq approxi\-mation is not violated in the simulation of the convective flow. This difference is attributed to the account of the quadratic summands for temperature. If tempera\-ture is considered to be a linear function with respect to the horizontal coordinates, the solutions obtained in this study coincide with those discussed in~\cite{ulg:bib:Aristov4}.



\begin{figure}[h!]



\begin{tabular}{p{.45\textwidth}p{.45\textwidth}}

\centerline{\includegraphics[width=0.45\textwidth]{fig_prosv1}} &

\centerline{\includegraphics[width=0.45\textwidth]{fig_prosv2}} \\

\caption{A qualitative form of the solutions  $T_1,$ $T_2,$ $P_1,$ $P_2$ of system (\ref{eq:Aristov_Prosviryakov:sys1}) when $\sigma\geq\frac{3}{2}$ ($1-\tau_1>0$,  $\tau_2>0$, $A>0$, $B>0$, $C>0$; $2-\tau_1<0$, $\tau_2<0$, $A<0$, $B<0$, $C<0$) \label{fig_prosv1}} &

\caption{A qualitative form of the solutions  $T_1,$ $T_2,$ $P_1,$ $P_2,$ $U,$ $V$ of system (\ref{eq:Aristov_Prosviryakov:sys1}) when $\frac{9}{10}<\sigma<\frac{3}{2}$ ($1-\tau_1>0$,  $\tau_2>0$, $A>0$, $B>0$, $C>0$; $2-\tau_1<0$, $\tau_2<0$, $A<0$, $B<0$, $C<0$) \label{fig_prosv2} }

\\

\centerline{\includegraphics[width=0.45\textwidth]{fig_prosv3}} &



\vspace{-2cm}

\caption{A qualitative form of the solutions  $T_1,$ $T_2,$ $P_1,$ $P_2,$ $U,$ $V$ of system (\ref{eq:Aristov_Prosviryakov:sys1}) when $\sigma\leq\frac{9}{10}$ ($1-\tau_1>0$,  $\tau_2>0$, $B>0$, $C>0$; $2-\tau_1<0$, $\tau_2<0$, $B<0$, $C<0$) \label{fig_prosv3} } \\





\end{tabular}



\vspace{-3cm}



%%\noindent

%%\begin{minipage}[h]{0.48\textwidth}

%%

%%\smallskip \footnotesize

%%Fig.~\ref{fig1}.

%%A qualitative form of the solutions  $T_1,$ $T_2,$ $P_1,$ $P_2$ of system (\ref{eq:Aristov_Prosviryakov:sys1}) when $\sigma\geq\frac{3}{2}$ ($1-\tau_1>0$,  $\tau_2>0$, $A>0$, $B>0$, $C>0$; $2-\tau_1<0$, $\tau_2<0$, $A<0$, $B<0$, $C<0$) \label{fig1} 

%%\end{minipage} \hfill

%%\begin{minipage}[h]{0.5\textwidth}

%%\centering

%%\includegraphics[width=0.98\textwidth]{fig1}

%%\end{minipage}

%\end{figure}

%

%

%\begin{figure}[h!]

%\noindent

%\begin{minipage}[h]{0.48\textwidth}

%\caption{A qualitative form of the solutions  $T_1,$ $T_2,$ $P_1,$ $P_2,$ $U,$ $V$ of system (\ref{eq:Aristov_Prosviryakov:sys1}) when $\frac{9}{10}<\sigma<frac{3}{2}$ ($1-\tau_1>0$,  $\tau_2>0$, $A>0$, $B>0$, $C>0$; $2-\tau_1<0$, $\tau_2<0$, $A<0$, $B<0$, $C<0$) 

%\label{fig2}

% }

%\smallskip \footnotesize

%Fig.~\ref{fig2}.

%A qualitative form of the solutions  $T_1,$ $T_2,$ $P_1,$ $P_2,$ $U,$ $V$ of system (\ref{eq:Aristov_Prosviryakov:sys1}) when $\frac{9}{10}<\sigma<\frac{3}{2}$ ($1-\tau_1>0$,  $\tau_2>0$, $A>0$, $B>0$, $C>0$; \centerline{$2-\tau_1<0$, $\tau_2<0$, $A<0$, $B<0$, $C<0$)} \label{fig2} 

%\end{minipage} \hfill

%\begin{minipage}[h]{0.5\textwidth}

%\centering

%\includegraphics[width=0.98\textwidth]{fig2}

%\end{minipage}

%\end{figure}

%

%

%

%\begin{figure}[h!]

%\noindent

%\begin{minipage}[h]{0.48\textwidth}

%\caption{A qualitative form of the solutions  $T_1,$ $T_2,$ $P_1,$ $P_2,$ $U,$ $V$ of system (\ref{eq:Aristov_Prosviryakov:sys1}) when $\sigma\leq\frac{9}{10}$ ($1-\tau_1>0$,  $\tau_2>0$, $B>0$, $C>0$; $2-\tau_1<0$, $\tau_2<0$, $B<0$, $C<0$) 

%\label{fig3}

% }

%\smallskip \footnotesize

%Fig.~\ref{fig3}.

%A qualitative form of the solutions  $T_1,$ $T_2,$ $P_1,$ $P_2,$ $U,$ $V$ of system (\ref{eq:Aristov_Prosviryakov:sys1}) when $\sigma\leq\frac{9}{10}$ ($1-\tau_1>0$,  $\tau_2>0$, $B>0$, $C>0$; 

%\centerline{$2-\tau_1<0$, $\tau_2<0$, $B<0$, $C<0$)}

% \label{fig3} 

%\end{minipage} \hfill

%\begin{minipage}[h]{0.5\textwidth}

%\centering

%\includegraphics[width=0.98\textwidth]{fig3}

%\end{minipage}

\end{figure}









\newpage















\Section{Solution of equation  with the Prandtl number different from one}

An important particular case of equation \eqref{eq:Aristov_Prosviryakov:trans4} was studied above; namely, when the values of kinematic viscosity and thermal diffusivity of the fluid coincide. 

Let us now write the solution of equation \eqref{eq:Aristov_Prosviryakov:trans4} in the general case:

\begin{multline}

\label{eq:Aristov_Prosviryakov:soltempPr}

\Theta_1=\text{\sf Ai}\Bigl(\xi+\frac{4\Omega z}{\sqrt{\nu\chi}}\sqrt[3]{\frac{\nu^2\chi^2} {16\Omega^2}}\Bigr)

\left(C_1\text{\sf Ai}\left(\xi\right)+C_2\text{\sf Bi}\left(\xi\right)\right)+

\\

+\text{\sf Bi}\Bigl(\xi+\frac{4\Omega z}{\sqrt{\nu\chi}}\sqrt[3]{\frac{\nu^2\chi^2}{16\Omega^2}}\Bigr)\left(C_3\text{\sf Ai}\left(\xi\right)+C_4\text{\sf Bi}\left(\xi\right)\right).

\end{multline}

Here $$\xi=\sqrt[3]{\nu^2\chi^2}\displaystyle\frac{ i\left(\nu-\chi\right)+2\sqrt{\nu\chi}\Omega z}{\nu\chi\sqrt[3]{16\Omega^2}}$$ is an auxiliary variable introduced in order to make the solution notation concise;  $\text{\sf Ai} \left(\xi\right)$   and $\text{\sf Bi}\left(\xi\right)$  are the Airy function and the associated function~\cite{ulg:bib:Polyanin}. Solution \eqref{eq:Aristov_Prosviryakov:soltempPr} is complex-valued, since it is written for arbitrary differences between $\nu$  and $\chi.$  This is done for the convenience of the solution analysis. Solution \eqref{eq:Aristov_Prosviryakov:soltempPr} being known, pressure  $S_1$ and velocity $W$ are computed similarly to the previous case.	



When the dissipative coefficients  $\nu$ and  $\chi$ of the nonisothermal fluid are equal, solution \eqref{eq:Aristov_Prosviryakov:soltempPr}  is reduced to formula \eqref{eq:Aristov_Prosviryakov:soltemp}. Studying the convergence of series \eqref{eq:Aristov_Prosviryakov:soltempPr}, we find that the solution expressed in terms of this series converges on the solution domain, i.\,e. layer thickness. 



The consideration of equation \eqref{eq:Aristov_Prosviryakov:trans4} with an arbitrary relation between $\nu$  and~$\chi$ offers solutions to system \eqref{eq:Aristov_Prosviryakov:sys1}, \eqref{eq:Aristov_Prosviryakov:sys2} in qualitative terms, see figs.~\ref{fig_prosv1} to~\ref{fig_prosv3}. In the simulation of the hydrodynamic fields by the linear combination of the Airy functions the behavior of the solutions is predictable from the very beginning, since these special functions satisfy the simplest differential equation having a point where the fluctuation of the solution is replaced by its exponential growth. It was shown above that the solution does not go off to infinity.



\smallskip



\Section{Conclusion}

The problem of simulating rotating fluid masses has a long history. To describe fluid motion (relative equilibrium), ordinary differential equa\-tions are used that enable one to obtain solutions simulating fluid flows by the theory of motion stability with the disturbance of the background (principal) flow. This paper discusses the exact solution to an overdetermined boundary value problem, which describes stationary flow at any scales where the Boussinesq approximation is valid. In this case, stationary dynamic fluid equilibria are chara\-cterized by the formation of counterflows, the sign of velocity being able to alternate several times, depending on the boundary conditions and the geometric dimensions of the fluid layer. A similar situation is true for temperature and pressure, namely, the formation of zones with the negative and positive values of temperature and pressure with respect to the reference value. 



\smallskip 

 



\AdditionalContent{Competing interests}{I declare that I have no competing interests.} 

\AdditionalContent{Author's Responsibilities}{I take full responsibility for submitting the final manuscript in print. I approved the final version of the manuscript.} 





\AdditionalContent{Funding}{The work was done within the state assignment from FASO Russia, theme No. AAAA-A18-118020790140-5.}









\begin{thebibliography}{28}



\Bibitem{ulg:bib:Newtoni}

\by Newton~I.

\book Opera quae exstant omnia

\bookinfo Faksimile-Neudruck der Ausgabe von Samuel Horsley, London 1779--1785 in f\"unf B\"anden. Band 1,~4

\lang In German

\zmath{0129.25102}

\publaddr Stuttgart-Bad Cannstatt

\publ Friedrich Frommann Verlag (G\"unther Holzboog)

\totalpages xx+592

\yr 1964



\Bibitem{ulg:bib:Turnbull}

\book The Correspondence of Isaac Newton

\bookvol 1. 1661-1675

\ed H.~W.~Turnbull

\publ Cambridge Univ. Press

\publaddr New York

\yr 1959

\totalpages xxxviii+468 





\Bibitem{ulg:bib:Poincare}

\by Poincar\'e~H.

\book Cin\'ematique et m\'ecanismes. Potentiel et m\'ecanique des fluides

\bookinfo Cours profess\'e \`a la Sorbonne

\lang In French

\zmath{30.0608.01}

\publaddr Paris

\publ G.~Carr\'e et C.~Naud

\totalpages iv+385 

\yr 1899



\Bibitem{ulg:bib:Lyapunov1}

\by Lyapunov~A.~M.

\paper On the stability of ellipsoidal equilibrium forms of rotating fluid

\inbook Collected works. \rm  Vol. III

\yr 1959

\publ Akad. Nauk SSSR

\publaddr Moscow

\pages 5--113

\zmath{0192.31203}

\lang In Russian



\Bibitem{ulg:bib:Lyapunov2}

\by Lyapunov~A.~M.

\paper On the equilibrium figures of rotating homogeneous liquid mass slightly different from ellipsoids

\inbook Collected works. \rm  Vol. IV

\yr 1959

\publ Akad. Nauk SSSR

\publaddr Moscow

\pages 5--645

\zmath{0192.31204}

\lang In Russian



\Bibitem{ulg:bib:Chandrasekhar}

\by Chandrasekhar~S.

\book Ellipsoidal figures of equilibrium

\zmath{0213.52304}

\publaddr New Haven, London

\publ Yale University Press 

\totalpages ix+252 

\yr 1969



\Bibitem{ulg:bib:Hagihara}

\by Hagihara~Y.

\book Theories of equilibrium figures of a rotating homogeneous fluid mass

\yr 1970

\serial NASA Special Publication 

\vol 186 

\publaddr Washington

\publ US Government Printing Office

\adsnasa{1970NASSP.186.....H}



\Bibitem{ulg:bib:Borisov}

\by Borisov~A.~V., Kilin~A.~A., Mamaev~I.~S.

\paper The Hamiltonian Dynamics of Self-gravitating Liquid and Gas Ellipsoid

\jour Regul. Chaotic Dyn. 

\yr 2009

\vol 14

\issue 2

\pages 179--217

\zmath{https://zbmath.org/?q=an:1229.70049}

\crossref{10.1134/S1560354709020014}



\Bibitem{ulg:bib:Appell}

\by Appell~P.

\book Trait\'e de m\'ecanique rationnelle

\bookinfo Tome IV, 1: Figures d'\'equilibre d'une masse liquide homog\`ene en rotation

\lang In French

\zmath{58.1300.02}

\totalpages viii+342 

\publaddr  Paris

\publ Gauthier-Villars 

\yr 1932







\Bibitem{ulg:bib:Betti}

\by Betti~E.

\paper Sopra i moti che conservano la figura ellissoidale a una massa fluida eterogenea

\lang In Italian

\zmath{13.0718.01}

\jour Brioschi Ann

\issueinfo  (2)~X

\pages 173--187

\yr 1879







\Bibitem{ulg:bib:Volterra}

\by Volterra~V.

\paper Sur la stratification d\'une masse fluide en \'equilibre

\jour Acta Math.

\yr 1903

\vol 27

\issue 1

\pages 105--124



\Bibitem{ulg:bib:Liouville}

\by Liouville~J.

\paper Sur la figure d\'une masse fluide homog\'ene, en \'equilibre et dou\'ee d\'un mouvement de rotation

\jour J. de l'\'Ecole Polytech.

\yr 1834

\vol 14

\pages 289–296



\Bibitem{ulg:bib:Dirichlet}

\by Dirichlet~G.~L.

\paper Untersuchungen \"uber ein Problem der Hydrodynamik (Aus dessen Nachlass hergestellt von Herrn R. Dedekind zu Z\"urich)

\jour J. Reine Angew. Math. (Crelle's Journal)

\yr 1861

\vol 58

\pages181--216



\Bibitem{ulg:bib:Arnold}

\by Arnold~V. I., Khesin~B.~A.

\book Topological Methods in Hydrodynamics

\yr 1999

\publ Springer

\publaddr New York

\totalpages 392



\Bibitem{ulg:bib:Dolzhansky}

\by Dolzhansky~F.~V.

\paper On the mechanical prototypes of fundamental hydrodynamic invariants and slow manifolds

\jour Physics–Uspekhi

\yr 2005

\vol 48

\issue 12

\pages 1205--1234

\crossref{10.1070/PU2005v048n12ABEH002375}



\Bibitem{ulg:bib:Hadamard}

\by Hadamard~J.

\paper Mouvement permanent lent d'une sph\`ere liquide et visqueuse dans un liquide visqueux

\zmath{42.0813.01}

\jour C.~R.~Acad. Sci., Paris 

\vol 152

\issue 25

\pages 1735--1738 

\yr 1911





\Bibitem{ulg:bib:Rybczynski}

\by Rybczi\'nski, W.

\paper \"Uber die fortschreitende Bewegung einer fl\"ussigen Kugel in einem z\"ahen Medium

\zmath{42.0811.02}

\jour Bull. Int. Acad. Sci. Cracovie. Ser. A

\vol 1

\pages 40--46 

\yr 1911



\Bibitem{ulg:bib:Aristov1}

\by Aristov~S.~N., Prosviryakov~E.~Yu. 

\paper A New Class of Exact Solutions for Three Dimensional Thermal Diffusion Equations

\jour Theor. Found. Chem. Eng.

\yr 2016

\vol 50

\issue 3

\pages 286--293

\crossref{10.1134/S0040579516030027}



\Bibitem{ulg:bib:Aristov2}

\by Aristov~S.~N., Knyazev~D.~V., Polyanin~A.~D.

\paper Exact solutions of the Navier-Stokes equations with the linear dependence of velocity components on two space variables

\jour Theor. Found. Chem. Eng.

\crossref{10.1134/S0040579509050066}

\yr 2009

\vol 43

\issue 5

\pages 642--662



\Bibitem{ulg:bib:Polyanin}

\by Polyanin~A.~D., Zaitsev~V.~F.

\book Handbook of ordinary differential equations. Exact solutions, methods, and problems

\zmath{06776009}

\publaddr Boca Raton, FL

\publ CRC Press 

\totalpages xxix+1456 

\yr 2018



\Bibitem{ulg:bib:Aristov3}

\by Aristov~S.~N., Prosviryakov~E.~Yu.

\paper Unsteady Layered Vortical Fluid Flows

\jour Fluid Dyn.

\crossref{10.1134/S0015462816020034}

\yr 2016

\vol 51

\issue 2

\pages 148--154



\Bibitem{ulg:bib:Aristov4}

\by Aristov~S.~N., Prosviryakov~E.~Yu., Spevak~L.~F.

\paper Unsteady-State B\'enard--Marangoni Convection  in Layered Viscous Incompressible Flows

\jour Theor. Found. Chem. Eng.

\crossref{10.1134/S0040579516020019}

\yr 2016

\vol 50

\issue 2

\pages 132--141



\Bibitem{ulg:bib:Aristov5}

\by Aristov~S.~N., Prosviryakov~E.~Yu.

\paper On one class of analytic solutions of the stationary axisymmetric convection B\'enard--Marangoni viscous incompressible fluid

\jour Vestn. Samar. Gos. Tekhn. Univ., Ser. Fiz.-Mat. Nauki \rm [J. Samara State Tech. Univ., Ser. Phys. Math. Sci.]

\yr 2013

\vol 3(32)

\pages 110--118

\mathnet{http://mi.mathnet.ru/vsgtu1205}

\crossref{10.14498/vsgtu1205}

\zmath{https://zbmath.org/?q=an:06968790}

\elib{http://elibrary.ru/item.asp?id=20710187}

\lang In Russian



\Bibitem{ulg:bib:Aristov6}

\by Aristov~S.~N., Privalova~V.~V., Prosviryakov~E.~Yu.

\paper Stationary nonisothermal Couette flow. Quadratic heating of the upper boundary of the fluid layer

\jour Nelin. Dinam.

\yr 2016

\vol 12

\issue 2

\pages 167--178

\mathnet{http://mi.mathnet.ru/nd519}

\elib{http://elibrary.ru/item.asp?id=26193555}

\crossref{10.20537/nd1602001}

\lang In Russian



\Bibitem{ulg:bib:Burmasheva1}

\by Burmasheva~N.~V., Prosviryakov~E.~Yu.

\paper A large-scale layered stationary convection of an incompressible viscous fluid under the action of shear stresses at the upper boundary. Velocity field investigation

\jour Vestn. Samar. Gos. Tekhn. Univ., Ser. Fiz.-Mat. Nauki \rm [J. Samara State Tech. Univ., Ser. Phys. Math. Sci.]

\yr 2017

\vol 21

\issue 1

\pages 180--196

\mathnet{http://mi.mathnet.ru/vsgtu1527}

\crossref{10.14498/vsgtu1527}

\elib{http://elibrary.ru/item.asp?id=29245104}

\lang In Russian



\Bibitem{ulg:bib:Burmasheva2}

\by Burmasheva~N.~V., Prosviryakov~E.~Yu.

\paper A large-scale layered stationary convection of a incompressible viscous fluid under the action of shear stresses at the upper boundary. Temperature and presure field investigation

\jour Vestn. Samar. Gos. Tekhn. Univ., Ser. Fiz.-Mat. Nauki \rm [J. Samara State Tech. Univ., Ser. Phys. Math. Sci.]

\yr 2017

\vol 21

\issue 4

\pages 736--751

\mathnet{http://mi.mathnet.ru/vsgtu1568}

\crossref{10.14498/vsgtu1568}

\zmath{https://zbmath.org/?q=an:06964885}

\elib{http://elibrary.ru/item.asp?id=32712835}

\lang In Russian



\Bibitem{ulg:bib:Gorshkov}

\by Gorshkov~A.~V., Prosviryakov~E.~Yu.

\paper Ekman Convective Layer Flow of a Viscous Incompressible Fluid

\jour Izv. Atmos. Ocean. Phys.

\yr 2018

\vol 54

\issue 2

\pages 89–195

\crossref{10.1134/S0001433818020081}





\Bibitem{ulg:bib:Aristov7}

\by Aristov~S.~N., Knyazev~D.~V.

\paper Localized convective flows in a nonuniformly heated liquid layer

\jour Fluid Dyn.

\yr 2014

\vol 49

\issue 5

\pages 565--575

\crossref{10.1134/S0015462814050020}











\end{thebibliography}





\renewcommand{\baselinestretch}{0.95}



\newpage



\selectlanguage{russian}



$\;$







\makerutitle



\vspace{-7mm}



\AdditionalContent{Конкурирующие интересы}{Конкурирующих интересов не имею.}



\AdditionalContent{Авторский вклад и ответственность}{Я несу полную ответственность за предоставление окончательной версии рукописи в печать. Окончательная версия 

рукописи мною одобрена.}

\AdditionalContent{Финансирование}{Работа выполнена в рамках государственного задания ФАНО, тема № АААА-А18-118020790140-5.}





\renewcommand{\baselinestretch}{0.9}







\renewcommand{\RuCitation}{\noindent\textbf{\CYRO\cyrb\cyrr\cyra\cyrz\cyre\cyrc\ \cyrd\cyrl\cyrya\ \cyrc\cyri\cyrt\cyri\cyrr\cyro\cyrv\cyra\cyrn\cyri\cyrya\\} {\RuAuthorInContents} {\RuTitleInContents}~/\!/ \textit{\RuJournalName}, \JournalYear.  \CYRT.~\JournalVolume, \textnumero~\JournalNumber.  \CYRS.~\thefirstpage--\thelastpage. \doiurl.}



\renewcommand{\EnCitation}{\noindent\textbf{Please cite this article in press as:\\} {\EnAuthorInContents} {\EnTitleInContents}, \textit{\EnJournalName}\/ [\ParallelJournalName],  \JournalYear, vol.~\JournalVolume, no.~\JournalNumber,  pp.~\thefirstpage--\thelastpage.  \midoiurl\ (In Russian).}









%\end{document}

 \newpage

%\input{preamble}

%

%\usepackage{samgtu-name}

%

%\usepackage{slashbox}



\vsgtu{1653} %registration number

\FirstPage{750}

\LastPage{761}

%

%\pdfcrop

%\draft



\ShortCommunication



%\begin{document}





\hfill \qrcode[hyperlink,height=.8in]{http://doi.org/10.14498/vsgtu1653}

\vspace{-.8in}





\udc{517.958:539.3(1)}

\msc{74F05, 74K20}



\rutitle{Динамическая устойчивость геометрически нерегулярной нагретой пологой цилиндрической оболочки в~сверхзвуковом потоке газа}

\rutitleshort{Динамическая устойчивость геометрически нерегулярной\dots}



\entitle{Dynamic stability of heated geometrically irregular cylindrical shell in supersonic gas flow}

\entitleshort{Dynamic stability of heated geometrically irregular cylindrical shell in supersonic gas flow}





\ruauthor{Г.~Н.~Белосточный, О.~А.~Мыльцина}{Б\,е\,л\,о\,с\,т\,о\,ч\,н\,ы\,й~Г.~Н.,\ М\,ы\,л\,ь\,ц\,и\,н\,а~О.~А.}

\enauthor{G.~N.~Belostochnyi, O.~A.~Myltcina}{B\,e\,l\,o\,s\,t\,o\,c\,h\,n\,y\,i~G.~N.,\ M\,y\,l\,t\,c\,i\,n\,a~O.~A.}



\ruaffil{\saratovuni.}



\enaffil{\engsaratovuni.}



\ruauthorsdetails{3}{% Количество авторов



\fullauthor{\rupid{98285}{Григорий Николаевич Белосточный}}

%\CAuthor % Автор-корреспондент

\orcid{0000-0003-4471-6599} % ORCID ID автора 

\regalia{доктор технических наук, профессор} 

\position{профессор} 

\dept{каф. математической теории упругости и биомеханики} 

\email{belostochny@mail.ru}

\fullauthor{\rupid{98283}{Ольга Анатольевна Мыльцина}}

\CAuthor % Автор-корреспондент

\orcid{0000-0003-4718-2772} % ORCID ID автора 

\regalia{кандидат физико-математических наук} 

\position{ассистент} 

\dept{каф. теории функций и стохастического анализа} 

\email{omyltcina@yandex.ru}

}



%Olga Myltcina https://orcid.org/0000-0003-4718-2772 

%Grigory Belostochnyihttps://orcid.org/0000-0001-6949-4174

%Acel Polienko https://orcid.org/0000-0003-4471-6599







\enauthorsdetails{3}{% Количество авторов

\fullauthor{\rupid{98285}{Grigory N. Belostochnyi}}

%\CAuthor % Автор-корреспондент

\orcid{0000-0003-4471-6599} % ORCID ID автора 

\regalia{Dr. Techn. Sci.} 

\position{Professor} 

\dept{Dept. of Mathematic Theory of Elasticity \&  Biomechanics} 

\email[n]{belostochny@mail.ru} 

\fullauthor{\rupid{98283}{Olga A. Myltcina}}

\CAuthor % Автор-корреспондент

\orcid{0000-0003-4718-2772} % ORCID ID автора 

\regalia{Cand. Phys. \& Math. Sci.} 

\position{Assistant} 

\dept{Dept. of Functions \& Approxmation Theory} 

\email[n]{omyltcina@yandex.ru}

}





\rukeywords{динамическая устойчивость, температура, пологие оболочки, сверхзвук, континуальность, обобщенные функции, поршневая теория, аэродинамика, критические скорости, ребра жесткости, демпфирование, кривизна, критерий Гурвица, изотропия}

\enkeywords{dynamic stability, temperature, flat shells, supersonic, continuity, generalized functions, piston theory, aerodynamics, critical velocities, ribs, damping, curvature, Routh–Hurwitz stability criterion, isotropy}



\ruabstract{На базе модели типа Лява рассматривается нагретая до постоянной температуры геометрически нерегулярная пологая цилиндрическая оболочка, обдуваемая сверхзвуковым потоком газа со стороны одной из ее основных поверхностей. За основу взята континуальная модель термоупругой системы в~виде тонкостенной оболочки, подкрепленной ребрами вдоль набегающего газового потока. Сингулярная система уравнений динамической термоустойчивости геометрически нерегулярной оболочки содержит слагаемые, учитывающие «растяжение-сжатие» и~сдвиг подкрепляющих элементов в~тангенциальной плоскости, тангенциальные усилия, вызванные нагревом оболочки, и~поперечную нагрузку, стандартным образом записанную по «поршневой теории».



Тангенциальные усилия предварительно определяются как решение сингулярных дифференциальных уравнений безмоментной термоупругости геометрически нерегулярной оболочки с~учетом краевых усилий.



Решение системы динамических уравнений термоупругости оболочки разыскивается в~виде суммы двойного тригонометрического ряда (для функции прогиба) с~переменными по временной координате коэффициентами. На основании метода Галеркина получена однородная система для коэффициентов аппроксимирующего ряда, которая сведена к~одному дифференциальному уравнению четвертого порядка. Решение приводится во втором приближении, что соответствует двум полуволнам в~направлении потока и~одной полуволне в~перпендикулярном направлении. На основании стандартных методов анализа динамической устойчивости тонкостенных конструкций определяются критические значения скоростей газового потока.



Количественные результаты приводятся в виде таблиц, иллюстрирующих влияние геометрических параметров термоупругой системы «обо\-лоч\-ка-ребра», температуры  на устойчивость геометрически нерегулярной цилиндрической оболочки в~сверхзвуковом потоке газа с~учетом демпфирования.}



\enabstract{On the basis of the Love model, a geometrically irregular heated cylindrical shell blown by a supersonic gas flow from one of its main surfaces is considered. The continuum model of a thermoelastic system in the form of a thin-walled shell supported by ribs along the incoming gas flow is taken as a basis. The singular system of equations for the dynamic thermal stability of a geometrically irregular shell contains terms that take into account the tension-compression and the shift of the reinforcing elements in the tangential plane, the tangential forces caused by the heating of the shell and the transverse load, as standard recorded by the piston theory. The solution of a singular system of differential equations in displacements, in the second approximation for the deflection function, is sought in the form of a double trigonometric series with time coordinate variables. 





Tangential forces are predefined as the solution of singular differential equations of non-moment thermoelasticity of a geometrically irregular shell taking into account boundary forces. 



The solution of the system of dynamic equations of thermoelasticity of the shell is sought in the form of the sum of the double trigonometric series (for the deflection function) with time coordinate variable coefficients. On the basis of the Galerkin method, a homogeneous system for the coefficients of the approximating series is obtained, which is reduced to one fourth-order differential equation. The solution is given in the second approximation, which corresponds to two half-waves in the direction of flow and one half-wave in the perpendicular direction. On the basis of standard methods of analysis of dynamic stability of thin-walled structures are determined critical values of the gas flow rate. 



The quantitative results are presented in the form of tables illustrating the influence of the geometrical parameters of the thermoelastic shell-edge system, temperature and damping on the stability of a geometrically irregular cylindrical shell in a supersonic gas flow.}



\newdate{myla}{11}{10}{2018}

\date{\displaydate{myla}}

\newdate{mylb}{10}{11}{2018}

\revisiondate{\displaydate{mylb}}

\newdate{mylc}{12}{11}{2018}

\accepteddate{\displaydate{mylc}}



\newdate{myle}{21}{12}{2018}

\renewcommand{\JournalDataOnline}{\displaydate{myle}}



\makerutitle







Геометрически нерегулярные тонкостенные упругие системы, обширный класс которых составляют ребристые оболочки и~пластинки, широко используются в~различных областях современной техники. Условия эксплуатации таких систем предусматривают совместное воздействие температурных полей и~высокоскоростных газовых потоков. Исследованию упругого поведения гладких пластин и~оболочек на основе атермической

теории посвящено большое число работ, полный перечень которых содержал бы десятки

наименований. Ограничимся некоторыми из них \cite{myl-bib:1, myl-bib:2, myl-bib:3,myl-bib:4,myl-bib:5,myl-bib:5a}.

Значительно меньше работ

посвящено  исследованиям совместного воздействия температурных факторов и сверхзвукового потока

на тонкостенные системы. Важные для практики результаты в этой области содержатся в работах \cite{myl-bib:6, myl-bib:7, myl-bib:8}.



%Работы, в которых анализируется влияние подкрепляющих  оболочку ребер под действием сверхзвукового потока на базе термической теории в открытой научной литературе отсутствуют. 

Работы, в которых анализируется влияние подкрепляющих  ребер на поведение оболочки, находящейся под действием сверхзвукового потока, на базе термической теории, в~открытой научной литературе отсутствуют. 

Это связано не с маловажностью проблемы, а прежде всего с чрезвычайной математической сложностью таких задач, решаемых на основе дискретной модели «оболочка–ребра». Использование континуальной модели типа «конструктивная анизотропия» мало пригодно для практических целей, так как количественные результаты, полученные на ее основе, далеки от физической реальности.



Обращение к  континуальным моделям на основе теории обобщенных функций, основные положения которых содержится в работах \cite{myl-bib:9, myl-bib:10, myl-bib:11, myl-bib:12, myl-bib:15,myl-bib:16,myl-bib:18}, позволило сводить решения задач статической и динамической термоупругости ребристых  оболочек к интегрированию систем сингулярных дифференциальных уравнений точными и~приближенными методами высшего анализа \cite{myl-bib:19, myl-bib:21}, что немаловажно для инженерной практики.



\smallskip



{\bf 1.} Система сингулярных дифференциальных уравнений динамической термоупругости пологой геометрически нерегулярной цилиндрической оболочки, стандартным образом отнесенной к декартовым координатам на основе континуальной модели, в которых учет тангенциальных усилий записывается в форме Рейснера \cite{myl-bib:22a}, в компонентах поля перемещений примет вид

\begin{gather}

u_{\bigcom 11}+\frac{1-\nu}{2}u_{\bigcom 22}+\frac{1+\nu}{2} v_{\bigcom 12}-k_{11}w_{\bigcom 1}+

\varepsilon_1\frac{1-\nu}{2} \sum_{i=1}^n \frac{h_i}{h}(u_{\bigcom 2}+v_{\bigcom 1} )_{\bigcom 2}\delta(x-x_i)=0,

\\

\frac{1+\nu}{2}u_{\bigcom 21}+ v_{\bigcom 22}+\frac{1-\nu}{2}v_{\bigcom 11}-\nu k_{11}w_{\bigcom 2}+\varepsilon_2\sum_{i=1}^n\frac{h_i}{h}a_i ( \nu u_{\bigcom 1}+v_{\bigcom 2})_{\bigcom 2} \delta(x-x_i)=0,

\end{gather}



\vspace{-12mm}



\begin{multline}\label{bel_eq1}

\nabla^2\nabla^2 w+\frac{B}{D}k^2_{11}w-k_{11}\frac{B}{D} ( u_{\bigcom 1}+\nu v_{\bigcom 2})- 

\\

- \frac{1}{D} \bigl[ 

( T^{11}_0 w_{\bigcom 1})_{\bigcom 1}+

( T^{12}_0 w_{\bigcom 1})_{\bigcom 2}+  

(T^{12}_0 w_{\bigcom 2})_{\bigcom 1}+ 

(T^{22}_0 w_{\bigcom 2})_{\bigcom 2}\bigr]-

\\

-

\frac{1}{D}\sum_{i=1}^n\frac{h_i}{h}a_i

(T^{22}_0 w_{\bigcom 2})\delta(x-x_i)+ 

\\

+\sum_{i=1}^n \Bigl( \frac{h_i}{h}\Bigr)^3 \Phi_{3i}a_iw_{\bigcom 2222}\delta(x-x_i)+

\\

+

2(1-\nu)\sum_{i=1}^n\Bigl( \frac{h_i}{h}\Bigr)^3 \Phi_{3i}a_iw_{\bigcom 122}\bigr|_{x_i}\frac{d\delta(x-x_i)}{dx}= \\

=\frac{\gamma h}{g D}w_{\bigcom tt}\biggl[ 1+\sum_{i=1}^n \frac{h_i}{h}a_i\delta(x-x_i)\biggr] 

-\frac{\mu} {h D}w_{\bigcom t}-p_0\frac{\varkappa \mu}{D}w_{\bigcom 2}, 

\end{multline}

где $p_0\frac{\varkappa \mu}{D}w_{\bigcom 2}$ "---  относительная интенсивность поперечной нагрузки, вызванная прогибом пластинки, стандартным образом записана согласно поршневой теории \cite{myl-bib:1,myl-bib:22, myl-bib:30}; $M={v_y}/{c_0}$ "--- число Маха; 

$v_y$ "--- невозмущенная скорость набегающего газового потока; 

$c_0$ "--- скорость звука на бесконечности; $u$, $v$, $w$ "---  компоненты поля (векторы) перемещений; 

${h_i}/{h}$ "--- относительная высота $i$-того ребра, число которых $n$; 

$a_i$ "--- ширина $i$-того ребра; 

$k_{11}$ "--- кривизна цилиндрической оболочки; 

$D=\frac{E h^3}{12(1-\nu)}$; $\nu$ "--- коэффициент Пуассона; 

$E$ "--- модуль Юнга, 

$\gamma$ "--- удельный вес; 

$g$ "--- интенсивность поля тяжести; 

$$\Phi_{3i}=1+3\frac{h_i}{h}+\Bigl( \frac{h_i}{h}\Bigr)^3 ;$$ 

$\mu$ "--- коэффициент демпфирования; 

$\delta(x-x_i)$ "--- обобщенная дельта"=функция Дирака~\cite{myl-bib:25};  

$\varepsilon_j$ $(j=0,1)$ "--- знаковые числа, равные 1 или 0; 

$ T^{ij}_0 $ "--- тангенциальные усилия, вызванные нагревом пластинки до постоянной температуры $\Theta_0$, которые определяются на основании решений системы сингулярных однородных дифференциальных уравнений, описывающих безмоментное состояние термоупругой системы в компонентах поля перемещений $\bar{u}^0(u^0,v^0,w^0)$:

 \begin{multline}

 u^0_{\bigcom 11}+\frac{1-\nu}{2}u^0_{\bigcom 22}+\frac{1+\nu}{2}v^0_{\bigcom 12}-k_{11}w^0_{\bigcom 1}+

\\

 +\frac{1-\nu}{2}\sum_{i=1}^{n}\frac{h_i}{h}a_i ( u^0_{\bigcom 2}+v^0_{\bigcom 1})_{\bigcom 2}\delta(x-x_i) =0,

 \end{multline}

 \begin{multline}

 \frac{1+\nu}{2}u^0_{\bigcom 12}+v^0,_{22}+\frac{1-\nu}{2}v^0_{\bigcom 11}-\nu k_{11}w^0_{\bigcom 2}+

\\

+ 

 \sum_{i=1}^{n}\frac{h_i}{h}a_i( v^0_{\bigcom 2}+\nu u^0_{\bigcom 1})_{\bigcom 2}\delta(x-x_i)=0,

 \end{multline}

\begin{equation}\label{eq:bel:momentnot}

k_{11}^2 w^0+k_{11} (u^0_{\bigcom 1}+\nu v^0_{\bigcom 2} ) =\alpha(1+\nu)k_{11}\Theta_0.

\end{equation}





При неоднородных краевых условиях

\begin{equation*}

\begin{array}{ll}

\text{$x=0$, $x=a$:}&

u^0_{\bigcom 2}+v^0_{\bigcom 1}=0, \quad u^0_{\bigcom 1}+\nu v^0_{\bigcom 2}-k_{11}w^0=\alpha(1+\nu \Theta_0), \quad w^0=0;

\\[2mm]

\text{$y=0$, $y=b$:}&

v^0=0, \quad u^0_{\bigcom 2}+ v^0_{\bigcom 1}=0, \quad w^0=0

\end{array}

\end{equation*}

решения сингулярной системы \eqref{eq:bel:momentnot} запишутся в элементарных функциях:

\begin{equation}\label{bel_gr1}

v^0=0, \quad u^0=\alpha(1+\nu)\Theta_0 x, \quad w^0=0,

\end{equation}

а тангенциальные усилия примут вид

\begin{equation*}

T^{11}_0=T^{12}_0=0, \quad T^{22}_0=-E h\alpha\Theta_0.

\end{equation*}



Решение системы \eqref{bel_eq1} (в~которой отсутствуют инерционные слагаемые в~тангенциальной плоскости \cite{myl-bib:1,myl-bib:6, myl-bib:7}), тождественно удовлетворяющее условиям шарнирного закрепления всех сторон пологой цилиндрической оболочки

\begin{equation*}

\begin{array}{ll}

\text{при $x=0$, $x=a$:} &

T^{12}=0, \quad T^{11}=0, \quad w=0 \quad M^{11}=0;

\\[2mm]

\text{при $y=0$, $y=b$:} &

v=0, \quad T^{12}=0, \quad w=0 \quad M^{22}=0,

\end{array}

\end{equation*}

зададим в виде

\begin{gather}\label{bel_appr1}

u(x,y,t)=\tilde{u}(t)u_1 ({x}/{a} )u_2 ({y}/{b} ), \quad 

v(x,y,t)=\tilde{v}(t)v_1({x}/{a} )v_2 ({y}/{b} ),

\\

w(x,y,t)=w_{11}(t)\sin \frac{\pi x}{a} \sin\frac{\pi y}{b} +w_{12}(t)\sin \frac{\pi x}{a} \sin\frac{2\pi y}{b},

\end{gather}

где 

\begin{gather}

u_1({x}/{a}) =( {x}/{a})^4-2( {x}/{a})^3+( {x}/{a})^2,\quad

u_2({y}/{b}) =-({2}/{3})( {y}/{b})^3+( {y/}{b})^2,\\ 

v_1({x}/{a}) =({x}/{a})^4-2( {x}/{a})^3+( {x}/{a})^2,

\quad 

v_2({y}/{b}) =({y}/{b})( {y}/{b}-1).

\end{gather}



Первые два уравнения системы \eqref{bel_eq1} на основании подстановок \eqref{bel_appr1} с последующим применением процедуры Галеркина \cite{myl-bib:27,  myl-bib:29}  перепишутся в виде 

\begin{gather}

e_{11}\tilde{u}(t)+e_{12}\tilde{v}(t)=k_{11}a\Bigl[ \pi \frac{b}{a}\mathit{I}_4 w_{11}(t)+\pi\frac{b}{a}\mathit{I}_5 w_{12}(t) \Bigr],

\\

\label{bel_eq7}

e_{21}\tilde{u}(t)+e_{22}\tilde{v}(t)=\nu k_{11} a \bigl[ \pi \mathit{I}_9 w_{11}(t)+2 \pi \mathit{I}_{10} w_{12}(t)\bigr].

\end{gather}

Здесь введены следующие обозначения:

\begin{gather}

e_{11}=\frac{b}{a}\mathit{I}_1+\frac{1-\nu}{2}\frac{a}{b}\mathit{I}_2, \quad 

e_{12}=\frac{1+\nu}{2}\mathit{I}_3,\\

e_{21}=\frac{1+\nu}{2}\mathit{I}_6,\quad 

e_{22}=\frac{a}{b}\mathit{I}_7+\frac{1-\nu}{2}\frac{b}{a}\mathit{I}_8;

\\

\mathit{I}_1=\int_0^1 \int_0^1 \tilde{u}_{1 \bigcom 11} \tilde{u}_1 \tilde{u}_2^2dX dY,\quad 

\mathit{I}_2=\int_0^1 \int_0^1 \tilde{u}_{2  \bigcom 22}\tilde{u}_2\tilde{u}_1^2dX dY,

\\

\mathit{I}_3=\int_0^1 \int_0^1 \tilde{v}_{1\bigcom 1}\tilde{u}_1\tilde{u}_2 \tilde{v}_2,_{2}dX dY, \dots,

\\

\mathit{I}_{10}=\int_0^1 \int_0^1 \tilde{v}_1 \tilde{v}_2 \sin (\pi X )\cos (2\pi Y )dX dY,\quad  X={x}/{a}, \quad  Y={y}/{b}.

\end{gather}



Разрешая неоднородную алгебраическую систему относительно коэффициентов $\tilde{u}(t)$ и $\tilde{v}(t)$, получим

\begin{equation}\label{bel_eq9}

\tilde{u}=k_{11}a\bigl(L_1 w_{11}(t)+L_2 w_{12}(t) \bigr),\quad 

\tilde{v}=k_{11}a\bigl(L_3 w_{11}(t)+L_4 w_{12}(t) \bigr),

\end{equation}

где $L_j$ "--- безразмерные величины:

\begin{gather}

L_1=\frac{\pi \frac{b}{a}\mathit{I}_4e_{22}-\nu \pi \mathit{I}_9 e_{12} }{e_{11} e_{22}-e_{12}e_{21}},\quad L_2=\frac{\pi \frac{b}{a}\mathit{I}_5e_{22}-\nu 2 \pi \mathit{I}_{10} e_{12} }{e_{11} e_{22}-e_{12}e_{21}}, 

\\

L_3=\frac{-\pi \frac{b}{a}\mathit{I}_4e_{21}+\nu \pi \mathit{I}_9 e_{11} }{e_{11} e_{22}-e_{12}e_{21}},\quad L_4=\frac{-\pi \frac{b}{a}\mathit{I}_5e_{21}+\nu 2 \pi \mathit{I}_{10} e_{11} }{e_{11} e_{22}-e_{12}e_{21}}. 

\end{gather}

Окончательно  функции $u(x,y,t)$ и $v(x,y,t)$ примут следующий вид:

\begin{gather}\label{bel_eq11}

u(x,y,t)=k_{11} a \bigl(L_1 w_{11}(t)+L_2 w_{12}(t) \bigr) u_1({x}/{a} )u_2({y}/{b}),

\\

v(x,y,t)=k_{11} a \bigl(L_3 w_{11}(t)+L_4 w_{12}(t) \bigr) v_1({x}/{a} )v_2({y}/{b} ).

\end{gather}



Из третьего уравнения системы \eqref{bel_eq1}  на основании аналогичных преобразований получим систему обыкновенных однородных дифференциальных уравнений относительно коэффициентов $w_{11}(t)$ и $w_{12}(t)$:

\begin{equation}\label{bel_eq12}

\xi_{11}\left[w_{11} \right]+\xi_{12} w_{12}=0, \quad  \xi_{21}w_{11}+\xi_{22} \left[ w_{12}\right] =0,

\end{equation}

где $\xi_{kl}$ "--- дифференциальные операторы:

\begin{gather}

\xi_{11}=\frac{\gamma h a^4}{g D}\biggl( 1+2 \sum\limits_{i=1}^n \tilde{\beta}_i^s\biggr)

\frac{d^2}{dt^2}+\tilde{\mu} \frac{d}{dt}+ f_{11}\Bigl( k_{11}, \Theta_0, \frac{h_i}{h}\Bigr), 

\\

\xi_{12}=\frac{p_0 \varkappa M a^3}{D} 4\pi \frac{a}{b} \mathit{I}_{13}-\tilde{f}_{12}, \quad \xi_{21}=\frac{p_0 \varkappa M a^3}{D} 2\pi \frac{a}{b} \mathit{I}_{16}-\tilde{f}_{11},

\\

\xi_{22}=\frac{\gamma h a^4}{g D}

\biggl( 1+2 \sum\limits_{i=1}^n \tilde{\beta}_i^s\biggr)\frac{d^2}{dt^2}+\tilde{\mu} \frac{d}{dt}+ f_{12}

\Bigl( k_{11},\Theta_0,\frac{h_i}{h}\Bigr), 

\end{gather}

в которых

\begin{multline}

f_{11}=\Bigl(\pi^2+\Bigl( \frac{\pi a}{b}\Bigr) ^2 \Bigr) ^2+

12

\Bigl( \frac{h_i}{h}\Bigr)^2(k_{11}a)^2-

48

\Bigl( \frac{a}{h}\Bigr)^2 (k_{11}a)^2\Bigl(L_1\mathit{I}_{11}+\nu \frac{a}{b}L_3 \mathit{I}_{12} \Bigr) -

\\

-12(1-\nu^2)\Bigl( \frac{\pi a}{b}\Bigr)^2 \Bigl(\frac{a}{h} \Bigr)^2

\biggl(1+\sum\limits_{i=1}^n \tilde{\beta}_i^s \biggr) \alpha \Theta_0+

\Bigl( \frac{\pi a}{b}\Bigr)^4 2\sum\limits_{i=1}^n \beta_i^s+

\\

+4(1-\nu)\pi^2\Bigl(\frac{\pi a}{b} \Bigr)^2  \sum\limits_{i=1}^n \beta_i^c,

\end{multline}

\begin{equation*}

\tilde{f}_{12}=48 \Bigl(\frac{a}{h} \Bigr)^2(k_{11}a)^2\Bigl(L_2 \mathit{I}_{11}+\nu\frac{a}{b}L_4\mathit{I}_{12}\Bigr),

\end{equation*}

\begin{equation*}

\tilde{f}_{11}=48 \Bigl(\frac{a}{h} \Bigr)^2(k_{11}a)^2\Bigl(L_1 \mathit{I}_{14}+\nu \frac{a}{b}L_3\mathit{I}_{15}\Bigr),

\end{equation*}

\begin{multline}

f_{12}=\Bigl(\pi^2+\Bigl( \frac{2\pi a}{b}\Bigr) ^2 \Bigr)^2-

12\Bigl( \frac{a}{h}\Bigr)^2(k_{11}a)^2-

48\Bigl( \frac{a}{h}\Bigr)^2 (k_{11}a)^2 \Bigl(L_2\mathit{I}_{14}+\nu \frac{a}{b}L_4 \mathit{I}_{15} \Bigr) -

\\

-12(1-\nu^2)\Bigl( \frac{2\pi a}{b}\Bigr)^2

\Bigl(\frac{a}{h} \Bigr)^2

\biggl(1+\sum\limits_{i=1}^n \tilde{\beta}_i^s \biggr) \alpha \Theta_0+\Bigl( \frac{2\pi a}{b}\Bigr)^4 2\sum\limits_{i=1}^n \beta_i^s+

\\

+4(1-\nu)\pi^2\Bigl(\frac{2\pi a}{b} \Bigr)^2  \sum\limits_{i=1}^n \beta_i^c, 

\end{multline}

\begin{equation*}

\beta_i^c=\Bigl( \frac{h_i}{h}\Bigr) ^3\frac{a_i}{a}\varPhi_{3i}\cos^2 \frac{\pi x_i}{a},

\quad 

\beta_i^s=\Bigl( \frac{h_i}{h}\Bigr) ^3\frac{a_i}{a}\varPhi_{3i}\sin^2 \frac{\pi x_i}{a},

\quad 

\tilde{\beta}_i^s= \frac{h_i}{h}\frac{a_i}{a}\sin^2 \frac{\pi x_i}{a}.

\end{equation*}



Система дифференциальных уравнений \eqref{bel_eq12} подстановкой 

\begin{equation*}

w_{11}=-\xi_{12}\varPhi(t), \quad w_{12}=\xi_{11}\left[\varPhi(t) \right] 

\end{equation*}

сводится к одному дифференциальному уравнению четвертого порядка, которое в случае отсутствия  демпфирования ($\mu=0$) содержит только четные производные от функции $\varPhi(t)$:

\begin{multline}\label{bel_eq15}

\biggl[ 

\frac{\gamma h a^4}{g D}\biggl( 1+2\sum\limits_{i=1}^n \tilde{\beta}_i^s\biggr) \biggr]^2

\frac{d^4\varPhi}{dt^4} +

\biggl[ \frac{\gamma h a^4}{g D}\biggl( 1+2\sum\limits_{i=1}^n \tilde{\beta}_i^s\biggr)

( f_{11}+f_{12})  \biggr]\frac{d^2\varPhi}{dt^2}+

\\

+\Bigl[ f_{11}f_{12}-8\pi \Bigl( \frac{a}{b}\Bigr)^2\mathit{I}_{13}\mathit{I}_{16}\tilde{M}+\tilde{M} 

\Bigl(4 \pi \frac{a}{b}\mathit{I}_{13}\tilde{f}_{11}+2 \pi \frac{a}{b}\mathit{I}_{16}\tilde{f}_{12} \Bigr) - \tilde{f}_{11}\tilde{f}_{12}\Bigr] \varPhi=0,

\end{multline}

где $\tilde{M}=\frac{p_0 \varkappa M a^3}{D}$.



Если $\mu\neq0$, то дифференциальное уравнение для функции $\varPhi(t)$ запишется в~виде

\begin{equation}\label{bel_eq16}

\frac{d^4\varPhi}{dt^4}+2 e_1\frac{d^3\varPhi}{dt^3}+e_3\frac{d^2\varPhi}{dt^2}+e_4\frac{d\varPhi}{dt}+e_5\varPhi=0,

\end{equation}

где

\begin{equation*}

e_1=\frac{\tilde{\mu} }{ \frac{\gamma h a^4}{g D}\left(1+2\sum_{i=1}^n\tilde{\beta}_i^s \right) }, \quad e_3=\frac{\frac{\gamma h a^4}{g D}\left(1+2\sum_{i=1}^n\tilde{\beta}_i^s \right)\left( f_{11}+f_{12}\right)+ \tilde{\mu}}{\left( \frac{\gamma h a^4}{g D}\left(1+2\sum_{i=1}^n\tilde{\beta}_i^s \right)\right)^2},

\end{equation*}

\begin{equation*}

e_4=\frac{\tilde{\mu}\left(f_{11}+f_{12} \right)  }{\left(  \frac{\gamma h a^4}{g D}\left(1+2\sum_{i=1}^n\tilde{\beta}_i^s \right)\right)^2},\quad e_5=\frac{f_{11} f_{12}-\xi_{21}\xi_{12}}{\left( \frac{\gamma h a^4}{g D}\left(1+2\sum_{i=1}^n\tilde{\beta}_i^s \right)\right)^2}, \quad \tilde{\mu}=\frac{\mu a^4 h}{D}.

\end{equation*}



Из условия, при котором хотя бы один корень из четырех  характеристического уравнения для дифференциального уравнения \eqref{bel_eq15} будет положительным, определяется наименьшее значение  относительной скорости потока газа ${v_y}/{c_0}$, при котором прогиб термоупругой системы растет. При учете демпфирования ($\mu\neq0$) возникает вопрос об устойчивости системы \eqref{bel_eq12}, которая сведена к одному дифференциальному уравнению \eqref{bel_eq16}. 

В~этом случае это значение скорости потока определяется на основании критерия Гурвица~\cite{myl-bib:32}:

\begin{equation*}

\Delta_1=2e_1>0, \quad \Delta_2=2e_1e_3-e_4>0, \quad \Delta_3=2e_1e_3e_4-e_4^2-4e_1^2e_5>0.

\end{equation*}



Кривизна пологой цилиндрической оболочки задавалась в виде $k_{11}\hm=-{4 \tilde{\delta}}/{a^2}$ \cite{myl-bib:33, myl-bib:34}, где $\tilde{\delta}$ "--- наибольший подъем оболочки над ее планом. Выражение для относительной кривизны $\tilde{k}_{11}$ принимает вид

\begin{equation*}

\tilde{k}_{11}=k_{11} a=-4\frac{\tilde{\delta}}{h}\frac{h}{a},

\end{equation*}

где ${\tilde{\delta}}/{h}\in[0, 5]$ \cite{myl-bib:33, myl-bib:34}.



\smallskip



{\bf 2.} Численные результаты приведены в~таблице. 

Вычисления проводились при следующих значениях параметров: 

${h}/{a}=0.01$, 

${a_i}/{a}=0.01$, 

$\mu=0.001$, 

${a}/{b}=1$.





\begin{table}[b!]

	

\small \centering



Значения параметра ${v_y}/{c_0}$  в зависимости от других параметров



[The values of the parameter ${v_y}/{c_0}$ depending on other parameters]

	

\smallskip

	\renewcommand{\tabcolsep}{0.35cm}

\begin{tabular}{c|c|c|c|c|c|c|c}

	\hline 

	\multicolumn{2}{c|} {\rule{0mm}{12pt}}&  \multicolumn{3}{c|} {${h_i}/{h}=2.5$}    &  \multicolumn{3}{c} {${h_i}/{h}=5$}   \\ 

	\hline 

 \backslashbox{$\tilde{k}_{11}$}{$n$} & 0 & 1 & 3 & 5 & 1 & 3 & 5 \\ 

	\hline 

		\multicolumn{8}{c} {\rule{0mm}{12pt}for $\Theta_0=0$,  $\varepsilon_1=\varepsilon_2=0$} \\

	\hline 	

	0.10 & 0.24 & 1.50 & 3.16 & 3.50 & 6.47 & 13.35 & 14.79 \\ 

	0.15 & 0.54 & 1.40 & 3.06 & 3.40 & 6.37 & 13.25 & 14.69 \\

	0.20 & 0.97 & 1.26 & 2.92 & 3.26 & 6.23 & 13.11 & 14.55 \\  

	\hline 

		\multicolumn{8}{c} {\rule{0mm}{12pt}for $\Theta_0=2$,  $\varepsilon_1=\varepsilon_2=0$} \\

	\hline 

	0.10 & 1.08 & 0.61 & 2.22 & 2.54 & 5.55 & 12.34 & 13.73 \\ 

	0.15 & 1.39 & 0.51 & 2.12 & 2.44 & 5.45 & 12.24 & 13.63 \\ 

	0.20 & 1.81 & 0.37 & 1.98 & 2.30 & 5.31 & 12.10 & 13.49 \\ 

	\hline 

		\multicolumn{8}{c} {\rule{0mm}{12pt}for $\Theta_0=4$,  $\varepsilon_1=\varepsilon_2=0$} \\

	\hline 

		0.10 & 1.93 & 0.27 & 1.28 & 1.58 & 4.63 & 11.32 & 12.67 \\ 

		0.15 & 2.23 & 0.37 & 1.18 & 1.48 & 4.53 & 11.22 & 12.57 \\ 

		0.20 & 2.66 & 0.51 & 1.04 & 1.34 & 4.39 & 11.08 & 12.43 \\ 

	\hline 

		\multicolumn{8}{c} {\rule{0mm}{12pt}for $\Theta_0=0$,  $\varepsilon_1=\varepsilon_2=1$} \\

	\hline 		

		0.10 & 0.24 & 1.50 & 3.16 & 3.51 & 6.48 & 13.37 & 14.80 \\ 

		0.15 & 0.54 & 1.40 & 3.07 & 3.42 & 6.39 & 13.28 & 14.72 \\ 

		0.20 & 0.97 & 1.27 & 2.94 & 3.29 & 6.26 & 13.16 & 14.60 \\ 

	\hline 

		\multicolumn{8}{c} {\rule{0mm}{12pt}for $\Theta_0=2$,  $\varepsilon_1=\varepsilon_2=1$} \\

	\hline 	

		0.10 & 1.08 & 0.61 & 2.23 & 2.55 & 5.56 & 12.35 & 13.74 \\ 

		0.15 & 1.39 & 0.52 & 2.13 & 2.46 & 5.47 & 12.26 & 13.66 \\ 

		0.20 & 1.81 & 0.38 & 2.00 & 2.33 & 5.34 & 12.14 & 13.54 \\ 

	\hline 

		\multicolumn{8}{c} {\rule{0mm}{12pt}for $\Theta_0=4$,  $\varepsilon_1=\varepsilon_2=1$} \\

	\hline 	

		0.10 & 1.93 & 0.26 & 1.29 & 1.58 & 4.64 & 11.33 & 12.68 \\ 

		0.15 & 2.23 & 0.36 & 1.19 & 1.49 & 4.55 & 11.25 & 12.60 \\ 

		0.20 & 2.66 & 0.49 & 1.06 & 1.36 & 4.42 & 11.13 & 12.48 \\ 

		\hline 

	\end{tabular}  

\end{table} 



Количественный анализ выявил закономерности влияния геометрических параметров  и температуры  на величины относительных скоростей потока,  начиная с которых прогибы термоупругой системы увеличиваются во времени, и 

сделать нижеследующий выводы.

 

\begin{enumerate}

	\item[1.] Учет в дифференциальных уравнениях системы \eqref{bel_eq1} слагаемых, учитывающих <<растяжение"=сжатие>> ребер и их <<сдвиг>> в~тангенциальной плоскости (слагаемых, содержащих знаковые числа), практически не влияет на величину относительной скорости ${v_y}/{c_0}$ газового потока. По этой причине удерживать эти слагаемые в уравнениях нет необходимости.

	\item[2.] Увеличение температуры ведет к уменьшению предельной скорости потока при прочих равных условиях.

	\item[3.] Существенно повышают динамическую устойчивость геометрически нерегулярной оболочки параметры ${h_i}/{h}$ и число подкрепляющих элементов~$n$. 

	\item[4.] Величина относительной предельной скорости потока мало чувствительна к изменению относительной кривизны пологой оболочки.

\end{enumerate} 





Перечисленные закономерности выявлены для случая двух полуволн в~направлении потока и одной полуволны в перпендикулярном направлении.  

Следует отметить, что при $k_{11}=0$ полученные решения аналитически преобразуются к решениям, приведенным в работах \cite{myl-bib:31, myl-bib:35}. В случае <<холодной>> ($\Theta_0=0$) гладкой пластины   ($k_{11}=0$, ${a_i}/{a}=0$) решения принимают вид, приведенный в~\cite{myl-bib:30}.



\AdditionalContent{Конкурирующие интересы}{Заявляем, что в отношении авторства и публикации этой статьи конфликта интересов не имеем.} 



\AdditionalContent{Авторский вклад и ответственность}{Все авторы принимали участие в разработке концепции статьи и в написании рукописи. Авторы несут ответственность за предоставление окончательной рукописи в печать. Окончательная версия рукописи была одобрена всеми авторами.} 



\AdditionalContent{Финансирование}{Результаты получены в рамках выполнения государственного задания Минобрнауки России № 9.8570.2017/8.9.}



\begin{thebibliography}{10}



\RBibitem{myl-bib:1}

\by Вольмир~А.~С.

\book Оболочки в потоке жидкости и газа

\yr 1979

\publ Наука

\publaddr М.

\totalpages 320



\RBibitem{myl-bib:2}

\by Амбарцумян~C.~А., Багдасарян~Ж.~Е.

\paper Об устойчивости ортотропных пластинок, обтекаемых сверхзвуковым потоком газа

\jour Изв. АН СССР, ОТН, Механика и машиностроение

\yr 1961

\issue 4

\pages 91--96



\RBibitem{myl-bib:3}

\by Болотин~В.~В., Новичков~Ю.~Н.

\paper Выпучивание и установившийся флайтер термически сжатых панелей, находящихся в сверхзвуковом потоке

\jour Инж. журн.

\yr 1961

\issue 2

\pages 82--96







\RBibitem{myl-bib:4}

\by Мовчан~А.~А.

\paper О колебаниях пластинки, движущейся в газе

\jour ПММ

\yr 1956

\vol 20

\issue 2

\pages 211--222





\RBibitem{myl-bib:5}

\by Дун~Мин-дэ,

\paper Об устойчивости упругой пластинки при сверхзвуковом обтекании

\jour Докл. АН СССР

\yr 1958

\vol 120

\issue 4

\pages 726--729

\mathnet{http://mi.mathnet.ru/dan23108}



\RBibitem{myl-bib:5a}

\by Веденеев~В.~В. 

\paper Высокочастотный флаттер прямоугольной пластины

\jour Изв. РАН. МЖГ

\yr 2006

\issue 4

\pages 173--181









\RBibitem{myl-bib:6}

\by Огибалов~П.~М., Грибанов~В.~Ф.

\book Термоустойчивость пластин и оболочек

\yr 1968

\publ МГУ

\publaddr М.

\totalpages 520







\RBibitem{myl-bib:7}

\by Огибалов~П.~М. 

\book Вопросы динамики и устойчивости оболочек

\yr 1963

\publ МГУ

\publaddr М.

\totalpages 417

\zmath{0115.41504}





\RBibitem{myl-bib:8}

\by Болотин~В.~В.

\paper Температурное выпучивание пластин и пологих оболочек в сверхзвуковом потоке газа 

\inbook Расчеты на прочность

\yr 1960

\volinfo Вып.~6

\publaddr М.

\publ Машгиз

\pages 190--216



\RBibitem{myl-bib:9}

\by Жилин~П.~А.

\paper Линейная теория ребристых оболочек

\jour Изв. АН СССР. МТТ

\yr 1970

\issue 4

\pages 150--166



 



\RBibitem{myl-bib:10}

\by Белосточный~Г.~Н., Ульянова~О.~И.

\paper Континуальная модель композиции из оболочек вращения с термочувствительной толщиной

\jour Изв. РАН. МТТ

\yr 2011

\issue 2

\pages 184--191

\elib{https://elibrary.ru/item.asp?id=15785027}



 

 \RBibitem{myl-bib:11}

 \by Белосточный~Г.~Н., Рассудов~В.~М.

 \paper Континуальная модель термочувствительной ортотропной системы «оболочка–ребра» с учетом влияния больших прогибов

 \inbook Механика деформируемых сред

 \yr 1983

\volinfo Вып.~8

\publaddr Саратов

\publ  Сарат. политехн. ин-т

 \pages 10--22



  

\RBibitem{myl-bib:12}

\by Белосточный~Г.~Н., Рассудов~В.~М.

\paper Континуальный подход к термоустойчивости упругих систем «пластинка–ребра» 

\inbook Прикладная теория упругости

\yr 1980

\publaddr Саратов

\publ  Сарат. политехн. ин-т

\pages 94--99





 









\RBibitem{myl-bib:15}

\by Жилин~П.~А.

\paper Общая теория ребристых оболочек. Прочность гидротурбин

\inbook Тр. ЦКТИ

\yr 1968

\volinfo Вып.~8

\publaddr Л.

\pages 46--70







\RBibitem{myl-bib:16}

\by Карпов~В.~В., Сальников~А.~Ю.

\paper Вариационный метод вывода нелинейных уравнений движения пологих ребристых оболочек 

\jour Вестн. гражд. инженеров

\yr 2008

\issue 4(17)

\pages 121--124

\elib{https://elibrary.ru/item.asp?id=11675045}







\RBibitem{myl-bib:18}

\by  Белосточный~Г.~Н., Мыльцина~О.~А.

\paper Уравнения термоупругости композиций из оболочек вращения 

\jour Вестник СГТУ

\yr 2011

\issue 4 (59)

\issueinfo Вып.~1

\pages 56--64



   

\RBibitem{myl-bib:19}

\by  Онанов~Г.~Г.

\paper Уравнения с сингулярными коэффициентами типа дельта-функции и~ее производных

\jour Докл. АН СССР

\yr 1970

\vol 191

\issue 5

\pages 997--1000

\mathnet{http://mi.mathnet.ru/dan35339}

\mathscinet{http://www.ams.org/mathscinet-getitem?mr=0470298}

\zmath{https://zbmath.org/?q=an:0215.20101}









\RBibitem{myl-bib:21}

\by  Белосточный~Г.~Н.

\paper  Аналитические методы определения замкнутых интегралов сингулярных дифференциальных уравнений термоупругости геометрически нерегулярных оболочек

\jour Докл. Акад. воен. наук

\yr 1999

\issue 1

\pages 14--26





\Bibitem{myl-bib:22a}

\by Geckeler~J.~W.

\paper Elastostatik

\lang In German

\zmath{54.0844.01}

\inbook Handbuch der Physik 

\bookvol 6

\pages 141--308

\yr 1928





\RBibitem{myl-bib:22}

\by Ильюшин~А.~А.

\paper  Закон плоских сечений в аэродинамики больших сверхзвуковых скоростей 

\jour ПММ

\yr 1956

\issue 6

\pages 733--755







 

\RBibitem{myl-bib:30}

\by Вольмир~А.~С. 

\book  Устойчивость деформируемых систем

\yr 1967

\publ Наука

\publaddr М.

\totalpages 984



 

 



 

\RBibitem{myl-bib:25}

\by Владимиров~В.~С. 

\book Обобщенные функции в математической физике

\yr 1979

\publ Наука

\publaddr М.

\totalpages 319

\zmath{0515.46033}









\RBibitem{myl-bib:27}

\by Канторович~Л.~В., Крылов~В.~И.

\book  Приближенные методы высшего анализа

\yr 1962

\publ Физматлит

\publaddr М.

\totalpages 708

\zmath{0100.33102}









\Bibitem{myl-bib:29}

\by Rektorys~K.

\book Variational methods in mathematics, science and engineering

\zmath{0481.49002}

\publaddr Dordrecht, Boston,  London

\publ D.~Reidel Publ. 

\totalpages 571 

\yr 1980

\moreref

\crossref{10.1007/978-94-011-6450-4}. \nofrills





\RBibitem{myl-bib:31}

\by Белосточный~Г.~Н., Рассудов~В.~М. 

\paper  Термоупругие системы типа «пластинка–ребра» в сверхзвуковом потоке газа  

\inbook Прикладная теория упругости

\publaddr Саратов

\volinfo Вып.~8

\publ Сарат. политехн. ин-т

\yr 1983

\pages 114--121







\RBibitem{myl-bib:32}

\by Егоров~К.~В.  

\book  Основы теории автоматического регулирования

\yr 1967

\publ Энергия

\publaddr М.

\totalpages 648





\RBibitem{myl-bib:33}

\by Рассудов~В.~М., Красюков~В.~П., Панкратов~Н.~Д.  

\book  Некоторые задачи термоупругости пластинок и пологих оболочек 

\yr 1973

\publ Сарат. ун-т

\publaddr Саратов

\totalpages 155





\RBibitem{myl-bib:34}

\by Назаров~А.~А.

\book Основы теории и методы расчета пологих оболочек

\yr 1966

\publ  Стройиздат

\publaddr Л., М.

\totalpages 304





\RBibitem{myl-bib:35}

\by Мыльцина~О.~А., Белосточный~Г.~Н. 

\paper  Устойчивость нагретой ортотропной геометрически нерегулярной пластинки в сверхзвуковом потоке газа 

\jour Вестник Пермского национального исследовательского политехнического университета. Механика

\yr 2017

\issue 4

\pages 109--120

\crossref{10.15593/perm.mech/2017.4.08}















\end{thebibliography}





\newpage



\makeentitle



\vspace{-5mm}





\AdditionalContent{Competing interests}{We declare that we have no conflicts of interest in the authorship and publication of this article.} 



\AdditionalContent{Authors' contributions and responsibilities}{Each author has participated in the article concept development and in the manuscript writing. The authors are absolutely responsible for submitting the final manuscript in print. Each author has approved the final version of manuscript.} 





\AdditionalContent{Funding}{The results have been obtained within the State Assignment of the Ministry of Education

and Science of the Russian Federation no. 9.8570.2017/8.9.}









\begin{thebibliography}{10}



\Bibitem{myl-bib:1:en}

\lang In Russian

\by Vol'mir~A.~S.

\book Obolochki v potoke zhidkosti i gaza \rm [The shells in flow liquid and gas] 

\publaddr Moscow

\publ Nauka

\yr 1979

\totalpages 320



\Bibitem{myl-bib:2:en}

\lang In Russian

\by Ambartsumyan~S.~A. Bagdasaryan~Zh.~E.

\paper Of the stability of orthotropic plates streamlined by a supersonic gas flow

\jour Izv. Akad. Nauk SSSR, OTN, Mekhanika i Mashinostroenie

\yr 1961

\issue 4

\pages 91--96



\Bibitem{myl-bib:3:en}

\lang In Russian

\by Bolotin~V.~V., Novichkov~Yu.~N. 

\paper Buckling and steady-state flayter thermally compressed panels in supersonic flow

\jour Inzhenernyi Zhurnal 

\yr 1961

\issue 2

\pages 82--96







\Bibitem{myl-bib:4:en}

\lang In Russian

\by Movchan~A.~A. 

\paper Oscillations of the plate moving in the gas

\jour Prikladnaia Matematika i Mekhanika

\yr 1956

\vol 20

\issue 2

\pages 211--222





\Bibitem{myl-bib:5:en}

\lang In Russian

\by Tung~Ming-teh,

\paper The stability of an elastic plate in a supersonic flow

\jour Dokl. Akad. Nauk SSSR

\yr 1958

\vol 120

\issue 4

\pages 726--729



\Bibitem{myl-bib:5a}

\by Vedeneev~V.~V.

\jour Fluid Dyn.

\crossref{10.1007/s10697-006-0083-2}

\yr 2006

\vol 41

\issue 4

\pages 641–648 

\paper High-frequency flutter of a rectangular plate









\Bibitem{myl-bib:6en}

\lang In Russian

\by Ogibalov~P.~M., Gribanov~V.~F. 

\book Termoustoichivost' plastin i obolochek \rm  [Thermostability of plates and shells]

\yr 1968

\publaddr Moscow

\publ Moscow State Univ.

\totalpages 520







\Bibitem{myl-bib:7en}

\lang In Russian

\by Ogibalov~P.~M.

\book Voprosy dinamiki i ustoichivosti obolochek\rm  [The problems of dynamics and stability of shells]

\publaddr Moscow

\publ Moscow State Univ.

\yr 1963

\totalpages 417







\Bibitem{myl-bib:8:en}

\lang In Russian

\by Bolotin~V.~V.

\paper Thermal buckling of plates and shallow shells in a supersonic gas flow

\inbook Raschety na prochnost'

\publaddr Moscow

\publ Mashgiz

\yr 1960

\volinfo Issue~6

\pages 190--216



\Bibitem{myl-bib:9:en}

\lang In Russian

\by Zhilin~P.~A.

\paper Linear theory of ribbed shells

\jour Izv. Akad. Nauk SSSR, Mekhanika Tverdogo Tela

\yr 1970

\issue 4

\pages 150--166



 



\Bibitem{myl-bib:10:en}

\paper Continuum model for a composition of shells of revolution with thermosensitive thickness

\by Belostochnyi~G.~N., Ul'yanova~O.~I.

\jour Mechanics of Solids

\yr 2011

\vol 46

\issue 2

\pages 184--191

\elib{https://elibrary.ru/item.asp?id=16982711}

\crossref{10.3103/S0025654411020051}



 

\Bibitem{myl-bib:11:en}

\lang In Russian

\by Belostochnyi~G.~N., Rassudov~V.~M.

\paper Continuum model of orthotropic heat-sensitive ``shell-ribs'' system taking into account the influence of large deflection

 \yr 1983

\volinfo Issue~8

\inbook Prikladnaia teoriia uprugosti

\publ Saratov Polytechnic Inst.

\publaddr Saratov 

 \pages 10--22



  

\Bibitem{myl-bib:12:en}

\lang In Russian

\by Belostochnyi~G.~N., Rassudov~V.~M.

\paper Continuum approach to the thermal stability of elastic ``plate-ribs'' systems

\inbook Prikladnaia teoriia uprugosti

\publ Saratov Polytechnic Inst.

\publaddr Saratov 

\yr 1980

\pages 94--99







\Bibitem{myl-bib:15:en}

\lang In Russian

\by Zhilin~P.~A.

\paper General theory of ribbed shells. The strength of turbines

\inbook  Proc. of the Central Turbine-Boiler Institute

\yr 1968

\volinfo Issue~8

\publaddr Leningrad

\pages 46--70







\Bibitem{myl-bib:16:en}

\lang In Russian

\by Karpov~V.~V., Sal'nikov~A.~Yu.

\paper Variational method output nonlinear equations of motion of shallow ribbed shells

\jour Vestnik Grazhdanskikh Inzhenerov

\yr 2008

\issue 4(17)

\pages 121--124





\Bibitem{myl-bib:18:en}

\lang In Russian

\by Belostochnyi~G.~N., Myltcina~O.~A.

\paper Thermoelasticity equations of shells compositions

\jour Vestnik Saratovskogo Tekhnicheskogo Universiteta

\yr 2011

\issue 4 (59)

\issueinfo Issue~1

\pages 56--64



   

\Bibitem{myl-bib:19:en}

\lang In Russian

\by Onanov~G.~G.

\paper Equations with singular coefficients of the type of the delta-function and its derivatives

\jour Dokl. Akad. Nauk SSSR

\yr 1970

\vol 191

\issue 5

\pages 997--1000









\Bibitem{myl-bib:21:en}

\lang In Russian

\by Belostochnyi~G.~N.

\paper Analytical methods for the determination of closed integrals of singular differential equations of thermoelasticity, geometrically irregular shells

\jour Doklady Akademii Voennykh Nauk

\yr 1999

\issue 1

\pages 14--26





\Bibitem{myl-bib:22a:en}

\by Geckeler~J.~W.

\paper Elastostatik

\lang In German

\zmath{54.0844.01}

\inbook Handbuch der Physik 

\bookvol 6

\pages 141--308

\yr 1928





\Bibitem{myl-bib:22:en}

\lang In Russian

\by Il'yushin~A.~A.

\paper Law of plane sections in aerodynamics of high supersonic velocity

\jour Prikladnaia Matematika i Mekhanika

\yr 1956

\issue 6

\pages 733--755







 

\Bibitem{myl-bib:30:en}

\lang In Russian

\by Vol'mir~A.~S.

\book  Ustoichivost' deformiruemykh sistem \rm [Stability of deformable systems]

\publaddr Moscow

\publ Nauka

\yr 1967

\totalpages 984



 

 



 

\Bibitem{myl-bib:25:en}

\by Vladimirov~V.~S.

\book Generalized functions in mathematical physics

\zmath{0515.46034}

\publaddr Moscow

\publ Mir Publ.

\totalpages  362 

\yr 1979









\Bibitem{myl-bib:27:en}

\by Kantorovich~L.~V., Krylov~V.~I.

\book Approximate methods of higher analysis

\zmath{0083.35301}

\publaddr Groningen

\publ P.~Noordhoff Ltd.

\totalpages xii+681

\yr 1958











\Bibitem{myl-bib:29:en}

\by Rektorys~K.

\book Variational methods in mathematics, science and engineering

\zmath{0481.49002}

\publaddr Dordrecht, Boston,  London

\publ D.~Reidel Publ. 

\totalpages 571 

\yr 1980

\moreref

\crossref{10.1007/978-94-011-6450-4}. \nofrills





\Bibitem{myl-bib:31:en}

\lang In Russian

\by Belostochnyi~G.~N., Rassudov~V.~M.

\paper Thermoelastic system of ``plate-ribs'' type  in a supersonic gas flow

\inbook Prikladnaia teoriia uprugosti

\publ Saratov Polytechnic Inst.

\publaddr Saratov 

\volinfo Issue~8

\yr 1983

\pages 114--121







\Bibitem{myl-bib:32:en}

\lang In Russian

\by Egorov~K.~V.

\book Osnovy teorii avtomaticheskogo regulirovaniia \rm [Fundamentals of the theory of automatic control]

\publaddr Moscow

\publ Energiia

\yr 1967

\totalpages 648





\Bibitem{myl-bib:33:en}

\lang In Russian

\by Rassudov~V.~M., Krasiukov~V.~P., Pankratov~N.~D.  

\book  Nekotorye zadachi termouprugosti plastinok i pologikh obolochek  \rm [Some problems of thermoelasticity of plates and flat covers]

\yr 1973

\publ Saratov Univ.

\publaddr Saratov

\totalpages 155

 





\Bibitem{myl-bib:34:en}

\lang In Russian

\by Nazarov~A.~A.

\book Osnovy teorii i metody rascheta pologikh obolochek \rm [Fundamentals of the Theory and Methods for Designing Shallow Shells]

\yr 1966

\publ  Stroiizdat

\publaddr Leningrad, Moscow

\totalpages 304





\Bibitem{myl-bib:35:en}

\lang In Russian

\by Myltcina~O.~A., Belostochnyi~G.~N.

\paper Stability of heated orthotropic geometrically irregular plate in a supersonic gas flow

\jour PNRPU Mechanics Bulletin

\yr 2017

\issue 4

\pages 109--120

\crossref{10.15593/perm.mech/2017.4.08}





\end{thebibliography}







%\end{document}

 \newpage

%\input{./preamble}
%
%\usepackage{samgtu-name}
%
\vsgtu{1624} %registration number
\FirstPage{762}
\LastPage{773}
%
%\pdfcrop
%\draft  
%
\ShortCommunication
%
%\begin{document} 




\hfill \qrcode[hyperlink,height=.8in]{http://doi.org/10.14498/vsgtu1624}
\vspace{-.9in}

\udc{539.3}
\msc{74C10}

\rutitle{Об одном подходе к определению предельной несущей~способности
	механических систем с~разупрочняющимися элементами}
\rutitleshort{Об одном подходе к определению предельной несущей способности
	механических систем\dots}

\entitle{One approach to determination of the ultimate load-bearing capacity of mechanical systems with softening elements}
\entitleshort{One approach to determination of the ultimate load-bearing capacity of mechanical systems\dots}

\ruauthor{В.~В.~Стружанов\affil{1}, А.~В.~Коркин\affil{2}, А.~Е.~Чайкин\affil{2}}
{С\,т\,р\,у\,ж\,а\,н\,о\,в~В.~В., К\,о\,р\,к\,и\,н~А.~В., Ч\,а\,й\,к\,и\,н~А.~Е.}
\enauthor{V.\,V.~Struzhanov\affil{1}, A.\,V.~Korkin\affil{2}, A.\,E.~Chaykin\affil{2}} {S\,t\,r\,u\,z\,h\,a\,n\,o\,v~V.~V., K\,o\,r\,k\,i\,n~A.~V., C\,h\,a\,y\,k\,i\,n~A.~E.}

\ruaffil{\ruimash.
\and 
\ruurfuien.}

\enaffil{\enimash.
\and 
\enurfuien.}

\ruauthorsdetails{3}{% Количество авторов
\fullauthor{\rupid{39278}{Валерий Владимирович Стружанов}}
\CAuthor % Автор-корреспондент
\orcid{0000-0002-3669-2032} % ORCID ID автора 
\regalia{доктор физико-математических наук, профессор} 
\position{главный научный сотрудник} 
\dept[p]{~лаб. микромеханики материалов} 
\email{stru@imach.uran.ru}
\fullauthor{\rupid{73471}{Александр Владимирович Коркин}}
\orcid{0000-0003-3533-4257} % ORCID ID автора 
\regalia{аспирант} 
\email{alexkorkin@list.ru}
\fullauthor{\rupid{143469}{Алексей Евгеньевич Чайкин}}
\orcid{0000-0001-5582-2384} % ORCID ID автора 
\regalia{аспирант}  
\email{chaykin.ae@yandex.ru}\vspace{-3mm}
}

\enauthorsdetails{3}{
\fullauthor{\enpid{39278}{Valery V. Struzhanov}}
\CAuthor % Сorresponding Author
\orcid{0000-0002-3669-2032} % ORCID ID avtora 
\regalia{Dr. Phys. \& Math. Sci., Professor} 
\position{Chief Reseacher} 
\dept{Lab. of Matherial Micromechanics}
\email[br]{stru@imach.uran.ru}
\fullauthor{\enpid{73471}{Aleksandr V. Korkin}}
\orcid{0000-0003-3533-4257} % ORCID ID avtora 
\regalia{Postgraduate Student} 
\email{alexkorkin@list.ru}
\fullauthor{\enpid{143469}{Aleksey E.  Chaykin}}
\orcid{0000-0001-5582-2384} % ORCID ID avtora 
\regalia{Postgraduate Student} 
\email{chaykin.ae@yandex.ru}
}

\rukeywords{потенциальная функция, вырожденные критические точки, единый потенциал, матрица Гессе, тонкостенный резервуар, предельное давление}
\enkeywords{potential function, degenerate critical points, unified potential, Hesse matrix, thin walled reservoir, limiting pressure}

\ruabstract{Изложены основные положения теории расчета предельных нагрузок, действующих на дискретные механические системы с~разупрочняющимися элементами. Методика опирается на численное определение вырожденных критических точек потенциальной функции системы, где происходит переход от устойчивости процесса нагружения к~неустойчивости (катастрофа, разрушение), и~позволяет избежать решения большого числа нелинейных уравнений равновесия. 
В~качестве примера применения методики решена задача об определении предельного внутреннего давления в~тонкостенном цилиндрическом резервуаре. При построении потенциальной функции системы использован специально построенный единый потенциал для плоского квадратного элемента материала в~условиях двухосного растяжения, описывающий все стадии деформирования, включая и~разупрочнение.}

\enabstract{The fundamental provisions of the limiting load calculation theory are presented for a discrete mechanical system with softening elements. The method is based on the numerical determination of degenerate critical points for the potential function of the system. At these points there is a transition from the stability of the loading process to instability such as a catastrophe or a failure. This approach helps to avoid solving a large number of nonlinear equilibrium equations. The problem of determining the limiting internal pressure in a thin walled cylindrical tank is solved as an example. A unified potential specially defined for a flat square element of material in biaxial tension is used in developing a potential function of the system. It describes all stages of deformation including the softening stage.}

\newdate{struzha}{11}{05}{2018}
\date{\displaydate{struzha}}
\newdate{struzhb}{11}{10}{2018}
\revisiondate{\displaydate{struzhb}}
\newdate{struzhc}{12}{11}{2018}
\accepteddate{\displaydate{struzhc}}

\newdate{struzhe}{11}{12}{2018}
\renewcommand{\JournalDataOnline}{\displaydate{struzhe}}
%\ForPeerReview
\makerutitle

\newpage

\Section[n]{Введение}
Вычисление предельной несущей способности механических систем основывается на расчете параметров равновесных состояний при постепенно возрастающей нагрузке. Затем эти параметры тем или иным способом сводят к некоторому скаляру и сравнивают его с предельной величиной, полученной в эксперименте. Например, на каждом шаге догружения определяют напряженное состояние элемента конструкции, находят соответствующий инвариант тензора напряжений, который сравнивают с его предельной величиной, взятой из эксперимента.

Возможен и другой подход, основанный на исследовании устойчивости положений равновесия, когда разрушение связывают с невозможностью равновесия, то есть с моментом потери устойчивости [\citen{bib1, bib2}]. Реализация данного подхода возможна, если учитывать физически неустойчивые состояния материала (разупрочнение) [\citen{bib2}--\citen{bib6}]. Но и в этом случае возникает необходимость в решении не удовлетворяющих условиям корректности по Адамару [\citen{bib8}]  больших систем нелинейных уравнений равновесия для определения критических точек потенциальной функции механической системы, что является нетривиальной задачей [\citen{bib9}]. Кроме того, выделение из них вырожденных критических точек, где происходит смена устойчивости на неустойчивость, требует применения нетрадиционного аппарата теории катастроф [\citen{bib10,bib11}]. Однако информация о параметрах равновесия (критических точках потенциальной функции системы) при возрастающих нагрузках, вообще говоря, является излишней. При определении предельной несущей способности нужно знать только вырожденные критические точки. В данной работе представлена методика, которая позволяет избежать аналитического решения нелинейных уравнений. 
Она основана на специальной числовой процедуре приближенного выбора вырожденных критических точек. В качестве примера рассмотрена задача о~разрушении цилиндрического резервуара под действием внутреннего давления. Для построения потенциальной функции разработана методика построения единого потенциала для элемента материала, находящегося в плоском напряженном состоянии и описывающего его свойства как при упрочнении, так и при разупрочнении.

\smallskip

\Section{Общие положения} Положение механической системы в пространстве определяется конечным набором обобщенных координат (параметров состояния)  $ q_i$, $i = 1, 2, \ldots, N $. 
Величины параметров состояния зависят от набора задаваемых параметров управления системой $ s_j$, 
$j = 1, 2, \ldots, M $ 
(внешние силы, перемещения некоторых точек системы, физико"=механические свойства элементов и т.~п.). 
Причем параметры состояния должны принимать такие значения, при которых механическая система находится в состоянии равновесия. 
Разрушение системы есть явление того же порядка, что и явление невозможности равновесия [\citen{bib1}]. 
Когда равновесие становится неустойчивым, то реализуется или скачкообразный переход в ближайшее устойчивое равновесие, или глобальная катастрофа (разрушение). Поэтому под предельными значениями параметров управления следует понимать такие их величины, при реализации которых нарушается устойчивость системы.

В тех случаях, когда механическая система является градиентной [\citen{bib11}], для определения предельных значений параметров управления можно воспользоваться следующей методикой. Градиентная система при квазистатическом нагружении с постоянной температурой описывается потенциальной функцией (лагранжианом) $ \Pi(q_i ,s_j) $, связывающей параметры состояния и~управления. 
Эта функция есть сумма потенциальных функций элементов системы, которые могут быть как вогнутыми, так и~выпукло"=вогнутыми. 
Выпуклые вниз (вогнутые) участки отвечают  устойчивым состояниям элемента, выпуклые (выпуклые вверх) "--- неустойчивым.  
Точки перехода от выпуклости к~вогнутости "--- пограничные точки, являющиеся в общем случае седловыми точками соответствующих потенциальных функций. 

Когда механическая система находится в~равновесии (устойчивом или неустойчивом), для связи параметров состояния и~управления в~этом равновесии можно использовать уравнение Лагранжа второго рода [\citen{bib12}]. 
Тогда уравнение равновесия имеет вид
\begin{equation}\label{S_eq1}
\frac{\partial\Pi(q_i,s_j)}{\partial q_i} = 0, \quad  i=1, 2, \ldots,N.
\end{equation}
Если функция   является выпукло"=вогнутой, то при заданных управляющих параметрах система \eqref{S_eq1} может иметь одно или более одного решения, как устойчивых, так и~неустойчивых. 
Следовательно, в~данном случае не выполняются условия Адамара [\citen{bib8}].

Все решения системы \eqref{S_eq1} для всевозможных параметров управления образуют совокупность критических точек потенциальной функции $ \Pi $.  
Известно [\citen{bib11}], что смена типа равновесия (устойчивого или неустойчивого) происходит в~вырожденных критических точках, которые определяются из совместного решения уравнений \eqref{S_eq1} и~уравнения, получающегося приравниванием к~нулю детерминанта матрицы устойчивости (матрицы Гессе потенциальной функции):
\begin{equation}\label{S_eq2}
\mathop{\rm det}\Bigl| \frac{\partial^{2}\Pi(q_i,s_j)}{\partial q_m \partial s_n}\Bigr|  = 0,  \quad m,n = 1, 2, \ldots, N.
\end{equation}
Очевидно, что уравнение \eqref{S_eq2} выделяет из множества критических точек вырожденные критические точки. 

Таким образом, для отыскания предельных значений параметров управления, при достижении которых система теряет устойчивость, сначала, как правило, вычисляются все критические точки (решения уравнений \eqref{S_eq1}), затем выделяются вырожденные точки и после подстановки их в уравнение \eqref{S_eq1} определяются соответствующие параметры управления. Те параметры управления, при которых в первый раз достигается вырожденная критическая точка, то есть происходит смена устойчивого состояния механической системы на неустойчивое, и будут предельными значениями внешних нагрузок.

При реализации данной системы нахождения предельных значений параметров управления необходимо искать решения большого числа нелинейных уравнений, что представляет достаточно сложную задачу. 
Однако задача о~вычислении предельных значений параметров управления не требует, вообще говоря, знания всех критических точек функции $ \Pi $ (параметров всех положений равновесия механической системы). Необходимы только вырожденные точки, удовлетворяющие уравнению \eqref{S_eq2}. Кроме того, для технических приложений можно использовать их приближенные значения, вычисленные с достаточной степенью точности. 

Вычисление приближенных значений вырожденных критических точек можно осуществить, применяя следующую численную процедуру. Рассмотрим евклидово пространство $ \mathbb R^M \times \mathbb R^N $ "--- декартово произведение евклидовых пространств размерности $ M $ и $ N $. 
Элементами этого пространства являются числа $ ( q_i,s_j) $. 
Задавая физически обоснованные пределы изменения параметров состояния и управления, получаем $ (M{+}N) $"=мерный куб в пространстве $\mathbb  R^M \times \mathbb R^N $, 
который разбиваем сеткой узлов на множество «кубиков» с~заданным шагом. 
Затем в~каждом узле сетки разбиения вычисляем значения компонент матрицы Гессе и~величину ее определителя. 
Таким образом, каждому узлу ставится в~соответствие значение детерминанта матрицы Гессе. 

После этого выделяем те узлы, в которых детерминант близок к нулю с достаточной степенью точности. 
Полученное множество $B$ точек пространства $ \mathbb R^M \times \mathbb R^N $ содержит в~себе все вырожденные критические точки потенциальной функции~$ \Pi $. 
Чтобы выделить из множества $B$ критические точки функции $ \Pi $, следует подставить элементы множества $B$ в~уравнения равновесия \eqref{S_eq1}.  
Те~из них, которые удовлетворяют этим уравнениям с~заданной степенью точности, будут вырожденными критическими точками. 
Параметры управления,  соответствующие ближайшей (по евкледовой метрике) к~началу координат вырожденной точке, образуют совокупность искомых предельных значений управлений. 
Если существует группа вырожденных точек, находящихся на одинаковом наименьшем расстоянии от начала координат, то получим различные сочетания значений параметров управления.

Таким образом, минуя стадию расчета параметров равновесных состояний, можно найти критические нагрузки, что достаточно для оценки прочности механических систем.

\smallskip

\Section{Модельный пример} Применим изложенную выше методику к определению предельных параметров управления (нагрузок) при растяжении модельной стержневой системы, изображенной на  рис.~\ref{Sfig1}. 
Стержень {\sl 2} "--- абсолютно упругий с жесткостью при растяжении, равной $c$. 
Стержень {\sl 1} изготовлен из разупрочняющегося материала.

\begin{figure}[b!]
\centering
		\includegraphics[width=0.4\textwidth]{fig_stru1}
		\caption{Растягиваемая стержневая система с разупрочняющимся элементом} \label{Sfig1}
		\smallskip \footnotesize
		[Figure~\ref{Sfig1}. 		A rod system with a softening element under tension]

\end{figure}

\smallskip

{\bf 2.1.}
При мягком нагружении растягивающая сила $ p $ "--- параметр управления, 
а~удлинение $ x $  стержня {\sl 1}  и перемещение $ u $ правого свободного конца   стержня~{\sl 2}  "--- параметры состояния (см.  рис.~\ref{Sfig1}). 
Сопротивление растяжению стержня {\sl 1} задано функцией $ g(x) $ 
(зависимость растягивающей силы от удлинения), обладающей как восходящей ветвью (упрочнение), так и~ветвью, падающей до нуля (стадия разупрочнения) [\citen{bib2, bib5}].  
В~этом случае его потенциальная функция, равная $\displaystyle \int ^{x}_{0} g(x)\, dx $, является выпукло"=вогнутой, а лагранжиан системы~[\citen{bib5}] 
$$
Q^p = \int^{x}_{0} g(x)  dx + \frac{c}{2} (u-x)^2 - \int ^{u}_{0} p  du,
$$	
где второе слагаемое "--- энергия упругих деформаций стержня {\sl 2}, последнее "--- работа внешней силы, взятая со знаком минус.

Уравнения равновесия имеют вид
\begin{equation}\label{S_eq3}
Q^{p}_{\bigcom x} = g(x) - c(u-x) = 0, \quad  Q^{p}_{\bigcom u} = c(u-x) - p = 0.
\end{equation}
Здесь запятой обозначены частные производные по соответствующим аргументам.
Приравнивая к нулю детерминант матрицы Гессе, получим уравнение 
\begin{equation}\label{S_eq4}
g_{\bigcom x} = 0.
\end{equation}

Отметим, что это уравнение определяет максимум функции  $ g(x) $. 
Его можно разрешить, используя следующую процедуру. 

Проведем разбиение оси $X $ узлами с~достаточно малым шагом. 
Подставляя координаты этих узлов в~уравнение \eqref{S_eq4}, находим среди них значение $ x' $, приближенно ему удовлетворяющее. 
Затем из уравнений \eqref{S_eq3} для данного $ x' $ определяем $ p' $ "--- предельную величину растягивающей силы.

\smallskip

{\bf 2.2.}
В случае жесткого нагружения, когда растяжение осуществляется заданным перемещением $ u $ (параметр управления), а $ x $ является единственным параметром состояния, лагранжиан системы есть функция $ Q^u $, равная двум первым слагаемым в выражении для $ Q^p $. 

Тогда имеем лишь одно уравнение равновесия (первое уравнение в~системе \eqref{S_eq3}). 
В~матрице Гессе только один элемент. 
Тогда равенство детерминанта нулю дает уравнение $ g'_{x} + c = 0 $. 
Его приближенное решение ищем так же, как и~решение уравнения \eqref{S_eq4}. 

В зависимости от величины параметра $ c $ возможны три варианта. 

В первом "--- решения не существует. Следовательно, деформирование системы идет равновесно вплоть до разделения стержня {\sl 1} на части при значении $ x = x^z $, когда падающая ветвь функции $ g(x) $ достигает оси $ X $. Предельное значение параметра управления получаем из уравнения равновесия после подстановки $ x = x^z $: 
$$ u^z = g(x^z)/c+ x^z.$$ 

Во втором варианте имеется только одно приближенное решение $ x =  x_1 $. 
После достижения параметром $ x $ значения этой величины система переходит в~неустойчивое состояние. 
Поэтому предельный параметр управления $$ u_1 = g(x_1)/c + x_1 .$$ 

И, наконец, в третьем варианте имеем два решения $ x_1$, $x_2$  $(x_1 > x_2) $. 
Тогда предельное значение  $ u = u_1 $.

Отметим, что такие же результаты были получены и в работах [\citen{bib5, bib9}] только после решения нелинейных уравнений равновесия при постепенном возрастании параметров управления и построения соответствующих многообразий катастроф и сепаратрис, состоящих из вырожденных критических точек.

\smallskip

\Section{Вогнуто"=выпуклый потенциал при двухосном растяжении тонкой пластины} 
Для реализации положений теории, изложенной выше, необходимо знать вогнуто"=выпуклый потенциал элементов системы. 
В~механике сплошной среды эта функция должна быть определена для элемента материала вплоть до разупрочнения и~разделения на фрагменты.

Рассмотрим плоский квадратный элемент материала с~единичными размерами, находящийся в плоско"=напряженном состоянии, причем касательные напряжения и~деформации отсутствуют. 
Деформирование осуществляется заданием нормальных деформаций $ \varepsilon_1$, $\varepsilon_2 $ (двуосное растяжение). Если материал находится в состоянии упругости, то его  потенциальная энергия 
\begin{equation}\label{S_eq5}
V^e(\varepsilon_1, \varepsilon_2) = \frac{1}{2}(\sigma_1\varepsilon_1 + \sigma_2\varepsilon_2) = \frac{E}{2(1-\nu^2)} (\varepsilon^2_1 + 2\nu \varepsilon_1 \varepsilon_2 + \varepsilon^2_2 ) ,
\end{equation}
где $ E $ "--- модуль Юнга, $ \nu $ "--- коэффициент Пуассона. 
Здесь использован закон Гука, связывающий нормальные напряжения $ \sigma_1$, $\sigma_2 $ и~нормальные деформации $ \varepsilon_1$, $\varepsilon_2 $ и записанный в векторно"=матричной форме:
\begin{equation}\label{S_eq6}
\begin{pmatrix}
\sigma_1 \\ \sigma_2
\end{pmatrix}
 = C 
\begin{pmatrix} \varepsilon_1 \\ \varepsilon_2
\end{pmatrix}.
\end{equation}
В этом выражении 
$$
C = \frac{E}{(1-\nu^2)} \begin{pmatrix} 1 & \nu \\ \nu & 1 \end{pmatrix}
$$
можно трактовать как матричный коэффициент пропорциональности.

Формула \eqref{S_eq5} определяет кривые второго порядка на плоскости $\varepsilon_1 \varepsilon_2 $ 
(линии уровня потенциала $ V^e $), а~именно центральные эллипсы, главные оси которых наклонены к~декартовым осям $ \varepsilon_1$, $\varepsilon_2 $ под углом $ \pi/4 $. 
Большая ось располагается во втором и~четвертом квадрантах, малая "--- в~первом и~третьем. 
Для полуосей  выполняется отношение 
$$ \dfrac{a^{2}}{\rho^2} = \frac{(1-\nu)}{(1+\nu)}, $$ 
где $ a $ "--- малая полуось, $ \rho $ "--- большая полуось. 
Точки, расположенные на малой оси, отвечают равномерному растяжению $ (\varepsilon_1 = \varepsilon_2) $ или сжатию $ (-\varepsilon_1 = -\varepsilon_2) $, а~точки на большой оси определяют чистый сдвиг   
$ (\varepsilon_1 = -\varepsilon_2$, $ -\varepsilon_1 = \varepsilon_2 ) $. 
Произвольная точка  $ (\varepsilon_1,\varepsilon_2) $ лежит на линии уровня с большой полуосью, равной
\begin{equation}\label{S_eq7}
\rho = \frac{1}{\sqrt{2}}\sqrt{{(\varepsilon_1 + \varepsilon_2)}^2 \frac{1+\nu}{1-\nu} + {(\varepsilon_1 - \varepsilon_2)}^2} .
\end{equation}

При чистом сдвиге $ (\varepsilon_1 = \varepsilon$, $\varepsilon_2 = - \varepsilon)$  $\rho = \gamma /\sqrt{2}$, где    $ \gamma = \varepsilon_1 - \varepsilon_2 = 2 \varepsilon $ "--- сдвиг.
После исчерпания упругости для описания свойств материала можно было бы использовать теорию малых упруго-пластических деформаций [\citen{bib13, bib14}], в~которой предполагается упругость объемной деформации и~пропорциональность девиаторов тензоров напряжений и деформаций. 
Однако после выхода на стадию разупрочнения первая гипотеза не выполняется, так как материал уже существенно поврежден множеством микротрещин и объемный модуль изменяется. 

Чтобы остаться в рамках деформационной теории, следует отказаться от первой гипотезы, а~вместо второй сохранить пропорциональность напряжений и~деформаций с некоторым уже переменным, зависящим от деформаций матричным коэффициентом. 
Например, в рассматриваемой задаче %после перехода материала  при активном нагружении 
этот матричный коэффициент в законе \eqref{S_eq6} можно принять в~виде $ C \psi(\varepsilon_1,\varepsilon_2)$, 
где $ \psi $ "--- некоторая скалярная функция, требующая специального определения.

Следствием принятого предположения является тот факт, что в~области неупругости линии уровня функции энергии деформаций $ V $ (потенциальной энергии) подобны линиям уровня потенциала $ V^{e} $, то есть представляют собой эллипсы с~тем же самым отношением главных осей. 
Необходимо только знать распределение значений потенциала по  этим линиям уровня.

Воспользуемся полной диаграммой для чистого сдвига $ \tau(\gamma) $ ($ \tau $ "--- касательное напряжение, $ \gamma $ "--- деформация сдвига), которая обладает восходящей и нисходящей до нуля ветвями. 
Такую диаграмму материала можно получить, используя опытную диаграмму кручения стержней  круглого поперечного сечения с~пересчетом на диаграмму материала [\citen{bib15}]. 
Тогда энергия, затраченная на деформирование, есть   
$$ V = \int ^{\gamma}_{0}\tau(\gamma)  d\gamma .$$

После взятия интеграла и замены $ \gamma $ на $ \sqrt{2} \rho $  с~учетом выражения \eqref{S_eq7} получаем искомый единый потенциал $ V(\varepsilon_1,\varepsilon_2) $. 

Пусть $ \tau = 0.1G   \gamma(1-5 \gamma)$, если $ 0 \le \gamma \le 0.2 $, где  
$ G = \frac{E}{2(1+\nu)} $ "---  модуль сдвига в упругости. 
Тогда
\begin{equation}\label{S_eq8}
V(\varepsilon_1,\varepsilon_2) = 
\left\lbrace  
\begin{array}{cl}
\frac{V^{e}}{2}\Bigl( 0.1 - \frac{1}{3}\sqrt{\frac{V^{e}}{G}}\Bigr), & (\varepsilon_1,\varepsilon_2) \in \Omega , \\
0.002 G/3 , & (\varepsilon_1,\varepsilon_2) \notin \Omega.
\end{array}
\right. 
\end{equation}


\begin{figure}[b!]
\begin{tabular}{p{0.35\textwidth}p{0.6\textwidth}}

\includegraphics[width=0.3\textwidth]{fig_stru2} & 

\vspace{-5.5cm}

~~~\,Связь между напряжениями и деформациями (определяющие соотношения) в данном случае будет задана законом
$$
\begin{pmatrix} \sigma_1  \\ 
\sigma_2 
\end{pmatrix}
=
\begin{pmatrix} \frac{\partial V}{\partial \varepsilon_1} \\ 
\frac{\partial V}{\partial \varepsilon_2} \end{pmatrix}
=
C \psi (\varepsilon_1,\varepsilon_2) 
\begin{pmatrix} \varepsilon_1 \\ 
\varepsilon_2 \end{pmatrix},
$$
где 
$$ 
\psi (\varepsilon_1,\varepsilon_2) = \frac{1}{2} \Bigl(0.1 - \frac{1}{2} \sqrt{\frac{V^e}{G}} \Bigr).  
$$


~~~\,Отметим, что напряжения являются компонентами градиента функции~$ V $.

\\

\vspace{-.8cm}
\caption{Качественный вид поверхности потенциала $ V(\varepsilon_1,\varepsilon_2) $ } \label{Sfig2}
		\smallskip \footnotesize
		[Figure~\ref{Sfig2}. Qualitative form of the \centerline{potential surface $ V(\varepsilon_1,\varepsilon_2) $]} & \\
\end{tabular}
\end{figure}

%\begin{figure}[b!]
%	\begin{center}
%		\includegraphics[width=0.3\textwidth]{fig_stru2}
%		\caption{Качественный вид поверхности потенциала $ V(\varepsilon_1,\varepsilon_2) $ } \label{Sfig2}
%		\smallskip \footnotesize
%		[Figure~\ref{Sfig2}. Qualitative form of the potential surface $ V(\varepsilon_1,\varepsilon_2) $]
%	\end{center}
%\end{figure}



Очевидно, что область $ \Omega $ ограничена эллипсом с~полуосями 
$ \rho = 0.141$; $a = \rho\sqrt{(1+\nu)/(1-\nu)} $.  
Качественный вид поверхности $ V(\varepsilon_1,\varepsilon_2) $ изображен на рис.~\ref{Sfig2}. 
Функция $ V $ описывает все состояния материала (выпуклость вниз "--- устойчивость, упрочнение; 
выпуклость вверх "--- неустойчивость, разупрочнение).





\smallskip


\Section{Разрушение цилиндрического резервуара} Рассмотрим тонкостенный резервуар в~виде длинной трубы с~заделанными абсолютно жесткими крышками торцами. 
Радиус трубы равен $ b $, длина $ l $, толщина стенки $ t$  $(t\ll b)$. 
Резервуар находится под действием  монотонно и квазистатически возрастающего при постоянной температуре внутреннего давления с~интенсивностью~$ p $ [\citen{bib16, bib17}]. 
Материал обладает свойством деформационного разупрочнения. 
Применим изложенную выше методику для расчета предельного значения давления. 

Каждый элемент трубы находится в~плоском напряженном состоянии при всестороннем растяжении: 
$ \varepsilon_1 = \varepsilon_z $ "--- продольная деформация, 
$ \varepsilon_2 = \varepsilon_{\theta} $ "--- тангенциальная деформация. Его потенциальная функция $ V $ задана выражением \eqref{S_eq8}. 
%где  $ \varepsilon_1 = \varepsilon_z , \varepsilon_1 = \varepsilon_{\theta} $. 

Механическая система (резервуар) обладает двумя параметрами состояния (обобщенными координатами) $ \varepsilon_z$, $\varepsilon_{\theta} $ (однородное напряженно"=деформированное состояние) и~одним параметром управления $ p $. 

Лагранжиан (потенциальная функция) этой системы
$$
W = V(\varepsilon_z , \varepsilon_{\theta})2\pi b l t - p \pi b^2 l \varepsilon_z - 2\pi p b^2 l \varepsilon_{\theta}.
$$
Здесь первое слагаемое "--- полная энергия деформаций стенок трубы, второе и третье "--- взятые со знаком минус соответственно работа сил давления на перемещении жестких крышек и~перемещении, равном увеличению радиуса. 
Тогда уравнения равновесия, определяющие критические точки функции $ W $, имеют вид
\begin{equation}\label{S_eq9}
\frac{\partial W}{\partial \varepsilon_z} = \frac{\partial V}{\partial \varepsilon_z} 2\pi b l t  - p \pi b^2 l = 0;
\quad  
\frac{\partial W}{\partial \varepsilon_{\theta}} = \frac{\partial V}{\partial \varepsilon_{\theta}} 2\pi b l t  - 2p \pi b^2 l = 0.
\end{equation}
Уравнения \eqref{S_eq9} сводятся к одному:
$$
2 \frac{\partial V}{\partial \varepsilon_z} = \frac{\partial V}{\partial \varepsilon_{\theta}}.
$$
Производя необходимые вычисления, находим, что
\begin{equation}\label{S_eq10}
\varepsilon_{\theta} = \frac{2-\nu}{1-2\nu} \varepsilon_{z},
\end{equation}
то есть в пространстве деформаций путь деформирования элемента материала является прямой линией, уравнение которой задается равенством \eqref{S_eq10}. 
Кроме этого, из уравнений \eqref{S_eq9} следует, что  $ \sigma_z = pb/2t$, $\sigma_{\theta} = pb/t $. 
Отсюда путь нагружения в~пространстве напряжений также является прямой линией. 
Таким образом, имеем ситуацию с пропорциональным нагружением элемента материала.

Определим, наконец, предельное значение внутреннего давления. 
Как следует из приведенных выше теоретических положений, для этого необходимо гессиан функции $ W $ приравнять к~нулю, найти все решения полученного уравнения, по крайней мере, приближенно, и выделить из них вырожденные критические точки потенциала $ W $, которые принадлежат множеству критических точек (решений уравнений \eqref{S_eq9}) потенциала. 
Имеем 
$$
\Bigl\lbrace 
\frac{\partial^{2} V}{\partial \varepsilon^{2}_{z}} \cdot \frac{\partial^{2} V}{\partial \varepsilon^{2}_{\theta}} - 
\Bigl( \frac{\partial^{2} V}{\partial \varepsilon_{z}\partial \varepsilon_{\theta}}\Bigr) 
\Bigr\rbrace 
4 \pi^{2} b^2 l^2 t = 0.
$$
После несложных преобразований получаем
\begin{multline}\label{S_eq11}
\bigl(0.05 - \alpha \beta - \alpha \beta^{-1}(\varepsilon_{z} + \nu \varepsilon_{\theta})^2\bigr) 
\bigl(0.05 - \alpha \beta - \alpha \beta^{-1}(\varepsilon_{\theta} + \nu \varepsilon_{z})^2\bigr)-
\\
 - \nu^2 \bigl( 0.05 - \alpha \beta - \alpha \beta^{-1}(\varepsilon_{z} + \nu \varepsilon_{\theta}) (\varepsilon_{\theta} + \nu \varepsilon_{z})\bigr)^2 = 0
,                                          
\end{multline}
где 
$ \alpha=0.25 (1-\nu)^{ - 1/2}$, 
$ \beta = (\varepsilon^{2}_{z} + 2\nu \varepsilon_{z} \varepsilon_{\theta} + \varepsilon^{2}_{\theta})^{1/2}$;  
$ \varepsilon_{z}$, $\varepsilon_{\theta} \ge 0$, $(\varepsilon_{z}, \varepsilon_{\theta}) \in \Omega $.

Для вычисления вырожденных критических точек функции $ W $ применим следующую численную процедуру. 




Построим на плоскости  $ \varepsilon_{z} \varepsilon_{\theta} $ квадрат $ \Omega_1 $  такой, что в нем $ \varepsilon_{z}$, $\varepsilon_{\theta} > 0$,     $(\varepsilon_{z}, \varepsilon_{\theta}) \in \Omega \subset  \Omega_1 $ (рис.~\ref{Sfig3}). 
Для определения стороны этого квадрата найдем на осях $ \varepsilon_{\theta}$ и $\varepsilon_{z} $ равные отрезки, внутри которых детерминант матрицы Гессе меняет знак. 
Получаем отрезки длиной $0.1$. 
Эту величину и берем за сторону квадрата $\Omega_1$. 
Отметим, что данный прием может быть использован при построении $M \times N$-мерного куба в общей теории. 



\begin{figure}[h!]
\centering
	\includegraphics[width=0.4\textwidth]{fig_stru3}
	\caption{Вид кривой, в точках которой гессиан равен нулю} \label{Sfig3} 
	\smallskip \footnotesize
		[Figure~\ref{Sfig3}. The curve where the Hessian is equal to zero at all points]
\end{figure}

Разобьем этот квадрат прямоугольной сеткой с~достаточно малым шагом (на рис. \ref{Sfig3} показана укрупненная сетка). В~каждом узле сетки по формуле \eqref{S_eq11} при  $ \nu = 0.3 $ вычисляем детерминант (гессиан) матрицы Гессе. 
Выделяем те из них, где детерминант мало отличен от нуля, что определяется заданной точностью вычислений.

Производя расчеты, получаем множество точек, расположенных на кривой {\sl 1} (рис. \ref{Sfig3}). 
Затем координаты точек кривой {\sl 1} подставляем в уравнения \eqref{S_eq9} и~выделяем те из них, которые удовлетворяют уравнениям с~заданной точностью. 
В~результате находим только одну точку (точка $A$, рис. \ref{Sfig3}). Точка $ A $ лежит на прямой \eqref{S_eq10} и~имеет координаты $ \varepsilon_{z}  = 0.018$, $\varepsilon_{\theta}  = 0.076 $ (с~точностью до округления). 
Это и~есть искомая вырожденная критическая точка потенциальной функции $ W $. 
Подставляя ее координаты в~уравнения равновесия, находим критическое внутреннее давление, при котором происходит потеря устойчивости системы (катастрофа, разрушение). 
Если  $ E = 2 \cdot 10^5 $~МПа, $ b = 100 $ мм,  $ t = 2 $ мм, то критическое (предельное) значение давления $ p = 18 $~МПа. Заметим, что величина предельного давления увеличивается пропорционально толщине стенки.


\pagebreak



\AdditionalContent{Конкурирующие интересы}{Мы не имеем конкурирующих интересов.} 
%\AdditionalContent{Конкурирующие интересы}{Конкурирующих интересов не имею.} 
\AdditionalContent{Авторский вклад и ответственность}{Все авторы принимали участие в разработке концепции статьи и в написании рукописи. Все авторы несут полную ответственность за представление окончательной рукописи в печать. Окончательная версия рукописи была одобрена всеми авторами.} 
%\AdditionalContent{Авторская ответственность}{Я несу полную ответственность за предоставление окончательной версии рукописи в печать. Окончательная версия рукописи мною одобрена.} 
%\AdditionalContent{Авторский вклад и ответственность}{Все  авторы  принимали  участие  в  разработке  концепции  \mbox{статьи}  и  в  написании  рукописи. Авторы несут  полную  ответственность  за  предоставление  окончательной рукописи в печать. Окончательная  версия рукописи была одобрена всеми авторами.} 
\AdditionalContent{Финансирование}{Исследование не имело финансирования.}
%\AdditionalContent{Финансирование}{Исследование выполнялось без финансирования.}
%\AdditionalContent{Финансирование}{Работа выполнена при поддержке РФФИ (проект № xx–xx–xxxxx_a).}
%\AdditionalContent{Благодарность}{Выразите свою благодарность.}
%\AdditionalContent{Благодарность}{Авторы благодарны рецензенту за тщательное прочтение статьи и~ценные предложения и комментарии.}

\begin{thebibliography}{10}

\RBibitem{bib1}
\by Седов~Л.~И.
\book Механика сплошной среды
\bookvol 1
\yr 1970
\publ Наука
\publaddr М.
\totalpages 492
%\misctransl
%\by Sedov ~L.~I.
%\book Mekhanika sploshnoy sredy \rm [Continuum Mechanics] vol. 1 
%\yr 1970
%\publ Nauka
%\publaddr Moscow
%\lang In Russian
%\totalpages 492

\RBibitem{bib2}
\by Стружанов~В.~В., Миронов~В.~И.
\book Деформационное разупрочнение материала в элементах конструкций
\yr 1995
\publ УрО РАН 
\publaddr Екатеринбург
\totalpages 190
%\misctransl
%\by Struzhanov ~V.~V., Mironov ~V.~I.
%\book Deformatsionnoye razuprochneniye materiala v elementakh konstruktsiy \rm [Deformational Weakening of the Material in Structural Elements]
%\yr 1995
%\publ UrO RAN
%\publaddr Ekaterinburg
%\lang In Russian
%\totalpages 190

\RBibitem{bib3}
\by Вильдеман~В.~Э., Чаусов~Н.~Г.
\paper Условия деформационного разупрочнения материала при растяжении образца специальной конфигурации
\jour Заводская лаборатория. Диагностика материалов
\yr 2007
\vol 73
\issue 10
\pages 55--59
\elib{https://elibrary.ru/item.asp?id=9907854}
%\misctransl
%\lang In Russian
%\by Vildeman ~V.~E., Chausov ~N.~G.
%\paper Usloviya deformatsionnogo razuprochneniya materiala pri rastyazhenii obraztsa spetsial'noy konfiguratsii \rm [Conditions of Deformational Weakening of the Material with the Specimen of a Special Configuration under Tension]
%\jour Zavodskaya laboratoriya. Diagnostika materialov.
%\yr 2007
%\vol 73
%\issue 10
%\pages 55--59

\RBibitem{bib4}
\by Вильдеман~В.~Э., Третьяков~М.~П.
\paper Испытания материалов с построением полных диаграмм деформирования
\jour Проблемы машиностроения и надежности машин
\yr 2013
\vol 
\issue 2
\pages 93--98
\elib{https://elibrary.ru/item.asp?id=19548776}
%\misctransl
%\lang In Russian
%\by Vildeman ~V.~E., Tretyakov ~M.~P.
%\paper Ispytaniya materialov s postroyeniyem polnykh diagramm deformirovaniya \rm [Material Tests and Development of Complete Deformation Diagrams]
%\jour Problemy mashinostroyeniya i nadezhnosti mashin
%\yr 2013
%\vol 
%\issue 2
%\pages 93--98

\RBibitem{bib5}
\by Стружанов~В.~В., Коркин~А.~В.
\paper Об устойчивости процесса растяжения одной стержневой системы с~разупрочняющимися элементами
\jour Вестник Уральского государственного университета путей сообщения
\yr 2016
\issue 3(31)
\pages 4--17
\elib{https://elibrary.ru/item.asp?id=27249016}
\crossref{10.20291/2079-0392-2016-3-4-17}
%\misctransl
%\lang In Russian
%\by Struzhanov ~V.~V., Korkin ~A.~V.
%\paper Ob ustoychivosti protsessa rastyazheniya odnoy sterzhnevoy sistemy s razuprochnyayushchimisya elementami \rm [On the Stability of a Single Rod System with Weakening Elements under Tension]
%\jour Vestnik Ural'skogo gos. un–ta putey soobshcheniya
%\yr 2016
%\vol 
%\issue 3(31)
%\pages 4--17

\RBibitem{bib6}
\by Андреева~Е.~А.
\paper Решение одномерных задач пластичности для разупрочняющегося материала
\jour Вестн. Сам. гос. техн. ун-та. Сер. Физ.-мат. науки
\yr 2008
\issue 2(17)
\pages 152--160
\mathnet{http://mi.mathnet.ru/vsgtu642}
\crossref{10.14498/vsgtu642}


%\RBibitem{bib7}
%\by Горбунов~С.~В., Радченко~В.~П.
%\paper Вариант решения краевой задачи неупругого деформирования разупрочняющейся среды и его экспериментальная проверка
%\jour Материалы Х Всероссийской конференции по механике деформируемого твёрдого тела (18–22 сентября 2017 г., Самара, Россия)
%\yr 2017
%\vol 1
%\issue
%\pages 179--181
%\misctransl
%\lang In Russian
%\by Gorbunov ~S.~V., Radchenko ~V.~P.
%\paper Variant resheniya krayevoy zadachi neuprugogo deformirovaniya razuprochnyayushcheysya sredy i yego eksperimental'naya proverka \rm [A Solution Case of the Boundary-Value Problem for Inelastic Deformation of a Weakening Medium and its Experimental Verification]
%\jour Materialy X  Vserossiyskoy konferentsii po mekhanike deformiruyemogo tvordogo tela.  
%\yr 2017
%\vol 1
%\issue
%\pages 179--181

\RBibitem{bib8}
\by Арсенин~В.~Я.
\book Методы математической физики и специальные функции
\yr 1974
\publ Наука
\publaddr М.
\totalpages 431
%\misctransl
%\by Arsenin ~V. ~Ya.
%\book Metody matematicheskoy fiziki i spetsial'nyye funktsii. \rm [Methods of Mathematical Physics and Special Functions]
%\yr 1974
%\publ Nauka
%\publaddr Moscow
%\lang In Russian
%\totalpages 431

\RBibitem{bib9}
\by Стружанов~В.~В., Коркин~А.~В.
\paper О решении нелинейных уравнений равновесия  одной растягиваемой стержневой системы с разупрочняющимися элементами методом простых итераций
\jour Вестник Уральского государственного университета путей сообщения
\yr 2017
\issue 2(34)
\pages 4--16
\crossref{10.20291/2079-0392-2017-2-4-16 }
\elib{https://elibrary.ru/item.asp?id=29409142}
%
%\misctransl
%\lang In Russian
%\by Struzhanov ~V.~V., Korkin ~A.~V.
%\paper O reshenii nelineynykh uravneniy ravnovesiya odnoy rastyagivayemoy sterzhnevoy sistemy s razuprochnyayushchimisya elementami metodom prostykh iteratsiy \rm [The solution of Nonlinear Equilibrium Equations for a Single Rod System with Weakening Elements under Tension Using the Simple Iterations Method]
%\jour Vestnik Ural'skogo gos. un-ta putey soobshcheniya
%\yr 2017
%\vol 
%\issue 2(34)
%\pages 4--16

\Bibitem{bib10}
\by Poston~T., Stewart~I.
\book Catastrophe Theory and Its Applications
\yr 1978
\publ Pitman
\serial Surveys and Reference Works in Mathematics
\vol 2
\zmath{0382.58006}
\publaddr London, San Francisco, Melbourne
\totalpages xviii+491 

\Bibitem{bib11}
\by Gilmore~R.
\book Catastrophe theory for scientists and engineers
\zmath{0497.58001}
\serial A Wiley-Interscience Publication
\publaddr New York 
\publ John Wiley \& Sons
\totalpages xvii+666 
\yr 1981

\Bibitem{bib12}
\by Pars~L.~A.
\book A treatise on analytical dynamics
\zmath{0125.12004}
\publaddr London
\publ Heinemann Educational Books 
\totalpages xxi+641 
\yr 1965
 
\RBibitem{bib13}
\by Ильюшин~А.~А.
\paper Об основах общей математической теории пластичности
\jour Вестник Московского университета. Серия 1. Математика. Механика
\yr 1961
\issue 3
\pages 31--36
\zmath{0100.20304}
%\misctransl
%\lang In Russian
%\by Ilyushin ~A.~A.
%\paper Ob osnovakh obshchey matematicheskoy teorii plastichnosti \rm [Fundamentals of the General Mathematical Theory of Plasticity]
%\jour Voprosy teorii plastichnosti. Izd-vo Akademii nauk SSSR
%\yr 1961
%\vol 
%\issue 
%\pages 3--29

\RBibitem{bib14}
\by Качанов~Л.~М.
\book Основы теории пластичности
\yr 1969
\publ Наука
\publaddr М.
\totalpages 420
%\misctransl
%\by Kachanov ~L.~M.
%\book Osnovy teorii plastichnosti \rm [Fundamentals of the Theory of Plasticity]
%\yr 1969
%\publ Nauka
%\publaddr Moscow
%\lang In Russian
%\totalpages 420

\RBibitem{bib15}
\by Стружанов~В.~В.
\paper Определение диаграммы деформирования материала с~падающей ветвью по диаграмме кручения цилиндрического образца
\jour Сиб. журн. индустр. матем.
\yr 2012
\vol 15
\issue 1
\pages 138--144
\mathnet{http://mi.mathnet.ru/sjim718}

\RBibitem{bib16}
\by Радченко~В.~П.
\book Введение в механику деформируемых систем
\yr 2009
\publ СамГТУ
\publaddr Самара
\totalpages 241
%\misctransl
%\by Radchenko ~V.~P.
%\book Vvedeniye v mekhaniku deformiruyemykh system \rm [Introduction to Mechanics of De-formation Systems]
%\yr 2009
%\publ Samarskiy gos. tekhn. un-t.
%\publaddr Samara
%\lang In Russian
%\totalpages 241

\Bibitem{bib17}
\by Southwell~R.~V.
\book An introduction to the theory of elasticity. For engineers and physicists
\zmath{62.1529.04}
\totalpages ix+510 
\publaddr Oxford
\publ Clarendon Press
\serial Oxford Engineering Science Series
\yr 1936


\end{thebibliography}


\newpage

\makeentitle

\vspace{-5mm}

\AdditionalContent{Competing interests}{We have no competing interests.} 
%\AdditionalContent{Competing interests}{I have ... (or) no competing interests.} 
%\AdditionalContent{Competing interests}{We have ... (or) no competing interests.} 
\AdditionalContent{Authors' contributions and responsibilities}{Each authors has participated in the article concept development  and in the manuscript writing. The authors are absolutely responsible for submitting the final manuscript in print. Each authors has approved the final version of manuscript.}
%\AdditionalContent{Author's Responsibilities}{I take full responsibility for submitting the final manuscript in print. I approved the final version of the manuscript.} 
%\AdditionalContent{Authors' contributions and responsibilities}{Each author has participated in the article concept development and in the manuscript writing. The authors are absolutely responsible for submitting the final manuscript in print. Each author has approved the final version of manuscript.} 

\AdditionalContent{Funding}{This research received no specific grant from any funding agency in the public, commercial, or not-for-profit sectors.}




\begin{thebibliography}{10}

\Bibitem{bib1:1}
\lang In Russian
\by Sedov~L.~I.
\book Mekhanika sploshnoi sredy \rm [Continuum Mechanics] 
\bookvol 1
\yr 1970
\publ Nauka
\publaddr Moscow
\totalpages 492

\Bibitem{bib2:1}
\lang In Russian
\by Struzhanov~V.~V., Mironov~V.~I.
\book Deformatsionnoe razuprochnenie materiala v elementakh konstruktsii \rm [Deformational Softening of Material in Structural Elements]
\yr 1995
\publ UrO RAN 
\publaddr Ekaterinburg
\totalpages 190

\Bibitem{bib3}
\lang In Russian 
\by Vil'deman~V.~E., Chausov~N.~G.
\paper Conditions of strain softening upon stretching of the specimen of special configuration
\jour Zavodskaia laboratoriia. Diagnostika materialov
\yr 2007
\vol 73
\issue 10
\pages 55--59


\Bibitem{bib4:en}
\paper Material testing  by plotting total deformation curves
\by Vil'deman~V.~E., Tretyakov~M.~P.
\jour J. Mach. Manuf. Reliab.
\yr  2013
\vol 42
\issue 2
\pages 166--170
\elib{https://elibrary.ru/item.asp?id=20435521}
\crossref{10.3103/S1052618813010159}



\Bibitem{bib5:1}
\lang In Russian
\by Struzhanov~V.~V., Korkin~A.~V.
\paper Regarding stretching process stability of one bar system with softening elements
\jour Vestnik Ural'skogo gosudarstvennogo universiteta putei soobshcheniia
\yr 2016
\issue 3(31)
\pages 4--17
\crossref{10.20291/2079-0392-2016-3-4-17}


\Bibitem{bib6}
\lang In Russian
\by Andreeva~E.~A.
\paper Solution of one-dimensional softening materials plasticity problems
\jour Vestn. Samar. Gos. Tekhn. Univ., Ser. Fiz.-Mat. Nauki \rm [J. Samara State Tech. Univ., Ser. Phys. Math. Sci.]
\yr 2008
\issue 2(17)
\pages 152--160
\crossref{10.14498/vsgtu642}



\Bibitem{bib8:1}
\lang In Russian
\by Arsenin~V.~Ya.
\book Metody matematicheskoi fiziki i spetsial'nye funktsii \rm [Methods of mathematical physics and special functions]
\yr 1974
\publ Nauka
\publaddr Moscow
\totalpages 431


\Bibitem{bib9:1}
\lang In Russian
\by Struzhanov~V.~V., Korkin~A.~V.
\paper On the solution of non-linear equations of trim of a single expandable frame structure with softening elements by the method of simple iterations
\jour Vestnik Ural'skogo gosudarstvennogo universiteta putei soobshcheniia
\yr 2017
\issue 2(34)
\pages 4--16
\crossref{10.20291/2079-0392-2017-2-4-16}


\Bibitem{bib10:1}
\by Poston~T., Stewart~I.
\book Catastrophe Theory and Its Applications
\yr 1978
\publ Pitman
\serial Surveys and Reference Works in Mathematics
\vol 2
\zmath{0382.58006}
\publaddr London, San Francisco, Melbourne
\totalpages xviii+491 

\Bibitem{bib11:1}
\by Gilmore~R.
\book Catastrophe theory for scientists and engineers
\zmath{0497.58001}
\serial A Wiley-Interscience Publication
\publaddr New York 
\publ John Wiley \& Sons
\totalpages xvii+666 
\yr 1981

\Bibitem{bib12:1}
\by Pars~L.~A.
\book A treatise on analytical dynamics
\zmath{0125.12004}
\publaddr London
\publ Heinemann Educational Books 
\totalpages xxi+641 
\yr 1965
 
\Bibitem{bib13:1}
\lang In Russian
\by Ilyushin~A.~A.
\paper On the foundations of the general mathematical theory of plasticity
\jour Vestn. Mosk. Univ., Ser.~I 
\yr 1961
\issue 3
\pages 31--36



\Bibitem{bib14:1}
\by Kachanov~L.~M.
\book Foundations of the theory of plasticity
\zmath{0231.73015}
\serial North-Holland Series in Applied Mathematics and Mechanics
\vol 12
\publaddr Amsterdam
\publ North-Holland Publ.
\totalpages xiii+482 
\yr 1971


\Bibitem{bib15:!}
\by Struzhanov~V.~V.
\paper The determination of the deformation diagram of a~material with a~falling branch using the torsion diagram of a~cylindrical sample
\jour Sib. Zh. Ind. Mat.
\yr 2012
\vol 15
\issue 1
\pages 138--144

\Bibitem{bib16:1}
\lang In Russian
\by Radchenko~V.~P.
\book Vvedenie v mekhaniku deformiruemykh sistem \rm [Introduction to the mechanics of deformable systems]
\yr 2009
\publ Samara State Technical Univ.
\publaddr Samara
\totalpages 241


\Bibitem{bib17:1}
\by Southwell~R.~V.
\book An introduction to the theory of elasticity. For engineers and physicists
\zmath{62.1529.04}
\totalpages ix+510 
\publaddr Oxford
\publ Clarendon Press
\serial Oxford Engineering Science Series
\yr 1936


\end{thebibliography}


%\end{document}
 \newpage

%\input{preamble}

%

\vsgtu{1639} %registration number

\FirstPage{774}

\LastPage{784}

%

%\usepackage{samgtu-name}

%

%\pdfcrop

%\draft

%

\ShortCommunication

%

%\begin{document}





$\;$

\vspace{1mm}



\hfill \qrcode[hyperlink,height=.8in]{https://doi.org/10.14498/vsgtu1639}

\vspace{-.8in}



\udc{517.955, 517.968.7}

\msc{26A33, 35K15, 35R11}



\rutitle{К проблеме единственности решения задачи Коши для~уравнения дробной диффузии с оператором Бесселя}

\rutitleshort{К проблеме единственности решения задачи Коши\dots}

\rutitlehack{К проблеме единственности решения \\ задачи Коши для~уравнения дробной диффузии с~оператором Бесселя}



\entitle{On the uniqueness of the solution of the Cauchy problem for the equation of fractional diffusion with~Bessel~operator}

\entitleshort{On the uniqueness of the solution of the Cauchy problem\dots}



\ruauthor{Ф.~Г.~Хуштова}{Х\,у\,ш\,т\,о\,в\,а~Ф.~Г.}

\enauthor{F.~G.~Khushtova}{K\,h\,u\,s\,h\,t\,o\,v\,a~F.~G.}



\ruaffil{\runiipma.}



\enaffil{\enniipma.}



\ruauthorsdetails{1}{% Количество авторов

\fullauthor{\rupid{53181}{Фатима Гидовна Хуштова}}

\CAuthor % Автор-корреспондент

\orcid{0000-0003-4088-3621} % ORCID ID автора 

%\regalia{без ученой степени} 

\position[n]{научный сотрудник} 

\dept{отдел дробного исчисления} 

\email{khushtova@yandex.ru}

}



\enauthorsdetails{1}{

\fullauthor{\enpid{53181}{Fatima G. Khushtova}}

\CAuthor % Сorresponding Author

\orcid{0000-0003-4088-3621} % ORCID ID

%\regalia{Cand. Phys. \& Math. Sci., Associate Professor} 

\position[n]{Researcher} 

\dept{Dept. of Fractional Calculus}

\email{khushtova@yandex.ru}

}



\rukeywords{уравнение дробной диффузии, оператор дробного дифференцирования, оператор Бесселя, задача Коши, единственность решения, условие Тихонова, функция Райта}

\enkeywords{fractional diffusion equation, fractional differentiation operator, Bessel operator, Cauchy problem, solution uniqueness, Tikhonov condition,	 Wright function}



\ruabstract{Рассматривается уравнение дробной диффузии с сингулярным оператором Бесселя, действующим по пространственной переменной, и оператором дробного дифференцирования Римана--Лиувилля, действующим по временной переменной. 

Когда порядок дробной производной равен единице, а особенность у оператора Бесселя отсутствует, рассматриваемое уравнение совпадает с~классическим уравнением теплопроводности. 

Ранее для уравнения дробной диффузии с оператором Бесселя было построено решение задачи Коши и доказана теорема единственности решения в классе функций экспоненциального роста. 



Построен пример, показывающий, что увеличение показателя степени в условии, гарантирующем единственность решения задачи Коши, влечет за собой неединственность решения. 

С помощью известных свойств функции Райта получены оценки для построенной функции.

Показывается, что она, будучи не равной тождественно нулю, удовлетворяет однородному уравнению и однородному условию Коши. 

}





\enabstract{In this paper, we consider fractional diffusion equation involving the Bessel operator acting with respect to a spatial variable and the Riemann-Liouville fractional differentiation operator acting with respect to a time variable. When the order of the fractional derivative is unity, and the singularity of the Bessel operator is absent, this equation coincides with the classical heat equation. Earlier, a solution of the Cauchy problem has been considered for the considered equation and a uniqueness theorem has been proved for a class of functions satisfying the analog of the Tikhonov condition. 



In this paper, we have constructed an example to show that the exponent (power) at the condition of the uniqueness of the solution to the Cauchy problem cannot be raised under. Its increase leads to a non-uniqueness of the solution. Using the well-known properties of the Wright function, we have obtained estimates for constructed function, which satisfies the homogeneous equation and the zero Cauchy condition.

}



\newdate{khusha}{28}{08}{2018}

\date{\displaydate{khusha}}

\newdate{khushb}{25}{10}{2018}

\revisiondate{\displaydate{khushb}}

\newdate{khushc}{12}{11}{2018}

\accepteddate{\displaydate{khushc}}



\newdate{khushe}{28}{11}{2018}

\renewcommand{\JournalDataOnline}{\displaydate{khushe}}



%\ForPeerReview

\makerutitle



\newpage



\Section[n]{Введение}

В области $\Omega=\{(x,y): -\infty<x<\infty,\ 0<y<T \}$ рассмотрим уравнение

\begin{equation}\label{GL3_1}

B_xu(x,y)-D^{\alpha}_{0y}u(x,y)=0, 

\end{equation}

где $D_{0y}^{\alpha}$ "--- оператор дробного дифференцирования в смысле  

Римана"--~Лиувилля  порядка $\alpha,$ определяемый равенствами~\cite{Drob, Pskhu}   

$$D_{0y}^{\alpha} \, g(y)= \frac{dg}{dy},$$ если $\alpha=1$, и

$$ 

D_{0y}^{\alpha}\,g(y)=\frac{1}{\Gamma(1-\alpha)}\frac{d}{dy}

\int _{0}^{y}\frac{g(t)}{(y-t)^{\alpha}}dt,

$$ 

если $0<\alpha<1;$

$$

B_x=x^{-b} \frac{d}{dx}\left(x^b \frac{d}{dx}\right)=\frac{d^2}{d x^2}+\frac{b}{x} \frac{d}{d x}

$$

"--- оператор Бесселя, $|b|<1.$



Уравнение \eqref{GL3_1} при $\alpha=1,$ то есть уравнение

\begin{equation}\label{UrAlpha-1}

u_{xx}(x,y)+\frac{b}{x}u_{x}(x,y)-u_{y}(x,y)=0,

\end{equation}

совпадает с %~названным И.~А.~Киприяновым 

$B$-параболическим уравнением.

Краевые задачи в ограниченной и~неограниченной областях для этого уравнения при различных значениях параметра $b$ рассматривали многие авторы. 

Например, в работе V.~Alexia\-des~\cite{Alexiades} для него решены первая, вторая и~третья краевые задачи в~области с~подвижной границей, в~работе D.~Calton\cite{Calton} для него  исследована задача Коши.

 

 Уравнение \eqref{UrAlpha-1} заменой $z=x^2/(4n)$ сводится к~уравнению

\begin{equation} \label{UrKep2}

u_{zz}(z,y)+\frac{m+1}{z}\,u_{z}(z,y)-\frac{n}{z} u_{y}(z,y)=0,\quad m=(b-1)/2,

\end{equation}

для которого S.~K\c{e}pi\'nski  \cite{Kepinski_2} исследовал краевую задачу в~полуполосе $x>0,$ $0<y<y_0.$

Уравнение \eqref{UrAlpha-1} было также объектом исследования  монографий С.~А.~Терсенова~\cite{Tersenov} и~М.~И.~Матийчука~\cite{Matiychuk}.



Дифференциальные уравнения, содержащие оператор Бесселя, наиболее подробно исследованы в~работах И.~А.~Киприянова и~его учеников \cite{Kipriyanov-Monographiya,  Kipriyanov}.

Отметим здесь, что очень важные результаты по применению операторов преобразования в~теории дифференциальных уравнений с~особенностями в~коэффициентах, в~частности с~операторами Бесселя, получены В.~В.~Катраховым и~С.~М.~Ситником~\cite{KatrakhovSitnik, Sitnik}.





В случае, когда $b=0,$ $0<\alpha<2,$ уравнение \eqref{GL3_1} совпадает с диффузионно"=волновым уравнением. 

Различные краевые задачи для него, а также для многомерных его обобщений подробно исследованы в работах многих авторов.  Приведем некоторые из них. 



А.~Н.~Кочубей \cite{Kochubei_2} исследовал задачу Коши для уравнения диффузии дробного порядка с~регуляризованной дробной производной (производной Капуто) и~эллиптическим оператором с~непрерывными ограниченными вещественными коэффициентами

$$

L=\sum\limits_{i,j=1}^{n}a_{i,j}(x)\frac{\partial^2}{\partial x_i \partial x_j}+

\sum\limits_{j=1}^{n}b_{j}(x)\frac{\partial}{\partial x_j}+c(x) u.

$$

В случае, когда $L=\Delta$ "--- оператор Лапласа, фундаментальное решение построено в терминах $H$-функции Фокса. 

Исследованию задачи Коши для диффузионного и диффузионно"=волнового уравнений дробного порядка с~производной Капуто  посвящены также  работы А.~Н.~Кочубея и С.~Д.~Эйдельмана \cite{Kochubei_1, Kochubei_3, Kochubei_4, Kochubei_5}.



А.~В.~Псху~\cite{Pskhu2} построил фундаментальное решение и исследовал задачу Коши  для многомерного диффузионно"=волнового уравнения с~оператором дробного дифференцирования Джрбашяна"--~Нерсесяна, который  включает в~себя как частный случай операторы дробного дифференцирования Римана"--~Лиувилля и Капуто.



В~работах  \cite{Pskhu-SibirJournal, Pskhu3} также построено фундаментальное решение и исследована задача Коши для уравнения дробной диффузии соответственно с~оператором дискретно распределенного дифференцирования и оператором Джрбашяна"--~Нерсесяна, действующим по многим переменным времени.

В работе \cite{Pskhu4} решена первая краевая задача в~нецилиндрической области для диффузионно"=волнового уравнения с~оператором Джрбашяна"--~Нерсесяна. 

Доказана теорема существования и~единственности исследуемой задачи, конструктивно построено ее решение. 



А.~А.~Ворошилов и А.~А.~Килбас \cite{VorKilbas} методом интегральных преобразований построили решение задачи Коши для уравнения

\begin{equation} \label{MnogPerem}

D^{\alpha}_{0t}\,u(x,t)=\lambda^2 \Delta_x u(x,t), \quad  x\in{\mathbb R}^m,\quad  t>0, 

\end{equation}

где $\Delta_x=\sum_{j=1}^{m}\partial^2/\partial x_{j}^{2},$ $n-1<\alpha <n,$ $n\in\mathbb{N}.$ 

При $0<\alpha\leqslant 1$ и $1<\alpha<2$ решения  выписаны в терминах $H$-функции Фокса.

В работе~\cite{VorKilbas2} было также  построено решение задачи Коши в~случае, когда в уравнении

\eqref{MnogPerem} вместо оператора Римана"--~Лиувилля содержится оператор Капуто. 



Задача Коши для уравнения \eqref{MnogPerem} при $\lambda=1,$ 

$0<\alpha\leqslant 1$  была ранее рассмотрена С.~Х.~Геккиевой~\cite{Gekk}, а~в~работе \cite{Gekkieva1} ею исследована первая краевая задача в~первом квадранте для уравнения \eqref{MnogPerem} при $m=\lambda=1,$ 

$0<\alpha\leqslant 1.$  







Уравнение диффузии дробного порядка и некоторые его обобщения рассматривали 

F.~Mainardi, G.~Pagnini, Yu.~Luchko, P.~Paradisi~\cite{Mainardi-1, Mainardi-3, Mainardi-2, Pagnini_1,   Pagnini_Paradisi}  и~многие другие. 

Более подробную библиографию по данному вопросу можно найти в~монографиях \cite{Pskhu, MathaiSaxenaHaubold, Podlubny}.

 

%Чтобы подготовить Вашу статью для рецензента, т.~е. обезличить её, воспользуйтесь командой \verb"\ForPeerReview" перед командой 

%\verb"\makerutitle".



\smallskip





\Section{Постановка задачи Коши и теорема единственности решения}

Пусть $\Omega^{+}=\Omega\cap\{x>0 \},$ $\Omega^{-}=\Omega\cap\{x<0\},$ 

$\bar{\Omega}$ "--- замыкание области $\Omega.$ 



\smallskip



{\sc\small Определение.}

{\it Регулярным  решением} уравнения \eqref{GL3_1} в области $\Omega$ 

будем называть функцию

$u=u(x,y),$ удовлетворяющую уравнению \eqref{GL3_1} в области

$\Omega^{+}\cup\Omega^{-},$ и такую, что $y^{\,1-\alpha}u\in{C(\bar{\Omega})},$ 

$|x|^b u_{x}\in{C(\Omega)},$ $u_{xx},$ $D_{\,0y}^{\,\alpha} u\in{C(\Omega^{+}\cup\Omega^{-})}.$



\smallskip



{\sc\small Задача Коши.}

{\it Найти регулярное в области $\Omega$ решение

уравнения} \eqref{GL3_1}, {\it удовлетворяющее условию

\begin{equation}\label{UslSMCoshi}

\lim\limits_{y \rightarrow 0} D_{0y}^{\alpha-1}u(x,y)=\varphi(x), \quad  -\infty<x<\infty,

\end{equation}

где $\varphi(x)$ "--- заданная функция}.



\smallskip



В работе \cite{KhushtovaVestn2} доказана следующая теорема единственности.



\smallskip



\begin{newthm}{Теорема}  

Существует не более одного регулярного решения задачи Коши {\rm \eqref{GL3_1}, \eqref{UslSMCoshi} } в классе функций$,$ удовлетворяющих для некоторой положительной постоянной $k$ условию

\begin{equation} \label{AnalogTikhonovaInfty}

\lim\limits_{|x|\to \infty}y^{1-\alpha}\,u(x,y)\exp{\big(-k\,|x|^{\,\frac{2}{2-\alpha}} \big)}=0.

\end{equation}  

\end{newthm}



\smallskip





Отметим, что условие \eqref{AnalogTikhonovaInfty} имеет место и в случае задачи Коши для уравнения дробной диффузии.  

Теоремы единственности решения задачи Коши в~классе ограниченных функций  и~в~классе функций экспоненциального роста для многомерного уравнения дробной диффузии с~оператором Капуто были доказаны в работе А.~Н.~Кочубея \cite{Kochubei_2},

теорема единственности решения задачи Коши в классе функций экспоненциального роста для многомерного диффузионно"=волнового уравнения с~оператором Джрбашяна"--~Нерсесяна  "--- в~работе А.~В.~Псху~\cite{Pskhu2}. 

В этих же работах построены примеры, показывающие, что увеличение показателя $2/(2-\alpha)$ ведет к~утрате единственности решения соответствующих задач. 

При $\alpha=1$ условие \eqref{AnalogTikhonovaInfty} совпадает с~хорошо известным условием Тихонова 

$$

\lim\limits_{|x|\to \infty}u(x,y)\exp{\big(-k\,|x|^{2} \big)}=0.

$$

Как известно, пример неединственности решения задачи Коши для уравнения теплопроводности был впервые построен А.~Н.~Тихоновым в~основополагающей работе~\cite{tikhonovUsl-1} (см.~также \cite{tikhonovUsl-2}).



В данной работе эти идеи развиваются применительно к~задаче \eqref{GL3_1}, \eqref{UslSMCoshi}.



\smallskip



\Section{Неулучшаемость показателя степени в условии единственности решения задачи Коши}

Рассмотрим функцию

\begin{equation} \label{T(z,zeta)}  

T_{\mu}(z,\zeta)=\sum\limits_{n=0}^{\infty}\frac{z^{\,\beta+n}}{\,n!\,\Gamma(1+\beta+n)}\,\phi\left(-\rho,\mu-\alpha\beta-\alpha n; -\zeta\right),

\end{equation}    

где  $\phi\,(-\rho,\delta;z)$ "--- {\it функция Райта}, определяемая степенным рядом \cite{Pskhu, GorenfloLuchkoMainardi}

$$

\phi\,(-\rho,\delta;z)=\sum\limits_{n=0}^{\infty}\frac{z^{n}}{\,n!\,\Gamma(\delta-\rho\, n)}.

$$



Далее будем считать, что $\rho\in(0;1).$ 

Для любых $\nu$, $\delta\in \mathbb{R}$ справедливы формулы  \cite[с.~26, 44]{Pskhu}

\begin{gather}\label{diff_x}

\frac{d}{dz}\phi(-\rho,\delta;-z)=-\phi(-\rho,\delta-\rho;-z),

\\

\label{diff_Right}

D_{ay}^{\nu} \,y^{\delta-1}\phi(-\rho,\delta;-c\,y^{-\rho})=y^{\delta-\nu-1}\phi(-\rho,\delta-\nu;-c\,y^{-\rho}).

\end{gather}



\smallskip



\begin{lemma}[1]Если $\delta\geqslant 1,$ то для любого положительного $z$ справедливо нера\-вен\-ство~{\rm \cite[с.~47]{Pskhu} }

\begin{equation}\label{Right_delta>1}

\phi\,(-\rho,\delta;-z)\leqslant \frac{1}{\Gamma\left( \delta \right)}

\exp\left[-(1-\rho)\,\rho^{\frac{\rho}{1-\rho}}\, z^{\frac{1}{1-\rho}} \right].

\end{equation}

\end{lemma}



\smallskip



\begin{lemma}[2]Если $\delta<1,$ то для любых положительных $x$ и $y,$ 

$a\in(0;x),$ $\omega\in(1/2; \omega_0),$ $\omega_0=\min\{1,1/(2\rho)\},$

справедливы неравенства~{\rm \cite[с.~49]{Pskhu}}

\begin{gather}\label{Right-0}

\Bigl|\,y^{\,\delta-1}\phi\Bigl(-\rho,\delta;-\frac{x}{y^{\rho}}\Bigr)\Bigr|

\leqslant 

\frac{\Gamma\left( (1-\delta)/\rho\right)}{\,\rho\, \pi \left[aC_{\rho}(\rho,\omega) \right]^{(1-\delta)/\rho}}\,

\phi\Bigl(-\rho,1;-\frac{x-a}{y^{\rho}}\Bigr),

\\

\label{Right-2}

\Bigl|y^{\,\delta-1}\phi\Bigl(-\rho,\delta;-\frac{x}{y^{\rho}}\Bigr)\Bigr|

\leqslant \frac{\Gamma\left( 1-\delta\right)}{\pi \left[yC_{\rho}(1,\omega) \right]^{1-\delta}},

\end{gather}

где 

$C_{\rho}(\rho,\omega)=\frac1{\rho\pi}{\cos \rho\omega\pi},$ $C_{\rho}(1,\omega)=-\frac1{\pi}{\cos \omega\pi}.

$

\end{lemma}



\smallskip



Как следует из \eqref{Right-2}, ряд \eqref{T(z,zeta)} сходится абсолютно для любых  $z\in\mathbb{C},$ $\zeta>0,$ $\rho, \alpha, \beta\in(0;1),$ $\mu\in\mathbb{R}.$



\smallskip



Докажем следующую лемму.



\smallskip

\begin{lemma}[3]Если $0<\alpha<\rho<1,$ $\mu\geqslant 0,$ то при любых $x\in\mathbb{R},$ $y>0$ справедливо неравенство

\begin{equation} \label{T_muOtsenka}  

\Bigl| y^{\mu-1} T_{\mu}\Bigl(\frac{x^2}{4y^{\alpha}},\frac{1}{y^{\rho}}\Bigr)\Bigr|

\leqslant 

\frac{C

|x|^{\frac{2[1-\rho(1-\beta)]}{2\rho-\alpha}} y^{\mu}}{\Gamma(\mu+1)}\exp\left[ k_1  |x|^{\frac{2\rho}{2\rho-\alpha}}-k_2   y^{-\frac{\rho}{1-\rho}}\right],

\end{equation}    

где $C,$ $k_1,$ $k_2$ "--- положительные постоянные$,$ зависящие только от 

$\alpha,$ $\beta$ и $\rho.$

\end{lemma}



\smallskip



\begin{proof}

Из \eqref{Right-0} следует оценка

$$

\Bigl| y^{-\alpha\beta-\alpha n-1} \phi\Bigl(-\rho,-\alpha\beta-\alpha n; -\frac{1}{y^{\rho}}\Bigr)\Bigr|

\leqslant

 \frac{\Gamma\left( 1/\rho+\varepsilon\beta+\varepsilon n\right)}{ \rho \pi\left(aC_{\rho} \right)^{1/\rho+\varepsilon\beta+\varepsilon n}}

\phi\Bigl(-\rho,1; -\frac{1-a}{y^{\rho}}\Bigr),

$$

где $\varepsilon=\alpha/\rho,$ $a\in(0;1),$ $C_{\rho}=C_{\rho}(\rho,\omega).$



С помощью последней оценки получим

\begin{multline}

\Bigl| y^{\,\mu-1} T_{\mu}\Bigl(\frac{x^2}{4y^{\alpha}},\frac{1}{y^{\rho}}\Bigr)\Bigr|=

\\

=\left|

 D_{0y}^{-\mu}

  \sum\limits_{n=0}^{\infty}\frac{|x|^{2\beta+2n}}{ 2^{2\beta+2n}n! \Gamma(1+\beta+n)}

y^{-\alpha\beta-\alpha n-1}\phi \Bigl(-\rho,-\alpha\beta-\alpha n; -\frac{1}{y^{\rho}}\Bigr)

 \right|\leqslant

\\

 \label{WrightRyadZ}  

\leqslant

\frac{ |x|^{2\beta}y^{\mu}\phi\bigl(-\rho,\mu+1; -\frac{1-a}{y^{\rho}}\bigr)}{2^{2\beta}\rho \pi \left(aC_{\rho} \right)^{1/\rho+\varepsilon\beta}}

\sum\limits_{n=0}^{\infty}

\frac{\Gamma\left(1/\rho+\varepsilon\beta+\varepsilon n\right)}{n!

\Gamma(1+\beta+n)}

\Bigl( \frac{x^2}{4\left(aC_{\rho} \right)^{\varepsilon}}\Bigr)^{n}.

\end{multline}



Степенной ряд, стоящий справа от неравенства \eqref{WrightRyadZ}, представляет собой 

обобщенную функцию Райта

$$

{_1\Psi_{1}}\left[ \,z \,\bigg|

\begin{array}{ll}

 \big(1/\rho+\varepsilon\beta, \,\varepsilon\,\big)\\

\big(\,1+\beta, 1\,\big)\\

\end{array}

\right]=

\sum\limits_{n=0}^{\infty}

\frac{\Gamma\left(1/\rho+\varepsilon\beta+\varepsilon n\right)}{n!

\,\Gamma(1+\beta+n)} z^{n}, \; \; z=\frac{x^{\,2}}{4(aC_{\rho})^{\varepsilon}}.

$$

Из ее асимптотических свойств при $z\to \infty$ следует~\cite{Wright}

$$

{_1\Psi_{1}}\left[ \,z\,\bigg|

\begin{array}{ll}

 \big(1/\rho+\varepsilon\beta, \,\varepsilon\,\big)\\

\big(\,1+\beta, 1\,\big)\\

\end{array}

\right]=Z^{\,1/\rho-\beta(1-\varepsilon)-1}\,e^{Z}

\left[ \sum\limits_{m=0}^{M-1}A_m Z^{-m}+O\left(Z^{-M}\right) \right],

$$

где

$Z=(2-\varepsilon)\,\varepsilon^{\,\frac{\varepsilon}{2-\varepsilon}}z^{\,\frac{1}{2-\varepsilon}},$  $\varepsilon<2,$ а коэффициенты $A_m,$ $m=0, 1, 2, ...,$ зависят только от $\rho,$ $\varepsilon$ и $\beta.$



Учитывая, что $a$ "--- произвольное число из интервала $(0;1),$ из последней оценки и~соотношений \eqref{WrightRyadZ}, \eqref{Right_delta>1} получаем \eqref{T_muOtsenka}. \hfill \end{proof}



\smallskip



Покажем, что функция $y^{\mu-1} T_{\mu}\bigl(\frac{x^2}{4y^{\alpha}},\frac{1}{y^{\rho}}\bigr)$ является решением уравнения \eqref{GL3_1}. 

Из \eqref{T(z,zeta)}, \eqref{diff_x}  и \eqref{diff_Right} имеем

\begin{multline}

B_x T_{\mu}\Bigl(\frac{x^2}{4y^{\alpha}},\frac{1}{y^{\rho}}\Bigr)=

\\

=y^{-\alpha}\sum\limits_{n=1}^{\infty}\frac{1}{(n-1)! \Gamma(\beta+n)}

\Bigl(\frac{x^2}{4y^{\alpha}}\Bigr)^{\beta+n-1}

\phi\Bigl(-\rho,\mu-\alpha\beta-\alpha n; -\frac{1}{y^{\rho}}\Bigr)=

\\

=y^{-\alpha}\sum\limits_{k=0}^{\infty}\frac{1}{k! \Gamma(1+\beta+k)}

\Bigl(\frac{x^2}{4y^{\alpha}}\Bigr)^{\beta+k}

\phi\Bigl(-\rho,\mu-\alpha-\alpha\beta-\alpha k; -\frac{1}{y^{\rho}}\Bigr)=

\\ 

=y^{-\alpha} T_{\mu-\alpha}\Bigl(\frac{x^2}{\,4y^{\alpha}},\frac{1}{y^{\rho}}\Bigr),

\end{multline}   

\begin{equation}

\label{DT(z,zeta)}  

D_{0y}^{\nu} y^{\mu-1} T_{\mu}\Bigl(\frac{x^2}{4y^{\alpha}},\frac{1}{y^{\rho}}\Bigr)=

y^{\mu-\nu-1} T_{\mu-\nu}\Bigl(\frac{x^2}{4y^{\alpha}},\frac{1}{y^{\rho}}\Bigr).

\end{equation}



Из последних двух равенств следует

\begin{equation} \label{BDyT}  

\left(B_x-D_{0y}^{\alpha}\right)y^{\mu-1} T_{\mu}\Bigl(\frac{x^2}{4y^{\alpha}}, \frac{1}{y^{\rho}}\Bigr)=0.

\end{equation}   





Из равенства \eqref{DT(z,zeta)} и оценки \eqref{T_muOtsenka} также имеем

\begin{equation} \label{D(alpha-1)T}  

\lim\limits_{y \rightarrow 0} D_{0y}^{\alpha-1} y^{\mu-1} T_{\mu}\Bigl(\frac{x^2}{4y^{\alpha}}, \frac{1}{y^{\rho}}\Bigr)=0.

\end{equation}   



Таким образом, из соотношений  \eqref{T_muOtsenka}, \eqref{BDyT} и \eqref{D(alpha-1)T} следует, что нетривиальная функция 

$$

u(x,y)=y^{\mu-1} T_{\mu}\Bigl(\frac{x^2}{4y^{\alpha}},\frac{1}{y^{\rho}}\Bigr)

$$

является решением однородного уравнения \eqref{GL3_1}, удовлетворяет однородному условию \eqref{UslSMCoshi} $(\varphi(x)\equiv 0)$ и для любого положительного сколь угодно малого $\delta,$ 

некоторых положительных постоянных $k$ и $C,$ а также параметра 

$\rho,$  выбранного из условия 

$\rho=\frac{\alpha}{2}\frac{2+\delta(2-\alpha)}{\alpha+\delta(2-\alpha)},$

имеет место оценка

$$

\left| u(x,y) \right|\leqslant C \exp\left(k |x|^{\delta+\frac{2}{2-\alpha}} \right).

$$



Последняя оценка показывает, что увеличение показателя  $2/(2-\alpha)$ в~условии \eqref{AnalogTikhonovaInfty} ведет к~неединственности решения задачи  \eqref{GL3_1},  \eqref{UslSMCoshi} и~в~этом смысле условие  \eqref{AnalogTikhonovaInfty} является необходимым.





\smallskip









\Section[N]{Заключение}

В данной работе показывается, что показатель степени в~условии, обеспечивающем единственность решения задачи Коши для уравнения дробной диффузии с оператором Бесселя, является предельным, и его увеличение приводит к нарушению единственности. 



\smallskip





\AdditionalContent{Конкурирующие интересы}{Конкурирующих интересов не имею.} 



\AdditionalContent{Авторская ответственность}{Я несу полную ответственность за предоставление окончательной версии рукописи в печать. Окончательная версия рукописи мною одобрена.} 





\AdditionalContent{Финансирование}{Исследование выполнялось без финансирования.}



\smallskip





\begin{thebibliography}{99}



\RBibitem{Drob}

\by Нахушев~А.~М.

\book Дробное исчисление и его применение

\yr 2003

\publ Физматлит

\publaddr М.

\totalpages 271

%\sortdate 1/1/2003

\zmath{https://zbmath.org/?q=an:1066.26005}

%\misctransl

%\by Nakhushev~A.~M.

%\book Drobnoe ischislenie i ego primenenie \rm [Fractional calculus and its applications]

%\yr 2003

%\publ Fizmatlit

%\publaddr Moscow

%\lang In Russian

%\totalpages 271

%%\sortdate 1/1/2003



\RBibitem{Pskhu}

\by Псху~А.~В.

\book Уравнения в частных производных дробного порядка

\yr 2005

\publ Наука

\publaddr М.

\totalpages 199

%\sortdate 1/1/2005

\zmath{https://zbmath.org/?q=an:1193.35245}

%\misctransl

%\by Pskhu~A.~V.

%\book Uravneniia v chastnykh proizvodnykh drobnogo poriadka \rm [Partial differential equations of fractional order]

%\yr 2005

%\publ Nauka

%\publaddr Moscow

%\lang In Russian

%\totalpages 199

%%\sortdate 1/1/2005



\Bibitem{Alexiades}

\by Alexiades~V.

\paper Generalized axially symmetric heat potentials and singular parabolic initial boundary value problems

\jour Arch. Rational Mech. Anal.

\yr 1982

\vol 79

\issue 4

\pages 325--350

%\citnumscopus 3

%\sortdate 1/1/1982

\crossref{10.1007/BF00250797}

\mathscinet{http://www.ams.org/mathscinet-getitem?mr=656798}

\zmath{https://zbmath.org/?q=an:0511.35050}

\scopus{http://www.scopus.com/record/display.url?origin=inward&eid=2-s2.0-0020402046}



\Bibitem{Calton}

\by Calton~D.

\paper Cauchy's problem for a singular parabolic partial differential equation

\jour J.~Diff. Equations

\yr 1970

\vol 8

\issue 2

\pages 250--257

%\citnumscopus 8

%\sortdate 1/1/1970

\crossref{10.1016/0022-0396(70)90004-5}

\mathscinet{http://www.ams.org/mathscinet-getitem?mr=271547}

\scopus{http://www.scopus.com/record/display.url?origin=inward&eid=2-s2.0-49849107701}



\Bibitem{Kepinski_2}

\by K\c{e}pi\'nski~S.

\paper \"Uber die Differentialgleichung $\frac{\partial^2 z }{\partial x^2}+\frac{m+1}{x}\frac{\partial z }{\partial x} -\frac{n}{x}\frac{\partial z }{\partial t}=0$

\jour Math. Ann.

\yr 1905

\vol 61

\issue 3

\pages 397--405

\lang In German

%\sortdate 1/1/1905

\mathscinet{http://www.ams.org/mathscinet-getitem?mr=1511354}

\zmath{https://zbmath.org/?q=an:36.0416.01}



\RBibitem{Tersenov}

\by Терсенов~С.~А.

\book Параболические уравнения с~меняющимся направлением времени

\yr 1985

\publ Наука

\publaddr М.

\totalpages 105

%\sortdate 1/1/1985

%\misctransl

%\by Tersenov~S.~A.

%\book Parabolicheskie uravneniia s~meniaiushchimsia napravleniem vremeni \rm [Parabolic Equations with a Changing Time Direction]

%\yr 1985

%\publ Nauka

%\publaddr Moscow

%\totalpages 105

%%\sortdate 1/1/1985



\RBibitem{Matiychuk}

\by Матiйчук~M.~I.

\book Параболiчнi сингулярнi крайовi задачi

\yr 1999

\publ Iн-т математики НАН Укра\"\iни

\publaddr Ки\"\iв

\totalpages 176

%%\sortdate 1/1/1999

%\misctransl

%\by Matiichuk~M.~I.

%\book Parabolichni singuliarni kraiovi zadachi \rm [Parabolic Singular Boundary-Value Problems]

%\yr 1999

%\publ Institute of Mathematics of the Ukrainian National Academy of Sciences

%\publaddr Kiev

%\lang In Ukrainian

%\totalpages 176

%%\sortdate 1/1/1999



\RBibitem{Kipriyanov-Monographiya}

\by Киприянов~И.~А.

\book Сингулярные эллиптические краевые задачи

\yr 1997

\publ Наука

\publaddr М.

\totalpages 208

%\sortdate 1/1/1997

%\misctransl

%\by Kipriyanov~I.~A.

%\book Singuliarnye ellipticheskie kraevye zadachi \rm [Singular Elliptic Boundary-Value Problems]

%\yr 1997

%\publ Nauka

%\publaddr Moscow

%\lang In Russian

%\totalpages 208

%%\sortdate 1/1/1997



\RBibitem{Kipriyanov}

\by Киприянов~И.~А., Катрахов~В.~В., Ляпин~В.~М.

\paper О краевых задачах в областях общего вида для сингулярных параболических систем уравнений

\jour Докл. АН СССР

\yr 1976

\vol 230

\issue 6

\pages 1271--1274

%\sortdate 1/1/1976

\mathnet{http://mi.mathnet.ru/rus/dan40691}

\mathscinet{http://www.ams.org/mathscinet-getitem?mr=0425366}

\zmath{https://zbmath.org/?q=an:0354.35056}

%\misctransl

%\by Kipriyanov~I.~A., Katrakhov~V.~V., Lyapin~V.~M.

%\paper On boundary value problems in domains of general type for singular parabolic systems of equations

%\jour Sov. Math., Dokl.

%\yr 1976

%\vol 17

%\pages 1461--1464

%%\sortdate 1/1/1976

%\mathscinet{http://www.ams.org/mathscinet-getitem?mr=0425366}

%\zmath{https://zbmath.org/?q=an:0354.35056}



\RBibitem{Sitnik}

\by Ситник~С.~М.

\yr 2016

\publaddr Воронеж

\thesis Применение операторов преобразования Бушмана--Эрдейи и их обобщений в~теории дифференциальных уравнений с~особенностями в коэффициентах

\thesisinfo Дис.~\dots\ д-ра физ.-мат. наук: 01.01.02

\totalpages 307

%\sortdate 1/1/2016

%\misctransl

%\by Sitnik~S.~M.

%\yr 2016

%\publaddr Voronezh

%\thesis Application of the Bushman--Erd\'elyi transformation operators and their generalizations in the theory of differential equations with singularities in the coefficients

%\thesisinfo Dr. Phys. Math. Sci. Thesis

%\lang In Russian

%\totalpages 307

%%\sortdate 1/1/2016



\RBibitem{KatrakhovSitnik}

\by Катрахов~В.~В., Ситник~С.~М.

\paper Метод операторов преобразования и краевые задачи для сингулярных эллиптических уравнений

\inbook Сингулярные дифференциальные уравнения

\serial СМФН

\yr 2018

\vol 64

\issue 2

\pages 211--426

\publ Российский университет дружбы народов

\publaddr М.

%\sortdate 1/1/2018

\mathnet{http://mi.mathnet.ru/rus/cmfd355}

\crossref{10.22363/2413-3639-2018-64-2-211-426}

%\misctransl

%\by Katrakhov~V.~V., Sitnik~S.~M.

%\paper The transmutation method and boundary-value problems for singular elliptic equations

%\inbook Singular differential equations

%\serial CMFD

%\yr 2018

%\vol 64

%\issue 2

%\pages 211--426

%\publ Peoples' Friendship University of Russia

%\publaddr Moscow

%\lang In Russian

%%\sortdate 1/1/2018

%\crossref{https://doi.org/10.22363/2413-3639-2018-64-2-211-426}



\RBibitem{Kochubei_2}

\by Кочубей~А.~Н.

\paper Диффузия дробного порядка

\jour Дифференц. уравнения

\yr 1990

\vol 26

\issue 4

\pages 660--670

%\sortdate 1/1/1990

\mathnet{http://mi.mathnet.ru/rus/de7144}

\mathscinet{http://www.ams.org/mathscinet-getitem?mr=1061448}

\zmath{https://zbmath.org/?q=an:0712.35049}

%\transl

%\by Kochubei~A.~N.

%\paper Diffusion of fractional order

%\jour Differ. Equ.

%\yr 1990

%\vol 26

%\issue 4

%\pages 485--492

%%\sortdate 1/1/1990

%\mathscinet{http://www.ams.org/mathscinet-getitem?mr=1061448}

%\zmath{https://zbmath.org/?q=an:0729.35064}



\RBibitem{Kochubei_1}

\by Кочубей~А.~Н.

\paper Задача Коши для эволюционных уравнений дробного порядка

\jour Дифференц. уравнения

\yr 1989

\vol 25

\issue 8

\pages 1359--1368

%\sortdate 1/1/1989

\mathnet{http://mi.mathnet.ru/rus/de6938}

\mathscinet{http://www.ams.org/mathscinet-getitem?mr=1014153}

%\transl

%\by Kochubei~A.~N.

%\paper The Cauchy problem for evolution equations of fractional order

%\jour Differ. Equ.

%\yr 1989

%\vol 25

%\issue 8

%\pages 967--974

%%\sortdate 1/1/1989

%\mathscinet{http://www.ams.org/mathscinet-getitem?mr=1014153}



\RBibitem{Kochubei_3}

\by Кочубей~А.~Н., Эйдельман~С.~Д.

\paper Задача Коши для эволюционных уравнений дробного порядка

\jour Докл. РАН

\yr 2004

\vol 394

\issue 2

\pages 159--161

%\sortdate 1/1/2004

\mathnet{http://mi.mathnet.ru/rus/dan1624}

\mathscinet{http://www.ams.org/mathscinet-getitem?mr=2089744}

\zmath{https://zbmath.org/?q=an:1129.35427}

%\misctransl

%\by Kochubeii~A.~N., \'Eidel'man~S.~D.

%\paper The Cauchy problem for evolution equations of fractional order

%\jour Dokl. Akad. Nauk

%\yr 2004

%\vol 394

%\issue 2

%\pages 159--161

%\lang In Russian

%%\sortdate 1/1/2004



\Bibitem{Kochubei_4}

\by Kochubei~A.~N.

\paper Asymptotic Properties of Solutions of the Fractional Diffusion-Wave Equation

\jour Fract. Calc. Appl. Anal.

\yr 2014

\vol 17

\issue 3

\pages 881--896

%\citnumscopus 9

%\sortdate 1/1/2014

\crossref{10.2478/s13540-014-0203-3}

\mathscinet{http://www.ams.org/mathscinet-getitem?mr=3260311}

\zmath{https://zbmath.org/?q=an:1323.35197}

\elib{http://elibrary.ru/item.asp?id=24541812}

\scopus{http://www.scopus.com/record/display.url?origin=inward&eid=2-s2.0-84904312827}



\Bibitem{Kochubei_5}

\by Kochubei~A.~N.

\paper Cauchy problem for fractional diffusion-wave equations with variable coefficients

\jour Applicable Analysis

\yr 2014

\vol 93

\issue 10

\pages 2211--2242

%\citnumscopus 7

%\sortdate 1/1/2014

\crossref{10.1080/00036811.2013.875162}

\mathscinet{http://www.ams.org/mathscinet-getitem?mr=3240385}

\zmath{https://zbmath.org/?q=an:1297.35274}

\elib{http://elibrary.ru/item.asp?id=24642244}

\scopus{http://www.scopus.com/record/display.url?origin=inward&eid=2-s2.0-84905387477}



\RBibitem{Pskhu2}

\by Псху~А.~В.

\paper Фундаментальное решение диффузионно-волнового уравнения дробного порядка

\jour Изв. РАН. Сер. матем.

\yr 2009

\vol 73

\issue 2

\pages 141--182

%\sortdate 1/1/2009

\mathnet{http://mi.mathnet.ru/rus/im2429}

\crossref{10.4213/im2429}

\mathscinet{http://www.ams.org/mathscinet-getitem?mr=2532450}

\zmath{https://zbmath.org/?q=an:1172.26001}

\adsnasa{http://adsabs.harvard.edu/cgi-bin/bib_query?2009IzMat..73..351P}

\elib{http://elibrary.ru/item.asp?id=20425204}

%\transl

%\by Pskhu~A.~V.

%\paper The fundamental solution of a diffusion-wave equation of fractional order

%\jour Izv. Math.

%\yr 2009

%\vol 73

%\issue 2

%\pages 351--392

%%\sortdate 1/1/2009

%\crossref{https://doi.org/10.1070/IM2009v073n02ABEH002450}

%\mathscinet{http://www.ams.org/mathscinet-getitem?mr=2532450}

%\zmath{https://zbmath.org/?q=an:1172.26001}

%\isi{http://gateway.isiknowledge.com/gateway/Gateway.cgi?GWVersion=2&SrcApp=PARTNER_APP&SrcAuth=LinksAMR&DestLinkType=FullRecord&DestApp=ALL_WOS&KeyUT=000266177900006}

%\elib{http://elibrary.ru/item.asp?id=13597908}

%\scopus{http://www.scopus.com/record/display.url?origin=inward&eid=2-s2.0-65349154473}



\RBibitem{Pskhu-SibirJournal}

\by Псху~А.~В.

\paper Уравнение дробной диффузии с оператором дискретно распределенного дифференцирования

\jour Сиб. электрон. матем. изв.

\yr 2016

\vol 13

\pages 1078--1098

%\sortdate 1/1/2016

\mathnet{http://mi.mathnet.ru/rus/semr736}

\crossref{10.17377/semi.2016.13.086}

\mathscinet{http://www.ams.org/mathscinet-getitem?mr=3580054}

\zmath{https://zbmath.org/?q=an:1370.35268}

\elib{http://elibrary.ru/item.asp?id=28127183}

%\misctransl

%\by Pskhu~A.~V.

%\paper Fractional diffusion equation with discretely distributed differentiation operator

%\jour Sib. \`Elektron. Mat. Izv.

%\yr 2016

%\vol 13

%\pages 1078--1098

%\lang In Russian

%%\sortdate 1/1/2016

%\crossref{https://doi.org/10.17377/semi.2016.13.086}



\Bibitem{Pskhu3}

\by Pskhu~A.~V.

\paper Multi-time fractional diffusion equation

\jour Eur. Phys. J.~Spec. Top.

\yr 2013

\vol 222

\issue 8

\pages 1939--1950

%\citnumscopus 2

%\sortdate 1/1/2013

\crossref{10.1140/epjst/e2013-01975-y}

\scopus{http://www.scopus.com/record/display.url?origin=inward&eid=2-s2.0-84885034281}



\RBibitem{Pskhu4}

\by Псху~А.~В.

\paper Первая краевая задача для дробного диффузионно-волнового уравнения в~нецилиндрической области

\jour Изв. РАН. Сер. матем.

\yr 2017

\vol 81

\issue 6

\pages 158--179

%\sortdate 1/1/2017

\mathnet{http://mi.mathnet.ru/rus/im8520}

\crossref{10.4213/im8520}

\mathscinet{http://www.ams.org/mathscinet-getitem?mr=3733319}

\zmath{https://zbmath.org/?q=an:06844352}

\adsnasa{http://adsabs.harvard.edu/cgi-bin/bib_query?2017IzMat..81.1212P}

\elib{http://elibrary.ru/item.asp?id=30737847}

%\transl

%\by Pskhu~A.~V.

%\paper The first boundary-value problem for a fractional diffusion-wave equation in a non-cylindrical domain

%\jour Izv. Math.

%\yr 2017

%\vol 81

%\issue 6

%\pages 1212--1233

%%\sortdate 1/1/2017

%\crossref{https://doi.org/10.1070/IM8520}

%\mathscinet{http://www.ams.org/mathscinet-getitem?mr=3733319}

%\zmath{https://zbmath.org/?q=an:06844352}

%\isi{http://gateway.isiknowledge.com/gateway/Gateway.cgi?GWVersion=2&SrcApp=PARTNER_APP&SrcAuth=LinksAMR&DestLinkType=FullRecord&DestApp=ALL_WOS&KeyUT=000418891300007}

%\scopus{http://www.scopus.com/record/display.url?origin=inward&eid=2-s2.0-85040983678}



\RBibitem{VorKilbas}

\by Ворошилов~А.~А., Килбас~А.~А.

\paper Задача типа Коши для диффузионно-волнового уравнения с частной производной Римана--Лиувилля

\jour Докл. РАН

\yr 2006

\vol 406

\issue 1

\pages 12--16

%\sortdate 1/1/2006

\mathnet{http://mi.mathnet.ru/rus/dan1045}

\mathscinet{http://www.ams.org/mathscinet-getitem?mr=2256848}

\zmath{https://zbmath.org/?q=an:1155.35396}

\elib{http://elibrary.ru/item.asp?id=9176502}

%\misctransl

%\by Voroshilov~A.~A., Kilbas~A.~A.

%\paper A Cauchy-type problem for a diffusion-wave equation with a Riemann-Liouville partial derivative

%\jour Dokl. Akad. Nauk

%\yr 2006

%\vol 406

%\issue 1

%\pages 12--16

%\lang In Russian

%%\sortdate 1/1/2006



\RBibitem{VorKilbas2}

\by Ворошилов~А.~А., Килбас~А.~А.

\paper Задача Коши для диффузионно-волнового уравнения с~частной производной Капуто

\jour Дифференц. уравнения

\yr 2006

\vol 42

\issue 5

\pages 599--609

%\sortdate 1/1/2006

\mathnet{http://mi.mathnet.ru/rus/de11489}

\mathscinet{http://www.ams.org/mathscinet-getitem?mr=2292158}

\zmath{https://zbmath.org/?q=an:1123.35302}

\elib{http://elibrary.ru/item.asp?id=9245260}

%\transl

%\by Voroshilov~A.~A., Kilbas~A.~A.

%\paper The Cauchy problem for the diffusion-wave equation with the Caputo partial derivative

%\jour Differ. Equ.

%\yr 2006

%\vol 42

%\issue 5

%\pages 638--649

%%\citnumscopus 11

%%\sortdate 1/1/2006

%\crossref{https://doi.org/10.1134/S0012266106050041}

%\mathscinet{http://www.ams.org/mathscinet-getitem?mr=2292158}

%\zmath{https://zbmath.org/?q=an:1123.35302}

%\elib{http://elibrary.ru/item.asp?id=27786841}

%\scopus{http://www.scopus.com/record/display.url?origin=inward&eid=2-s2.0-33745304762}



\RBibitem{Gekk}

\by Геккиева~С.~Х.

\paper Задача Коши для обобщенного уравнения переноса с дробной по времени производной

\jour Доклады Адыгской (Черкесской) Международной академии наук

\yr 2000

\vol 5

\issue 1

\pages 16--19

%\sortdate 1/1/2000

%\misctransl

%\by Gekkieva~S.~Kh.

%\paper The Cauchy problem for the generalized transport equation with time-fractional derivative

%\jour Dokl. Adyg. (Cherkess.) Mezdunar. Akad. Nauk

%\yr 2000

%\vol 5

%\issue 1

%\pages 16--19

%\lang In Russian

%%\sortdate 1/1/2000



\RBibitem{Gekkieva1}

\by Геккиева~С.~Х.

\paper Краевая задача для обобщенного уравнения переноса с дробной производной в полубесконечной области

\jour Изв. КБНЦ РАН

\yr 2002

\issue 1(8)

\pages 6--8

%\sortdate 1/1/2002

%\misctransl

%\by Gekkieva~S.~Kh.

%\paper A boundary value problem for the generalized transfer equation with a fractional derivative in a semi-infinite domain

%\jour Izv. KBNTs RAN

%\yr 2002

%\issue 1(8)

%\pages 6--8

%\lang In Russian

%%\sortdate 1/1/2002



\Bibitem{Mainardi-1}

\by Mainardi~F.

\paper The fundamental solutions for the fractional diffusion-wave equation

\jour Appl. Math. Letters

\yr 1996

\vol 6

\issue 1

\pages 23--28

%\citnumscopus 343

%\sortdate 1/1/1996

\crossref{10.1016/0893-9659(96)00089-4}

\mathscinet{http://www.ams.org/mathscinet-getitem?mr=1419811}

\zmath{https://zbmath.org/?q=an:0879.35036}

\scopus{http://www.scopus.com/record/display.url?origin=inward&eid=2-s2.0-30244460855}



\Bibitem{Mainardi-3}

\by Mainardi~F.

\paper The time fractional diffusion-wave equation

\jour Radiophys. Quantum Electron.

\yr 1995

\vol 38

\issue 1-2

\pages 13--24

%\citnumscopus 39

%\sortdate 1/1/1995

\crossref{10.1007/BF01051854}

\mathscinet{http://www.ams.org/mathscinet-getitem?mr=1427165}

\scopus{http://www.scopus.com/record/display.url?origin=inward&eid=2-s2.0-0029427754}



\Bibitem{Mainardi-2}

\by Mainardi~F., Luchko~Yu., Pagnini~G.

\paper The fundamental solution of the space-time fractional diffusion equation

\jour Fract. Calc. Appl. Anal.

\yr 2001

\vol 4

\issue 2

\pages 153--192

\arxiv \href{https://arxiv.org/abs/cond-mat/0702419}{cond-mat/0702419} [cond-mat.stat-mech]

%\sortdate 1/1/2001

\mathscinet{http://www.ams.org/mathscinet-getitem?mr=1829592}

\zmath{https://zbmath.org/?q=an:1054.35156}



\Bibitem{Pagnini_1}

\by Pagnini~G.

\paper The M-Wright function as a generalization of the Gaussian density for fractional diffusion processes

\jour Fract. Calc. Appl. Anal.

\yr 2013

\vol 16

\issue 2

\pages 436--453

%\citnumscopus 22

%\sortdate 1/1/2013

\crossref{10.2478/s13540-013-0027-6}

\mathscinet{http://www.ams.org/mathscinet-getitem?mr=3033698}

\zmath{https://zbmath.org/?q=an:1312.33061}

\elib{http://elibrary.ru/item.asp?id=28682912}

\scopus{http://www.scopus.com/record/display.url?origin=inward&eid=2-s2.0-84879620364}



\Bibitem{Pagnini_Paradisi}

\by Pagnini~G., Paradisi~P.

\paper A stochastic solution with Gaussian stationary increments of the symmetric space-time fractional diffusion equation

\jour Fract. Calc. Appl. Anal.

\yr 2016

\vol 19

\issue 2

\pages 408--440

%\citnumscopus 9

%\sortdate 1/1/2016

\crossref{10.1515/fca-2016-0022}

\mathscinet{http://www.ams.org/mathscinet-getitem?mr=3513003}

\zmath{https://zbmath.org/?q=an:1341.60073}

\elib{http://elibrary.ru/item.asp?id=27104131}

\scopus{http://www.scopus.com/record/display.url?origin=inward&eid=2-s2.0-84969800118}



\Bibitem{Podlubny}

\by Podlubny~I.

\book Fractional differential equations. An introduction to fractional derivatives, fractional differential equations, to methods of their solution and some of their applications

\serial Mathematics in Science and Engineering

\yr 1999

\vol 198

\publ Academic Press

\publaddr San Diego, CA

\totalpages xxiv+340

%\sortdate 1/1/1999

\mathscinet{http://www.ams.org/mathscinet-getitem?mr=1658022}

\zmath{https://zbmath.org/?q=an:0924.34008}



\Bibitem{MathaiSaxenaHaubold}

\by Mathai~A.~M., Saxena~R.~K., Haubold~H.~J.

\book The $H$-function. Theory and Applications

\yr 2010

\publ Springer

\publaddr Dordrecht

\totalpages xiv+268~pp \nofrills

%\citnumscopus 411

%\sortdate 1/1/2010

\crossref{10.1007/978-1-4419-0916-9}

\mathscinet{http://www.ams.org/mathscinet-getitem?mr=2562766}

\zmath{https://zbmath.org/?q=an:1181.33001}

\scopus{http://www.scopus.com/record/display.url?origin=inward&eid=2-s2.0-84891408876}



\RBibitem{KhushtovaVestn2}

\by Хуштова~Ф.~Г.

\paper Задача Коши для уравнения параболического типа~с~оператором Бесселя и~частной производной~Римана--Лиувилля

\jour Вестн. Сам. гос. техн. ун-та. Сер. Физ.-мат. науки

\yr 2016

\vol 20

\issue 1

\pages 74--84

%\sortdate 1/1/2016

\mathnet{http://mi.mathnet.ru/rus/vsgtu1455}

\crossref{10.14498/vsgtu1455}

\zmath{https://zbmath.org/?q=an:06964474}

\elib{http://elibrary.ru/item.asp?id=26898117}

%\misctransl

%\by Khushtova~F.~G.

%\paper Cauchy problem for a parabolic equation with Bessel operator and Riemann--Liouville partial derivative

%\jour Vestn. Samar. Gos. Tekhn. Univ., Ser. Fiz.-Mat. Nauki \rm [J. Samara State Tech. Univ., Ser. Phys. Math. Sci.]

%\yr 2016

%\vol 20

%\issue 1

%\pages 74--84

%%\sortdate 1/1/2016



\Bibitem{GorenfloLuchkoMainardi}

\by Gorenflo~R., Luchko~Y., Mainardi~F.

\paper Analytical properties and applications of the Wright function

\jour Fract. Calc. Appl. Anal.

\yr 1999

\vol 2

\issue 4

\pages 383--414

\arxiv \href{https://arxiv.org/abs/math-ph/0701069}{math-ph/0701069}

%\sortdate 1/1/1999

\mathscinet{http://www.ams.org/mathscinet-getitem?mr=1752379}

\zmath{https://zbmath.org/?q=an:1027.33006}



\Bibitem{tikhonovUsl-1}

\by Tychonoff~A.

\paper Th\'eor\`emes d'unicit\'e pour l'\'equation de la chaleur

\jour Mat. Sb.

\yr 1935

\vol 42

\issue 2

\pages 199--216

\lang In French

%\sortdate 1/1/1935

\mathnet{http://mi.mathnet.ru/rus/sm6410}

\zmath{https://zbmath.org/?q=an:0012.35501|61.1203.05}



\RBibitem{tikhonovUsl-2}

\by Тихонов~А.~Н.

\paper Теоремы единственности для уравнения теплопроводности

\inbook Собрание научных трудов: в десяти томах

\yr 2009

\pages 371--382

\publ Наука

\publaddr М.

\volinfo Т.~II. Математика. Ч.~2. Вычислительная математика 1956--1979. Математическая физика 1933--1948. Ред.-сост. Т.~А.~Сушкевич, А.~В.~Гулин

%\sortdate 1/1/2009

\zmath{https://zbmath.org/?q=an:1200.01055}

%\misctransl

%\by Tikhonov~A.~N.

%\paper Uniqueness theorems for the heat equation

%\inbook Collection of scientific works. In ten volumes

%\yr 2009

%\pages 371--382

%\publ Nauka

%\publaddr Moscow

%\lang In Russian and French

%\volinfo Vol. 2: Mathematics. Part 2: Numerical mathematics 1956–1979. Mathematical physics 1933–1948. Edited by T.~A.~Sushkevich and A.~V.~Gulin

%%\sortdate 1/1/2009

%\zmath{https://zbmath.org/?q=an:1200.01055}



\Bibitem{Wright}

\by Wright~E.~M.

\paper The asymptotic expansion of the generalized hypergeometric function

\jour Proc. Lond. Math. Soc., II. Ser.

\yr 1940

\vol 46

\issue 1

\pages 389--408

%\citnumscopus 136

%\sortdate 1/1/1940

\crossref{10.1112/plms/s2-46.1.389}

\mathscinet{http://www.ams.org/mathscinet-getitem?mr=3876}

\zmath{https://zbmath.org/?q=an:0025.40402|66.1243.01}

\scopus{http://www.scopus.com/record/display.url?origin=inward&eid=2-s2.0-84963078619}



\end{thebibliography}



	





\newpage



\makeentitle



\vspace{-6mm}



%\AdditionalContent{Competing interests}{Specify what conflicts of interest are available.} 

\AdditionalContent{Competing interests}{I have no competing interests.} 

%\AdditionalContent{Competing interests}{We have ... (or) no competing interests.} 

%\AdditionalContent{Authors' contributions and responsibilities}{What is the author's responsibility for each author?}

\AdditionalContent{Author's Responsibilities}{I take full responsibility for submitting the final ma\-nuscript in print. I approved the final version of the manuscript.} 

%\AdditionalContent{Authors' contributions and responsibilities}{Each author has participated in the article concept development and in the manuscript writing. The authors are absolutely responsible for submitting the final manuscript in print. Each author has approved the final version of manuscript.} 



%\AdditionalContent{Funding}{Indicate the source(s) of your funding.}



\AdditionalContent{Funding}{This research received no specific grant from any funding agency in the public,

commercial, or not-for-profit sectors.}



\newpage



\begin{thebibliography}{10}



\Bibitem{Drob:2}

\lang In Russian

\by Nakhushev~A.~M.

\book Drobnoe ischislenie i ego primenenie \rm [Fractional calculus and its applications]

\publaddr Moscow

\publ Fizmatlit 

\totalpages 271

\yr 2003





\Bibitem{Pskhu:w}

\lang In Russian

\by Pskhu~A.~V.

\book Uravneniia v chastnykh proizvodnykh drobnogo poriadka \rm [Partial differential equations of fractional order]

\publaddr Moscow

\publ Nauka 

\totalpages 199 

\yr 2005







\Bibitem{Alexiades:w}

\by Alexiades~V.

\paper Generalized axially symmetric heat potentials and singular parabolic initial boundary value problems

\jour  Arch. Rational Mech. Anal. 

\yr 1982

\vol 79

\issue 4

\pages 325--350

\crossref{10.1007/BF00250797}







\Bibitem{Calton:w}

\by Calton~D.

\paper Cauchy's problem for a singular parabolic partial differential equation

\jour J.~Diff. Equations

\yr 1970

\vol 8

\issue 2

\pages 250--257

\crossref{10.1016/0022-0396(70)90004-5}







\Bibitem{Kepinski_2:w}

\by K\c{e}pi\'nski~S.

\paper \"Uber die Differentialgleichung $\frac{\partial^2 z }{\partial x^2}+\frac{m+1}{x}\frac{\partial z }{\partial x} -\frac{n}{x}\frac{\partial z }{\partial t}=0$

\jour   Math. Ann.

\yr 1905

\vol 61

\issue 3

\pages 397--405

\zmath{36.0416.01}

\lang In German



\Bibitem{Tersenov:en}

\lang In Russian

\by Tersenov~S.~A.

\book Parabolicheskie uravneniia s~meniaiushchimsia napravleniem vremeni \rm [Parabolic Equations with a Changing Time Direction]

\yr 1985

\publ Nauka

\publaddr Moscow

\totalpages 105



\Bibitem{Matiychuk:en}

\by Matiichuk~M.~I.

\book  Parabolichni singuliarni kraiovi zadachi \rm [Parabolic Singular Boundary-Value Problems]

\yr 1999

\publ Institute of Mathematics of the Ukrainian National  Academy of Sciences

\publaddr Kiev

\totalpages 176

\lang In Ukrainian





\Bibitem{Kipriyanov-Monographiya:en}

\lang In Russian

\by Kipriyanov~I.~A.

\book  Singuliarnye ellipticheskie kraevye zadachi \rm [Singular Elliptic Boundary-Value Problems]

\yr 1997

\publ Nauka

\publaddr Moscow

\totalpages 208



\Bibitem{Kipriyanov:en}

\by Kipriyanov~I.~A., Katrakhov~V.~V., Lyapin~V.~M.

\paper On boundary value problems in domains of general type for singular parabolic systems of equations

\jour Sov. Math., Dokl. 

\vol 17

\yr 1976

\pages 1461--1464 

\mathscinet{http://www.ams.org/mathscinet-getitem?mr=0425366}

\zmath{https://zbmath.org/?q=an:0354.35056}





\Bibitem{Sitnik:en}

\lang In Russian

\by Sitnik~S.~M. 

\thesis Application of the Bushman--Erd\'elyi transformation operators and their generalizations in the theory of differential equations with singularities in the coefficients

\thesisinfo Dr. Phys. Math. Sci. Thesis

\yr 2016

\publaddr Voronezh

\totalpages 307







\Bibitem{KatrakhovSitnik:en}

\lang In Russian

\by Katrakhov~V.~V., Sitnik~S.~M.

\paper The transmutation method and boundary-value problems for singular elliptic equations

\inbook Singular differential equations

\serial CMFD

\yr 2018

\vol 64

\issue 2

\pages 211--426

\publ Peoples' Friendship University of Russia

\publaddr Moscow

\crossref{10.22363/2413-3639-2018-64-2-211-426}







\Bibitem{Kochubei_2:en}

\by Kochubei~A.~N.

\paper Diffusion of fractional order

\jour Differ. Equ.

\yr 1990

\vol 26

\issue 4

\pages 485--492







\Bibitem{Kochubei_1:w}

\by Kochubei~A.~N.

\paper The Cauchy problem for evolution equations of fractional order

\jour Differ. Equ.

\yr 1989

\vol 25

\issue 8

\pages 967--974



\Bibitem{Kochubei_3:e}

\lang In Russian

\by Kochubeii~A.~N., \'Eidel'man~S.~D.

\paper The Cauchy problem for evolution equations of fractional order

\jour Dokl. Akad. Nauk

\yr 2004

\vol 394

\issue 2

\pages 159--161







\Bibitem{Kochubei_4:e}

\by Kochubei~A.~N.

\paper Asymptotic Properties of Solutions of the Fractional Diffusion-Wave Equation

\jour Fract. Calc. Appl. Anal.

\yr 2014

\vol 17

\issue 3

\pages 881--896

\crossref{10.2478/s13540-014-0203-3}





\Bibitem{Kochubei_5:e}

\by Kochubei~A.~N.

\paper Cauchy problem for fractional diffusion-wave equations with variable coefficients

\jour  Applicable Analysis

\yr 2014

\vol 93

\issue 10

\pages 2211--2242

\crossref{10.1080/00036811.2013.875162}







\Bibitem{Pskhu2:en}

\paper The fundamental solution of a diffusion-wave equation of fractional order

\by Pskhu~A.~V.

\jour Izv. Math.

\yr 2009

\vol 73

\issue 2

\pages 351--392

\crossref{10.1070/IM2009v073n02ABEH002450}

\isi{http://gateway.isiknowledge.com/gateway/Gateway.cgi?GWVersion=2&SrcApp=PARTNER_APP&SrcAuth=LinksAMR&DestLinkType=FullRecord&DestApp=ALL_WOS&KeyUT=000266177900006}

\elib{http://elibrary.ru/item.asp?id=13597908}

\scopus{http://www.scopus.com/record/display.url?origin=inward&eid=2-s2.0-65349154473}





\Bibitem{Pskhu-SibirJournal:en}

\lang In Russian

\by Pskhu~A.~V.

\paper Fractional diffusion equation with discretely distributed differentiation operator

\jour Sib. \`Elektron. Mat. Izv.

\yr 2016

\vol 13

\pages 1078--1098

\crossref{10.17377/semi.2016.13.086}







\Bibitem{Pskhu3:en}

\by Pskhu~A.~V.

\paper Multi-time fractional diffusion equation

\jour Eur. Phys. J.~Spec. Top.

\yr 2013

\vol 222

\issue 8

\pages 1939--1950

\crossref{10.1140/epjst/e2013-01975-y}





\Bibitem{Pskhu4:en}

\paper The first boundary-value problem for a fractional diffusion-wave equation in a non-cylindrical domain

\by Pskhu~A.~V.

\jour Izv. Math.

\yr 2017

\vol 81

\issue 6

\pages 1212--1233

\crossref{10.1070/IM8520}

\isi{http://gateway.isiknowledge.com/gateway/Gateway.cgi?GWVersion=2&SrcApp=PARTNER_APP&SrcAuth=LinksAMR&DestLinkType=FullRecord&DestApp=ALL_WOS&KeyUT=000418891300007}

\scopus{http://www.scopus.com/record/display.url?origin=inward&eid=2-s2.0-85040983678}







\Bibitem{VorKilbas:a} 

\lang In Russian

\by Voroshilov~A.~A., Kilbas~A.~A.

\paper A Cauchy-type problem for a diffusion-wave equation with a

Riemann-Liouville partial derivative

\jour Dokl. Akad. Nauk

\yr 2006

\vol 406

\issue 1

\pages 12--16



\Bibitem{VorKilbas2:en} 

\by Voroshilov~A.~A., Kilbas~A.~A.

\paper The Cauchy problem for the diffusion-wave equation with the Caputo partial derivative

\mathscinet{http://www.ams.org/mathscinet-getitem?mr=2292158}

\jour Differ. Equ.

\yr 2006

\vol 42

\issue 5

\pages 638--649

\crossref{10.1134/S0012266106050041}



\Bibitem{Gekk:en} 

\lang In Russian

\by Gekkieva~S.~Kh.

\paper The Cauchy problem for the generalized transport equation with time-fractional derivative

\jour Dokl. Adyg. (Cherkess.) Mezdunar. Akad. Nauk

\yr 2000

\vol 5

\issue 1

\pages 16--19







\Bibitem{Gekkieva1:en}

\lang In Russian

\by Gekkieva~S.~Kh.

\paper A boundary value problem for the generalized transfer equation with a fractional derivative in a semi-infinite domain

\jour Izv. KBNTs RAN

\yr 2002

\issue 1(8)

\pages 6--8











  





\Bibitem{Mainardi-1:en}

\by Mainardi~F.

\paper The fundamental solutions for the fractional diffusion-wave equation

\jour Appl.  Math. Letters

\yr 1996

\vol 6

\issue 1

\pages 23--28

\crossref{10.1016/0893-9659(96)00089-4}



\Bibitem{Mainardi-3}

\by Mainardi~F.

\paper  The time fractional diffusion-wave equation

\jour Radiophys. Quantum Electron.

\yr 1995

\vol 38

\issue 1-2

\pages 13--24

\crossref{10.1007/BF01051854}



\Bibitem{Mainardi-2:en}

\by Mainardi~F., Luchko~Yu., Pagnini~G.

\paper The fundamental solution of the space-time fractional diffusion equation

\jour Fract. Calc. Appl. Anal.

\yr 2001

\vol 4

\issue 2

\pages 153--192

\arxiv  	\href{https://arxiv.org/abs/cond-mat/0702419}{cond-mat/0702419} [cond-mat.stat-mech]















\Bibitem{Pagnini_1:en}

\by Pagnini~G.

\paper The M-Wright function as a generalization of the Gaussian density for fractional diffusion processes

\jour Fract. Calc. Appl. Anal.

\yr 2013

\vol 16

\issue 2

\pages 436--453

\crossref{10.2478/s13540-013-0027-6}





\Bibitem{Pagnini_Paradisi:en}

\by Pagnini~G., Paradisi~P.

\paper A stochastic solution with Gaussian stationary increments of the symmetric space-time fractional diffusion equation

\jour Fract. Calc. Appl. Anal.

\yr 2016

\vol 19

\issue 2

\pages 408--440

\crossref{10.1515/fca-2016-0022}





\Bibitem{Podlubny:en}

\by Podlubny~I.

\book Fractional differential equations. An introduction to fractional derivatives, fractional differential equations, to methods of their solution and some of their applications

\yr 1999

\publ Academic Press

\publaddr  San Diego, CA

\zmath{0924.34008}

\serial Mathematics in Science and Engineering

\vol 198

\totalpages xxiv+340





\Bibitem{MathaiSaxenaHaubold:en}

\by Mathai~A.~M., Saxena~R.~K., Haubold~H.~J.

\book The $H$-function. Theory and Applications

\yr 2010

\publ Springer

\publaddr  Dordrecht

\totalpages xiv+268~pp \nofrills

\crossref{10.1007/978-1-4419-0916-9}

\zmath{1181.33001}









\Bibitem{KhushtovaVestn2}

\by Khushtova~F.~G.

\paper Cauchy problem for a parabolic equation with Bessel operator and Riemann--Liouville partial derivative

\jour Vestn. Samar. Gos. Tekhn. Univ., Ser. Fiz.-Mat. Nauki \rm [J. Samara State Tech. Univ., Ser. Phys. Math. Sci.]

\yr 2016

\vol 20

\issue 1

\pages 74--84

\crossref{10.14498/vsgtu1455}





\Bibitem{GorenfloLuchkoMainardi:en}

\by Gorenflo~R., Luchko~Y., Mainardi~F.

\paper Analytical properties and applications of the Wright  function

\jour Fract. Calc. Appl. Anal.

\yr 1999

\vol 2

\issue 4

\pages 383--414

\arxiv \href{https://arxiv.org/abs/math-ph/0701069}{math-ph/0701069}







\Bibitem{tikhonovUsl-1:en}

\by Tychonoff~A.

\paper Th\'eor\`emes d'unicit\'e pour l'\'equation de la chaleur

\jour Mat. Sb.

\yr 1935

\vol 42

\issue 2

\pages 199--216

\lang In French

\mathnet{http://mi.mathnet.ru/msb6410}

\zmath{https://zbmath.org/?q=an:0012.35501|61.1203.05}









\Bibitem{tikhonovUsl-2:en}

\by Tikhonov~A.~N.

\paper Uniqueness theorems for the heat equation

\inbook Collection of scientific works. In ten volumes

\volinfo Vol. 2: Mathematics. Part 2: Numerical mathematics 1956–1979. Mathematical physics 1933–1948. Edited by T.~A.~Sushkevich and  A.~V.~Gulin

\yr 2009

\publ Nauka

\publaddr Moscow

\pages 371--382

\zmath{1200.01055}

\lang In Russian and French





\Bibitem{Wright:en}

\by Wright~E.~M.

\paper The asymptotic expansion of the generalized hypergeometric function

\jour Proc. Lond. Math. Soc., II. Ser. 

\yr 1940

\vol 46

\issue 1

\pages 389--408

\crossref{10.1112/plms/s2-46.1.389}

\zmath{0025.40402|66.1243.01}











\end{thebibliography}



%\end{document}

 \newpage

































%%%%

\input{preamble}





\usepackage{rotating_pr}

\usepackage{desclist}

\usepackage{cmap}

\usepackage{textcomp}

\usepackage{nccboxes}

\usepackage{euscript}

\usepackage{mathrsfs} 

\usepackage{multicol}

\usepackage{nccfoots}

\usepackage{mathrsfs}

\usepackage{calrsfs}

\allowdisplaybreaks[4]

\usepackage{qrcode}

\usepackage{afterpage}

\usepackage{wasysym}

\usepackage{wrapfig}

\usepackage{slashbox}



\usepackage{multirow}

\newcommand{\specialcell}[2][c]{%

  \begin{tabular}[#1]{@{}c@{}}#2\end{tabular}}





\usepackage{pgf,pgfplots,tikz}

\usetikzlibrary{arrows,snakes,backgrounds,patterns,shapes.geometric,calc}



\nofiles



\oek % для генерации pdf, отдаваемого в oek.rsl.ru

%\samgtupubl % для генерации pdf, отдаваемого в типографию СамГТУ

%\pdfcrop % обрезка pdf для размещения на math-net.ru

%\draft % На корректуре





%\DeclareMathOperator{\tr}{tr}

%\DeclareMathOperator{\ink}{Ink}

%\DeclareMathOperator{\erf}{erf}







%\makeatletter

%\newcommand{\PersonRadRef}{\@langrustrue\def\refname{{\bf\small Избранные  труды \rupid{19071}{О.~А.~Репина} за последние 5 лет}}}

%\makeatother



\DeclareMathOperator{\tr}{tr}

\DeclareMathOperator{\arcsh}{arcsh}









\usepackage{listings}

\usepackage{color}

\definecolor{dkgreen}{rgb}{0,0.6,0}

\definecolor{gray}{rgb}{0.5,0.5,0.5}

\definecolor{mauve}{rgb}{0.58,0,0.82}



\renewcommand*{\lstlistingname}{Листинг}

\usepackage{caption}

%\DeclareCaptionFont{white}{\color{white}} 

\DeclareCaptionFormat{listing}{\parbox{\textwidth}{\centering \small  #1.~#3}}

\captionsetup[lstlisting]{format=listing}



\lstset{ %

  inputencoding=utf8,

  language=R,                     % the language of the code

  basicstyle=\footnotesize,       % the size of the fonts that are used for the code

  numbers=left,                   % where to put the line-numbers

  numberstyle=\tiny\color{gray},  % the style that is used for the line-numbers

  stepnumber=1,                   % the step between two line-numbers. If it's 1, each line

                                  % will be numbered

  numbersep=5pt,                  % how far the line-numbers are from the code

  backgroundcolor=\color{white},  % choose the background color. You must add \usepackage{color}

  showspaces=false,               % show spaces adding particular underscores

  showstringspaces=false,         % underline spaces within strings

  showtabs=false,                 % show tabs within strings adding particular underscores

  frame=single,                   % adds a frame around the code

  rulecolor=\color{black},        % if not set, the frame-color may be changed on line-breaks within not-black text (e.g. commens (green here))

  tabsize=2,                      % sets default tabsize to 2 spaces

  captionpos=t,                   % sets the caption-position to bottom

  breaklines=true,                % sets automatic line breaking

  breakatwhitespace=false,        % sets if automatic breaks should only happen at whitespace

  %title=\lstname,                 % show the filename of files included with \lstinputlisting;

                                  % also try caption instead of title

  keywordstyle=\color{blue},      % keyword style

  commentstyle=\color{dkgreen},   % comment style

  stringstyle=\color{mauve},      % string literal style

  %escapeinside={\%*}{*)},         % if you want to add a comment within your code

  morekeywords={*,in,+,-,=}            % if you want to add more keywords to the set

} 







\begin{document}







\setcounter{page}{601}

\input{specialpages.tex} 







\input{./01du/partition.tex}

\input{./02mdtt/partition.tex}

\input{./03matmod/partition.tex}

\input{./04short/partition.tex}



%\input{./polieces.tex}



\end{document}



\newpage



$\;$





\chead[\fancyplain{}{}]{\fancyplain{}{}}

  \rhead[]{\fancyplain{}{}}

  \lhead[\fancyplain{}{}]{}

  \cfoot[]{}

  \lfoot[]{}

  \rfoot[]{}

\renewcommand{\plainheadrulewidth}{0pt}

\renewcommand{\headrulewidth}{0pt}





  \pagestyle{fancyplain}



  \thispagestyle{plain}





\newpage



$\;$





\chead[\fancyplain{}{}]{\fancyplain{}{}}

  \rhead[]{\fancyplain{}{}}

  \lhead[\fancyplain{}{}]{}

  \cfoot[]{}

  \lfoot[]{}

  \rfoot[]{}

\renewcommand{\plainheadrulewidth}{0pt}

\renewcommand{\headrulewidth}{0pt}





  \pagestyle{fancyplain}



  \thispagestyle{plain}





%\newpage

%

%$\;$

%

%

%\chead[\fancyplain{}{}]{\fancyplain{}{}}

%  \rhead[]{\fancyplain{}{}}

%  \lhead[\fancyplain{}{}]{}

%  \cfoot[]{}

%  \lfoot[]{}

%  \rfoot[]{}

%\renewcommand{\plainheadrulewidth}{0pt}

%\renewcommand{\headrulewidth}{0pt}

%

%

%  \pagestyle{fancyplain}

%

%  \thispagestyle{plain}







\end{document}





\Russian



\addtocontents{toc}{\bigskip\par\hspace{-7mm}}

\addtocontents{etoc}{\bigskip\par\hspace{-7mm}}



$\;$



\hfill \qrcode[hyperlink,height=.8in]{http://doi.org/10.14498/vsgtu1651}

\vspace{-1.5in}





%$$\quad $$ \vskip.6in \par \hskip-7mm  {\huge  Краткие сообщения}

$$\quad $$ \vskip.6in \par \hskip-7mm  {\huge  Short Communications}



%$$\quad $$ \vskip.6in \par \hskip-7mm  {\huge  Mathematical Modeling,\\ Numerical Methods \\[1mm] and Software Complexes}



\addtocontents{toc}{{\large\bf Краткие сообщения}\bigskip}

\addtocontents{etoc}{{\large\bf Short Communications}\bigskip}



\vspace{-5mm}



%\Partition{Математическое моделирование}{Mathematical Modeling}



%\input{preamble}

%\usepackage{samgtu-name}

%

%

%

\renewcommand{\RuCitation}{\noindent\textbf{\CYRO\cyrb\cyrr\cyra\cyrz\cyre\cyrc\ \cyrd\cyrl\cyrya\ \cyrc\cyri\cyrt\cyri\cyrr\cyro\cyrv\cyra\cyrn\cyri\cyrya\\}\selectlanguage{english}{\EnAuthorInContents} {\EnTitleInContents}, \textit{\EnJournalName}\/ [\ParallelJournalName],  \JournalYear, vol.~\JournalVolume, no.~\JournalNumber,  pp.~\thefirstpage--\thelastpage. \doiurl. \selectlanguage{russian}\smallskip}



\renewcommand{\EnCitation}{\noindent\textbf{Please cite this article in press as:\\} {\EnAuthorInContents} {\EnTitleInContents}, \textit{\EnJournalName}\/ [\ParallelJournalName],  \JournalYear, vol.~\JournalVolume, no.~\JournalNumber,  pp.~\thefirstpage--\thelastpage.  \midoiurl. }

%

\vsgtu{1651} %registration number

\FirstPage{735}

\LastPage{749}

%

%

%

%\pdfcrop

%\draft

%

\ShortCommunication

%

%\begin{document}

%

%

%

%$\;$

%\vspace{1mm}

%

%\hfill \qrcode[hyperlink,height=.8in]{http://doi.org/10.14498/vsgtu1651}

%\vspace{-.8in}





\renewcommand{\baselinestretch}{0.98}



\udc{532.51, 517.958:531.3-324}

\msc{76F02, 76F45, 76M45, 76R05, 76U05}



\rutitle{Динамические равновесия неизотермической жидкости}

\entitle{Dynamic equilibria of a nonisothermal fluid}



\ruauthor{Е.~Ю.~Просвиряков\affil{1,2}}{П\,р\,о\,с\,в\,и\,р\,я\,к\,о\,в~Е.~Ю.}

\enauthor{E.~Yu.~Prosviryakov\affil{1,2}}{P\,r\,o\,s\,v\,i\,r\,y\,a\,k\,o\,v~E.~Yu.}





\ruaffil{\ruaffid{2665}{Институт машиноведения УрО РАН,\\

Россия, 620049, Екатеринбург, ул. Комсомольская, 34}.

\and

\ruaffid{7370}{Уральский федеральный университет \\ им.~первого Президента России Б.~Н.~Ельцина, \\ 

Россия, 620002, Екатеринбург, ул. Мира, 19}.}



\enaffil{\enaffid{2665}{Institute of Engineering Science, Urals Branch, Russian Academy of Sciences,\\

34, Komsomolskaya st., Ekaterinburg, 620049, Russian Federation}.

\and

\enaffid{7370}{Ural Federal University named after the First President of Russia B.~N.~Yeltsin, \\ 19, Mira st., Ekaterinburg, 620002, Russian Federation}.

}



\enauthorsdetails{1}{

\fullauthor{\enpid{41426}{Evgeny Yu. Prosviryakov}}

\CAuthor 

\orcid{0000-0002-2349-7801} 

\regalia{Dr. Phys. \& Math. Sci., Professor} 

\position{Head of Sector} 

\dept{Sector of Nonlinear Vortex Hydrodynamics}

\email[br]{evgen_pros@mail.ru}

}



\ruauthorsdetails{1}{

\fullauthor{\rupid{41426}{Евгений Юрьевич Просвиряков}} 

\CAuthor 

\orcid{0000-0002-2349-7801}

\regalia{доктор физико-математических наук} 

\position{заведующий сектором} 

\dept{сектор нелинейной вихревой гидродинамики}

\email{evgen_pros@mail.ru}

}





\enkeywords{dynamic equilibrium,  rotating fluid flow, exact solution, Boussinesq approximation, counterflows, increased velocities}



\rukeywords{динамическое  равновесие, вращающийся поток жидкости, точное решение, аппроксимация Буссинеска, противотечение, увеличенные скорости}



\ruabstract{В рамках точности приближения Буссинеска рассмотрены стационарные динамические равновесия вращающейся массы неизотермической жидкости. Показано, что в этом случае в жидкости наблюдается конечное число противотечений и усиление скоростей по сравнению с заданными на границе значениями, а также формирование зон положительного и отрицательного давления и температуры. }



\enabstract{In this paper, stationary dynamic equilibria of the rotating mass of a nonisothermal fluid are discussed within the accuracy limits of the Boussinesq approximation. It is demonstrated that, in this case, a fluid exhibits a finite number of counterflows, higher values of velocities than those specified on the boundary and the formation of zones of positive and negative pressures and temperatures.}



\newdate{privalovaa}{07}{10}{2018}

\date{\displaydate{privalovaa}}

\newdate{privalovab}{10}{11}{2018}

\revisiondate{\displaydate{privalovab}}

\newdate{privalovac}{12}{11}{2018}

\accepteddate{\displaydate{privalovac}}



\newdate{privalovae}{13}{12}{2018}

\renewcommand{\JournalDataOnline}{\displaydate{privalovae}}





%\ForPeerReview

\makeentitle





\Section[n]{Introduction}

The study of the motion of a rotating self-gravitating fluid was started as long ago as by Isaak Newton~[\citen{ulg:bib:Newtoni},~\citen{ulg:bib:Turnbull}], and it is currently far from being completed. Newton proved to be the first to endeavour to obtain, from the motion equations, a solution simulating the shape of the Earth. Newton’s pioneering study gave rise to numerous investigations in this field of mechanics and physics and a great impetus to the development of the mathematics used practically in all the sections of mathematics. The development of the potential theory, the creation of the theory of special functions and the formation of the rapidly developing theory of integrable systems should be mentioned first.



Practically all the most prominent scientists, such as Appel, Basset, Betti, Gagen, Helmholtz, Dedekind, Dirichlet, Kirchhoff, Lipschitz, Lame, McLoren, Poincare, Riemann, Chandrasekhar, Jacoby, took part in the construction of the theory~[\citen{ulg:bib:Poincare}--\citen{ulg:bib:Dirichlet}]. This list of discoverers of exact solutions can be extended~[\citen{ulg:bib:Poincare}--\citen{ulg:bib:Dirichlet}]. The history of discovering the solutions can be found in the monographs and surveys~[\citen{ulg:bib:Poincare}--\citen{ulg:bib:Dirichlet}]. It should be noted that the scientific breakthrough in the simulation of rotating equilibrium figures was made by the Russian scientists Steklov, Lyapunov, Ovsyannikov and Zel'dovich.



Note that the development of the theory of motion stability described by systems of ordinary differential equations was initiated by studying the stability of equilibrium figures and the motion of an elastic body partially filled with a fluid~[\citen{ulg:bib:Lyapunov1},~\citen{ulg:bib:Lyapunov2}]. It is the work in these fields that offered a formulation to the well-known Lyapunov theorems of stability and instability, among other things, from the first approximation~[\citen{ulg:bib:Lyapunov1},~\citen{ulg:bib:Lyapunov2}]. Original research has currently been done on the stability of Dirichlet and Jacobi ellipsoids~\cite{ulg:bib:Borisov}.



Ovsyannikov’s exact gas-hydrodynamic solution became the first solution    si\-mulating gas cloud motion in the class of linear velocities~\cite{ulg:bib:Borisov}. Afterwards models of Dyson and Fujimoto were proposed~\cite{ulg:bib:Borisov}.



There is another substantial model, which influenced the development of the theory of rotating fluids. It is the generalized solid proposed for research by Arnold~\cite{ulg:bib:Arnold}. In the survey made by Dolzhansky~\cite{ulg:bib:Dolzhansky} there is a detailed analysis of this approach with some generalizations for different classes of fluids.



As a rule, dynamic fluid equilibria are simulated with the application of two mathematical models. One model is Lagrangian fluid simulation, i.e. the use of the ordinary differential equations of Lagrange, Hamilton and Jacoby to describe the relative equilibrium (solid-state) of the continuum. The other model consists in using the Euler equations with potential forces to analyze the structures of a relative fluid equilibrium. Thus, in the simulation of dynamic fluid equilibria no account is taken of dissipation and temperature effect. In other words, an isothermal perfect incompressible fluid is dealt with. 

	

It is apparent that this deficiency in studying fluid flows can be compensated if this problem is analyzed with the application of the Oberbeck--Boussinesq system of equations. Note that it is not only the Boussinesq approximation that has not been used before to describe dynamic fluid equilibria, but also the characteristic scale of solutions. This statement needs to be clarified. Despite Chandrasekhar~\cite{ulg:bib:Chandrasekhar} used the Boussinesq approximation to simulate astrophysical objects, he used the virial method of physical field expansion for solving his problems. The application of this body of mathematics enabled one to reduce the Oberbeck--Boussinesq system to a system of ordinary differential equations and to use the classical problem statements when interpreting results. Similar reasoning holds for~[\citen{ulg:bib:Arnold},~\citen{ulg:bib:Dolzhansky}].

	

There are studies dealing with large-scale fluid or gas motions~\cite{ulg:bib:Chandrasekhar}, when it is necessary to take into account the inhomogeneity of the gravitational field, and with scales for which surface tension forces are principally important, e. g. the Hadamard-Rybczynski problem~[\citen{ulg:bib:Hadamard},~\citen{ulg:bib:Rybczynski}]. In both extreme cases, the principal equilibrium shape is spherical. The consideration of rotation was always viewed as some disturbance of the principal flow, however it was not viewed as a forming factor for equilibrium structures.



Note that the very possibility of the existence of equilibrium figures in the region of intermediate scales, when neither the inhomogeneity of the gravitational field nor the surface tension forces are essential, does not seem to have been studied. At any rate, the authors are unfamiliar with works viewing the problem this way.



Thus, we intend to ascertain the possible existence of equilibrium structures in the region of the so-called intermediate scales, when the Boussinesq approximation is obviously applicable to the simulation of a heat-conducting viscous incompres\-sible fluid.





\smallskip











\Section{Problem statement}

In the Boussinesq approximation~[\citen{ulg:bib:Chandrasekhar},~\citen{ulg:bib:Aristov1}] in the Cartesian coordinates (the $Oz$  axis is directed upwards), the system of equations describing the heat convection of a viscous incompressible fluid are written as

\begin{gather}

\frac{\partial \bf{V}}{\partial t}+\left(\bf{V}\cdot\nabla\right)\bf{V}=-\nabla \textit{P} + \nu \Delta \bf{V}+g \beta \textit{T} \bold{k},

\\    

      \label{eq:Aristov_Prosviryakov:NSE}

\frac{\partial T}{\partial t}+\bf{V}\cdot\nabla \textit{T}=\chi \Delta \textit{T},

\\     

\nabla \cdot \bf{V}=0.

 \end{gather}

The system of equations~\eqref{eq:Aristov_Prosviryakov:NSE} involves the following symbols: $\bf{V}=$ $\left(V_x , V_y , V_z\right)$  is the flow velocity vector;  $P$  is the pressure deviation from hydrostatic, which is related to the constant average fluid density  $\rho;$  $T$ is the deviation from the average temperature;   $\nu$  and $\chi$ are the coefficients of kinematic viscosity and thermal diffusivity of the fluid, respectively; $\beta$ is the temperature coefficient of fluid volumetric expansion; $\bold{k}$  is the unit vector directed vertically upwards, $\nabla$  is the Hamilton operator,  $\Delta$ is the Laplace operator (Laplacian).



The solution of system~\eqref{eq:Aristov_Prosviryakov:NSE} is sought in the approximation to the large scale. For large-scale flows, the characteristic horizontal scale greatly exceeds in magnitude the characteristic dimension during the entire fluid motion. Free surface curvature can in this scale be neglected, the vertical velocity $V_z$  being assumed zero. Thus, system~\eqref{eq:Aristov_Prosviryakov:NSE} is transformed to the following system:

\begin{gather}

\frac{\partial V_x}{\partial t}+V_x \frac{\partial V_x}{\partial x}+V_y \frac{\partial V_x}{\partial y}=-\frac{\partial P}{\partial x} + \nu \Bigl(\frac{\partial^2 V_x}{\partial x^2}+\frac{\partial^2 V_x}{\partial y^2}+\frac{\partial^2 V_x}{\partial z^2}\Bigr),

\\       

\frac{\partial V_y}{\partial t}+V_x \frac{\partial V_y}{\partial x}+V_y \frac{\partial V_y}{\partial y}=-\frac{\partial P}{\partial y} + \nu \Bigl(\frac{\partial^2 V_y}{\partial x^2}+\frac{\partial^2 V_y}{\partial y^2}+\frac{\partial^2 V_y}{\partial z^2}\Bigr),

\\

       \label{eq:Aristov_Prosviryakov:system}

\frac{\partial P}{\partial z}=g\beta T,

\\

\frac{\partial T}{\partial t}+V_x \frac{\partial T}{\partial x}+V_y \frac{\partial T}{\partial y}=\chi \Bigl(\frac{\partial^2 T}{\partial x^2}+\frac{\partial^2 T}{\partial y^2}+\frac{\partial^2 V_x}{\partial z^2}\Bigr),

\\

\frac{\partial V_x}{\partial x}+\frac{\partial V_y}{\partial y}=0.

 \end{gather}

 

Obviously, system \eqref{eq:Aristov_Prosviryakov:system} is nonlinear and overdetermined, four unknown functions of velocity and temperature being required from five equations. The hydrodynamic and temperature fields are computed as~\cite{ulg:bib:Aristov1}

\begin{gather}

V_x=U+xu_1+y u_2,

\quad 

V_y=V+xv_1+y v_2,

\\

       \label{eq:Aristov_Prosviryakov:solution}

P=P_0+xP_1+yP_2+\frac{x^2}{2}P_{11}+\frac{y^2}{2}P_{22}+xyP_{12},

\\

T=T_0+xT_1+yT_2+\frac{x^2}{2}T_{11}+\frac{y^2}{2}T_{22}+xyT_{12}.     

 \end{gather}

 

For horizontal coordinates, all the functions in the formulae of system \eqref{eq:Aristov_Prosviryakov:solution}  depend on the variable $z$  ant time $t.$ Substituting relation \eqref{eq:Aristov_Prosviryakov:solution} into system \eqref{eq:Aristov_Prosviryakov:system}, we arrive at the system

\begin{gather}

\hat{L}U+Uu_1+Vu_2+P_1=0,

\\

\hat{L}u_1+u^{2}_1+v_1u_2+P_{11}=0,

\\

\hat{L}u_2+u_1u_2+u_2v_2+P_{12}=0,

\\

\hat{L}V+Uv_1+Vv_2+P_{2}=0,

\\

\hat{L}v_1+u_1v_1+v_1v_2+P_{12}=0,

\\

\hat{L}v_2+u_2v_1+v^{2}_2+P_{22}=0,

\\

       \label{eq:Aristov_Prosviryakov:main system }

\hat{M}T_0+UT_1+VT_2-\chi\left(T_{11}+T_{22}\right)=0,

\\

\hat{M}T_1+UT_{11}+VT_{12}+u_1T_1+v_1T_2=0,

\\

\hat{M}T_2+UT_{12}+VT_{22}+u_2T_1+v_2T_2=0,

\\

\hat{M}T_{11}+2u_1T_{11}+2v_1T_{12}=0,

\\

\hat{M}T_{22}+2u_2T_{12}+2v_2T_{22}=0,

\\

\hat{M}T_{12}+u_1T_{12}+u_2T_{11}+v_1T_{22}+v_2T_{12}=0,

\\

u_1+v_2=0,

\\

\frac{\partial P_0}{\partial z}=g\beta T_0,

\quad 

\frac{\partial P_1}{\partial z}=g\beta T_1,

\quad 

\frac{\partial P_2}{\partial z}=g\beta T_2,

\\

\frac{\partial P_{11}}{\partial z}=g\beta T_{11},

\quad 

\frac{\partial P_{12}}{\partial z}=g\beta T_{12},

\quad 

\frac{\partial P_{22}}{\partial z}=g\beta T_{22}.

 \end{gather}

Here, the following differential parabolic operators in partial derivatives are intro\-duced:

$$\hat{L}=\frac{\partial }{\partial t}-\nu\frac{\partial^2 }{\partial z^2},\quad 

\hat{M}=\frac{\partial }{\partial t}-\chi\frac{\partial^2 }{\partial z^2}. 

$$





\smallskip





\Section{Analyzing the solvability of system}

In what follows we do not discuss the solvability of system \eqref{eq:Aristov_Prosviryakov:main system } as a whole, but restrict our analysis to one special case. We will seek a solution to the system of equations when the boundaries of the fluid layer rotate. In~\cite{ulg:bib:Aristov2} the validity of equalities for the velocity gradients in the simulation of axisymmetric flows within class \eqref{eq:Aristov_Prosviryakov:solution} was demonstrated,

\begin{gather}

       \label{eq:Aristov_Prosviryakov:sol1}

u_1=0,

\quad 

v_2=0,       

\quad 

u_2=-v_1.       

 \end{gather}

Next the lower boundary of the layer  $z=0$ is assumed to be immobile. The upper boundary describable by the equation $z=h$  rotates with constant angular velocity  $\Omega$. Using equalities \eqref{eq:Aristov_Prosviryakov:sol1}, we obtain the following solutions to system~\eqref{eq:Aristov_Prosviryakov:main system }, which determine a number of functions in class \eqref{eq:Aristov_Prosviryakov:solution} for the coordinates  $x$ and  $y$:

\begin{gather}

u_2=\frac{\Omega z}{h}, 

\quad 

v_1=-\frac{\Omega z}{h},

\quad 

P_{12}=0, 

\\

\label{eq:Aristov_Prosviryakov:sol2}

 P_{11}=P_{22}=-u_2v_1=\frac{\Omega^2 z^2}{h^2},

\quad 

T_{12}=0, 

\\ g\beta T_{11}=g\beta T_{22}=\frac{2\Omega^2 z}{h^2}.

\end{gather}



For the homogeneous summands in expressions \eqref{eq:Aristov_Prosviryakov:solution}, due to solution \eqref{eq:Aristov_Prosviryakov:sol2}, system \eqref{eq:Aristov_Prosviryakov:main system } yields two closed subsystems. First system has of the form

\begin{gather}

\hat{L}U+\frac{\Omega z}{h}V+P_1=0, 

\\

\hat{L}V-\frac{\Omega z}{h}U+P_{2}=0,

\\

\label{eq:Aristov_Prosviryakov:sys1}

\hat{M}T_1+\frac{2\Omega^2 z}{g\beta h^2}U-\frac{\Omega z}{h}T_2=0, 

\\ \hat{M}T_2+\frac{2\Omega^2 z}{g\beta h^2}V+\frac{\Omega z}{h}T_1=0, 

\\

\frac{\partial P_1}{\partial z}=g\beta T_1, 

\quad  

\frac{\partial P_2}{\partial z}=g\beta T_2.

\end{gather}



Second system has of the form

\begin{gather}

\label{eq:Aristov_Prosviryakov:sys2}

\hat{M}T_0+UT_1+VT_2=\frac{4\chi\Omega^2 z}{g\beta h^2},

\\

\frac{\partial P_0}{\partial z}=g\beta T_0.

\end{gather}



Note that the integration of system \eqref{eq:Aristov_Prosviryakov:sys2} is possible only after the integration of system \eqref{eq:Aristov_Prosviryakov:sys1}. In order to find a solution to system \eqref{eq:Aristov_Prosviryakov:sys1}, we introduce the following complex functions:

$$

W=U+iV,\quad 

S=P_1+iP_2,\quad 

\Theta=g\beta\left(T_1+iT_2\right).$$

Here  $i$ is an imaginary unit. Thus, system \eqref{eq:Aristov_Prosviryakov:sys1} is in this case transformed to the following form:

 \begin{gather}

(\hat{L}-i\Omega z) W+S=0,

\\

       \label{eq:Aristov_Prosviryakov:sys1complex}

(\hat{M}+i\Omega z) \Theta+2\Omega^2 z W=0,       

\\

\frac{d S}{dz}=\Theta.      

 \end{gather}

 

To reduce the order of the system, the following transformation is made:

 \begin{gather}

       \label{eq:Aristov_Prosviryakov:trans}

S=S_1+i\Omega z W,

\quad 

\Theta=\Theta_1+2i\Omega W.    

\end{gather}

Substituting \eqref{eq:Aristov_Prosviryakov:trans} into system \eqref{eq:Aristov_Prosviryakov:sys1complex}, we obtain the equalities

 \begin{gather}

\nu\frac{d^2 W}{dz^2}+S_1=0,

\\

       \label{eq:Aristov_Prosviryakov:trans2}

\chi\frac{d^2 \Theta_1}{d z^2}+i\Omega \Bigl(z \Theta_1- 2\frac{\chi}{\nu}S_1\Bigr)=0,       

\\

\frac{d S_1}{dz}=\Theta_1-i\Omega z^2 \frac{d}{dz}\Bigl(\frac{W}{z}\Bigr). 

 \end{gather}

 

For further transformations, we use the differential identity

$$

\frac{d^2 W}{dz^2}=\frac{1}{z}\frac{d}{dz}\left(z^2\frac{d}{dz}\Bigl(\frac{W}{z}\Bigr)\right).

$$

This identity holds for any twice-differentiable function. Consequently, by differentiating the last equation in system \eqref{eq:Aristov_Prosviryakov:trans2}, we arrive at a system for determining the functions  $\Theta_1$  and  $S_1,$

 \begin{gather}

       \label{eq:Aristov_Prosviryakov:trans3}

\chi\frac{d^2 \Theta_1}{d z^2}+i\Omega \left(z \Theta_1- 2\frac{\chi}{\nu}S_1\right)=0,

\quad 

\nu\frac{d^2 S_1}{dz^2}=i\Omega z S_1 +\nu \frac{d \Theta_1}{dz}.          

 \end{gather}

 

Taking into account that equations \eqref{eq:Aristov_Prosviryakov:trans2} offer an explicit expression for  $W$, we obtain an integral of the initial system with automatic order reduction, system \eqref{eq:Aristov_Prosviryakov:trans3} being transformed to the operator equation

$$

\left[\Bigl(\nu\frac{d^2}{dz^2}-i\Omega z\Bigr)\Bigl(\chi\frac{d^2}{dz^2}+i\Omega z\Bigr)-2i\Omega\chi\frac{d}{dz}\right]\Theta_1=0.

$$

Simplifying the latter equality, we derive the following fourth-order ordinary differential equation with complex coefficients:

 \begin{gather}

       \label{eq:Aristov_Prosviryakov:trans4}

\Bigl[\nu\chi\frac{d^4}{dz^4}+i\Omega\left(\nu-\chi\right)\frac{d^2}{dz^2}z+\Omega^2 z^2\Bigr]\Theta_1=0.    

 \end{gather}



\smallskip







\Section{Solving equation when the dissipative coefficients coincide}

 We assume in equation \eqref{eq:Aristov_Prosviryakov:trans4} that $\nu=\chi$  (the Prandtl number equals one). In this case, equation \eqref{eq:Aristov_Prosviryakov:trans4} acquires the form

  \begin{gather}

       \label{eq:Aristov_Prosviryakov:Pr1}

\frac{d^4 \Theta_1}{dz^4}+\left(36\sigma^3 z\right)^2 \Theta_1=0,    

 \end{gather}

where  $\sigma=\frac{1}{6}\sqrt[3]{\frac{\Omega}{6\nu h}}$ is the dispersion relation. For the convenience of computations, in order to illustrate the physical meaning of the solutions and to make the notation concise, we introduce a new variable  $Z=\frac{z}{h \sigma}.$ The general solution of equation \eqref{eq:Aristov_Prosviryakov:Pr1} is written as

\begin{multline}

\label{eq:Aristov_Prosviryakov:soltemp}

\Theta_1={^{3}_{0}F} \Bigl[\Bigl\{\frac{1}{2}, \frac{2}{3}, \frac{5}{6}\Bigr\},-Z^6\Bigr]C_1 +

{^{3}_{0}F} 

\Bigl[\Bigl\{\frac{2}{3}, \frac{5}{6}, \frac{7}{6}\Bigr\},-Z^6\Bigr] C_2 Z +

\\

+

{^{3}_{0}F} \Bigl[\Bigl\{\frac{5}{6}, \frac{7}{6}, \frac{4}{3}\Bigr\},-Z^6\Bigr] C_3 Z^2  +

{^{3}_{0}F} \Bigl[\Bigl\{\frac{7}{6}, \frac{4}{3},\frac{3}{2}\Bigr\},-Z^6\Bigr] C_4 Z^3  . 

\end{multline}

Here,  ${^{3}_{0}F}=F\left(a; b; Z\right)=F\left(a_1, a_2, \ldots, a_p; b_1, b_2, \ldots, b_q; Z\right)$ is a generalized hypergeometric function~\cite{ulg:bib:Polyanin}. The choice of using the generalized hypergeometric function for writing the solutions is dictated by notation conciseness. Note that there is another form of writing the solution of equation \eqref{eq:Aristov_Prosviryakov:Pr1}, which is based on using the Mittag--Leffler functions~\cite{ulg:bib:Polyanin}. Note that each particular solution \eqref{eq:Aristov_Prosviryakov:soltemp} is a function of a real variable and that the general solution is considered in the complex plane. Thus, at least one integration constant is a complex number.



We now turn to finding the generalized pressure $S_1$  from the differential equation \eqref{eq:Aristov_Prosviryakov:trans2} 

 \begin{gather}

       \label{eq:Aristov_Prosviryakov:eqpr}

S_1=-\frac{i}{2\sigma}\frac{d^2 \Theta_1}{d z^2}+\frac{z}{2}\Theta_1.

\end{gather}



The substitution of solution \eqref{eq:Aristov_Prosviryakov:soltemp} into equation \eqref{eq:Aristov_Prosviryakov:eqpr} gives

\begin{multline}

\label{eq:Aristov_Prosviryakov:solstress}

S_1= {^{3}_{0}F} 

\Bigl[\Bigl\{\frac{1}{2}, \frac{2}{3}, \frac{5}{6}\Bigr\},-Z^6\Bigr] \frac{C_1}{2}Z+

{^{3}_{0}F} 

\Bigl[\Bigl\{\frac{2}{3}, \frac{5}{6}, \frac{7}{6}\Bigr\},-Z^6\Bigr] \frac{C_2}{2}Z^2+

\\

+{^{3}_{0}F} 

\Bigl[\Bigl\{\frac{5}{6}, \frac{7}{6}, \frac{4}{3},\Bigr\},Z\Bigr] \sqrt[3]{36\sigma^2} \frac{2\cdot 5C_3}{6!}Z^3+

{^{3}_{0}F} 

\Bigl[\Bigl\{\frac{7}{6}, \frac{4}{3}, \frac{3}{2}\Bigr\},Z\Bigr] \sigma \frac{2\cdot 5C_4}{6!}Z^4-

\\

-\frac{1}{72\sigma}i\left\{ {^{3}_{0}F} 

\Bigl[\Bigl\{\frac{5}{6}, \frac{7}{6}, \frac{4}{3}\Bigr\},Z\Bigr] \sqrt[3]{676\sigma^2} C_3 Z^3 

-{^{3}_{0}F} 

\Bigl[\Bigl\{\frac{7}{6}, \frac{4}{3}, \frac{3}{2}\Bigr\},Z\Bigr] 6 \sigma C_4 Z+

\right.

\\

+{^{3}_{0}F} 

\Bigl[\Bigl\{\frac{5}{6}, \frac{7}{6}, \frac{4}{3}\Bigr\},Z\Bigr] 3 \sigma^2 C_1 Z^4 +

\\

+{^{3}_{0}F} 

\Bigl[\Bigl\{\frac{5}{3}, \frac{11}{6}, \frac{13}{6},\Bigr\},Z\Bigr] \frac{2^4\cdot 3^3\sigma^2}{7!}\sqrt[3]{\frac{3\sigma}{4}} C_2 Z^5

+

\\

+{^{3}_{0}F} 

\Bigl[\Bigl\{\frac{5}{3}, \frac{11}{6}, \frac{13}{6},\Bigr\},Z\Bigr] \frac{2^4\cdot 3^4\sigma^2}{7!}\sqrt[3]{6\sigma} C_2 Z^5+

\\

+{^{3}_{0}F} 

\Bigl[\Bigl\{\frac{11}{6}, \frac{13}{6}, \frac{7}{3} \Bigr\},Z\Bigr] \frac{2^2\cdot 3^4\sigma^2}{7!}\sqrt[3]{\frac{9\sigma^2}{2}} C_3 Z^6+

\\

+{^{3}_{0}F} \Bigl[\Bigl\{\frac{13}{6}, \frac{7}{3}, \frac{5}{2} \Bigr\},Z\Bigr] \frac{2\cdot 5\cdot 11\sigma^3}{7!} C_4 Z^7-

\\

-{^{3}_{0}F} \Bigl[\Bigl\{\frac{5}{2}, \frac{8}{3}, \frac{17}{6} \Bigr\},Z\Bigr] \frac{2^7\cdot 3^3\cdot 7\sigma^4}{11!} C_1 Z^{10}-

\\

-{^{3}_{0}F} \Bigl[\Bigl\{\frac{8}{3}, \frac{17}{6}, \frac{19}{6} \Bigr\},Z\Bigr] \frac{2^9\cdot 3^6\sigma^4}{13!}\sqrt[3]{\frac{3\sigma}{4}} C_2 Z^{11}-

\\

-{^{3}_{0}F} \Bigl[\Bigl\{\frac{17}{6}, \frac{19}{6}, \frac{10}{3} \Bigr\},Z\Bigr] \frac{2^8\cdot 3^5\cdot 5\sigma^4}{14!}\sqrt[3]{\frac{9\sigma^2}{2}} C_3 Z^{12}-

\\

\left.

-{^{3}_{0}F} \Bigl[\Bigl\{\frac{19}{6}, \frac{10}{3}, \frac{7}{2} \Bigr\},Z\Bigr] \frac{2^7\cdot 3^2\cdot 5\cdot 11\sigma^5}{14!} C_4 Z^{13}

\right\}.

\end{multline}





Complex velocity $W$  can be expressed in two ways. Firstly, it can be obtained by integrating the first equation in system \eqref{eq:Aristov_Prosviryakov:sys1complex}, but it is a second-order equation. Hence, we obtain a solution depending on six integration constants, whereas system \eqref{eq:Aristov_Prosviryakov:trans2} is a fifth-order equation system. Thus our problem is complicated by additional analysis of overdetermination. A correct expression for velocity $W$  can be obtained by the integration of the third equation in system \eqref{eq:Aristov_Prosviryakov:sys1complex},

\begin{multline}

\label{eq:Aristov_Prosviryakov:solvelo}

\Omega W=C_5 Z+ \left({^{3}_{0}F} \Bigl[\Bigl\{\frac{2}{3}, \frac{5}{6}, \frac{7}{6} \Bigr\},Z\Bigr]-1\right) \frac{i}{2} \sqrt[3]{\frac{\sigma}{36}} C_2 Z-

\\

-{^{4}_{1}F} \Bigl[\Bigl\{-\frac{1}{6} \Bigr\}, \Bigl\{\frac{1}{2},\frac{2}{3}, \frac{5}{6}, \frac{5}{6} \Bigr\},Z\Bigr]\frac{2i \Gamma\left(-\frac{1}{6}\right)}{4!\Gamma\left(\frac{5}{6}\right)} C_1 Z^2+

\\

+{^{4}_{1}F} \Bigl[\Bigl\{-\frac{1}{6} \Bigr\},\Bigl\{\frac{5}{6},\frac{7}{6}, \frac{4}{3}, \frac{3}{2} \Bigr\},Z\Bigr]\frac{2\cdot 5\Gamma\left(-\frac{1}{6}\right)}{6!\Gamma\left(\frac{5}{6}\right)} C_4+

\\

+{^{4}_{1}F} \Bigl[\Bigl\{\frac{1}{6} \Bigr\},\Bigl\{\frac{5}{6},\frac{7}{6}, \frac{7}{6}, \frac{4}{3} \Bigr\},Z\Bigr]\frac{2\cdot 5i \Gamma\left(\frac{1}{6}\right)}{6!\Gamma\left(\frac{7}{6}\right)} \sqrt[3]{\frac{\sigma^2}{6}} C_3 Z^2+

\\

+{^{4}_{1}F} \Bigl[\Bigl\{\frac{1}{3} \Bigr\},\Bigl\{\frac{7}{6},\frac{4}{3}, \frac{4}{3}, \frac{3}{2} \Bigr\},Z\Bigr]\frac{2^4\cdot 3\cdot 5\cdot 7i \sigma \Gamma\left(\frac{1}{3}\right)}{9!\Gamma\left(\frac{4}{3}\right)} C_4 Z^3-

\\

-{^{4}_{1}F} \Bigl[\Bigl\{\frac{1}{3} \Bigr\},\Bigl\{\frac{4}{3},\frac{3}{2}, \frac{5}{3}, \frac{11}{6} \Bigr\},Z\Bigr]\frac{2^2\cdot 5\sigma \Gamma\left(\frac{1}{3}\right)}{6!\Gamma\left(\frac{4}{3}\right)} C_1 Z^3-

\\

-{^{4}_{1}F} \Bigl[\Bigl\{\frac{1}{2} \Bigr\},\Bigl\{\frac{3}{2},\frac{5}{3}, \frac{11}{6}, \frac{13}{6} \Bigr\},Z\Bigr]\frac{\sigma }{4!}\sqrt[3]{\frac{\sigma}{36}} C_2 Z^4-

\\

-{^{4}_{1}F} \Bigl[\Bigl\{\frac{2}{3} \Bigr\},\Bigl\{\frac{5}{3},\frac{11}{6}, \frac{13}{6}, \frac{7}{3} \Bigr\},Z\Bigr] \frac{2^4\cdot 3^2\cdot 7\sigma \Gamma\left(\frac{2}{3}\right)}{9!\Gamma\left(\frac{5}{3}\right)}\sqrt[3]{\frac{\sigma}{6}} C_3 Z^5-

\\

-{^{4}_{1}F} \Bigl[\Bigl\{\frac{5}{6} \Bigr\},\Bigl\{\frac{3}{2},\frac{5}{3}, \frac{11}{6}, \frac{11}{6} \Bigr\},Z\Bigr] \frac{2^3\cdot 3^2\cdot 7i\sigma^2 \Gamma\left(\frac{5}{6}\right)}{9!\Gamma\left(\frac{11}{6}\right)} C_1 Z^6-

\\

-{^{4}_{1}F} \Bigl[\Bigl\{\frac{5}{6} \Bigr\},\Bigl\{\frac{11}{6},\frac{13}{6}, \frac{7}{3}, \frac{5}{2} \Bigr\},Z\Bigr] \frac{2^3\cdot 5^2\cdot 11\cdot 13\cdot 83\sigma^2 \Gamma\left(\frac{5}{6}\right)}{13! \Gamma\left(\frac{11}{6}\right)} C_4 Z^6-

\\ 

-{^{4}_{1}F} \Bigl[\Bigl\{\frac{7}{6} \Bigr\},\Bigl\{\frac{11}{6},\frac{13}{6}, \frac{13}{6}, \frac{7}{3} \Bigr\},Z\Bigr] \frac{2\cdot 3^2i\sigma^2 \Gamma\left(\frac{7}{6}\right)}{9! \Gamma\left(\frac{13}{6}\right)}\sqrt[3]{\frac{\sigma^2}{6}} C_3 Z^8-

\\

-{^{4}_{1}F} \Bigl[\Bigl\{\frac{4}{3} \Bigr\},\Bigl\{\frac{13}{6},\frac{7}{3}, \frac{7}{3}, \frac{5}{2} \Bigr\},Z\Bigr] \frac{2^3\cdot 5^2\cdot 11\cdot 13 i\sigma^3 \Gamma\left(\frac{4}{3}\right)}{13!\Gamma\left(\frac{7}{3}\right)} C_4 Z^9+

\\

+{^{4}_{1}F} \Bigl[\Bigl\{\frac{4}{3} \Bigr\},\Bigl\{\frac{7}{3},\frac{5}{2}, \frac{8}{3}, \frac{17}{6} \Bigr\},Z\Bigr] \frac{2^3\cdot 3\cdot 5\cdot 7\sigma^3 \Gamma\left(\frac{4}{3}\right)}{11! \Gamma\left(\frac{7}{3}\right)} C_1 Z^9+

\\

+{^{4}_{1}F} \Bigl[\Bigl\{\frac{3}{2} \Bigr\},\Bigl\{\frac{5}{2},\frac{8}{3}, \frac{17}{6}, \frac{19}{6} \Bigr\},Z\Bigr] \frac{2^7\cdot 3^5\sigma^3}{13!} \sqrt[3]{\frac{\sigma}{36}} C_2 Z^{10}+

\\

+{^{4}_{1}F} \Bigl[\Bigl\{\frac{5}{3} \Bigr\},\Bigl\{\frac{8}{3},\frac{17}{6}, \frac{19}{6}, \frac{10}{3} \Bigr\},Z\Bigr] \frac{2^3\cdot 3^4\cdot 5\sigma^3 \Gamma\left(\frac{5}{3}\right)} {13! \Gamma\left(\frac{8}{3}\right)} \sqrt[3]{\frac{\sigma^2}{6}} C_3 Z^{11}+

\\

+{^{4}_{1}F} \Bigl[\Bigl\{\frac{11}{6} \Bigr\},\Bigl\{\frac{17}{6},\frac{19}{6}, \frac{10}{3}, \frac{7}{2} \Bigr\},Z\Bigr] \frac{2^6\cdot 5\cdot 11\sigma^4 \Gamma\left(\frac{11}{6}\right)} {14! \Gamma\left(\frac{17}{6}\right)} C_4 Z^{12}-

\\

-{^{4}_{1}F} \Bigl[\Bigl\{\frac{7}{3} \Bigr\},\Bigl\{\frac{10}{3},\frac{7}{2}, \frac{11}{3}, \frac{23}{6} \Bigr\},Z\Bigr] \frac{2^6\cdot 3\cdot 7^2\cdot 13\sigma^5 \Gamma\left(\frac{7}{3}\right)} {17! \Gamma\left(\frac{10}{3}\right)} C_1 Z^{15}-

\\

-{^{4}_{1}F} \Bigl[\Bigl\{\frac{5}{2} \Bigr\},\Bigl\{\frac{7}{2},\frac{11}{3}, \frac{23}{6}, \frac{25}{6} \Bigr\},Z\Bigr] \frac{2^8\cdot 3^7\cdot 7\sigma^5 } {19!} \sqrt[3]{\frac{\sigma}{36}} C_2 Z^{16}-

\\

-{^{4}_{1}F} \Bigl[\Bigl\{\frac{8}{3} \Bigr\},\Bigl\{\frac{11}{3},\frac{23}{6}, \frac{25}{6}, \frac{13}{3} \Bigr\},Z\Bigr] \frac{2^7\cdot 3^6\cdot 5\sigma^5 \Gamma\left(\frac{8}{3}\right)C_3 Z^{17}} {19! \Gamma\left(\frac{11}{3}\right)} \sqrt[3]{\frac{\sigma^2}{6}} -

\\

-{^{4}_{1}F} \Bigl[\Bigl\{\frac{8}{3} \Bigr\},\Bigl\{\frac{23}{6},\frac{25}{6}, \frac{13}{3}, \frac{9}{2} \Bigr\},Z\Bigr] \frac{2^8\cdot 3^2\cdot 5^2\cdot 11\cdot 17\sigma^6 \Gamma\left(\frac{17}{6}\right)} {21! \Gamma\left(\frac{23}{6}\right)} \sqrt[3]{\frac{\sigma^2}{6}} C_4 Z^{18}.

\end{multline}

Here $\Gamma\left(\lambda\right)= \displaystyle\int_{0}^{+\infty}r^{\lambda-1}\exp\left(-r\right)dr$ is the Euler gamma function~\cite{ulg:bib:Polyanin}. Solutions~\eqref{eq:Aristov_Prosviryakov:soltemp}, \eqref{eq:Aristov_Prosviryakov:solstress},~\eqref{eq:Aristov_Prosviryakov:solvelo} are written as a linear  combination of generalized hypergeometric functions in the form of infinite power series. Consequently, examination of solution convergence is an urgent problem. Using the properties of convergence of these functions, we can state that the convergence radius for each solution equals infinity~\cite{ulg:bib:Polyanin}. It means that the solution of our problem always converges.



Let us now formulate the boundary conditions to determine the integration constants of equation system \eqref{eq:Aristov_Prosviryakov:sys1}. Without any claim to solution generality and without studying the shapes of the fluid, we give a simple example of such solutions. Let the simplest and most physically meaningful solution be a flow between two infinite disks rotating with free velocities. This is possible, since the solution can always be shifted along the $Oz$-axis, an appropriate distance between the disks being chosen. The disks must be nonuniformly heated to a proper temperature (the temperature must be related to the radius by a well-defined law). Assume that the lower boundary is isothermal (the temperature of the lower disk is zero); the adhesion condition holds on both boundaries. Thus, the following equalities are valid:

$$\text{when} \ \ z=0, \ \ U=V=0, \ \ T_1=T_2=0;$$

$$\text{when} \ \ z=h, \ \ U=A\cos\alpha, \ \ V=A\sin\alpha, \ \ T_1=B, \ \ T_2=C, \ \ P_1=\tau_1, \ \ P_2=\tau_2.$$



The conditions for determining the integration constants, written on the boundary, do not determine the integration constants uniquely. The system of linear algebraic equations for determining the integration constants is underdetermined, since there are few boundary conditions. To close the system of equations, we specify the fluid flow rate conditions as

$$\int _{0}^{h}Udz =Q_1,\quad  \int_{0}^{h}Vdz =Q_2.$$



When the flow rate is specified, it is generally assumed that  $Q_1=Q_2=0$~\cite{ulg:bib:Aristov2}. It is a physically justified convention~[\citen{ulg:bib:Aristov1},~\citen{ulg:bib:Aristov2}]. However, it is possible to obtain a solution by specifying a nonzero flow rate. Next, not limiting the generality of the reasoning, we will study the solutions simulating fluid motions with the zero flow rate. 



For the above-introduced complex velocities, pressure and temperature ~\eqref{eq:Aristov_Prosviryakov:soltemp}, \eqref{eq:Aristov_Prosviryakov:solstress},~\eqref{eq:Aristov_Prosviryakov:solvelo}, in view of transformations \eqref{eq:Aristov_Prosviryakov:trans}, the following boundary conditions are valid:

\begin{gather}

\label{eq:Aristov_Prosviryakov:bond}

\text{when} \ \ z=0: \ \ W=0, \ \ \Theta_1=0;

\\

\text{when} \ \ z=h: \ \ W=A\exp\left(i\alpha\right), \ \ \Theta_1=B+2A\Omega\sin\alpha+i\left(C-2A\Omega\cos\alpha\right),

\\

\int_{0}^{h}Wdz =0.

\end{gather}



With the boundary conditions \eqref{eq:Aristov_Prosviryakov:bond}, it follows from equality \eqref{eq:Aristov_Prosviryakov:soltemp} that  $C_1=0$, and from solution \eqref{eq:Aristov_Prosviryakov:solvelo} that  $C_4=0$. Thus, the expressions for the determination of hydrodynamic fields become much simpler. Remember that, in the computation of the integration constants, the variable $z=h\sigma Z$  is to be replaced.











\Section{Analysis of the solutions}

 Using conditions \eqref{eq:Aristov_Prosviryakov:bond}, it is easy, but rather cumbersome, to compute the complex  $C_2,$ $C_3,$ and $C_5$  which are related to the dispersion relation   $\sigma$ (kinematic viscosity) by the fractional linear law. Analyzing conditions \eqref{eq:Aristov_Prosviryakov:bond}, we find that the boundary value problem under study is meaningful at any positive value of  $\sigma.$ Nevertheless, the structure and behavior of the solutions depend considerably on the value of  $\sigma.$



It is well known that, to analyze the polynomial solutions of system \eqref{eq:Aristov_Prosviryakov:sys1}, it would suffice to localize the roots of the linear combination of the generalized hypergeometric functions~\cite{ulg:bib:Polyanin}. Any hypergeometric function is known to be monotonically increasing with a zero value at the origin. Therefore, the appearance of other zero values is possible only when there is a combination of functions with at least one negative weight coefficient~\cite{ulg:bib:Polyanin}. The appearance of negative coefficients depends on $\sigma$  and the function domain, which also depends on the dispersion relation $Z\in\left[0, \frac{1}{\sigma}\right].$  Temperature is characterized by three types of solutions. 

When  $\sigma\geq\frac{3}{2},$ the function describing temperature is a strictly monotonically increasing or decreasing function (fig.~\ref{fig_prosv1}). 

If the evaluation  $\frac{9}{10}<\sigma<\frac{3}{2}$  is valid for the dispersion relation, there is exactly one zero value of temperature inside the fluid layer (fig.~\ref{fig_prosv2}). With the further decrease, temperature gets a finite number of zeros (fig.~\ref{fig_prosv3}). Consequently, the fluid layer is divided into the regions with positive and negative temperature, which are characterized by alternating signs. The analysis of pressure is similar to that of temperature. 





Velocities are characterized by only two qualitative behaviors of the functions (figs.~\ref{fig_prosv2} and ~\ref{fig_prosv3}). Thus, at any  $\sigma,$ there appear fluid counterflows. It was shown in~[\citen{ulg:bib:Aristov3}--\citen{ulg:bib:Gorshkov}] that the appearance of counterflows is a purely nonlinear effect, i.\,e., the viscous forces are comparable to the convective acceleration, and even without Coriolis forces. However, when  $\sigma$ tends to zero, counterflow zones are formed, whose thickness decreases with  $\sigma.$ Other values of the dispersion relation are characteristic of velocity, and they are easy to compute. However, the values obtained for the temperature and pressure fields are a good reference point for refining the characteristic dispersion numbers of velocities.





When analyzing the solutions (fig.~\ref{fig_prosv3}), we can see their ``fluctuation''. In~\cite{ulg:bib:Aristov7} localized convective flows in terms were studied in the axisymmetric formulation. Solution ``damping'' was observed there. We use a wider class of solutions. There may be not only ``damping'' (when $\sigma$  is close to one), but also ``resonance'' initiation. Yet, the solution does not become singular, and the Boussinesq approxi\-mation is not violated in the simulation of the convective flow. This difference is attributed to the account of the quadratic summands for temperature. If tempera\-ture is considered to be a linear function with respect to the horizontal coordinates, the solutions obtained in this study coincide with those discussed in~\cite{ulg:bib:Aristov4}.



\begin{figure}[h!]



\begin{tabular}{p{.45\textwidth}p{.45\textwidth}}

\centerline{\includegraphics[width=0.45\textwidth]{fig_prosv1}} &

\centerline{\includegraphics[width=0.45\textwidth]{fig_prosv2}} \\

\caption{A qualitative form of the solutions  $T_1,$ $T_2,$ $P_1,$ $P_2$ of system (\ref{eq:Aristov_Prosviryakov:sys1}) when $\sigma\geq\frac{3}{2}$ ($1-\tau_1>0$,  $\tau_2>0$, $A>0$, $B>0$, $C>0$; $2-\tau_1<0$, $\tau_2<0$, $A<0$, $B<0$, $C<0$) \label{fig_prosv1}} &

\caption{A qualitative form of the solutions  $T_1,$ $T_2,$ $P_1,$ $P_2,$ $U,$ $V$ of system (\ref{eq:Aristov_Prosviryakov:sys1}) when $\frac{9}{10}<\sigma<\frac{3}{2}$ ($1-\tau_1>0$,  $\tau_2>0$, $A>0$, $B>0$, $C>0$; $2-\tau_1<0$, $\tau_2<0$, $A<0$, $B<0$, $C<0$) \label{fig_prosv2} }

\\

\centerline{\includegraphics[width=0.45\textwidth]{fig_prosv3}} &



\vspace{-2cm}

\caption{A qualitative form of the solutions  $T_1,$ $T_2,$ $P_1,$ $P_2,$ $U,$ $V$ of system (\ref{eq:Aristov_Prosviryakov:sys1}) when $\sigma\leq\frac{9}{10}$ ($1-\tau_1>0$,  $\tau_2>0$, $B>0$, $C>0$; $2-\tau_1<0$, $\tau_2<0$, $B<0$, $C<0$) \label{fig_prosv3} } \\





\end{tabular}



\vspace{-3cm}



%%\noindent

%%\begin{minipage}[h]{0.48\textwidth}

%%

%%\smallskip \footnotesize

%%Fig.~\ref{fig1}.

%%A qualitative form of the solutions  $T_1,$ $T_2,$ $P_1,$ $P_2$ of system (\ref{eq:Aristov_Prosviryakov:sys1}) when $\sigma\geq\frac{3}{2}$ ($1-\tau_1>0$,  $\tau_2>0$, $A>0$, $B>0$, $C>0$; $2-\tau_1<0$, $\tau_2<0$, $A<0$, $B<0$, $C<0$) \label{fig1} 

%%\end{minipage} \hfill

%%\begin{minipage}[h]{0.5\textwidth}

%%\centering

%%\includegraphics[width=0.98\textwidth]{fig1}

%%\end{minipage}

%\end{figure}

%

%

%\begin{figure}[h!]

%\noindent

%\begin{minipage}[h]{0.48\textwidth}

%\caption{A qualitative form of the solutions  $T_1,$ $T_2,$ $P_1,$ $P_2,$ $U,$ $V$ of system (\ref{eq:Aristov_Prosviryakov:sys1}) when $\frac{9}{10}<\sigma<frac{3}{2}$ ($1-\tau_1>0$,  $\tau_2>0$, $A>0$, $B>0$, $C>0$; $2-\tau_1<0$, $\tau_2<0$, $A<0$, $B<0$, $C<0$) 

%\label{fig2}

% }

%\smallskip \footnotesize

%Fig.~\ref{fig2}.

%A qualitative form of the solutions  $T_1,$ $T_2,$ $P_1,$ $P_2,$ $U,$ $V$ of system (\ref{eq:Aristov_Prosviryakov:sys1}) when $\frac{9}{10}<\sigma<\frac{3}{2}$ ($1-\tau_1>0$,  $\tau_2>0$, $A>0$, $B>0$, $C>0$; \centerline{$2-\tau_1<0$, $\tau_2<0$, $A<0$, $B<0$, $C<0$)} \label{fig2} 

%\end{minipage} \hfill

%\begin{minipage}[h]{0.5\textwidth}

%\centering

%\includegraphics[width=0.98\textwidth]{fig2}

%\end{minipage}

%\end{figure}

%

%

%

%\begin{figure}[h!]

%\noindent

%\begin{minipage}[h]{0.48\textwidth}

%\caption{A qualitative form of the solutions  $T_1,$ $T_2,$ $P_1,$ $P_2,$ $U,$ $V$ of system (\ref{eq:Aristov_Prosviryakov:sys1}) when $\sigma\leq\frac{9}{10}$ ($1-\tau_1>0$,  $\tau_2>0$, $B>0$, $C>0$; $2-\tau_1<0$, $\tau_2<0$, $B<0$, $C<0$) 

%\label{fig3}

% }

%\smallskip \footnotesize

%Fig.~\ref{fig3}.

%A qualitative form of the solutions  $T_1,$ $T_2,$ $P_1,$ $P_2,$ $U,$ $V$ of system (\ref{eq:Aristov_Prosviryakov:sys1}) when $\sigma\leq\frac{9}{10}$ ($1-\tau_1>0$,  $\tau_2>0$, $B>0$, $C>0$; 

%\centerline{$2-\tau_1<0$, $\tau_2<0$, $B<0$, $C<0$)}

% \label{fig3} 

%\end{minipage} \hfill

%\begin{minipage}[h]{0.5\textwidth}

%\centering

%\includegraphics[width=0.98\textwidth]{fig3}

%\end{minipage}

\end{figure}









\newpage















\Section{Solution of equation  with the Prandtl number different from one}

An important particular case of equation \eqref{eq:Aristov_Prosviryakov:trans4} was studied above; namely, when the values of kinematic viscosity and thermal diffusivity of the fluid coincide. 

Let us now write the solution of equation \eqref{eq:Aristov_Prosviryakov:trans4} in the general case:

\begin{multline}

\label{eq:Aristov_Prosviryakov:soltempPr}

\Theta_1=\text{\sf Ai}\Bigl(\xi+\frac{4\Omega z}{\sqrt{\nu\chi}}\sqrt[3]{\frac{\nu^2\chi^2} {16\Omega^2}}\Bigr)

\left(C_1\text{\sf Ai}\left(\xi\right)+C_2\text{\sf Bi}\left(\xi\right)\right)+

\\

+\text{\sf Bi}\Bigl(\xi+\frac{4\Omega z}{\sqrt{\nu\chi}}\sqrt[3]{\frac{\nu^2\chi^2}{16\Omega^2}}\Bigr)\left(C_3\text{\sf Ai}\left(\xi\right)+C_4\text{\sf Bi}\left(\xi\right)\right).

\end{multline}

Here $$\xi=\sqrt[3]{\nu^2\chi^2}\displaystyle\frac{ i\left(\nu-\chi\right)+2\sqrt{\nu\chi}\Omega z}{\nu\chi\sqrt[3]{16\Omega^2}}$$ is an auxiliary variable introduced in order to make the solution notation concise;  $\text{\sf Ai} \left(\xi\right)$   and $\text{\sf Bi}\left(\xi\right)$  are the Airy function and the associated function~\cite{ulg:bib:Polyanin}. Solution \eqref{eq:Aristov_Prosviryakov:soltempPr} is complex-valued, since it is written for arbitrary differences between $\nu$  and $\chi.$  This is done for the convenience of the solution analysis. Solution \eqref{eq:Aristov_Prosviryakov:soltempPr} being known, pressure  $S_1$ and velocity $W$ are computed similarly to the previous case.	



When the dissipative coefficients  $\nu$ and  $\chi$ of the nonisothermal fluid are equal, solution \eqref{eq:Aristov_Prosviryakov:soltempPr}  is reduced to formula \eqref{eq:Aristov_Prosviryakov:soltemp}. Studying the convergence of series \eqref{eq:Aristov_Prosviryakov:soltempPr}, we find that the solution expressed in terms of this series converges on the solution domain, i.\,e. layer thickness. 



The consideration of equation \eqref{eq:Aristov_Prosviryakov:trans4} with an arbitrary relation between $\nu$  and~$\chi$ offers solutions to system \eqref{eq:Aristov_Prosviryakov:sys1}, \eqref{eq:Aristov_Prosviryakov:sys2} in qualitative terms, see figs.~\ref{fig_prosv1} to~\ref{fig_prosv3}. In the simulation of the hydrodynamic fields by the linear combination of the Airy functions the behavior of the solutions is predictable from the very beginning, since these special functions satisfy the simplest differential equation having a point where the fluctuation of the solution is replaced by its exponential growth. It was shown above that the solution does not go off to infinity.



\smallskip



\Section{Conclusion}

The problem of simulating rotating fluid masses has a long history. To describe fluid motion (relative equilibrium), ordinary differential equa\-tions are used that enable one to obtain solutions simulating fluid flows by the theory of motion stability with the disturbance of the background (principal) flow. This paper discusses the exact solution to an overdetermined boundary value problem, which describes stationary flow at any scales where the Boussinesq approximation is valid. In this case, stationary dynamic fluid equilibria are chara\-cterized by the formation of counterflows, the sign of velocity being able to alternate several times, depending on the boundary conditions and the geometric dimensions of the fluid layer. A similar situation is true for temperature and pressure, namely, the formation of zones with the negative and positive values of temperature and pressure with respect to the reference value. 



\smallskip 

 



\AdditionalContent{Competing interests}{I declare that I have no competing interests.} 

\AdditionalContent{Author's Responsibilities}{I take full responsibility for submitting the final manuscript in print. I approved the final version of the manuscript.} 





\AdditionalContent{Funding}{The work was done within the state assignment from FASO Russia, theme No. AAAA-A18-118020790140-5.}









\begin{thebibliography}{28}



\Bibitem{ulg:bib:Newtoni}

\by Newton~I.

\book Opera quae exstant omnia

\bookinfo Faksimile-Neudruck der Ausgabe von Samuel Horsley, London 1779--1785 in f\"unf B\"anden. Band 1,~4

\lang In German

\zmath{0129.25102}

\publaddr Stuttgart-Bad Cannstatt

\publ Friedrich Frommann Verlag (G\"unther Holzboog)

\totalpages xx+592

\yr 1964



\Bibitem{ulg:bib:Turnbull}

\book The Correspondence of Isaac Newton

\bookvol 1. 1661-1675

\ed H.~W.~Turnbull

\publ Cambridge Univ. Press

\publaddr New York

\yr 1959

\totalpages xxxviii+468 





\Bibitem{ulg:bib:Poincare}

\by Poincar\'e~H.

\book Cin\'ematique et m\'ecanismes. Potentiel et m\'ecanique des fluides

\bookinfo Cours profess\'e \`a la Sorbonne

\lang In French

\zmath{30.0608.01}

\publaddr Paris

\publ G.~Carr\'e et C.~Naud

\totalpages iv+385 

\yr 1899



\Bibitem{ulg:bib:Lyapunov1}

\by Lyapunov~A.~M.

\paper On the stability of ellipsoidal equilibrium forms of rotating fluid

\inbook Collected works. \rm  Vol. III

\yr 1959

\publ Akad. Nauk SSSR

\publaddr Moscow

\pages 5--113

\zmath{0192.31203}

\lang In Russian



\Bibitem{ulg:bib:Lyapunov2}

\by Lyapunov~A.~M.

\paper On the equilibrium figures of rotating homogeneous liquid mass slightly different from ellipsoids

\inbook Collected works. \rm  Vol. IV

\yr 1959

\publ Akad. Nauk SSSR

\publaddr Moscow

\pages 5--645

\zmath{0192.31204}

\lang In Russian



\Bibitem{ulg:bib:Chandrasekhar}

\by Chandrasekhar~S.

\book Ellipsoidal figures of equilibrium

\zmath{0213.52304}

\publaddr New Haven, London

\publ Yale University Press 

\totalpages ix+252 

\yr 1969



\Bibitem{ulg:bib:Hagihara}

\by Hagihara~Y.

\book Theories of equilibrium figures of a rotating homogeneous fluid mass

\yr 1970

\serial NASA Special Publication 

\vol 186 

\publaddr Washington

\publ US Government Printing Office

\adsnasa{1970NASSP.186.....H}



\Bibitem{ulg:bib:Borisov}

\by Borisov~A.~V., Kilin~A.~A., Mamaev~I.~S.

\paper The Hamiltonian Dynamics of Self-gravitating Liquid and Gas Ellipsoid

\jour Regul. Chaotic Dyn. 

\yr 2009

\vol 14

\issue 2

\pages 179--217

\zmath{https://zbmath.org/?q=an:1229.70049}

\crossref{10.1134/S1560354709020014}



\Bibitem{ulg:bib:Appell}

\by Appell~P.

\book Trait\'e de m\'ecanique rationnelle

\bookinfo Tome IV, 1: Figures d'\'equilibre d'une masse liquide homog\`ene en rotation

\lang In French

\zmath{58.1300.02}

\totalpages viii+342 

\publaddr  Paris

\publ Gauthier-Villars 

\yr 1932







\Bibitem{ulg:bib:Betti}

\by Betti~E.

\paper Sopra i moti che conservano la figura ellissoidale a una massa fluida eterogenea

\lang In Italian

\zmath{13.0718.01}

\jour Brioschi Ann

\issueinfo  (2)~X

\pages 173--187

\yr 1879







\Bibitem{ulg:bib:Volterra}

\by Volterra~V.

\paper Sur la stratification d\'une masse fluide en \'equilibre

\jour Acta Math.

\yr 1903

\vol 27

\issue 1

\pages 105--124



\Bibitem{ulg:bib:Liouville}

\by Liouville~J.

\paper Sur la figure d\'une masse fluide homog\'ene, en \'equilibre et dou\'ee d\'un mouvement de rotation

\jour J. de l'\'Ecole Polytech.

\yr 1834

\vol 14

\pages 289–296



\Bibitem{ulg:bib:Dirichlet}

\by Dirichlet~G.~L.

\paper Untersuchungen \"uber ein Problem der Hydrodynamik (Aus dessen Nachlass hergestellt von Herrn R. Dedekind zu Z\"urich)

\jour J. Reine Angew. Math. (Crelle's Journal)

\yr 1861

\vol 58

\pages181--216



\Bibitem{ulg:bib:Arnold}

\by Arnold~V. I., Khesin~B.~A.

\book Topological Methods in Hydrodynamics

\yr 1999

\publ Springer

\publaddr New York

\totalpages 392



\Bibitem{ulg:bib:Dolzhansky}

\by Dolzhansky~F.~V.

\paper On the mechanical prototypes of fundamental hydrodynamic invariants and slow manifolds

\jour Physics–Uspekhi

\yr 2005

\vol 48

\issue 12

\pages 1205--1234

\crossref{10.1070/PU2005v048n12ABEH002375}



\Bibitem{ulg:bib:Hadamard}

\by Hadamard~J.

\paper Mouvement permanent lent d'une sph\`ere liquide et visqueuse dans un liquide visqueux

\zmath{42.0813.01}

\jour C.~R.~Acad. Sci., Paris 

\vol 152

\issue 25

\pages 1735--1738 

\yr 1911





\Bibitem{ulg:bib:Rybczynski}

\by Rybczi\'nski, W.

\paper \"Uber die fortschreitende Bewegung einer fl\"ussigen Kugel in einem z\"ahen Medium

\zmath{42.0811.02}

\jour Bull. Int. Acad. Sci. Cracovie. Ser. A

\vol 1

\pages 40--46 

\yr 1911



\Bibitem{ulg:bib:Aristov1}

\by Aristov~S.~N., Prosviryakov~E.~Yu. 

\paper A New Class of Exact Solutions for Three Dimensional Thermal Diffusion Equations

\jour Theor. Found. Chem. Eng.

\yr 2016

\vol 50

\issue 3

\pages 286--293

\crossref{10.1134/S0040579516030027}



\Bibitem{ulg:bib:Aristov2}

\by Aristov~S.~N., Knyazev~D.~V., Polyanin~A.~D.

\paper Exact solutions of the Navier-Stokes equations with the linear dependence of velocity components on two space variables

\jour Theor. Found. Chem. Eng.

\crossref{10.1134/S0040579509050066}

\yr 2009

\vol 43

\issue 5

\pages 642--662



\Bibitem{ulg:bib:Polyanin}

\by Polyanin~A.~D., Zaitsev~V.~F.

\book Handbook of ordinary differential equations. Exact solutions, methods, and problems

\zmath{06776009}

\publaddr Boca Raton, FL

\publ CRC Press 

\totalpages xxix+1456 

\yr 2018



\Bibitem{ulg:bib:Aristov3}

\by Aristov~S.~N., Prosviryakov~E.~Yu.

\paper Unsteady Layered Vortical Fluid Flows

\jour Fluid Dyn.

\crossref{10.1134/S0015462816020034}

\yr 2016

\vol 51

\issue 2

\pages 148--154



\Bibitem{ulg:bib:Aristov4}

\by Aristov~S.~N., Prosviryakov~E.~Yu., Spevak~L.~F.

\paper Unsteady-State B\'enard--Marangoni Convection  in Layered Viscous Incompressible Flows

\jour Theor. Found. Chem. Eng.

\crossref{10.1134/S0040579516020019}

\yr 2016

\vol 50

\issue 2

\pages 132--141



\Bibitem{ulg:bib:Aristov5}

\by Aristov~S.~N., Prosviryakov~E.~Yu.

\paper On one class of analytic solutions of the stationary axisymmetric convection B\'enard--Marangoni viscous incompressible fluid

\jour Vestn. Samar. Gos. Tekhn. Univ., Ser. Fiz.-Mat. Nauki \rm [J. Samara State Tech. Univ., Ser. Phys. Math. Sci.]

\yr 2013

\vol 3(32)

\pages 110--118

\mathnet{http://mi.mathnet.ru/vsgtu1205}

\crossref{10.14498/vsgtu1205}

\zmath{https://zbmath.org/?q=an:06968790}

\elib{http://elibrary.ru/item.asp?id=20710187}

\lang In Russian



\Bibitem{ulg:bib:Aristov6}

\by Aristov~S.~N., Privalova~V.~V., Prosviryakov~E.~Yu.

\paper Stationary nonisothermal Couette flow. Quadratic heating of the upper boundary of the fluid layer

\jour Nelin. Dinam.

\yr 2016

\vol 12

\issue 2

\pages 167--178

\mathnet{http://mi.mathnet.ru/nd519}

\elib{http://elibrary.ru/item.asp?id=26193555}

\crossref{10.20537/nd1602001}

\lang In Russian



\Bibitem{ulg:bib:Burmasheva1}

\by Burmasheva~N.~V., Prosviryakov~E.~Yu.

\paper A large-scale layered stationary convection of an incompressible viscous fluid under the action of shear stresses at the upper boundary. Velocity field investigation

\jour Vestn. Samar. Gos. Tekhn. Univ., Ser. Fiz.-Mat. Nauki \rm [J. Samara State Tech. Univ., Ser. Phys. Math. Sci.]

\yr 2017

\vol 21

\issue 1

\pages 180--196

\mathnet{http://mi.mathnet.ru/vsgtu1527}

\crossref{10.14498/vsgtu1527}

\elib{http://elibrary.ru/item.asp?id=29245104}

\lang In Russian



\Bibitem{ulg:bib:Burmasheva2}

\by Burmasheva~N.~V., Prosviryakov~E.~Yu.

\paper A large-scale layered stationary convection of a incompressible viscous fluid under the action of shear stresses at the upper boundary. Temperature and presure field investigation

\jour Vestn. Samar. Gos. Tekhn. Univ., Ser. Fiz.-Mat. Nauki \rm [J. Samara State Tech. Univ., Ser. Phys. Math. Sci.]

\yr 2017

\vol 21

\issue 4

\pages 736--751

\mathnet{http://mi.mathnet.ru/vsgtu1568}

\crossref{10.14498/vsgtu1568}

\zmath{https://zbmath.org/?q=an:06964885}

\elib{http://elibrary.ru/item.asp?id=32712835}

\lang In Russian



\Bibitem{ulg:bib:Gorshkov}

\by Gorshkov~A.~V., Prosviryakov~E.~Yu.

\paper Ekman Convective Layer Flow of a Viscous Incompressible Fluid

\jour Izv. Atmos. Ocean. Phys.

\yr 2018

\vol 54

\issue 2

\pages 89–195

\crossref{10.1134/S0001433818020081}





\Bibitem{ulg:bib:Aristov7}

\by Aristov~S.~N., Knyazev~D.~V.

\paper Localized convective flows in a nonuniformly heated liquid layer

\jour Fluid Dyn.

\yr 2014

\vol 49

\issue 5

\pages 565--575

\crossref{10.1134/S0015462814050020}











\end{thebibliography}





\renewcommand{\baselinestretch}{0.95}



\newpage



\selectlanguage{russian}



$\;$







\makerutitle



\vspace{-7mm}



\AdditionalContent{Конкурирующие интересы}{Конкурирующих интересов не имею.}



\AdditionalContent{Авторский вклад и ответственность}{Я несу полную ответственность за предоставление окончательной версии рукописи в печать. Окончательная версия 

рукописи мною одобрена.}

\AdditionalContent{Финансирование}{Работа выполнена в рамках государственного задания ФАНО, тема № АААА-А18-118020790140-5.}





\renewcommand{\baselinestretch}{0.9}







\renewcommand{\RuCitation}{\noindent\textbf{\CYRO\cyrb\cyrr\cyra\cyrz\cyre\cyrc\ \cyrd\cyrl\cyrya\ \cyrc\cyri\cyrt\cyri\cyrr\cyro\cyrv\cyra\cyrn\cyri\cyrya\\} {\RuAuthorInContents} {\RuTitleInContents}~/\!/ \textit{\RuJournalName}, \JournalYear.  \CYRT.~\JournalVolume, \textnumero~\JournalNumber.  \CYRS.~\thefirstpage--\thelastpage. \doiurl.}



\renewcommand{\EnCitation}{\noindent\textbf{Please cite this article in press as:\\} {\EnAuthorInContents} {\EnTitleInContents}, \textit{\EnJournalName}\/ [\ParallelJournalName],  \JournalYear, vol.~\JournalVolume, no.~\JournalNumber,  pp.~\thefirstpage--\thelastpage.  \midoiurl\ (In Russian).}









%\end{document}

 \newpage

%\input{preamble}

%

%\usepackage{samgtu-name}

%

%\usepackage{slashbox}



\vsgtu{1653} %registration number

\FirstPage{750}

\LastPage{761}

%

%\pdfcrop

%\draft



\ShortCommunication



%\begin{document}





\hfill \qrcode[hyperlink,height=.8in]{http://doi.org/10.14498/vsgtu1653}

\vspace{-.8in}





\udc{517.958:539.3(1)}

\msc{74F05, 74K20}



\rutitle{Динамическая устойчивость геометрически нерегулярной нагретой пологой цилиндрической оболочки в~сверхзвуковом потоке газа}

\rutitleshort{Динамическая устойчивость геометрически нерегулярной\dots}



\entitle{Dynamic stability of heated geometrically irregular cylindrical shell in supersonic gas flow}

\entitleshort{Dynamic stability of heated geometrically irregular cylindrical shell in supersonic gas flow}





\ruauthor{Г.~Н.~Белосточный, О.~А.~Мыльцина}{Б\,е\,л\,о\,с\,т\,о\,ч\,н\,ы\,й~Г.~Н.,\ М\,ы\,л\,ь\,ц\,и\,н\,а~О.~А.}

\enauthor{G.~N.~Belostochnyi, O.~A.~Myltcina}{B\,e\,l\,o\,s\,t\,o\,c\,h\,n\,y\,i~G.~N.,\ M\,y\,l\,t\,c\,i\,n\,a~O.~A.}



\ruaffil{\saratovuni.}



\enaffil{\engsaratovuni.}



\ruauthorsdetails{3}{% Количество авторов



\fullauthor{\rupid{98285}{Григорий Николаевич Белосточный}}

%\CAuthor % Автор-корреспондент

\orcid{0000-0003-4471-6599} % ORCID ID автора 

\regalia{доктор технических наук, профессор} 

\position{профессор} 

\dept{каф. математической теории упругости и биомеханики} 

\email{belostochny@mail.ru}

\fullauthor{\rupid{98283}{Ольга Анатольевна Мыльцина}}

\CAuthor % Автор-корреспондент

\orcid{0000-0003-4718-2772} % ORCID ID автора 

\regalia{кандидат физико-математических наук} 

\position{ассистент} 

\dept{каф. теории функций и стохастического анализа} 

\email{omyltcina@yandex.ru}

}



%Olga Myltcina https://orcid.org/0000-0003-4718-2772 

%Grigory Belostochnyihttps://orcid.org/0000-0001-6949-4174

%Acel Polienko https://orcid.org/0000-0003-4471-6599







\enauthorsdetails{3}{% Количество авторов

\fullauthor{\rupid{98285}{Grigory N. Belostochnyi}}

%\CAuthor % Автор-корреспондент

\orcid{0000-0003-4471-6599} % ORCID ID автора 

\regalia{Dr. Techn. Sci.} 

\position{Professor} 

\dept{Dept. of Mathematic Theory of Elasticity \&  Biomechanics} 

\email[n]{belostochny@mail.ru} 

\fullauthor{\rupid{98283}{Olga A. Myltcina}}

\CAuthor % Автор-корреспондент

\orcid{0000-0003-4718-2772} % ORCID ID автора 

\regalia{Cand. Phys. \& Math. Sci.} 

\position{Assistant} 

\dept{Dept. of Functions \& Approxmation Theory} 

\email[n]{omyltcina@yandex.ru}

}





\rukeywords{динамическая устойчивость, температура, пологие оболочки, сверхзвук, континуальность, обобщенные функции, поршневая теория, аэродинамика, критические скорости, ребра жесткости, демпфирование, кривизна, критерий Гурвица, изотропия}

\enkeywords{dynamic stability, temperature, flat shells, supersonic, continuity, generalized functions, piston theory, aerodynamics, critical velocities, ribs, damping, curvature, Routh–Hurwitz stability criterion, isotropy}



\ruabstract{На базе модели типа Лява рассматривается нагретая до постоянной температуры геометрически нерегулярная пологая цилиндрическая оболочка, обдуваемая сверхзвуковым потоком газа со стороны одной из ее основных поверхностей. За основу взята континуальная модель термоупругой системы в~виде тонкостенной оболочки, подкрепленной ребрами вдоль набегающего газового потока. Сингулярная система уравнений динамической термоустойчивости геометрически нерегулярной оболочки содержит слагаемые, учитывающие «растяжение-сжатие» и~сдвиг подкрепляющих элементов в~тангенциальной плоскости, тангенциальные усилия, вызванные нагревом оболочки, и~поперечную нагрузку, стандартным образом записанную по «поршневой теории».



Тангенциальные усилия предварительно определяются как решение сингулярных дифференциальных уравнений безмоментной термоупругости геометрически нерегулярной оболочки с~учетом краевых усилий.



Решение системы динамических уравнений термоупругости оболочки разыскивается в~виде суммы двойного тригонометрического ряда (для функции прогиба) с~переменными по временной координате коэффициентами. На основании метода Галеркина получена однородная система для коэффициентов аппроксимирующего ряда, которая сведена к~одному дифференциальному уравнению четвертого порядка. Решение приводится во втором приближении, что соответствует двум полуволнам в~направлении потока и~одной полуволне в~перпендикулярном направлении. На основании стандартных методов анализа динамической устойчивости тонкостенных конструкций определяются критические значения скоростей газового потока.



Количественные результаты приводятся в виде таблиц, иллюстрирующих влияние геометрических параметров термоупругой системы «обо\-лоч\-ка-ребра», температуры  на устойчивость геометрически нерегулярной цилиндрической оболочки в~сверхзвуковом потоке газа с~учетом демпфирования.}



\enabstract{On the basis of the Love model, a geometrically irregular heated cylindrical shell blown by a supersonic gas flow from one of its main surfaces is considered. The continuum model of a thermoelastic system in the form of a thin-walled shell supported by ribs along the incoming gas flow is taken as a basis. The singular system of equations for the dynamic thermal stability of a geometrically irregular shell contains terms that take into account the tension-compression and the shift of the reinforcing elements in the tangential plane, the tangential forces caused by the heating of the shell and the transverse load, as standard recorded by the piston theory. The solution of a singular system of differential equations in displacements, in the second approximation for the deflection function, is sought in the form of a double trigonometric series with time coordinate variables. 





Tangential forces are predefined as the solution of singular differential equations of non-moment thermoelasticity of a geometrically irregular shell taking into account boundary forces. 



The solution of the system of dynamic equations of thermoelasticity of the shell is sought in the form of the sum of the double trigonometric series (for the deflection function) with time coordinate variable coefficients. On the basis of the Galerkin method, a homogeneous system for the coefficients of the approximating series is obtained, which is reduced to one fourth-order differential equation. The solution is given in the second approximation, which corresponds to two half-waves in the direction of flow and one half-wave in the perpendicular direction. On the basis of standard methods of analysis of dynamic stability of thin-walled structures are determined critical values of the gas flow rate. 



The quantitative results are presented in the form of tables illustrating the influence of the geometrical parameters of the thermoelastic shell-edge system, temperature and damping on the stability of a geometrically irregular cylindrical shell in a supersonic gas flow.}



\newdate{myla}{11}{10}{2018}

\date{\displaydate{myla}}

\newdate{mylb}{10}{11}{2018}

\revisiondate{\displaydate{mylb}}

\newdate{mylc}{12}{11}{2018}

\accepteddate{\displaydate{mylc}}



\newdate{myle}{21}{12}{2018}

\renewcommand{\JournalDataOnline}{\displaydate{myle}}



\makerutitle







Геометрически нерегулярные тонкостенные упругие системы, обширный класс которых составляют ребристые оболочки и~пластинки, широко используются в~различных областях современной техники. Условия эксплуатации таких систем предусматривают совместное воздействие температурных полей и~высокоскоростных газовых потоков. Исследованию упругого поведения гладких пластин и~оболочек на основе атермической

теории посвящено большое число работ, полный перечень которых содержал бы десятки

наименований. Ограничимся некоторыми из них \cite{myl-bib:1, myl-bib:2, myl-bib:3,myl-bib:4,myl-bib:5,myl-bib:5a}.

Значительно меньше работ

посвящено  исследованиям совместного воздействия температурных факторов и сверхзвукового потока

на тонкостенные системы. Важные для практики результаты в этой области содержатся в работах \cite{myl-bib:6, myl-bib:7, myl-bib:8}.



%Работы, в которых анализируется влияние подкрепляющих  оболочку ребер под действием сверхзвукового потока на базе термической теории в открытой научной литературе отсутствуют. 

Работы, в которых анализируется влияние подкрепляющих  ребер на поведение оболочки, находящейся под действием сверхзвукового потока, на базе термической теории, в~открытой научной литературе отсутствуют. 

Это связано не с маловажностью проблемы, а прежде всего с чрезвычайной математической сложностью таких задач, решаемых на основе дискретной модели «оболочка–ребра». Использование континуальной модели типа «конструктивная анизотропия» мало пригодно для практических целей, так как количественные результаты, полученные на ее основе, далеки от физической реальности.



Обращение к  континуальным моделям на основе теории обобщенных функций, основные положения которых содержится в работах \cite{myl-bib:9, myl-bib:10, myl-bib:11, myl-bib:12, myl-bib:15,myl-bib:16,myl-bib:18}, позволило сводить решения задач статической и динамической термоупругости ребристых  оболочек к интегрированию систем сингулярных дифференциальных уравнений точными и~приближенными методами высшего анализа \cite{myl-bib:19, myl-bib:21}, что немаловажно для инженерной практики.



\smallskip



{\bf 1.} Система сингулярных дифференциальных уравнений динамической термоупругости пологой геометрически нерегулярной цилиндрической оболочки, стандартным образом отнесенной к декартовым координатам на основе континуальной модели, в которых учет тангенциальных усилий записывается в форме Рейснера \cite{myl-bib:22a}, в компонентах поля перемещений примет вид

\begin{gather}

u_{\bigcom 11}+\frac{1-\nu}{2}u_{\bigcom 22}+\frac{1+\nu}{2} v_{\bigcom 12}-k_{11}w_{\bigcom 1}+

\varepsilon_1\frac{1-\nu}{2} \sum_{i=1}^n \frac{h_i}{h}(u_{\bigcom 2}+v_{\bigcom 1} )_{\bigcom 2}\delta(x-x_i)=0,

\\

\frac{1+\nu}{2}u_{\bigcom 21}+ v_{\bigcom 22}+\frac{1-\nu}{2}v_{\bigcom 11}-\nu k_{11}w_{\bigcom 2}+\varepsilon_2\sum_{i=1}^n\frac{h_i}{h}a_i ( \nu u_{\bigcom 1}+v_{\bigcom 2})_{\bigcom 2} \delta(x-x_i)=0,

\end{gather}



\vspace{-12mm}



\begin{multline}\label{bel_eq1}

\nabla^2\nabla^2 w+\frac{B}{D}k^2_{11}w-k_{11}\frac{B}{D} ( u_{\bigcom 1}+\nu v_{\bigcom 2})- 

\\

- \frac{1}{D} \bigl[ 

( T^{11}_0 w_{\bigcom 1})_{\bigcom 1}+

( T^{12}_0 w_{\bigcom 1})_{\bigcom 2}+  

(T^{12}_0 w_{\bigcom 2})_{\bigcom 1}+ 

(T^{22}_0 w_{\bigcom 2})_{\bigcom 2}\bigr]-

\\

-

\frac{1}{D}\sum_{i=1}^n\frac{h_i}{h}a_i

(T^{22}_0 w_{\bigcom 2})\delta(x-x_i)+ 

\\

+\sum_{i=1}^n \Bigl( \frac{h_i}{h}\Bigr)^3 \Phi_{3i}a_iw_{\bigcom 2222}\delta(x-x_i)+

\\

+

2(1-\nu)\sum_{i=1}^n\Bigl( \frac{h_i}{h}\Bigr)^3 \Phi_{3i}a_iw_{\bigcom 122}\bigr|_{x_i}\frac{d\delta(x-x_i)}{dx}= \\

=\frac{\gamma h}{g D}w_{\bigcom tt}\biggl[ 1+\sum_{i=1}^n \frac{h_i}{h}a_i\delta(x-x_i)\biggr] 

-\frac{\mu} {h D}w_{\bigcom t}-p_0\frac{\varkappa \mu}{D}w_{\bigcom 2}, 

\end{multline}

где $p_0\frac{\varkappa \mu}{D}w_{\bigcom 2}$ "---  относительная интенсивность поперечной нагрузки, вызванная прогибом пластинки, стандартным образом записана согласно поршневой теории \cite{myl-bib:1,myl-bib:22, myl-bib:30}; $M={v_y}/{c_0}$ "--- число Маха; 

$v_y$ "--- невозмущенная скорость набегающего газового потока; 

$c_0$ "--- скорость звука на бесконечности; $u$, $v$, $w$ "---  компоненты поля (векторы) перемещений; 

${h_i}/{h}$ "--- относительная высота $i$-того ребра, число которых $n$; 

$a_i$ "--- ширина $i$-того ребра; 

$k_{11}$ "--- кривизна цилиндрической оболочки; 

$D=\frac{E h^3}{12(1-\nu)}$; $\nu$ "--- коэффициент Пуассона; 

$E$ "--- модуль Юнга, 

$\gamma$ "--- удельный вес; 

$g$ "--- интенсивность поля тяжести; 

$$\Phi_{3i}=1+3\frac{h_i}{h}+\Bigl( \frac{h_i}{h}\Bigr)^3 ;$$ 

$\mu$ "--- коэффициент демпфирования; 

$\delta(x-x_i)$ "--- обобщенная дельта"=функция Дирака~\cite{myl-bib:25};  

$\varepsilon_j$ $(j=0,1)$ "--- знаковые числа, равные 1 или 0; 

$ T^{ij}_0 $ "--- тангенциальные усилия, вызванные нагревом пластинки до постоянной температуры $\Theta_0$, которые определяются на основании решений системы сингулярных однородных дифференциальных уравнений, описывающих безмоментное состояние термоупругой системы в компонентах поля перемещений $\bar{u}^0(u^0,v^0,w^0)$:

 \begin{multline}

 u^0_{\bigcom 11}+\frac{1-\nu}{2}u^0_{\bigcom 22}+\frac{1+\nu}{2}v^0_{\bigcom 12}-k_{11}w^0_{\bigcom 1}+

\\

 +\frac{1-\nu}{2}\sum_{i=1}^{n}\frac{h_i}{h}a_i ( u^0_{\bigcom 2}+v^0_{\bigcom 1})_{\bigcom 2}\delta(x-x_i) =0,

 \end{multline}

 \begin{multline}

 \frac{1+\nu}{2}u^0_{\bigcom 12}+v^0,_{22}+\frac{1-\nu}{2}v^0_{\bigcom 11}-\nu k_{11}w^0_{\bigcom 2}+

\\

+ 

 \sum_{i=1}^{n}\frac{h_i}{h}a_i( v^0_{\bigcom 2}+\nu u^0_{\bigcom 1})_{\bigcom 2}\delta(x-x_i)=0,

 \end{multline}

\begin{equation}\label{eq:bel:momentnot}

k_{11}^2 w^0+k_{11} (u^0_{\bigcom 1}+\nu v^0_{\bigcom 2} ) =\alpha(1+\nu)k_{11}\Theta_0.

\end{equation}





При неоднородных краевых условиях

\begin{equation*}

\begin{array}{ll}

\text{$x=0$, $x=a$:}&

u^0_{\bigcom 2}+v^0_{\bigcom 1}=0, \quad u^0_{\bigcom 1}+\nu v^0_{\bigcom 2}-k_{11}w^0=\alpha(1+\nu \Theta_0), \quad w^0=0;

\\[2mm]

\text{$y=0$, $y=b$:}&

v^0=0, \quad u^0_{\bigcom 2}+ v^0_{\bigcom 1}=0, \quad w^0=0

\end{array}

\end{equation*}

решения сингулярной системы \eqref{eq:bel:momentnot} запишутся в элементарных функциях:

\begin{equation}\label{bel_gr1}

v^0=0, \quad u^0=\alpha(1+\nu)\Theta_0 x, \quad w^0=0,

\end{equation}

а тангенциальные усилия примут вид

\begin{equation*}

T^{11}_0=T^{12}_0=0, \quad T^{22}_0=-E h\alpha\Theta_0.

\end{equation*}



Решение системы \eqref{bel_eq1} (в~которой отсутствуют инерционные слагаемые в~тангенциальной плоскости \cite{myl-bib:1,myl-bib:6, myl-bib:7}), тождественно удовлетворяющее условиям шарнирного закрепления всех сторон пологой цилиндрической оболочки

\begin{equation*}

\begin{array}{ll}

\text{при $x=0$, $x=a$:} &

T^{12}=0, \quad T^{11}=0, \quad w=0 \quad M^{11}=0;

\\[2mm]

\text{при $y=0$, $y=b$:} &

v=0, \quad T^{12}=0, \quad w=0 \quad M^{22}=0,

\end{array}

\end{equation*}

зададим в виде

\begin{gather}\label{bel_appr1}

u(x,y,t)=\tilde{u}(t)u_1 ({x}/{a} )u_2 ({y}/{b} ), \quad 

v(x,y,t)=\tilde{v}(t)v_1({x}/{a} )v_2 ({y}/{b} ),

\\

w(x,y,t)=w_{11}(t)\sin \frac{\pi x}{a} \sin\frac{\pi y}{b} +w_{12}(t)\sin \frac{\pi x}{a} \sin\frac{2\pi y}{b},

\end{gather}

где 

\begin{gather}

u_1({x}/{a}) =( {x}/{a})^4-2( {x}/{a})^3+( {x}/{a})^2,\quad

u_2({y}/{b}) =-({2}/{3})( {y}/{b})^3+( {y/}{b})^2,\\ 

v_1({x}/{a}) =({x}/{a})^4-2( {x}/{a})^3+( {x}/{a})^2,

\quad 

v_2({y}/{b}) =({y}/{b})( {y}/{b}-1).

\end{gather}



Первые два уравнения системы \eqref{bel_eq1} на основании подстановок \eqref{bel_appr1} с последующим применением процедуры Галеркина \cite{myl-bib:27,  myl-bib:29}  перепишутся в виде 

\begin{gather}

e_{11}\tilde{u}(t)+e_{12}\tilde{v}(t)=k_{11}a\Bigl[ \pi \frac{b}{a}\mathit{I}_4 w_{11}(t)+\pi\frac{b}{a}\mathit{I}_5 w_{12}(t) \Bigr],

\\

\label{bel_eq7}

e_{21}\tilde{u}(t)+e_{22}\tilde{v}(t)=\nu k_{11} a \bigl[ \pi \mathit{I}_9 w_{11}(t)+2 \pi \mathit{I}_{10} w_{12}(t)\bigr].

\end{gather}

Здесь введены следующие обозначения:

\begin{gather}

e_{11}=\frac{b}{a}\mathit{I}_1+\frac{1-\nu}{2}\frac{a}{b}\mathit{I}_2, \quad 

e_{12}=\frac{1+\nu}{2}\mathit{I}_3,\\

e_{21}=\frac{1+\nu}{2}\mathit{I}_6,\quad 

e_{22}=\frac{a}{b}\mathit{I}_7+\frac{1-\nu}{2}\frac{b}{a}\mathit{I}_8;

\\

\mathit{I}_1=\int_0^1 \int_0^1 \tilde{u}_{1 \bigcom 11} \tilde{u}_1 \tilde{u}_2^2dX dY,\quad 

\mathit{I}_2=\int_0^1 \int_0^1 \tilde{u}_{2  \bigcom 22}\tilde{u}_2\tilde{u}_1^2dX dY,

\\

\mathit{I}_3=\int_0^1 \int_0^1 \tilde{v}_{1\bigcom 1}\tilde{u}_1\tilde{u}_2 \tilde{v}_2,_{2}dX dY, \dots,

\\

\mathit{I}_{10}=\int_0^1 \int_0^1 \tilde{v}_1 \tilde{v}_2 \sin (\pi X )\cos (2\pi Y )dX dY,\quad  X={x}/{a}, \quad  Y={y}/{b}.

\end{gather}



Разрешая неоднородную алгебраическую систему относительно коэффициентов $\tilde{u}(t)$ и $\tilde{v}(t)$, получим

\begin{equation}\label{bel_eq9}

\tilde{u}=k_{11}a\bigl(L_1 w_{11}(t)+L_2 w_{12}(t) \bigr),\quad 

\tilde{v}=k_{11}a\bigl(L_3 w_{11}(t)+L_4 w_{12}(t) \bigr),

\end{equation}

где $L_j$ "--- безразмерные величины:

\begin{gather}

L_1=\frac{\pi \frac{b}{a}\mathit{I}_4e_{22}-\nu \pi \mathit{I}_9 e_{12} }{e_{11} e_{22}-e_{12}e_{21}},\quad L_2=\frac{\pi \frac{b}{a}\mathit{I}_5e_{22}-\nu 2 \pi \mathit{I}_{10} e_{12} }{e_{11} e_{22}-e_{12}e_{21}}, 

\\

L_3=\frac{-\pi \frac{b}{a}\mathit{I}_4e_{21}+\nu \pi \mathit{I}_9 e_{11} }{e_{11} e_{22}-e_{12}e_{21}},\quad L_4=\frac{-\pi \frac{b}{a}\mathit{I}_5e_{21}+\nu 2 \pi \mathit{I}_{10} e_{11} }{e_{11} e_{22}-e_{12}e_{21}}. 

\end{gather}

Окончательно  функции $u(x,y,t)$ и $v(x,y,t)$ примут следующий вид:

\begin{gather}\label{bel_eq11}

u(x,y,t)=k_{11} a \bigl(L_1 w_{11}(t)+L_2 w_{12}(t) \bigr) u_1({x}/{a} )u_2({y}/{b}),

\\

v(x,y,t)=k_{11} a \bigl(L_3 w_{11}(t)+L_4 w_{12}(t) \bigr) v_1({x}/{a} )v_2({y}/{b} ).

\end{gather}



Из третьего уравнения системы \eqref{bel_eq1}  на основании аналогичных преобразований получим систему обыкновенных однородных дифференциальных уравнений относительно коэффициентов $w_{11}(t)$ и $w_{12}(t)$:

\begin{equation}\label{bel_eq12}

\xi_{11}\left[w_{11} \right]+\xi_{12} w_{12}=0, \quad  \xi_{21}w_{11}+\xi_{22} \left[ w_{12}\right] =0,

\end{equation}

где $\xi_{kl}$ "--- дифференциальные операторы:

\begin{gather}

\xi_{11}=\frac{\gamma h a^4}{g D}\biggl( 1+2 \sum\limits_{i=1}^n \tilde{\beta}_i^s\biggr)

\frac{d^2}{dt^2}+\tilde{\mu} \frac{d}{dt}+ f_{11}\Bigl( k_{11}, \Theta_0, \frac{h_i}{h}\Bigr), 

\\

\xi_{12}=\frac{p_0 \varkappa M a^3}{D} 4\pi \frac{a}{b} \mathit{I}_{13}-\tilde{f}_{12}, \quad \xi_{21}=\frac{p_0 \varkappa M a^3}{D} 2\pi \frac{a}{b} \mathit{I}_{16}-\tilde{f}_{11},

\\

\xi_{22}=\frac{\gamma h a^4}{g D}

\biggl( 1+2 \sum\limits_{i=1}^n \tilde{\beta}_i^s\biggr)\frac{d^2}{dt^2}+\tilde{\mu} \frac{d}{dt}+ f_{12}

\Bigl( k_{11},\Theta_0,\frac{h_i}{h}\Bigr), 

\end{gather}

в которых

\begin{multline}

f_{11}=\Bigl(\pi^2+\Bigl( \frac{\pi a}{b}\Bigr) ^2 \Bigr) ^2+

12

\Bigl( \frac{h_i}{h}\Bigr)^2(k_{11}a)^2-

48

\Bigl( \frac{a}{h}\Bigr)^2 (k_{11}a)^2\Bigl(L_1\mathit{I}_{11}+\nu \frac{a}{b}L_3 \mathit{I}_{12} \Bigr) -

\\

-12(1-\nu^2)\Bigl( \frac{\pi a}{b}\Bigr)^2 \Bigl(\frac{a}{h} \Bigr)^2

\biggl(1+\sum\limits_{i=1}^n \tilde{\beta}_i^s \biggr) \alpha \Theta_0+

\Bigl( \frac{\pi a}{b}\Bigr)^4 2\sum\limits_{i=1}^n \beta_i^s+

\\

+4(1-\nu)\pi^2\Bigl(\frac{\pi a}{b} \Bigr)^2  \sum\limits_{i=1}^n \beta_i^c,

\end{multline}

\begin{equation*}

\tilde{f}_{12}=48 \Bigl(\frac{a}{h} \Bigr)^2(k_{11}a)^2\Bigl(L_2 \mathit{I}_{11}+\nu\frac{a}{b}L_4\mathit{I}_{12}\Bigr),

\end{equation*}

\begin{equation*}

\tilde{f}_{11}=48 \Bigl(\frac{a}{h} \Bigr)^2(k_{11}a)^2\Bigl(L_1 \mathit{I}_{14}+\nu \frac{a}{b}L_3\mathit{I}_{15}\Bigr),

\end{equation*}

\begin{multline}

f_{12}=\Bigl(\pi^2+\Bigl( \frac{2\pi a}{b}\Bigr) ^2 \Bigr)^2-

12\Bigl( \frac{a}{h}\Bigr)^2(k_{11}a)^2-

48\Bigl( \frac{a}{h}\Bigr)^2 (k_{11}a)^2 \Bigl(L_2\mathit{I}_{14}+\nu \frac{a}{b}L_4 \mathit{I}_{15} \Bigr) -

\\

-12(1-\nu^2)\Bigl( \frac{2\pi a}{b}\Bigr)^2

\Bigl(\frac{a}{h} \Bigr)^2

\biggl(1+\sum\limits_{i=1}^n \tilde{\beta}_i^s \biggr) \alpha \Theta_0+\Bigl( \frac{2\pi a}{b}\Bigr)^4 2\sum\limits_{i=1}^n \beta_i^s+

\\

+4(1-\nu)\pi^2\Bigl(\frac{2\pi a}{b} \Bigr)^2  \sum\limits_{i=1}^n \beta_i^c, 

\end{multline}

\begin{equation*}

\beta_i^c=\Bigl( \frac{h_i}{h}\Bigr) ^3\frac{a_i}{a}\varPhi_{3i}\cos^2 \frac{\pi x_i}{a},

\quad 

\beta_i^s=\Bigl( \frac{h_i}{h}\Bigr) ^3\frac{a_i}{a}\varPhi_{3i}\sin^2 \frac{\pi x_i}{a},

\quad 

\tilde{\beta}_i^s= \frac{h_i}{h}\frac{a_i}{a}\sin^2 \frac{\pi x_i}{a}.

\end{equation*}



Система дифференциальных уравнений \eqref{bel_eq12} подстановкой 

\begin{equation*}

w_{11}=-\xi_{12}\varPhi(t), \quad w_{12}=\xi_{11}\left[\varPhi(t) \right] 

\end{equation*}

сводится к одному дифференциальному уравнению четвертого порядка, которое в случае отсутствия  демпфирования ($\mu=0$) содержит только четные производные от функции $\varPhi(t)$:

\begin{multline}\label{bel_eq15}

\biggl[ 

\frac{\gamma h a^4}{g D}\biggl( 1+2\sum\limits_{i=1}^n \tilde{\beta}_i^s\biggr) \biggr]^2

\frac{d^4\varPhi}{dt^4} +

\biggl[ \frac{\gamma h a^4}{g D}\biggl( 1+2\sum\limits_{i=1}^n \tilde{\beta}_i^s\biggr)

( f_{11}+f_{12})  \biggr]\frac{d^2\varPhi}{dt^2}+

\\

+\Bigl[ f_{11}f_{12}-8\pi \Bigl( \frac{a}{b}\Bigr)^2\mathit{I}_{13}\mathit{I}_{16}\tilde{M}+\tilde{M} 

\Bigl(4 \pi \frac{a}{b}\mathit{I}_{13}\tilde{f}_{11}+2 \pi \frac{a}{b}\mathit{I}_{16}\tilde{f}_{12} \Bigr) - \tilde{f}_{11}\tilde{f}_{12}\Bigr] \varPhi=0,

\end{multline}

где $\tilde{M}=\frac{p_0 \varkappa M a^3}{D}$.



Если $\mu\neq0$, то дифференциальное уравнение для функции $\varPhi(t)$ запишется в~виде

\begin{equation}\label{bel_eq16}

\frac{d^4\varPhi}{dt^4}+2 e_1\frac{d^3\varPhi}{dt^3}+e_3\frac{d^2\varPhi}{dt^2}+e_4\frac{d\varPhi}{dt}+e_5\varPhi=0,

\end{equation}

где

\begin{equation*}

e_1=\frac{\tilde{\mu} }{ \frac{\gamma h a^4}{g D}\left(1+2\sum_{i=1}^n\tilde{\beta}_i^s \right) }, \quad e_3=\frac{\frac{\gamma h a^4}{g D}\left(1+2\sum_{i=1}^n\tilde{\beta}_i^s \right)\left( f_{11}+f_{12}\right)+ \tilde{\mu}}{\left( \frac{\gamma h a^4}{g D}\left(1+2\sum_{i=1}^n\tilde{\beta}_i^s \right)\right)^2},

\end{equation*}

\begin{equation*}

e_4=\frac{\tilde{\mu}\left(f_{11}+f_{12} \right)  }{\left(  \frac{\gamma h a^4}{g D}\left(1+2\sum_{i=1}^n\tilde{\beta}_i^s \right)\right)^2},\quad e_5=\frac{f_{11} f_{12}-\xi_{21}\xi_{12}}{\left( \frac{\gamma h a^4}{g D}\left(1+2\sum_{i=1}^n\tilde{\beta}_i^s \right)\right)^2}, \quad \tilde{\mu}=\frac{\mu a^4 h}{D}.

\end{equation*}



Из условия, при котором хотя бы один корень из четырех  характеристического уравнения для дифференциального уравнения \eqref{bel_eq15} будет положительным, определяется наименьшее значение  относительной скорости потока газа ${v_y}/{c_0}$, при котором прогиб термоупругой системы растет. При учете демпфирования ($\mu\neq0$) возникает вопрос об устойчивости системы \eqref{bel_eq12}, которая сведена к одному дифференциальному уравнению \eqref{bel_eq16}. 

В~этом случае это значение скорости потока определяется на основании критерия Гурвица~\cite{myl-bib:32}:

\begin{equation*}

\Delta_1=2e_1>0, \quad \Delta_2=2e_1e_3-e_4>0, \quad \Delta_3=2e_1e_3e_4-e_4^2-4e_1^2e_5>0.

\end{equation*}



Кривизна пологой цилиндрической оболочки задавалась в виде $k_{11}\hm=-{4 \tilde{\delta}}/{a^2}$ \cite{myl-bib:33, myl-bib:34}, где $\tilde{\delta}$ "--- наибольший подъем оболочки над ее планом. Выражение для относительной кривизны $\tilde{k}_{11}$ принимает вид

\begin{equation*}

\tilde{k}_{11}=k_{11} a=-4\frac{\tilde{\delta}}{h}\frac{h}{a},

\end{equation*}

где ${\tilde{\delta}}/{h}\in[0, 5]$ \cite{myl-bib:33, myl-bib:34}.



\smallskip



{\bf 2.} Численные результаты приведены в~таблице. 

Вычисления проводились при следующих значениях параметров: 

${h}/{a}=0.01$, 

${a_i}/{a}=0.01$, 

$\mu=0.001$, 

${a}/{b}=1$.





\begin{table}[b!]

	

\small \centering



Значения параметра ${v_y}/{c_0}$  в зависимости от других параметров



[The values of the parameter ${v_y}/{c_0}$ depending on other parameters]

	

\smallskip

	\renewcommand{\tabcolsep}{0.35cm}

\begin{tabular}{c|c|c|c|c|c|c|c}

	\hline 

	\multicolumn{2}{c|} {\rule{0mm}{12pt}}&  \multicolumn{3}{c|} {${h_i}/{h}=2.5$}    &  \multicolumn{3}{c} {${h_i}/{h}=5$}   \\ 

	\hline 

 \backslashbox{$\tilde{k}_{11}$}{$n$} & 0 & 1 & 3 & 5 & 1 & 3 & 5 \\ 

	\hline 

		\multicolumn{8}{c} {\rule{0mm}{12pt}for $\Theta_0=0$,  $\varepsilon_1=\varepsilon_2=0$} \\

	\hline 	

	0.10 & 0.24 & 1.50 & 3.16 & 3.50 & 6.47 & 13.35 & 14.79 \\ 

	0.15 & 0.54 & 1.40 & 3.06 & 3.40 & 6.37 & 13.25 & 14.69 \\

	0.20 & 0.97 & 1.26 & 2.92 & 3.26 & 6.23 & 13.11 & 14.55 \\  

	\hline 

		\multicolumn{8}{c} {\rule{0mm}{12pt}for $\Theta_0=2$,  $\varepsilon_1=\varepsilon_2=0$} \\

	\hline 

	0.10 & 1.08 & 0.61 & 2.22 & 2.54 & 5.55 & 12.34 & 13.73 \\ 

	0.15 & 1.39 & 0.51 & 2.12 & 2.44 & 5.45 & 12.24 & 13.63 \\ 

	0.20 & 1.81 & 0.37 & 1.98 & 2.30 & 5.31 & 12.10 & 13.49 \\ 

	\hline 

		\multicolumn{8}{c} {\rule{0mm}{12pt}for $\Theta_0=4$,  $\varepsilon_1=\varepsilon_2=0$} \\

	\hline 

		0.10 & 1.93 & 0.27 & 1.28 & 1.58 & 4.63 & 11.32 & 12.67 \\ 

		0.15 & 2.23 & 0.37 & 1.18 & 1.48 & 4.53 & 11.22 & 12.57 \\ 

		0.20 & 2.66 & 0.51 & 1.04 & 1.34 & 4.39 & 11.08 & 12.43 \\ 

	\hline 

		\multicolumn{8}{c} {\rule{0mm}{12pt}for $\Theta_0=0$,  $\varepsilon_1=\varepsilon_2=1$} \\

	\hline 		

		0.10 & 0.24 & 1.50 & 3.16 & 3.51 & 6.48 & 13.37 & 14.80 \\ 

		0.15 & 0.54 & 1.40 & 3.07 & 3.42 & 6.39 & 13.28 & 14.72 \\ 

		0.20 & 0.97 & 1.27 & 2.94 & 3.29 & 6.26 & 13.16 & 14.60 \\ 

	\hline 

		\multicolumn{8}{c} {\rule{0mm}{12pt}for $\Theta_0=2$,  $\varepsilon_1=\varepsilon_2=1$} \\

	\hline 	

		0.10 & 1.08 & 0.61 & 2.23 & 2.55 & 5.56 & 12.35 & 13.74 \\ 

		0.15 & 1.39 & 0.52 & 2.13 & 2.46 & 5.47 & 12.26 & 13.66 \\ 

		0.20 & 1.81 & 0.38 & 2.00 & 2.33 & 5.34 & 12.14 & 13.54 \\ 

	\hline 

		\multicolumn{8}{c} {\rule{0mm}{12pt}for $\Theta_0=4$,  $\varepsilon_1=\varepsilon_2=1$} \\

	\hline 	

		0.10 & 1.93 & 0.26 & 1.29 & 1.58 & 4.64 & 11.33 & 12.68 \\ 

		0.15 & 2.23 & 0.36 & 1.19 & 1.49 & 4.55 & 11.25 & 12.60 \\ 

		0.20 & 2.66 & 0.49 & 1.06 & 1.36 & 4.42 & 11.13 & 12.48 \\ 

		\hline 

	\end{tabular}  

\end{table} 



Количественный анализ выявил закономерности влияния геометрических параметров  и температуры  на величины относительных скоростей потока,  начиная с которых прогибы термоупругой системы увеличиваются во времени, и 

сделать нижеследующий выводы.

 

\begin{enumerate}

	\item[1.] Учет в дифференциальных уравнениях системы \eqref{bel_eq1} слагаемых, учитывающих <<растяжение"=сжатие>> ребер и их <<сдвиг>> в~тангенциальной плоскости (слагаемых, содержащих знаковые числа), практически не влияет на величину относительной скорости ${v_y}/{c_0}$ газового потока. По этой причине удерживать эти слагаемые в уравнениях нет необходимости.

	\item[2.] Увеличение температуры ведет к уменьшению предельной скорости потока при прочих равных условиях.

	\item[3.] Существенно повышают динамическую устойчивость геометрически нерегулярной оболочки параметры ${h_i}/{h}$ и число подкрепляющих элементов~$n$. 

	\item[4.] Величина относительной предельной скорости потока мало чувствительна к изменению относительной кривизны пологой оболочки.

\end{enumerate} 





Перечисленные закономерности выявлены для случая двух полуволн в~направлении потока и одной полуволны в перпендикулярном направлении.  

Следует отметить, что при $k_{11}=0$ полученные решения аналитически преобразуются к решениям, приведенным в работах \cite{myl-bib:31, myl-bib:35}. В случае <<холодной>> ($\Theta_0=0$) гладкой пластины   ($k_{11}=0$, ${a_i}/{a}=0$) решения принимают вид, приведенный в~\cite{myl-bib:30}.



\AdditionalContent{Конкурирующие интересы}{Заявляем, что в отношении авторства и публикации этой статьи конфликта интересов не имеем.} 



\AdditionalContent{Авторский вклад и ответственность}{Все авторы принимали участие в разработке концепции статьи и в написании рукописи. Авторы несут ответственность за предоставление окончательной рукописи в печать. Окончательная версия рукописи была одобрена всеми авторами.} 



\AdditionalContent{Финансирование}{Результаты получены в рамках выполнения государственного задания Минобрнауки России № 9.8570.2017/8.9.}



\begin{thebibliography}{10}



\RBibitem{myl-bib:1}

\by Вольмир~А.~С.

\book Оболочки в потоке жидкости и газа

\yr 1979

\publ Наука

\publaddr М.

\totalpages 320



\RBibitem{myl-bib:2}

\by Амбарцумян~C.~А., Багдасарян~Ж.~Е.

\paper Об устойчивости ортотропных пластинок, обтекаемых сверхзвуковым потоком газа

\jour Изв. АН СССР, ОТН, Механика и машиностроение

\yr 1961

\issue 4

\pages 91--96



\RBibitem{myl-bib:3}

\by Болотин~В.~В., Новичков~Ю.~Н.

\paper Выпучивание и установившийся флайтер термически сжатых панелей, находящихся в сверхзвуковом потоке

\jour Инж. журн.

\yr 1961

\issue 2

\pages 82--96







\RBibitem{myl-bib:4}

\by Мовчан~А.~А.

\paper О колебаниях пластинки, движущейся в газе

\jour ПММ

\yr 1956

\vol 20

\issue 2

\pages 211--222





\RBibitem{myl-bib:5}

\by Дун~Мин-дэ,

\paper Об устойчивости упругой пластинки при сверхзвуковом обтекании

\jour Докл. АН СССР

\yr 1958

\vol 120

\issue 4

\pages 726--729

\mathnet{http://mi.mathnet.ru/dan23108}



\RBibitem{myl-bib:5a}

\by Веденеев~В.~В. 

\paper Высокочастотный флаттер прямоугольной пластины

\jour Изв. РАН. МЖГ

\yr 2006

\issue 4

\pages 173--181









\RBibitem{myl-bib:6}

\by Огибалов~П.~М., Грибанов~В.~Ф.

\book Термоустойчивость пластин и оболочек

\yr 1968

\publ МГУ

\publaddr М.

\totalpages 520







\RBibitem{myl-bib:7}

\by Огибалов~П.~М. 

\book Вопросы динамики и устойчивости оболочек

\yr 1963

\publ МГУ

\publaddr М.

\totalpages 417

\zmath{0115.41504}





\RBibitem{myl-bib:8}

\by Болотин~В.~В.

\paper Температурное выпучивание пластин и пологих оболочек в сверхзвуковом потоке газа 

\inbook Расчеты на прочность

\yr 1960

\volinfo Вып.~6

\publaddr М.

\publ Машгиз

\pages 190--216



\RBibitem{myl-bib:9}

\by Жилин~П.~А.

\paper Линейная теория ребристых оболочек

\jour Изв. АН СССР. МТТ

\yr 1970

\issue 4

\pages 150--166



 



\RBibitem{myl-bib:10}

\by Белосточный~Г.~Н., Ульянова~О.~И.

\paper Континуальная модель композиции из оболочек вращения с термочувствительной толщиной

\jour Изв. РАН. МТТ

\yr 2011

\issue 2

\pages 184--191

\elib{https://elibrary.ru/item.asp?id=15785027}



 

 \RBibitem{myl-bib:11}

 \by Белосточный~Г.~Н., Рассудов~В.~М.

 \paper Континуальная модель термочувствительной ортотропной системы «оболочка–ребра» с учетом влияния больших прогибов

 \inbook Механика деформируемых сред

 \yr 1983

\volinfo Вып.~8

\publaddr Саратов

\publ  Сарат. политехн. ин-т

 \pages 10--22



  

\RBibitem{myl-bib:12}

\by Белосточный~Г.~Н., Рассудов~В.~М.

\paper Континуальный подход к термоустойчивости упругих систем «пластинка–ребра» 

\inbook Прикладная теория упругости

\yr 1980

\publaddr Саратов

\publ  Сарат. политехн. ин-т

\pages 94--99





 









\RBibitem{myl-bib:15}

\by Жилин~П.~А.

\paper Общая теория ребристых оболочек. Прочность гидротурбин

\inbook Тр. ЦКТИ

\yr 1968

\volinfo Вып.~8

\publaddr Л.

\pages 46--70







\RBibitem{myl-bib:16}

\by Карпов~В.~В., Сальников~А.~Ю.

\paper Вариационный метод вывода нелинейных уравнений движения пологих ребристых оболочек 

\jour Вестн. гражд. инженеров

\yr 2008

\issue 4(17)

\pages 121--124

\elib{https://elibrary.ru/item.asp?id=11675045}







\RBibitem{myl-bib:18}

\by  Белосточный~Г.~Н., Мыльцина~О.~А.

\paper Уравнения термоупругости композиций из оболочек вращения 

\jour Вестник СГТУ

\yr 2011

\issue 4 (59)

\issueinfo Вып.~1

\pages 56--64



   

\RBibitem{myl-bib:19}

\by  Онанов~Г.~Г.

\paper Уравнения с сингулярными коэффициентами типа дельта-функции и~ее производных

\jour Докл. АН СССР

\yr 1970

\vol 191

\issue 5

\pages 997--1000

\mathnet{http://mi.mathnet.ru/dan35339}

\mathscinet{http://www.ams.org/mathscinet-getitem?mr=0470298}

\zmath{https://zbmath.org/?q=an:0215.20101}









\RBibitem{myl-bib:21}

\by  Белосточный~Г.~Н.

\paper  Аналитические методы определения замкнутых интегралов сингулярных дифференциальных уравнений термоупругости геометрически нерегулярных оболочек

\jour Докл. Акад. воен. наук

\yr 1999

\issue 1

\pages 14--26





\Bibitem{myl-bib:22a}

\by Geckeler~J.~W.

\paper Elastostatik

\lang In German

\zmath{54.0844.01}

\inbook Handbuch der Physik 

\bookvol 6

\pages 141--308

\yr 1928





\RBibitem{myl-bib:22}

\by Ильюшин~А.~А.

\paper  Закон плоских сечений в аэродинамики больших сверхзвуковых скоростей 

\jour ПММ

\yr 1956

\issue 6

\pages 733--755







 

\RBibitem{myl-bib:30}

\by Вольмир~А.~С. 

\book  Устойчивость деформируемых систем

\yr 1967

\publ Наука

\publaddr М.

\totalpages 984



 

 



 

\RBibitem{myl-bib:25}

\by Владимиров~В.~С. 

\book Обобщенные функции в математической физике

\yr 1979

\publ Наука

\publaddr М.

\totalpages 319

\zmath{0515.46033}









\RBibitem{myl-bib:27}

\by Канторович~Л.~В., Крылов~В.~И.

\book  Приближенные методы высшего анализа

\yr 1962

\publ Физматлит

\publaddr М.

\totalpages 708

\zmath{0100.33102}









\Bibitem{myl-bib:29}

\by Rektorys~K.

\book Variational methods in mathematics, science and engineering

\zmath{0481.49002}

\publaddr Dordrecht, Boston,  London

\publ D.~Reidel Publ. 

\totalpages 571 

\yr 1980

\moreref

\crossref{10.1007/978-94-011-6450-4}. \nofrills





\RBibitem{myl-bib:31}

\by Белосточный~Г.~Н., Рассудов~В.~М. 

\paper  Термоупругие системы типа «пластинка–ребра» в сверхзвуковом потоке газа  

\inbook Прикладная теория упругости

\publaddr Саратов

\volinfo Вып.~8

\publ Сарат. политехн. ин-т

\yr 1983

\pages 114--121







\RBibitem{myl-bib:32}

\by Егоров~К.~В.  

\book  Основы теории автоматического регулирования

\yr 1967

\publ Энергия

\publaddr М.

\totalpages 648





\RBibitem{myl-bib:33}

\by Рассудов~В.~М., Красюков~В.~П., Панкратов~Н.~Д.  

\book  Некоторые задачи термоупругости пластинок и пологих оболочек 

\yr 1973

\publ Сарат. ун-т

\publaddr Саратов

\totalpages 155





\RBibitem{myl-bib:34}

\by Назаров~А.~А.

\book Основы теории и методы расчета пологих оболочек

\yr 1966

\publ  Стройиздат

\publaddr Л., М.

\totalpages 304





\RBibitem{myl-bib:35}

\by Мыльцина~О.~А., Белосточный~Г.~Н. 

\paper  Устойчивость нагретой ортотропной геометрически нерегулярной пластинки в сверхзвуковом потоке газа 

\jour Вестник Пермского национального исследовательского политехнического университета. Механика

\yr 2017

\issue 4

\pages 109--120

\crossref{10.15593/perm.mech/2017.4.08}















\end{thebibliography}





\newpage



\makeentitle



\vspace{-5mm}





\AdditionalContent{Competing interests}{We declare that we have no conflicts of interest in the authorship and publication of this article.} 



\AdditionalContent{Authors' contributions and responsibilities}{Each author has participated in the article concept development and in the manuscript writing. The authors are absolutely responsible for submitting the final manuscript in print. Each author has approved the final version of manuscript.} 





\AdditionalContent{Funding}{The results have been obtained within the State Assignment of the Ministry of Education

and Science of the Russian Federation no. 9.8570.2017/8.9.}









\begin{thebibliography}{10}



\Bibitem{myl-bib:1:en}

\lang In Russian

\by Vol'mir~A.~S.

\book Obolochki v potoke zhidkosti i gaza \rm [The shells in flow liquid and gas] 

\publaddr Moscow

\publ Nauka

\yr 1979

\totalpages 320



\Bibitem{myl-bib:2:en}

\lang In Russian

\by Ambartsumyan~S.~A. Bagdasaryan~Zh.~E.

\paper Of the stability of orthotropic plates streamlined by a supersonic gas flow

\jour Izv. Akad. Nauk SSSR, OTN, Mekhanika i Mashinostroenie

\yr 1961

\issue 4

\pages 91--96



\Bibitem{myl-bib:3:en}

\lang In Russian

\by Bolotin~V.~V., Novichkov~Yu.~N. 

\paper Buckling and steady-state flayter thermally compressed panels in supersonic flow

\jour Inzhenernyi Zhurnal 

\yr 1961

\issue 2

\pages 82--96







\Bibitem{myl-bib:4:en}

\lang In Russian

\by Movchan~A.~A. 

\paper Oscillations of the plate moving in the gas

\jour Prikladnaia Matematika i Mekhanika

\yr 1956

\vol 20

\issue 2

\pages 211--222





\Bibitem{myl-bib:5:en}

\lang In Russian

\by Tung~Ming-teh,

\paper The stability of an elastic plate in a supersonic flow

\jour Dokl. Akad. Nauk SSSR

\yr 1958

\vol 120

\issue 4

\pages 726--729



\Bibitem{myl-bib:5a}

\by Vedeneev~V.~V.

\jour Fluid Dyn.

\crossref{10.1007/s10697-006-0083-2}

\yr 2006

\vol 41

\issue 4

\pages 641–648 

\paper High-frequency flutter of a rectangular plate









\Bibitem{myl-bib:6en}

\lang In Russian

\by Ogibalov~P.~M., Gribanov~V.~F. 

\book Termoustoichivost' plastin i obolochek \rm  [Thermostability of plates and shells]

\yr 1968

\publaddr Moscow

\publ Moscow State Univ.

\totalpages 520







\Bibitem{myl-bib:7en}

\lang In Russian

\by Ogibalov~P.~M.

\book Voprosy dinamiki i ustoichivosti obolochek\rm  [The problems of dynamics and stability of shells]

\publaddr Moscow

\publ Moscow State Univ.

\yr 1963

\totalpages 417







\Bibitem{myl-bib:8:en}

\lang In Russian

\by Bolotin~V.~V.

\paper Thermal buckling of plates and shallow shells in a supersonic gas flow

\inbook Raschety na prochnost'

\publaddr Moscow

\publ Mashgiz

\yr 1960

\volinfo Issue~6

\pages 190--216



\Bibitem{myl-bib:9:en}

\lang In Russian

\by Zhilin~P.~A.

\paper Linear theory of ribbed shells

\jour Izv. Akad. Nauk SSSR, Mekhanika Tverdogo Tela

\yr 1970

\issue 4

\pages 150--166



 



\Bibitem{myl-bib:10:en}

\paper Continuum model for a composition of shells of revolution with thermosensitive thickness

\by Belostochnyi~G.~N., Ul'yanova~O.~I.

\jour Mechanics of Solids

\yr 2011

\vol 46

\issue 2

\pages 184--191

\elib{https://elibrary.ru/item.asp?id=16982711}

\crossref{10.3103/S0025654411020051}



 

\Bibitem{myl-bib:11:en}

\lang In Russian

\by Belostochnyi~G.~N., Rassudov~V.~M.

\paper Continuum model of orthotropic heat-sensitive ``shell-ribs'' system taking into account the influence of large deflection

 \yr 1983

\volinfo Issue~8

\inbook Prikladnaia teoriia uprugosti

\publ Saratov Polytechnic Inst.

\publaddr Saratov 

 \pages 10--22



  

\Bibitem{myl-bib:12:en}

\lang In Russian

\by Belostochnyi~G.~N., Rassudov~V.~M.

\paper Continuum approach to the thermal stability of elastic ``plate-ribs'' systems

\inbook Prikladnaia teoriia uprugosti

\publ Saratov Polytechnic Inst.

\publaddr Saratov 

\yr 1980

\pages 94--99







\Bibitem{myl-bib:15:en}

\lang In Russian

\by Zhilin~P.~A.

\paper General theory of ribbed shells. The strength of turbines

\inbook  Proc. of the Central Turbine-Boiler Institute

\yr 1968

\volinfo Issue~8

\publaddr Leningrad

\pages 46--70







\Bibitem{myl-bib:16:en}

\lang In Russian

\by Karpov~V.~V., Sal'nikov~A.~Yu.

\paper Variational method output nonlinear equations of motion of shallow ribbed shells

\jour Vestnik Grazhdanskikh Inzhenerov

\yr 2008

\issue 4(17)

\pages 121--124





\Bibitem{myl-bib:18:en}

\lang In Russian

\by Belostochnyi~G.~N., Myltcina~O.~A.

\paper Thermoelasticity equations of shells compositions

\jour Vestnik Saratovskogo Tekhnicheskogo Universiteta

\yr 2011

\issue 4 (59)

\issueinfo Issue~1

\pages 56--64



   

\Bibitem{myl-bib:19:en}

\lang In Russian

\by Onanov~G.~G.

\paper Equations with singular coefficients of the type of the delta-function and its derivatives

\jour Dokl. Akad. Nauk SSSR

\yr 1970

\vol 191

\issue 5

\pages 997--1000









\Bibitem{myl-bib:21:en}

\lang In Russian

\by Belostochnyi~G.~N.

\paper Analytical methods for the determination of closed integrals of singular differential equations of thermoelasticity, geometrically irregular shells

\jour Doklady Akademii Voennykh Nauk

\yr 1999

\issue 1

\pages 14--26





\Bibitem{myl-bib:22a:en}

\by Geckeler~J.~W.

\paper Elastostatik

\lang In German

\zmath{54.0844.01}

\inbook Handbuch der Physik 

\bookvol 6

\pages 141--308

\yr 1928





\Bibitem{myl-bib:22:en}

\lang In Russian

\by Il'yushin~A.~A.

\paper Law of plane sections in aerodynamics of high supersonic velocity

\jour Prikladnaia Matematika i Mekhanika

\yr 1956

\issue 6

\pages 733--755







 

\Bibitem{myl-bib:30:en}

\lang In Russian

\by Vol'mir~A.~S.

\book  Ustoichivost' deformiruemykh sistem \rm [Stability of deformable systems]

\publaddr Moscow

\publ Nauka

\yr 1967

\totalpages 984



 

 



 

\Bibitem{myl-bib:25:en}

\by Vladimirov~V.~S.

\book Generalized functions in mathematical physics

\zmath{0515.46034}

\publaddr Moscow

\publ Mir Publ.

\totalpages  362 

\yr 1979









\Bibitem{myl-bib:27:en}

\by Kantorovich~L.~V., Krylov~V.~I.

\book Approximate methods of higher analysis

\zmath{0083.35301}

\publaddr Groningen

\publ P.~Noordhoff Ltd.

\totalpages xii+681

\yr 1958











\Bibitem{myl-bib:29:en}

\by Rektorys~K.

\book Variational methods in mathematics, science and engineering

\zmath{0481.49002}

\publaddr Dordrecht, Boston,  London

\publ D.~Reidel Publ. 

\totalpages 571 

\yr 1980

\moreref

\crossref{10.1007/978-94-011-6450-4}. \nofrills





\Bibitem{myl-bib:31:en}

\lang In Russian

\by Belostochnyi~G.~N., Rassudov~V.~M.

\paper Thermoelastic system of ``plate-ribs'' type  in a supersonic gas flow

\inbook Prikladnaia teoriia uprugosti

\publ Saratov Polytechnic Inst.

\publaddr Saratov 

\volinfo Issue~8

\yr 1983

\pages 114--121







\Bibitem{myl-bib:32:en}

\lang In Russian

\by Egorov~K.~V.

\book Osnovy teorii avtomaticheskogo regulirovaniia \rm [Fundamentals of the theory of automatic control]

\publaddr Moscow

\publ Energiia

\yr 1967

\totalpages 648





\Bibitem{myl-bib:33:en}

\lang In Russian

\by Rassudov~V.~M., Krasiukov~V.~P., Pankratov~N.~D.  

\book  Nekotorye zadachi termouprugosti plastinok i pologikh obolochek  \rm [Some problems of thermoelasticity of plates and flat covers]

\yr 1973

\publ Saratov Univ.

\publaddr Saratov

\totalpages 155

 





\Bibitem{myl-bib:34:en}

\lang In Russian

\by Nazarov~A.~A.

\book Osnovy teorii i metody rascheta pologikh obolochek \rm [Fundamentals of the Theory and Methods for Designing Shallow Shells]

\yr 1966

\publ  Stroiizdat

\publaddr Leningrad, Moscow

\totalpages 304





\Bibitem{myl-bib:35:en}

\lang In Russian

\by Myltcina~O.~A., Belostochnyi~G.~N.

\paper Stability of heated orthotropic geometrically irregular plate in a supersonic gas flow

\jour PNRPU Mechanics Bulletin

\yr 2017

\issue 4

\pages 109--120

\crossref{10.15593/perm.mech/2017.4.08}





\end{thebibliography}







%\end{document}

 \newpage

%\input{./preamble}
%
%\usepackage{samgtu-name}
%
\vsgtu{1624} %registration number
\FirstPage{762}
\LastPage{773}
%
%\pdfcrop
%\draft  
%
\ShortCommunication
%
%\begin{document} 




\hfill \qrcode[hyperlink,height=.8in]{http://doi.org/10.14498/vsgtu1624}
\vspace{-.9in}

\udc{539.3}
\msc{74C10}

\rutitle{Об одном подходе к определению предельной несущей~способности
	механических систем с~разупрочняющимися элементами}
\rutitleshort{Об одном подходе к определению предельной несущей способности
	механических систем\dots}

\entitle{One approach to determination of the ultimate load-bearing capacity of mechanical systems with softening elements}
\entitleshort{One approach to determination of the ultimate load-bearing capacity of mechanical systems\dots}

\ruauthor{В.~В.~Стружанов\affil{1}, А.~В.~Коркин\affil{2}, А.~Е.~Чайкин\affil{2}}
{С\,т\,р\,у\,ж\,а\,н\,о\,в~В.~В., К\,о\,р\,к\,и\,н~А.~В., Ч\,а\,й\,к\,и\,н~А.~Е.}
\enauthor{V.\,V.~Struzhanov\affil{1}, A.\,V.~Korkin\affil{2}, A.\,E.~Chaykin\affil{2}} {S\,t\,r\,u\,z\,h\,a\,n\,o\,v~V.~V., K\,o\,r\,k\,i\,n~A.~V., C\,h\,a\,y\,k\,i\,n~A.~E.}

\ruaffil{\ruimash.
\and 
\ruurfuien.}

\enaffil{\enimash.
\and 
\enurfuien.}

\ruauthorsdetails{3}{% Количество авторов
\fullauthor{\rupid{39278}{Валерий Владимирович Стружанов}}
\CAuthor % Автор-корреспондент
\orcid{0000-0002-3669-2032} % ORCID ID автора 
\regalia{доктор физико-математических наук, профессор} 
\position{главный научный сотрудник} 
\dept[p]{~лаб. микромеханики материалов} 
\email{stru@imach.uran.ru}
\fullauthor{\rupid{73471}{Александр Владимирович Коркин}}
\orcid{0000-0003-3533-4257} % ORCID ID автора 
\regalia{аспирант} 
\email{alexkorkin@list.ru}
\fullauthor{\rupid{143469}{Алексей Евгеньевич Чайкин}}
\orcid{0000-0001-5582-2384} % ORCID ID автора 
\regalia{аспирант}  
\email{chaykin.ae@yandex.ru}\vspace{-3mm}
}

\enauthorsdetails{3}{
\fullauthor{\enpid{39278}{Valery V. Struzhanov}}
\CAuthor % Сorresponding Author
\orcid{0000-0002-3669-2032} % ORCID ID avtora 
\regalia{Dr. Phys. \& Math. Sci., Professor} 
\position{Chief Reseacher} 
\dept{Lab. of Matherial Micromechanics}
\email[br]{stru@imach.uran.ru}
\fullauthor{\enpid{73471}{Aleksandr V. Korkin}}
\orcid{0000-0003-3533-4257} % ORCID ID avtora 
\regalia{Postgraduate Student} 
\email{alexkorkin@list.ru}
\fullauthor{\enpid{143469}{Aleksey E.  Chaykin}}
\orcid{0000-0001-5582-2384} % ORCID ID avtora 
\regalia{Postgraduate Student} 
\email{chaykin.ae@yandex.ru}
}

\rukeywords{потенциальная функция, вырожденные критические точки, единый потенциал, матрица Гессе, тонкостенный резервуар, предельное давление}
\enkeywords{potential function, degenerate critical points, unified potential, Hesse matrix, thin walled reservoir, limiting pressure}

\ruabstract{Изложены основные положения теории расчета предельных нагрузок, действующих на дискретные механические системы с~разупрочняющимися элементами. Методика опирается на численное определение вырожденных критических точек потенциальной функции системы, где происходит переход от устойчивости процесса нагружения к~неустойчивости (катастрофа, разрушение), и~позволяет избежать решения большого числа нелинейных уравнений равновесия. 
В~качестве примера применения методики решена задача об определении предельного внутреннего давления в~тонкостенном цилиндрическом резервуаре. При построении потенциальной функции системы использован специально построенный единый потенциал для плоского квадратного элемента материала в~условиях двухосного растяжения, описывающий все стадии деформирования, включая и~разупрочнение.}

\enabstract{The fundamental provisions of the limiting load calculation theory are presented for a discrete mechanical system with softening elements. The method is based on the numerical determination of degenerate critical points for the potential function of the system. At these points there is a transition from the stability of the loading process to instability such as a catastrophe or a failure. This approach helps to avoid solving a large number of nonlinear equilibrium equations. The problem of determining the limiting internal pressure in a thin walled cylindrical tank is solved as an example. A unified potential specially defined for a flat square element of material in biaxial tension is used in developing a potential function of the system. It describes all stages of deformation including the softening stage.}

\newdate{struzha}{11}{05}{2018}
\date{\displaydate{struzha}}
\newdate{struzhb}{11}{10}{2018}
\revisiondate{\displaydate{struzhb}}
\newdate{struzhc}{12}{11}{2018}
\accepteddate{\displaydate{struzhc}}

\newdate{struzhe}{11}{12}{2018}
\renewcommand{\JournalDataOnline}{\displaydate{struzhe}}
%\ForPeerReview
\makerutitle

\newpage

\Section[n]{Введение}
Вычисление предельной несущей способности механических систем основывается на расчете параметров равновесных состояний при постепенно возрастающей нагрузке. Затем эти параметры тем или иным способом сводят к некоторому скаляру и сравнивают его с предельной величиной, полученной в эксперименте. Например, на каждом шаге догружения определяют напряженное состояние элемента конструкции, находят соответствующий инвариант тензора напряжений, который сравнивают с его предельной величиной, взятой из эксперимента.

Возможен и другой подход, основанный на исследовании устойчивости положений равновесия, когда разрушение связывают с невозможностью равновесия, то есть с моментом потери устойчивости [\citen{bib1, bib2}]. Реализация данного подхода возможна, если учитывать физически неустойчивые состояния материала (разупрочнение) [\citen{bib2}--\citen{bib6}]. Но и в этом случае возникает необходимость в решении не удовлетворяющих условиям корректности по Адамару [\citen{bib8}]  больших систем нелинейных уравнений равновесия для определения критических точек потенциальной функции механической системы, что является нетривиальной задачей [\citen{bib9}]. Кроме того, выделение из них вырожденных критических точек, где происходит смена устойчивости на неустойчивость, требует применения нетрадиционного аппарата теории катастроф [\citen{bib10,bib11}]. Однако информация о параметрах равновесия (критических точках потенциальной функции системы) при возрастающих нагрузках, вообще говоря, является излишней. При определении предельной несущей способности нужно знать только вырожденные критические точки. В данной работе представлена методика, которая позволяет избежать аналитического решения нелинейных уравнений. 
Она основана на специальной числовой процедуре приближенного выбора вырожденных критических точек. В качестве примера рассмотрена задача о~разрушении цилиндрического резервуара под действием внутреннего давления. Для построения потенциальной функции разработана методика построения единого потенциала для элемента материала, находящегося в плоском напряженном состоянии и описывающего его свойства как при упрочнении, так и при разупрочнении.

\smallskip

\Section{Общие положения} Положение механической системы в пространстве определяется конечным набором обобщенных координат (параметров состояния)  $ q_i$, $i = 1, 2, \ldots, N $. 
Величины параметров состояния зависят от набора задаваемых параметров управления системой $ s_j$, 
$j = 1, 2, \ldots, M $ 
(внешние силы, перемещения некоторых точек системы, физико"=механические свойства элементов и т.~п.). 
Причем параметры состояния должны принимать такие значения, при которых механическая система находится в состоянии равновесия. 
Разрушение системы есть явление того же порядка, что и явление невозможности равновесия [\citen{bib1}]. 
Когда равновесие становится неустойчивым, то реализуется или скачкообразный переход в ближайшее устойчивое равновесие, или глобальная катастрофа (разрушение). Поэтому под предельными значениями параметров управления следует понимать такие их величины, при реализации которых нарушается устойчивость системы.

В тех случаях, когда механическая система является градиентной [\citen{bib11}], для определения предельных значений параметров управления можно воспользоваться следующей методикой. Градиентная система при квазистатическом нагружении с постоянной температурой описывается потенциальной функцией (лагранжианом) $ \Pi(q_i ,s_j) $, связывающей параметры состояния и~управления. 
Эта функция есть сумма потенциальных функций элементов системы, которые могут быть как вогнутыми, так и~выпукло"=вогнутыми. 
Выпуклые вниз (вогнутые) участки отвечают  устойчивым состояниям элемента, выпуклые (выпуклые вверх) "--- неустойчивым.  
Точки перехода от выпуклости к~вогнутости "--- пограничные точки, являющиеся в общем случае седловыми точками соответствующих потенциальных функций. 

Когда механическая система находится в~равновесии (устойчивом или неустойчивом), для связи параметров состояния и~управления в~этом равновесии можно использовать уравнение Лагранжа второго рода [\citen{bib12}]. 
Тогда уравнение равновесия имеет вид
\begin{equation}\label{S_eq1}
\frac{\partial\Pi(q_i,s_j)}{\partial q_i} = 0, \quad  i=1, 2, \ldots,N.
\end{equation}
Если функция   является выпукло"=вогнутой, то при заданных управляющих параметрах система \eqref{S_eq1} может иметь одно или более одного решения, как устойчивых, так и~неустойчивых. 
Следовательно, в~данном случае не выполняются условия Адамара [\citen{bib8}].

Все решения системы \eqref{S_eq1} для всевозможных параметров управления образуют совокупность критических точек потенциальной функции $ \Pi $.  
Известно [\citen{bib11}], что смена типа равновесия (устойчивого или неустойчивого) происходит в~вырожденных критических точках, которые определяются из совместного решения уравнений \eqref{S_eq1} и~уравнения, получающегося приравниванием к~нулю детерминанта матрицы устойчивости (матрицы Гессе потенциальной функции):
\begin{equation}\label{S_eq2}
\mathop{\rm det}\Bigl| \frac{\partial^{2}\Pi(q_i,s_j)}{\partial q_m \partial s_n}\Bigr|  = 0,  \quad m,n = 1, 2, \ldots, N.
\end{equation}
Очевидно, что уравнение \eqref{S_eq2} выделяет из множества критических точек вырожденные критические точки. 

Таким образом, для отыскания предельных значений параметров управления, при достижении которых система теряет устойчивость, сначала, как правило, вычисляются все критические точки (решения уравнений \eqref{S_eq1}), затем выделяются вырожденные точки и после подстановки их в уравнение \eqref{S_eq1} определяются соответствующие параметры управления. Те параметры управления, при которых в первый раз достигается вырожденная критическая точка, то есть происходит смена устойчивого состояния механической системы на неустойчивое, и будут предельными значениями внешних нагрузок.

При реализации данной системы нахождения предельных значений параметров управления необходимо искать решения большого числа нелинейных уравнений, что представляет достаточно сложную задачу. 
Однако задача о~вычислении предельных значений параметров управления не требует, вообще говоря, знания всех критических точек функции $ \Pi $ (параметров всех положений равновесия механической системы). Необходимы только вырожденные точки, удовлетворяющие уравнению \eqref{S_eq2}. Кроме того, для технических приложений можно использовать их приближенные значения, вычисленные с достаточной степенью точности. 

Вычисление приближенных значений вырожденных критических точек можно осуществить, применяя следующую численную процедуру. Рассмотрим евклидово пространство $ \mathbb R^M \times \mathbb R^N $ "--- декартово произведение евклидовых пространств размерности $ M $ и $ N $. 
Элементами этого пространства являются числа $ ( q_i,s_j) $. 
Задавая физически обоснованные пределы изменения параметров состояния и управления, получаем $ (M{+}N) $"=мерный куб в пространстве $\mathbb  R^M \times \mathbb R^N $, 
который разбиваем сеткой узлов на множество «кубиков» с~заданным шагом. 
Затем в~каждом узле сетки разбиения вычисляем значения компонент матрицы Гессе и~величину ее определителя. 
Таким образом, каждому узлу ставится в~соответствие значение детерминанта матрицы Гессе. 

После этого выделяем те узлы, в которых детерминант близок к нулю с достаточной степенью точности. 
Полученное множество $B$ точек пространства $ \mathbb R^M \times \mathbb R^N $ содержит в~себе все вырожденные критические точки потенциальной функции~$ \Pi $. 
Чтобы выделить из множества $B$ критические точки функции $ \Pi $, следует подставить элементы множества $B$ в~уравнения равновесия \eqref{S_eq1}.  
Те~из них, которые удовлетворяют этим уравнениям с~заданной степенью точности, будут вырожденными критическими точками. 
Параметры управления,  соответствующие ближайшей (по евкледовой метрике) к~началу координат вырожденной точке, образуют совокупность искомых предельных значений управлений. 
Если существует группа вырожденных точек, находящихся на одинаковом наименьшем расстоянии от начала координат, то получим различные сочетания значений параметров управления.

Таким образом, минуя стадию расчета параметров равновесных состояний, можно найти критические нагрузки, что достаточно для оценки прочности механических систем.

\smallskip

\Section{Модельный пример} Применим изложенную выше методику к определению предельных параметров управления (нагрузок) при растяжении модельной стержневой системы, изображенной на  рис.~\ref{Sfig1}. 
Стержень {\sl 2} "--- абсолютно упругий с жесткостью при растяжении, равной $c$. 
Стержень {\sl 1} изготовлен из разупрочняющегося материала.

\begin{figure}[b!]
\centering
		\includegraphics[width=0.4\textwidth]{fig_stru1}
		\caption{Растягиваемая стержневая система с разупрочняющимся элементом} \label{Sfig1}
		\smallskip \footnotesize
		[Figure~\ref{Sfig1}. 		A rod system with a softening element under tension]

\end{figure}

\smallskip

{\bf 2.1.}
При мягком нагружении растягивающая сила $ p $ "--- параметр управления, 
а~удлинение $ x $  стержня {\sl 1}  и перемещение $ u $ правого свободного конца   стержня~{\sl 2}  "--- параметры состояния (см.  рис.~\ref{Sfig1}). 
Сопротивление растяжению стержня {\sl 1} задано функцией $ g(x) $ 
(зависимость растягивающей силы от удлинения), обладающей как восходящей ветвью (упрочнение), так и~ветвью, падающей до нуля (стадия разупрочнения) [\citen{bib2, bib5}].  
В~этом случае его потенциальная функция, равная $\displaystyle \int ^{x}_{0} g(x)\, dx $, является выпукло"=вогнутой, а лагранжиан системы~[\citen{bib5}] 
$$
Q^p = \int^{x}_{0} g(x)  dx + \frac{c}{2} (u-x)^2 - \int ^{u}_{0} p  du,
$$	
где второе слагаемое "--- энергия упругих деформаций стержня {\sl 2}, последнее "--- работа внешней силы, взятая со знаком минус.

Уравнения равновесия имеют вид
\begin{equation}\label{S_eq3}
Q^{p}_{\bigcom x} = g(x) - c(u-x) = 0, \quad  Q^{p}_{\bigcom u} = c(u-x) - p = 0.
\end{equation}
Здесь запятой обозначены частные производные по соответствующим аргументам.
Приравнивая к нулю детерминант матрицы Гессе, получим уравнение 
\begin{equation}\label{S_eq4}
g_{\bigcom x} = 0.
\end{equation}

Отметим, что это уравнение определяет максимум функции  $ g(x) $. 
Его можно разрешить, используя следующую процедуру. 

Проведем разбиение оси $X $ узлами с~достаточно малым шагом. 
Подставляя координаты этих узлов в~уравнение \eqref{S_eq4}, находим среди них значение $ x' $, приближенно ему удовлетворяющее. 
Затем из уравнений \eqref{S_eq3} для данного $ x' $ определяем $ p' $ "--- предельную величину растягивающей силы.

\smallskip

{\bf 2.2.}
В случае жесткого нагружения, когда растяжение осуществляется заданным перемещением $ u $ (параметр управления), а $ x $ является единственным параметром состояния, лагранжиан системы есть функция $ Q^u $, равная двум первым слагаемым в выражении для $ Q^p $. 

Тогда имеем лишь одно уравнение равновесия (первое уравнение в~системе \eqref{S_eq3}). 
В~матрице Гессе только один элемент. 
Тогда равенство детерминанта нулю дает уравнение $ g'_{x} + c = 0 $. 
Его приближенное решение ищем так же, как и~решение уравнения \eqref{S_eq4}. 

В зависимости от величины параметра $ c $ возможны три варианта. 

В первом "--- решения не существует. Следовательно, деформирование системы идет равновесно вплоть до разделения стержня {\sl 1} на части при значении $ x = x^z $, когда падающая ветвь функции $ g(x) $ достигает оси $ X $. Предельное значение параметра управления получаем из уравнения равновесия после подстановки $ x = x^z $: 
$$ u^z = g(x^z)/c+ x^z.$$ 

Во втором варианте имеется только одно приближенное решение $ x =  x_1 $. 
После достижения параметром $ x $ значения этой величины система переходит в~неустойчивое состояние. 
Поэтому предельный параметр управления $$ u_1 = g(x_1)/c + x_1 .$$ 

И, наконец, в третьем варианте имеем два решения $ x_1$, $x_2$  $(x_1 > x_2) $. 
Тогда предельное значение  $ u = u_1 $.

Отметим, что такие же результаты были получены и в работах [\citen{bib5, bib9}] только после решения нелинейных уравнений равновесия при постепенном возрастании параметров управления и построения соответствующих многообразий катастроф и сепаратрис, состоящих из вырожденных критических точек.

\smallskip

\Section{Вогнуто"=выпуклый потенциал при двухосном растяжении тонкой пластины} 
Для реализации положений теории, изложенной выше, необходимо знать вогнуто"=выпуклый потенциал элементов системы. 
В~механике сплошной среды эта функция должна быть определена для элемента материала вплоть до разупрочнения и~разделения на фрагменты.

Рассмотрим плоский квадратный элемент материала с~единичными размерами, находящийся в плоско"=напряженном состоянии, причем касательные напряжения и~деформации отсутствуют. 
Деформирование осуществляется заданием нормальных деформаций $ \varepsilon_1$, $\varepsilon_2 $ (двуосное растяжение). Если материал находится в состоянии упругости, то его  потенциальная энергия 
\begin{equation}\label{S_eq5}
V^e(\varepsilon_1, \varepsilon_2) = \frac{1}{2}(\sigma_1\varepsilon_1 + \sigma_2\varepsilon_2) = \frac{E}{2(1-\nu^2)} (\varepsilon^2_1 + 2\nu \varepsilon_1 \varepsilon_2 + \varepsilon^2_2 ) ,
\end{equation}
где $ E $ "--- модуль Юнга, $ \nu $ "--- коэффициент Пуассона. 
Здесь использован закон Гука, связывающий нормальные напряжения $ \sigma_1$, $\sigma_2 $ и~нормальные деформации $ \varepsilon_1$, $\varepsilon_2 $ и записанный в векторно"=матричной форме:
\begin{equation}\label{S_eq6}
\begin{pmatrix}
\sigma_1 \\ \sigma_2
\end{pmatrix}
 = C 
\begin{pmatrix} \varepsilon_1 \\ \varepsilon_2
\end{pmatrix}.
\end{equation}
В этом выражении 
$$
C = \frac{E}{(1-\nu^2)} \begin{pmatrix} 1 & \nu \\ \nu & 1 \end{pmatrix}
$$
можно трактовать как матричный коэффициент пропорциональности.

Формула \eqref{S_eq5} определяет кривые второго порядка на плоскости $\varepsilon_1 \varepsilon_2 $ 
(линии уровня потенциала $ V^e $), а~именно центральные эллипсы, главные оси которых наклонены к~декартовым осям $ \varepsilon_1$, $\varepsilon_2 $ под углом $ \pi/4 $. 
Большая ось располагается во втором и~четвертом квадрантах, малая "--- в~первом и~третьем. 
Для полуосей  выполняется отношение 
$$ \dfrac{a^{2}}{\rho^2} = \frac{(1-\nu)}{(1+\nu)}, $$ 
где $ a $ "--- малая полуось, $ \rho $ "--- большая полуось. 
Точки, расположенные на малой оси, отвечают равномерному растяжению $ (\varepsilon_1 = \varepsilon_2) $ или сжатию $ (-\varepsilon_1 = -\varepsilon_2) $, а~точки на большой оси определяют чистый сдвиг   
$ (\varepsilon_1 = -\varepsilon_2$, $ -\varepsilon_1 = \varepsilon_2 ) $. 
Произвольная точка  $ (\varepsilon_1,\varepsilon_2) $ лежит на линии уровня с большой полуосью, равной
\begin{equation}\label{S_eq7}
\rho = \frac{1}{\sqrt{2}}\sqrt{{(\varepsilon_1 + \varepsilon_2)}^2 \frac{1+\nu}{1-\nu} + {(\varepsilon_1 - \varepsilon_2)}^2} .
\end{equation}

При чистом сдвиге $ (\varepsilon_1 = \varepsilon$, $\varepsilon_2 = - \varepsilon)$  $\rho = \gamma /\sqrt{2}$, где    $ \gamma = \varepsilon_1 - \varepsilon_2 = 2 \varepsilon $ "--- сдвиг.
После исчерпания упругости для описания свойств материала можно было бы использовать теорию малых упруго-пластических деформаций [\citen{bib13, bib14}], в~которой предполагается упругость объемной деформации и~пропорциональность девиаторов тензоров напряжений и деформаций. 
Однако после выхода на стадию разупрочнения первая гипотеза не выполняется, так как материал уже существенно поврежден множеством микротрещин и объемный модуль изменяется. 

Чтобы остаться в рамках деформационной теории, следует отказаться от первой гипотезы, а~вместо второй сохранить пропорциональность напряжений и~деформаций с некоторым уже переменным, зависящим от деформаций матричным коэффициентом. 
Например, в рассматриваемой задаче %после перехода материала  при активном нагружении 
этот матричный коэффициент в законе \eqref{S_eq6} можно принять в~виде $ C \psi(\varepsilon_1,\varepsilon_2)$, 
где $ \psi $ "--- некоторая скалярная функция, требующая специального определения.

Следствием принятого предположения является тот факт, что в~области неупругости линии уровня функции энергии деформаций $ V $ (потенциальной энергии) подобны линиям уровня потенциала $ V^{e} $, то есть представляют собой эллипсы с~тем же самым отношением главных осей. 
Необходимо только знать распределение значений потенциала по  этим линиям уровня.

Воспользуемся полной диаграммой для чистого сдвига $ \tau(\gamma) $ ($ \tau $ "--- касательное напряжение, $ \gamma $ "--- деформация сдвига), которая обладает восходящей и нисходящей до нуля ветвями. 
Такую диаграмму материала можно получить, используя опытную диаграмму кручения стержней  круглого поперечного сечения с~пересчетом на диаграмму материала [\citen{bib15}]. 
Тогда энергия, затраченная на деформирование, есть   
$$ V = \int ^{\gamma}_{0}\tau(\gamma)  d\gamma .$$

После взятия интеграла и замены $ \gamma $ на $ \sqrt{2} \rho $  с~учетом выражения \eqref{S_eq7} получаем искомый единый потенциал $ V(\varepsilon_1,\varepsilon_2) $. 

Пусть $ \tau = 0.1G   \gamma(1-5 \gamma)$, если $ 0 \le \gamma \le 0.2 $, где  
$ G = \frac{E}{2(1+\nu)} $ "---  модуль сдвига в упругости. 
Тогда
\begin{equation}\label{S_eq8}
V(\varepsilon_1,\varepsilon_2) = 
\left\lbrace  
\begin{array}{cl}
\frac{V^{e}}{2}\Bigl( 0.1 - \frac{1}{3}\sqrt{\frac{V^{e}}{G}}\Bigr), & (\varepsilon_1,\varepsilon_2) \in \Omega , \\
0.002 G/3 , & (\varepsilon_1,\varepsilon_2) \notin \Omega.
\end{array}
\right. 
\end{equation}


\begin{figure}[b!]
\begin{tabular}{p{0.35\textwidth}p{0.6\textwidth}}

\includegraphics[width=0.3\textwidth]{fig_stru2} & 

\vspace{-5.5cm}

~~~\,Связь между напряжениями и деформациями (определяющие соотношения) в данном случае будет задана законом
$$
\begin{pmatrix} \sigma_1  \\ 
\sigma_2 
\end{pmatrix}
=
\begin{pmatrix} \frac{\partial V}{\partial \varepsilon_1} \\ 
\frac{\partial V}{\partial \varepsilon_2} \end{pmatrix}
=
C \psi (\varepsilon_1,\varepsilon_2) 
\begin{pmatrix} \varepsilon_1 \\ 
\varepsilon_2 \end{pmatrix},
$$
где 
$$ 
\psi (\varepsilon_1,\varepsilon_2) = \frac{1}{2} \Bigl(0.1 - \frac{1}{2} \sqrt{\frac{V^e}{G}} \Bigr).  
$$


~~~\,Отметим, что напряжения являются компонентами градиента функции~$ V $.

\\

\vspace{-.8cm}
\caption{Качественный вид поверхности потенциала $ V(\varepsilon_1,\varepsilon_2) $ } \label{Sfig2}
		\smallskip \footnotesize
		[Figure~\ref{Sfig2}. Qualitative form of the \centerline{potential surface $ V(\varepsilon_1,\varepsilon_2) $]} & \\
\end{tabular}
\end{figure}

%\begin{figure}[b!]
%	\begin{center}
%		\includegraphics[width=0.3\textwidth]{fig_stru2}
%		\caption{Качественный вид поверхности потенциала $ V(\varepsilon_1,\varepsilon_2) $ } \label{Sfig2}
%		\smallskip \footnotesize
%		[Figure~\ref{Sfig2}. Qualitative form of the potential surface $ V(\varepsilon_1,\varepsilon_2) $]
%	\end{center}
%\end{figure}



Очевидно, что область $ \Omega $ ограничена эллипсом с~полуосями 
$ \rho = 0.141$; $a = \rho\sqrt{(1+\nu)/(1-\nu)} $.  
Качественный вид поверхности $ V(\varepsilon_1,\varepsilon_2) $ изображен на рис.~\ref{Sfig2}. 
Функция $ V $ описывает все состояния материала (выпуклость вниз "--- устойчивость, упрочнение; 
выпуклость вверх "--- неустойчивость, разупрочнение).





\smallskip


\Section{Разрушение цилиндрического резервуара} Рассмотрим тонкостенный резервуар в~виде длинной трубы с~заделанными абсолютно жесткими крышками торцами. 
Радиус трубы равен $ b $, длина $ l $, толщина стенки $ t$  $(t\ll b)$. 
Резервуар находится под действием  монотонно и квазистатически возрастающего при постоянной температуре внутреннего давления с~интенсивностью~$ p $ [\citen{bib16, bib17}]. 
Материал обладает свойством деформационного разупрочнения. 
Применим изложенную выше методику для расчета предельного значения давления. 

Каждый элемент трубы находится в~плоском напряженном состоянии при всестороннем растяжении: 
$ \varepsilon_1 = \varepsilon_z $ "--- продольная деформация, 
$ \varepsilon_2 = \varepsilon_{\theta} $ "--- тангенциальная деформация. Его потенциальная функция $ V $ задана выражением \eqref{S_eq8}. 
%где  $ \varepsilon_1 = \varepsilon_z , \varepsilon_1 = \varepsilon_{\theta} $. 

Механическая система (резервуар) обладает двумя параметрами состояния (обобщенными координатами) $ \varepsilon_z$, $\varepsilon_{\theta} $ (однородное напряженно"=деформированное состояние) и~одним параметром управления $ p $. 

Лагранжиан (потенциальная функция) этой системы
$$
W = V(\varepsilon_z , \varepsilon_{\theta})2\pi b l t - p \pi b^2 l \varepsilon_z - 2\pi p b^2 l \varepsilon_{\theta}.
$$
Здесь первое слагаемое "--- полная энергия деформаций стенок трубы, второе и третье "--- взятые со знаком минус соответственно работа сил давления на перемещении жестких крышек и~перемещении, равном увеличению радиуса. 
Тогда уравнения равновесия, определяющие критические точки функции $ W $, имеют вид
\begin{equation}\label{S_eq9}
\frac{\partial W}{\partial \varepsilon_z} = \frac{\partial V}{\partial \varepsilon_z} 2\pi b l t  - p \pi b^2 l = 0;
\quad  
\frac{\partial W}{\partial \varepsilon_{\theta}} = \frac{\partial V}{\partial \varepsilon_{\theta}} 2\pi b l t  - 2p \pi b^2 l = 0.
\end{equation}
Уравнения \eqref{S_eq9} сводятся к одному:
$$
2 \frac{\partial V}{\partial \varepsilon_z} = \frac{\partial V}{\partial \varepsilon_{\theta}}.
$$
Производя необходимые вычисления, находим, что
\begin{equation}\label{S_eq10}
\varepsilon_{\theta} = \frac{2-\nu}{1-2\nu} \varepsilon_{z},
\end{equation}
то есть в пространстве деформаций путь деформирования элемента материала является прямой линией, уравнение которой задается равенством \eqref{S_eq10}. 
Кроме этого, из уравнений \eqref{S_eq9} следует, что  $ \sigma_z = pb/2t$, $\sigma_{\theta} = pb/t $. 
Отсюда путь нагружения в~пространстве напряжений также является прямой линией. 
Таким образом, имеем ситуацию с пропорциональным нагружением элемента материала.

Определим, наконец, предельное значение внутреннего давления. 
Как следует из приведенных выше теоретических положений, для этого необходимо гессиан функции $ W $ приравнять к~нулю, найти все решения полученного уравнения, по крайней мере, приближенно, и выделить из них вырожденные критические точки потенциала $ W $, которые принадлежат множеству критических точек (решений уравнений \eqref{S_eq9}) потенциала. 
Имеем 
$$
\Bigl\lbrace 
\frac{\partial^{2} V}{\partial \varepsilon^{2}_{z}} \cdot \frac{\partial^{2} V}{\partial \varepsilon^{2}_{\theta}} - 
\Bigl( \frac{\partial^{2} V}{\partial \varepsilon_{z}\partial \varepsilon_{\theta}}\Bigr) 
\Bigr\rbrace 
4 \pi^{2} b^2 l^2 t = 0.
$$
После несложных преобразований получаем
\begin{multline}\label{S_eq11}
\bigl(0.05 - \alpha \beta - \alpha \beta^{-1}(\varepsilon_{z} + \nu \varepsilon_{\theta})^2\bigr) 
\bigl(0.05 - \alpha \beta - \alpha \beta^{-1}(\varepsilon_{\theta} + \nu \varepsilon_{z})^2\bigr)-
\\
 - \nu^2 \bigl( 0.05 - \alpha \beta - \alpha \beta^{-1}(\varepsilon_{z} + \nu \varepsilon_{\theta}) (\varepsilon_{\theta} + \nu \varepsilon_{z})\bigr)^2 = 0
,                                          
\end{multline}
где 
$ \alpha=0.25 (1-\nu)^{ - 1/2}$, 
$ \beta = (\varepsilon^{2}_{z} + 2\nu \varepsilon_{z} \varepsilon_{\theta} + \varepsilon^{2}_{\theta})^{1/2}$;  
$ \varepsilon_{z}$, $\varepsilon_{\theta} \ge 0$, $(\varepsilon_{z}, \varepsilon_{\theta}) \in \Omega $.

Для вычисления вырожденных критических точек функции $ W $ применим следующую численную процедуру. 




Построим на плоскости  $ \varepsilon_{z} \varepsilon_{\theta} $ квадрат $ \Omega_1 $  такой, что в нем $ \varepsilon_{z}$, $\varepsilon_{\theta} > 0$,     $(\varepsilon_{z}, \varepsilon_{\theta}) \in \Omega \subset  \Omega_1 $ (рис.~\ref{Sfig3}). 
Для определения стороны этого квадрата найдем на осях $ \varepsilon_{\theta}$ и $\varepsilon_{z} $ равные отрезки, внутри которых детерминант матрицы Гессе меняет знак. 
Получаем отрезки длиной $0.1$. 
Эту величину и берем за сторону квадрата $\Omega_1$. 
Отметим, что данный прием может быть использован при построении $M \times N$-мерного куба в общей теории. 



\begin{figure}[h!]
\centering
	\includegraphics[width=0.4\textwidth]{fig_stru3}
	\caption{Вид кривой, в точках которой гессиан равен нулю} \label{Sfig3} 
	\smallskip \footnotesize
		[Figure~\ref{Sfig3}. The curve where the Hessian is equal to zero at all points]
\end{figure}

Разобьем этот квадрат прямоугольной сеткой с~достаточно малым шагом (на рис. \ref{Sfig3} показана укрупненная сетка). В~каждом узле сетки по формуле \eqref{S_eq11} при  $ \nu = 0.3 $ вычисляем детерминант (гессиан) матрицы Гессе. 
Выделяем те из них, где детерминант мало отличен от нуля, что определяется заданной точностью вычислений.

Производя расчеты, получаем множество точек, расположенных на кривой {\sl 1} (рис. \ref{Sfig3}). 
Затем координаты точек кривой {\sl 1} подставляем в уравнения \eqref{S_eq9} и~выделяем те из них, которые удовлетворяют уравнениям с~заданной точностью. 
В~результате находим только одну точку (точка $A$, рис. \ref{Sfig3}). Точка $ A $ лежит на прямой \eqref{S_eq10} и~имеет координаты $ \varepsilon_{z}  = 0.018$, $\varepsilon_{\theta}  = 0.076 $ (с~точностью до округления). 
Это и~есть искомая вырожденная критическая точка потенциальной функции $ W $. 
Подставляя ее координаты в~уравнения равновесия, находим критическое внутреннее давление, при котором происходит потеря устойчивости системы (катастрофа, разрушение). 
Если  $ E = 2 \cdot 10^5 $~МПа, $ b = 100 $ мм,  $ t = 2 $ мм, то критическое (предельное) значение давления $ p = 18 $~МПа. Заметим, что величина предельного давления увеличивается пропорционально толщине стенки.


\pagebreak



\AdditionalContent{Конкурирующие интересы}{Мы не имеем конкурирующих интересов.} 
%\AdditionalContent{Конкурирующие интересы}{Конкурирующих интересов не имею.} 
\AdditionalContent{Авторский вклад и ответственность}{Все авторы принимали участие в разработке концепции статьи и в написании рукописи. Все авторы несут полную ответственность за представление окончательной рукописи в печать. Окончательная версия рукописи была одобрена всеми авторами.} 
%\AdditionalContent{Авторская ответственность}{Я несу полную ответственность за предоставление окончательной версии рукописи в печать. Окончательная версия рукописи мною одобрена.} 
%\AdditionalContent{Авторский вклад и ответственность}{Все  авторы  принимали  участие  в  разработке  концепции  \mbox{статьи}  и  в  написании  рукописи. Авторы несут  полную  ответственность  за  предоставление  окончательной рукописи в печать. Окончательная  версия рукописи была одобрена всеми авторами.} 
\AdditionalContent{Финансирование}{Исследование не имело финансирования.}
%\AdditionalContent{Финансирование}{Исследование выполнялось без финансирования.}
%\AdditionalContent{Финансирование}{Работа выполнена при поддержке РФФИ (проект № xx–xx–xxxxx_a).}
%\AdditionalContent{Благодарность}{Выразите свою благодарность.}
%\AdditionalContent{Благодарность}{Авторы благодарны рецензенту за тщательное прочтение статьи и~ценные предложения и комментарии.}

\begin{thebibliography}{10}

\RBibitem{bib1}
\by Седов~Л.~И.
\book Механика сплошной среды
\bookvol 1
\yr 1970
\publ Наука
\publaddr М.
\totalpages 492
%\misctransl
%\by Sedov ~L.~I.
%\book Mekhanika sploshnoy sredy \rm [Continuum Mechanics] vol. 1 
%\yr 1970
%\publ Nauka
%\publaddr Moscow
%\lang In Russian
%\totalpages 492

\RBibitem{bib2}
\by Стружанов~В.~В., Миронов~В.~И.
\book Деформационное разупрочнение материала в элементах конструкций
\yr 1995
\publ УрО РАН 
\publaddr Екатеринбург
\totalpages 190
%\misctransl
%\by Struzhanov ~V.~V., Mironov ~V.~I.
%\book Deformatsionnoye razuprochneniye materiala v elementakh konstruktsiy \rm [Deformational Weakening of the Material in Structural Elements]
%\yr 1995
%\publ UrO RAN
%\publaddr Ekaterinburg
%\lang In Russian
%\totalpages 190

\RBibitem{bib3}
\by Вильдеман~В.~Э., Чаусов~Н.~Г.
\paper Условия деформационного разупрочнения материала при растяжении образца специальной конфигурации
\jour Заводская лаборатория. Диагностика материалов
\yr 2007
\vol 73
\issue 10
\pages 55--59
\elib{https://elibrary.ru/item.asp?id=9907854}
%\misctransl
%\lang In Russian
%\by Vildeman ~V.~E., Chausov ~N.~G.
%\paper Usloviya deformatsionnogo razuprochneniya materiala pri rastyazhenii obraztsa spetsial'noy konfiguratsii \rm [Conditions of Deformational Weakening of the Material with the Specimen of a Special Configuration under Tension]
%\jour Zavodskaya laboratoriya. Diagnostika materialov.
%\yr 2007
%\vol 73
%\issue 10
%\pages 55--59

\RBibitem{bib4}
\by Вильдеман~В.~Э., Третьяков~М.~П.
\paper Испытания материалов с построением полных диаграмм деформирования
\jour Проблемы машиностроения и надежности машин
\yr 2013
\vol 
\issue 2
\pages 93--98
\elib{https://elibrary.ru/item.asp?id=19548776}
%\misctransl
%\lang In Russian
%\by Vildeman ~V.~E., Tretyakov ~M.~P.
%\paper Ispytaniya materialov s postroyeniyem polnykh diagramm deformirovaniya \rm [Material Tests and Development of Complete Deformation Diagrams]
%\jour Problemy mashinostroyeniya i nadezhnosti mashin
%\yr 2013
%\vol 
%\issue 2
%\pages 93--98

\RBibitem{bib5}
\by Стружанов~В.~В., Коркин~А.~В.
\paper Об устойчивости процесса растяжения одной стержневой системы с~разупрочняющимися элементами
\jour Вестник Уральского государственного университета путей сообщения
\yr 2016
\issue 3(31)
\pages 4--17
\elib{https://elibrary.ru/item.asp?id=27249016}
\crossref{10.20291/2079-0392-2016-3-4-17}
%\misctransl
%\lang In Russian
%\by Struzhanov ~V.~V., Korkin ~A.~V.
%\paper Ob ustoychivosti protsessa rastyazheniya odnoy sterzhnevoy sistemy s razuprochnyayushchimisya elementami \rm [On the Stability of a Single Rod System with Weakening Elements under Tension]
%\jour Vestnik Ural'skogo gos. un–ta putey soobshcheniya
%\yr 2016
%\vol 
%\issue 3(31)
%\pages 4--17

\RBibitem{bib6}
\by Андреева~Е.~А.
\paper Решение одномерных задач пластичности для разупрочняющегося материала
\jour Вестн. Сам. гос. техн. ун-та. Сер. Физ.-мат. науки
\yr 2008
\issue 2(17)
\pages 152--160
\mathnet{http://mi.mathnet.ru/vsgtu642}
\crossref{10.14498/vsgtu642}


%\RBibitem{bib7}
%\by Горбунов~С.~В., Радченко~В.~П.
%\paper Вариант решения краевой задачи неупругого деформирования разупрочняющейся среды и его экспериментальная проверка
%\jour Материалы Х Всероссийской конференции по механике деформируемого твёрдого тела (18–22 сентября 2017 г., Самара, Россия)
%\yr 2017
%\vol 1
%\issue
%\pages 179--181
%\misctransl
%\lang In Russian
%\by Gorbunov ~S.~V., Radchenko ~V.~P.
%\paper Variant resheniya krayevoy zadachi neuprugogo deformirovaniya razuprochnyayushcheysya sredy i yego eksperimental'naya proverka \rm [A Solution Case of the Boundary-Value Problem for Inelastic Deformation of a Weakening Medium and its Experimental Verification]
%\jour Materialy X  Vserossiyskoy konferentsii po mekhanike deformiruyemogo tvordogo tela.  
%\yr 2017
%\vol 1
%\issue
%\pages 179--181

\RBibitem{bib8}
\by Арсенин~В.~Я.
\book Методы математической физики и специальные функции
\yr 1974
\publ Наука
\publaddr М.
\totalpages 431
%\misctransl
%\by Arsenin ~V. ~Ya.
%\book Metody matematicheskoy fiziki i spetsial'nyye funktsii. \rm [Methods of Mathematical Physics and Special Functions]
%\yr 1974
%\publ Nauka
%\publaddr Moscow
%\lang In Russian
%\totalpages 431

\RBibitem{bib9}
\by Стружанов~В.~В., Коркин~А.~В.
\paper О решении нелинейных уравнений равновесия  одной растягиваемой стержневой системы с разупрочняющимися элементами методом простых итераций
\jour Вестник Уральского государственного университета путей сообщения
\yr 2017
\issue 2(34)
\pages 4--16
\crossref{10.20291/2079-0392-2017-2-4-16 }
\elib{https://elibrary.ru/item.asp?id=29409142}
%
%\misctransl
%\lang In Russian
%\by Struzhanov ~V.~V., Korkin ~A.~V.
%\paper O reshenii nelineynykh uravneniy ravnovesiya odnoy rastyagivayemoy sterzhnevoy sistemy s razuprochnyayushchimisya elementami metodom prostykh iteratsiy \rm [The solution of Nonlinear Equilibrium Equations for a Single Rod System with Weakening Elements under Tension Using the Simple Iterations Method]
%\jour Vestnik Ural'skogo gos. un-ta putey soobshcheniya
%\yr 2017
%\vol 
%\issue 2(34)
%\pages 4--16

\Bibitem{bib10}
\by Poston~T., Stewart~I.
\book Catastrophe Theory and Its Applications
\yr 1978
\publ Pitman
\serial Surveys and Reference Works in Mathematics
\vol 2
\zmath{0382.58006}
\publaddr London, San Francisco, Melbourne
\totalpages xviii+491 

\Bibitem{bib11}
\by Gilmore~R.
\book Catastrophe theory for scientists and engineers
\zmath{0497.58001}
\serial A Wiley-Interscience Publication
\publaddr New York 
\publ John Wiley \& Sons
\totalpages xvii+666 
\yr 1981

\Bibitem{bib12}
\by Pars~L.~A.
\book A treatise on analytical dynamics
\zmath{0125.12004}
\publaddr London
\publ Heinemann Educational Books 
\totalpages xxi+641 
\yr 1965
 
\RBibitem{bib13}
\by Ильюшин~А.~А.
\paper Об основах общей математической теории пластичности
\jour Вестник Московского университета. Серия 1. Математика. Механика
\yr 1961
\issue 3
\pages 31--36
\zmath{0100.20304}
%\misctransl
%\lang In Russian
%\by Ilyushin ~A.~A.
%\paper Ob osnovakh obshchey matematicheskoy teorii plastichnosti \rm [Fundamentals of the General Mathematical Theory of Plasticity]
%\jour Voprosy teorii plastichnosti. Izd-vo Akademii nauk SSSR
%\yr 1961
%\vol 
%\issue 
%\pages 3--29

\RBibitem{bib14}
\by Качанов~Л.~М.
\book Основы теории пластичности
\yr 1969
\publ Наука
\publaddr М.
\totalpages 420
%\misctransl
%\by Kachanov ~L.~M.
%\book Osnovy teorii plastichnosti \rm [Fundamentals of the Theory of Plasticity]
%\yr 1969
%\publ Nauka
%\publaddr Moscow
%\lang In Russian
%\totalpages 420

\RBibitem{bib15}
\by Стружанов~В.~В.
\paper Определение диаграммы деформирования материала с~падающей ветвью по диаграмме кручения цилиндрического образца
\jour Сиб. журн. индустр. матем.
\yr 2012
\vol 15
\issue 1
\pages 138--144
\mathnet{http://mi.mathnet.ru/sjim718}

\RBibitem{bib16}
\by Радченко~В.~П.
\book Введение в механику деформируемых систем
\yr 2009
\publ СамГТУ
\publaddr Самара
\totalpages 241
%\misctransl
%\by Radchenko ~V.~P.
%\book Vvedeniye v mekhaniku deformiruyemykh system \rm [Introduction to Mechanics of De-formation Systems]
%\yr 2009
%\publ Samarskiy gos. tekhn. un-t.
%\publaddr Samara
%\lang In Russian
%\totalpages 241

\Bibitem{bib17}
\by Southwell~R.~V.
\book An introduction to the theory of elasticity. For engineers and physicists
\zmath{62.1529.04}
\totalpages ix+510 
\publaddr Oxford
\publ Clarendon Press
\serial Oxford Engineering Science Series
\yr 1936


\end{thebibliography}


\newpage

\makeentitle

\vspace{-5mm}

\AdditionalContent{Competing interests}{We have no competing interests.} 
%\AdditionalContent{Competing interests}{I have ... (or) no competing interests.} 
%\AdditionalContent{Competing interests}{We have ... (or) no competing interests.} 
\AdditionalContent{Authors' contributions and responsibilities}{Each authors has participated in the article concept development  and in the manuscript writing. The authors are absolutely responsible for submitting the final manuscript in print. Each authors has approved the final version of manuscript.}
%\AdditionalContent{Author's Responsibilities}{I take full responsibility for submitting the final manuscript in print. I approved the final version of the manuscript.} 
%\AdditionalContent{Authors' contributions and responsibilities}{Each author has participated in the article concept development and in the manuscript writing. The authors are absolutely responsible for submitting the final manuscript in print. Each author has approved the final version of manuscript.} 

\AdditionalContent{Funding}{This research received no specific grant from any funding agency in the public, commercial, or not-for-profit sectors.}




\begin{thebibliography}{10}

\Bibitem{bib1:1}
\lang In Russian
\by Sedov~L.~I.
\book Mekhanika sploshnoi sredy \rm [Continuum Mechanics] 
\bookvol 1
\yr 1970
\publ Nauka
\publaddr Moscow
\totalpages 492

\Bibitem{bib2:1}
\lang In Russian
\by Struzhanov~V.~V., Mironov~V.~I.
\book Deformatsionnoe razuprochnenie materiala v elementakh konstruktsii \rm [Deformational Softening of Material in Structural Elements]
\yr 1995
\publ UrO RAN 
\publaddr Ekaterinburg
\totalpages 190

\Bibitem{bib3}
\lang In Russian 
\by Vil'deman~V.~E., Chausov~N.~G.
\paper Conditions of strain softening upon stretching of the specimen of special configuration
\jour Zavodskaia laboratoriia. Diagnostika materialov
\yr 2007
\vol 73
\issue 10
\pages 55--59


\Bibitem{bib4:en}
\paper Material testing  by plotting total deformation curves
\by Vil'deman~V.~E., Tretyakov~M.~P.
\jour J. Mach. Manuf. Reliab.
\yr  2013
\vol 42
\issue 2
\pages 166--170
\elib{https://elibrary.ru/item.asp?id=20435521}
\crossref{10.3103/S1052618813010159}



\Bibitem{bib5:1}
\lang In Russian
\by Struzhanov~V.~V., Korkin~A.~V.
\paper Regarding stretching process stability of one bar system with softening elements
\jour Vestnik Ural'skogo gosudarstvennogo universiteta putei soobshcheniia
\yr 2016
\issue 3(31)
\pages 4--17
\crossref{10.20291/2079-0392-2016-3-4-17}


\Bibitem{bib6}
\lang In Russian
\by Andreeva~E.~A.
\paper Solution of one-dimensional softening materials plasticity problems
\jour Vestn. Samar. Gos. Tekhn. Univ., Ser. Fiz.-Mat. Nauki \rm [J. Samara State Tech. Univ., Ser. Phys. Math. Sci.]
\yr 2008
\issue 2(17)
\pages 152--160
\crossref{10.14498/vsgtu642}



\Bibitem{bib8:1}
\lang In Russian
\by Arsenin~V.~Ya.
\book Metody matematicheskoi fiziki i spetsial'nye funktsii \rm [Methods of mathematical physics and special functions]
\yr 1974
\publ Nauka
\publaddr Moscow
\totalpages 431


\Bibitem{bib9:1}
\lang In Russian
\by Struzhanov~V.~V., Korkin~A.~V.
\paper On the solution of non-linear equations of trim of a single expandable frame structure with softening elements by the method of simple iterations
\jour Vestnik Ural'skogo gosudarstvennogo universiteta putei soobshcheniia
\yr 2017
\issue 2(34)
\pages 4--16
\crossref{10.20291/2079-0392-2017-2-4-16}


\Bibitem{bib10:1}
\by Poston~T., Stewart~I.
\book Catastrophe Theory and Its Applications
\yr 1978
\publ Pitman
\serial Surveys and Reference Works in Mathematics
\vol 2
\zmath{0382.58006}
\publaddr London, San Francisco, Melbourne
\totalpages xviii+491 

\Bibitem{bib11:1}
\by Gilmore~R.
\book Catastrophe theory for scientists and engineers
\zmath{0497.58001}
\serial A Wiley-Interscience Publication
\publaddr New York 
\publ John Wiley \& Sons
\totalpages xvii+666 
\yr 1981

\Bibitem{bib12:1}
\by Pars~L.~A.
\book A treatise on analytical dynamics
\zmath{0125.12004}
\publaddr London
\publ Heinemann Educational Books 
\totalpages xxi+641 
\yr 1965
 
\Bibitem{bib13:1}
\lang In Russian
\by Ilyushin~A.~A.
\paper On the foundations of the general mathematical theory of plasticity
\jour Vestn. Mosk. Univ., Ser.~I 
\yr 1961
\issue 3
\pages 31--36



\Bibitem{bib14:1}
\by Kachanov~L.~M.
\book Foundations of the theory of plasticity
\zmath{0231.73015}
\serial North-Holland Series in Applied Mathematics and Mechanics
\vol 12
\publaddr Amsterdam
\publ North-Holland Publ.
\totalpages xiii+482 
\yr 1971


\Bibitem{bib15:!}
\by Struzhanov~V.~V.
\paper The determination of the deformation diagram of a~material with a~falling branch using the torsion diagram of a~cylindrical sample
\jour Sib. Zh. Ind. Mat.
\yr 2012
\vol 15
\issue 1
\pages 138--144

\Bibitem{bib16:1}
\lang In Russian
\by Radchenko~V.~P.
\book Vvedenie v mekhaniku deformiruemykh sistem \rm [Introduction to the mechanics of deformable systems]
\yr 2009
\publ Samara State Technical Univ.
\publaddr Samara
\totalpages 241


\Bibitem{bib17:1}
\by Southwell~R.~V.
\book An introduction to the theory of elasticity. For engineers and physicists
\zmath{62.1529.04}
\totalpages ix+510 
\publaddr Oxford
\publ Clarendon Press
\serial Oxford Engineering Science Series
\yr 1936


\end{thebibliography}


%\end{document}
 \newpage

%\input{preamble}

%

\vsgtu{1639} %registration number

\FirstPage{774}

\LastPage{784}

%

%\usepackage{samgtu-name}

%

%\pdfcrop

%\draft

%

\ShortCommunication

%

%\begin{document}





$\;$

\vspace{1mm}



\hfill \qrcode[hyperlink,height=.8in]{https://doi.org/10.14498/vsgtu1639}

\vspace{-.8in}



\udc{517.955, 517.968.7}

\msc{26A33, 35K15, 35R11}



\rutitle{К проблеме единственности решения задачи Коши для~уравнения дробной диффузии с оператором Бесселя}

\rutitleshort{К проблеме единственности решения задачи Коши\dots}

\rutitlehack{К проблеме единственности решения \\ задачи Коши для~уравнения дробной диффузии с~оператором Бесселя}



\entitle{On the uniqueness of the solution of the Cauchy problem for the equation of fractional diffusion with~Bessel~operator}

\entitleshort{On the uniqueness of the solution of the Cauchy problem\dots}



\ruauthor{Ф.~Г.~Хуштова}{Х\,у\,ш\,т\,о\,в\,а~Ф.~Г.}

\enauthor{F.~G.~Khushtova}{K\,h\,u\,s\,h\,t\,o\,v\,a~F.~G.}



\ruaffil{\runiipma.}



\enaffil{\enniipma.}



\ruauthorsdetails{1}{% Количество авторов

\fullauthor{\rupid{53181}{Фатима Гидовна Хуштова}}

\CAuthor % Автор-корреспондент

\orcid{0000-0003-4088-3621} % ORCID ID автора 

%\regalia{без ученой степени} 

\position[n]{научный сотрудник} 

\dept{отдел дробного исчисления} 

\email{khushtova@yandex.ru}

}



\enauthorsdetails{1}{

\fullauthor{\enpid{53181}{Fatima G. Khushtova}}

\CAuthor % Сorresponding Author

\orcid{0000-0003-4088-3621} % ORCID ID

%\regalia{Cand. Phys. \& Math. Sci., Associate Professor} 

\position[n]{Researcher} 

\dept{Dept. of Fractional Calculus}

\email{khushtova@yandex.ru}

}



\rukeywords{уравнение дробной диффузии, оператор дробного дифференцирования, оператор Бесселя, задача Коши, единственность решения, условие Тихонова, функция Райта}

\enkeywords{fractional diffusion equation, fractional differentiation operator, Bessel operator, Cauchy problem, solution uniqueness, Tikhonov condition,	 Wright function}



\ruabstract{Рассматривается уравнение дробной диффузии с сингулярным оператором Бесселя, действующим по пространственной переменной, и оператором дробного дифференцирования Римана--Лиувилля, действующим по временной переменной. 

Когда порядок дробной производной равен единице, а особенность у оператора Бесселя отсутствует, рассматриваемое уравнение совпадает с~классическим уравнением теплопроводности. 

Ранее для уравнения дробной диффузии с оператором Бесселя было построено решение задачи Коши и доказана теорема единственности решения в классе функций экспоненциального роста. 



Построен пример, показывающий, что увеличение показателя степени в условии, гарантирующем единственность решения задачи Коши, влечет за собой неединственность решения. 

С помощью известных свойств функции Райта получены оценки для построенной функции.

Показывается, что она, будучи не равной тождественно нулю, удовлетворяет однородному уравнению и однородному условию Коши. 

}





\enabstract{In this paper, we consider fractional diffusion equation involving the Bessel operator acting with respect to a spatial variable and the Riemann-Liouville fractional differentiation operator acting with respect to a time variable. When the order of the fractional derivative is unity, and the singularity of the Bessel operator is absent, this equation coincides with the classical heat equation. Earlier, a solution of the Cauchy problem has been considered for the considered equation and a uniqueness theorem has been proved for a class of functions satisfying the analog of the Tikhonov condition. 



In this paper, we have constructed an example to show that the exponent (power) at the condition of the uniqueness of the solution to the Cauchy problem cannot be raised under. Its increase leads to a non-uniqueness of the solution. Using the well-known properties of the Wright function, we have obtained estimates for constructed function, which satisfies the homogeneous equation and the zero Cauchy condition.

}



\newdate{khusha}{28}{08}{2018}

\date{\displaydate{khusha}}

\newdate{khushb}{25}{10}{2018}

\revisiondate{\displaydate{khushb}}

\newdate{khushc}{12}{11}{2018}

\accepteddate{\displaydate{khushc}}



\newdate{khushe}{28}{11}{2018}

\renewcommand{\JournalDataOnline}{\displaydate{khushe}}



%\ForPeerReview

\makerutitle



\newpage



\Section[n]{Введение}

В области $\Omega=\{(x,y): -\infty<x<\infty,\ 0<y<T \}$ рассмотрим уравнение

\begin{equation}\label{GL3_1}

B_xu(x,y)-D^{\alpha}_{0y}u(x,y)=0, 

\end{equation}

где $D_{0y}^{\alpha}$ "--- оператор дробного дифференцирования в смысле  

Римана"--~Лиувилля  порядка $\alpha,$ определяемый равенствами~\cite{Drob, Pskhu}   

$$D_{0y}^{\alpha} \, g(y)= \frac{dg}{dy},$$ если $\alpha=1$, и

$$ 

D_{0y}^{\alpha}\,g(y)=\frac{1}{\Gamma(1-\alpha)}\frac{d}{dy}

\int _{0}^{y}\frac{g(t)}{(y-t)^{\alpha}}dt,

$$ 

если $0<\alpha<1;$

$$

B_x=x^{-b} \frac{d}{dx}\left(x^b \frac{d}{dx}\right)=\frac{d^2}{d x^2}+\frac{b}{x} \frac{d}{d x}

$$

"--- оператор Бесселя, $|b|<1.$



Уравнение \eqref{GL3_1} при $\alpha=1,$ то есть уравнение

\begin{equation}\label{UrAlpha-1}

u_{xx}(x,y)+\frac{b}{x}u_{x}(x,y)-u_{y}(x,y)=0,

\end{equation}

совпадает с %~названным И.~А.~Киприяновым 

$B$-параболическим уравнением.

Краевые задачи в ограниченной и~неограниченной областях для этого уравнения при различных значениях параметра $b$ рассматривали многие авторы. 

Например, в работе V.~Alexia\-des~\cite{Alexiades} для него решены первая, вторая и~третья краевые задачи в~области с~подвижной границей, в~работе D.~Calton\cite{Calton} для него  исследована задача Коши.

 

 Уравнение \eqref{UrAlpha-1} заменой $z=x^2/(4n)$ сводится к~уравнению

\begin{equation} \label{UrKep2}

u_{zz}(z,y)+\frac{m+1}{z}\,u_{z}(z,y)-\frac{n}{z} u_{y}(z,y)=0,\quad m=(b-1)/2,

\end{equation}

для которого S.~K\c{e}pi\'nski  \cite{Kepinski_2} исследовал краевую задачу в~полуполосе $x>0,$ $0<y<y_0.$

Уравнение \eqref{UrAlpha-1} было также объектом исследования  монографий С.~А.~Терсенова~\cite{Tersenov} и~М.~И.~Матийчука~\cite{Matiychuk}.



Дифференциальные уравнения, содержащие оператор Бесселя, наиболее подробно исследованы в~работах И.~А.~Киприянова и~его учеников \cite{Kipriyanov-Monographiya,  Kipriyanov}.

Отметим здесь, что очень важные результаты по применению операторов преобразования в~теории дифференциальных уравнений с~особенностями в~коэффициентах, в~частности с~операторами Бесселя, получены В.~В.~Катраховым и~С.~М.~Ситником~\cite{KatrakhovSitnik, Sitnik}.





В случае, когда $b=0,$ $0<\alpha<2,$ уравнение \eqref{GL3_1} совпадает с диффузионно"=волновым уравнением. 

Различные краевые задачи для него, а также для многомерных его обобщений подробно исследованы в работах многих авторов.  Приведем некоторые из них. 



А.~Н.~Кочубей \cite{Kochubei_2} исследовал задачу Коши для уравнения диффузии дробного порядка с~регуляризованной дробной производной (производной Капуто) и~эллиптическим оператором с~непрерывными ограниченными вещественными коэффициентами

$$

L=\sum\limits_{i,j=1}^{n}a_{i,j}(x)\frac{\partial^2}{\partial x_i \partial x_j}+

\sum\limits_{j=1}^{n}b_{j}(x)\frac{\partial}{\partial x_j}+c(x) u.

$$

В случае, когда $L=\Delta$ "--- оператор Лапласа, фундаментальное решение построено в терминах $H$-функции Фокса. 

Исследованию задачи Коши для диффузионного и диффузионно"=волнового уравнений дробного порядка с~производной Капуто  посвящены также  работы А.~Н.~Кочубея и С.~Д.~Эйдельмана \cite{Kochubei_1, Kochubei_3, Kochubei_4, Kochubei_5}.



А.~В.~Псху~\cite{Pskhu2} построил фундаментальное решение и исследовал задачу Коши  для многомерного диффузионно"=волнового уравнения с~оператором дробного дифференцирования Джрбашяна"--~Нерсесяна, который  включает в~себя как частный случай операторы дробного дифференцирования Римана"--~Лиувилля и Капуто.



В~работах  \cite{Pskhu-SibirJournal, Pskhu3} также построено фундаментальное решение и исследована задача Коши для уравнения дробной диффузии соответственно с~оператором дискретно распределенного дифференцирования и оператором Джрбашяна"--~Нерсесяна, действующим по многим переменным времени.

В работе \cite{Pskhu4} решена первая краевая задача в~нецилиндрической области для диффузионно"=волнового уравнения с~оператором Джрбашяна"--~Нерсесяна. 

Доказана теорема существования и~единственности исследуемой задачи, конструктивно построено ее решение. 



А.~А.~Ворошилов и А.~А.~Килбас \cite{VorKilbas} методом интегральных преобразований построили решение задачи Коши для уравнения

\begin{equation} \label{MnogPerem}

D^{\alpha}_{0t}\,u(x,t)=\lambda^2 \Delta_x u(x,t), \quad  x\in{\mathbb R}^m,\quad  t>0, 

\end{equation}

где $\Delta_x=\sum_{j=1}^{m}\partial^2/\partial x_{j}^{2},$ $n-1<\alpha <n,$ $n\in\mathbb{N}.$ 

При $0<\alpha\leqslant 1$ и $1<\alpha<2$ решения  выписаны в терминах $H$-функции Фокса.

В работе~\cite{VorKilbas2} было также  построено решение задачи Коши в~случае, когда в уравнении

\eqref{MnogPerem} вместо оператора Римана"--~Лиувилля содержится оператор Капуто. 



Задача Коши для уравнения \eqref{MnogPerem} при $\lambda=1,$ 

$0<\alpha\leqslant 1$  была ранее рассмотрена С.~Х.~Геккиевой~\cite{Gekk}, а~в~работе \cite{Gekkieva1} ею исследована первая краевая задача в~первом квадранте для уравнения \eqref{MnogPerem} при $m=\lambda=1,$ 

$0<\alpha\leqslant 1.$  







Уравнение диффузии дробного порядка и некоторые его обобщения рассматривали 

F.~Mainardi, G.~Pagnini, Yu.~Luchko, P.~Paradisi~\cite{Mainardi-1, Mainardi-3, Mainardi-2, Pagnini_1,   Pagnini_Paradisi}  и~многие другие. 

Более подробную библиографию по данному вопросу можно найти в~монографиях \cite{Pskhu, MathaiSaxenaHaubold, Podlubny}.

 

%Чтобы подготовить Вашу статью для рецензента, т.~е. обезличить её, воспользуйтесь командой \verb"\ForPeerReview" перед командой 

%\verb"\makerutitle".



\smallskip





\Section{Постановка задачи Коши и теорема единственности решения}

Пусть $\Omega^{+}=\Omega\cap\{x>0 \},$ $\Omega^{-}=\Omega\cap\{x<0\},$ 

$\bar{\Omega}$ "--- замыкание области $\Omega.$ 



\smallskip



{\sc\small Определение.}

{\it Регулярным  решением} уравнения \eqref{GL3_1} в области $\Omega$ 

будем называть функцию

$u=u(x,y),$ удовлетворяющую уравнению \eqref{GL3_1} в области

$\Omega^{+}\cup\Omega^{-},$ и такую, что $y^{\,1-\alpha}u\in{C(\bar{\Omega})},$ 

$|x|^b u_{x}\in{C(\Omega)},$ $u_{xx},$ $D_{\,0y}^{\,\alpha} u\in{C(\Omega^{+}\cup\Omega^{-})}.$



\smallskip



{\sc\small Задача Коши.}

{\it Найти регулярное в области $\Omega$ решение

уравнения} \eqref{GL3_1}, {\it удовлетворяющее условию

\begin{equation}\label{UslSMCoshi}

\lim\limits_{y \rightarrow 0} D_{0y}^{\alpha-1}u(x,y)=\varphi(x), \quad  -\infty<x<\infty,

\end{equation}

где $\varphi(x)$ "--- заданная функция}.



\smallskip



В работе \cite{KhushtovaVestn2} доказана следующая теорема единственности.



\smallskip



\begin{newthm}{Теорема}  

Существует не более одного регулярного решения задачи Коши {\rm \eqref{GL3_1}, \eqref{UslSMCoshi} } в классе функций$,$ удовлетворяющих для некоторой положительной постоянной $k$ условию

\begin{equation} \label{AnalogTikhonovaInfty}

\lim\limits_{|x|\to \infty}y^{1-\alpha}\,u(x,y)\exp{\big(-k\,|x|^{\,\frac{2}{2-\alpha}} \big)}=0.

\end{equation}  

\end{newthm}



\smallskip





Отметим, что условие \eqref{AnalogTikhonovaInfty} имеет место и в случае задачи Коши для уравнения дробной диффузии.  

Теоремы единственности решения задачи Коши в~классе ограниченных функций  и~в~классе функций экспоненциального роста для многомерного уравнения дробной диффузии с~оператором Капуто были доказаны в работе А.~Н.~Кочубея \cite{Kochubei_2},

теорема единственности решения задачи Коши в классе функций экспоненциального роста для многомерного диффузионно"=волнового уравнения с~оператором Джрбашяна"--~Нерсесяна  "--- в~работе А.~В.~Псху~\cite{Pskhu2}. 

В этих же работах построены примеры, показывающие, что увеличение показателя $2/(2-\alpha)$ ведет к~утрате единственности решения соответствующих задач. 

При $\alpha=1$ условие \eqref{AnalogTikhonovaInfty} совпадает с~хорошо известным условием Тихонова 

$$

\lim\limits_{|x|\to \infty}u(x,y)\exp{\big(-k\,|x|^{2} \big)}=0.

$$

Как известно, пример неединственности решения задачи Коши для уравнения теплопроводности был впервые построен А.~Н.~Тихоновым в~основополагающей работе~\cite{tikhonovUsl-1} (см.~также \cite{tikhonovUsl-2}).



В данной работе эти идеи развиваются применительно к~задаче \eqref{GL3_1}, \eqref{UslSMCoshi}.



\smallskip



\Section{Неулучшаемость показателя степени в условии единственности решения задачи Коши}

Рассмотрим функцию

\begin{equation} \label{T(z,zeta)}  

T_{\mu}(z,\zeta)=\sum\limits_{n=0}^{\infty}\frac{z^{\,\beta+n}}{\,n!\,\Gamma(1+\beta+n)}\,\phi\left(-\rho,\mu-\alpha\beta-\alpha n; -\zeta\right),

\end{equation}    

где  $\phi\,(-\rho,\delta;z)$ "--- {\it функция Райта}, определяемая степенным рядом \cite{Pskhu, GorenfloLuchkoMainardi}

$$

\phi\,(-\rho,\delta;z)=\sum\limits_{n=0}^{\infty}\frac{z^{n}}{\,n!\,\Gamma(\delta-\rho\, n)}.

$$



Далее будем считать, что $\rho\in(0;1).$ 

Для любых $\nu$, $\delta\in \mathbb{R}$ справедливы формулы  \cite[с.~26, 44]{Pskhu}

\begin{gather}\label{diff_x}

\frac{d}{dz}\phi(-\rho,\delta;-z)=-\phi(-\rho,\delta-\rho;-z),

\\

\label{diff_Right}

D_{ay}^{\nu} \,y^{\delta-1}\phi(-\rho,\delta;-c\,y^{-\rho})=y^{\delta-\nu-1}\phi(-\rho,\delta-\nu;-c\,y^{-\rho}).

\end{gather}



\smallskip



\begin{lemma}[1]Если $\delta\geqslant 1,$ то для любого положительного $z$ справедливо нера\-вен\-ство~{\rm \cite[с.~47]{Pskhu} }

\begin{equation}\label{Right_delta>1}

\phi\,(-\rho,\delta;-z)\leqslant \frac{1}{\Gamma\left( \delta \right)}

\exp\left[-(1-\rho)\,\rho^{\frac{\rho}{1-\rho}}\, z^{\frac{1}{1-\rho}} \right].

\end{equation}

\end{lemma}



\smallskip



\begin{lemma}[2]Если $\delta<1,$ то для любых положительных $x$ и $y,$ 

$a\in(0;x),$ $\omega\in(1/2; \omega_0),$ $\omega_0=\min\{1,1/(2\rho)\},$

справедливы неравенства~{\rm \cite[с.~49]{Pskhu}}

\begin{gather}\label{Right-0}

\Bigl|\,y^{\,\delta-1}\phi\Bigl(-\rho,\delta;-\frac{x}{y^{\rho}}\Bigr)\Bigr|

\leqslant 

\frac{\Gamma\left( (1-\delta)/\rho\right)}{\,\rho\, \pi \left[aC_{\rho}(\rho,\omega) \right]^{(1-\delta)/\rho}}\,

\phi\Bigl(-\rho,1;-\frac{x-a}{y^{\rho}}\Bigr),

\\

\label{Right-2}

\Bigl|y^{\,\delta-1}\phi\Bigl(-\rho,\delta;-\frac{x}{y^{\rho}}\Bigr)\Bigr|

\leqslant \frac{\Gamma\left( 1-\delta\right)}{\pi \left[yC_{\rho}(1,\omega) \right]^{1-\delta}},

\end{gather}

где 

$C_{\rho}(\rho,\omega)=\frac1{\rho\pi}{\cos \rho\omega\pi},$ $C_{\rho}(1,\omega)=-\frac1{\pi}{\cos \omega\pi}.

$

\end{lemma}



\smallskip



Как следует из \eqref{Right-2}, ряд \eqref{T(z,zeta)} сходится абсолютно для любых  $z\in\mathbb{C},$ $\zeta>0,$ $\rho, \alpha, \beta\in(0;1),$ $\mu\in\mathbb{R}.$



\smallskip



Докажем следующую лемму.



\smallskip

\begin{lemma}[3]Если $0<\alpha<\rho<1,$ $\mu\geqslant 0,$ то при любых $x\in\mathbb{R},$ $y>0$ справедливо неравенство

\begin{equation} \label{T_muOtsenka}  

\Bigl| y^{\mu-1} T_{\mu}\Bigl(\frac{x^2}{4y^{\alpha}},\frac{1}{y^{\rho}}\Bigr)\Bigr|

\leqslant 

\frac{C

|x|^{\frac{2[1-\rho(1-\beta)]}{2\rho-\alpha}} y^{\mu}}{\Gamma(\mu+1)}\exp\left[ k_1  |x|^{\frac{2\rho}{2\rho-\alpha}}-k_2   y^{-\frac{\rho}{1-\rho}}\right],

\end{equation}    

где $C,$ $k_1,$ $k_2$ "--- положительные постоянные$,$ зависящие только от 

$\alpha,$ $\beta$ и $\rho.$

\end{lemma}



\smallskip



\begin{proof}

Из \eqref{Right-0} следует оценка

$$

\Bigl| y^{-\alpha\beta-\alpha n-1} \phi\Bigl(-\rho,-\alpha\beta-\alpha n; -\frac{1}{y^{\rho}}\Bigr)\Bigr|

\leqslant

 \frac{\Gamma\left( 1/\rho+\varepsilon\beta+\varepsilon n\right)}{ \rho \pi\left(aC_{\rho} \right)^{1/\rho+\varepsilon\beta+\varepsilon n}}

\phi\Bigl(-\rho,1; -\frac{1-a}{y^{\rho}}\Bigr),

$$

где $\varepsilon=\alpha/\rho,$ $a\in(0;1),$ $C_{\rho}=C_{\rho}(\rho,\omega).$



С помощью последней оценки получим

\begin{multline}

\Bigl| y^{\,\mu-1} T_{\mu}\Bigl(\frac{x^2}{4y^{\alpha}},\frac{1}{y^{\rho}}\Bigr)\Bigr|=

\\

=\left|

 D_{0y}^{-\mu}

  \sum\limits_{n=0}^{\infty}\frac{|x|^{2\beta+2n}}{ 2^{2\beta+2n}n! \Gamma(1+\beta+n)}

y^{-\alpha\beta-\alpha n-1}\phi \Bigl(-\rho,-\alpha\beta-\alpha n; -\frac{1}{y^{\rho}}\Bigr)

 \right|\leqslant

\\

 \label{WrightRyadZ}  

\leqslant

\frac{ |x|^{2\beta}y^{\mu}\phi\bigl(-\rho,\mu+1; -\frac{1-a}{y^{\rho}}\bigr)}{2^{2\beta}\rho \pi \left(aC_{\rho} \right)^{1/\rho+\varepsilon\beta}}

\sum\limits_{n=0}^{\infty}

\frac{\Gamma\left(1/\rho+\varepsilon\beta+\varepsilon n\right)}{n!

\Gamma(1+\beta+n)}

\Bigl( \frac{x^2}{4\left(aC_{\rho} \right)^{\varepsilon}}\Bigr)^{n}.

\end{multline}



Степенной ряд, стоящий справа от неравенства \eqref{WrightRyadZ}, представляет собой 

обобщенную функцию Райта

$$

{_1\Psi_{1}}\left[ \,z \,\bigg|

\begin{array}{ll}

 \big(1/\rho+\varepsilon\beta, \,\varepsilon\,\big)\\

\big(\,1+\beta, 1\,\big)\\

\end{array}

\right]=

\sum\limits_{n=0}^{\infty}

\frac{\Gamma\left(1/\rho+\varepsilon\beta+\varepsilon n\right)}{n!

\,\Gamma(1+\beta+n)} z^{n}, \; \; z=\frac{x^{\,2}}{4(aC_{\rho})^{\varepsilon}}.

$$

Из ее асимптотических свойств при $z\to \infty$ следует~\cite{Wright}

$$

{_1\Psi_{1}}\left[ \,z\,\bigg|

\begin{array}{ll}

 \big(1/\rho+\varepsilon\beta, \,\varepsilon\,\big)\\

\big(\,1+\beta, 1\,\big)\\

\end{array}

\right]=Z^{\,1/\rho-\beta(1-\varepsilon)-1}\,e^{Z}

\left[ \sum\limits_{m=0}^{M-1}A_m Z^{-m}+O\left(Z^{-M}\right) \right],

$$

где

$Z=(2-\varepsilon)\,\varepsilon^{\,\frac{\varepsilon}{2-\varepsilon}}z^{\,\frac{1}{2-\varepsilon}},$  $\varepsilon<2,$ а коэффициенты $A_m,$ $m=0, 1, 2, ...,$ зависят только от $\rho,$ $\varepsilon$ и $\beta.$



Учитывая, что $a$ "--- произвольное число из интервала $(0;1),$ из последней оценки и~соотношений \eqref{WrightRyadZ}, \eqref{Right_delta>1} получаем \eqref{T_muOtsenka}. \hfill \end{proof}



\smallskip



Покажем, что функция $y^{\mu-1} T_{\mu}\bigl(\frac{x^2}{4y^{\alpha}},\frac{1}{y^{\rho}}\bigr)$ является решением уравнения \eqref{GL3_1}. 

Из \eqref{T(z,zeta)}, \eqref{diff_x}  и \eqref{diff_Right} имеем

\begin{multline}

B_x T_{\mu}\Bigl(\frac{x^2}{4y^{\alpha}},\frac{1}{y^{\rho}}\Bigr)=

\\

=y^{-\alpha}\sum\limits_{n=1}^{\infty}\frac{1}{(n-1)! \Gamma(\beta+n)}

\Bigl(\frac{x^2}{4y^{\alpha}}\Bigr)^{\beta+n-1}

\phi\Bigl(-\rho,\mu-\alpha\beta-\alpha n; -\frac{1}{y^{\rho}}\Bigr)=

\\

=y^{-\alpha}\sum\limits_{k=0}^{\infty}\frac{1}{k! \Gamma(1+\beta+k)}

\Bigl(\frac{x^2}{4y^{\alpha}}\Bigr)^{\beta+k}

\phi\Bigl(-\rho,\mu-\alpha-\alpha\beta-\alpha k; -\frac{1}{y^{\rho}}\Bigr)=

\\ 

=y^{-\alpha} T_{\mu-\alpha}\Bigl(\frac{x^2}{\,4y^{\alpha}},\frac{1}{y^{\rho}}\Bigr),

\end{multline}   

\begin{equation}

\label{DT(z,zeta)}  

D_{0y}^{\nu} y^{\mu-1} T_{\mu}\Bigl(\frac{x^2}{4y^{\alpha}},\frac{1}{y^{\rho}}\Bigr)=

y^{\mu-\nu-1} T_{\mu-\nu}\Bigl(\frac{x^2}{4y^{\alpha}},\frac{1}{y^{\rho}}\Bigr).

\end{equation}



Из последних двух равенств следует

\begin{equation} \label{BDyT}  

\left(B_x-D_{0y}^{\alpha}\right)y^{\mu-1} T_{\mu}\Bigl(\frac{x^2}{4y^{\alpha}}, \frac{1}{y^{\rho}}\Bigr)=0.

\end{equation}   





Из равенства \eqref{DT(z,zeta)} и оценки \eqref{T_muOtsenka} также имеем

\begin{equation} \label{D(alpha-1)T}  

\lim\limits_{y \rightarrow 0} D_{0y}^{\alpha-1} y^{\mu-1} T_{\mu}\Bigl(\frac{x^2}{4y^{\alpha}}, \frac{1}{y^{\rho}}\Bigr)=0.

\end{equation}   



Таким образом, из соотношений  \eqref{T_muOtsenka}, \eqref{BDyT} и \eqref{D(alpha-1)T} следует, что нетривиальная функция 

$$

u(x,y)=y^{\mu-1} T_{\mu}\Bigl(\frac{x^2}{4y^{\alpha}},\frac{1}{y^{\rho}}\Bigr)

$$

является решением однородного уравнения \eqref{GL3_1}, удовлетворяет однородному условию \eqref{UslSMCoshi} $(\varphi(x)\equiv 0)$ и для любого положительного сколь угодно малого $\delta,$ 

некоторых положительных постоянных $k$ и $C,$ а также параметра 

$\rho,$  выбранного из условия 

$\rho=\frac{\alpha}{2}\frac{2+\delta(2-\alpha)}{\alpha+\delta(2-\alpha)},$

имеет место оценка

$$

\left| u(x,y) \right|\leqslant C \exp\left(k |x|^{\delta+\frac{2}{2-\alpha}} \right).

$$



Последняя оценка показывает, что увеличение показателя  $2/(2-\alpha)$ в~условии \eqref{AnalogTikhonovaInfty} ведет к~неединственности решения задачи  \eqref{GL3_1},  \eqref{UslSMCoshi} и~в~этом смысле условие  \eqref{AnalogTikhonovaInfty} является необходимым.





\smallskip









\Section[N]{Заключение}

В данной работе показывается, что показатель степени в~условии, обеспечивающем единственность решения задачи Коши для уравнения дробной диффузии с оператором Бесселя, является предельным, и его увеличение приводит к нарушению единственности. 



\smallskip





\AdditionalContent{Конкурирующие интересы}{Конкурирующих интересов не имею.} 



\AdditionalContent{Авторская ответственность}{Я несу полную ответственность за предоставление окончательной версии рукописи в печать. Окончательная версия рукописи мною одобрена.} 





\AdditionalContent{Финансирование}{Исследование выполнялось без финансирования.}



\smallskip





\begin{thebibliography}{99}



\RBibitem{Drob}

\by Нахушев~А.~М.

\book Дробное исчисление и его применение

\yr 2003

\publ Физматлит

\publaddr М.

\totalpages 271

%\sortdate 1/1/2003

\zmath{https://zbmath.org/?q=an:1066.26005}

%\misctransl

%\by Nakhushev~A.~M.

%\book Drobnoe ischislenie i ego primenenie \rm [Fractional calculus and its applications]

%\yr 2003

%\publ Fizmatlit

%\publaddr Moscow

%\lang In Russian

%\totalpages 271

%%\sortdate 1/1/2003



\RBibitem{Pskhu}

\by Псху~А.~В.

\book Уравнения в частных производных дробного порядка

\yr 2005

\publ Наука

\publaddr М.

\totalpages 199

%\sortdate 1/1/2005

\zmath{https://zbmath.org/?q=an:1193.35245}

%\misctransl

%\by Pskhu~A.~V.

%\book Uravneniia v chastnykh proizvodnykh drobnogo poriadka \rm [Partial differential equations of fractional order]

%\yr 2005

%\publ Nauka

%\publaddr Moscow

%\lang In Russian

%\totalpages 199

%%\sortdate 1/1/2005



\Bibitem{Alexiades}

\by Alexiades~V.

\paper Generalized axially symmetric heat potentials and singular parabolic initial boundary value problems

\jour Arch. Rational Mech. Anal.

\yr 1982

\vol 79

\issue 4

\pages 325--350

%\citnumscopus 3

%\sortdate 1/1/1982

\crossref{10.1007/BF00250797}

\mathscinet{http://www.ams.org/mathscinet-getitem?mr=656798}

\zmath{https://zbmath.org/?q=an:0511.35050}

\scopus{http://www.scopus.com/record/display.url?origin=inward&eid=2-s2.0-0020402046}



\Bibitem{Calton}

\by Calton~D.

\paper Cauchy's problem for a singular parabolic partial differential equation

\jour J.~Diff. Equations

\yr 1970

\vol 8

\issue 2

\pages 250--257

%\citnumscopus 8

%\sortdate 1/1/1970

\crossref{10.1016/0022-0396(70)90004-5}

\mathscinet{http://www.ams.org/mathscinet-getitem?mr=271547}

\scopus{http://www.scopus.com/record/display.url?origin=inward&eid=2-s2.0-49849107701}



\Bibitem{Kepinski_2}

\by K\c{e}pi\'nski~S.

\paper \"Uber die Differentialgleichung $\frac{\partial^2 z }{\partial x^2}+\frac{m+1}{x}\frac{\partial z }{\partial x} -\frac{n}{x}\frac{\partial z }{\partial t}=0$

\jour Math. Ann.

\yr 1905

\vol 61

\issue 3

\pages 397--405

\lang In German

%\sortdate 1/1/1905

\mathscinet{http://www.ams.org/mathscinet-getitem?mr=1511354}

\zmath{https://zbmath.org/?q=an:36.0416.01}



\RBibitem{Tersenov}

\by Терсенов~С.~А.

\book Параболические уравнения с~меняющимся направлением времени

\yr 1985

\publ Наука

\publaddr М.

\totalpages 105

%\sortdate 1/1/1985

%\misctransl

%\by Tersenov~S.~A.

%\book Parabolicheskie uravneniia s~meniaiushchimsia napravleniem vremeni \rm [Parabolic Equations with a Changing Time Direction]

%\yr 1985

%\publ Nauka

%\publaddr Moscow

%\totalpages 105

%%\sortdate 1/1/1985



\RBibitem{Matiychuk}

\by Матiйчук~M.~I.

\book Параболiчнi сингулярнi крайовi задачi

\yr 1999

\publ Iн-т математики НАН Укра\"\iни

\publaddr Ки\"\iв

\totalpages 176

%%\sortdate 1/1/1999

%\misctransl

%\by Matiichuk~M.~I.

%\book Parabolichni singuliarni kraiovi zadachi \rm [Parabolic Singular Boundary-Value Problems]

%\yr 1999

%\publ Institute of Mathematics of the Ukrainian National Academy of Sciences

%\publaddr Kiev

%\lang In Ukrainian

%\totalpages 176

%%\sortdate 1/1/1999



\RBibitem{Kipriyanov-Monographiya}

\by Киприянов~И.~А.

\book Сингулярные эллиптические краевые задачи

\yr 1997

\publ Наука

\publaddr М.

\totalpages 208

%\sortdate 1/1/1997

%\misctransl

%\by Kipriyanov~I.~A.

%\book Singuliarnye ellipticheskie kraevye zadachi \rm [Singular Elliptic Boundary-Value Problems]

%\yr 1997

%\publ Nauka

%\publaddr Moscow

%\lang In Russian

%\totalpages 208

%%\sortdate 1/1/1997



\RBibitem{Kipriyanov}

\by Киприянов~И.~А., Катрахов~В.~В., Ляпин~В.~М.

\paper О краевых задачах в областях общего вида для сингулярных параболических систем уравнений

\jour Докл. АН СССР

\yr 1976

\vol 230

\issue 6

\pages 1271--1274

%\sortdate 1/1/1976

\mathnet{http://mi.mathnet.ru/rus/dan40691}

\mathscinet{http://www.ams.org/mathscinet-getitem?mr=0425366}

\zmath{https://zbmath.org/?q=an:0354.35056}

%\misctransl

%\by Kipriyanov~I.~A., Katrakhov~V.~V., Lyapin~V.~M.

%\paper On boundary value problems in domains of general type for singular parabolic systems of equations

%\jour Sov. Math., Dokl.

%\yr 1976

%\vol 17

%\pages 1461--1464

%%\sortdate 1/1/1976

%\mathscinet{http://www.ams.org/mathscinet-getitem?mr=0425366}

%\zmath{https://zbmath.org/?q=an:0354.35056}



\RBibitem{Sitnik}

\by Ситник~С.~М.

\yr 2016

\publaddr Воронеж

\thesis Применение операторов преобразования Бушмана--Эрдейи и их обобщений в~теории дифференциальных уравнений с~особенностями в коэффициентах

\thesisinfo Дис.~\dots\ д-ра физ.-мат. наук: 01.01.02

\totalpages 307

%\sortdate 1/1/2016

%\misctransl

%\by Sitnik~S.~M.

%\yr 2016

%\publaddr Voronezh

%\thesis Application of the Bushman--Erd\'elyi transformation operators and their generalizations in the theory of differential equations with singularities in the coefficients

%\thesisinfo Dr. Phys. Math. Sci. Thesis

%\lang In Russian

%\totalpages 307

%%\sortdate 1/1/2016



\RBibitem{KatrakhovSitnik}

\by Катрахов~В.~В., Ситник~С.~М.

\paper Метод операторов преобразования и краевые задачи для сингулярных эллиптических уравнений

\inbook Сингулярные дифференциальные уравнения

\serial СМФН

\yr 2018

\vol 64

\issue 2

\pages 211--426

\publ Российский университет дружбы народов

\publaddr М.

%\sortdate 1/1/2018

\mathnet{http://mi.mathnet.ru/rus/cmfd355}

\crossref{10.22363/2413-3639-2018-64-2-211-426}

%\misctransl

%\by Katrakhov~V.~V., Sitnik~S.~M.

%\paper The transmutation method and boundary-value problems for singular elliptic equations

%\inbook Singular differential equations

%\serial CMFD

%\yr 2018

%\vol 64

%\issue 2

%\pages 211--426

%\publ Peoples' Friendship University of Russia

%\publaddr Moscow

%\lang In Russian

%%\sortdate 1/1/2018

%\crossref{https://doi.org/10.22363/2413-3639-2018-64-2-211-426}



\RBibitem{Kochubei_2}

\by Кочубей~А.~Н.

\paper Диффузия дробного порядка

\jour Дифференц. уравнения

\yr 1990

\vol 26

\issue 4

\pages 660--670

%\sortdate 1/1/1990

\mathnet{http://mi.mathnet.ru/rus/de7144}

\mathscinet{http://www.ams.org/mathscinet-getitem?mr=1061448}

\zmath{https://zbmath.org/?q=an:0712.35049}

%\transl

%\by Kochubei~A.~N.

%\paper Diffusion of fractional order

%\jour Differ. Equ.

%\yr 1990

%\vol 26

%\issue 4

%\pages 485--492

%%\sortdate 1/1/1990

%\mathscinet{http://www.ams.org/mathscinet-getitem?mr=1061448}

%\zmath{https://zbmath.org/?q=an:0729.35064}



\RBibitem{Kochubei_1}

\by Кочубей~А.~Н.

\paper Задача Коши для эволюционных уравнений дробного порядка

\jour Дифференц. уравнения

\yr 1989

\vol 25

\issue 8

\pages 1359--1368

%\sortdate 1/1/1989

\mathnet{http://mi.mathnet.ru/rus/de6938}

\mathscinet{http://www.ams.org/mathscinet-getitem?mr=1014153}

%\transl

%\by Kochubei~A.~N.

%\paper The Cauchy problem for evolution equations of fractional order

%\jour Differ. Equ.

%\yr 1989

%\vol 25

%\issue 8

%\pages 967--974

%%\sortdate 1/1/1989

%\mathscinet{http://www.ams.org/mathscinet-getitem?mr=1014153}



\RBibitem{Kochubei_3}

\by Кочубей~А.~Н., Эйдельман~С.~Д.

\paper Задача Коши для эволюционных уравнений дробного порядка

\jour Докл. РАН

\yr 2004

\vol 394

\issue 2

\pages 159--161

%\sortdate 1/1/2004

\mathnet{http://mi.mathnet.ru/rus/dan1624}

\mathscinet{http://www.ams.org/mathscinet-getitem?mr=2089744}

\zmath{https://zbmath.org/?q=an:1129.35427}

%\misctransl

%\by Kochubeii~A.~N., \'Eidel'man~S.~D.

%\paper The Cauchy problem for evolution equations of fractional order

%\jour Dokl. Akad. Nauk

%\yr 2004

%\vol 394

%\issue 2

%\pages 159--161

%\lang In Russian

%%\sortdate 1/1/2004



\Bibitem{Kochubei_4}

\by Kochubei~A.~N.

\paper Asymptotic Properties of Solutions of the Fractional Diffusion-Wave Equation

\jour Fract. Calc. Appl. Anal.

\yr 2014

\vol 17

\issue 3

\pages 881--896

%\citnumscopus 9

%\sortdate 1/1/2014

\crossref{10.2478/s13540-014-0203-3}

\mathscinet{http://www.ams.org/mathscinet-getitem?mr=3260311}

\zmath{https://zbmath.org/?q=an:1323.35197}

\elib{http://elibrary.ru/item.asp?id=24541812}

\scopus{http://www.scopus.com/record/display.url?origin=inward&eid=2-s2.0-84904312827}



\Bibitem{Kochubei_5}

\by Kochubei~A.~N.

\paper Cauchy problem for fractional diffusion-wave equations with variable coefficients

\jour Applicable Analysis

\yr 2014

\vol 93

\issue 10

\pages 2211--2242

%\citnumscopus 7

%\sortdate 1/1/2014

\crossref{10.1080/00036811.2013.875162}

\mathscinet{http://www.ams.org/mathscinet-getitem?mr=3240385}

\zmath{https://zbmath.org/?q=an:1297.35274}

\elib{http://elibrary.ru/item.asp?id=24642244}

\scopus{http://www.scopus.com/record/display.url?origin=inward&eid=2-s2.0-84905387477}



\RBibitem{Pskhu2}

\by Псху~А.~В.

\paper Фундаментальное решение диффузионно-волнового уравнения дробного порядка

\jour Изв. РАН. Сер. матем.

\yr 2009

\vol 73

\issue 2

\pages 141--182

%\sortdate 1/1/2009

\mathnet{http://mi.mathnet.ru/rus/im2429}

\crossref{10.4213/im2429}

\mathscinet{http://www.ams.org/mathscinet-getitem?mr=2532450}

\zmath{https://zbmath.org/?q=an:1172.26001}

\adsnasa{http://adsabs.harvard.edu/cgi-bin/bib_query?2009IzMat..73..351P}

\elib{http://elibrary.ru/item.asp?id=20425204}

%\transl

%\by Pskhu~A.~V.

%\paper The fundamental solution of a diffusion-wave equation of fractional order

%\jour Izv. Math.

%\yr 2009

%\vol 73

%\issue 2

%\pages 351--392

%%\sortdate 1/1/2009

%\crossref{https://doi.org/10.1070/IM2009v073n02ABEH002450}

%\mathscinet{http://www.ams.org/mathscinet-getitem?mr=2532450}

%\zmath{https://zbmath.org/?q=an:1172.26001}

%\isi{http://gateway.isiknowledge.com/gateway/Gateway.cgi?GWVersion=2&SrcApp=PARTNER_APP&SrcAuth=LinksAMR&DestLinkType=FullRecord&DestApp=ALL_WOS&KeyUT=000266177900006}

%\elib{http://elibrary.ru/item.asp?id=13597908}

%\scopus{http://www.scopus.com/record/display.url?origin=inward&eid=2-s2.0-65349154473}



\RBibitem{Pskhu-SibirJournal}

\by Псху~А.~В.

\paper Уравнение дробной диффузии с оператором дискретно распределенного дифференцирования

\jour Сиб. электрон. матем. изв.

\yr 2016

\vol 13

\pages 1078--1098

%\sortdate 1/1/2016

\mathnet{http://mi.mathnet.ru/rus/semr736}

\crossref{10.17377/semi.2016.13.086}

\mathscinet{http://www.ams.org/mathscinet-getitem?mr=3580054}

\zmath{https://zbmath.org/?q=an:1370.35268}

\elib{http://elibrary.ru/item.asp?id=28127183}

%\misctransl

%\by Pskhu~A.~V.

%\paper Fractional diffusion equation with discretely distributed differentiation operator

%\jour Sib. \`Elektron. Mat. Izv.

%\yr 2016

%\vol 13

%\pages 1078--1098

%\lang In Russian

%%\sortdate 1/1/2016

%\crossref{https://doi.org/10.17377/semi.2016.13.086}



\Bibitem{Pskhu3}

\by Pskhu~A.~V.

\paper Multi-time fractional diffusion equation

\jour Eur. Phys. J.~Spec. Top.

\yr 2013

\vol 222

\issue 8

\pages 1939--1950

%\citnumscopus 2

%\sortdate 1/1/2013

\crossref{10.1140/epjst/e2013-01975-y}

\scopus{http://www.scopus.com/record/display.url?origin=inward&eid=2-s2.0-84885034281}



\RBibitem{Pskhu4}

\by Псху~А.~В.

\paper Первая краевая задача для дробного диффузионно-волнового уравнения в~нецилиндрической области

\jour Изв. РАН. Сер. матем.

\yr 2017

\vol 81

\issue 6

\pages 158--179

%\sortdate 1/1/2017

\mathnet{http://mi.mathnet.ru/rus/im8520}

\crossref{10.4213/im8520}

\mathscinet{http://www.ams.org/mathscinet-getitem?mr=3733319}

\zmath{https://zbmath.org/?q=an:06844352}

\adsnasa{http://adsabs.harvard.edu/cgi-bin/bib_query?2017IzMat..81.1212P}

\elib{http://elibrary.ru/item.asp?id=30737847}

%\transl

%\by Pskhu~A.~V.

%\paper The first boundary-value problem for a fractional diffusion-wave equation in a non-cylindrical domain

%\jour Izv. Math.

%\yr 2017

%\vol 81

%\issue 6

%\pages 1212--1233

%%\sortdate 1/1/2017

%\crossref{https://doi.org/10.1070/IM8520}

%\mathscinet{http://www.ams.org/mathscinet-getitem?mr=3733319}

%\zmath{https://zbmath.org/?q=an:06844352}

%\isi{http://gateway.isiknowledge.com/gateway/Gateway.cgi?GWVersion=2&SrcApp=PARTNER_APP&SrcAuth=LinksAMR&DestLinkType=FullRecord&DestApp=ALL_WOS&KeyUT=000418891300007}

%\scopus{http://www.scopus.com/record/display.url?origin=inward&eid=2-s2.0-85040983678}



\RBibitem{VorKilbas}

\by Ворошилов~А.~А., Килбас~А.~А.

\paper Задача типа Коши для диффузионно-волнового уравнения с частной производной Римана--Лиувилля

\jour Докл. РАН

\yr 2006

\vol 406

\issue 1

\pages 12--16

%\sortdate 1/1/2006

\mathnet{http://mi.mathnet.ru/rus/dan1045}

\mathscinet{http://www.ams.org/mathscinet-getitem?mr=2256848}

\zmath{https://zbmath.org/?q=an:1155.35396}

\elib{http://elibrary.ru/item.asp?id=9176502}

%\misctransl

%\by Voroshilov~A.~A., Kilbas~A.~A.

%\paper A Cauchy-type problem for a diffusion-wave equation with a Riemann-Liouville partial derivative

%\jour Dokl. Akad. Nauk

%\yr 2006

%\vol 406

%\issue 1

%\pages 12--16

%\lang In Russian

%%\sortdate 1/1/2006



\RBibitem{VorKilbas2}

\by Ворошилов~А.~А., Килбас~А.~А.

\paper Задача Коши для диффузионно-волнового уравнения с~частной производной Капуто

\jour Дифференц. уравнения

\yr 2006

\vol 42

\issue 5

\pages 599--609

%\sortdate 1/1/2006

\mathnet{http://mi.mathnet.ru/rus/de11489}

\mathscinet{http://www.ams.org/mathscinet-getitem?mr=2292158}

\zmath{https://zbmath.org/?q=an:1123.35302}

\elib{http://elibrary.ru/item.asp?id=9245260}

%\transl

%\by Voroshilov~A.~A., Kilbas~A.~A.

%\paper The Cauchy problem for the diffusion-wave equation with the Caputo partial derivative

%\jour Differ. Equ.

%\yr 2006

%\vol 42

%\issue 5

%\pages 638--649

%%\citnumscopus 11

%%\sortdate 1/1/2006

%\crossref{https://doi.org/10.1134/S0012266106050041}

%\mathscinet{http://www.ams.org/mathscinet-getitem?mr=2292158}

%\zmath{https://zbmath.org/?q=an:1123.35302}

%\elib{http://elibrary.ru/item.asp?id=27786841}

%\scopus{http://www.scopus.com/record/display.url?origin=inward&eid=2-s2.0-33745304762}



\RBibitem{Gekk}

\by Геккиева~С.~Х.

\paper Задача Коши для обобщенного уравнения переноса с дробной по времени производной

\jour Доклады Адыгской (Черкесской) Международной академии наук

\yr 2000

\vol 5

\issue 1

\pages 16--19

%\sortdate 1/1/2000

%\misctransl

%\by Gekkieva~S.~Kh.

%\paper The Cauchy problem for the generalized transport equation with time-fractional derivative

%\jour Dokl. Adyg. (Cherkess.) Mezdunar. Akad. Nauk

%\yr 2000

%\vol 5

%\issue 1

%\pages 16--19

%\lang In Russian

%%\sortdate 1/1/2000



\RBibitem{Gekkieva1}

\by Геккиева~С.~Х.

\paper Краевая задача для обобщенного уравнения переноса с дробной производной в полубесконечной области

\jour Изв. КБНЦ РАН

\yr 2002

\issue 1(8)

\pages 6--8

%\sortdate 1/1/2002

%\misctransl

%\by Gekkieva~S.~Kh.

%\paper A boundary value problem for the generalized transfer equation with a fractional derivative in a semi-infinite domain

%\jour Izv. KBNTs RAN

%\yr 2002

%\issue 1(8)

%\pages 6--8

%\lang In Russian

%%\sortdate 1/1/2002



\Bibitem{Mainardi-1}

\by Mainardi~F.

\paper The fundamental solutions for the fractional diffusion-wave equation

\jour Appl. Math. Letters

\yr 1996

\vol 6

\issue 1

\pages 23--28

%\citnumscopus 343

%\sortdate 1/1/1996

\crossref{10.1016/0893-9659(96)00089-4}

\mathscinet{http://www.ams.org/mathscinet-getitem?mr=1419811}

\zmath{https://zbmath.org/?q=an:0879.35036}

\scopus{http://www.scopus.com/record/display.url?origin=inward&eid=2-s2.0-30244460855}



\Bibitem{Mainardi-3}

\by Mainardi~F.

\paper The time fractional diffusion-wave equation

\jour Radiophys. Quantum Electron.

\yr 1995

\vol 38

\issue 1-2

\pages 13--24

%\citnumscopus 39

%\sortdate 1/1/1995

\crossref{10.1007/BF01051854}

\mathscinet{http://www.ams.org/mathscinet-getitem?mr=1427165}

\scopus{http://www.scopus.com/record/display.url?origin=inward&eid=2-s2.0-0029427754}



\Bibitem{Mainardi-2}

\by Mainardi~F., Luchko~Yu., Pagnini~G.

\paper The fundamental solution of the space-time fractional diffusion equation

\jour Fract. Calc. Appl. Anal.

\yr 2001

\vol 4

\issue 2

\pages 153--192

\arxiv \href{https://arxiv.org/abs/cond-mat/0702419}{cond-mat/0702419} [cond-mat.stat-mech]

%\sortdate 1/1/2001

\mathscinet{http://www.ams.org/mathscinet-getitem?mr=1829592}

\zmath{https://zbmath.org/?q=an:1054.35156}



\Bibitem{Pagnini_1}

\by Pagnini~G.

\paper The M-Wright function as a generalization of the Gaussian density for fractional diffusion processes

\jour Fract. Calc. Appl. Anal.

\yr 2013

\vol 16

\issue 2

\pages 436--453

%\citnumscopus 22

%\sortdate 1/1/2013

\crossref{10.2478/s13540-013-0027-6}

\mathscinet{http://www.ams.org/mathscinet-getitem?mr=3033698}

\zmath{https://zbmath.org/?q=an:1312.33061}

\elib{http://elibrary.ru/item.asp?id=28682912}

\scopus{http://www.scopus.com/record/display.url?origin=inward&eid=2-s2.0-84879620364}



\Bibitem{Pagnini_Paradisi}

\by Pagnini~G., Paradisi~P.

\paper A stochastic solution with Gaussian stationary increments of the symmetric space-time fractional diffusion equation

\jour Fract. Calc. Appl. Anal.

\yr 2016

\vol 19

\issue 2

\pages 408--440

%\citnumscopus 9

%\sortdate 1/1/2016

\crossref{10.1515/fca-2016-0022}

\mathscinet{http://www.ams.org/mathscinet-getitem?mr=3513003}

\zmath{https://zbmath.org/?q=an:1341.60073}

\elib{http://elibrary.ru/item.asp?id=27104131}

\scopus{http://www.scopus.com/record/display.url?origin=inward&eid=2-s2.0-84969800118}



\Bibitem{Podlubny}

\by Podlubny~I.

\book Fractional differential equations. An introduction to fractional derivatives, fractional differential equations, to methods of their solution and some of their applications

\serial Mathematics in Science and Engineering

\yr 1999

\vol 198

\publ Academic Press

\publaddr San Diego, CA

\totalpages xxiv+340

%\sortdate 1/1/1999

\mathscinet{http://www.ams.org/mathscinet-getitem?mr=1658022}

\zmath{https://zbmath.org/?q=an:0924.34008}



\Bibitem{MathaiSaxenaHaubold}

\by Mathai~A.~M., Saxena~R.~K., Haubold~H.~J.

\book The $H$-function. Theory and Applications

\yr 2010

\publ Springer

\publaddr Dordrecht

\totalpages xiv+268~pp \nofrills

%\citnumscopus 411

%\sortdate 1/1/2010

\crossref{10.1007/978-1-4419-0916-9}

\mathscinet{http://www.ams.org/mathscinet-getitem?mr=2562766}

\zmath{https://zbmath.org/?q=an:1181.33001}

\scopus{http://www.scopus.com/record/display.url?origin=inward&eid=2-s2.0-84891408876}



\RBibitem{KhushtovaVestn2}

\by Хуштова~Ф.~Г.

\paper Задача Коши для уравнения параболического типа~с~оператором Бесселя и~частной производной~Римана--Лиувилля

\jour Вестн. Сам. гос. техн. ун-та. Сер. Физ.-мат. науки

\yr 2016

\vol 20

\issue 1

\pages 74--84

%\sortdate 1/1/2016

\mathnet{http://mi.mathnet.ru/rus/vsgtu1455}

\crossref{10.14498/vsgtu1455}

\zmath{https://zbmath.org/?q=an:06964474}

\elib{http://elibrary.ru/item.asp?id=26898117}

%\misctransl

%\by Khushtova~F.~G.

%\paper Cauchy problem for a parabolic equation with Bessel operator and Riemann--Liouville partial derivative

%\jour Vestn. Samar. Gos. Tekhn. Univ., Ser. Fiz.-Mat. Nauki \rm [J. Samara State Tech. Univ., Ser. Phys. Math. Sci.]

%\yr 2016

%\vol 20

%\issue 1

%\pages 74--84

%%\sortdate 1/1/2016



\Bibitem{GorenfloLuchkoMainardi}

\by Gorenflo~R., Luchko~Y., Mainardi~F.

\paper Analytical properties and applications of the Wright function

\jour Fract. Calc. Appl. Anal.

\yr 1999

\vol 2

\issue 4

\pages 383--414

\arxiv \href{https://arxiv.org/abs/math-ph/0701069}{math-ph/0701069}

%\sortdate 1/1/1999

\mathscinet{http://www.ams.org/mathscinet-getitem?mr=1752379}

\zmath{https://zbmath.org/?q=an:1027.33006}



\Bibitem{tikhonovUsl-1}

\by Tychonoff~A.

\paper Th\'eor\`emes d'unicit\'e pour l'\'equation de la chaleur

\jour Mat. Sb.

\yr 1935

\vol 42

\issue 2

\pages 199--216

\lang In French

%\sortdate 1/1/1935

\mathnet{http://mi.mathnet.ru/rus/sm6410}

\zmath{https://zbmath.org/?q=an:0012.35501|61.1203.05}



\RBibitem{tikhonovUsl-2}

\by Тихонов~А.~Н.

\paper Теоремы единственности для уравнения теплопроводности

\inbook Собрание научных трудов: в десяти томах

\yr 2009

\pages 371--382

\publ Наука

\publaddr М.

\volinfo Т.~II. Математика. Ч.~2. Вычислительная математика 1956--1979. Математическая физика 1933--1948. Ред.-сост. Т.~А.~Сушкевич, А.~В.~Гулин

%\sortdate 1/1/2009

\zmath{https://zbmath.org/?q=an:1200.01055}

%\misctransl

%\by Tikhonov~A.~N.

%\paper Uniqueness theorems for the heat equation

%\inbook Collection of scientific works. In ten volumes

%\yr 2009

%\pages 371--382

%\publ Nauka

%\publaddr Moscow

%\lang In Russian and French

%\volinfo Vol. 2: Mathematics. Part 2: Numerical mathematics 1956–1979. Mathematical physics 1933–1948. Edited by T.~A.~Sushkevich and A.~V.~Gulin

%%\sortdate 1/1/2009

%\zmath{https://zbmath.org/?q=an:1200.01055}



\Bibitem{Wright}

\by Wright~E.~M.

\paper The asymptotic expansion of the generalized hypergeometric function

\jour Proc. Lond. Math. Soc., II. Ser.

\yr 1940

\vol 46

\issue 1

\pages 389--408

%\citnumscopus 136

%\sortdate 1/1/1940

\crossref{10.1112/plms/s2-46.1.389}

\mathscinet{http://www.ams.org/mathscinet-getitem?mr=3876}

\zmath{https://zbmath.org/?q=an:0025.40402|66.1243.01}

\scopus{http://www.scopus.com/record/display.url?origin=inward&eid=2-s2.0-84963078619}



\end{thebibliography}



	





\newpage



\makeentitle



\vspace{-6mm}



%\AdditionalContent{Competing interests}{Specify what conflicts of interest are available.} 

\AdditionalContent{Competing interests}{I have no competing interests.} 

%\AdditionalContent{Competing interests}{We have ... (or) no competing interests.} 

%\AdditionalContent{Authors' contributions and responsibilities}{What is the author's responsibility for each author?}

\AdditionalContent{Author's Responsibilities}{I take full responsibility for submitting the final ma\-nuscript in print. I approved the final version of the manuscript.} 

%\AdditionalContent{Authors' contributions and responsibilities}{Each author has participated in the article concept development and in the manuscript writing. The authors are absolutely responsible for submitting the final manuscript in print. Each author has approved the final version of manuscript.} 



%\AdditionalContent{Funding}{Indicate the source(s) of your funding.}



\AdditionalContent{Funding}{This research received no specific grant from any funding agency in the public,

commercial, or not-for-profit sectors.}



\newpage



\begin{thebibliography}{10}



\Bibitem{Drob:2}

\lang In Russian

\by Nakhushev~A.~M.

\book Drobnoe ischislenie i ego primenenie \rm [Fractional calculus and its applications]

\publaddr Moscow

\publ Fizmatlit 

\totalpages 271

\yr 2003





\Bibitem{Pskhu:w}

\lang In Russian

\by Pskhu~A.~V.

\book Uravneniia v chastnykh proizvodnykh drobnogo poriadka \rm [Partial differential equations of fractional order]

\publaddr Moscow

\publ Nauka 

\totalpages 199 

\yr 2005







\Bibitem{Alexiades:w}

\by Alexiades~V.

\paper Generalized axially symmetric heat potentials and singular parabolic initial boundary value problems

\jour  Arch. Rational Mech. Anal. 

\yr 1982

\vol 79

\issue 4

\pages 325--350

\crossref{10.1007/BF00250797}







\Bibitem{Calton:w}

\by Calton~D.

\paper Cauchy's problem for a singular parabolic partial differential equation

\jour J.~Diff. Equations

\yr 1970

\vol 8

\issue 2

\pages 250--257

\crossref{10.1016/0022-0396(70)90004-5}







\Bibitem{Kepinski_2:w}

\by K\c{e}pi\'nski~S.

\paper \"Uber die Differentialgleichung $\frac{\partial^2 z }{\partial x^2}+\frac{m+1}{x}\frac{\partial z }{\partial x} -\frac{n}{x}\frac{\partial z }{\partial t}=0$

\jour   Math. Ann.

\yr 1905

\vol 61

\issue 3

\pages 397--405

\zmath{36.0416.01}

\lang In German



\Bibitem{Tersenov:en}

\lang In Russian

\by Tersenov~S.~A.

\book Parabolicheskie uravneniia s~meniaiushchimsia napravleniem vremeni \rm [Parabolic Equations with a Changing Time Direction]

\yr 1985

\publ Nauka

\publaddr Moscow

\totalpages 105



\Bibitem{Matiychuk:en}

\by Matiichuk~M.~I.

\book  Parabolichni singuliarni kraiovi zadachi \rm [Parabolic Singular Boundary-Value Problems]

\yr 1999

\publ Institute of Mathematics of the Ukrainian National  Academy of Sciences

\publaddr Kiev

\totalpages 176

\lang In Ukrainian





\Bibitem{Kipriyanov-Monographiya:en}

\lang In Russian

\by Kipriyanov~I.~A.

\book  Singuliarnye ellipticheskie kraevye zadachi \rm [Singular Elliptic Boundary-Value Problems]

\yr 1997

\publ Nauka

\publaddr Moscow

\totalpages 208



\Bibitem{Kipriyanov:en}

\by Kipriyanov~I.~A., Katrakhov~V.~V., Lyapin~V.~M.

\paper On boundary value problems in domains of general type for singular parabolic systems of equations

\jour Sov. Math., Dokl. 

\vol 17

\yr 1976

\pages 1461--1464 

\mathscinet{http://www.ams.org/mathscinet-getitem?mr=0425366}

\zmath{https://zbmath.org/?q=an:0354.35056}





\Bibitem{Sitnik:en}

\lang In Russian

\by Sitnik~S.~M. 

\thesis Application of the Bushman--Erd\'elyi transformation operators and their generalizations in the theory of differential equations with singularities in the coefficients

\thesisinfo Dr. Phys. Math. Sci. Thesis

\yr 2016

\publaddr Voronezh

\totalpages 307







\Bibitem{KatrakhovSitnik:en}

\lang In Russian

\by Katrakhov~V.~V., Sitnik~S.~M.

\paper The transmutation method and boundary-value problems for singular elliptic equations

\inbook Singular differential equations

\serial CMFD

\yr 2018

\vol 64

\issue 2

\pages 211--426

\publ Peoples' Friendship University of Russia

\publaddr Moscow

\crossref{10.22363/2413-3639-2018-64-2-211-426}







\Bibitem{Kochubei_2:en}

\by Kochubei~A.~N.

\paper Diffusion of fractional order

\jour Differ. Equ.

\yr 1990

\vol 26

\issue 4

\pages 485--492







\Bibitem{Kochubei_1:w}

\by Kochubei~A.~N.

\paper The Cauchy problem for evolution equations of fractional order

\jour Differ. Equ.

\yr 1989

\vol 25

\issue 8

\pages 967--974



\Bibitem{Kochubei_3:e}

\lang In Russian

\by Kochubeii~A.~N., \'Eidel'man~S.~D.

\paper The Cauchy problem for evolution equations of fractional order

\jour Dokl. Akad. Nauk

\yr 2004

\vol 394

\issue 2

\pages 159--161







\Bibitem{Kochubei_4:e}

\by Kochubei~A.~N.

\paper Asymptotic Properties of Solutions of the Fractional Diffusion-Wave Equation

\jour Fract. Calc. Appl. Anal.

\yr 2014

\vol 17

\issue 3

\pages 881--896

\crossref{10.2478/s13540-014-0203-3}





\Bibitem{Kochubei_5:e}

\by Kochubei~A.~N.

\paper Cauchy problem for fractional diffusion-wave equations with variable coefficients

\jour  Applicable Analysis

\yr 2014

\vol 93

\issue 10

\pages 2211--2242

\crossref{10.1080/00036811.2013.875162}







\Bibitem{Pskhu2:en}

\paper The fundamental solution of a diffusion-wave equation of fractional order

\by Pskhu~A.~V.

\jour Izv. Math.

\yr 2009

\vol 73

\issue 2

\pages 351--392

\crossref{10.1070/IM2009v073n02ABEH002450}

\isi{http://gateway.isiknowledge.com/gateway/Gateway.cgi?GWVersion=2&SrcApp=PARTNER_APP&SrcAuth=LinksAMR&DestLinkType=FullRecord&DestApp=ALL_WOS&KeyUT=000266177900006}

\elib{http://elibrary.ru/item.asp?id=13597908}

\scopus{http://www.scopus.com/record/display.url?origin=inward&eid=2-s2.0-65349154473}





\Bibitem{Pskhu-SibirJournal:en}

\lang In Russian

\by Pskhu~A.~V.

\paper Fractional diffusion equation with discretely distributed differentiation operator

\jour Sib. \`Elektron. Mat. Izv.

\yr 2016

\vol 13

\pages 1078--1098

\crossref{10.17377/semi.2016.13.086}







\Bibitem{Pskhu3:en}

\by Pskhu~A.~V.

\paper Multi-time fractional diffusion equation

\jour Eur. Phys. J.~Spec. Top.

\yr 2013

\vol 222

\issue 8

\pages 1939--1950

\crossref{10.1140/epjst/e2013-01975-y}





\Bibitem{Pskhu4:en}

\paper The first boundary-value problem for a fractional diffusion-wave equation in a non-cylindrical domain

\by Pskhu~A.~V.

\jour Izv. Math.

\yr 2017

\vol 81

\issue 6

\pages 1212--1233

\crossref{10.1070/IM8520}

\isi{http://gateway.isiknowledge.com/gateway/Gateway.cgi?GWVersion=2&SrcApp=PARTNER_APP&SrcAuth=LinksAMR&DestLinkType=FullRecord&DestApp=ALL_WOS&KeyUT=000418891300007}

\scopus{http://www.scopus.com/record/display.url?origin=inward&eid=2-s2.0-85040983678}







\Bibitem{VorKilbas:a} 

\lang In Russian

\by Voroshilov~A.~A., Kilbas~A.~A.

\paper A Cauchy-type problem for a diffusion-wave equation with a

Riemann-Liouville partial derivative

\jour Dokl. Akad. Nauk

\yr 2006

\vol 406

\issue 1

\pages 12--16



\Bibitem{VorKilbas2:en} 

\by Voroshilov~A.~A., Kilbas~A.~A.

\paper The Cauchy problem for the diffusion-wave equation with the Caputo partial derivative

\mathscinet{http://www.ams.org/mathscinet-getitem?mr=2292158}

\jour Differ. Equ.

\yr 2006

\vol 42

\issue 5

\pages 638--649

\crossref{10.1134/S0012266106050041}



\Bibitem{Gekk:en} 

\lang In Russian

\by Gekkieva~S.~Kh.

\paper The Cauchy problem for the generalized transport equation with time-fractional derivative

\jour Dokl. Adyg. (Cherkess.) Mezdunar. Akad. Nauk

\yr 2000

\vol 5

\issue 1

\pages 16--19







\Bibitem{Gekkieva1:en}

\lang In Russian

\by Gekkieva~S.~Kh.

\paper A boundary value problem for the generalized transfer equation with a fractional derivative in a semi-infinite domain

\jour Izv. KBNTs RAN

\yr 2002

\issue 1(8)

\pages 6--8











  





\Bibitem{Mainardi-1:en}

\by Mainardi~F.

\paper The fundamental solutions for the fractional diffusion-wave equation

\jour Appl.  Math. Letters

\yr 1996

\vol 6

\issue 1

\pages 23--28

\crossref{10.1016/0893-9659(96)00089-4}



\Bibitem{Mainardi-3}

\by Mainardi~F.

\paper  The time fractional diffusion-wave equation

\jour Radiophys. Quantum Electron.

\yr 1995

\vol 38

\issue 1-2

\pages 13--24

\crossref{10.1007/BF01051854}



\Bibitem{Mainardi-2:en}

\by Mainardi~F., Luchko~Yu., Pagnini~G.

\paper The fundamental solution of the space-time fractional diffusion equation

\jour Fract. Calc. Appl. Anal.

\yr 2001

\vol 4

\issue 2

\pages 153--192

\arxiv  	\href{https://arxiv.org/abs/cond-mat/0702419}{cond-mat/0702419} [cond-mat.stat-mech]















\Bibitem{Pagnini_1:en}

\by Pagnini~G.

\paper The M-Wright function as a generalization of the Gaussian density for fractional diffusion processes

\jour Fract. Calc. Appl. Anal.

\yr 2013

\vol 16

\issue 2

\pages 436--453

\crossref{10.2478/s13540-013-0027-6}





\Bibitem{Pagnini_Paradisi:en}

\by Pagnini~G., Paradisi~P.

\paper A stochastic solution with Gaussian stationary increments of the symmetric space-time fractional diffusion equation

\jour Fract. Calc. Appl. Anal.

\yr 2016

\vol 19

\issue 2

\pages 408--440

\crossref{10.1515/fca-2016-0022}





\Bibitem{Podlubny:en}

\by Podlubny~I.

\book Fractional differential equations. An introduction to fractional derivatives, fractional differential equations, to methods of their solution and some of their applications

\yr 1999

\publ Academic Press

\publaddr  San Diego, CA

\zmath{0924.34008}

\serial Mathematics in Science and Engineering

\vol 198

\totalpages xxiv+340





\Bibitem{MathaiSaxenaHaubold:en}

\by Mathai~A.~M., Saxena~R.~K., Haubold~H.~J.

\book The $H$-function. Theory and Applications

\yr 2010

\publ Springer

\publaddr  Dordrecht

\totalpages xiv+268~pp \nofrills

\crossref{10.1007/978-1-4419-0916-9}

\zmath{1181.33001}









\Bibitem{KhushtovaVestn2}

\by Khushtova~F.~G.

\paper Cauchy problem for a parabolic equation with Bessel operator and Riemann--Liouville partial derivative

\jour Vestn. Samar. Gos. Tekhn. Univ., Ser. Fiz.-Mat. Nauki \rm [J. Samara State Tech. Univ., Ser. Phys. Math. Sci.]

\yr 2016

\vol 20

\issue 1

\pages 74--84

\crossref{10.14498/vsgtu1455}





\Bibitem{GorenfloLuchkoMainardi:en}

\by Gorenflo~R., Luchko~Y., Mainardi~F.

\paper Analytical properties and applications of the Wright  function

\jour Fract. Calc. Appl. Anal.

\yr 1999

\vol 2

\issue 4

\pages 383--414

\arxiv \href{https://arxiv.org/abs/math-ph/0701069}{math-ph/0701069}







\Bibitem{tikhonovUsl-1:en}

\by Tychonoff~A.

\paper Th\'eor\`emes d'unicit\'e pour l'\'equation de la chaleur

\jour Mat. Sb.

\yr 1935

\vol 42

\issue 2

\pages 199--216

\lang In French

\mathnet{http://mi.mathnet.ru/msb6410}

\zmath{https://zbmath.org/?q=an:0012.35501|61.1203.05}









\Bibitem{tikhonovUsl-2:en}

\by Tikhonov~A.~N.

\paper Uniqueness theorems for the heat equation

\inbook Collection of scientific works. In ten volumes

\volinfo Vol. 2: Mathematics. Part 2: Numerical mathematics 1956–1979. Mathematical physics 1933–1948. Edited by T.~A.~Sushkevich and  A.~V.~Gulin

\yr 2009

\publ Nauka

\publaddr Moscow

\pages 371--382

\zmath{1200.01055}

\lang In Russian and French





\Bibitem{Wright:en}

\by Wright~E.~M.

\paper The asymptotic expansion of the generalized hypergeometric function

\jour Proc. Lond. Math. Soc., II. Ser. 

\yr 1940

\vol 46

\issue 1

\pages 389--408

\crossref{10.1112/plms/s2-46.1.389}

\zmath{0025.40402|66.1243.01}











\end{thebibliography}



%\end{document}

 \newpage

































%%%%

\input{preamble}





\usepackage{rotating_pr}

\usepackage{desclist}

\usepackage{cmap}

\usepackage{textcomp}

\usepackage{nccboxes}

\usepackage{euscript}

\usepackage{mathrsfs} 

\usepackage{multicol}

\usepackage{nccfoots}

\usepackage{mathrsfs}

\usepackage{calrsfs}

\allowdisplaybreaks[4]

\usepackage{qrcode}

\usepackage{afterpage}

\usepackage{wasysym}

\usepackage{wrapfig}

\usepackage{slashbox}



\usepackage{multirow}

\newcommand{\specialcell}[2][c]{%

  \begin{tabular}[#1]{@{}c@{}}#2\end{tabular}}





\usepackage{pgf,pgfplots,tikz}

\usetikzlibrary{arrows,snakes,backgrounds,patterns,shapes.geometric,calc}



\nofiles



\oek % для генерации pdf, отдаваемого в oek.rsl.ru

%\samgtupubl % для генерации pdf, отдаваемого в типографию СамГТУ

%\pdfcrop % обрезка pdf для размещения на math-net.ru

%\draft % На корректуре





%\DeclareMathOperator{\tr}{tr}

%\DeclareMathOperator{\ink}{Ink}

%\DeclareMathOperator{\erf}{erf}







%\makeatletter

%\newcommand{\PersonRadRef}{\@langrustrue\def\refname{{\bf\small Избранные  труды \rupid{19071}{О.~А.~Репина} за последние 5 лет}}}

%\makeatother



\DeclareMathOperator{\tr}{tr}

\DeclareMathOperator{\arcsh}{arcsh}









\usepackage{listings}

\usepackage{color}

\definecolor{dkgreen}{rgb}{0,0.6,0}

\definecolor{gray}{rgb}{0.5,0.5,0.5}

\definecolor{mauve}{rgb}{0.58,0,0.82}



\renewcommand*{\lstlistingname}{Листинг}

\usepackage{caption}

%\DeclareCaptionFont{white}{\color{white}} 

\DeclareCaptionFormat{listing}{\parbox{\textwidth}{\centering \small  #1.~#3}}

\captionsetup[lstlisting]{format=listing}



\lstset{ %

  inputencoding=utf8,

  language=R,                     % the language of the code

  basicstyle=\footnotesize,       % the size of the fonts that are used for the code

  numbers=left,                   % where to put the line-numbers

  numberstyle=\tiny\color{gray},  % the style that is used for the line-numbers

  stepnumber=1,                   % the step between two line-numbers. If it's 1, each line

                                  % will be numbered

  numbersep=5pt,                  % how far the line-numbers are from the code

  backgroundcolor=\color{white},  % choose the background color. You must add \usepackage{color}

  showspaces=false,               % show spaces adding particular underscores

  showstringspaces=false,         % underline spaces within strings

  showtabs=false,                 % show tabs within strings adding particular underscores

  frame=single,                   % adds a frame around the code

  rulecolor=\color{black},        % if not set, the frame-color may be changed on line-breaks within not-black text (e.g. commens (green here))

  tabsize=2,                      % sets default tabsize to 2 spaces

  captionpos=t,                   % sets the caption-position to bottom

  breaklines=true,                % sets automatic line breaking

  breakatwhitespace=false,        % sets if automatic breaks should only happen at whitespace

  %title=\lstname,                 % show the filename of files included with \lstinputlisting;

                                  % also try caption instead of title

  keywordstyle=\color{blue},      % keyword style

  commentstyle=\color{dkgreen},   % comment style

  stringstyle=\color{mauve},      % string literal style

  %escapeinside={\%*}{*)},         % if you want to add a comment within your code

  morekeywords={*,in,+,-,=}            % if you want to add more keywords to the set

} 







\begin{document}







\setcounter{page}{601}

\input{specialpages.tex} 







\input{./01du/partition.tex}

\input{./02mdtt/partition.tex}

\input{./03matmod/partition.tex}

\input{./04short/partition.tex}



%\input{./polieces.tex}



\end{document}



\newpage



$\;$





\chead[\fancyplain{}{}]{\fancyplain{}{}}

  \rhead[]{\fancyplain{}{}}

  \lhead[\fancyplain{}{}]{}

  \cfoot[]{}

  \lfoot[]{}

  \rfoot[]{}

\renewcommand{\plainheadrulewidth}{0pt}

\renewcommand{\headrulewidth}{0pt}





  \pagestyle{fancyplain}



  \thispagestyle{plain}





\newpage



$\;$





\chead[\fancyplain{}{}]{\fancyplain{}{}}

  \rhead[]{\fancyplain{}{}}

  \lhead[\fancyplain{}{}]{}

  \cfoot[]{}

  \lfoot[]{}

  \rfoot[]{}

\renewcommand{\plainheadrulewidth}{0pt}

\renewcommand{\headrulewidth}{0pt}





  \pagestyle{fancyplain}



  \thispagestyle{plain}





%\newpage

%

%$\;$

%

%

%\chead[\fancyplain{}{}]{\fancyplain{}{}}

%  \rhead[]{\fancyplain{}{}}

%  \lhead[\fancyplain{}{}]{}

%  \cfoot[]{}

%  \lfoot[]{}

%  \rfoot[]{}

%\renewcommand{\plainheadrulewidth}{0pt}

%\renewcommand{\headrulewidth}{0pt}

%

%

%  \pagestyle{fancyplain}

%

%  \thispagestyle{plain}







\end{document}





\Russian



\addtocontents{toc}{\bigskip\par\hspace{-7mm}}

\addtocontents{etoc}{\bigskip\par\hspace{-7mm}}



$\;$



\hfill \qrcode[hyperlink,height=.8in]{http://doi.org/10.14498/vsgtu1651}

\vspace{-1.5in}





%$$\quad $$ \vskip.6in \par \hskip-7mm  {\huge  Краткие сообщения}

$$\quad $$ \vskip.6in \par \hskip-7mm  {\huge  Short Communications}



%$$\quad $$ \vskip.6in \par \hskip-7mm  {\huge  Mathematical Modeling,\\ Numerical Methods \\[1mm] and Software Complexes}



\addtocontents{toc}{{\large\bf Краткие сообщения}\bigskip}

\addtocontents{etoc}{{\large\bf Short Communications}\bigskip}



\vspace{-5mm}



%\Partition{Математическое моделирование}{Mathematical Modeling}



%\input{preamble}

%\usepackage{samgtu-name}

%

%

%

\renewcommand{\RuCitation}{\noindent\textbf{\CYRO\cyrb\cyrr\cyra\cyrz\cyre\cyrc\ \cyrd\cyrl\cyrya\ \cyrc\cyri\cyrt\cyri\cyrr\cyro\cyrv\cyra\cyrn\cyri\cyrya\\}\selectlanguage{english}{\EnAuthorInContents} {\EnTitleInContents}, \textit{\EnJournalName}\/ [\ParallelJournalName],  \JournalYear, vol.~\JournalVolume, no.~\JournalNumber,  pp.~\thefirstpage--\thelastpage. \doiurl. \selectlanguage{russian}\smallskip}



\renewcommand{\EnCitation}{\noindent\textbf{Please cite this article in press as:\\} {\EnAuthorInContents} {\EnTitleInContents}, \textit{\EnJournalName}\/ [\ParallelJournalName],  \JournalYear, vol.~\JournalVolume, no.~\JournalNumber,  pp.~\thefirstpage--\thelastpage.  \midoiurl. }

%

\vsgtu{1651} %registration number

\FirstPage{735}

\LastPage{749}

%

%

%

%\pdfcrop

%\draft

%

\ShortCommunication

%

%\begin{document}

%

%

%

%$\;$

%\vspace{1mm}

%

%\hfill \qrcode[hyperlink,height=.8in]{http://doi.org/10.14498/vsgtu1651}

%\vspace{-.8in}





\renewcommand{\baselinestretch}{0.98}



\udc{532.51, 517.958:531.3-324}

\msc{76F02, 76F45, 76M45, 76R05, 76U05}



\rutitle{Динамические равновесия неизотермической жидкости}

\entitle{Dynamic equilibria of a nonisothermal fluid}



\ruauthor{Е.~Ю.~Просвиряков\affil{1,2}}{П\,р\,о\,с\,в\,и\,р\,я\,к\,о\,в~Е.~Ю.}

\enauthor{E.~Yu.~Prosviryakov\affil{1,2}}{P\,r\,o\,s\,v\,i\,r\,y\,a\,k\,o\,v~E.~Yu.}





\ruaffil{\ruaffid{2665}{Институт машиноведения УрО РАН,\\

Россия, 620049, Екатеринбург, ул. Комсомольская, 34}.

\and

\ruaffid{7370}{Уральский федеральный университет \\ им.~первого Президента России Б.~Н.~Ельцина, \\ 

Россия, 620002, Екатеринбург, ул. Мира, 19}.}



\enaffil{\enaffid{2665}{Institute of Engineering Science, Urals Branch, Russian Academy of Sciences,\\

34, Komsomolskaya st., Ekaterinburg, 620049, Russian Federation}.

\and

\enaffid{7370}{Ural Federal University named after the First President of Russia B.~N.~Yeltsin, \\ 19, Mira st., Ekaterinburg, 620002, Russian Federation}.

}



\enauthorsdetails{1}{

\fullauthor{\enpid{41426}{Evgeny Yu. Prosviryakov}}

\CAuthor 

\orcid{0000-0002-2349-7801} 

\regalia{Dr. Phys. \& Math. Sci., Professor} 

\position{Head of Sector} 

\dept{Sector of Nonlinear Vortex Hydrodynamics}

\email[br]{evgen_pros@mail.ru}

}



\ruauthorsdetails{1}{

\fullauthor{\rupid{41426}{Евгений Юрьевич Просвиряков}} 

\CAuthor 

\orcid{0000-0002-2349-7801}

\regalia{доктор физико-математических наук} 

\position{заведующий сектором} 

\dept{сектор нелинейной вихревой гидродинамики}

\email{evgen_pros@mail.ru}

}





\enkeywords{dynamic equilibrium,  rotating fluid flow, exact solution, Boussinesq approximation, counterflows, increased velocities}



\rukeywords{динамическое  равновесие, вращающийся поток жидкости, точное решение, аппроксимация Буссинеска, противотечение, увеличенные скорости}



\ruabstract{В рамках точности приближения Буссинеска рассмотрены стационарные динамические равновесия вращающейся массы неизотермической жидкости. Показано, что в этом случае в жидкости наблюдается конечное число противотечений и усиление скоростей по сравнению с заданными на границе значениями, а также формирование зон положительного и отрицательного давления и температуры. }



\enabstract{In this paper, stationary dynamic equilibria of the rotating mass of a nonisothermal fluid are discussed within the accuracy limits of the Boussinesq approximation. It is demonstrated that, in this case, a fluid exhibits a finite number of counterflows, higher values of velocities than those specified on the boundary and the formation of zones of positive and negative pressures and temperatures.}



\newdate{privalovaa}{07}{10}{2018}

\date{\displaydate{privalovaa}}

\newdate{privalovab}{10}{11}{2018}

\revisiondate{\displaydate{privalovab}}

\newdate{privalovac}{12}{11}{2018}

\accepteddate{\displaydate{privalovac}}



\newdate{privalovae}{13}{12}{2018}

\renewcommand{\JournalDataOnline}{\displaydate{privalovae}}





%\ForPeerReview

\makeentitle





\Section[n]{Introduction}

The study of the motion of a rotating self-gravitating fluid was started as long ago as by Isaak Newton~[\citen{ulg:bib:Newtoni},~\citen{ulg:bib:Turnbull}], and it is currently far from being completed. Newton proved to be the first to endeavour to obtain, from the motion equations, a solution simulating the shape of the Earth. Newton’s pioneering study gave rise to numerous investigations in this field of mechanics and physics and a great impetus to the development of the mathematics used practically in all the sections of mathematics. The development of the potential theory, the creation of the theory of special functions and the formation of the rapidly developing theory of integrable systems should be mentioned first.



Practically all the most prominent scientists, such as Appel, Basset, Betti, Gagen, Helmholtz, Dedekind, Dirichlet, Kirchhoff, Lipschitz, Lame, McLoren, Poincare, Riemann, Chandrasekhar, Jacoby, took part in the construction of the theory~[\citen{ulg:bib:Poincare}--\citen{ulg:bib:Dirichlet}]. This list of discoverers of exact solutions can be extended~[\citen{ulg:bib:Poincare}--\citen{ulg:bib:Dirichlet}]. The history of discovering the solutions can be found in the monographs and surveys~[\citen{ulg:bib:Poincare}--\citen{ulg:bib:Dirichlet}]. It should be noted that the scientific breakthrough in the simulation of rotating equilibrium figures was made by the Russian scientists Steklov, Lyapunov, Ovsyannikov and Zel'dovich.



Note that the development of the theory of motion stability described by systems of ordinary differential equations was initiated by studying the stability of equilibrium figures and the motion of an elastic body partially filled with a fluid~[\citen{ulg:bib:Lyapunov1},~\citen{ulg:bib:Lyapunov2}]. It is the work in these fields that offered a formulation to the well-known Lyapunov theorems of stability and instability, among other things, from the first approximation~[\citen{ulg:bib:Lyapunov1},~\citen{ulg:bib:Lyapunov2}]. Original research has currently been done on the stability of Dirichlet and Jacobi ellipsoids~\cite{ulg:bib:Borisov}.



Ovsyannikov’s exact gas-hydrodynamic solution became the first solution    si\-mulating gas cloud motion in the class of linear velocities~\cite{ulg:bib:Borisov}. Afterwards models of Dyson and Fujimoto were proposed~\cite{ulg:bib:Borisov}.



There is another substantial model, which influenced the development of the theory of rotating fluids. It is the generalized solid proposed for research by Arnold~\cite{ulg:bib:Arnold}. In the survey made by Dolzhansky~\cite{ulg:bib:Dolzhansky} there is a detailed analysis of this approach with some generalizations for different classes of fluids.



As a rule, dynamic fluid equilibria are simulated with the application of two mathematical models. One model is Lagrangian fluid simulation, i.e. the use of the ordinary differential equations of Lagrange, Hamilton and Jacoby to describe the relative equilibrium (solid-state) of the continuum. The other model consists in using the Euler equations with potential forces to analyze the structures of a relative fluid equilibrium. Thus, in the simulation of dynamic fluid equilibria no account is taken of dissipation and temperature effect. In other words, an isothermal perfect incompressible fluid is dealt with. 

	

It is apparent that this deficiency in studying fluid flows can be compensated if this problem is analyzed with the application of the Oberbeck--Boussinesq system of equations. Note that it is not only the Boussinesq approximation that has not been used before to describe dynamic fluid equilibria, but also the characteristic scale of solutions. This statement needs to be clarified. Despite Chandrasekhar~\cite{ulg:bib:Chandrasekhar} used the Boussinesq approximation to simulate astrophysical objects, he used the virial method of physical field expansion for solving his problems. The application of this body of mathematics enabled one to reduce the Oberbeck--Boussinesq system to a system of ordinary differential equations and to use the classical problem statements when interpreting results. Similar reasoning holds for~[\citen{ulg:bib:Arnold},~\citen{ulg:bib:Dolzhansky}].

	

There are studies dealing with large-scale fluid or gas motions~\cite{ulg:bib:Chandrasekhar}, when it is necessary to take into account the inhomogeneity of the gravitational field, and with scales for which surface tension forces are principally important, e. g. the Hadamard-Rybczynski problem~[\citen{ulg:bib:Hadamard},~\citen{ulg:bib:Rybczynski}]. In both extreme cases, the principal equilibrium shape is spherical. The consideration of rotation was always viewed as some disturbance of the principal flow, however it was not viewed as a forming factor for equilibrium structures.



Note that the very possibility of the existence of equilibrium figures in the region of intermediate scales, when neither the inhomogeneity of the gravitational field nor the surface tension forces are essential, does not seem to have been studied. At any rate, the authors are unfamiliar with works viewing the problem this way.



Thus, we intend to ascertain the possible existence of equilibrium structures in the region of the so-called intermediate scales, when the Boussinesq approximation is obviously applicable to the simulation of a heat-conducting viscous incompres\-sible fluid.





\smallskip











\Section{Problem statement}

In the Boussinesq approximation~[\citen{ulg:bib:Chandrasekhar},~\citen{ulg:bib:Aristov1}] in the Cartesian coordinates (the $Oz$  axis is directed upwards), the system of equations describing the heat convection of a viscous incompressible fluid are written as

\begin{gather}

\frac{\partial \bf{V}}{\partial t}+\left(\bf{V}\cdot\nabla\right)\bf{V}=-\nabla \textit{P} + \nu \Delta \bf{V}+g \beta \textit{T} \bold{k},

\\    

      \label{eq:Aristov_Prosviryakov:NSE}

\frac{\partial T}{\partial t}+\bf{V}\cdot\nabla \textit{T}=\chi \Delta \textit{T},

\\     

\nabla \cdot \bf{V}=0.

 \end{gather}

The system of equations~\eqref{eq:Aristov_Prosviryakov:NSE} involves the following symbols: $\bf{V}=$ $\left(V_x , V_y , V_z\right)$  is the flow velocity vector;  $P$  is the pressure deviation from hydrostatic, which is related to the constant average fluid density  $\rho;$  $T$ is the deviation from the average temperature;   $\nu$  and $\chi$ are the coefficients of kinematic viscosity and thermal diffusivity of the fluid, respectively; $\beta$ is the temperature coefficient of fluid volumetric expansion; $\bold{k}$  is the unit vector directed vertically upwards, $\nabla$  is the Hamilton operator,  $\Delta$ is the Laplace operator (Laplacian).



The solution of system~\eqref{eq:Aristov_Prosviryakov:NSE} is sought in the approximation to the large scale. For large-scale flows, the characteristic horizontal scale greatly exceeds in magnitude the characteristic dimension during the entire fluid motion. Free surface curvature can in this scale be neglected, the vertical velocity $V_z$  being assumed zero. Thus, system~\eqref{eq:Aristov_Prosviryakov:NSE} is transformed to the following system:

\begin{gather}

\frac{\partial V_x}{\partial t}+V_x \frac{\partial V_x}{\partial x}+V_y \frac{\partial V_x}{\partial y}=-\frac{\partial P}{\partial x} + \nu \Bigl(\frac{\partial^2 V_x}{\partial x^2}+\frac{\partial^2 V_x}{\partial y^2}+\frac{\partial^2 V_x}{\partial z^2}\Bigr),

\\       

\frac{\partial V_y}{\partial t}+V_x \frac{\partial V_y}{\partial x}+V_y \frac{\partial V_y}{\partial y}=-\frac{\partial P}{\partial y} + \nu \Bigl(\frac{\partial^2 V_y}{\partial x^2}+\frac{\partial^2 V_y}{\partial y^2}+\frac{\partial^2 V_y}{\partial z^2}\Bigr),

\\

       \label{eq:Aristov_Prosviryakov:system}

\frac{\partial P}{\partial z}=g\beta T,

\\

\frac{\partial T}{\partial t}+V_x \frac{\partial T}{\partial x}+V_y \frac{\partial T}{\partial y}=\chi \Bigl(\frac{\partial^2 T}{\partial x^2}+\frac{\partial^2 T}{\partial y^2}+\frac{\partial^2 V_x}{\partial z^2}\Bigr),

\\

\frac{\partial V_x}{\partial x}+\frac{\partial V_y}{\partial y}=0.

 \end{gather}

 

Obviously, system \eqref{eq:Aristov_Prosviryakov:system} is nonlinear and overdetermined, four unknown functions of velocity and temperature being required from five equations. The hydrodynamic and temperature fields are computed as~\cite{ulg:bib:Aristov1}

\begin{gather}

V_x=U+xu_1+y u_2,

\quad 

V_y=V+xv_1+y v_2,

\\

       \label{eq:Aristov_Prosviryakov:solution}

P=P_0+xP_1+yP_2+\frac{x^2}{2}P_{11}+\frac{y^2}{2}P_{22}+xyP_{12},

\\

T=T_0+xT_1+yT_2+\frac{x^2}{2}T_{11}+\frac{y^2}{2}T_{22}+xyT_{12}.     

 \end{gather}

 

For horizontal coordinates, all the functions in the formulae of system \eqref{eq:Aristov_Prosviryakov:solution}  depend on the variable $z$  ant time $t.$ Substituting relation \eqref{eq:Aristov_Prosviryakov:solution} into system \eqref{eq:Aristov_Prosviryakov:system}, we arrive at the system

\begin{gather}

\hat{L}U+Uu_1+Vu_2+P_1=0,

\\

\hat{L}u_1+u^{2}_1+v_1u_2+P_{11}=0,

\\

\hat{L}u_2+u_1u_2+u_2v_2+P_{12}=0,

\\

\hat{L}V+Uv_1+Vv_2+P_{2}=0,

\\

\hat{L}v_1+u_1v_1+v_1v_2+P_{12}=0,

\\

\hat{L}v_2+u_2v_1+v^{2}_2+P_{22}=0,

\\

       \label{eq:Aristov_Prosviryakov:main system }

\hat{M}T_0+UT_1+VT_2-\chi\left(T_{11}+T_{22}\right)=0,

\\

\hat{M}T_1+UT_{11}+VT_{12}+u_1T_1+v_1T_2=0,

\\

\hat{M}T_2+UT_{12}+VT_{22}+u_2T_1+v_2T_2=0,

\\

\hat{M}T_{11}+2u_1T_{11}+2v_1T_{12}=0,

\\

\hat{M}T_{22}+2u_2T_{12}+2v_2T_{22}=0,

\\

\hat{M}T_{12}+u_1T_{12}+u_2T_{11}+v_1T_{22}+v_2T_{12}=0,

\\

u_1+v_2=0,

\\

\frac{\partial P_0}{\partial z}=g\beta T_0,

\quad 

\frac{\partial P_1}{\partial z}=g\beta T_1,

\quad 

\frac{\partial P_2}{\partial z}=g\beta T_2,

\\

\frac{\partial P_{11}}{\partial z}=g\beta T_{11},

\quad 

\frac{\partial P_{12}}{\partial z}=g\beta T_{12},

\quad 

\frac{\partial P_{22}}{\partial z}=g\beta T_{22}.

 \end{gather}

Here, the following differential parabolic operators in partial derivatives are intro\-duced:

$$\hat{L}=\frac{\partial }{\partial t}-\nu\frac{\partial^2 }{\partial z^2},\quad 

\hat{M}=\frac{\partial }{\partial t}-\chi\frac{\partial^2 }{\partial z^2}. 

$$





\smallskip





\Section{Analyzing the solvability of system}

In what follows we do not discuss the solvability of system \eqref{eq:Aristov_Prosviryakov:main system } as a whole, but restrict our analysis to one special case. We will seek a solution to the system of equations when the boundaries of the fluid layer rotate. In~\cite{ulg:bib:Aristov2} the validity of equalities for the velocity gradients in the simulation of axisymmetric flows within class \eqref{eq:Aristov_Prosviryakov:solution} was demonstrated,

\begin{gather}

       \label{eq:Aristov_Prosviryakov:sol1}

u_1=0,

\quad 

v_2=0,       

\quad 

u_2=-v_1.       

 \end{gather}

Next the lower boundary of the layer  $z=0$ is assumed to be immobile. The upper boundary describable by the equation $z=h$  rotates with constant angular velocity  $\Omega$. Using equalities \eqref{eq:Aristov_Prosviryakov:sol1}, we obtain the following solutions to system~\eqref{eq:Aristov_Prosviryakov:main system }, which determine a number of functions in class \eqref{eq:Aristov_Prosviryakov:solution} for the coordinates  $x$ and  $y$:

\begin{gather}

u_2=\frac{\Omega z}{h}, 

\quad 

v_1=-\frac{\Omega z}{h},

\quad 

P_{12}=0, 

\\

\label{eq:Aristov_Prosviryakov:sol2}

 P_{11}=P_{22}=-u_2v_1=\frac{\Omega^2 z^2}{h^2},

\quad 

T_{12}=0, 

\\ g\beta T_{11}=g\beta T_{22}=\frac{2\Omega^2 z}{h^2}.

\end{gather}



For the homogeneous summands in expressions \eqref{eq:Aristov_Prosviryakov:solution}, due to solution \eqref{eq:Aristov_Prosviryakov:sol2}, system \eqref{eq:Aristov_Prosviryakov:main system } yields two closed subsystems. First system has of the form

\begin{gather}

\hat{L}U+\frac{\Omega z}{h}V+P_1=0, 

\\

\hat{L}V-\frac{\Omega z}{h}U+P_{2}=0,

\\

\label{eq:Aristov_Prosviryakov:sys1}

\hat{M}T_1+\frac{2\Omega^2 z}{g\beta h^2}U-\frac{\Omega z}{h}T_2=0, 

\\ \hat{M}T_2+\frac{2\Omega^2 z}{g\beta h^2}V+\frac{\Omega z}{h}T_1=0, 

\\

\frac{\partial P_1}{\partial z}=g\beta T_1, 

\quad  

\frac{\partial P_2}{\partial z}=g\beta T_2.

\end{gather}



Second system has of the form

\begin{gather}

\label{eq:Aristov_Prosviryakov:sys2}

\hat{M}T_0+UT_1+VT_2=\frac{4\chi\Omega^2 z}{g\beta h^2},

\\

\frac{\partial P_0}{\partial z}=g\beta T_0.

\end{gather}



Note that the integration of system \eqref{eq:Aristov_Prosviryakov:sys2} is possible only after the integration of system \eqref{eq:Aristov_Prosviryakov:sys1}. In order to find a solution to system \eqref{eq:Aristov_Prosviryakov:sys1}, we introduce the following complex functions:

$$

W=U+iV,\quad 

S=P_1+iP_2,\quad 

\Theta=g\beta\left(T_1+iT_2\right).$$

Here  $i$ is an imaginary unit. Thus, system \eqref{eq:Aristov_Prosviryakov:sys1} is in this case transformed to the following form:

 \begin{gather}

(\hat{L}-i\Omega z) W+S=0,

\\

       \label{eq:Aristov_Prosviryakov:sys1complex}

(\hat{M}+i\Omega z) \Theta+2\Omega^2 z W=0,       

\\

\frac{d S}{dz}=\Theta.      

 \end{gather}

 

To reduce the order of the system, the following transformation is made:

 \begin{gather}

       \label{eq:Aristov_Prosviryakov:trans}

S=S_1+i\Omega z W,

\quad 

\Theta=\Theta_1+2i\Omega W.    

\end{gather}

Substituting \eqref{eq:Aristov_Prosviryakov:trans} into system \eqref{eq:Aristov_Prosviryakov:sys1complex}, we obtain the equalities

 \begin{gather}

\nu\frac{d^2 W}{dz^2}+S_1=0,

\\

       \label{eq:Aristov_Prosviryakov:trans2}

\chi\frac{d^2 \Theta_1}{d z^2}+i\Omega \Bigl(z \Theta_1- 2\frac{\chi}{\nu}S_1\Bigr)=0,       

\\

\frac{d S_1}{dz}=\Theta_1-i\Omega z^2 \frac{d}{dz}\Bigl(\frac{W}{z}\Bigr). 

 \end{gather}

 

For further transformations, we use the differential identity

$$

\frac{d^2 W}{dz^2}=\frac{1}{z}\frac{d}{dz}\left(z^2\frac{d}{dz}\Bigl(\frac{W}{z}\Bigr)\right).

$$

This identity holds for any twice-differentiable function. Consequently, by differentiating the last equation in system \eqref{eq:Aristov_Prosviryakov:trans2}, we arrive at a system for determining the functions  $\Theta_1$  and  $S_1,$

 \begin{gather}

       \label{eq:Aristov_Prosviryakov:trans3}

\chi\frac{d^2 \Theta_1}{d z^2}+i\Omega \left(z \Theta_1- 2\frac{\chi}{\nu}S_1\right)=0,

\quad 

\nu\frac{d^2 S_1}{dz^2}=i\Omega z S_1 +\nu \frac{d \Theta_1}{dz}.          

 \end{gather}

 

Taking into account that equations \eqref{eq:Aristov_Prosviryakov:trans2} offer an explicit expression for  $W$, we obtain an integral of the initial system with automatic order reduction, system \eqref{eq:Aristov_Prosviryakov:trans3} being transformed to the operator equation

$$

\left[\Bigl(\nu\frac{d^2}{dz^2}-i\Omega z\Bigr)\Bigl(\chi\frac{d^2}{dz^2}+i\Omega z\Bigr)-2i\Omega\chi\frac{d}{dz}\right]\Theta_1=0.

$$

Simplifying the latter equality, we derive the following fourth-order ordinary differential equation with complex coefficients:

 \begin{gather}

       \label{eq:Aristov_Prosviryakov:trans4}

\Bigl[\nu\chi\frac{d^4}{dz^4}+i\Omega\left(\nu-\chi\right)\frac{d^2}{dz^2}z+\Omega^2 z^2\Bigr]\Theta_1=0.    

 \end{gather}



\smallskip







\Section{Solving equation when the dissipative coefficients coincide}

 We assume in equation \eqref{eq:Aristov_Prosviryakov:trans4} that $\nu=\chi$  (the Prandtl number equals one). In this case, equation \eqref{eq:Aristov_Prosviryakov:trans4} acquires the form

  \begin{gather}

       \label{eq:Aristov_Prosviryakov:Pr1}

\frac{d^4 \Theta_1}{dz^4}+\left(36\sigma^3 z\right)^2 \Theta_1=0,    

 \end{gather}

where  $\sigma=\frac{1}{6}\sqrt[3]{\frac{\Omega}{6\nu h}}$ is the dispersion relation. For the convenience of computations, in order to illustrate the physical meaning of the solutions and to make the notation concise, we introduce a new variable  $Z=\frac{z}{h \sigma}.$ The general solution of equation \eqref{eq:Aristov_Prosviryakov:Pr1} is written as

\begin{multline}

\label{eq:Aristov_Prosviryakov:soltemp}

\Theta_1={^{3}_{0}F} \Bigl[\Bigl\{\frac{1}{2}, \frac{2}{3}, \frac{5}{6}\Bigr\},-Z^6\Bigr]C_1 +

{^{3}_{0}F} 

\Bigl[\Bigl\{\frac{2}{3}, \frac{5}{6}, \frac{7}{6}\Bigr\},-Z^6\Bigr] C_2 Z +

\\

+

{^{3}_{0}F} \Bigl[\Bigl\{\frac{5}{6}, \frac{7}{6}, \frac{4}{3}\Bigr\},-Z^6\Bigr] C_3 Z^2  +

{^{3}_{0}F} \Bigl[\Bigl\{\frac{7}{6}, \frac{4}{3},\frac{3}{2}\Bigr\},-Z^6\Bigr] C_4 Z^3  . 

\end{multline}

Here,  ${^{3}_{0}F}=F\left(a; b; Z\right)=F\left(a_1, a_2, \ldots, a_p; b_1, b_2, \ldots, b_q; Z\right)$ is a generalized hypergeometric function~\cite{ulg:bib:Polyanin}. The choice of using the generalized hypergeometric function for writing the solutions is dictated by notation conciseness. Note that there is another form of writing the solution of equation \eqref{eq:Aristov_Prosviryakov:Pr1}, which is based on using the Mittag--Leffler functions~\cite{ulg:bib:Polyanin}. Note that each particular solution \eqref{eq:Aristov_Prosviryakov:soltemp} is a function of a real variable and that the general solution is considered in the complex plane. Thus, at least one integration constant is a complex number.



We now turn to finding the generalized pressure $S_1$  from the differential equation \eqref{eq:Aristov_Prosviryakov:trans2} 

 \begin{gather}

       \label{eq:Aristov_Prosviryakov:eqpr}

S_1=-\frac{i}{2\sigma}\frac{d^2 \Theta_1}{d z^2}+\frac{z}{2}\Theta_1.

\end{gather}



The substitution of solution \eqref{eq:Aristov_Prosviryakov:soltemp} into equation \eqref{eq:Aristov_Prosviryakov:eqpr} gives

\begin{multline}

\label{eq:Aristov_Prosviryakov:solstress}

S_1= {^{3}_{0}F} 

\Bigl[\Bigl\{\frac{1}{2}, \frac{2}{3}, \frac{5}{6}\Bigr\},-Z^6\Bigr] \frac{C_1}{2}Z+

{^{3}_{0}F} 

\Bigl[\Bigl\{\frac{2}{3}, \frac{5}{6}, \frac{7}{6}\Bigr\},-Z^6\Bigr] \frac{C_2}{2}Z^2+

\\

+{^{3}_{0}F} 

\Bigl[\Bigl\{\frac{5}{6}, \frac{7}{6}, \frac{4}{3},\Bigr\},Z\Bigr] \sqrt[3]{36\sigma^2} \frac{2\cdot 5C_3}{6!}Z^3+

{^{3}_{0}F} 

\Bigl[\Bigl\{\frac{7}{6}, \frac{4}{3}, \frac{3}{2}\Bigr\},Z\Bigr] \sigma \frac{2\cdot 5C_4}{6!}Z^4-

\\

-\frac{1}{72\sigma}i\left\{ {^{3}_{0}F} 

\Bigl[\Bigl\{\frac{5}{6}, \frac{7}{6}, \frac{4}{3}\Bigr\},Z\Bigr] \sqrt[3]{676\sigma^2} C_3 Z^3 

-{^{3}_{0}F} 

\Bigl[\Bigl\{\frac{7}{6}, \frac{4}{3}, \frac{3}{2}\Bigr\},Z\Bigr] 6 \sigma C_4 Z+

\right.

\\

+{^{3}_{0}F} 

\Bigl[\Bigl\{\frac{5}{6}, \frac{7}{6}, \frac{4}{3}\Bigr\},Z\Bigr] 3 \sigma^2 C_1 Z^4 +

\\

+{^{3}_{0}F} 

\Bigl[\Bigl\{\frac{5}{3}, \frac{11}{6}, \frac{13}{6},\Bigr\},Z\Bigr] \frac{2^4\cdot 3^3\sigma^2}{7!}\sqrt[3]{\frac{3\sigma}{4}} C_2 Z^5

+

\\

+{^{3}_{0}F} 

\Bigl[\Bigl\{\frac{5}{3}, \frac{11}{6}, \frac{13}{6},\Bigr\},Z\Bigr] \frac{2^4\cdot 3^4\sigma^2}{7!}\sqrt[3]{6\sigma} C_2 Z^5+

\\

+{^{3}_{0}F} 

\Bigl[\Bigl\{\frac{11}{6}, \frac{13}{6}, \frac{7}{3} \Bigr\},Z\Bigr] \frac{2^2\cdot 3^4\sigma^2}{7!}\sqrt[3]{\frac{9\sigma^2}{2}} C_3 Z^6+

\\

+{^{3}_{0}F} \Bigl[\Bigl\{\frac{13}{6}, \frac{7}{3}, \frac{5}{2} \Bigr\},Z\Bigr] \frac{2\cdot 5\cdot 11\sigma^3}{7!} C_4 Z^7-

\\

-{^{3}_{0}F} \Bigl[\Bigl\{\frac{5}{2}, \frac{8}{3}, \frac{17}{6} \Bigr\},Z\Bigr] \frac{2^7\cdot 3^3\cdot 7\sigma^4}{11!} C_1 Z^{10}-

\\

-{^{3}_{0}F} \Bigl[\Bigl\{\frac{8}{3}, \frac{17}{6}, \frac{19}{6} \Bigr\},Z\Bigr] \frac{2^9\cdot 3^6\sigma^4}{13!}\sqrt[3]{\frac{3\sigma}{4}} C_2 Z^{11}-

\\

-{^{3}_{0}F} \Bigl[\Bigl\{\frac{17}{6}, \frac{19}{6}, \frac{10}{3} \Bigr\},Z\Bigr] \frac{2^8\cdot 3^5\cdot 5\sigma^4}{14!}\sqrt[3]{\frac{9\sigma^2}{2}} C_3 Z^{12}-

\\

\left.

-{^{3}_{0}F} \Bigl[\Bigl\{\frac{19}{6}, \frac{10}{3}, \frac{7}{2} \Bigr\},Z\Bigr] \frac{2^7\cdot 3^2\cdot 5\cdot 11\sigma^5}{14!} C_4 Z^{13}

\right\}.

\end{multline}





Complex velocity $W$  can be expressed in two ways. Firstly, it can be obtained by integrating the first equation in system \eqref{eq:Aristov_Prosviryakov:sys1complex}, but it is a second-order equation. Hence, we obtain a solution depending on six integration constants, whereas system \eqref{eq:Aristov_Prosviryakov:trans2} is a fifth-order equation system. Thus our problem is complicated by additional analysis of overdetermination. A correct expression for velocity $W$  can be obtained by the integration of the third equation in system \eqref{eq:Aristov_Prosviryakov:sys1complex},

\begin{multline}

\label{eq:Aristov_Prosviryakov:solvelo}

\Omega W=C_5 Z+ \left({^{3}_{0}F} \Bigl[\Bigl\{\frac{2}{3}, \frac{5}{6}, \frac{7}{6} \Bigr\},Z\Bigr]-1\right) \frac{i}{2} \sqrt[3]{\frac{\sigma}{36}} C_2 Z-

\\

-{^{4}_{1}F} \Bigl[\Bigl\{-\frac{1}{6} \Bigr\}, \Bigl\{\frac{1}{2},\frac{2}{3}, \frac{5}{6}, \frac{5}{6} \Bigr\},Z\Bigr]\frac{2i \Gamma\left(-\frac{1}{6}\right)}{4!\Gamma\left(\frac{5}{6}\right)} C_1 Z^2+

\\

+{^{4}_{1}F} \Bigl[\Bigl\{-\frac{1}{6} \Bigr\},\Bigl\{\frac{5}{6},\frac{7}{6}, \frac{4}{3}, \frac{3}{2} \Bigr\},Z\Bigr]\frac{2\cdot 5\Gamma\left(-\frac{1}{6}\right)}{6!\Gamma\left(\frac{5}{6}\right)} C_4+

\\

+{^{4}_{1}F} \Bigl[\Bigl\{\frac{1}{6} \Bigr\},\Bigl\{\frac{5}{6},\frac{7}{6}, \frac{7}{6}, \frac{4}{3} \Bigr\},Z\Bigr]\frac{2\cdot 5i \Gamma\left(\frac{1}{6}\right)}{6!\Gamma\left(\frac{7}{6}\right)} \sqrt[3]{\frac{\sigma^2}{6}} C_3 Z^2+

\\

+{^{4}_{1}F} \Bigl[\Bigl\{\frac{1}{3} \Bigr\},\Bigl\{\frac{7}{6},\frac{4}{3}, \frac{4}{3}, \frac{3}{2} \Bigr\},Z\Bigr]\frac{2^4\cdot 3\cdot 5\cdot 7i \sigma \Gamma\left(\frac{1}{3}\right)}{9!\Gamma\left(\frac{4}{3}\right)} C_4 Z^3-

\\

-{^{4}_{1}F} \Bigl[\Bigl\{\frac{1}{3} \Bigr\},\Bigl\{\frac{4}{3},\frac{3}{2}, \frac{5}{3}, \frac{11}{6} \Bigr\},Z\Bigr]\frac{2^2\cdot 5\sigma \Gamma\left(\frac{1}{3}\right)}{6!\Gamma\left(\frac{4}{3}\right)} C_1 Z^3-

\\

-{^{4}_{1}F} \Bigl[\Bigl\{\frac{1}{2} \Bigr\},\Bigl\{\frac{3}{2},\frac{5}{3}, \frac{11}{6}, \frac{13}{6} \Bigr\},Z\Bigr]\frac{\sigma }{4!}\sqrt[3]{\frac{\sigma}{36}} C_2 Z^4-

\\

-{^{4}_{1}F} \Bigl[\Bigl\{\frac{2}{3} \Bigr\},\Bigl\{\frac{5}{3},\frac{11}{6}, \frac{13}{6}, \frac{7}{3} \Bigr\},Z\Bigr] \frac{2^4\cdot 3^2\cdot 7\sigma \Gamma\left(\frac{2}{3}\right)}{9!\Gamma\left(\frac{5}{3}\right)}\sqrt[3]{\frac{\sigma}{6}} C_3 Z^5-

\\

-{^{4}_{1}F} \Bigl[\Bigl\{\frac{5}{6} \Bigr\},\Bigl\{\frac{3}{2},\frac{5}{3}, \frac{11}{6}, \frac{11}{6} \Bigr\},Z\Bigr] \frac{2^3\cdot 3^2\cdot 7i\sigma^2 \Gamma\left(\frac{5}{6}\right)}{9!\Gamma\left(\frac{11}{6}\right)} C_1 Z^6-

\\

-{^{4}_{1}F} \Bigl[\Bigl\{\frac{5}{6} \Bigr\},\Bigl\{\frac{11}{6},\frac{13}{6}, \frac{7}{3}, \frac{5}{2} \Bigr\},Z\Bigr] \frac{2^3\cdot 5^2\cdot 11\cdot 13\cdot 83\sigma^2 \Gamma\left(\frac{5}{6}\right)}{13! \Gamma\left(\frac{11}{6}\right)} C_4 Z^6-

\\ 

-{^{4}_{1}F} \Bigl[\Bigl\{\frac{7}{6} \Bigr\},\Bigl\{\frac{11}{6},\frac{13}{6}, \frac{13}{6}, \frac{7}{3} \Bigr\},Z\Bigr] \frac{2\cdot 3^2i\sigma^2 \Gamma\left(\frac{7}{6}\right)}{9! \Gamma\left(\frac{13}{6}\right)}\sqrt[3]{\frac{\sigma^2}{6}} C_3 Z^8-

\\

-{^{4}_{1}F} \Bigl[\Bigl\{\frac{4}{3} \Bigr\},\Bigl\{\frac{13}{6},\frac{7}{3}, \frac{7}{3}, \frac{5}{2} \Bigr\},Z\Bigr] \frac{2^3\cdot 5^2\cdot 11\cdot 13 i\sigma^3 \Gamma\left(\frac{4}{3}\right)}{13!\Gamma\left(\frac{7}{3}\right)} C_4 Z^9+

\\

+{^{4}_{1}F} \Bigl[\Bigl\{\frac{4}{3} \Bigr\},\Bigl\{\frac{7}{3},\frac{5}{2}, \frac{8}{3}, \frac{17}{6} \Bigr\},Z\Bigr] \frac{2^3\cdot 3\cdot 5\cdot 7\sigma^3 \Gamma\left(\frac{4}{3}\right)}{11! \Gamma\left(\frac{7}{3}\right)} C_1 Z^9+

\\

+{^{4}_{1}F} \Bigl[\Bigl\{\frac{3}{2} \Bigr\},\Bigl\{\frac{5}{2},\frac{8}{3}, \frac{17}{6}, \frac{19}{6} \Bigr\},Z\Bigr] \frac{2^7\cdot 3^5\sigma^3}{13!} \sqrt[3]{\frac{\sigma}{36}} C_2 Z^{10}+

\\

+{^{4}_{1}F} \Bigl[\Bigl\{\frac{5}{3} \Bigr\},\Bigl\{\frac{8}{3},\frac{17}{6}, \frac{19}{6}, \frac{10}{3} \Bigr\},Z\Bigr] \frac{2^3\cdot 3^4\cdot 5\sigma^3 \Gamma\left(\frac{5}{3}\right)} {13! \Gamma\left(\frac{8}{3}\right)} \sqrt[3]{\frac{\sigma^2}{6}} C_3 Z^{11}+

\\

+{^{4}_{1}F} \Bigl[\Bigl\{\frac{11}{6} \Bigr\},\Bigl\{\frac{17}{6},\frac{19}{6}, \frac{10}{3}, \frac{7}{2} \Bigr\},Z\Bigr] \frac{2^6\cdot 5\cdot 11\sigma^4 \Gamma\left(\frac{11}{6}\right)} {14! \Gamma\left(\frac{17}{6}\right)} C_4 Z^{12}-

\\

-{^{4}_{1}F} \Bigl[\Bigl\{\frac{7}{3} \Bigr\},\Bigl\{\frac{10}{3},\frac{7}{2}, \frac{11}{3}, \frac{23}{6} \Bigr\},Z\Bigr] \frac{2^6\cdot 3\cdot 7^2\cdot 13\sigma^5 \Gamma\left(\frac{7}{3}\right)} {17! \Gamma\left(\frac{10}{3}\right)} C_1 Z^{15}-

\\

-{^{4}_{1}F} \Bigl[\Bigl\{\frac{5}{2} \Bigr\},\Bigl\{\frac{7}{2},\frac{11}{3}, \frac{23}{6}, \frac{25}{6} \Bigr\},Z\Bigr] \frac{2^8\cdot 3^7\cdot 7\sigma^5 } {19!} \sqrt[3]{\frac{\sigma}{36}} C_2 Z^{16}-

\\

-{^{4}_{1}F} \Bigl[\Bigl\{\frac{8}{3} \Bigr\},\Bigl\{\frac{11}{3},\frac{23}{6}, \frac{25}{6}, \frac{13}{3} \Bigr\},Z\Bigr] \frac{2^7\cdot 3^6\cdot 5\sigma^5 \Gamma\left(\frac{8}{3}\right)C_3 Z^{17}} {19! \Gamma\left(\frac{11}{3}\right)} \sqrt[3]{\frac{\sigma^2}{6}} -

\\

-{^{4}_{1}F} \Bigl[\Bigl\{\frac{8}{3} \Bigr\},\Bigl\{\frac{23}{6},\frac{25}{6}, \frac{13}{3}, \frac{9}{2} \Bigr\},Z\Bigr] \frac{2^8\cdot 3^2\cdot 5^2\cdot 11\cdot 17\sigma^6 \Gamma\left(\frac{17}{6}\right)} {21! \Gamma\left(\frac{23}{6}\right)} \sqrt[3]{\frac{\sigma^2}{6}} C_4 Z^{18}.

\end{multline}

Here $\Gamma\left(\lambda\right)= \displaystyle\int_{0}^{+\infty}r^{\lambda-1}\exp\left(-r\right)dr$ is the Euler gamma function~\cite{ulg:bib:Polyanin}. Solutions~\eqref{eq:Aristov_Prosviryakov:soltemp}, \eqref{eq:Aristov_Prosviryakov:solstress},~\eqref{eq:Aristov_Prosviryakov:solvelo} are written as a linear  combination of generalized hypergeometric functions in the form of infinite power series. Consequently, examination of solution convergence is an urgent problem. Using the properties of convergence of these functions, we can state that the convergence radius for each solution equals infinity~\cite{ulg:bib:Polyanin}. It means that the solution of our problem always converges.



Let us now formulate the boundary conditions to determine the integration constants of equation system \eqref{eq:Aristov_Prosviryakov:sys1}. Without any claim to solution generality and without studying the shapes of the fluid, we give a simple example of such solutions. Let the simplest and most physically meaningful solution be a flow between two infinite disks rotating with free velocities. This is possible, since the solution can always be shifted along the $Oz$-axis, an appropriate distance between the disks being chosen. The disks must be nonuniformly heated to a proper temperature (the temperature must be related to the radius by a well-defined law). Assume that the lower boundary is isothermal (the temperature of the lower disk is zero); the adhesion condition holds on both boundaries. Thus, the following equalities are valid:

$$\text{when} \ \ z=0, \ \ U=V=0, \ \ T_1=T_2=0;$$

$$\text{when} \ \ z=h, \ \ U=A\cos\alpha, \ \ V=A\sin\alpha, \ \ T_1=B, \ \ T_2=C, \ \ P_1=\tau_1, \ \ P_2=\tau_2.$$



The conditions for determining the integration constants, written on the boundary, do not determine the integration constants uniquely. The system of linear algebraic equations for determining the integration constants is underdetermined, since there are few boundary conditions. To close the system of equations, we specify the fluid flow rate conditions as

$$\int _{0}^{h}Udz =Q_1,\quad  \int_{0}^{h}Vdz =Q_2.$$



When the flow rate is specified, it is generally assumed that  $Q_1=Q_2=0$~\cite{ulg:bib:Aristov2}. It is a physically justified convention~[\citen{ulg:bib:Aristov1},~\citen{ulg:bib:Aristov2}]. However, it is possible to obtain a solution by specifying a nonzero flow rate. Next, not limiting the generality of the reasoning, we will study the solutions simulating fluid motions with the zero flow rate. 



For the above-introduced complex velocities, pressure and temperature ~\eqref{eq:Aristov_Prosviryakov:soltemp}, \eqref{eq:Aristov_Prosviryakov:solstress},~\eqref{eq:Aristov_Prosviryakov:solvelo}, in view of transformations \eqref{eq:Aristov_Prosviryakov:trans}, the following boundary conditions are valid:

\begin{gather}

\label{eq:Aristov_Prosviryakov:bond}

\text{when} \ \ z=0: \ \ W=0, \ \ \Theta_1=0;

\\

\text{when} \ \ z=h: \ \ W=A\exp\left(i\alpha\right), \ \ \Theta_1=B+2A\Omega\sin\alpha+i\left(C-2A\Omega\cos\alpha\right),

\\

\int_{0}^{h}Wdz =0.

\end{gather}



With the boundary conditions \eqref{eq:Aristov_Prosviryakov:bond}, it follows from equality \eqref{eq:Aristov_Prosviryakov:soltemp} that  $C_1=0$, and from solution \eqref{eq:Aristov_Prosviryakov:solvelo} that  $C_4=0$. Thus, the expressions for the determination of hydrodynamic fields become much simpler. Remember that, in the computation of the integration constants, the variable $z=h\sigma Z$  is to be replaced.











\Section{Analysis of the solutions}

 Using conditions \eqref{eq:Aristov_Prosviryakov:bond}, it is easy, but rather cumbersome, to compute the complex  $C_2,$ $C_3,$ and $C_5$  which are related to the dispersion relation   $\sigma$ (kinematic viscosity) by the fractional linear law. Analyzing conditions \eqref{eq:Aristov_Prosviryakov:bond}, we find that the boundary value problem under study is meaningful at any positive value of  $\sigma.$ Nevertheless, the structure and behavior of the solutions depend considerably on the value of  $\sigma.$



It is well known that, to analyze the polynomial solutions of system \eqref{eq:Aristov_Prosviryakov:sys1}, it would suffice to localize the roots of the linear combination of the generalized hypergeometric functions~\cite{ulg:bib:Polyanin}. Any hypergeometric function is known to be monotonically increasing with a zero value at the origin. Therefore, the appearance of other zero values is possible only when there is a combination of functions with at least one negative weight coefficient~\cite{ulg:bib:Polyanin}. The appearance of negative coefficients depends on $\sigma$  and the function domain, which also depends on the dispersion relation $Z\in\left[0, \frac{1}{\sigma}\right].$  Temperature is characterized by three types of solutions. 

When  $\sigma\geq\frac{3}{2},$ the function describing temperature is a strictly monotonically increasing or decreasing function (fig.~\ref{fig_prosv1}). 

If the evaluation  $\frac{9}{10}<\sigma<\frac{3}{2}$  is valid for the dispersion relation, there is exactly one zero value of temperature inside the fluid layer (fig.~\ref{fig_prosv2}). With the further decrease, temperature gets a finite number of zeros (fig.~\ref{fig_prosv3}). Consequently, the fluid layer is divided into the regions with positive and negative temperature, which are characterized by alternating signs. The analysis of pressure is similar to that of temperature. 





Velocities are characterized by only two qualitative behaviors of the functions (figs.~\ref{fig_prosv2} and ~\ref{fig_prosv3}). Thus, at any  $\sigma,$ there appear fluid counterflows. It was shown in~[\citen{ulg:bib:Aristov3}--\citen{ulg:bib:Gorshkov}] that the appearance of counterflows is a purely nonlinear effect, i.\,e., the viscous forces are comparable to the convective acceleration, and even without Coriolis forces. However, when  $\sigma$ tends to zero, counterflow zones are formed, whose thickness decreases with  $\sigma.$ Other values of the dispersion relation are characteristic of velocity, and they are easy to compute. However, the values obtained for the temperature and pressure fields are a good reference point for refining the characteristic dispersion numbers of velocities.





When analyzing the solutions (fig.~\ref{fig_prosv3}), we can see their ``fluctuation''. In~\cite{ulg:bib:Aristov7} localized convective flows in terms were studied in the axisymmetric formulation. Solution ``damping'' was observed there. We use a wider class of solutions. There may be not only ``damping'' (when $\sigma$  is close to one), but also ``resonance'' initiation. Yet, the solution does not become singular, and the Boussinesq approxi\-mation is not violated in the simulation of the convective flow. This difference is attributed to the account of the quadratic summands for temperature. If tempera\-ture is considered to be a linear function with respect to the horizontal coordinates, the solutions obtained in this study coincide with those discussed in~\cite{ulg:bib:Aristov4}.



\begin{figure}[h!]



\begin{tabular}{p{.45\textwidth}p{.45\textwidth}}

\centerline{\includegraphics[width=0.45\textwidth]{fig_prosv1}} &

\centerline{\includegraphics[width=0.45\textwidth]{fig_prosv2}} \\

\caption{A qualitative form of the solutions  $T_1,$ $T_2,$ $P_1,$ $P_2$ of system (\ref{eq:Aristov_Prosviryakov:sys1}) when $\sigma\geq\frac{3}{2}$ ($1-\tau_1>0$,  $\tau_2>0$, $A>0$, $B>0$, $C>0$; $2-\tau_1<0$, $\tau_2<0$, $A<0$, $B<0$, $C<0$) \label{fig_prosv1}} &

\caption{A qualitative form of the solutions  $T_1,$ $T_2,$ $P_1,$ $P_2,$ $U,$ $V$ of system (\ref{eq:Aristov_Prosviryakov:sys1}) when $\frac{9}{10}<\sigma<\frac{3}{2}$ ($1-\tau_1>0$,  $\tau_2>0$, $A>0$, $B>0$, $C>0$; $2-\tau_1<0$, $\tau_2<0$, $A<0$, $B<0$, $C<0$) \label{fig_prosv2} }

\\

\centerline{\includegraphics[width=0.45\textwidth]{fig_prosv3}} &



\vspace{-2cm}

\caption{A qualitative form of the solutions  $T_1,$ $T_2,$ $P_1,$ $P_2,$ $U,$ $V$ of system (\ref{eq:Aristov_Prosviryakov:sys1}) when $\sigma\leq\frac{9}{10}$ ($1-\tau_1>0$,  $\tau_2>0$, $B>0$, $C>0$; $2-\tau_1<0$, $\tau_2<0$, $B<0$, $C<0$) \label{fig_prosv3} } \\





\end{tabular}



\vspace{-3cm}



%%\noindent

%%\begin{minipage}[h]{0.48\textwidth}

%%

%%\smallskip \footnotesize

%%Fig.~\ref{fig1}.

%%A qualitative form of the solutions  $T_1,$ $T_2,$ $P_1,$ $P_2$ of system (\ref{eq:Aristov_Prosviryakov:sys1}) when $\sigma\geq\frac{3}{2}$ ($1-\tau_1>0$,  $\tau_2>0$, $A>0$, $B>0$, $C>0$; $2-\tau_1<0$, $\tau_2<0$, $A<0$, $B<0$, $C<0$) \label{fig1} 

%%\end{minipage} \hfill

%%\begin{minipage}[h]{0.5\textwidth}

%%\centering

%%\includegraphics[width=0.98\textwidth]{fig1}

%%\end{minipage}

%\end{figure}

%

%

%\begin{figure}[h!]

%\noindent

%\begin{minipage}[h]{0.48\textwidth}

%\caption{A qualitative form of the solutions  $T_1,$ $T_2,$ $P_1,$ $P_2,$ $U,$ $V$ of system (\ref{eq:Aristov_Prosviryakov:sys1}) when $\frac{9}{10}<\sigma<frac{3}{2}$ ($1-\tau_1>0$,  $\tau_2>0$, $A>0$, $B>0$, $C>0$; $2-\tau_1<0$, $\tau_2<0$, $A<0$, $B<0$, $C<0$) 

%\label{fig2}

% }

%\smallskip \footnotesize

%Fig.~\ref{fig2}.

%A qualitative form of the solutions  $T_1,$ $T_2,$ $P_1,$ $P_2,$ $U,$ $V$ of system (\ref{eq:Aristov_Prosviryakov:sys1}) when $\frac{9}{10}<\sigma<\frac{3}{2}$ ($1-\tau_1>0$,  $\tau_2>0$, $A>0$, $B>0$, $C>0$; \centerline{$2-\tau_1<0$, $\tau_2<0$, $A<0$, $B<0$, $C<0$)} \label{fig2} 

%\end{minipage} \hfill

%\begin{minipage}[h]{0.5\textwidth}

%\centering

%\includegraphics[width=0.98\textwidth]{fig2}

%\end{minipage}

%\end{figure}

%

%

%

%\begin{figure}[h!]

%\noindent

%\begin{minipage}[h]{0.48\textwidth}

%\caption{A qualitative form of the solutions  $T_1,$ $T_2,$ $P_1,$ $P_2,$ $U,$ $V$ of system (\ref{eq:Aristov_Prosviryakov:sys1}) when $\sigma\leq\frac{9}{10}$ ($1-\tau_1>0$,  $\tau_2>0$, $B>0$, $C>0$; $2-\tau_1<0$, $\tau_2<0$, $B<0$, $C<0$) 

%\label{fig3}

% }

%\smallskip \footnotesize

%Fig.~\ref{fig3}.

%A qualitative form of the solutions  $T_1,$ $T_2,$ $P_1,$ $P_2,$ $U,$ $V$ of system (\ref{eq:Aristov_Prosviryakov:sys1}) when $\sigma\leq\frac{9}{10}$ ($1-\tau_1>0$,  $\tau_2>0$, $B>0$, $C>0$; 

%\centerline{$2-\tau_1<0$, $\tau_2<0$, $B<0$, $C<0$)}

% \label{fig3} 

%\end{minipage} \hfill

%\begin{minipage}[h]{0.5\textwidth}

%\centering

%\includegraphics[width=0.98\textwidth]{fig3}

%\end{minipage}

\end{figure}









\newpage















\Section{Solution of equation  with the Prandtl number different from one}

An important particular case of equation \eqref{eq:Aristov_Prosviryakov:trans4} was studied above; namely, when the values of kinematic viscosity and thermal diffusivity of the fluid coincide. 

Let us now write the solution of equation \eqref{eq:Aristov_Prosviryakov:trans4} in the general case:

\begin{multline}

\label{eq:Aristov_Prosviryakov:soltempPr}

\Theta_1=\text{\sf Ai}\Bigl(\xi+\frac{4\Omega z}{\sqrt{\nu\chi}}\sqrt[3]{\frac{\nu^2\chi^2} {16\Omega^2}}\Bigr)

\left(C_1\text{\sf Ai}\left(\xi\right)+C_2\text{\sf Bi}\left(\xi\right)\right)+

\\

+\text{\sf Bi}\Bigl(\xi+\frac{4\Omega z}{\sqrt{\nu\chi}}\sqrt[3]{\frac{\nu^2\chi^2}{16\Omega^2}}\Bigr)\left(C_3\text{\sf Ai}\left(\xi\right)+C_4\text{\sf Bi}\left(\xi\right)\right).

\end{multline}

Here $$\xi=\sqrt[3]{\nu^2\chi^2}\displaystyle\frac{ i\left(\nu-\chi\right)+2\sqrt{\nu\chi}\Omega z}{\nu\chi\sqrt[3]{16\Omega^2}}$$ is an auxiliary variable introduced in order to make the solution notation concise;  $\text{\sf Ai} \left(\xi\right)$   and $\text{\sf Bi}\left(\xi\right)$  are the Airy function and the associated function~\cite{ulg:bib:Polyanin}. Solution \eqref{eq:Aristov_Prosviryakov:soltempPr} is complex-valued, since it is written for arbitrary differences between $\nu$  and $\chi.$  This is done for the convenience of the solution analysis. Solution \eqref{eq:Aristov_Prosviryakov:soltempPr} being known, pressure  $S_1$ and velocity $W$ are computed similarly to the previous case.	



When the dissipative coefficients  $\nu$ and  $\chi$ of the nonisothermal fluid are equal, solution \eqref{eq:Aristov_Prosviryakov:soltempPr}  is reduced to formula \eqref{eq:Aristov_Prosviryakov:soltemp}. Studying the convergence of series \eqref{eq:Aristov_Prosviryakov:soltempPr}, we find that the solution expressed in terms of this series converges on the solution domain, i.\,e. layer thickness. 



The consideration of equation \eqref{eq:Aristov_Prosviryakov:trans4} with an arbitrary relation between $\nu$  and~$\chi$ offers solutions to system \eqref{eq:Aristov_Prosviryakov:sys1}, \eqref{eq:Aristov_Prosviryakov:sys2} in qualitative terms, see figs.~\ref{fig_prosv1} to~\ref{fig_prosv3}. In the simulation of the hydrodynamic fields by the linear combination of the Airy functions the behavior of the solutions is predictable from the very beginning, since these special functions satisfy the simplest differential equation having a point where the fluctuation of the solution is replaced by its exponential growth. It was shown above that the solution does not go off to infinity.



\smallskip



\Section{Conclusion}

The problem of simulating rotating fluid masses has a long history. To describe fluid motion (relative equilibrium), ordinary differential equa\-tions are used that enable one to obtain solutions simulating fluid flows by the theory of motion stability with the disturbance of the background (principal) flow. This paper discusses the exact solution to an overdetermined boundary value problem, which describes stationary flow at any scales where the Boussinesq approximation is valid. In this case, stationary dynamic fluid equilibria are chara\-cterized by the formation of counterflows, the sign of velocity being able to alternate several times, depending on the boundary conditions and the geometric dimensions of the fluid layer. A similar situation is true for temperature and pressure, namely, the formation of zones with the negative and positive values of temperature and pressure with respect to the reference value. 



\smallskip 

 



\AdditionalContent{Competing interests}{I declare that I have no competing interests.} 

\AdditionalContent{Author's Responsibilities}{I take full responsibility for submitting the final manuscript in print. I approved the final version of the manuscript.} 





\AdditionalContent{Funding}{The work was done within the state assignment from FASO Russia, theme No. AAAA-A18-118020790140-5.}









\begin{thebibliography}{28}



\Bibitem{ulg:bib:Newtoni}

\by Newton~I.

\book Opera quae exstant omnia

\bookinfo Faksimile-Neudruck der Ausgabe von Samuel Horsley, London 1779--1785 in f\"unf B\"anden. Band 1,~4

\lang In German

\zmath{0129.25102}

\publaddr Stuttgart-Bad Cannstatt

\publ Friedrich Frommann Verlag (G\"unther Holzboog)

\totalpages xx+592

\yr 1964



\Bibitem{ulg:bib:Turnbull}

\book The Correspondence of Isaac Newton

\bookvol 1. 1661-1675

\ed H.~W.~Turnbull

\publ Cambridge Univ. Press

\publaddr New York

\yr 1959

\totalpages xxxviii+468 





\Bibitem{ulg:bib:Poincare}

\by Poincar\'e~H.

\book Cin\'ematique et m\'ecanismes. Potentiel et m\'ecanique des fluides

\bookinfo Cours profess\'e \`a la Sorbonne

\lang In French

\zmath{30.0608.01}

\publaddr Paris

\publ G.~Carr\'e et C.~Naud

\totalpages iv+385 

\yr 1899



\Bibitem{ulg:bib:Lyapunov1}

\by Lyapunov~A.~M.

\paper On the stability of ellipsoidal equilibrium forms of rotating fluid

\inbook Collected works. \rm  Vol. III

\yr 1959

\publ Akad. Nauk SSSR

\publaddr Moscow

\pages 5--113

\zmath{0192.31203}

\lang In Russian



\Bibitem{ulg:bib:Lyapunov2}

\by Lyapunov~A.~M.

\paper On the equilibrium figures of rotating homogeneous liquid mass slightly different from ellipsoids

\inbook Collected works. \rm  Vol. IV

\yr 1959

\publ Akad. Nauk SSSR

\publaddr Moscow

\pages 5--645

\zmath{0192.31204}

\lang In Russian



\Bibitem{ulg:bib:Chandrasekhar}

\by Chandrasekhar~S.

\book Ellipsoidal figures of equilibrium

\zmath{0213.52304}

\publaddr New Haven, London

\publ Yale University Press 

\totalpages ix+252 

\yr 1969



\Bibitem{ulg:bib:Hagihara}

\by Hagihara~Y.

\book Theories of equilibrium figures of a rotating homogeneous fluid mass

\yr 1970

\serial NASA Special Publication 

\vol 186 

\publaddr Washington

\publ US Government Printing Office

\adsnasa{1970NASSP.186.....H}



\Bibitem{ulg:bib:Borisov}

\by Borisov~A.~V., Kilin~A.~A., Mamaev~I.~S.

\paper The Hamiltonian Dynamics of Self-gravitating Liquid and Gas Ellipsoid

\jour Regul. Chaotic Dyn. 

\yr 2009

\vol 14

\issue 2

\pages 179--217

\zmath{https://zbmath.org/?q=an:1229.70049}

\crossref{10.1134/S1560354709020014}



\Bibitem{ulg:bib:Appell}

\by Appell~P.

\book Trait\'e de m\'ecanique rationnelle

\bookinfo Tome IV, 1: Figures d'\'equilibre d'une masse liquide homog\`ene en rotation

\lang In French

\zmath{58.1300.02}

\totalpages viii+342 

\publaddr  Paris

\publ Gauthier-Villars 

\yr 1932







\Bibitem{ulg:bib:Betti}

\by Betti~E.

\paper Sopra i moti che conservano la figura ellissoidale a una massa fluida eterogenea

\lang In Italian

\zmath{13.0718.01}

\jour Brioschi Ann

\issueinfo  (2)~X

\pages 173--187

\yr 1879







\Bibitem{ulg:bib:Volterra}

\by Volterra~V.

\paper Sur la stratification d\'une masse fluide en \'equilibre

\jour Acta Math.

\yr 1903

\vol 27

\issue 1

\pages 105--124



\Bibitem{ulg:bib:Liouville}

\by Liouville~J.

\paper Sur la figure d\'une masse fluide homog\'ene, en \'equilibre et dou\'ee d\'un mouvement de rotation

\jour J. de l'\'Ecole Polytech.

\yr 1834

\vol 14

\pages 289–296



\Bibitem{ulg:bib:Dirichlet}

\by Dirichlet~G.~L.

\paper Untersuchungen \"uber ein Problem der Hydrodynamik (Aus dessen Nachlass hergestellt von Herrn R. Dedekind zu Z\"urich)

\jour J. Reine Angew. Math. (Crelle's Journal)

\yr 1861

\vol 58

\pages181--216



\Bibitem{ulg:bib:Arnold}

\by Arnold~V. I., Khesin~B.~A.

\book Topological Methods in Hydrodynamics

\yr 1999

\publ Springer

\publaddr New York

\totalpages 392



\Bibitem{ulg:bib:Dolzhansky}

\by Dolzhansky~F.~V.

\paper On the mechanical prototypes of fundamental hydrodynamic invariants and slow manifolds

\jour Physics–Uspekhi

\yr 2005

\vol 48

\issue 12

\pages 1205--1234

\crossref{10.1070/PU2005v048n12ABEH002375}



\Bibitem{ulg:bib:Hadamard}

\by Hadamard~J.

\paper Mouvement permanent lent d'une sph\`ere liquide et visqueuse dans un liquide visqueux

\zmath{42.0813.01}

\jour C.~R.~Acad. Sci., Paris 

\vol 152

\issue 25

\pages 1735--1738 

\yr 1911





\Bibitem{ulg:bib:Rybczynski}

\by Rybczi\'nski, W.

\paper \"Uber die fortschreitende Bewegung einer fl\"ussigen Kugel in einem z\"ahen Medium

\zmath{42.0811.02}

\jour Bull. Int. Acad. Sci. Cracovie. Ser. A

\vol 1

\pages 40--46 

\yr 1911



\Bibitem{ulg:bib:Aristov1}

\by Aristov~S.~N., Prosviryakov~E.~Yu. 

\paper A New Class of Exact Solutions for Three Dimensional Thermal Diffusion Equations

\jour Theor. Found. Chem. Eng.

\yr 2016

\vol 50

\issue 3

\pages 286--293

\crossref{10.1134/S0040579516030027}



\Bibitem{ulg:bib:Aristov2}

\by Aristov~S.~N., Knyazev~D.~V., Polyanin~A.~D.

\paper Exact solutions of the Navier-Stokes equations with the linear dependence of velocity components on two space variables

\jour Theor. Found. Chem. Eng.

\crossref{10.1134/S0040579509050066}

\yr 2009

\vol 43

\issue 5

\pages 642--662



\Bibitem{ulg:bib:Polyanin}

\by Polyanin~A.~D., Zaitsev~V.~F.

\book Handbook of ordinary differential equations. Exact solutions, methods, and problems

\zmath{06776009}

\publaddr Boca Raton, FL

\publ CRC Press 

\totalpages xxix+1456 

\yr 2018



\Bibitem{ulg:bib:Aristov3}

\by Aristov~S.~N., Prosviryakov~E.~Yu.

\paper Unsteady Layered Vortical Fluid Flows

\jour Fluid Dyn.

\crossref{10.1134/S0015462816020034}

\yr 2016

\vol 51

\issue 2

\pages 148--154



\Bibitem{ulg:bib:Aristov4}

\by Aristov~S.~N., Prosviryakov~E.~Yu., Spevak~L.~F.

\paper Unsteady-State B\'enard--Marangoni Convection  in Layered Viscous Incompressible Flows

\jour Theor. Found. Chem. Eng.

\crossref{10.1134/S0040579516020019}

\yr 2016

\vol 50

\issue 2

\pages 132--141



\Bibitem{ulg:bib:Aristov5}

\by Aristov~S.~N., Prosviryakov~E.~Yu.

\paper On one class of analytic solutions of the stationary axisymmetric convection B\'enard--Marangoni viscous incompressible fluid

\jour Vestn. Samar. Gos. Tekhn. Univ., Ser. Fiz.-Mat. Nauki \rm [J. Samara State Tech. Univ., Ser. Phys. Math. Sci.]

\yr 2013

\vol 3(32)

\pages 110--118

\mathnet{http://mi.mathnet.ru/vsgtu1205}

\crossref{10.14498/vsgtu1205}

\zmath{https://zbmath.org/?q=an:06968790}

\elib{http://elibrary.ru/item.asp?id=20710187}

\lang In Russian



\Bibitem{ulg:bib:Aristov6}

\by Aristov~S.~N., Privalova~V.~V., Prosviryakov~E.~Yu.

\paper Stationary nonisothermal Couette flow. Quadratic heating of the upper boundary of the fluid layer

\jour Nelin. Dinam.

\yr 2016

\vol 12

\issue 2

\pages 167--178

\mathnet{http://mi.mathnet.ru/nd519}

\elib{http://elibrary.ru/item.asp?id=26193555}

\crossref{10.20537/nd1602001}

\lang In Russian



\Bibitem{ulg:bib:Burmasheva1}

\by Burmasheva~N.~V., Prosviryakov~E.~Yu.

\paper A large-scale layered stationary convection of an incompressible viscous fluid under the action of shear stresses at the upper boundary. Velocity field investigation

\jour Vestn. Samar. Gos. Tekhn. Univ., Ser. Fiz.-Mat. Nauki \rm [J. Samara State Tech. Univ., Ser. Phys. Math. Sci.]

\yr 2017

\vol 21

\issue 1

\pages 180--196

\mathnet{http://mi.mathnet.ru/vsgtu1527}

\crossref{10.14498/vsgtu1527}

\elib{http://elibrary.ru/item.asp?id=29245104}

\lang In Russian



\Bibitem{ulg:bib:Burmasheva2}

\by Burmasheva~N.~V., Prosviryakov~E.~Yu.

\paper A large-scale layered stationary convection of a incompressible viscous fluid under the action of shear stresses at the upper boundary. Temperature and presure field investigation

\jour Vestn. Samar. Gos. Tekhn. Univ., Ser. Fiz.-Mat. Nauki \rm [J. Samara State Tech. Univ., Ser. Phys. Math. Sci.]

\yr 2017

\vol 21

\issue 4

\pages 736--751

\mathnet{http://mi.mathnet.ru/vsgtu1568}

\crossref{10.14498/vsgtu1568}

\zmath{https://zbmath.org/?q=an:06964885}

\elib{http://elibrary.ru/item.asp?id=32712835}

\lang In Russian



\Bibitem{ulg:bib:Gorshkov}

\by Gorshkov~A.~V., Prosviryakov~E.~Yu.

\paper Ekman Convective Layer Flow of a Viscous Incompressible Fluid

\jour Izv. Atmos. Ocean. Phys.

\yr 2018

\vol 54

\issue 2

\pages 89–195

\crossref{10.1134/S0001433818020081}





\Bibitem{ulg:bib:Aristov7}

\by Aristov~S.~N., Knyazev~D.~V.

\paper Localized convective flows in a nonuniformly heated liquid layer

\jour Fluid Dyn.

\yr 2014

\vol 49

\issue 5

\pages 565--575

\crossref{10.1134/S0015462814050020}











\end{thebibliography}





\renewcommand{\baselinestretch}{0.95}



\newpage



\selectlanguage{russian}



$\;$







\makerutitle



\vspace{-7mm}



\AdditionalContent{Конкурирующие интересы}{Конкурирующих интересов не имею.}



\AdditionalContent{Авторский вклад и ответственность}{Я несу полную ответственность за предоставление окончательной версии рукописи в печать. Окончательная версия 

рукописи мною одобрена.}

\AdditionalContent{Финансирование}{Работа выполнена в рамках государственного задания ФАНО, тема № АААА-А18-118020790140-5.}





\renewcommand{\baselinestretch}{0.9}







\renewcommand{\RuCitation}{\noindent\textbf{\CYRO\cyrb\cyrr\cyra\cyrz\cyre\cyrc\ \cyrd\cyrl\cyrya\ \cyrc\cyri\cyrt\cyri\cyrr\cyro\cyrv\cyra\cyrn\cyri\cyrya\\} {\RuAuthorInContents} {\RuTitleInContents}~/\!/ \textit{\RuJournalName}, \JournalYear.  \CYRT.~\JournalVolume, \textnumero~\JournalNumber.  \CYRS.~\thefirstpage--\thelastpage. \doiurl.}



\renewcommand{\EnCitation}{\noindent\textbf{Please cite this article in press as:\\} {\EnAuthorInContents} {\EnTitleInContents}, \textit{\EnJournalName}\/ [\ParallelJournalName],  \JournalYear, vol.~\JournalVolume, no.~\JournalNumber,  pp.~\thefirstpage--\thelastpage.  \midoiurl\ (In Russian).}









%\end{document}

 \newpage

%\input{preamble}

%

%\usepackage{samgtu-name}

%

%\usepackage{slashbox}



\vsgtu{1653} %registration number

\FirstPage{750}

\LastPage{761}

%

%\pdfcrop

%\draft



\ShortCommunication



%\begin{document}





\hfill \qrcode[hyperlink,height=.8in]{http://doi.org/10.14498/vsgtu1653}

\vspace{-.8in}





\udc{517.958:539.3(1)}

\msc{74F05, 74K20}



\rutitle{Динамическая устойчивость геометрически нерегулярной нагретой пологой цилиндрической оболочки в~сверхзвуковом потоке газа}

\rutitleshort{Динамическая устойчивость геометрически нерегулярной\dots}



\entitle{Dynamic stability of heated geometrically irregular cylindrical shell in supersonic gas flow}

\entitleshort{Dynamic stability of heated geometrically irregular cylindrical shell in supersonic gas flow}





\ruauthor{Г.~Н.~Белосточный, О.~А.~Мыльцина}{Б\,е\,л\,о\,с\,т\,о\,ч\,н\,ы\,й~Г.~Н.,\ М\,ы\,л\,ь\,ц\,и\,н\,а~О.~А.}

\enauthor{G.~N.~Belostochnyi, O.~A.~Myltcina}{B\,e\,l\,o\,s\,t\,o\,c\,h\,n\,y\,i~G.~N.,\ M\,y\,l\,t\,c\,i\,n\,a~O.~A.}



\ruaffil{\saratovuni.}



\enaffil{\engsaratovuni.}



\ruauthorsdetails{3}{% Количество авторов



\fullauthor{\rupid{98285}{Григорий Николаевич Белосточный}}

%\CAuthor % Автор-корреспондент

\orcid{0000-0003-4471-6599} % ORCID ID автора 

\regalia{доктор технических наук, профессор} 

\position{профессор} 

\dept{каф. математической теории упругости и биомеханики} 

\email{belostochny@mail.ru}

\fullauthor{\rupid{98283}{Ольга Анатольевна Мыльцина}}

\CAuthor % Автор-корреспондент

\orcid{0000-0003-4718-2772} % ORCID ID автора 

\regalia{кандидат физико-математических наук} 

\position{ассистент} 

\dept{каф. теории функций и стохастического анализа} 

\email{omyltcina@yandex.ru}

}



%Olga Myltcina https://orcid.org/0000-0003-4718-2772 

%Grigory Belostochnyihttps://orcid.org/0000-0001-6949-4174

%Acel Polienko https://orcid.org/0000-0003-4471-6599







\enauthorsdetails{3}{% Количество авторов

\fullauthor{\rupid{98285}{Grigory N. Belostochnyi}}

%\CAuthor % Автор-корреспондент

\orcid{0000-0003-4471-6599} % ORCID ID автора 

\regalia{Dr. Techn. Sci.} 

\position{Professor} 

\dept{Dept. of Mathematic Theory of Elasticity \&  Biomechanics} 

\email[n]{belostochny@mail.ru} 

\fullauthor{\rupid{98283}{Olga A. Myltcina}}

\CAuthor % Автор-корреспондент

\orcid{0000-0003-4718-2772} % ORCID ID автора 

\regalia{Cand. Phys. \& Math. Sci.} 

\position{Assistant} 

\dept{Dept. of Functions \& Approxmation Theory} 

\email[n]{omyltcina@yandex.ru}

}





\rukeywords{динамическая устойчивость, температура, пологие оболочки, сверхзвук, континуальность, обобщенные функции, поршневая теория, аэродинамика, критические скорости, ребра жесткости, демпфирование, кривизна, критерий Гурвица, изотропия}

\enkeywords{dynamic stability, temperature, flat shells, supersonic, continuity, generalized functions, piston theory, aerodynamics, critical velocities, ribs, damping, curvature, Routh–Hurwitz stability criterion, isotropy}



\ruabstract{На базе модели типа Лява рассматривается нагретая до постоянной температуры геометрически нерегулярная пологая цилиндрическая оболочка, обдуваемая сверхзвуковым потоком газа со стороны одной из ее основных поверхностей. За основу взята континуальная модель термоупругой системы в~виде тонкостенной оболочки, подкрепленной ребрами вдоль набегающего газового потока. Сингулярная система уравнений динамической термоустойчивости геометрически нерегулярной оболочки содержит слагаемые, учитывающие «растяжение-сжатие» и~сдвиг подкрепляющих элементов в~тангенциальной плоскости, тангенциальные усилия, вызванные нагревом оболочки, и~поперечную нагрузку, стандартным образом записанную по «поршневой теории».



Тангенциальные усилия предварительно определяются как решение сингулярных дифференциальных уравнений безмоментной термоупругости геометрически нерегулярной оболочки с~учетом краевых усилий.



Решение системы динамических уравнений термоупругости оболочки разыскивается в~виде суммы двойного тригонометрического ряда (для функции прогиба) с~переменными по временной координате коэффициентами. На основании метода Галеркина получена однородная система для коэффициентов аппроксимирующего ряда, которая сведена к~одному дифференциальному уравнению четвертого порядка. Решение приводится во втором приближении, что соответствует двум полуволнам в~направлении потока и~одной полуволне в~перпендикулярном направлении. На основании стандартных методов анализа динамической устойчивости тонкостенных конструкций определяются критические значения скоростей газового потока.



Количественные результаты приводятся в виде таблиц, иллюстрирующих влияние геометрических параметров термоупругой системы «обо\-лоч\-ка-ребра», температуры  на устойчивость геометрически нерегулярной цилиндрической оболочки в~сверхзвуковом потоке газа с~учетом демпфирования.}



\enabstract{On the basis of the Love model, a geometrically irregular heated cylindrical shell blown by a supersonic gas flow from one of its main surfaces is considered. The continuum model of a thermoelastic system in the form of a thin-walled shell supported by ribs along the incoming gas flow is taken as a basis. The singular system of equations for the dynamic thermal stability of a geometrically irregular shell contains terms that take into account the tension-compression and the shift of the reinforcing elements in the tangential plane, the tangential forces caused by the heating of the shell and the transverse load, as standard recorded by the piston theory. The solution of a singular system of differential equations in displacements, in the second approximation for the deflection function, is sought in the form of a double trigonometric series with time coordinate variables. 





Tangential forces are predefined as the solution of singular differential equations of non-moment thermoelasticity of a geometrically irregular shell taking into account boundary forces. 



The solution of the system of dynamic equations of thermoelasticity of the shell is sought in the form of the sum of the double trigonometric series (for the deflection function) with time coordinate variable coefficients. On the basis of the Galerkin method, a homogeneous system for the coefficients of the approximating series is obtained, which is reduced to one fourth-order differential equation. The solution is given in the second approximation, which corresponds to two half-waves in the direction of flow and one half-wave in the perpendicular direction. On the basis of standard methods of analysis of dynamic stability of thin-walled structures are determined critical values of the gas flow rate. 



The quantitative results are presented in the form of tables illustrating the influence of the geometrical parameters of the thermoelastic shell-edge system, temperature and damping on the stability of a geometrically irregular cylindrical shell in a supersonic gas flow.}



\newdate{myla}{11}{10}{2018}

\date{\displaydate{myla}}

\newdate{mylb}{10}{11}{2018}

\revisiondate{\displaydate{mylb}}

\newdate{mylc}{12}{11}{2018}

\accepteddate{\displaydate{mylc}}



\newdate{myle}{21}{12}{2018}

\renewcommand{\JournalDataOnline}{\displaydate{myle}}



\makerutitle







Геометрически нерегулярные тонкостенные упругие системы, обширный класс которых составляют ребристые оболочки и~пластинки, широко используются в~различных областях современной техники. Условия эксплуатации таких систем предусматривают совместное воздействие температурных полей и~высокоскоростных газовых потоков. Исследованию упругого поведения гладких пластин и~оболочек на основе атермической

теории посвящено большое число работ, полный перечень которых содержал бы десятки

наименований. Ограничимся некоторыми из них \cite{myl-bib:1, myl-bib:2, myl-bib:3,myl-bib:4,myl-bib:5,myl-bib:5a}.

Значительно меньше работ

посвящено  исследованиям совместного воздействия температурных факторов и сверхзвукового потока

на тонкостенные системы. Важные для практики результаты в этой области содержатся в работах \cite{myl-bib:6, myl-bib:7, myl-bib:8}.



%Работы, в которых анализируется влияние подкрепляющих  оболочку ребер под действием сверхзвукового потока на базе термической теории в открытой научной литературе отсутствуют. 

Работы, в которых анализируется влияние подкрепляющих  ребер на поведение оболочки, находящейся под действием сверхзвукового потока, на базе термической теории, в~открытой научной литературе отсутствуют. 

Это связано не с маловажностью проблемы, а прежде всего с чрезвычайной математической сложностью таких задач, решаемых на основе дискретной модели «оболочка–ребра». Использование континуальной модели типа «конструктивная анизотропия» мало пригодно для практических целей, так как количественные результаты, полученные на ее основе, далеки от физической реальности.



Обращение к  континуальным моделям на основе теории обобщенных функций, основные положения которых содержится в работах \cite{myl-bib:9, myl-bib:10, myl-bib:11, myl-bib:12, myl-bib:15,myl-bib:16,myl-bib:18}, позволило сводить решения задач статической и динамической термоупругости ребристых  оболочек к интегрированию систем сингулярных дифференциальных уравнений точными и~приближенными методами высшего анализа \cite{myl-bib:19, myl-bib:21}, что немаловажно для инженерной практики.



\smallskip



{\bf 1.} Система сингулярных дифференциальных уравнений динамической термоупругости пологой геометрически нерегулярной цилиндрической оболочки, стандартным образом отнесенной к декартовым координатам на основе континуальной модели, в которых учет тангенциальных усилий записывается в форме Рейснера \cite{myl-bib:22a}, в компонентах поля перемещений примет вид

\begin{gather}

u_{\bigcom 11}+\frac{1-\nu}{2}u_{\bigcom 22}+\frac{1+\nu}{2} v_{\bigcom 12}-k_{11}w_{\bigcom 1}+

\varepsilon_1\frac{1-\nu}{2} \sum_{i=1}^n \frac{h_i}{h}(u_{\bigcom 2}+v_{\bigcom 1} )_{\bigcom 2}\delta(x-x_i)=0,

\\

\frac{1+\nu}{2}u_{\bigcom 21}+ v_{\bigcom 22}+\frac{1-\nu}{2}v_{\bigcom 11}-\nu k_{11}w_{\bigcom 2}+\varepsilon_2\sum_{i=1}^n\frac{h_i}{h}a_i ( \nu u_{\bigcom 1}+v_{\bigcom 2})_{\bigcom 2} \delta(x-x_i)=0,

\end{gather}



\vspace{-12mm}



\begin{multline}\label{bel_eq1}

\nabla^2\nabla^2 w+\frac{B}{D}k^2_{11}w-k_{11}\frac{B}{D} ( u_{\bigcom 1}+\nu v_{\bigcom 2})- 

\\

- \frac{1}{D} \bigl[ 

( T^{11}_0 w_{\bigcom 1})_{\bigcom 1}+

( T^{12}_0 w_{\bigcom 1})_{\bigcom 2}+  

(T^{12}_0 w_{\bigcom 2})_{\bigcom 1}+ 

(T^{22}_0 w_{\bigcom 2})_{\bigcom 2}\bigr]-

\\

-

\frac{1}{D}\sum_{i=1}^n\frac{h_i}{h}a_i

(T^{22}_0 w_{\bigcom 2})\delta(x-x_i)+ 

\\

+\sum_{i=1}^n \Bigl( \frac{h_i}{h}\Bigr)^3 \Phi_{3i}a_iw_{\bigcom 2222}\delta(x-x_i)+

\\

+

2(1-\nu)\sum_{i=1}^n\Bigl( \frac{h_i}{h}\Bigr)^3 \Phi_{3i}a_iw_{\bigcom 122}\bigr|_{x_i}\frac{d\delta(x-x_i)}{dx}= \\

=\frac{\gamma h}{g D}w_{\bigcom tt}\biggl[ 1+\sum_{i=1}^n \frac{h_i}{h}a_i\delta(x-x_i)\biggr] 

-\frac{\mu} {h D}w_{\bigcom t}-p_0\frac{\varkappa \mu}{D}w_{\bigcom 2}, 

\end{multline}

где $p_0\frac{\varkappa \mu}{D}w_{\bigcom 2}$ "---  относительная интенсивность поперечной нагрузки, вызванная прогибом пластинки, стандартным образом записана согласно поршневой теории \cite{myl-bib:1,myl-bib:22, myl-bib:30}; $M={v_y}/{c_0}$ "--- число Маха; 

$v_y$ "--- невозмущенная скорость набегающего газового потока; 

$c_0$ "--- скорость звука на бесконечности; $u$, $v$, $w$ "---  компоненты поля (векторы) перемещений; 

${h_i}/{h}$ "--- относительная высота $i$-того ребра, число которых $n$; 

$a_i$ "--- ширина $i$-того ребра; 

$k_{11}$ "--- кривизна цилиндрической оболочки; 

$D=\frac{E h^3}{12(1-\nu)}$; $\nu$ "--- коэффициент Пуассона; 

$E$ "--- модуль Юнга, 

$\gamma$ "--- удельный вес; 

$g$ "--- интенсивность поля тяжести; 

$$\Phi_{3i}=1+3\frac{h_i}{h}+\Bigl( \frac{h_i}{h}\Bigr)^3 ;$$ 

$\mu$ "--- коэффициент демпфирования; 

$\delta(x-x_i)$ "--- обобщенная дельта"=функция Дирака~\cite{myl-bib:25};  

$\varepsilon_j$ $(j=0,1)$ "--- знаковые числа, равные 1 или 0; 

$ T^{ij}_0 $ "--- тангенциальные усилия, вызванные нагревом пластинки до постоянной температуры $\Theta_0$, которые определяются на основании решений системы сингулярных однородных дифференциальных уравнений, описывающих безмоментное состояние термоупругой системы в компонентах поля перемещений $\bar{u}^0(u^0,v^0,w^0)$:

 \begin{multline}

 u^0_{\bigcom 11}+\frac{1-\nu}{2}u^0_{\bigcom 22}+\frac{1+\nu}{2}v^0_{\bigcom 12}-k_{11}w^0_{\bigcom 1}+

\\

 +\frac{1-\nu}{2}\sum_{i=1}^{n}\frac{h_i}{h}a_i ( u^0_{\bigcom 2}+v^0_{\bigcom 1})_{\bigcom 2}\delta(x-x_i) =0,

 \end{multline}

 \begin{multline}

 \frac{1+\nu}{2}u^0_{\bigcom 12}+v^0,_{22}+\frac{1-\nu}{2}v^0_{\bigcom 11}-\nu k_{11}w^0_{\bigcom 2}+

\\

+ 

 \sum_{i=1}^{n}\frac{h_i}{h}a_i( v^0_{\bigcom 2}+\nu u^0_{\bigcom 1})_{\bigcom 2}\delta(x-x_i)=0,

 \end{multline}

\begin{equation}\label{eq:bel:momentnot}

k_{11}^2 w^0+k_{11} (u^0_{\bigcom 1}+\nu v^0_{\bigcom 2} ) =\alpha(1+\nu)k_{11}\Theta_0.

\end{equation}





При неоднородных краевых условиях

\begin{equation*}

\begin{array}{ll}

\text{$x=0$, $x=a$:}&

u^0_{\bigcom 2}+v^0_{\bigcom 1}=0, \quad u^0_{\bigcom 1}+\nu v^0_{\bigcom 2}-k_{11}w^0=\alpha(1+\nu \Theta_0), \quad w^0=0;

\\[2mm]

\text{$y=0$, $y=b$:}&

v^0=0, \quad u^0_{\bigcom 2}+ v^0_{\bigcom 1}=0, \quad w^0=0

\end{array}

\end{equation*}

решения сингулярной системы \eqref{eq:bel:momentnot} запишутся в элементарных функциях:

\begin{equation}\label{bel_gr1}

v^0=0, \quad u^0=\alpha(1+\nu)\Theta_0 x, \quad w^0=0,

\end{equation}

а тангенциальные усилия примут вид

\begin{equation*}

T^{11}_0=T^{12}_0=0, \quad T^{22}_0=-E h\alpha\Theta_0.

\end{equation*}



Решение системы \eqref{bel_eq1} (в~которой отсутствуют инерционные слагаемые в~тангенциальной плоскости \cite{myl-bib:1,myl-bib:6, myl-bib:7}), тождественно удовлетворяющее условиям шарнирного закрепления всех сторон пологой цилиндрической оболочки

\begin{equation*}

\begin{array}{ll}

\text{при $x=0$, $x=a$:} &

T^{12}=0, \quad T^{11}=0, \quad w=0 \quad M^{11}=0;

\\[2mm]

\text{при $y=0$, $y=b$:} &

v=0, \quad T^{12}=0, \quad w=0 \quad M^{22}=0,

\end{array}

\end{equation*}

зададим в виде

\begin{gather}\label{bel_appr1}

u(x,y,t)=\tilde{u}(t)u_1 ({x}/{a} )u_2 ({y}/{b} ), \quad 

v(x,y,t)=\tilde{v}(t)v_1({x}/{a} )v_2 ({y}/{b} ),

\\

w(x,y,t)=w_{11}(t)\sin \frac{\pi x}{a} \sin\frac{\pi y}{b} +w_{12}(t)\sin \frac{\pi x}{a} \sin\frac{2\pi y}{b},

\end{gather}

где 

\begin{gather}

u_1({x}/{a}) =( {x}/{a})^4-2( {x}/{a})^3+( {x}/{a})^2,\quad

u_2({y}/{b}) =-({2}/{3})( {y}/{b})^3+( {y/}{b})^2,\\ 

v_1({x}/{a}) =({x}/{a})^4-2( {x}/{a})^3+( {x}/{a})^2,

\quad 

v_2({y}/{b}) =({y}/{b})( {y}/{b}-1).

\end{gather}



Первые два уравнения системы \eqref{bel_eq1} на основании подстановок \eqref{bel_appr1} с последующим применением процедуры Галеркина \cite{myl-bib:27,  myl-bib:29}  перепишутся в виде 

\begin{gather}

e_{11}\tilde{u}(t)+e_{12}\tilde{v}(t)=k_{11}a\Bigl[ \pi \frac{b}{a}\mathit{I}_4 w_{11}(t)+\pi\frac{b}{a}\mathit{I}_5 w_{12}(t) \Bigr],

\\

\label{bel_eq7}

e_{21}\tilde{u}(t)+e_{22}\tilde{v}(t)=\nu k_{11} a \bigl[ \pi \mathit{I}_9 w_{11}(t)+2 \pi \mathit{I}_{10} w_{12}(t)\bigr].

\end{gather}

Здесь введены следующие обозначения:

\begin{gather}

e_{11}=\frac{b}{a}\mathit{I}_1+\frac{1-\nu}{2}\frac{a}{b}\mathit{I}_2, \quad 

e_{12}=\frac{1+\nu}{2}\mathit{I}_3,\\

e_{21}=\frac{1+\nu}{2}\mathit{I}_6,\quad 

e_{22}=\frac{a}{b}\mathit{I}_7+\frac{1-\nu}{2}\frac{b}{a}\mathit{I}_8;

\\

\mathit{I}_1=\int_0^1 \int_0^1 \tilde{u}_{1 \bigcom 11} \tilde{u}_1 \tilde{u}_2^2dX dY,\quad 

\mathit{I}_2=\int_0^1 \int_0^1 \tilde{u}_{2  \bigcom 22}\tilde{u}_2\tilde{u}_1^2dX dY,

\\

\mathit{I}_3=\int_0^1 \int_0^1 \tilde{v}_{1\bigcom 1}\tilde{u}_1\tilde{u}_2 \tilde{v}_2,_{2}dX dY, \dots,

\\

\mathit{I}_{10}=\int_0^1 \int_0^1 \tilde{v}_1 \tilde{v}_2 \sin (\pi X )\cos (2\pi Y )dX dY,\quad  X={x}/{a}, \quad  Y={y}/{b}.

\end{gather}



Разрешая неоднородную алгебраическую систему относительно коэффициентов $\tilde{u}(t)$ и $\tilde{v}(t)$, получим

\begin{equation}\label{bel_eq9}

\tilde{u}=k_{11}a\bigl(L_1 w_{11}(t)+L_2 w_{12}(t) \bigr),\quad 

\tilde{v}=k_{11}a\bigl(L_3 w_{11}(t)+L_4 w_{12}(t) \bigr),

\end{equation}

где $L_j$ "--- безразмерные величины:

\begin{gather}

L_1=\frac{\pi \frac{b}{a}\mathit{I}_4e_{22}-\nu \pi \mathit{I}_9 e_{12} }{e_{11} e_{22}-e_{12}e_{21}},\quad L_2=\frac{\pi \frac{b}{a}\mathit{I}_5e_{22}-\nu 2 \pi \mathit{I}_{10} e_{12} }{e_{11} e_{22}-e_{12}e_{21}}, 

\\

L_3=\frac{-\pi \frac{b}{a}\mathit{I}_4e_{21}+\nu \pi \mathit{I}_9 e_{11} }{e_{11} e_{22}-e_{12}e_{21}},\quad L_4=\frac{-\pi \frac{b}{a}\mathit{I}_5e_{21}+\nu 2 \pi \mathit{I}_{10} e_{11} }{e_{11} e_{22}-e_{12}e_{21}}. 

\end{gather}

Окончательно  функции $u(x,y,t)$ и $v(x,y,t)$ примут следующий вид:

\begin{gather}\label{bel_eq11}

u(x,y,t)=k_{11} a \bigl(L_1 w_{11}(t)+L_2 w_{12}(t) \bigr) u_1({x}/{a} )u_2({y}/{b}),

\\

v(x,y,t)=k_{11} a \bigl(L_3 w_{11}(t)+L_4 w_{12}(t) \bigr) v_1({x}/{a} )v_2({y}/{b} ).

\end{gather}



Из третьего уравнения системы \eqref{bel_eq1}  на основании аналогичных преобразований получим систему обыкновенных однородных дифференциальных уравнений относительно коэффициентов $w_{11}(t)$ и $w_{12}(t)$:

\begin{equation}\label{bel_eq12}

\xi_{11}\left[w_{11} \right]+\xi_{12} w_{12}=0, \quad  \xi_{21}w_{11}+\xi_{22} \left[ w_{12}\right] =0,

\end{equation}

где $\xi_{kl}$ "--- дифференциальные операторы:

\begin{gather}

\xi_{11}=\frac{\gamma h a^4}{g D}\biggl( 1+2 \sum\limits_{i=1}^n \tilde{\beta}_i^s\biggr)

\frac{d^2}{dt^2}+\tilde{\mu} \frac{d}{dt}+ f_{11}\Bigl( k_{11}, \Theta_0, \frac{h_i}{h}\Bigr), 

\\

\xi_{12}=\frac{p_0 \varkappa M a^3}{D} 4\pi \frac{a}{b} \mathit{I}_{13}-\tilde{f}_{12}, \quad \xi_{21}=\frac{p_0 \varkappa M a^3}{D} 2\pi \frac{a}{b} \mathit{I}_{16}-\tilde{f}_{11},

\\

\xi_{22}=\frac{\gamma h a^4}{g D}

\biggl( 1+2 \sum\limits_{i=1}^n \tilde{\beta}_i^s\biggr)\frac{d^2}{dt^2}+\tilde{\mu} \frac{d}{dt}+ f_{12}

\Bigl( k_{11},\Theta_0,\frac{h_i}{h}\Bigr), 

\end{gather}

в которых

\begin{multline}

f_{11}=\Bigl(\pi^2+\Bigl( \frac{\pi a}{b}\Bigr) ^2 \Bigr) ^2+

12

\Bigl( \frac{h_i}{h}\Bigr)^2(k_{11}a)^2-

48

\Bigl( \frac{a}{h}\Bigr)^2 (k_{11}a)^2\Bigl(L_1\mathit{I}_{11}+\nu \frac{a}{b}L_3 \mathit{I}_{12} \Bigr) -

\\

-12(1-\nu^2)\Bigl( \frac{\pi a}{b}\Bigr)^2 \Bigl(\frac{a}{h} \Bigr)^2

\biggl(1+\sum\limits_{i=1}^n \tilde{\beta}_i^s \biggr) \alpha \Theta_0+

\Bigl( \frac{\pi a}{b}\Bigr)^4 2\sum\limits_{i=1}^n \beta_i^s+

\\

+4(1-\nu)\pi^2\Bigl(\frac{\pi a}{b} \Bigr)^2  \sum\limits_{i=1}^n \beta_i^c,

\end{multline}

\begin{equation*}

\tilde{f}_{12}=48 \Bigl(\frac{a}{h} \Bigr)^2(k_{11}a)^2\Bigl(L_2 \mathit{I}_{11}+\nu\frac{a}{b}L_4\mathit{I}_{12}\Bigr),

\end{equation*}

\begin{equation*}

\tilde{f}_{11}=48 \Bigl(\frac{a}{h} \Bigr)^2(k_{11}a)^2\Bigl(L_1 \mathit{I}_{14}+\nu \frac{a}{b}L_3\mathit{I}_{15}\Bigr),

\end{equation*}

\begin{multline}

f_{12}=\Bigl(\pi^2+\Bigl( \frac{2\pi a}{b}\Bigr) ^2 \Bigr)^2-

12\Bigl( \frac{a}{h}\Bigr)^2(k_{11}a)^2-

48\Bigl( \frac{a}{h}\Bigr)^2 (k_{11}a)^2 \Bigl(L_2\mathit{I}_{14}+\nu \frac{a}{b}L_4 \mathit{I}_{15} \Bigr) -

\\

-12(1-\nu^2)\Bigl( \frac{2\pi a}{b}\Bigr)^2

\Bigl(\frac{a}{h} \Bigr)^2

\biggl(1+\sum\limits_{i=1}^n \tilde{\beta}_i^s \biggr) \alpha \Theta_0+\Bigl( \frac{2\pi a}{b}\Bigr)^4 2\sum\limits_{i=1}^n \beta_i^s+

\\

+4(1-\nu)\pi^2\Bigl(\frac{2\pi a}{b} \Bigr)^2  \sum\limits_{i=1}^n \beta_i^c, 

\end{multline}

\begin{equation*}

\beta_i^c=\Bigl( \frac{h_i}{h}\Bigr) ^3\frac{a_i}{a}\varPhi_{3i}\cos^2 \frac{\pi x_i}{a},

\quad 

\beta_i^s=\Bigl( \frac{h_i}{h}\Bigr) ^3\frac{a_i}{a}\varPhi_{3i}\sin^2 \frac{\pi x_i}{a},

\quad 

\tilde{\beta}_i^s= \frac{h_i}{h}\frac{a_i}{a}\sin^2 \frac{\pi x_i}{a}.

\end{equation*}



Система дифференциальных уравнений \eqref{bel_eq12} подстановкой 

\begin{equation*}

w_{11}=-\xi_{12}\varPhi(t), \quad w_{12}=\xi_{11}\left[\varPhi(t) \right] 

\end{equation*}

сводится к одному дифференциальному уравнению четвертого порядка, которое в случае отсутствия  демпфирования ($\mu=0$) содержит только четные производные от функции $\varPhi(t)$:

\begin{multline}\label{bel_eq15}

\biggl[ 

\frac{\gamma h a^4}{g D}\biggl( 1+2\sum\limits_{i=1}^n \tilde{\beta}_i^s\biggr) \biggr]^2

\frac{d^4\varPhi}{dt^4} +

\biggl[ \frac{\gamma h a^4}{g D}\biggl( 1+2\sum\limits_{i=1}^n \tilde{\beta}_i^s\biggr)

( f_{11}+f_{12})  \biggr]\frac{d^2\varPhi}{dt^2}+

\\

+\Bigl[ f_{11}f_{12}-8\pi \Bigl( \frac{a}{b}\Bigr)^2\mathit{I}_{13}\mathit{I}_{16}\tilde{M}+\tilde{M} 

\Bigl(4 \pi \frac{a}{b}\mathit{I}_{13}\tilde{f}_{11}+2 \pi \frac{a}{b}\mathit{I}_{16}\tilde{f}_{12} \Bigr) - \tilde{f}_{11}\tilde{f}_{12}\Bigr] \varPhi=0,

\end{multline}

где $\tilde{M}=\frac{p_0 \varkappa M a^3}{D}$.



Если $\mu\neq0$, то дифференциальное уравнение для функции $\varPhi(t)$ запишется в~виде

\begin{equation}\label{bel_eq16}

\frac{d^4\varPhi}{dt^4}+2 e_1\frac{d^3\varPhi}{dt^3}+e_3\frac{d^2\varPhi}{dt^2}+e_4\frac{d\varPhi}{dt}+e_5\varPhi=0,

\end{equation}

где

\begin{equation*}

e_1=\frac{\tilde{\mu} }{ \frac{\gamma h a^4}{g D}\left(1+2\sum_{i=1}^n\tilde{\beta}_i^s \right) }, \quad e_3=\frac{\frac{\gamma h a^4}{g D}\left(1+2\sum_{i=1}^n\tilde{\beta}_i^s \right)\left( f_{11}+f_{12}\right)+ \tilde{\mu}}{\left( \frac{\gamma h a^4}{g D}\left(1+2\sum_{i=1}^n\tilde{\beta}_i^s \right)\right)^2},

\end{equation*}

\begin{equation*}

e_4=\frac{\tilde{\mu}\left(f_{11}+f_{12} \right)  }{\left(  \frac{\gamma h a^4}{g D}\left(1+2\sum_{i=1}^n\tilde{\beta}_i^s \right)\right)^2},\quad e_5=\frac{f_{11} f_{12}-\xi_{21}\xi_{12}}{\left( \frac{\gamma h a^4}{g D}\left(1+2\sum_{i=1}^n\tilde{\beta}_i^s \right)\right)^2}, \quad \tilde{\mu}=\frac{\mu a^4 h}{D}.

\end{equation*}



Из условия, при котором хотя бы один корень из четырех  характеристического уравнения для дифференциального уравнения \eqref{bel_eq15} будет положительным, определяется наименьшее значение  относительной скорости потока газа ${v_y}/{c_0}$, при котором прогиб термоупругой системы растет. При учете демпфирования ($\mu\neq0$) возникает вопрос об устойчивости системы \eqref{bel_eq12}, которая сведена к одному дифференциальному уравнению \eqref{bel_eq16}. 

В~этом случае это значение скорости потока определяется на основании критерия Гурвица~\cite{myl-bib:32}:

\begin{equation*}

\Delta_1=2e_1>0, \quad \Delta_2=2e_1e_3-e_4>0, \quad \Delta_3=2e_1e_3e_4-e_4^2-4e_1^2e_5>0.

\end{equation*}



Кривизна пологой цилиндрической оболочки задавалась в виде $k_{11}\hm=-{4 \tilde{\delta}}/{a^2}$ \cite{myl-bib:33, myl-bib:34}, где $\tilde{\delta}$ "--- наибольший подъем оболочки над ее планом. Выражение для относительной кривизны $\tilde{k}_{11}$ принимает вид

\begin{equation*}

\tilde{k}_{11}=k_{11} a=-4\frac{\tilde{\delta}}{h}\frac{h}{a},

\end{equation*}

где ${\tilde{\delta}}/{h}\in[0, 5]$ \cite{myl-bib:33, myl-bib:34}.



\smallskip



{\bf 2.} Численные результаты приведены в~таблице. 

Вычисления проводились при следующих значениях параметров: 

${h}/{a}=0.01$, 

${a_i}/{a}=0.01$, 

$\mu=0.001$, 

${a}/{b}=1$.





\begin{table}[b!]

	

\small \centering



Значения параметра ${v_y}/{c_0}$  в зависимости от других параметров



[The values of the parameter ${v_y}/{c_0}$ depending on other parameters]

	

\smallskip

	\renewcommand{\tabcolsep}{0.35cm}

\begin{tabular}{c|c|c|c|c|c|c|c}

	\hline 

	\multicolumn{2}{c|} {\rule{0mm}{12pt}}&  \multicolumn{3}{c|} {${h_i}/{h}=2.5$}    &  \multicolumn{3}{c} {${h_i}/{h}=5$}   \\ 

	\hline 

 \backslashbox{$\tilde{k}_{11}$}{$n$} & 0 & 1 & 3 & 5 & 1 & 3 & 5 \\ 

	\hline 

		\multicolumn{8}{c} {\rule{0mm}{12pt}for $\Theta_0=0$,  $\varepsilon_1=\varepsilon_2=0$} \\

	\hline 	

	0.10 & 0.24 & 1.50 & 3.16 & 3.50 & 6.47 & 13.35 & 14.79 \\ 

	0.15 & 0.54 & 1.40 & 3.06 & 3.40 & 6.37 & 13.25 & 14.69 \\

	0.20 & 0.97 & 1.26 & 2.92 & 3.26 & 6.23 & 13.11 & 14.55 \\  

	\hline 

		\multicolumn{8}{c} {\rule{0mm}{12pt}for $\Theta_0=2$,  $\varepsilon_1=\varepsilon_2=0$} \\

	\hline 

	0.10 & 1.08 & 0.61 & 2.22 & 2.54 & 5.55 & 12.34 & 13.73 \\ 

	0.15 & 1.39 & 0.51 & 2.12 & 2.44 & 5.45 & 12.24 & 13.63 \\ 

	0.20 & 1.81 & 0.37 & 1.98 & 2.30 & 5.31 & 12.10 & 13.49 \\ 

	\hline 

		\multicolumn{8}{c} {\rule{0mm}{12pt}for $\Theta_0=4$,  $\varepsilon_1=\varepsilon_2=0$} \\

	\hline 

		0.10 & 1.93 & 0.27 & 1.28 & 1.58 & 4.63 & 11.32 & 12.67 \\ 

		0.15 & 2.23 & 0.37 & 1.18 & 1.48 & 4.53 & 11.22 & 12.57 \\ 

		0.20 & 2.66 & 0.51 & 1.04 & 1.34 & 4.39 & 11.08 & 12.43 \\ 

	\hline 

		\multicolumn{8}{c} {\rule{0mm}{12pt}for $\Theta_0=0$,  $\varepsilon_1=\varepsilon_2=1$} \\

	\hline 		

		0.10 & 0.24 & 1.50 & 3.16 & 3.51 & 6.48 & 13.37 & 14.80 \\ 

		0.15 & 0.54 & 1.40 & 3.07 & 3.42 & 6.39 & 13.28 & 14.72 \\ 

		0.20 & 0.97 & 1.27 & 2.94 & 3.29 & 6.26 & 13.16 & 14.60 \\ 

	\hline 

		\multicolumn{8}{c} {\rule{0mm}{12pt}for $\Theta_0=2$,  $\varepsilon_1=\varepsilon_2=1$} \\

	\hline 	

		0.10 & 1.08 & 0.61 & 2.23 & 2.55 & 5.56 & 12.35 & 13.74 \\ 

		0.15 & 1.39 & 0.52 & 2.13 & 2.46 & 5.47 & 12.26 & 13.66 \\ 

		0.20 & 1.81 & 0.38 & 2.00 & 2.33 & 5.34 & 12.14 & 13.54 \\ 

	\hline 

		\multicolumn{8}{c} {\rule{0mm}{12pt}for $\Theta_0=4$,  $\varepsilon_1=\varepsilon_2=1$} \\

	\hline 	

		0.10 & 1.93 & 0.26 & 1.29 & 1.58 & 4.64 & 11.33 & 12.68 \\ 

		0.15 & 2.23 & 0.36 & 1.19 & 1.49 & 4.55 & 11.25 & 12.60 \\ 

		0.20 & 2.66 & 0.49 & 1.06 & 1.36 & 4.42 & 11.13 & 12.48 \\ 

		\hline 

	\end{tabular}  

\end{table} 



Количественный анализ выявил закономерности влияния геометрических параметров  и температуры  на величины относительных скоростей потока,  начиная с которых прогибы термоупругой системы увеличиваются во времени, и 

сделать нижеследующий выводы.

 

\begin{enumerate}

	\item[1.] Учет в дифференциальных уравнениях системы \eqref{bel_eq1} слагаемых, учитывающих <<растяжение"=сжатие>> ребер и их <<сдвиг>> в~тангенциальной плоскости (слагаемых, содержащих знаковые числа), практически не влияет на величину относительной скорости ${v_y}/{c_0}$ газового потока. По этой причине удерживать эти слагаемые в уравнениях нет необходимости.

	\item[2.] Увеличение температуры ведет к уменьшению предельной скорости потока при прочих равных условиях.

	\item[3.] Существенно повышают динамическую устойчивость геометрически нерегулярной оболочки параметры ${h_i}/{h}$ и число подкрепляющих элементов~$n$. 

	\item[4.] Величина относительной предельной скорости потока мало чувствительна к изменению относительной кривизны пологой оболочки.

\end{enumerate} 





Перечисленные закономерности выявлены для случая двух полуволн в~направлении потока и одной полуволны в перпендикулярном направлении.  

Следует отметить, что при $k_{11}=0$ полученные решения аналитически преобразуются к решениям, приведенным в работах \cite{myl-bib:31, myl-bib:35}. В случае <<холодной>> ($\Theta_0=0$) гладкой пластины   ($k_{11}=0$, ${a_i}/{a}=0$) решения принимают вид, приведенный в~\cite{myl-bib:30}.



\AdditionalContent{Конкурирующие интересы}{Заявляем, что в отношении авторства и публикации этой статьи конфликта интересов не имеем.} 



\AdditionalContent{Авторский вклад и ответственность}{Все авторы принимали участие в разработке концепции статьи и в написании рукописи. Авторы несут ответственность за предоставление окончательной рукописи в печать. Окончательная версия рукописи была одобрена всеми авторами.} 



\AdditionalContent{Финансирование}{Результаты получены в рамках выполнения государственного задания Минобрнауки России № 9.8570.2017/8.9.}



\begin{thebibliography}{10}



\RBibitem{myl-bib:1}

\by Вольмир~А.~С.

\book Оболочки в потоке жидкости и газа

\yr 1979

\publ Наука

\publaddr М.

\totalpages 320



\RBibitem{myl-bib:2}

\by Амбарцумян~C.~А., Багдасарян~Ж.~Е.

\paper Об устойчивости ортотропных пластинок, обтекаемых сверхзвуковым потоком газа

\jour Изв. АН СССР, ОТН, Механика и машиностроение

\yr 1961

\issue 4

\pages 91--96



\RBibitem{myl-bib:3}

\by Болотин~В.~В., Новичков~Ю.~Н.

\paper Выпучивание и установившийся флайтер термически сжатых панелей, находящихся в сверхзвуковом потоке

\jour Инж. журн.

\yr 1961

\issue 2

\pages 82--96







\RBibitem{myl-bib:4}

\by Мовчан~А.~А.

\paper О колебаниях пластинки, движущейся в газе

\jour ПММ

\yr 1956

\vol 20

\issue 2

\pages 211--222





\RBibitem{myl-bib:5}

\by Дун~Мин-дэ,

\paper Об устойчивости упругой пластинки при сверхзвуковом обтекании

\jour Докл. АН СССР

\yr 1958

\vol 120

\issue 4

\pages 726--729

\mathnet{http://mi.mathnet.ru/dan23108}



\RBibitem{myl-bib:5a}

\by Веденеев~В.~В. 

\paper Высокочастотный флаттер прямоугольной пластины

\jour Изв. РАН. МЖГ

\yr 2006

\issue 4

\pages 173--181









\RBibitem{myl-bib:6}

\by Огибалов~П.~М., Грибанов~В.~Ф.

\book Термоустойчивость пластин и оболочек

\yr 1968

\publ МГУ

\publaddr М.

\totalpages 520







\RBibitem{myl-bib:7}

\by Огибалов~П.~М. 

\book Вопросы динамики и устойчивости оболочек

\yr 1963

\publ МГУ

\publaddr М.

\totalpages 417

\zmath{0115.41504}





\RBibitem{myl-bib:8}

\by Болотин~В.~В.

\paper Температурное выпучивание пластин и пологих оболочек в сверхзвуковом потоке газа 

\inbook Расчеты на прочность

\yr 1960

\volinfo Вып.~6

\publaddr М.

\publ Машгиз

\pages 190--216



\RBibitem{myl-bib:9}

\by Жилин~П.~А.

\paper Линейная теория ребристых оболочек

\jour Изв. АН СССР. МТТ

\yr 1970

\issue 4

\pages 150--166



 



\RBibitem{myl-bib:10}

\by Белосточный~Г.~Н., Ульянова~О.~И.

\paper Континуальная модель композиции из оболочек вращения с термочувствительной толщиной

\jour Изв. РАН. МТТ

\yr 2011

\issue 2

\pages 184--191

\elib{https://elibrary.ru/item.asp?id=15785027}



 

 \RBibitem{myl-bib:11}

 \by Белосточный~Г.~Н., Рассудов~В.~М.

 \paper Континуальная модель термочувствительной ортотропной системы «оболочка–ребра» с учетом влияния больших прогибов

 \inbook Механика деформируемых сред

 \yr 1983

\volinfo Вып.~8

\publaddr Саратов

\publ  Сарат. политехн. ин-т

 \pages 10--22



  

\RBibitem{myl-bib:12}

\by Белосточный~Г.~Н., Рассудов~В.~М.

\paper Континуальный подход к термоустойчивости упругих систем «пластинка–ребра» 

\inbook Прикладная теория упругости

\yr 1980

\publaddr Саратов

\publ  Сарат. политехн. ин-т

\pages 94--99





 









\RBibitem{myl-bib:15}

\by Жилин~П.~А.

\paper Общая теория ребристых оболочек. Прочность гидротурбин

\inbook Тр. ЦКТИ

\yr 1968

\volinfo Вып.~8

\publaddr Л.

\pages 46--70







\RBibitem{myl-bib:16}

\by Карпов~В.~В., Сальников~А.~Ю.

\paper Вариационный метод вывода нелинейных уравнений движения пологих ребристых оболочек 

\jour Вестн. гражд. инженеров

\yr 2008

\issue 4(17)

\pages 121--124

\elib{https://elibrary.ru/item.asp?id=11675045}







\RBibitem{myl-bib:18}

\by  Белосточный~Г.~Н., Мыльцина~О.~А.

\paper Уравнения термоупругости композиций из оболочек вращения 

\jour Вестник СГТУ

\yr 2011

\issue 4 (59)

\issueinfo Вып.~1

\pages 56--64



   

\RBibitem{myl-bib:19}

\by  Онанов~Г.~Г.

\paper Уравнения с сингулярными коэффициентами типа дельта-функции и~ее производных

\jour Докл. АН СССР

\yr 1970

\vol 191

\issue 5

\pages 997--1000

\mathnet{http://mi.mathnet.ru/dan35339}

\mathscinet{http://www.ams.org/mathscinet-getitem?mr=0470298}

\zmath{https://zbmath.org/?q=an:0215.20101}









\RBibitem{myl-bib:21}

\by  Белосточный~Г.~Н.

\paper  Аналитические методы определения замкнутых интегралов сингулярных дифференциальных уравнений термоупругости геометрически нерегулярных оболочек

\jour Докл. Акад. воен. наук

\yr 1999

\issue 1

\pages 14--26





\Bibitem{myl-bib:22a}

\by Geckeler~J.~W.

\paper Elastostatik

\lang In German

\zmath{54.0844.01}

\inbook Handbuch der Physik 

\bookvol 6

\pages 141--308

\yr 1928





\RBibitem{myl-bib:22}

\by Ильюшин~А.~А.

\paper  Закон плоских сечений в аэродинамики больших сверхзвуковых скоростей 

\jour ПММ

\yr 1956

\issue 6

\pages 733--755







 

\RBibitem{myl-bib:30}

\by Вольмир~А.~С. 

\book  Устойчивость деформируемых систем

\yr 1967

\publ Наука

\publaddr М.

\totalpages 984



 

 



 

\RBibitem{myl-bib:25}

\by Владимиров~В.~С. 

\book Обобщенные функции в математической физике

\yr 1979

\publ Наука

\publaddr М.

\totalpages 319

\zmath{0515.46033}









\RBibitem{myl-bib:27}

\by Канторович~Л.~В., Крылов~В.~И.

\book  Приближенные методы высшего анализа

\yr 1962

\publ Физматлит

\publaddr М.

\totalpages 708

\zmath{0100.33102}









\Bibitem{myl-bib:29}

\by Rektorys~K.

\book Variational methods in mathematics, science and engineering

\zmath{0481.49002}

\publaddr Dordrecht, Boston,  London

\publ D.~Reidel Publ. 

\totalpages 571 

\yr 1980

\moreref

\crossref{10.1007/978-94-011-6450-4}. \nofrills





\RBibitem{myl-bib:31}

\by Белосточный~Г.~Н., Рассудов~В.~М. 

\paper  Термоупругие системы типа «пластинка–ребра» в сверхзвуковом потоке газа  

\inbook Прикладная теория упругости

\publaddr Саратов

\volinfo Вып.~8

\publ Сарат. политехн. ин-т

\yr 1983

\pages 114--121







\RBibitem{myl-bib:32}

\by Егоров~К.~В.  

\book  Основы теории автоматического регулирования

\yr 1967

\publ Энергия

\publaddr М.

\totalpages 648





\RBibitem{myl-bib:33}

\by Рассудов~В.~М., Красюков~В.~П., Панкратов~Н.~Д.  

\book  Некоторые задачи термоупругости пластинок и пологих оболочек 

\yr 1973

\publ Сарат. ун-т

\publaddr Саратов

\totalpages 155





\RBibitem{myl-bib:34}

\by Назаров~А.~А.

\book Основы теории и методы расчета пологих оболочек

\yr 1966

\publ  Стройиздат

\publaddr Л., М.

\totalpages 304





\RBibitem{myl-bib:35}

\by Мыльцина~О.~А., Белосточный~Г.~Н. 

\paper  Устойчивость нагретой ортотропной геометрически нерегулярной пластинки в сверхзвуковом потоке газа 

\jour Вестник Пермского национального исследовательского политехнического университета. Механика

\yr 2017

\issue 4

\pages 109--120

\crossref{10.15593/perm.mech/2017.4.08}















\end{thebibliography}





\newpage



\makeentitle



\vspace{-5mm}





\AdditionalContent{Competing interests}{We declare that we have no conflicts of interest in the authorship and publication of this article.} 



\AdditionalContent{Authors' contributions and responsibilities}{Each author has participated in the article concept development and in the manuscript writing. The authors are absolutely responsible for submitting the final manuscript in print. Each author has approved the final version of manuscript.} 





\AdditionalContent{Funding}{The results have been obtained within the State Assignment of the Ministry of Education

and Science of the Russian Federation no. 9.8570.2017/8.9.}









\begin{thebibliography}{10}



\Bibitem{myl-bib:1:en}

\lang In Russian

\by Vol'mir~A.~S.

\book Obolochki v potoke zhidkosti i gaza \rm [The shells in flow liquid and gas] 

\publaddr Moscow

\publ Nauka

\yr 1979

\totalpages 320



\Bibitem{myl-bib:2:en}

\lang In Russian

\by Ambartsumyan~S.~A. Bagdasaryan~Zh.~E.

\paper Of the stability of orthotropic plates streamlined by a supersonic gas flow

\jour Izv. Akad. Nauk SSSR, OTN, Mekhanika i Mashinostroenie

\yr 1961

\issue 4

\pages 91--96



\Bibitem{myl-bib:3:en}

\lang In Russian

\by Bolotin~V.~V., Novichkov~Yu.~N. 

\paper Buckling and steady-state flayter thermally compressed panels in supersonic flow

\jour Inzhenernyi Zhurnal 

\yr 1961

\issue 2

\pages 82--96







\Bibitem{myl-bib:4:en}

\lang In Russian

\by Movchan~A.~A. 

\paper Oscillations of the plate moving in the gas

\jour Prikladnaia Matematika i Mekhanika

\yr 1956

\vol 20

\issue 2

\pages 211--222





\Bibitem{myl-bib:5:en}

\lang In Russian

\by Tung~Ming-teh,

\paper The stability of an elastic plate in a supersonic flow

\jour Dokl. Akad. Nauk SSSR

\yr 1958

\vol 120

\issue 4

\pages 726--729



\Bibitem{myl-bib:5a}

\by Vedeneev~V.~V.

\jour Fluid Dyn.

\crossref{10.1007/s10697-006-0083-2}

\yr 2006

\vol 41

\issue 4

\pages 641–648 

\paper High-frequency flutter of a rectangular plate









\Bibitem{myl-bib:6en}

\lang In Russian

\by Ogibalov~P.~M., Gribanov~V.~F. 

\book Termoustoichivost' plastin i obolochek \rm  [Thermostability of plates and shells]

\yr 1968

\publaddr Moscow

\publ Moscow State Univ.

\totalpages 520







\Bibitem{myl-bib:7en}

\lang In Russian

\by Ogibalov~P.~M.

\book Voprosy dinamiki i ustoichivosti obolochek\rm  [The problems of dynamics and stability of shells]

\publaddr Moscow

\publ Moscow State Univ.

\yr 1963

\totalpages 417







\Bibitem{myl-bib:8:en}

\lang In Russian

\by Bolotin~V.~V.

\paper Thermal buckling of plates and shallow shells in a supersonic gas flow

\inbook Raschety na prochnost'

\publaddr Moscow

\publ Mashgiz

\yr 1960

\volinfo Issue~6

\pages 190--216



\Bibitem{myl-bib:9:en}

\lang In Russian

\by Zhilin~P.~A.

\paper Linear theory of ribbed shells

\jour Izv. Akad. Nauk SSSR, Mekhanika Tverdogo Tela

\yr 1970

\issue 4

\pages 150--166



 



\Bibitem{myl-bib:10:en}

\paper Continuum model for a composition of shells of revolution with thermosensitive thickness

\by Belostochnyi~G.~N., Ul'yanova~O.~I.

\jour Mechanics of Solids

\yr 2011

\vol 46

\issue 2

\pages 184--191

\elib{https://elibrary.ru/item.asp?id=16982711}

\crossref{10.3103/S0025654411020051}



 

\Bibitem{myl-bib:11:en}

\lang In Russian

\by Belostochnyi~G.~N., Rassudov~V.~M.

\paper Continuum model of orthotropic heat-sensitive ``shell-ribs'' system taking into account the influence of large deflection

 \yr 1983

\volinfo Issue~8

\inbook Prikladnaia teoriia uprugosti

\publ Saratov Polytechnic Inst.

\publaddr Saratov 

 \pages 10--22



  

\Bibitem{myl-bib:12:en}

\lang In Russian

\by Belostochnyi~G.~N., Rassudov~V.~M.

\paper Continuum approach to the thermal stability of elastic ``plate-ribs'' systems

\inbook Prikladnaia teoriia uprugosti

\publ Saratov Polytechnic Inst.

\publaddr Saratov 

\yr 1980

\pages 94--99







\Bibitem{myl-bib:15:en}

\lang In Russian

\by Zhilin~P.~A.

\paper General theory of ribbed shells. The strength of turbines

\inbook  Proc. of the Central Turbine-Boiler Institute

\yr 1968

\volinfo Issue~8

\publaddr Leningrad

\pages 46--70







\Bibitem{myl-bib:16:en}

\lang In Russian

\by Karpov~V.~V., Sal'nikov~A.~Yu.

\paper Variational method output nonlinear equations of motion of shallow ribbed shells

\jour Vestnik Grazhdanskikh Inzhenerov

\yr 2008

\issue 4(17)

\pages 121--124





\Bibitem{myl-bib:18:en}

\lang In Russian

\by Belostochnyi~G.~N., Myltcina~O.~A.

\paper Thermoelasticity equations of shells compositions

\jour Vestnik Saratovskogo Tekhnicheskogo Universiteta

\yr 2011

\issue 4 (59)

\issueinfo Issue~1

\pages 56--64



   

\Bibitem{myl-bib:19:en}

\lang In Russian

\by Onanov~G.~G.

\paper Equations with singular coefficients of the type of the delta-function and its derivatives

\jour Dokl. Akad. Nauk SSSR

\yr 1970

\vol 191

\issue 5

\pages 997--1000









\Bibitem{myl-bib:21:en}

\lang In Russian

\by Belostochnyi~G.~N.

\paper Analytical methods for the determination of closed integrals of singular differential equations of thermoelasticity, geometrically irregular shells

\jour Doklady Akademii Voennykh Nauk

\yr 1999

\issue 1

\pages 14--26





\Bibitem{myl-bib:22a:en}

\by Geckeler~J.~W.

\paper Elastostatik

\lang In German

\zmath{54.0844.01}

\inbook Handbuch der Physik 

\bookvol 6

\pages 141--308

\yr 1928





\Bibitem{myl-bib:22:en}

\lang In Russian

\by Il'yushin~A.~A.

\paper Law of plane sections in aerodynamics of high supersonic velocity

\jour Prikladnaia Matematika i Mekhanika

\yr 1956

\issue 6

\pages 733--755







 

\Bibitem{myl-bib:30:en}

\lang In Russian

\by Vol'mir~A.~S.

\book  Ustoichivost' deformiruemykh sistem \rm [Stability of deformable systems]

\publaddr Moscow

\publ Nauka

\yr 1967

\totalpages 984



 

 



 

\Bibitem{myl-bib:25:en}

\by Vladimirov~V.~S.

\book Generalized functions in mathematical physics

\zmath{0515.46034}

\publaddr Moscow

\publ Mir Publ.

\totalpages  362 

\yr 1979









\Bibitem{myl-bib:27:en}

\by Kantorovich~L.~V., Krylov~V.~I.

\book Approximate methods of higher analysis

\zmath{0083.35301}

\publaddr Groningen

\publ P.~Noordhoff Ltd.

\totalpages xii+681

\yr 1958











\Bibitem{myl-bib:29:en}

\by Rektorys~K.

\book Variational methods in mathematics, science and engineering

\zmath{0481.49002}

\publaddr Dordrecht, Boston,  London

\publ D.~Reidel Publ. 

\totalpages 571 

\yr 1980

\moreref

\crossref{10.1007/978-94-011-6450-4}. \nofrills





\Bibitem{myl-bib:31:en}

\lang In Russian

\by Belostochnyi~G.~N., Rassudov~V.~M.

\paper Thermoelastic system of ``plate-ribs'' type  in a supersonic gas flow

\inbook Prikladnaia teoriia uprugosti

\publ Saratov Polytechnic Inst.

\publaddr Saratov 

\volinfo Issue~8

\yr 1983

\pages 114--121







\Bibitem{myl-bib:32:en}

\lang In Russian

\by Egorov~K.~V.

\book Osnovy teorii avtomaticheskogo regulirovaniia \rm [Fundamentals of the theory of automatic control]

\publaddr Moscow

\publ Energiia

\yr 1967

\totalpages 648





\Bibitem{myl-bib:33:en}

\lang In Russian

\by Rassudov~V.~M., Krasiukov~V.~P., Pankratov~N.~D.  

\book  Nekotorye zadachi termouprugosti plastinok i pologikh obolochek  \rm [Some problems of thermoelasticity of plates and flat covers]

\yr 1973

\publ Saratov Univ.

\publaddr Saratov

\totalpages 155

 





\Bibitem{myl-bib:34:en}

\lang In Russian

\by Nazarov~A.~A.

\book Osnovy teorii i metody rascheta pologikh obolochek \rm [Fundamentals of the Theory and Methods for Designing Shallow Shells]

\yr 1966

\publ  Stroiizdat

\publaddr Leningrad, Moscow

\totalpages 304





\Bibitem{myl-bib:35:en}

\lang In Russian

\by Myltcina~O.~A., Belostochnyi~G.~N.

\paper Stability of heated orthotropic geometrically irregular plate in a supersonic gas flow

\jour PNRPU Mechanics Bulletin

\yr 2017

\issue 4

\pages 109--120

\crossref{10.15593/perm.mech/2017.4.08}





\end{thebibliography}







%\end{document}

 \newpage

%\input{./preamble}
%
%\usepackage{samgtu-name}
%
\vsgtu{1624} %registration number
\FirstPage{762}
\LastPage{773}
%
%\pdfcrop
%\draft  
%
\ShortCommunication
%
%\begin{document} 




\hfill \qrcode[hyperlink,height=.8in]{http://doi.org/10.14498/vsgtu1624}
\vspace{-.9in}

\udc{539.3}
\msc{74C10}

\rutitle{Об одном подходе к определению предельной несущей~способности
	механических систем с~разупрочняющимися элементами}
\rutitleshort{Об одном подходе к определению предельной несущей способности
	механических систем\dots}

\entitle{One approach to determination of the ultimate load-bearing capacity of mechanical systems with softening elements}
\entitleshort{One approach to determination of the ultimate load-bearing capacity of mechanical systems\dots}

\ruauthor{В.~В.~Стружанов\affil{1}, А.~В.~Коркин\affil{2}, А.~Е.~Чайкин\affil{2}}
{С\,т\,р\,у\,ж\,а\,н\,о\,в~В.~В., К\,о\,р\,к\,и\,н~А.~В., Ч\,а\,й\,к\,и\,н~А.~Е.}
\enauthor{V.\,V.~Struzhanov\affil{1}, A.\,V.~Korkin\affil{2}, A.\,E.~Chaykin\affil{2}} {S\,t\,r\,u\,z\,h\,a\,n\,o\,v~V.~V., K\,o\,r\,k\,i\,n~A.~V., C\,h\,a\,y\,k\,i\,n~A.~E.}

\ruaffil{\ruimash.
\and 
\ruurfuien.}

\enaffil{\enimash.
\and 
\enurfuien.}

\ruauthorsdetails{3}{% Количество авторов
\fullauthor{\rupid{39278}{Валерий Владимирович Стружанов}}
\CAuthor % Автор-корреспондент
\orcid{0000-0002-3669-2032} % ORCID ID автора 
\regalia{доктор физико-математических наук, профессор} 
\position{главный научный сотрудник} 
\dept[p]{~лаб. микромеханики материалов} 
\email{stru@imach.uran.ru}
\fullauthor{\rupid{73471}{Александр Владимирович Коркин}}
\orcid{0000-0003-3533-4257} % ORCID ID автора 
\regalia{аспирант} 
\email{alexkorkin@list.ru}
\fullauthor{\rupid{143469}{Алексей Евгеньевич Чайкин}}
\orcid{0000-0001-5582-2384} % ORCID ID автора 
\regalia{аспирант}  
\email{chaykin.ae@yandex.ru}\vspace{-3mm}
}

\enauthorsdetails{3}{
\fullauthor{\enpid{39278}{Valery V. Struzhanov}}
\CAuthor % Сorresponding Author
\orcid{0000-0002-3669-2032} % ORCID ID avtora 
\regalia{Dr. Phys. \& Math. Sci., Professor} 
\position{Chief Reseacher} 
\dept{Lab. of Matherial Micromechanics}
\email[br]{stru@imach.uran.ru}
\fullauthor{\enpid{73471}{Aleksandr V. Korkin}}
\orcid{0000-0003-3533-4257} % ORCID ID avtora 
\regalia{Postgraduate Student} 
\email{alexkorkin@list.ru}
\fullauthor{\enpid{143469}{Aleksey E.  Chaykin}}
\orcid{0000-0001-5582-2384} % ORCID ID avtora 
\regalia{Postgraduate Student} 
\email{chaykin.ae@yandex.ru}
}

\rukeywords{потенциальная функция, вырожденные критические точки, единый потенциал, матрица Гессе, тонкостенный резервуар, предельное давление}
\enkeywords{potential function, degenerate critical points, unified potential, Hesse matrix, thin walled reservoir, limiting pressure}

\ruabstract{Изложены основные положения теории расчета предельных нагрузок, действующих на дискретные механические системы с~разупрочняющимися элементами. Методика опирается на численное определение вырожденных критических точек потенциальной функции системы, где происходит переход от устойчивости процесса нагружения к~неустойчивости (катастрофа, разрушение), и~позволяет избежать решения большого числа нелинейных уравнений равновесия. 
В~качестве примера применения методики решена задача об определении предельного внутреннего давления в~тонкостенном цилиндрическом резервуаре. При построении потенциальной функции системы использован специально построенный единый потенциал для плоского квадратного элемента материала в~условиях двухосного растяжения, описывающий все стадии деформирования, включая и~разупрочнение.}

\enabstract{The fundamental provisions of the limiting load calculation theory are presented for a discrete mechanical system with softening elements. The method is based on the numerical determination of degenerate critical points for the potential function of the system. At these points there is a transition from the stability of the loading process to instability such as a catastrophe or a failure. This approach helps to avoid solving a large number of nonlinear equilibrium equations. The problem of determining the limiting internal pressure in a thin walled cylindrical tank is solved as an example. A unified potential specially defined for a flat square element of material in biaxial tension is used in developing a potential function of the system. It describes all stages of deformation including the softening stage.}

\newdate{struzha}{11}{05}{2018}
\date{\displaydate{struzha}}
\newdate{struzhb}{11}{10}{2018}
\revisiondate{\displaydate{struzhb}}
\newdate{struzhc}{12}{11}{2018}
\accepteddate{\displaydate{struzhc}}

\newdate{struzhe}{11}{12}{2018}
\renewcommand{\JournalDataOnline}{\displaydate{struzhe}}
%\ForPeerReview
\makerutitle

\newpage

\Section[n]{Введение}
Вычисление предельной несущей способности механических систем основывается на расчете параметров равновесных состояний при постепенно возрастающей нагрузке. Затем эти параметры тем или иным способом сводят к некоторому скаляру и сравнивают его с предельной величиной, полученной в эксперименте. Например, на каждом шаге догружения определяют напряженное состояние элемента конструкции, находят соответствующий инвариант тензора напряжений, который сравнивают с его предельной величиной, взятой из эксперимента.

Возможен и другой подход, основанный на исследовании устойчивости положений равновесия, когда разрушение связывают с невозможностью равновесия, то есть с моментом потери устойчивости [\citen{bib1, bib2}]. Реализация данного подхода возможна, если учитывать физически неустойчивые состояния материала (разупрочнение) [\citen{bib2}--\citen{bib6}]. Но и в этом случае возникает необходимость в решении не удовлетворяющих условиям корректности по Адамару [\citen{bib8}]  больших систем нелинейных уравнений равновесия для определения критических точек потенциальной функции механической системы, что является нетривиальной задачей [\citen{bib9}]. Кроме того, выделение из них вырожденных критических точек, где происходит смена устойчивости на неустойчивость, требует применения нетрадиционного аппарата теории катастроф [\citen{bib10,bib11}]. Однако информация о параметрах равновесия (критических точках потенциальной функции системы) при возрастающих нагрузках, вообще говоря, является излишней. При определении предельной несущей способности нужно знать только вырожденные критические точки. В данной работе представлена методика, которая позволяет избежать аналитического решения нелинейных уравнений. 
Она основана на специальной числовой процедуре приближенного выбора вырожденных критических точек. В качестве примера рассмотрена задача о~разрушении цилиндрического резервуара под действием внутреннего давления. Для построения потенциальной функции разработана методика построения единого потенциала для элемента материала, находящегося в плоском напряженном состоянии и описывающего его свойства как при упрочнении, так и при разупрочнении.

\smallskip

\Section{Общие положения} Положение механической системы в пространстве определяется конечным набором обобщенных координат (параметров состояния)  $ q_i$, $i = 1, 2, \ldots, N $. 
Величины параметров состояния зависят от набора задаваемых параметров управления системой $ s_j$, 
$j = 1, 2, \ldots, M $ 
(внешние силы, перемещения некоторых точек системы, физико"=механические свойства элементов и т.~п.). 
Причем параметры состояния должны принимать такие значения, при которых механическая система находится в состоянии равновесия. 
Разрушение системы есть явление того же порядка, что и явление невозможности равновесия [\citen{bib1}]. 
Когда равновесие становится неустойчивым, то реализуется или скачкообразный переход в ближайшее устойчивое равновесие, или глобальная катастрофа (разрушение). Поэтому под предельными значениями параметров управления следует понимать такие их величины, при реализации которых нарушается устойчивость системы.

В тех случаях, когда механическая система является градиентной [\citen{bib11}], для определения предельных значений параметров управления можно воспользоваться следующей методикой. Градиентная система при квазистатическом нагружении с постоянной температурой описывается потенциальной функцией (лагранжианом) $ \Pi(q_i ,s_j) $, связывающей параметры состояния и~управления. 
Эта функция есть сумма потенциальных функций элементов системы, которые могут быть как вогнутыми, так и~выпукло"=вогнутыми. 
Выпуклые вниз (вогнутые) участки отвечают  устойчивым состояниям элемента, выпуклые (выпуклые вверх) "--- неустойчивым.  
Точки перехода от выпуклости к~вогнутости "--- пограничные точки, являющиеся в общем случае седловыми точками соответствующих потенциальных функций. 

Когда механическая система находится в~равновесии (устойчивом или неустойчивом), для связи параметров состояния и~управления в~этом равновесии можно использовать уравнение Лагранжа второго рода [\citen{bib12}]. 
Тогда уравнение равновесия имеет вид
\begin{equation}\label{S_eq1}
\frac{\partial\Pi(q_i,s_j)}{\partial q_i} = 0, \quad  i=1, 2, \ldots,N.
\end{equation}
Если функция   является выпукло"=вогнутой, то при заданных управляющих параметрах система \eqref{S_eq1} может иметь одно или более одного решения, как устойчивых, так и~неустойчивых. 
Следовательно, в~данном случае не выполняются условия Адамара [\citen{bib8}].

Все решения системы \eqref{S_eq1} для всевозможных параметров управления образуют совокупность критических точек потенциальной функции $ \Pi $.  
Известно [\citen{bib11}], что смена типа равновесия (устойчивого или неустойчивого) происходит в~вырожденных критических точках, которые определяются из совместного решения уравнений \eqref{S_eq1} и~уравнения, получающегося приравниванием к~нулю детерминанта матрицы устойчивости (матрицы Гессе потенциальной функции):
\begin{equation}\label{S_eq2}
\mathop{\rm det}\Bigl| \frac{\partial^{2}\Pi(q_i,s_j)}{\partial q_m \partial s_n}\Bigr|  = 0,  \quad m,n = 1, 2, \ldots, N.
\end{equation}
Очевидно, что уравнение \eqref{S_eq2} выделяет из множества критических точек вырожденные критические точки. 

Таким образом, для отыскания предельных значений параметров управления, при достижении которых система теряет устойчивость, сначала, как правило, вычисляются все критические точки (решения уравнений \eqref{S_eq1}), затем выделяются вырожденные точки и после подстановки их в уравнение \eqref{S_eq1} определяются соответствующие параметры управления. Те параметры управления, при которых в первый раз достигается вырожденная критическая точка, то есть происходит смена устойчивого состояния механической системы на неустойчивое, и будут предельными значениями внешних нагрузок.

При реализации данной системы нахождения предельных значений параметров управления необходимо искать решения большого числа нелинейных уравнений, что представляет достаточно сложную задачу. 
Однако задача о~вычислении предельных значений параметров управления не требует, вообще говоря, знания всех критических точек функции $ \Pi $ (параметров всех положений равновесия механической системы). Необходимы только вырожденные точки, удовлетворяющие уравнению \eqref{S_eq2}. Кроме того, для технических приложений можно использовать их приближенные значения, вычисленные с достаточной степенью точности. 

Вычисление приближенных значений вырожденных критических точек можно осуществить, применяя следующую численную процедуру. Рассмотрим евклидово пространство $ \mathbb R^M \times \mathbb R^N $ "--- декартово произведение евклидовых пространств размерности $ M $ и $ N $. 
Элементами этого пространства являются числа $ ( q_i,s_j) $. 
Задавая физически обоснованные пределы изменения параметров состояния и управления, получаем $ (M{+}N) $"=мерный куб в пространстве $\mathbb  R^M \times \mathbb R^N $, 
который разбиваем сеткой узлов на множество «кубиков» с~заданным шагом. 
Затем в~каждом узле сетки разбиения вычисляем значения компонент матрицы Гессе и~величину ее определителя. 
Таким образом, каждому узлу ставится в~соответствие значение детерминанта матрицы Гессе. 

После этого выделяем те узлы, в которых детерминант близок к нулю с достаточной степенью точности. 
Полученное множество $B$ точек пространства $ \mathbb R^M \times \mathbb R^N $ содержит в~себе все вырожденные критические точки потенциальной функции~$ \Pi $. 
Чтобы выделить из множества $B$ критические точки функции $ \Pi $, следует подставить элементы множества $B$ в~уравнения равновесия \eqref{S_eq1}.  
Те~из них, которые удовлетворяют этим уравнениям с~заданной степенью точности, будут вырожденными критическими точками. 
Параметры управления,  соответствующие ближайшей (по евкледовой метрике) к~началу координат вырожденной точке, образуют совокупность искомых предельных значений управлений. 
Если существует группа вырожденных точек, находящихся на одинаковом наименьшем расстоянии от начала координат, то получим различные сочетания значений параметров управления.

Таким образом, минуя стадию расчета параметров равновесных состояний, можно найти критические нагрузки, что достаточно для оценки прочности механических систем.

\smallskip

\Section{Модельный пример} Применим изложенную выше методику к определению предельных параметров управления (нагрузок) при растяжении модельной стержневой системы, изображенной на  рис.~\ref{Sfig1}. 
Стержень {\sl 2} "--- абсолютно упругий с жесткостью при растяжении, равной $c$. 
Стержень {\sl 1} изготовлен из разупрочняющегося материала.

\begin{figure}[b!]
\centering
		\includegraphics[width=0.4\textwidth]{fig_stru1}
		\caption{Растягиваемая стержневая система с разупрочняющимся элементом} \label{Sfig1}
		\smallskip \footnotesize
		[Figure~\ref{Sfig1}. 		A rod system with a softening element under tension]

\end{figure}

\smallskip

{\bf 2.1.}
При мягком нагружении растягивающая сила $ p $ "--- параметр управления, 
а~удлинение $ x $  стержня {\sl 1}  и перемещение $ u $ правого свободного конца   стержня~{\sl 2}  "--- параметры состояния (см.  рис.~\ref{Sfig1}). 
Сопротивление растяжению стержня {\sl 1} задано функцией $ g(x) $ 
(зависимость растягивающей силы от удлинения), обладающей как восходящей ветвью (упрочнение), так и~ветвью, падающей до нуля (стадия разупрочнения) [\citen{bib2, bib5}].  
В~этом случае его потенциальная функция, равная $\displaystyle \int ^{x}_{0} g(x)\, dx $, является выпукло"=вогнутой, а лагранжиан системы~[\citen{bib5}] 
$$
Q^p = \int^{x}_{0} g(x)  dx + \frac{c}{2} (u-x)^2 - \int ^{u}_{0} p  du,
$$	
где второе слагаемое "--- энергия упругих деформаций стержня {\sl 2}, последнее "--- работа внешней силы, взятая со знаком минус.

Уравнения равновесия имеют вид
\begin{equation}\label{S_eq3}
Q^{p}_{\bigcom x} = g(x) - c(u-x) = 0, \quad  Q^{p}_{\bigcom u} = c(u-x) - p = 0.
\end{equation}
Здесь запятой обозначены частные производные по соответствующим аргументам.
Приравнивая к нулю детерминант матрицы Гессе, получим уравнение 
\begin{equation}\label{S_eq4}
g_{\bigcom x} = 0.
\end{equation}

Отметим, что это уравнение определяет максимум функции  $ g(x) $. 
Его можно разрешить, используя следующую процедуру. 

Проведем разбиение оси $X $ узлами с~достаточно малым шагом. 
Подставляя координаты этих узлов в~уравнение \eqref{S_eq4}, находим среди них значение $ x' $, приближенно ему удовлетворяющее. 
Затем из уравнений \eqref{S_eq3} для данного $ x' $ определяем $ p' $ "--- предельную величину растягивающей силы.

\smallskip

{\bf 2.2.}
В случае жесткого нагружения, когда растяжение осуществляется заданным перемещением $ u $ (параметр управления), а $ x $ является единственным параметром состояния, лагранжиан системы есть функция $ Q^u $, равная двум первым слагаемым в выражении для $ Q^p $. 

Тогда имеем лишь одно уравнение равновесия (первое уравнение в~системе \eqref{S_eq3}). 
В~матрице Гессе только один элемент. 
Тогда равенство детерминанта нулю дает уравнение $ g'_{x} + c = 0 $. 
Его приближенное решение ищем так же, как и~решение уравнения \eqref{S_eq4}. 

В зависимости от величины параметра $ c $ возможны три варианта. 

В первом "--- решения не существует. Следовательно, деформирование системы идет равновесно вплоть до разделения стержня {\sl 1} на части при значении $ x = x^z $, когда падающая ветвь функции $ g(x) $ достигает оси $ X $. Предельное значение параметра управления получаем из уравнения равновесия после подстановки $ x = x^z $: 
$$ u^z = g(x^z)/c+ x^z.$$ 

Во втором варианте имеется только одно приближенное решение $ x =  x_1 $. 
После достижения параметром $ x $ значения этой величины система переходит в~неустойчивое состояние. 
Поэтому предельный параметр управления $$ u_1 = g(x_1)/c + x_1 .$$ 

И, наконец, в третьем варианте имеем два решения $ x_1$, $x_2$  $(x_1 > x_2) $. 
Тогда предельное значение  $ u = u_1 $.

Отметим, что такие же результаты были получены и в работах [\citen{bib5, bib9}] только после решения нелинейных уравнений равновесия при постепенном возрастании параметров управления и построения соответствующих многообразий катастроф и сепаратрис, состоящих из вырожденных критических точек.

\smallskip

\Section{Вогнуто"=выпуклый потенциал при двухосном растяжении тонкой пластины} 
Для реализации положений теории, изложенной выше, необходимо знать вогнуто"=выпуклый потенциал элементов системы. 
В~механике сплошной среды эта функция должна быть определена для элемента материала вплоть до разупрочнения и~разделения на фрагменты.

Рассмотрим плоский квадратный элемент материала с~единичными размерами, находящийся в плоско"=напряженном состоянии, причем касательные напряжения и~деформации отсутствуют. 
Деформирование осуществляется заданием нормальных деформаций $ \varepsilon_1$, $\varepsilon_2 $ (двуосное растяжение). Если материал находится в состоянии упругости, то его  потенциальная энергия 
\begin{equation}\label{S_eq5}
V^e(\varepsilon_1, \varepsilon_2) = \frac{1}{2}(\sigma_1\varepsilon_1 + \sigma_2\varepsilon_2) = \frac{E}{2(1-\nu^2)} (\varepsilon^2_1 + 2\nu \varepsilon_1 \varepsilon_2 + \varepsilon^2_2 ) ,
\end{equation}
где $ E $ "--- модуль Юнга, $ \nu $ "--- коэффициент Пуассона. 
Здесь использован закон Гука, связывающий нормальные напряжения $ \sigma_1$, $\sigma_2 $ и~нормальные деформации $ \varepsilon_1$, $\varepsilon_2 $ и записанный в векторно"=матричной форме:
\begin{equation}\label{S_eq6}
\begin{pmatrix}
\sigma_1 \\ \sigma_2
\end{pmatrix}
 = C 
\begin{pmatrix} \varepsilon_1 \\ \varepsilon_2
\end{pmatrix}.
\end{equation}
В этом выражении 
$$
C = \frac{E}{(1-\nu^2)} \begin{pmatrix} 1 & \nu \\ \nu & 1 \end{pmatrix}
$$
можно трактовать как матричный коэффициент пропорциональности.

Формула \eqref{S_eq5} определяет кривые второго порядка на плоскости $\varepsilon_1 \varepsilon_2 $ 
(линии уровня потенциала $ V^e $), а~именно центральные эллипсы, главные оси которых наклонены к~декартовым осям $ \varepsilon_1$, $\varepsilon_2 $ под углом $ \pi/4 $. 
Большая ось располагается во втором и~четвертом квадрантах, малая "--- в~первом и~третьем. 
Для полуосей  выполняется отношение 
$$ \dfrac{a^{2}}{\rho^2} = \frac{(1-\nu)}{(1+\nu)}, $$ 
где $ a $ "--- малая полуось, $ \rho $ "--- большая полуось. 
Точки, расположенные на малой оси, отвечают равномерному растяжению $ (\varepsilon_1 = \varepsilon_2) $ или сжатию $ (-\varepsilon_1 = -\varepsilon_2) $, а~точки на большой оси определяют чистый сдвиг   
$ (\varepsilon_1 = -\varepsilon_2$, $ -\varepsilon_1 = \varepsilon_2 ) $. 
Произвольная точка  $ (\varepsilon_1,\varepsilon_2) $ лежит на линии уровня с большой полуосью, равной
\begin{equation}\label{S_eq7}
\rho = \frac{1}{\sqrt{2}}\sqrt{{(\varepsilon_1 + \varepsilon_2)}^2 \frac{1+\nu}{1-\nu} + {(\varepsilon_1 - \varepsilon_2)}^2} .
\end{equation}

При чистом сдвиге $ (\varepsilon_1 = \varepsilon$, $\varepsilon_2 = - \varepsilon)$  $\rho = \gamma /\sqrt{2}$, где    $ \gamma = \varepsilon_1 - \varepsilon_2 = 2 \varepsilon $ "--- сдвиг.
После исчерпания упругости для описания свойств материала можно было бы использовать теорию малых упруго-пластических деформаций [\citen{bib13, bib14}], в~которой предполагается упругость объемной деформации и~пропорциональность девиаторов тензоров напряжений и деформаций. 
Однако после выхода на стадию разупрочнения первая гипотеза не выполняется, так как материал уже существенно поврежден множеством микротрещин и объемный модуль изменяется. 

Чтобы остаться в рамках деформационной теории, следует отказаться от первой гипотезы, а~вместо второй сохранить пропорциональность напряжений и~деформаций с некоторым уже переменным, зависящим от деформаций матричным коэффициентом. 
Например, в рассматриваемой задаче %после перехода материала  при активном нагружении 
этот матричный коэффициент в законе \eqref{S_eq6} можно принять в~виде $ C \psi(\varepsilon_1,\varepsilon_2)$, 
где $ \psi $ "--- некоторая скалярная функция, требующая специального определения.

Следствием принятого предположения является тот факт, что в~области неупругости линии уровня функции энергии деформаций $ V $ (потенциальной энергии) подобны линиям уровня потенциала $ V^{e} $, то есть представляют собой эллипсы с~тем же самым отношением главных осей. 
Необходимо только знать распределение значений потенциала по  этим линиям уровня.

Воспользуемся полной диаграммой для чистого сдвига $ \tau(\gamma) $ ($ \tau $ "--- касательное напряжение, $ \gamma $ "--- деформация сдвига), которая обладает восходящей и нисходящей до нуля ветвями. 
Такую диаграмму материала можно получить, используя опытную диаграмму кручения стержней  круглого поперечного сечения с~пересчетом на диаграмму материала [\citen{bib15}]. 
Тогда энергия, затраченная на деформирование, есть   
$$ V = \int ^{\gamma}_{0}\tau(\gamma)  d\gamma .$$

После взятия интеграла и замены $ \gamma $ на $ \sqrt{2} \rho $  с~учетом выражения \eqref{S_eq7} получаем искомый единый потенциал $ V(\varepsilon_1,\varepsilon_2) $. 

Пусть $ \tau = 0.1G   \gamma(1-5 \gamma)$, если $ 0 \le \gamma \le 0.2 $, где  
$ G = \frac{E}{2(1+\nu)} $ "---  модуль сдвига в упругости. 
Тогда
\begin{equation}\label{S_eq8}
V(\varepsilon_1,\varepsilon_2) = 
\left\lbrace  
\begin{array}{cl}
\frac{V^{e}}{2}\Bigl( 0.1 - \frac{1}{3}\sqrt{\frac{V^{e}}{G}}\Bigr), & (\varepsilon_1,\varepsilon_2) \in \Omega , \\
0.002 G/3 , & (\varepsilon_1,\varepsilon_2) \notin \Omega.
\end{array}
\right. 
\end{equation}


\begin{figure}[b!]
\begin{tabular}{p{0.35\textwidth}p{0.6\textwidth}}

\includegraphics[width=0.3\textwidth]{fig_stru2} & 

\vspace{-5.5cm}

~~~\,Связь между напряжениями и деформациями (определяющие соотношения) в данном случае будет задана законом
$$
\begin{pmatrix} \sigma_1  \\ 
\sigma_2 
\end{pmatrix}
=
\begin{pmatrix} \frac{\partial V}{\partial \varepsilon_1} \\ 
\frac{\partial V}{\partial \varepsilon_2} \end{pmatrix}
=
C \psi (\varepsilon_1,\varepsilon_2) 
\begin{pmatrix} \varepsilon_1 \\ 
\varepsilon_2 \end{pmatrix},
$$
где 
$$ 
\psi (\varepsilon_1,\varepsilon_2) = \frac{1}{2} \Bigl(0.1 - \frac{1}{2} \sqrt{\frac{V^e}{G}} \Bigr).  
$$


~~~\,Отметим, что напряжения являются компонентами градиента функции~$ V $.

\\

\vspace{-.8cm}
\caption{Качественный вид поверхности потенциала $ V(\varepsilon_1,\varepsilon_2) $ } \label{Sfig2}
		\smallskip \footnotesize
		[Figure~\ref{Sfig2}. Qualitative form of the \centerline{potential surface $ V(\varepsilon_1,\varepsilon_2) $]} & \\
\end{tabular}
\end{figure}

%\begin{figure}[b!]
%	\begin{center}
%		\includegraphics[width=0.3\textwidth]{fig_stru2}
%		\caption{Качественный вид поверхности потенциала $ V(\varepsilon_1,\varepsilon_2) $ } \label{Sfig2}
%		\smallskip \footnotesize
%		[Figure~\ref{Sfig2}. Qualitative form of the potential surface $ V(\varepsilon_1,\varepsilon_2) $]
%	\end{center}
%\end{figure}



Очевидно, что область $ \Omega $ ограничена эллипсом с~полуосями 
$ \rho = 0.141$; $a = \rho\sqrt{(1+\nu)/(1-\nu)} $.  
Качественный вид поверхности $ V(\varepsilon_1,\varepsilon_2) $ изображен на рис.~\ref{Sfig2}. 
Функция $ V $ описывает все состояния материала (выпуклость вниз "--- устойчивость, упрочнение; 
выпуклость вверх "--- неустойчивость, разупрочнение).





\smallskip


\Section{Разрушение цилиндрического резервуара} Рассмотрим тонкостенный резервуар в~виде длинной трубы с~заделанными абсолютно жесткими крышками торцами. 
Радиус трубы равен $ b $, длина $ l $, толщина стенки $ t$  $(t\ll b)$. 
Резервуар находится под действием  монотонно и квазистатически возрастающего при постоянной температуре внутреннего давления с~интенсивностью~$ p $ [\citen{bib16, bib17}]. 
Материал обладает свойством деформационного разупрочнения. 
Применим изложенную выше методику для расчета предельного значения давления. 

Каждый элемент трубы находится в~плоском напряженном состоянии при всестороннем растяжении: 
$ \varepsilon_1 = \varepsilon_z $ "--- продольная деформация, 
$ \varepsilon_2 = \varepsilon_{\theta} $ "--- тангенциальная деформация. Его потенциальная функция $ V $ задана выражением \eqref{S_eq8}. 
%где  $ \varepsilon_1 = \varepsilon_z , \varepsilon_1 = \varepsilon_{\theta} $. 

Механическая система (резервуар) обладает двумя параметрами состояния (обобщенными координатами) $ \varepsilon_z$, $\varepsilon_{\theta} $ (однородное напряженно"=деформированное состояние) и~одним параметром управления $ p $. 

Лагранжиан (потенциальная функция) этой системы
$$
W = V(\varepsilon_z , \varepsilon_{\theta})2\pi b l t - p \pi b^2 l \varepsilon_z - 2\pi p b^2 l \varepsilon_{\theta}.
$$
Здесь первое слагаемое "--- полная энергия деформаций стенок трубы, второе и третье "--- взятые со знаком минус соответственно работа сил давления на перемещении жестких крышек и~перемещении, равном увеличению радиуса. 
Тогда уравнения равновесия, определяющие критические точки функции $ W $, имеют вид
\begin{equation}\label{S_eq9}
\frac{\partial W}{\partial \varepsilon_z} = \frac{\partial V}{\partial \varepsilon_z} 2\pi b l t  - p \pi b^2 l = 0;
\quad  
\frac{\partial W}{\partial \varepsilon_{\theta}} = \frac{\partial V}{\partial \varepsilon_{\theta}} 2\pi b l t  - 2p \pi b^2 l = 0.
\end{equation}
Уравнения \eqref{S_eq9} сводятся к одному:
$$
2 \frac{\partial V}{\partial \varepsilon_z} = \frac{\partial V}{\partial \varepsilon_{\theta}}.
$$
Производя необходимые вычисления, находим, что
\begin{equation}\label{S_eq10}
\varepsilon_{\theta} = \frac{2-\nu}{1-2\nu} \varepsilon_{z},
\end{equation}
то есть в пространстве деформаций путь деформирования элемента материала является прямой линией, уравнение которой задается равенством \eqref{S_eq10}. 
Кроме этого, из уравнений \eqref{S_eq9} следует, что  $ \sigma_z = pb/2t$, $\sigma_{\theta} = pb/t $. 
Отсюда путь нагружения в~пространстве напряжений также является прямой линией. 
Таким образом, имеем ситуацию с пропорциональным нагружением элемента материала.

Определим, наконец, предельное значение внутреннего давления. 
Как следует из приведенных выше теоретических положений, для этого необходимо гессиан функции $ W $ приравнять к~нулю, найти все решения полученного уравнения, по крайней мере, приближенно, и выделить из них вырожденные критические точки потенциала $ W $, которые принадлежат множеству критических точек (решений уравнений \eqref{S_eq9}) потенциала. 
Имеем 
$$
\Bigl\lbrace 
\frac{\partial^{2} V}{\partial \varepsilon^{2}_{z}} \cdot \frac{\partial^{2} V}{\partial \varepsilon^{2}_{\theta}} - 
\Bigl( \frac{\partial^{2} V}{\partial \varepsilon_{z}\partial \varepsilon_{\theta}}\Bigr) 
\Bigr\rbrace 
4 \pi^{2} b^2 l^2 t = 0.
$$
После несложных преобразований получаем
\begin{multline}\label{S_eq11}
\bigl(0.05 - \alpha \beta - \alpha \beta^{-1}(\varepsilon_{z} + \nu \varepsilon_{\theta})^2\bigr) 
\bigl(0.05 - \alpha \beta - \alpha \beta^{-1}(\varepsilon_{\theta} + \nu \varepsilon_{z})^2\bigr)-
\\
 - \nu^2 \bigl( 0.05 - \alpha \beta - \alpha \beta^{-1}(\varepsilon_{z} + \nu \varepsilon_{\theta}) (\varepsilon_{\theta} + \nu \varepsilon_{z})\bigr)^2 = 0
,                                          
\end{multline}
где 
$ \alpha=0.25 (1-\nu)^{ - 1/2}$, 
$ \beta = (\varepsilon^{2}_{z} + 2\nu \varepsilon_{z} \varepsilon_{\theta} + \varepsilon^{2}_{\theta})^{1/2}$;  
$ \varepsilon_{z}$, $\varepsilon_{\theta} \ge 0$, $(\varepsilon_{z}, \varepsilon_{\theta}) \in \Omega $.

Для вычисления вырожденных критических точек функции $ W $ применим следующую численную процедуру. 




Построим на плоскости  $ \varepsilon_{z} \varepsilon_{\theta} $ квадрат $ \Omega_1 $  такой, что в нем $ \varepsilon_{z}$, $\varepsilon_{\theta} > 0$,     $(\varepsilon_{z}, \varepsilon_{\theta}) \in \Omega \subset  \Omega_1 $ (рис.~\ref{Sfig3}). 
Для определения стороны этого квадрата найдем на осях $ \varepsilon_{\theta}$ и $\varepsilon_{z} $ равные отрезки, внутри которых детерминант матрицы Гессе меняет знак. 
Получаем отрезки длиной $0.1$. 
Эту величину и берем за сторону квадрата $\Omega_1$. 
Отметим, что данный прием может быть использован при построении $M \times N$-мерного куба в общей теории. 



\begin{figure}[h!]
\centering
	\includegraphics[width=0.4\textwidth]{fig_stru3}
	\caption{Вид кривой, в точках которой гессиан равен нулю} \label{Sfig3} 
	\smallskip \footnotesize
		[Figure~\ref{Sfig3}. The curve where the Hessian is equal to zero at all points]
\end{figure}

Разобьем этот квадрат прямоугольной сеткой с~достаточно малым шагом (на рис. \ref{Sfig3} показана укрупненная сетка). В~каждом узле сетки по формуле \eqref{S_eq11} при  $ \nu = 0.3 $ вычисляем детерминант (гессиан) матрицы Гессе. 
Выделяем те из них, где детерминант мало отличен от нуля, что определяется заданной точностью вычислений.

Производя расчеты, получаем множество точек, расположенных на кривой {\sl 1} (рис. \ref{Sfig3}). 
Затем координаты точек кривой {\sl 1} подставляем в уравнения \eqref{S_eq9} и~выделяем те из них, которые удовлетворяют уравнениям с~заданной точностью. 
В~результате находим только одну точку (точка $A$, рис. \ref{Sfig3}). Точка $ A $ лежит на прямой \eqref{S_eq10} и~имеет координаты $ \varepsilon_{z}  = 0.018$, $\varepsilon_{\theta}  = 0.076 $ (с~точностью до округления). 
Это и~есть искомая вырожденная критическая точка потенциальной функции $ W $. 
Подставляя ее координаты в~уравнения равновесия, находим критическое внутреннее давление, при котором происходит потеря устойчивости системы (катастрофа, разрушение). 
Если  $ E = 2 \cdot 10^5 $~МПа, $ b = 100 $ мм,  $ t = 2 $ мм, то критическое (предельное) значение давления $ p = 18 $~МПа. Заметим, что величина предельного давления увеличивается пропорционально толщине стенки.


\pagebreak



\AdditionalContent{Конкурирующие интересы}{Мы не имеем конкурирующих интересов.} 
%\AdditionalContent{Конкурирующие интересы}{Конкурирующих интересов не имею.} 
\AdditionalContent{Авторский вклад и ответственность}{Все авторы принимали участие в разработке концепции статьи и в написании рукописи. Все авторы несут полную ответственность за представление окончательной рукописи в печать. Окончательная версия рукописи была одобрена всеми авторами.} 
%\AdditionalContent{Авторская ответственность}{Я несу полную ответственность за предоставление окончательной версии рукописи в печать. Окончательная версия рукописи мною одобрена.} 
%\AdditionalContent{Авторский вклад и ответственность}{Все  авторы  принимали  участие  в  разработке  концепции  \mbox{статьи}  и  в  написании  рукописи. Авторы несут  полную  ответственность  за  предоставление  окончательной рукописи в печать. Окончательная  версия рукописи была одобрена всеми авторами.} 
\AdditionalContent{Финансирование}{Исследование не имело финансирования.}
%\AdditionalContent{Финансирование}{Исследование выполнялось без финансирования.}
%\AdditionalContent{Финансирование}{Работа выполнена при поддержке РФФИ (проект № xx–xx–xxxxx_a).}
%\AdditionalContent{Благодарность}{Выразите свою благодарность.}
%\AdditionalContent{Благодарность}{Авторы благодарны рецензенту за тщательное прочтение статьи и~ценные предложения и комментарии.}

\begin{thebibliography}{10}

\RBibitem{bib1}
\by Седов~Л.~И.
\book Механика сплошной среды
\bookvol 1
\yr 1970
\publ Наука
\publaddr М.
\totalpages 492
%\misctransl
%\by Sedov ~L.~I.
%\book Mekhanika sploshnoy sredy \rm [Continuum Mechanics] vol. 1 
%\yr 1970
%\publ Nauka
%\publaddr Moscow
%\lang In Russian
%\totalpages 492

\RBibitem{bib2}
\by Стружанов~В.~В., Миронов~В.~И.
\book Деформационное разупрочнение материала в элементах конструкций
\yr 1995
\publ УрО РАН 
\publaddr Екатеринбург
\totalpages 190
%\misctransl
%\by Struzhanov ~V.~V., Mironov ~V.~I.
%\book Deformatsionnoye razuprochneniye materiala v elementakh konstruktsiy \rm [Deformational Weakening of the Material in Structural Elements]
%\yr 1995
%\publ UrO RAN
%\publaddr Ekaterinburg
%\lang In Russian
%\totalpages 190

\RBibitem{bib3}
\by Вильдеман~В.~Э., Чаусов~Н.~Г.
\paper Условия деформационного разупрочнения материала при растяжении образца специальной конфигурации
\jour Заводская лаборатория. Диагностика материалов
\yr 2007
\vol 73
\issue 10
\pages 55--59
\elib{https://elibrary.ru/item.asp?id=9907854}
%\misctransl
%\lang In Russian
%\by Vildeman ~V.~E., Chausov ~N.~G.
%\paper Usloviya deformatsionnogo razuprochneniya materiala pri rastyazhenii obraztsa spetsial'noy konfiguratsii \rm [Conditions of Deformational Weakening of the Material with the Specimen of a Special Configuration under Tension]
%\jour Zavodskaya laboratoriya. Diagnostika materialov.
%\yr 2007
%\vol 73
%\issue 10
%\pages 55--59

\RBibitem{bib4}
\by Вильдеман~В.~Э., Третьяков~М.~П.
\paper Испытания материалов с построением полных диаграмм деформирования
\jour Проблемы машиностроения и надежности машин
\yr 2013
\vol 
\issue 2
\pages 93--98
\elib{https://elibrary.ru/item.asp?id=19548776}
%\misctransl
%\lang In Russian
%\by Vildeman ~V.~E., Tretyakov ~M.~P.
%\paper Ispytaniya materialov s postroyeniyem polnykh diagramm deformirovaniya \rm [Material Tests and Development of Complete Deformation Diagrams]
%\jour Problemy mashinostroyeniya i nadezhnosti mashin
%\yr 2013
%\vol 
%\issue 2
%\pages 93--98

\RBibitem{bib5}
\by Стружанов~В.~В., Коркин~А.~В.
\paper Об устойчивости процесса растяжения одной стержневой системы с~разупрочняющимися элементами
\jour Вестник Уральского государственного университета путей сообщения
\yr 2016
\issue 3(31)
\pages 4--17
\elib{https://elibrary.ru/item.asp?id=27249016}
\crossref{10.20291/2079-0392-2016-3-4-17}
%\misctransl
%\lang In Russian
%\by Struzhanov ~V.~V., Korkin ~A.~V.
%\paper Ob ustoychivosti protsessa rastyazheniya odnoy sterzhnevoy sistemy s razuprochnyayushchimisya elementami \rm [On the Stability of a Single Rod System with Weakening Elements under Tension]
%\jour Vestnik Ural'skogo gos. un–ta putey soobshcheniya
%\yr 2016
%\vol 
%\issue 3(31)
%\pages 4--17

\RBibitem{bib6}
\by Андреева~Е.~А.
\paper Решение одномерных задач пластичности для разупрочняющегося материала
\jour Вестн. Сам. гос. техн. ун-та. Сер. Физ.-мат. науки
\yr 2008
\issue 2(17)
\pages 152--160
\mathnet{http://mi.mathnet.ru/vsgtu642}
\crossref{10.14498/vsgtu642}


%\RBibitem{bib7}
%\by Горбунов~С.~В., Радченко~В.~П.
%\paper Вариант решения краевой задачи неупругого деформирования разупрочняющейся среды и его экспериментальная проверка
%\jour Материалы Х Всероссийской конференции по механике деформируемого твёрдого тела (18–22 сентября 2017 г., Самара, Россия)
%\yr 2017
%\vol 1
%\issue
%\pages 179--181
%\misctransl
%\lang In Russian
%\by Gorbunov ~S.~V., Radchenko ~V.~P.
%\paper Variant resheniya krayevoy zadachi neuprugogo deformirovaniya razuprochnyayushcheysya sredy i yego eksperimental'naya proverka \rm [A Solution Case of the Boundary-Value Problem for Inelastic Deformation of a Weakening Medium and its Experimental Verification]
%\jour Materialy X  Vserossiyskoy konferentsii po mekhanike deformiruyemogo tvordogo tela.  
%\yr 2017
%\vol 1
%\issue
%\pages 179--181

\RBibitem{bib8}
\by Арсенин~В.~Я.
\book Методы математической физики и специальные функции
\yr 1974
\publ Наука
\publaddr М.
\totalpages 431
%\misctransl
%\by Arsenin ~V. ~Ya.
%\book Metody matematicheskoy fiziki i spetsial'nyye funktsii. \rm [Methods of Mathematical Physics and Special Functions]
%\yr 1974
%\publ Nauka
%\publaddr Moscow
%\lang In Russian
%\totalpages 431

\RBibitem{bib9}
\by Стружанов~В.~В., Коркин~А.~В.
\paper О решении нелинейных уравнений равновесия  одной растягиваемой стержневой системы с разупрочняющимися элементами методом простых итераций
\jour Вестник Уральского государственного университета путей сообщения
\yr 2017
\issue 2(34)
\pages 4--16
\crossref{10.20291/2079-0392-2017-2-4-16 }
\elib{https://elibrary.ru/item.asp?id=29409142}
%
%\misctransl
%\lang In Russian
%\by Struzhanov ~V.~V., Korkin ~A.~V.
%\paper O reshenii nelineynykh uravneniy ravnovesiya odnoy rastyagivayemoy sterzhnevoy sistemy s razuprochnyayushchimisya elementami metodom prostykh iteratsiy \rm [The solution of Nonlinear Equilibrium Equations for a Single Rod System with Weakening Elements under Tension Using the Simple Iterations Method]
%\jour Vestnik Ural'skogo gos. un-ta putey soobshcheniya
%\yr 2017
%\vol 
%\issue 2(34)
%\pages 4--16

\Bibitem{bib10}
\by Poston~T., Stewart~I.
\book Catastrophe Theory and Its Applications
\yr 1978
\publ Pitman
\serial Surveys and Reference Works in Mathematics
\vol 2
\zmath{0382.58006}
\publaddr London, San Francisco, Melbourne
\totalpages xviii+491 

\Bibitem{bib11}
\by Gilmore~R.
\book Catastrophe theory for scientists and engineers
\zmath{0497.58001}
\serial A Wiley-Interscience Publication
\publaddr New York 
\publ John Wiley \& Sons
\totalpages xvii+666 
\yr 1981

\Bibitem{bib12}
\by Pars~L.~A.
\book A treatise on analytical dynamics
\zmath{0125.12004}
\publaddr London
\publ Heinemann Educational Books 
\totalpages xxi+641 
\yr 1965
 
\RBibitem{bib13}
\by Ильюшин~А.~А.
\paper Об основах общей математической теории пластичности
\jour Вестник Московского университета. Серия 1. Математика. Механика
\yr 1961
\issue 3
\pages 31--36
\zmath{0100.20304}
%\misctransl
%\lang In Russian
%\by Ilyushin ~A.~A.
%\paper Ob osnovakh obshchey matematicheskoy teorii plastichnosti \rm [Fundamentals of the General Mathematical Theory of Plasticity]
%\jour Voprosy teorii plastichnosti. Izd-vo Akademii nauk SSSR
%\yr 1961
%\vol 
%\issue 
%\pages 3--29

\RBibitem{bib14}
\by Качанов~Л.~М.
\book Основы теории пластичности
\yr 1969
\publ Наука
\publaddr М.
\totalpages 420
%\misctransl
%\by Kachanov ~L.~M.
%\book Osnovy teorii plastichnosti \rm [Fundamentals of the Theory of Plasticity]
%\yr 1969
%\publ Nauka
%\publaddr Moscow
%\lang In Russian
%\totalpages 420

\RBibitem{bib15}
\by Стружанов~В.~В.
\paper Определение диаграммы деформирования материала с~падающей ветвью по диаграмме кручения цилиндрического образца
\jour Сиб. журн. индустр. матем.
\yr 2012
\vol 15
\issue 1
\pages 138--144
\mathnet{http://mi.mathnet.ru/sjim718}

\RBibitem{bib16}
\by Радченко~В.~П.
\book Введение в механику деформируемых систем
\yr 2009
\publ СамГТУ
\publaddr Самара
\totalpages 241
%\misctransl
%\by Radchenko ~V.~P.
%\book Vvedeniye v mekhaniku deformiruyemykh system \rm [Introduction to Mechanics of De-formation Systems]
%\yr 2009
%\publ Samarskiy gos. tekhn. un-t.
%\publaddr Samara
%\lang In Russian
%\totalpages 241

\Bibitem{bib17}
\by Southwell~R.~V.
\book An introduction to the theory of elasticity. For engineers and physicists
\zmath{62.1529.04}
\totalpages ix+510 
\publaddr Oxford
\publ Clarendon Press
\serial Oxford Engineering Science Series
\yr 1936


\end{thebibliography}


\newpage

\makeentitle

\vspace{-5mm}

\AdditionalContent{Competing interests}{We have no competing interests.} 
%\AdditionalContent{Competing interests}{I have ... (or) no competing interests.} 
%\AdditionalContent{Competing interests}{We have ... (or) no competing interests.} 
\AdditionalContent{Authors' contributions and responsibilities}{Each authors has participated in the article concept development  and in the manuscript writing. The authors are absolutely responsible for submitting the final manuscript in print. Each authors has approved the final version of manuscript.}
%\AdditionalContent{Author's Responsibilities}{I take full responsibility for submitting the final manuscript in print. I approved the final version of the manuscript.} 
%\AdditionalContent{Authors' contributions and responsibilities}{Each author has participated in the article concept development and in the manuscript writing. The authors are absolutely responsible for submitting the final manuscript in print. Each author has approved the final version of manuscript.} 

\AdditionalContent{Funding}{This research received no specific grant from any funding agency in the public, commercial, or not-for-profit sectors.}




\begin{thebibliography}{10}

\Bibitem{bib1:1}
\lang In Russian
\by Sedov~L.~I.
\book Mekhanika sploshnoi sredy \rm [Continuum Mechanics] 
\bookvol 1
\yr 1970
\publ Nauka
\publaddr Moscow
\totalpages 492

\Bibitem{bib2:1}
\lang In Russian
\by Struzhanov~V.~V., Mironov~V.~I.
\book Deformatsionnoe razuprochnenie materiala v elementakh konstruktsii \rm [Deformational Softening of Material in Structural Elements]
\yr 1995
\publ UrO RAN 
\publaddr Ekaterinburg
\totalpages 190

\Bibitem{bib3}
\lang In Russian 
\by Vil'deman~V.~E., Chausov~N.~G.
\paper Conditions of strain softening upon stretching of the specimen of special configuration
\jour Zavodskaia laboratoriia. Diagnostika materialov
\yr 2007
\vol 73
\issue 10
\pages 55--59


\Bibitem{bib4:en}
\paper Material testing  by plotting total deformation curves
\by Vil'deman~V.~E., Tretyakov~M.~P.
\jour J. Mach. Manuf. Reliab.
\yr  2013
\vol 42
\issue 2
\pages 166--170
\elib{https://elibrary.ru/item.asp?id=20435521}
\crossref{10.3103/S1052618813010159}



\Bibitem{bib5:1}
\lang In Russian
\by Struzhanov~V.~V., Korkin~A.~V.
\paper Regarding stretching process stability of one bar system with softening elements
\jour Vestnik Ural'skogo gosudarstvennogo universiteta putei soobshcheniia
\yr 2016
\issue 3(31)
\pages 4--17
\crossref{10.20291/2079-0392-2016-3-4-17}


\Bibitem{bib6}
\lang In Russian
\by Andreeva~E.~A.
\paper Solution of one-dimensional softening materials plasticity problems
\jour Vestn. Samar. Gos. Tekhn. Univ., Ser. Fiz.-Mat. Nauki \rm [J. Samara State Tech. Univ., Ser. Phys. Math. Sci.]
\yr 2008
\issue 2(17)
\pages 152--160
\crossref{10.14498/vsgtu642}



\Bibitem{bib8:1}
\lang In Russian
\by Arsenin~V.~Ya.
\book Metody matematicheskoi fiziki i spetsial'nye funktsii \rm [Methods of mathematical physics and special functions]
\yr 1974
\publ Nauka
\publaddr Moscow
\totalpages 431


\Bibitem{bib9:1}
\lang In Russian
\by Struzhanov~V.~V., Korkin~A.~V.
\paper On the solution of non-linear equations of trim of a single expandable frame structure with softening elements by the method of simple iterations
\jour Vestnik Ural'skogo gosudarstvennogo universiteta putei soobshcheniia
\yr 2017
\issue 2(34)
\pages 4--16
\crossref{10.20291/2079-0392-2017-2-4-16}


\Bibitem{bib10:1}
\by Poston~T., Stewart~I.
\book Catastrophe Theory and Its Applications
\yr 1978
\publ Pitman
\serial Surveys and Reference Works in Mathematics
\vol 2
\zmath{0382.58006}
\publaddr London, San Francisco, Melbourne
\totalpages xviii+491 

\Bibitem{bib11:1}
\by Gilmore~R.
\book Catastrophe theory for scientists and engineers
\zmath{0497.58001}
\serial A Wiley-Interscience Publication
\publaddr New York 
\publ John Wiley \& Sons
\totalpages xvii+666 
\yr 1981

\Bibitem{bib12:1}
\by Pars~L.~A.
\book A treatise on analytical dynamics
\zmath{0125.12004}
\publaddr London
\publ Heinemann Educational Books 
\totalpages xxi+641 
\yr 1965
 
\Bibitem{bib13:1}
\lang In Russian
\by Ilyushin~A.~A.
\paper On the foundations of the general mathematical theory of plasticity
\jour Vestn. Mosk. Univ., Ser.~I 
\yr 1961
\issue 3
\pages 31--36



\Bibitem{bib14:1}
\by Kachanov~L.~M.
\book Foundations of the theory of plasticity
\zmath{0231.73015}
\serial North-Holland Series in Applied Mathematics and Mechanics
\vol 12
\publaddr Amsterdam
\publ North-Holland Publ.
\totalpages xiii+482 
\yr 1971


\Bibitem{bib15:!}
\by Struzhanov~V.~V.
\paper The determination of the deformation diagram of a~material with a~falling branch using the torsion diagram of a~cylindrical sample
\jour Sib. Zh. Ind. Mat.
\yr 2012
\vol 15
\issue 1
\pages 138--144

\Bibitem{bib16:1}
\lang In Russian
\by Radchenko~V.~P.
\book Vvedenie v mekhaniku deformiruemykh sistem \rm [Introduction to the mechanics of deformable systems]
\yr 2009
\publ Samara State Technical Univ.
\publaddr Samara
\totalpages 241


\Bibitem{bib17:1}
\by Southwell~R.~V.
\book An introduction to the theory of elasticity. For engineers and physicists
\zmath{62.1529.04}
\totalpages ix+510 
\publaddr Oxford
\publ Clarendon Press
\serial Oxford Engineering Science Series
\yr 1936


\end{thebibliography}


%\end{document}
 \newpage

%\input{preamble}

%

\vsgtu{1639} %registration number

\FirstPage{774}

\LastPage{784}

%

%\usepackage{samgtu-name}

%

%\pdfcrop

%\draft

%

\ShortCommunication

%

%\begin{document}





$\;$

\vspace{1mm}



\hfill \qrcode[hyperlink,height=.8in]{https://doi.org/10.14498/vsgtu1639}

\vspace{-.8in}



\udc{517.955, 517.968.7}

\msc{26A33, 35K15, 35R11}



\rutitle{К проблеме единственности решения задачи Коши для~уравнения дробной диффузии с оператором Бесселя}

\rutitleshort{К проблеме единственности решения задачи Коши\dots}

\rutitlehack{К проблеме единственности решения \\ задачи Коши для~уравнения дробной диффузии с~оператором Бесселя}



\entitle{On the uniqueness of the solution of the Cauchy problem for the equation of fractional diffusion with~Bessel~operator}

\entitleshort{On the uniqueness of the solution of the Cauchy problem\dots}



\ruauthor{Ф.~Г.~Хуштова}{Х\,у\,ш\,т\,о\,в\,а~Ф.~Г.}

\enauthor{F.~G.~Khushtova}{K\,h\,u\,s\,h\,t\,o\,v\,a~F.~G.}



\ruaffil{\runiipma.}



\enaffil{\enniipma.}



\ruauthorsdetails{1}{% Количество авторов

\fullauthor{\rupid{53181}{Фатима Гидовна Хуштова}}

\CAuthor % Автор-корреспондент

\orcid{0000-0003-4088-3621} % ORCID ID автора 

%\regalia{без ученой степени} 

\position[n]{научный сотрудник} 

\dept{отдел дробного исчисления} 

\email{khushtova@yandex.ru}

}



\enauthorsdetails{1}{

\fullauthor{\enpid{53181}{Fatima G. Khushtova}}

\CAuthor % Сorresponding Author

\orcid{0000-0003-4088-3621} % ORCID ID

%\regalia{Cand. Phys. \& Math. Sci., Associate Professor} 

\position[n]{Researcher} 

\dept{Dept. of Fractional Calculus}

\email{khushtova@yandex.ru}

}



\rukeywords{уравнение дробной диффузии, оператор дробного дифференцирования, оператор Бесселя, задача Коши, единственность решения, условие Тихонова, функция Райта}

\enkeywords{fractional diffusion equation, fractional differentiation operator, Bessel operator, Cauchy problem, solution uniqueness, Tikhonov condition,	 Wright function}



\ruabstract{Рассматривается уравнение дробной диффузии с сингулярным оператором Бесселя, действующим по пространственной переменной, и оператором дробного дифференцирования Римана--Лиувилля, действующим по временной переменной. 

Когда порядок дробной производной равен единице, а особенность у оператора Бесселя отсутствует, рассматриваемое уравнение совпадает с~классическим уравнением теплопроводности. 

Ранее для уравнения дробной диффузии с оператором Бесселя было построено решение задачи Коши и доказана теорема единственности решения в классе функций экспоненциального роста. 



Построен пример, показывающий, что увеличение показателя степени в условии, гарантирующем единственность решения задачи Коши, влечет за собой неединственность решения. 

С помощью известных свойств функции Райта получены оценки для построенной функции.

Показывается, что она, будучи не равной тождественно нулю, удовлетворяет однородному уравнению и однородному условию Коши. 

}





\enabstract{In this paper, we consider fractional diffusion equation involving the Bessel operator acting with respect to a spatial variable and the Riemann-Liouville fractional differentiation operator acting with respect to a time variable. When the order of the fractional derivative is unity, and the singularity of the Bessel operator is absent, this equation coincides with the classical heat equation. Earlier, a solution of the Cauchy problem has been considered for the considered equation and a uniqueness theorem has been proved for a class of functions satisfying the analog of the Tikhonov condition. 



In this paper, we have constructed an example to show that the exponent (power) at the condition of the uniqueness of the solution to the Cauchy problem cannot be raised under. Its increase leads to a non-uniqueness of the solution. Using the well-known properties of the Wright function, we have obtained estimates for constructed function, which satisfies the homogeneous equation and the zero Cauchy condition.

}



\newdate{khusha}{28}{08}{2018}

\date{\displaydate{khusha}}

\newdate{khushb}{25}{10}{2018}

\revisiondate{\displaydate{khushb}}

\newdate{khushc}{12}{11}{2018}

\accepteddate{\displaydate{khushc}}



\newdate{khushe}{28}{11}{2018}

\renewcommand{\JournalDataOnline}{\displaydate{khushe}}



%\ForPeerReview

\makerutitle



\newpage



\Section[n]{Введение}

В области $\Omega=\{(x,y): -\infty<x<\infty,\ 0<y<T \}$ рассмотрим уравнение

\begin{equation}\label{GL3_1}

B_xu(x,y)-D^{\alpha}_{0y}u(x,y)=0, 

\end{equation}

где $D_{0y}^{\alpha}$ "--- оператор дробного дифференцирования в смысле  

Римана"--~Лиувилля  порядка $\alpha,$ определяемый равенствами~\cite{Drob, Pskhu}   

$$D_{0y}^{\alpha} \, g(y)= \frac{dg}{dy},$$ если $\alpha=1$, и

$$ 

D_{0y}^{\alpha}\,g(y)=\frac{1}{\Gamma(1-\alpha)}\frac{d}{dy}

\int _{0}^{y}\frac{g(t)}{(y-t)^{\alpha}}dt,

$$ 

если $0<\alpha<1;$

$$

B_x=x^{-b} \frac{d}{dx}\left(x^b \frac{d}{dx}\right)=\frac{d^2}{d x^2}+\frac{b}{x} \frac{d}{d x}

$$

"--- оператор Бесселя, $|b|<1.$



Уравнение \eqref{GL3_1} при $\alpha=1,$ то есть уравнение

\begin{equation}\label{UrAlpha-1}

u_{xx}(x,y)+\frac{b}{x}u_{x}(x,y)-u_{y}(x,y)=0,

\end{equation}

совпадает с %~названным И.~А.~Киприяновым 

$B$-параболическим уравнением.

Краевые задачи в ограниченной и~неограниченной областях для этого уравнения при различных значениях параметра $b$ рассматривали многие авторы. 

Например, в работе V.~Alexia\-des~\cite{Alexiades} для него решены первая, вторая и~третья краевые задачи в~области с~подвижной границей, в~работе D.~Calton\cite{Calton} для него  исследована задача Коши.

 

 Уравнение \eqref{UrAlpha-1} заменой $z=x^2/(4n)$ сводится к~уравнению

\begin{equation} \label{UrKep2}

u_{zz}(z,y)+\frac{m+1}{z}\,u_{z}(z,y)-\frac{n}{z} u_{y}(z,y)=0,\quad m=(b-1)/2,

\end{equation}

для которого S.~K\c{e}pi\'nski  \cite{Kepinski_2} исследовал краевую задачу в~полуполосе $x>0,$ $0<y<y_0.$

Уравнение \eqref{UrAlpha-1} было также объектом исследования  монографий С.~А.~Терсенова~\cite{Tersenov} и~М.~И.~Матийчука~\cite{Matiychuk}.



Дифференциальные уравнения, содержащие оператор Бесселя, наиболее подробно исследованы в~работах И.~А.~Киприянова и~его учеников \cite{Kipriyanov-Monographiya,  Kipriyanov}.

Отметим здесь, что очень важные результаты по применению операторов преобразования в~теории дифференциальных уравнений с~особенностями в~коэффициентах, в~частности с~операторами Бесселя, получены В.~В.~Катраховым и~С.~М.~Ситником~\cite{KatrakhovSitnik, Sitnik}.





В случае, когда $b=0,$ $0<\alpha<2,$ уравнение \eqref{GL3_1} совпадает с диффузионно"=волновым уравнением. 

Различные краевые задачи для него, а также для многомерных его обобщений подробно исследованы в работах многих авторов.  Приведем некоторые из них. 



А.~Н.~Кочубей \cite{Kochubei_2} исследовал задачу Коши для уравнения диффузии дробного порядка с~регуляризованной дробной производной (производной Капуто) и~эллиптическим оператором с~непрерывными ограниченными вещественными коэффициентами

$$

L=\sum\limits_{i,j=1}^{n}a_{i,j}(x)\frac{\partial^2}{\partial x_i \partial x_j}+

\sum\limits_{j=1}^{n}b_{j}(x)\frac{\partial}{\partial x_j}+c(x) u.

$$

В случае, когда $L=\Delta$ "--- оператор Лапласа, фундаментальное решение построено в терминах $H$-функции Фокса. 

Исследованию задачи Коши для диффузионного и диффузионно"=волнового уравнений дробного порядка с~производной Капуто  посвящены также  работы А.~Н.~Кочубея и С.~Д.~Эйдельмана \cite{Kochubei_1, Kochubei_3, Kochubei_4, Kochubei_5}.



А.~В.~Псху~\cite{Pskhu2} построил фундаментальное решение и исследовал задачу Коши  для многомерного диффузионно"=волнового уравнения с~оператором дробного дифференцирования Джрбашяна"--~Нерсесяна, который  включает в~себя как частный случай операторы дробного дифференцирования Римана"--~Лиувилля и Капуто.



В~работах  \cite{Pskhu-SibirJournal, Pskhu3} также построено фундаментальное решение и исследована задача Коши для уравнения дробной диффузии соответственно с~оператором дискретно распределенного дифференцирования и оператором Джрбашяна"--~Нерсесяна, действующим по многим переменным времени.

В работе \cite{Pskhu4} решена первая краевая задача в~нецилиндрической области для диффузионно"=волнового уравнения с~оператором Джрбашяна"--~Нерсесяна. 

Доказана теорема существования и~единственности исследуемой задачи, конструктивно построено ее решение. 



А.~А.~Ворошилов и А.~А.~Килбас \cite{VorKilbas} методом интегральных преобразований построили решение задачи Коши для уравнения

\begin{equation} \label{MnogPerem}

D^{\alpha}_{0t}\,u(x,t)=\lambda^2 \Delta_x u(x,t), \quad  x\in{\mathbb R}^m,\quad  t>0, 

\end{equation}

где $\Delta_x=\sum_{j=1}^{m}\partial^2/\partial x_{j}^{2},$ $n-1<\alpha <n,$ $n\in\mathbb{N}.$ 

При $0<\alpha\leqslant 1$ и $1<\alpha<2$ решения  выписаны в терминах $H$-функции Фокса.

В работе~\cite{VorKilbas2} было также  построено решение задачи Коши в~случае, когда в уравнении

\eqref{MnogPerem} вместо оператора Римана"--~Лиувилля содержится оператор Капуто. 



Задача Коши для уравнения \eqref{MnogPerem} при $\lambda=1,$ 

$0<\alpha\leqslant 1$  была ранее рассмотрена С.~Х.~Геккиевой~\cite{Gekk}, а~в~работе \cite{Gekkieva1} ею исследована первая краевая задача в~первом квадранте для уравнения \eqref{MnogPerem} при $m=\lambda=1,$ 

$0<\alpha\leqslant 1.$  







Уравнение диффузии дробного порядка и некоторые его обобщения рассматривали 

F.~Mainardi, G.~Pagnini, Yu.~Luchko, P.~Paradisi~\cite{Mainardi-1, Mainardi-3, Mainardi-2, Pagnini_1,   Pagnini_Paradisi}  и~многие другие. 

Более подробную библиографию по данному вопросу можно найти в~монографиях \cite{Pskhu, MathaiSaxenaHaubold, Podlubny}.

 

%Чтобы подготовить Вашу статью для рецензента, т.~е. обезличить её, воспользуйтесь командой \verb"\ForPeerReview" перед командой 

%\verb"\makerutitle".



\smallskip





\Section{Постановка задачи Коши и теорема единственности решения}

Пусть $\Omega^{+}=\Omega\cap\{x>0 \},$ $\Omega^{-}=\Omega\cap\{x<0\},$ 

$\bar{\Omega}$ "--- замыкание области $\Omega.$ 



\smallskip



{\sc\small Определение.}

{\it Регулярным  решением} уравнения \eqref{GL3_1} в области $\Omega$ 

будем называть функцию

$u=u(x,y),$ удовлетворяющую уравнению \eqref{GL3_1} в области

$\Omega^{+}\cup\Omega^{-},$ и такую, что $y^{\,1-\alpha}u\in{C(\bar{\Omega})},$ 

$|x|^b u_{x}\in{C(\Omega)},$ $u_{xx},$ $D_{\,0y}^{\,\alpha} u\in{C(\Omega^{+}\cup\Omega^{-})}.$



\smallskip



{\sc\small Задача Коши.}

{\it Найти регулярное в области $\Omega$ решение

уравнения} \eqref{GL3_1}, {\it удовлетворяющее условию

\begin{equation}\label{UslSMCoshi}

\lim\limits_{y \rightarrow 0} D_{0y}^{\alpha-1}u(x,y)=\varphi(x), \quad  -\infty<x<\infty,

\end{equation}

где $\varphi(x)$ "--- заданная функция}.



\smallskip



В работе \cite{KhushtovaVestn2} доказана следующая теорема единственности.



\smallskip



\begin{newthm}{Теорема}  

Существует не более одного регулярного решения задачи Коши {\rm \eqref{GL3_1}, \eqref{UslSMCoshi} } в классе функций$,$ удовлетворяющих для некоторой положительной постоянной $k$ условию

\begin{equation} \label{AnalogTikhonovaInfty}

\lim\limits_{|x|\to \infty}y^{1-\alpha}\,u(x,y)\exp{\big(-k\,|x|^{\,\frac{2}{2-\alpha}} \big)}=0.

\end{equation}  

\end{newthm}



\smallskip





Отметим, что условие \eqref{AnalogTikhonovaInfty} имеет место и в случае задачи Коши для уравнения дробной диффузии.  

Теоремы единственности решения задачи Коши в~классе ограниченных функций  и~в~классе функций экспоненциального роста для многомерного уравнения дробной диффузии с~оператором Капуто были доказаны в работе А.~Н.~Кочубея \cite{Kochubei_2},

теорема единственности решения задачи Коши в классе функций экспоненциального роста для многомерного диффузионно"=волнового уравнения с~оператором Джрбашяна"--~Нерсесяна  "--- в~работе А.~В.~Псху~\cite{Pskhu2}. 

В этих же работах построены примеры, показывающие, что увеличение показателя $2/(2-\alpha)$ ведет к~утрате единственности решения соответствующих задач. 

При $\alpha=1$ условие \eqref{AnalogTikhonovaInfty} совпадает с~хорошо известным условием Тихонова 

$$

\lim\limits_{|x|\to \infty}u(x,y)\exp{\big(-k\,|x|^{2} \big)}=0.

$$

Как известно, пример неединственности решения задачи Коши для уравнения теплопроводности был впервые построен А.~Н.~Тихоновым в~основополагающей работе~\cite{tikhonovUsl-1} (см.~также \cite{tikhonovUsl-2}).



В данной работе эти идеи развиваются применительно к~задаче \eqref{GL3_1}, \eqref{UslSMCoshi}.



\smallskip



\Section{Неулучшаемость показателя степени в условии единственности решения задачи Коши}

Рассмотрим функцию

\begin{equation} \label{T(z,zeta)}  

T_{\mu}(z,\zeta)=\sum\limits_{n=0}^{\infty}\frac{z^{\,\beta+n}}{\,n!\,\Gamma(1+\beta+n)}\,\phi\left(-\rho,\mu-\alpha\beta-\alpha n; -\zeta\right),

\end{equation}    

где  $\phi\,(-\rho,\delta;z)$ "--- {\it функция Райта}, определяемая степенным рядом \cite{Pskhu, GorenfloLuchkoMainardi}

$$

\phi\,(-\rho,\delta;z)=\sum\limits_{n=0}^{\infty}\frac{z^{n}}{\,n!\,\Gamma(\delta-\rho\, n)}.

$$



Далее будем считать, что $\rho\in(0;1).$ 

Для любых $\nu$, $\delta\in \mathbb{R}$ справедливы формулы  \cite[с.~26, 44]{Pskhu}

\begin{gather}\label{diff_x}

\frac{d}{dz}\phi(-\rho,\delta;-z)=-\phi(-\rho,\delta-\rho;-z),

\\

\label{diff_Right}

D_{ay}^{\nu} \,y^{\delta-1}\phi(-\rho,\delta;-c\,y^{-\rho})=y^{\delta-\nu-1}\phi(-\rho,\delta-\nu;-c\,y^{-\rho}).

\end{gather}



\smallskip



\begin{lemma}[1]Если $\delta\geqslant 1,$ то для любого положительного $z$ справедливо нера\-вен\-ство~{\rm \cite[с.~47]{Pskhu} }

\begin{equation}\label{Right_delta>1}

\phi\,(-\rho,\delta;-z)\leqslant \frac{1}{\Gamma\left( \delta \right)}

\exp\left[-(1-\rho)\,\rho^{\frac{\rho}{1-\rho}}\, z^{\frac{1}{1-\rho}} \right].

\end{equation}

\end{lemma}



\smallskip



\begin{lemma}[2]Если $\delta<1,$ то для любых положительных $x$ и $y,$ 

$a\in(0;x),$ $\omega\in(1/2; \omega_0),$ $\omega_0=\min\{1,1/(2\rho)\},$

справедливы неравенства~{\rm \cite[с.~49]{Pskhu}}

\begin{gather}\label{Right-0}

\Bigl|\,y^{\,\delta-1}\phi\Bigl(-\rho,\delta;-\frac{x}{y^{\rho}}\Bigr)\Bigr|

\leqslant 

\frac{\Gamma\left( (1-\delta)/\rho\right)}{\,\rho\, \pi \left[aC_{\rho}(\rho,\omega) \right]^{(1-\delta)/\rho}}\,

\phi\Bigl(-\rho,1;-\frac{x-a}{y^{\rho}}\Bigr),

\\

\label{Right-2}

\Bigl|y^{\,\delta-1}\phi\Bigl(-\rho,\delta;-\frac{x}{y^{\rho}}\Bigr)\Bigr|

\leqslant \frac{\Gamma\left( 1-\delta\right)}{\pi \left[yC_{\rho}(1,\omega) \right]^{1-\delta}},

\end{gather}

где 

$C_{\rho}(\rho,\omega)=\frac1{\rho\pi}{\cos \rho\omega\pi},$ $C_{\rho}(1,\omega)=-\frac1{\pi}{\cos \omega\pi}.

$

\end{lemma}



\smallskip



Как следует из \eqref{Right-2}, ряд \eqref{T(z,zeta)} сходится абсолютно для любых  $z\in\mathbb{C},$ $\zeta>0,$ $\rho, \alpha, \beta\in(0;1),$ $\mu\in\mathbb{R}.$



\smallskip



Докажем следующую лемму.



\smallskip

\begin{lemma}[3]Если $0<\alpha<\rho<1,$ $\mu\geqslant 0,$ то при любых $x\in\mathbb{R},$ $y>0$ справедливо неравенство

\begin{equation} \label{T_muOtsenka}  

\Bigl| y^{\mu-1} T_{\mu}\Bigl(\frac{x^2}{4y^{\alpha}},\frac{1}{y^{\rho}}\Bigr)\Bigr|

\leqslant 

\frac{C

|x|^{\frac{2[1-\rho(1-\beta)]}{2\rho-\alpha}} y^{\mu}}{\Gamma(\mu+1)}\exp\left[ k_1  |x|^{\frac{2\rho}{2\rho-\alpha}}-k_2   y^{-\frac{\rho}{1-\rho}}\right],

\end{equation}    

где $C,$ $k_1,$ $k_2$ "--- положительные постоянные$,$ зависящие только от 

$\alpha,$ $\beta$ и $\rho.$

\end{lemma}



\smallskip



\begin{proof}

Из \eqref{Right-0} следует оценка

$$

\Bigl| y^{-\alpha\beta-\alpha n-1} \phi\Bigl(-\rho,-\alpha\beta-\alpha n; -\frac{1}{y^{\rho}}\Bigr)\Bigr|

\leqslant

 \frac{\Gamma\left( 1/\rho+\varepsilon\beta+\varepsilon n\right)}{ \rho \pi\left(aC_{\rho} \right)^{1/\rho+\varepsilon\beta+\varepsilon n}}

\phi\Bigl(-\rho,1; -\frac{1-a}{y^{\rho}}\Bigr),

$$

где $\varepsilon=\alpha/\rho,$ $a\in(0;1),$ $C_{\rho}=C_{\rho}(\rho,\omega).$



С помощью последней оценки получим

\begin{multline}

\Bigl| y^{\,\mu-1} T_{\mu}\Bigl(\frac{x^2}{4y^{\alpha}},\frac{1}{y^{\rho}}\Bigr)\Bigr|=

\\

=\left|

 D_{0y}^{-\mu}

  \sum\limits_{n=0}^{\infty}\frac{|x|^{2\beta+2n}}{ 2^{2\beta+2n}n! \Gamma(1+\beta+n)}

y^{-\alpha\beta-\alpha n-1}\phi \Bigl(-\rho,-\alpha\beta-\alpha n; -\frac{1}{y^{\rho}}\Bigr)

 \right|\leqslant

\\

 \label{WrightRyadZ}  

\leqslant

\frac{ |x|^{2\beta}y^{\mu}\phi\bigl(-\rho,\mu+1; -\frac{1-a}{y^{\rho}}\bigr)}{2^{2\beta}\rho \pi \left(aC_{\rho} \right)^{1/\rho+\varepsilon\beta}}

\sum\limits_{n=0}^{\infty}

\frac{\Gamma\left(1/\rho+\varepsilon\beta+\varepsilon n\right)}{n!

\Gamma(1+\beta+n)}

\Bigl( \frac{x^2}{4\left(aC_{\rho} \right)^{\varepsilon}}\Bigr)^{n}.

\end{multline}



Степенной ряд, стоящий справа от неравенства \eqref{WrightRyadZ}, представляет собой 

обобщенную функцию Райта

$$

{_1\Psi_{1}}\left[ \,z \,\bigg|

\begin{array}{ll}

 \big(1/\rho+\varepsilon\beta, \,\varepsilon\,\big)\\

\big(\,1+\beta, 1\,\big)\\

\end{array}

\right]=

\sum\limits_{n=0}^{\infty}

\frac{\Gamma\left(1/\rho+\varepsilon\beta+\varepsilon n\right)}{n!

\,\Gamma(1+\beta+n)} z^{n}, \; \; z=\frac{x^{\,2}}{4(aC_{\rho})^{\varepsilon}}.

$$

Из ее асимптотических свойств при $z\to \infty$ следует~\cite{Wright}

$$

{_1\Psi_{1}}\left[ \,z\,\bigg|

\begin{array}{ll}

 \big(1/\rho+\varepsilon\beta, \,\varepsilon\,\big)\\

\big(\,1+\beta, 1\,\big)\\

\end{array}

\right]=Z^{\,1/\rho-\beta(1-\varepsilon)-1}\,e^{Z}

\left[ \sum\limits_{m=0}^{M-1}A_m Z^{-m}+O\left(Z^{-M}\right) \right],

$$

где

$Z=(2-\varepsilon)\,\varepsilon^{\,\frac{\varepsilon}{2-\varepsilon}}z^{\,\frac{1}{2-\varepsilon}},$  $\varepsilon<2,$ а коэффициенты $A_m,$ $m=0, 1, 2, ...,$ зависят только от $\rho,$ $\varepsilon$ и $\beta.$



Учитывая, что $a$ "--- произвольное число из интервала $(0;1),$ из последней оценки и~соотношений \eqref{WrightRyadZ}, \eqref{Right_delta>1} получаем \eqref{T_muOtsenka}. \hfill \end{proof}



\smallskip



Покажем, что функция $y^{\mu-1} T_{\mu}\bigl(\frac{x^2}{4y^{\alpha}},\frac{1}{y^{\rho}}\bigr)$ является решением уравнения \eqref{GL3_1}. 

Из \eqref{T(z,zeta)}, \eqref{diff_x}  и \eqref{diff_Right} имеем

\begin{multline}

B_x T_{\mu}\Bigl(\frac{x^2}{4y^{\alpha}},\frac{1}{y^{\rho}}\Bigr)=

\\

=y^{-\alpha}\sum\limits_{n=1}^{\infty}\frac{1}{(n-1)! \Gamma(\beta+n)}

\Bigl(\frac{x^2}{4y^{\alpha}}\Bigr)^{\beta+n-1}

\phi\Bigl(-\rho,\mu-\alpha\beta-\alpha n; -\frac{1}{y^{\rho}}\Bigr)=

\\

=y^{-\alpha}\sum\limits_{k=0}^{\infty}\frac{1}{k! \Gamma(1+\beta+k)}

\Bigl(\frac{x^2}{4y^{\alpha}}\Bigr)^{\beta+k}

\phi\Bigl(-\rho,\mu-\alpha-\alpha\beta-\alpha k; -\frac{1}{y^{\rho}}\Bigr)=

\\ 

=y^{-\alpha} T_{\mu-\alpha}\Bigl(\frac{x^2}{\,4y^{\alpha}},\frac{1}{y^{\rho}}\Bigr),

\end{multline}   

\begin{equation}

\label{DT(z,zeta)}  

D_{0y}^{\nu} y^{\mu-1} T_{\mu}\Bigl(\frac{x^2}{4y^{\alpha}},\frac{1}{y^{\rho}}\Bigr)=

y^{\mu-\nu-1} T_{\mu-\nu}\Bigl(\frac{x^2}{4y^{\alpha}},\frac{1}{y^{\rho}}\Bigr).

\end{equation}



Из последних двух равенств следует

\begin{equation} \label{BDyT}  

\left(B_x-D_{0y}^{\alpha}\right)y^{\mu-1} T_{\mu}\Bigl(\frac{x^2}{4y^{\alpha}}, \frac{1}{y^{\rho}}\Bigr)=0.

\end{equation}   





Из равенства \eqref{DT(z,zeta)} и оценки \eqref{T_muOtsenka} также имеем

\begin{equation} \label{D(alpha-1)T}  

\lim\limits_{y \rightarrow 0} D_{0y}^{\alpha-1} y^{\mu-1} T_{\mu}\Bigl(\frac{x^2}{4y^{\alpha}}, \frac{1}{y^{\rho}}\Bigr)=0.

\end{equation}   



Таким образом, из соотношений  \eqref{T_muOtsenka}, \eqref{BDyT} и \eqref{D(alpha-1)T} следует, что нетривиальная функция 

$$

u(x,y)=y^{\mu-1} T_{\mu}\Bigl(\frac{x^2}{4y^{\alpha}},\frac{1}{y^{\rho}}\Bigr)

$$

является решением однородного уравнения \eqref{GL3_1}, удовлетворяет однородному условию \eqref{UslSMCoshi} $(\varphi(x)\equiv 0)$ и для любого положительного сколь угодно малого $\delta,$ 

некоторых положительных постоянных $k$ и $C,$ а также параметра 

$\rho,$  выбранного из условия 

$\rho=\frac{\alpha}{2}\frac{2+\delta(2-\alpha)}{\alpha+\delta(2-\alpha)},$

имеет место оценка

$$

\left| u(x,y) \right|\leqslant C \exp\left(k |x|^{\delta+\frac{2}{2-\alpha}} \right).

$$



Последняя оценка показывает, что увеличение показателя  $2/(2-\alpha)$ в~условии \eqref{AnalogTikhonovaInfty} ведет к~неединственности решения задачи  \eqref{GL3_1},  \eqref{UslSMCoshi} и~в~этом смысле условие  \eqref{AnalogTikhonovaInfty} является необходимым.





\smallskip









\Section[N]{Заключение}

В данной работе показывается, что показатель степени в~условии, обеспечивающем единственность решения задачи Коши для уравнения дробной диффузии с оператором Бесселя, является предельным, и его увеличение приводит к нарушению единственности. 



\smallskip





\AdditionalContent{Конкурирующие интересы}{Конкурирующих интересов не имею.} 



\AdditionalContent{Авторская ответственность}{Я несу полную ответственность за предоставление окончательной версии рукописи в печать. Окончательная версия рукописи мною одобрена.} 





\AdditionalContent{Финансирование}{Исследование выполнялось без финансирования.}



\smallskip





\begin{thebibliography}{99}



\RBibitem{Drob}

\by Нахушев~А.~М.

\book Дробное исчисление и его применение

\yr 2003

\publ Физматлит

\publaddr М.

\totalpages 271

%\sortdate 1/1/2003

\zmath{https://zbmath.org/?q=an:1066.26005}

%\misctransl

%\by Nakhushev~A.~M.

%\book Drobnoe ischislenie i ego primenenie \rm [Fractional calculus and its applications]

%\yr 2003

%\publ Fizmatlit

%\publaddr Moscow

%\lang In Russian

%\totalpages 271

%%\sortdate 1/1/2003



\RBibitem{Pskhu}

\by Псху~А.~В.

\book Уравнения в частных производных дробного порядка

\yr 2005

\publ Наука

\publaddr М.

\totalpages 199

%\sortdate 1/1/2005

\zmath{https://zbmath.org/?q=an:1193.35245}

%\misctransl

%\by Pskhu~A.~V.

%\book Uravneniia v chastnykh proizvodnykh drobnogo poriadka \rm [Partial differential equations of fractional order]

%\yr 2005

%\publ Nauka

%\publaddr Moscow

%\lang In Russian

%\totalpages 199

%%\sortdate 1/1/2005



\Bibitem{Alexiades}

\by Alexiades~V.

\paper Generalized axially symmetric heat potentials and singular parabolic initial boundary value problems

\jour Arch. Rational Mech. Anal.

\yr 1982

\vol 79

\issue 4

\pages 325--350

%\citnumscopus 3

%\sortdate 1/1/1982

\crossref{10.1007/BF00250797}

\mathscinet{http://www.ams.org/mathscinet-getitem?mr=656798}

\zmath{https://zbmath.org/?q=an:0511.35050}

\scopus{http://www.scopus.com/record/display.url?origin=inward&eid=2-s2.0-0020402046}



\Bibitem{Calton}

\by Calton~D.

\paper Cauchy's problem for a singular parabolic partial differential equation

\jour J.~Diff. Equations

\yr 1970

\vol 8

\issue 2

\pages 250--257

%\citnumscopus 8

%\sortdate 1/1/1970

\crossref{10.1016/0022-0396(70)90004-5}

\mathscinet{http://www.ams.org/mathscinet-getitem?mr=271547}

\scopus{http://www.scopus.com/record/display.url?origin=inward&eid=2-s2.0-49849107701}



\Bibitem{Kepinski_2}

\by K\c{e}pi\'nski~S.

\paper \"Uber die Differentialgleichung $\frac{\partial^2 z }{\partial x^2}+\frac{m+1}{x}\frac{\partial z }{\partial x} -\frac{n}{x}\frac{\partial z }{\partial t}=0$

\jour Math. Ann.

\yr 1905

\vol 61

\issue 3

\pages 397--405

\lang In German

%\sortdate 1/1/1905

\mathscinet{http://www.ams.org/mathscinet-getitem?mr=1511354}

\zmath{https://zbmath.org/?q=an:36.0416.01}



\RBibitem{Tersenov}

\by Терсенов~С.~А.

\book Параболические уравнения с~меняющимся направлением времени

\yr 1985

\publ Наука

\publaddr М.

\totalpages 105

%\sortdate 1/1/1985

%\misctransl

%\by Tersenov~S.~A.

%\book Parabolicheskie uravneniia s~meniaiushchimsia napravleniem vremeni \rm [Parabolic Equations with a Changing Time Direction]

%\yr 1985

%\publ Nauka

%\publaddr Moscow

%\totalpages 105

%%\sortdate 1/1/1985



\RBibitem{Matiychuk}

\by Матiйчук~M.~I.

\book Параболiчнi сингулярнi крайовi задачi

\yr 1999

\publ Iн-т математики НАН Укра\"\iни

\publaddr Ки\"\iв

\totalpages 176

%%\sortdate 1/1/1999

%\misctransl

%\by Matiichuk~M.~I.

%\book Parabolichni singuliarni kraiovi zadachi \rm [Parabolic Singular Boundary-Value Problems]

%\yr 1999

%\publ Institute of Mathematics of the Ukrainian National Academy of Sciences

%\publaddr Kiev

%\lang In Ukrainian

%\totalpages 176

%%\sortdate 1/1/1999



\RBibitem{Kipriyanov-Monographiya}

\by Киприянов~И.~А.

\book Сингулярные эллиптические краевые задачи

\yr 1997

\publ Наука

\publaddr М.

\totalpages 208

%\sortdate 1/1/1997

%\misctransl

%\by Kipriyanov~I.~A.

%\book Singuliarnye ellipticheskie kraevye zadachi \rm [Singular Elliptic Boundary-Value Problems]

%\yr 1997

%\publ Nauka

%\publaddr Moscow

%\lang In Russian

%\totalpages 208

%%\sortdate 1/1/1997



\RBibitem{Kipriyanov}

\by Киприянов~И.~А., Катрахов~В.~В., Ляпин~В.~М.

\paper О краевых задачах в областях общего вида для сингулярных параболических систем уравнений

\jour Докл. АН СССР

\yr 1976

\vol 230

\issue 6

\pages 1271--1274

%\sortdate 1/1/1976

\mathnet{http://mi.mathnet.ru/rus/dan40691}

\mathscinet{http://www.ams.org/mathscinet-getitem?mr=0425366}

\zmath{https://zbmath.org/?q=an:0354.35056}

%\misctransl

%\by Kipriyanov~I.~A., Katrakhov~V.~V., Lyapin~V.~M.

%\paper On boundary value problems in domains of general type for singular parabolic systems of equations

%\jour Sov. Math., Dokl.

%\yr 1976

%\vol 17

%\pages 1461--1464

%%\sortdate 1/1/1976

%\mathscinet{http://www.ams.org/mathscinet-getitem?mr=0425366}

%\zmath{https://zbmath.org/?q=an:0354.35056}



\RBibitem{Sitnik}

\by Ситник~С.~М.

\yr 2016

\publaddr Воронеж

\thesis Применение операторов преобразования Бушмана--Эрдейи и их обобщений в~теории дифференциальных уравнений с~особенностями в коэффициентах

\thesisinfo Дис.~\dots\ д-ра физ.-мат. наук: 01.01.02

\totalpages 307

%\sortdate 1/1/2016

%\misctransl

%\by Sitnik~S.~M.

%\yr 2016

%\publaddr Voronezh

%\thesis Application of the Bushman--Erd\'elyi transformation operators and their generalizations in the theory of differential equations with singularities in the coefficients

%\thesisinfo Dr. Phys. Math. Sci. Thesis

%\lang In Russian

%\totalpages 307

%%\sortdate 1/1/2016



\RBibitem{KatrakhovSitnik}

\by Катрахов~В.~В., Ситник~С.~М.

\paper Метод операторов преобразования и краевые задачи для сингулярных эллиптических уравнений

\inbook Сингулярные дифференциальные уравнения

\serial СМФН

\yr 2018

\vol 64

\issue 2

\pages 211--426

\publ Российский университет дружбы народов

\publaddr М.

%\sortdate 1/1/2018

\mathnet{http://mi.mathnet.ru/rus/cmfd355}

\crossref{10.22363/2413-3639-2018-64-2-211-426}

%\misctransl

%\by Katrakhov~V.~V., Sitnik~S.~M.

%\paper The transmutation method and boundary-value problems for singular elliptic equations

%\inbook Singular differential equations

%\serial CMFD

%\yr 2018

%\vol 64

%\issue 2

%\pages 211--426

%\publ Peoples' Friendship University of Russia

%\publaddr Moscow

%\lang In Russian

%%\sortdate 1/1/2018

%\crossref{https://doi.org/10.22363/2413-3639-2018-64-2-211-426}



\RBibitem{Kochubei_2}

\by Кочубей~А.~Н.

\paper Диффузия дробного порядка

\jour Дифференц. уравнения

\yr 1990

\vol 26

\issue 4

\pages 660--670

%\sortdate 1/1/1990

\mathnet{http://mi.mathnet.ru/rus/de7144}

\mathscinet{http://www.ams.org/mathscinet-getitem?mr=1061448}

\zmath{https://zbmath.org/?q=an:0712.35049}

%\transl

%\by Kochubei~A.~N.

%\paper Diffusion of fractional order

%\jour Differ. Equ.

%\yr 1990

%\vol 26

%\issue 4

%\pages 485--492

%%\sortdate 1/1/1990

%\mathscinet{http://www.ams.org/mathscinet-getitem?mr=1061448}

%\zmath{https://zbmath.org/?q=an:0729.35064}



\RBibitem{Kochubei_1}

\by Кочубей~А.~Н.

\paper Задача Коши для эволюционных уравнений дробного порядка

\jour Дифференц. уравнения

\yr 1989

\vol 25

\issue 8

\pages 1359--1368

%\sortdate 1/1/1989

\mathnet{http://mi.mathnet.ru/rus/de6938}

\mathscinet{http://www.ams.org/mathscinet-getitem?mr=1014153}

%\transl

%\by Kochubei~A.~N.

%\paper The Cauchy problem for evolution equations of fractional order

%\jour Differ. Equ.

%\yr 1989

%\vol 25

%\issue 8

%\pages 967--974

%%\sortdate 1/1/1989

%\mathscinet{http://www.ams.org/mathscinet-getitem?mr=1014153}



\RBibitem{Kochubei_3}

\by Кочубей~А.~Н., Эйдельман~С.~Д.

\paper Задача Коши для эволюционных уравнений дробного порядка

\jour Докл. РАН

\yr 2004

\vol 394

\issue 2

\pages 159--161

%\sortdate 1/1/2004

\mathnet{http://mi.mathnet.ru/rus/dan1624}

\mathscinet{http://www.ams.org/mathscinet-getitem?mr=2089744}

\zmath{https://zbmath.org/?q=an:1129.35427}

%\misctransl

%\by Kochubeii~A.~N., \'Eidel'man~S.~D.

%\paper The Cauchy problem for evolution equations of fractional order

%\jour Dokl. Akad. Nauk

%\yr 2004

%\vol 394

%\issue 2

%\pages 159--161

%\lang In Russian

%%\sortdate 1/1/2004



\Bibitem{Kochubei_4}

\by Kochubei~A.~N.

\paper Asymptotic Properties of Solutions of the Fractional Diffusion-Wave Equation

\jour Fract. Calc. Appl. Anal.

\yr 2014

\vol 17

\issue 3

\pages 881--896

%\citnumscopus 9

%\sortdate 1/1/2014

\crossref{10.2478/s13540-014-0203-3}

\mathscinet{http://www.ams.org/mathscinet-getitem?mr=3260311}

\zmath{https://zbmath.org/?q=an:1323.35197}

\elib{http://elibrary.ru/item.asp?id=24541812}

\scopus{http://www.scopus.com/record/display.url?origin=inward&eid=2-s2.0-84904312827}



\Bibitem{Kochubei_5}

\by Kochubei~A.~N.

\paper Cauchy problem for fractional diffusion-wave equations with variable coefficients

\jour Applicable Analysis

\yr 2014

\vol 93

\issue 10

\pages 2211--2242

%\citnumscopus 7

%\sortdate 1/1/2014

\crossref{10.1080/00036811.2013.875162}

\mathscinet{http://www.ams.org/mathscinet-getitem?mr=3240385}

\zmath{https://zbmath.org/?q=an:1297.35274}

\elib{http://elibrary.ru/item.asp?id=24642244}

\scopus{http://www.scopus.com/record/display.url?origin=inward&eid=2-s2.0-84905387477}



\RBibitem{Pskhu2}

\by Псху~А.~В.

\paper Фундаментальное решение диффузионно-волнового уравнения дробного порядка

\jour Изв. РАН. Сер. матем.

\yr 2009

\vol 73

\issue 2

\pages 141--182

%\sortdate 1/1/2009

\mathnet{http://mi.mathnet.ru/rus/im2429}

\crossref{10.4213/im2429}

\mathscinet{http://www.ams.org/mathscinet-getitem?mr=2532450}

\zmath{https://zbmath.org/?q=an:1172.26001}

\adsnasa{http://adsabs.harvard.edu/cgi-bin/bib_query?2009IzMat..73..351P}

\elib{http://elibrary.ru/item.asp?id=20425204}

%\transl

%\by Pskhu~A.~V.

%\paper The fundamental solution of a diffusion-wave equation of fractional order

%\jour Izv. Math.

%\yr 2009

%\vol 73

%\issue 2

%\pages 351--392

%%\sortdate 1/1/2009

%\crossref{https://doi.org/10.1070/IM2009v073n02ABEH002450}

%\mathscinet{http://www.ams.org/mathscinet-getitem?mr=2532450}

%\zmath{https://zbmath.org/?q=an:1172.26001}

%\isi{http://gateway.isiknowledge.com/gateway/Gateway.cgi?GWVersion=2&SrcApp=PARTNER_APP&SrcAuth=LinksAMR&DestLinkType=FullRecord&DestApp=ALL_WOS&KeyUT=000266177900006}

%\elib{http://elibrary.ru/item.asp?id=13597908}

%\scopus{http://www.scopus.com/record/display.url?origin=inward&eid=2-s2.0-65349154473}



\RBibitem{Pskhu-SibirJournal}

\by Псху~А.~В.

\paper Уравнение дробной диффузии с оператором дискретно распределенного дифференцирования

\jour Сиб. электрон. матем. изв.

\yr 2016

\vol 13

\pages 1078--1098

%\sortdate 1/1/2016

\mathnet{http://mi.mathnet.ru/rus/semr736}

\crossref{10.17377/semi.2016.13.086}

\mathscinet{http://www.ams.org/mathscinet-getitem?mr=3580054}

\zmath{https://zbmath.org/?q=an:1370.35268}

\elib{http://elibrary.ru/item.asp?id=28127183}

%\misctransl

%\by Pskhu~A.~V.

%\paper Fractional diffusion equation with discretely distributed differentiation operator

%\jour Sib. \`Elektron. Mat. Izv.

%\yr 2016

%\vol 13

%\pages 1078--1098

%\lang In Russian

%%\sortdate 1/1/2016

%\crossref{https://doi.org/10.17377/semi.2016.13.086}



\Bibitem{Pskhu3}

\by Pskhu~A.~V.

\paper Multi-time fractional diffusion equation

\jour Eur. Phys. J.~Spec. Top.

\yr 2013

\vol 222

\issue 8

\pages 1939--1950

%\citnumscopus 2

%\sortdate 1/1/2013

\crossref{10.1140/epjst/e2013-01975-y}

\scopus{http://www.scopus.com/record/display.url?origin=inward&eid=2-s2.0-84885034281}



\RBibitem{Pskhu4}

\by Псху~А.~В.

\paper Первая краевая задача для дробного диффузионно-волнового уравнения в~нецилиндрической области

\jour Изв. РАН. Сер. матем.

\yr 2017

\vol 81

\issue 6

\pages 158--179

%\sortdate 1/1/2017

\mathnet{http://mi.mathnet.ru/rus/im8520}

\crossref{10.4213/im8520}

\mathscinet{http://www.ams.org/mathscinet-getitem?mr=3733319}

\zmath{https://zbmath.org/?q=an:06844352}

\adsnasa{http://adsabs.harvard.edu/cgi-bin/bib_query?2017IzMat..81.1212P}

\elib{http://elibrary.ru/item.asp?id=30737847}

%\transl

%\by Pskhu~A.~V.

%\paper The first boundary-value problem for a fractional diffusion-wave equation in a non-cylindrical domain

%\jour Izv. Math.

%\yr 2017

%\vol 81

%\issue 6

%\pages 1212--1233

%%\sortdate 1/1/2017

%\crossref{https://doi.org/10.1070/IM8520}

%\mathscinet{http://www.ams.org/mathscinet-getitem?mr=3733319}

%\zmath{https://zbmath.org/?q=an:06844352}

%\isi{http://gateway.isiknowledge.com/gateway/Gateway.cgi?GWVersion=2&SrcApp=PARTNER_APP&SrcAuth=LinksAMR&DestLinkType=FullRecord&DestApp=ALL_WOS&KeyUT=000418891300007}

%\scopus{http://www.scopus.com/record/display.url?origin=inward&eid=2-s2.0-85040983678}



\RBibitem{VorKilbas}

\by Ворошилов~А.~А., Килбас~А.~А.

\paper Задача типа Коши для диффузионно-волнового уравнения с частной производной Римана--Лиувилля

\jour Докл. РАН

\yr 2006

\vol 406

\issue 1

\pages 12--16

%\sortdate 1/1/2006

\mathnet{http://mi.mathnet.ru/rus/dan1045}

\mathscinet{http://www.ams.org/mathscinet-getitem?mr=2256848}

\zmath{https://zbmath.org/?q=an:1155.35396}

\elib{http://elibrary.ru/item.asp?id=9176502}

%\misctransl

%\by Voroshilov~A.~A., Kilbas~A.~A.

%\paper A Cauchy-type problem for a diffusion-wave equation with a Riemann-Liouville partial derivative

%\jour Dokl. Akad. Nauk

%\yr 2006

%\vol 406

%\issue 1

%\pages 12--16

%\lang In Russian

%%\sortdate 1/1/2006



\RBibitem{VorKilbas2}

\by Ворошилов~А.~А., Килбас~А.~А.

\paper Задача Коши для диффузионно-волнового уравнения с~частной производной Капуто

\jour Дифференц. уравнения

\yr 2006

\vol 42

\issue 5

\pages 599--609

%\sortdate 1/1/2006

\mathnet{http://mi.mathnet.ru/rus/de11489}

\mathscinet{http://www.ams.org/mathscinet-getitem?mr=2292158}

\zmath{https://zbmath.org/?q=an:1123.35302}

\elib{http://elibrary.ru/item.asp?id=9245260}

%\transl

%\by Voroshilov~A.~A., Kilbas~A.~A.

%\paper The Cauchy problem for the diffusion-wave equation with the Caputo partial derivative

%\jour Differ. Equ.

%\yr 2006

%\vol 42

%\issue 5

%\pages 638--649

%%\citnumscopus 11

%%\sortdate 1/1/2006

%\crossref{https://doi.org/10.1134/S0012266106050041}

%\mathscinet{http://www.ams.org/mathscinet-getitem?mr=2292158}

%\zmath{https://zbmath.org/?q=an:1123.35302}

%\elib{http://elibrary.ru/item.asp?id=27786841}

%\scopus{http://www.scopus.com/record/display.url?origin=inward&eid=2-s2.0-33745304762}



\RBibitem{Gekk}

\by Геккиева~С.~Х.

\paper Задача Коши для обобщенного уравнения переноса с дробной по времени производной

\jour Доклады Адыгской (Черкесской) Международной академии наук

\yr 2000

\vol 5

\issue 1

\pages 16--19

%\sortdate 1/1/2000

%\misctransl

%\by Gekkieva~S.~Kh.

%\paper The Cauchy problem for the generalized transport equation with time-fractional derivative

%\jour Dokl. Adyg. (Cherkess.) Mezdunar. Akad. Nauk

%\yr 2000

%\vol 5

%\issue 1

%\pages 16--19

%\lang In Russian

%%\sortdate 1/1/2000



\RBibitem{Gekkieva1}

\by Геккиева~С.~Х.

\paper Краевая задача для обобщенного уравнения переноса с дробной производной в полубесконечной области

\jour Изв. КБНЦ РАН

\yr 2002

\issue 1(8)

\pages 6--8

%\sortdate 1/1/2002

%\misctransl

%\by Gekkieva~S.~Kh.

%\paper A boundary value problem for the generalized transfer equation with a fractional derivative in a semi-infinite domain

%\jour Izv. KBNTs RAN

%\yr 2002

%\issue 1(8)

%\pages 6--8

%\lang In Russian

%%\sortdate 1/1/2002



\Bibitem{Mainardi-1}

\by Mainardi~F.

\paper The fundamental solutions for the fractional diffusion-wave equation

\jour Appl. Math. Letters

\yr 1996

\vol 6

\issue 1

\pages 23--28

%\citnumscopus 343

%\sortdate 1/1/1996

\crossref{10.1016/0893-9659(96)00089-4}

\mathscinet{http://www.ams.org/mathscinet-getitem?mr=1419811}

\zmath{https://zbmath.org/?q=an:0879.35036}

\scopus{http://www.scopus.com/record/display.url?origin=inward&eid=2-s2.0-30244460855}



\Bibitem{Mainardi-3}

\by Mainardi~F.

\paper The time fractional diffusion-wave equation

\jour Radiophys. Quantum Electron.

\yr 1995

\vol 38

\issue 1-2

\pages 13--24

%\citnumscopus 39

%\sortdate 1/1/1995

\crossref{10.1007/BF01051854}

\mathscinet{http://www.ams.org/mathscinet-getitem?mr=1427165}

\scopus{http://www.scopus.com/record/display.url?origin=inward&eid=2-s2.0-0029427754}



\Bibitem{Mainardi-2}

\by Mainardi~F., Luchko~Yu., Pagnini~G.

\paper The fundamental solution of the space-time fractional diffusion equation

\jour Fract. Calc. Appl. Anal.

\yr 2001

\vol 4

\issue 2

\pages 153--192

\arxiv \href{https://arxiv.org/abs/cond-mat/0702419}{cond-mat/0702419} [cond-mat.stat-mech]

%\sortdate 1/1/2001

\mathscinet{http://www.ams.org/mathscinet-getitem?mr=1829592}

\zmath{https://zbmath.org/?q=an:1054.35156}



\Bibitem{Pagnini_1}

\by Pagnini~G.

\paper The M-Wright function as a generalization of the Gaussian density for fractional diffusion processes

\jour Fract. Calc. Appl. Anal.

\yr 2013

\vol 16

\issue 2

\pages 436--453

%\citnumscopus 22

%\sortdate 1/1/2013

\crossref{10.2478/s13540-013-0027-6}

\mathscinet{http://www.ams.org/mathscinet-getitem?mr=3033698}

\zmath{https://zbmath.org/?q=an:1312.33061}

\elib{http://elibrary.ru/item.asp?id=28682912}

\scopus{http://www.scopus.com/record/display.url?origin=inward&eid=2-s2.0-84879620364}



\Bibitem{Pagnini_Paradisi}

\by Pagnini~G., Paradisi~P.

\paper A stochastic solution with Gaussian stationary increments of the symmetric space-time fractional diffusion equation

\jour Fract. Calc. Appl. Anal.

\yr 2016

\vol 19

\issue 2

\pages 408--440

%\citnumscopus 9

%\sortdate 1/1/2016

\crossref{10.1515/fca-2016-0022}

\mathscinet{http://www.ams.org/mathscinet-getitem?mr=3513003}

\zmath{https://zbmath.org/?q=an:1341.60073}

\elib{http://elibrary.ru/item.asp?id=27104131}

\scopus{http://www.scopus.com/record/display.url?origin=inward&eid=2-s2.0-84969800118}



\Bibitem{Podlubny}

\by Podlubny~I.

\book Fractional differential equations. An introduction to fractional derivatives, fractional differential equations, to methods of their solution and some of their applications

\serial Mathematics in Science and Engineering

\yr 1999

\vol 198

\publ Academic Press

\publaddr San Diego, CA

\totalpages xxiv+340

%\sortdate 1/1/1999

\mathscinet{http://www.ams.org/mathscinet-getitem?mr=1658022}

\zmath{https://zbmath.org/?q=an:0924.34008}



\Bibitem{MathaiSaxenaHaubold}

\by Mathai~A.~M., Saxena~R.~K., Haubold~H.~J.

\book The $H$-function. Theory and Applications

\yr 2010

\publ Springer

\publaddr Dordrecht

\totalpages xiv+268~pp \nofrills

%\citnumscopus 411

%\sortdate 1/1/2010

\crossref{10.1007/978-1-4419-0916-9}

\mathscinet{http://www.ams.org/mathscinet-getitem?mr=2562766}

\zmath{https://zbmath.org/?q=an:1181.33001}

\scopus{http://www.scopus.com/record/display.url?origin=inward&eid=2-s2.0-84891408876}



\RBibitem{KhushtovaVestn2}

\by Хуштова~Ф.~Г.

\paper Задача Коши для уравнения параболического типа~с~оператором Бесселя и~частной производной~Римана--Лиувилля

\jour Вестн. Сам. гос. техн. ун-та. Сер. Физ.-мат. науки

\yr 2016

\vol 20

\issue 1

\pages 74--84

%\sortdate 1/1/2016

\mathnet{http://mi.mathnet.ru/rus/vsgtu1455}

\crossref{10.14498/vsgtu1455}

\zmath{https://zbmath.org/?q=an:06964474}

\elib{http://elibrary.ru/item.asp?id=26898117}

%\misctransl

%\by Khushtova~F.~G.

%\paper Cauchy problem for a parabolic equation with Bessel operator and Riemann--Liouville partial derivative

%\jour Vestn. Samar. Gos. Tekhn. Univ., Ser. Fiz.-Mat. Nauki \rm [J. Samara State Tech. Univ., Ser. Phys. Math. Sci.]

%\yr 2016

%\vol 20

%\issue 1

%\pages 74--84

%%\sortdate 1/1/2016



\Bibitem{GorenfloLuchkoMainardi}

\by Gorenflo~R., Luchko~Y., Mainardi~F.

\paper Analytical properties and applications of the Wright function

\jour Fract. Calc. Appl. Anal.

\yr 1999

\vol 2

\issue 4

\pages 383--414

\arxiv \href{https://arxiv.org/abs/math-ph/0701069}{math-ph/0701069}

%\sortdate 1/1/1999

\mathscinet{http://www.ams.org/mathscinet-getitem?mr=1752379}

\zmath{https://zbmath.org/?q=an:1027.33006}



\Bibitem{tikhonovUsl-1}

\by Tychonoff~A.

\paper Th\'eor\`emes d'unicit\'e pour l'\'equation de la chaleur

\jour Mat. Sb.

\yr 1935

\vol 42

\issue 2

\pages 199--216

\lang In French

%\sortdate 1/1/1935

\mathnet{http://mi.mathnet.ru/rus/sm6410}

\zmath{https://zbmath.org/?q=an:0012.35501|61.1203.05}



\RBibitem{tikhonovUsl-2}

\by Тихонов~А.~Н.

\paper Теоремы единственности для уравнения теплопроводности

\inbook Собрание научных трудов: в десяти томах

\yr 2009

\pages 371--382

\publ Наука

\publaddr М.

\volinfo Т.~II. Математика. Ч.~2. Вычислительная математика 1956--1979. Математическая физика 1933--1948. Ред.-сост. Т.~А.~Сушкевич, А.~В.~Гулин

%\sortdate 1/1/2009

\zmath{https://zbmath.org/?q=an:1200.01055}

%\misctransl

%\by Tikhonov~A.~N.

%\paper Uniqueness theorems for the heat equation

%\inbook Collection of scientific works. In ten volumes

%\yr 2009

%\pages 371--382

%\publ Nauka

%\publaddr Moscow

%\lang In Russian and French

%\volinfo Vol. 2: Mathematics. Part 2: Numerical mathematics 1956–1979. Mathematical physics 1933–1948. Edited by T.~A.~Sushkevich and A.~V.~Gulin

%%\sortdate 1/1/2009

%\zmath{https://zbmath.org/?q=an:1200.01055}



\Bibitem{Wright}

\by Wright~E.~M.

\paper The asymptotic expansion of the generalized hypergeometric function

\jour Proc. Lond. Math. Soc., II. Ser.

\yr 1940

\vol 46

\issue 1

\pages 389--408

%\citnumscopus 136

%\sortdate 1/1/1940

\crossref{10.1112/plms/s2-46.1.389}

\mathscinet{http://www.ams.org/mathscinet-getitem?mr=3876}

\zmath{https://zbmath.org/?q=an:0025.40402|66.1243.01}

\scopus{http://www.scopus.com/record/display.url?origin=inward&eid=2-s2.0-84963078619}



\end{thebibliography}



	





\newpage



\makeentitle



\vspace{-6mm}



%\AdditionalContent{Competing interests}{Specify what conflicts of interest are available.} 

\AdditionalContent{Competing interests}{I have no competing interests.} 

%\AdditionalContent{Competing interests}{We have ... (or) no competing interests.} 

%\AdditionalContent{Authors' contributions and responsibilities}{What is the author's responsibility for each author?}

\AdditionalContent{Author's Responsibilities}{I take full responsibility for submitting the final ma\-nuscript in print. I approved the final version of the manuscript.} 

%\AdditionalContent{Authors' contributions and responsibilities}{Each author has participated in the article concept development and in the manuscript writing. The authors are absolutely responsible for submitting the final manuscript in print. Each author has approved the final version of manuscript.} 



%\AdditionalContent{Funding}{Indicate the source(s) of your funding.}



\AdditionalContent{Funding}{This research received no specific grant from any funding agency in the public,

commercial, or not-for-profit sectors.}



\newpage



\begin{thebibliography}{10}



\Bibitem{Drob:2}

\lang In Russian

\by Nakhushev~A.~M.

\book Drobnoe ischislenie i ego primenenie \rm [Fractional calculus and its applications]

\publaddr Moscow

\publ Fizmatlit 

\totalpages 271

\yr 2003





\Bibitem{Pskhu:w}

\lang In Russian

\by Pskhu~A.~V.

\book Uravneniia v chastnykh proizvodnykh drobnogo poriadka \rm [Partial differential equations of fractional order]

\publaddr Moscow

\publ Nauka 

\totalpages 199 

\yr 2005







\Bibitem{Alexiades:w}

\by Alexiades~V.

\paper Generalized axially symmetric heat potentials and singular parabolic initial boundary value problems

\jour  Arch. Rational Mech. Anal. 

\yr 1982

\vol 79

\issue 4

\pages 325--350

\crossref{10.1007/BF00250797}







\Bibitem{Calton:w}

\by Calton~D.

\paper Cauchy's problem for a singular parabolic partial differential equation

\jour J.~Diff. Equations

\yr 1970

\vol 8

\issue 2

\pages 250--257

\crossref{10.1016/0022-0396(70)90004-5}







\Bibitem{Kepinski_2:w}

\by K\c{e}pi\'nski~S.

\paper \"Uber die Differentialgleichung $\frac{\partial^2 z }{\partial x^2}+\frac{m+1}{x}\frac{\partial z }{\partial x} -\frac{n}{x}\frac{\partial z }{\partial t}=0$

\jour   Math. Ann.

\yr 1905

\vol 61

\issue 3

\pages 397--405

\zmath{36.0416.01}

\lang In German



\Bibitem{Tersenov:en}

\lang In Russian

\by Tersenov~S.~A.

\book Parabolicheskie uravneniia s~meniaiushchimsia napravleniem vremeni \rm [Parabolic Equations with a Changing Time Direction]

\yr 1985

\publ Nauka

\publaddr Moscow

\totalpages 105



\Bibitem{Matiychuk:en}

\by Matiichuk~M.~I.

\book  Parabolichni singuliarni kraiovi zadachi \rm [Parabolic Singular Boundary-Value Problems]

\yr 1999

\publ Institute of Mathematics of the Ukrainian National  Academy of Sciences

\publaddr Kiev

\totalpages 176

\lang In Ukrainian





\Bibitem{Kipriyanov-Monographiya:en}

\lang In Russian

\by Kipriyanov~I.~A.

\book  Singuliarnye ellipticheskie kraevye zadachi \rm [Singular Elliptic Boundary-Value Problems]

\yr 1997

\publ Nauka

\publaddr Moscow

\totalpages 208



\Bibitem{Kipriyanov:en}

\by Kipriyanov~I.~A., Katrakhov~V.~V., Lyapin~V.~M.

\paper On boundary value problems in domains of general type for singular parabolic systems of equations

\jour Sov. Math., Dokl. 

\vol 17

\yr 1976

\pages 1461--1464 

\mathscinet{http://www.ams.org/mathscinet-getitem?mr=0425366}

\zmath{https://zbmath.org/?q=an:0354.35056}





\Bibitem{Sitnik:en}

\lang In Russian

\by Sitnik~S.~M. 

\thesis Application of the Bushman--Erd\'elyi transformation operators and their generalizations in the theory of differential equations with singularities in the coefficients

\thesisinfo Dr. Phys. Math. Sci. Thesis

\yr 2016

\publaddr Voronezh

\totalpages 307







\Bibitem{KatrakhovSitnik:en}

\lang In Russian

\by Katrakhov~V.~V., Sitnik~S.~M.

\paper The transmutation method and boundary-value problems for singular elliptic equations

\inbook Singular differential equations

\serial CMFD

\yr 2018

\vol 64

\issue 2

\pages 211--426

\publ Peoples' Friendship University of Russia

\publaddr Moscow

\crossref{10.22363/2413-3639-2018-64-2-211-426}







\Bibitem{Kochubei_2:en}

\by Kochubei~A.~N.

\paper Diffusion of fractional order

\jour Differ. Equ.

\yr 1990

\vol 26

\issue 4

\pages 485--492







\Bibitem{Kochubei_1:w}

\by Kochubei~A.~N.

\paper The Cauchy problem for evolution equations of fractional order

\jour Differ. Equ.

\yr 1989

\vol 25

\issue 8

\pages 967--974



\Bibitem{Kochubei_3:e}

\lang In Russian

\by Kochubeii~A.~N., \'Eidel'man~S.~D.

\paper The Cauchy problem for evolution equations of fractional order

\jour Dokl. Akad. Nauk

\yr 2004

\vol 394

\issue 2

\pages 159--161







\Bibitem{Kochubei_4:e}

\by Kochubei~A.~N.

\paper Asymptotic Properties of Solutions of the Fractional Diffusion-Wave Equation

\jour Fract. Calc. Appl. Anal.

\yr 2014

\vol 17

\issue 3

\pages 881--896

\crossref{10.2478/s13540-014-0203-3}





\Bibitem{Kochubei_5:e}

\by Kochubei~A.~N.

\paper Cauchy problem for fractional diffusion-wave equations with variable coefficients

\jour  Applicable Analysis

\yr 2014

\vol 93

\issue 10

\pages 2211--2242

\crossref{10.1080/00036811.2013.875162}







\Bibitem{Pskhu2:en}

\paper The fundamental solution of a diffusion-wave equation of fractional order

\by Pskhu~A.~V.

\jour Izv. Math.

\yr 2009

\vol 73

\issue 2

\pages 351--392

\crossref{10.1070/IM2009v073n02ABEH002450}

\isi{http://gateway.isiknowledge.com/gateway/Gateway.cgi?GWVersion=2&SrcApp=PARTNER_APP&SrcAuth=LinksAMR&DestLinkType=FullRecord&DestApp=ALL_WOS&KeyUT=000266177900006}

\elib{http://elibrary.ru/item.asp?id=13597908}

\scopus{http://www.scopus.com/record/display.url?origin=inward&eid=2-s2.0-65349154473}





\Bibitem{Pskhu-SibirJournal:en}

\lang In Russian

\by Pskhu~A.~V.

\paper Fractional diffusion equation with discretely distributed differentiation operator

\jour Sib. \`Elektron. Mat. Izv.

\yr 2016

\vol 13

\pages 1078--1098

\crossref{10.17377/semi.2016.13.086}







\Bibitem{Pskhu3:en}

\by Pskhu~A.~V.

\paper Multi-time fractional diffusion equation

\jour Eur. Phys. J.~Spec. Top.

\yr 2013

\vol 222

\issue 8

\pages 1939--1950

\crossref{10.1140/epjst/e2013-01975-y}





\Bibitem{Pskhu4:en}

\paper The first boundary-value problem for a fractional diffusion-wave equation in a non-cylindrical domain

\by Pskhu~A.~V.

\jour Izv. Math.

\yr 2017

\vol 81

\issue 6

\pages 1212--1233

\crossref{10.1070/IM8520}

\isi{http://gateway.isiknowledge.com/gateway/Gateway.cgi?GWVersion=2&SrcApp=PARTNER_APP&SrcAuth=LinksAMR&DestLinkType=FullRecord&DestApp=ALL_WOS&KeyUT=000418891300007}

\scopus{http://www.scopus.com/record/display.url?origin=inward&eid=2-s2.0-85040983678}







\Bibitem{VorKilbas:a} 

\lang In Russian

\by Voroshilov~A.~A., Kilbas~A.~A.

\paper A Cauchy-type problem for a diffusion-wave equation with a

Riemann-Liouville partial derivative

\jour Dokl. Akad. Nauk

\yr 2006

\vol 406

\issue 1

\pages 12--16



\Bibitem{VorKilbas2:en} 

\by Voroshilov~A.~A., Kilbas~A.~A.

\paper The Cauchy problem for the diffusion-wave equation with the Caputo partial derivative

\mathscinet{http://www.ams.org/mathscinet-getitem?mr=2292158}

\jour Differ. Equ.

\yr 2006

\vol 42

\issue 5

\pages 638--649

\crossref{10.1134/S0012266106050041}



\Bibitem{Gekk:en} 

\lang In Russian

\by Gekkieva~S.~Kh.

\paper The Cauchy problem for the generalized transport equation with time-fractional derivative

\jour Dokl. Adyg. (Cherkess.) Mezdunar. Akad. Nauk

\yr 2000

\vol 5

\issue 1

\pages 16--19







\Bibitem{Gekkieva1:en}

\lang In Russian

\by Gekkieva~S.~Kh.

\paper A boundary value problem for the generalized transfer equation with a fractional derivative in a semi-infinite domain

\jour Izv. KBNTs RAN

\yr 2002

\issue 1(8)

\pages 6--8











  





\Bibitem{Mainardi-1:en}

\by Mainardi~F.

\paper The fundamental solutions for the fractional diffusion-wave equation

\jour Appl.  Math. Letters

\yr 1996

\vol 6

\issue 1

\pages 23--28

\crossref{10.1016/0893-9659(96)00089-4}



\Bibitem{Mainardi-3}

\by Mainardi~F.

\paper  The time fractional diffusion-wave equation

\jour Radiophys. Quantum Electron.

\yr 1995

\vol 38

\issue 1-2

\pages 13--24

\crossref{10.1007/BF01051854}



\Bibitem{Mainardi-2:en}

\by Mainardi~F., Luchko~Yu., Pagnini~G.

\paper The fundamental solution of the space-time fractional diffusion equation

\jour Fract. Calc. Appl. Anal.

\yr 2001

\vol 4

\issue 2

\pages 153--192

\arxiv  	\href{https://arxiv.org/abs/cond-mat/0702419}{cond-mat/0702419} [cond-mat.stat-mech]















\Bibitem{Pagnini_1:en}

\by Pagnini~G.

\paper The M-Wright function as a generalization of the Gaussian density for fractional diffusion processes

\jour Fract. Calc. Appl. Anal.

\yr 2013

\vol 16

\issue 2

\pages 436--453

\crossref{10.2478/s13540-013-0027-6}





\Bibitem{Pagnini_Paradisi:en}

\by Pagnini~G., Paradisi~P.

\paper A stochastic solution with Gaussian stationary increments of the symmetric space-time fractional diffusion equation

\jour Fract. Calc. Appl. Anal.

\yr 2016

\vol 19

\issue 2

\pages 408--440

\crossref{10.1515/fca-2016-0022}





\Bibitem{Podlubny:en}

\by Podlubny~I.

\book Fractional differential equations. An introduction to fractional derivatives, fractional differential equations, to methods of their solution and some of their applications

\yr 1999

\publ Academic Press

\publaddr  San Diego, CA

\zmath{0924.34008}

\serial Mathematics in Science and Engineering

\vol 198

\totalpages xxiv+340





\Bibitem{MathaiSaxenaHaubold:en}

\by Mathai~A.~M., Saxena~R.~K., Haubold~H.~J.

\book The $H$-function. Theory and Applications

\yr 2010

\publ Springer

\publaddr  Dordrecht

\totalpages xiv+268~pp \nofrills

\crossref{10.1007/978-1-4419-0916-9}

\zmath{1181.33001}









\Bibitem{KhushtovaVestn2}

\by Khushtova~F.~G.

\paper Cauchy problem for a parabolic equation with Bessel operator and Riemann--Liouville partial derivative

\jour Vestn. Samar. Gos. Tekhn. Univ., Ser. Fiz.-Mat. Nauki \rm [J. Samara State Tech. Univ., Ser. Phys. Math. Sci.]

\yr 2016

\vol 20

\issue 1

\pages 74--84

\crossref{10.14498/vsgtu1455}





\Bibitem{GorenfloLuchkoMainardi:en}

\by Gorenflo~R., Luchko~Y., Mainardi~F.

\paper Analytical properties and applications of the Wright  function

\jour Fract. Calc. Appl. Anal.

\yr 1999

\vol 2

\issue 4

\pages 383--414

\arxiv \href{https://arxiv.org/abs/math-ph/0701069}{math-ph/0701069}







\Bibitem{tikhonovUsl-1:en}

\by Tychonoff~A.

\paper Th\'eor\`emes d'unicit\'e pour l'\'equation de la chaleur

\jour Mat. Sb.

\yr 1935

\vol 42

\issue 2

\pages 199--216

\lang In French

\mathnet{http://mi.mathnet.ru/msb6410}

\zmath{https://zbmath.org/?q=an:0012.35501|61.1203.05}









\Bibitem{tikhonovUsl-2:en}

\by Tikhonov~A.~N.

\paper Uniqueness theorems for the heat equation

\inbook Collection of scientific works. In ten volumes

\volinfo Vol. 2: Mathematics. Part 2: Numerical mathematics 1956–1979. Mathematical physics 1933–1948. Edited by T.~A.~Sushkevich and  A.~V.~Gulin

\yr 2009

\publ Nauka

\publaddr Moscow

\pages 371--382

\zmath{1200.01055}

\lang In Russian and French





\Bibitem{Wright:en}

\by Wright~E.~M.

\paper The asymptotic expansion of the generalized hypergeometric function

\jour Proc. Lond. Math. Soc., II. Ser. 

\yr 1940

\vol 46

\issue 1

\pages 389--408

\crossref{10.1112/plms/s2-46.1.389}

\zmath{0025.40402|66.1243.01}











\end{thebibliography}



%\end{document}

 \newpage



































%%

%\documentclass[11pt,a4paper,twoside]{article}

\usepackage[utf8]{inputenc}

\usepackage[T1,T2A]{fontenc}

\usepackage[english.ancient,russian.ancient]{babel}

\usepackage{graphicx}

\usepackage{array}

\usepackage[doi]{samgtu-bib}

\usepackage{samgtu}

\usepackage{samgtu-name}

\usepackage{mathtools}

\mathtoolsset{showonlyrefs}

\usepackage{desclist}

\usepackage{euscript}

\usepackage{mathrsfs} 

\usepackage{multicol}

\usepackage{mathrsfs}

\usepackage{cite}

\usepackage{cmap}

\usepackage{qrcode}

\allowdisplaybreaks[4]

\usepackage[nodayofweek]{datetime}

\usepackage[thinlines]{easybmat}



\renewcommand{\JournalYear}{2018}

\renewcommand{\JournalVolume}{22}

\renewcommand{\JournalNumber}{4}



\newdate{JournalDateSign}{29}{12}{2018}% Дата подписания Вестника в печать

\renewcommand{\JournalDateSign}{\displaydate{JournalDateSign}}

\newdate{JournalDatePrint}{16}{01}{2019}% Дата выхода Вестника из печати

\renewcommand{\JournalDatePrint}{\displaydate{JournalDatePrint}}

%\renewcommand{\JournalUPL}{16.56} % Усл. печ. л.

%\renewcommand{\JournalUIL}{16.53} % Уч.-изд. л.

\renewcommand{\JournalUPL}{15.24} % Усл. печ. л.

\renewcommand{\JournalUIL}{15.21} % Уч.-изд. л.

\renewcommand{\JournalRegNumber}{220/18} % Регистрационный номер

\renewcommand{\JournalOrderNumber}{2} % Номер заказа



%\renewcommand{\JournalRMonth}{Март} % Месяц выхода Вестника

%\renewcommand{\JournalEMonth}{March} % Месяц выхода Вестника

%\renewcommand{\JournalRMonth}{Июнь} % Месяц выхода Вестника

%\renewcommand{\JournalEMonth}{June} % Месяц выхода Вестника

%\renewcommand{\JournalRMonth}{Сентябрь} % Месяц выхода Вестника

%\renewcommand{\JournalEMonth}{September} % Месяц выхода Вестника

\renewcommand{\JournalRMonth}{Декабрь} % Месяц выхода Вестника

\renewcommand{\JournalEMonth}{December} % Месяц выхода Вестника



\newcommand{\totop}[1]{\raisebox{6pt}[1pt][2pt]{#1}}

\newcommand{\tottop}[1]{\raisebox{12pt}[1pt][2pt]{#1}} 

\newcommand{\totttop}[1]{\raisebox{18pt}[1pt][2pt]{#1}} 

\newcommand{\tottttop}[1]{\raisebox{24pt}[1pt][2pt]{#1}} 

\newcommand{\totttttop}[1]{\raisebox{30pt}[1pt][2pt]{#1}} 

\newcommand{\tobot}[1]{\raisebox{-6pt}[1pt][2pt]{#1}} 

\newcommand{\tobbot}[1]{\raisebox{-12pt}[1pt][2pt]{#1}} 

\newcommand{\tobbbot}[1]{\raisebox{-18pt}[1pt][2pt]{#1}}



\graphicspath{{./00PIC/}}

%

%

%

%

%\pdfcrop % обрезка pdf для размещения на math-net.ru

%\draft % На корректуре

%

%

%

\FirstPage{388}

\LastPage{397}

%

%

%\begin{document}



\newpage





\selectlanguage{russian}



\phantomsection



\addcontentsline{toc}{section}{Редакционная политика журнала}

\addcontentsline{etoc}{section}{Editorial Policy}







\vsgtu{}

\chead[\fancyplain{}{\selectlanguage{russian}\therutitle}]{\fancyplain{}{\selectlanguage{russian} \therutitle}}

\rhead[]{\fancyplain{\RuJournal}{}}

\lhead[\fancyplain{\RuJournal}{}]{}

\cfoot[]{}

\lfoot[\thepage]{}

\rfoot[]{\thepage}

\renewcommand{\plainheadrulewidth}{0.7pt}

\renewcommand{\headrulewidth}{0.7pt}

\pagestyle{fancyplain}

\thispagestyle{plain}

\normalfont



$\;$



\bigskip



\rutitle{Редакционная политика журнала}





\begin{center}



{\large \bf Редакционная политика журнала «Вестник Самарского государственного технического университета. Серия:~Физико-математические науки»\footnote{Настоящая редакционная политика опубликована на сайте журнала и доступна по ссылке \url{http://www.mathnet.ru/rus/vsgtu/edpolicy}. Всегда проверяйте наличие новой редакции политики журнала на сайте!}}



\smallskip



{\small Редакционная коллегия журнала}



\smallskip



{\small (01 июня 2017 г.)}



\end{center}





\renewcommand{\abstractname}{Общие сведения о журнале}





\small



\begin{abstract}

Журнал «Вестник Самарского государственного технического университета. Серия: Физико"=математические науки» (параллельное название на английском языке "--- Journal of Samara State Technical University, Ser. Physical and Mathematical Sciences) является рецензируемым научным журналом Самарского государственного технического университета и~издаётся с~1996~г.



Самарский государственный технический университет сегодня "--- это один из одиннадцати ведущих опорных университетов России, который является крупнейшим научно-образовательным центром Поволжья с~более чем 100-летней историей.



Журнал долгое время позиционировал себя как издание, предназначенное для публикации новых научных знаний, полученных в российских научных школах. Однако в~настоящее время журнал ставит своей целью открытое распространение научных знаний среди российских и зарубежных учёных, поэтому он также ориентируется и~на зарубежных учёных, работающих в приоритетных научных направлениях Самарского государственного технического университета.



Журнал зарегистрирован в~Федеральной службе по надзору в сфере связи, информационных технологий и массовых коммуникаций (Роскомнадзор), свидетельство \href{https://rkn.gov.ru/mass-communications/reestr/media/?id=83010}{ПИ № ФС 77–66685 от 27.07.2016}. С~2011 года журнал выходит ежеквартально; объём номера "--- 200~c.; язык публикаций "--- русский, английский. Журнал издаётся в~печатной и~электронной формах.



Редакция журнала принимает и оценивает рукописи представленных статей независимо от расы, пола, национальности, происхождения, гражданства (подданства), рода занятий, места работы и проживания автора, а также от его политических, философских, религиозных и иных взглядов.



Представляемая в~журнал рукопись статьи должна быть законченным научным исследованием, нигде ранее не публиковавшимся и~не представленным к~публикации в~других изданиях.



Рукопись статьи должна содержать новые научные результаты по приоритетным направлениям Самарского государственного технического университета, таким как «Дифференциальные уравнения и математическая физика», «Механика деформируемого твёрдого тела», «Математическое моделирование, численные методы и~комплексы программ».



Журнал издается на средства издателя. Все публикации в~журнале бесплатны. Все публикации в~электронном виде распространяются бесплатно.



Целевую аудиторию журнала составляют учёные и исследователи, чьи научные интересы лежат в указанных направлениях: «Дифференциальные уравнения и математическая физика», «Механика деформируемого твёрдого тела», «Математическое моделирование, численные методы и комплексы программ».



\end{abstract}





\bigskip





\noindent

\textbf{1. Цели и задачи}



\begin{enumerate}



\item[1.] Журнал долгое время позиционировал себя как издание, предназначенное для публикации новых научных знаний, полученных в российских научных школах. Однако в~настоящее время журнал ставит своей целью открытое распространение научных знаний среди российских и зарубежных учёных, поэтому он также ориентируется и на зарубежных учёных, работающих в приоритетных научных направлениях Самарского государственного технического университета.



\item[2.] Журнал публикует оригинальные статьи и заказные обзоры по следующим направлениям: «Дифференциальные уравнения и~математическая физика», «Механика деформируемого твёрдого тела», «Математическое моделирование, численные методы и комплексы программ».



\item[3.] Журнал особое внимание уделяет следующим задачам:

\begin{itemize}

\item привлечение высококвалифицированных авторов, редакторов и рецензентов по профилю журнала;

\item обеспечение высокого научного уровня рецензирования всех поступающих материалов;

\item соответствие политики журнала международным стандартам  (стандартам \href{http://publicationethics.org/}{COPE}).

\end{itemize}



\item[4.] Целевую аудиторию журнала составляют учёные и~исследователи, чьи научные интересы лежат в~указанных направлениях: «Дифференциальные уравнения и математическая физика», «Механика деформируемого твёрдого тела», «Математическое моделирование, численные методы и комплексы программ». 



\end{enumerate}



\noindent

\textbf{2. Политика разделов журнала}





\begin{enumerate}



\item[1.] Журнал публикует оригинальные статьи, заказные обзоры и~краткие сообщения по следующим направлениям: «Дифференциальные уравнения и математическая физика», «Механика деформируемого твёрдого тела», «Математическое моделирование, численные методы и комплексы программ».

Журнал может публиковать и другие материалы (биографии, письма к редактору, комментарии и т.п.).



\item[2.] 

Журнал не публикует материалы рекламного характера.

\item[3.] При наличии в редакционном портфеле нескольких работ одного автора в~текущем номере может быть опубликована только одна работа (по выбору редакционной коллегии), вторая "--- в~следующем номере и~так далее. В~исключительных случаях редакционная коллегия может принять решение об опубликовании не более двух работ одного автора. 

\item[4.] 

Весь контент журнала индексируется в~базах данных.



\end{enumerate}





\textbf{2.1. Научные статьи}



\hypertarget{sectionPoliciesOriginalArticles}{}



\begin{enumerate}



\item[1.] В научных статьях должны быть представлены оригинальные исследования с~чётко сформулированными гипотезами, целями и~задачами. Статьи должны содержать новые подходы и новое понимание рассматриваемых проблем.



\item[2.] Все рукописи должны быть подготовлены с помощью \LaTeXe\ (при этом должен быть использован \href{https://www.overleaf.com/read/bzytwtpwfzjj}{наш шаблон}) и отправлены редактору с помощью \href{http://www.mathnet.ru/php/esubmission.phtml?jrnid=vsgtu&wshow=snote&option_lang=rus}{онлайн-службы представления рукописей}.



\item[3.]

Как правило, текст статьи не должен превышать 7500 слов или 30 страниц (без библиографического списка). Статья может включать 8 рисунков или таблиц. Статья должна содержать не менее 10 и не более 50 ссылок. Некоторые рукописи могут содержат больше слов и/или ссылок (по разрешению редактора), другие – меньше.



\item[4.] Все научные статьи должны быть структурированы (например, в формате IMRaD: введение, методы, результаты и дискуссия).



\item[5.] При подготовке рукописи авторы должны пользоваться \href{http://www.mathnet.ru/rus/vsgtu/authornotes#manuscriptPreparingRequirements}{руководством для написания статей}.



\item[6.] Представляемая в~журнал рукопись статьи должна быть законченным научным исследованием, нигде ранее не публиковавшимся и~не представленным к~публикации в~других изданиях. 



\item[7.] Все научные статьи отправляются на рецензирование.



\end{enumerate}





{\bf 2.2. Обзоры}



\begin{enumerate}



\item[1.] Обзор "--- попытка одного или нескольких авторов обобщить текущее состояние исследования по определенной теме (предметной области). В~идеале автор ищет все, что относится к~теме обзора, а~затем сортирует все это и формирует представление о~текущем «состоянии дел». 



\item[2.] Обзор должен содержать информацию об:

\begin{itemize}

\item основных научных публикациях по теме обзора;

\item последние крупные достижения и открытия;

\item значительные пробелы в исследованиях;

\item основные спорные вопросы по теме обзора;

\item основные научные вопросы и перспективы их решения.

\end{itemize}



\item[3.]  Обзоры могут сильно различаться по длине. Обзор может варьироваться от 8000 до 40000 слов или от~30 до 90 страниц (включая библиографический список, таблицы, рисунки).



\item[4.]  Для публикации принимаются обзоры только  от экспертов в~конкретной предметной области.



\item[5.]  Все обзоры отправляются на рецензирование.



\end{enumerate}





{\bf 2.3. Краткие сообщения}



\begin{enumerate} 

\item[1.] Краткое сообщение предназначено для печати небольшого, но независимого результата, вносящего значительный вклад в~область исследования.

\item[2.] Краткое сообщение не должно содержать предварительных результатов, если только эти результаты не представляют исключительный интерес и особенно актуальны.

\item[3.] К подаче и~структуре краткого сообщения предъявляются такие же требования, как и~к~\hyperlink{sectionPoliciesOriginalArticles}{научным статьям}.

\item[4.] Размер краткого сообщения не должен превышать 2500 слов или 10 страниц (без библиографического списка). Оно может включать два рисунка или две таблицы. Краткое сообщение должно содержать не менее 8 ссылок.



\item[5.] Краткие сообщения также отправляются на рецензирование.



\end{enumerate}



\pagebreak



{\bf 2.4. Другие материалы (биографии, письма к редактору, комментарии}



{\bf и т.п.)}



\begin{enumerate}

\item[1.] Размер таких материалов не должен превышать 2500 слов или 10 страниц.

\item[2.] Такие материалы не рецензируются. 

\end{enumerate}



\smallskip



\noindent

{\bf 3. Периодичность}



\smallskip



С~2011 года журнал выходит ежеквартально; объём номера "--- 200~c.; язык публикаций "--- русский, английский.



\smallskip



\smallskip



\smallskip



\noindent

{\bf 4. Политика свободного доступа}



\begin{enumerate}

\item[1.] Журнал сразу предоставляет открытый доступ к~своему контенту исходя из следующего принципа: свободный открытый доступ к результатам исследований способствует увеличению глобального обмена знаниями.



\item[2.] Все публикации журнала в электронном виде распространяются бесплатно и~без ограничений.



\item[3.] Весь контент журнала распространяется на условиях лицензии \href{https://creativecommons.org/licenses/by/4.0/deed.ru}{Creative Co\-m\-mons Attribution 4.0 Inter\-na\-tio\-nal} (\href{https://creativecommons.org/licenses/by/4.0/deed.ru}{CC BY 4.0}; \url{https://creativecommons.org/licenses/by/4.0/deed.ru}).



\item[4.] По решению редакционной коллегии и издателя весь контент журнала, принятый к публикации и/или опубликованный до 01 января 2017, также распространяется по лицензии \href{https://creativecommons.org/licenses/by/4.0/deed.ru}{CC BY 4.0}. 



\item[5.] Если авторы опубликованного ранее контента не согласны с условиями лицензии \href{https://creativecommons.org/licenses/by/4.0/deed.ru}{CC BY 4.0}, то просим их незамедлительно сообщить об этом в редакционную коллегию по e-mail \href{mailto:vsgtu@samgtu.ru}{vsgtu@samgtu.ru}.



\end{enumerate}



\noindent

{\bf 5. Авторские права и политика архивирования}



\smallskip



Авторы, публикующиеся в данном журнале, соглашаются со следующим: 

\begin{enumerate}

\item[1.]  Весь контент журнала распространяется по лицензии \href{https://creativecommons.org/licenses/by/4.0/deed.ru}{CC BY 4.0}. Согласно \href{https://creativecommons.org/licenses/by/4.0/deed.ru}{CC BY 4.0} авторы сохраняют права на свои статьи, но при этом разрешают всем беспрепятственно скачивать, повторно использовать, перепечатывать, изменять, распространять и/или копировать их при условии упоминания их авторства. На всё вышеперечисленное разрешения от авторов или издателя не требуется.



\item[2.]  Авторы сохраняют право заключать отдельные контрактные договорённости, касающиеся неэксклюзивного распространения версии работы в опубликованном здесь виде (например, размещение ее в институтском хранилище, публикацию в книге, перевод на другой язык), со ссылкой на её оригинальную публикацию в~этом журнале.



\item[3.] Согласно классификации \href{http://www.sherpa.ac.uk/romeo/index.php}{SHERPA/RoMEO} журнал относится к так называемым «зеленым» журналам. Авторам разрешено архивировать препринты и постпринты своих работ. Авторы имеют право размещать свою работу в сети Интернет (например, в~институтском хранилище или персональном сайте) до и~во время процесса рассмотрения её в~редакционной коллегии журнала.



\end{enumerate}



\smallskip



\noindent

{\bf 6. Рецензирование}



\begin{enumerate}



\item[1.] Сначала рукопись читается и~проверяется редактором на соответствие руководству для авторов (первоначальная оценка рукописи).



\item[2.]  Несоблюдение изложенных в \href{http://www.mathnet.ru/rus/vsgtu/authornotes#manuscriptPreparingRequirements}{правилах для авторов} критериев может привести к возврату рукописи для исправления без рецензирования.



\item[3.]  Статьи, которые считаются подходящими, направляются как минимум двум независимым рецензентам для оценки научного качества работы.



\item[4.]  Все рецензенты являются признанными специалистами по тематике рецензируемых материалов и~имеют в~течение последних трёх лет публикации по тематике рецензируемой статьи.



\item[5.]  В~журнале действует процесс одностороннего слепого рецензирования или двойного слепого рецензирования, если это возможно сделать для конкретной статьи.



\item[6.] При подготовке рецензии рецензенты обязаны пользоваться инструкциями журнала по рецензированию и~этическими нормами для рецензентов. 



\item[7.] Рецензии хранятся в~редакции журнала в~течение пяти лет.



\item[8.] Рукописи \href{https://goo.gl/wDhKbS}{оцениваются} в соответствии со следующими аспектами:

\begin{itemize}

\item Представлены ли в документе новые концепции (теории), идеи или данные? 

\item Являются ли научные методы и допущения обоснованными и четко обозначенными? 

\item Являются ли результаты достаточными, чтобы обосновать выводы? 

\item Является ли описание экспериментов и~расчётов достаточно полным и точным, чтобы они могли быть воспроизведены другими учеными (возможна ли прослеживаемость результатов)?

\item Придают ли авторы должное внимание аналогичным работам, чётко ли указывается их собственный новый/оригинальный вклад в~развитие проблемы?

\item Насколько точно отражает название статьи её содержание?

\item Отражает ли аннотация основные идеи и результаты исследования?

\item Является ли рукопись статьи хорошо структурированной и~понятной?

\item Насколько хорош и точен язык статьи?

\item Правильно ли определены и используются математические формулы, символы, аббревиатуры и единицы?

\item Должны ли быть какие-либо части документа (текст, формулы, рисунки, таблицы) разъяснены, сокращены, объединены или устранены?

\item Удовлетворителен ли библиографический список?

\end{itemize}



\item[9.] Все вышеперечисленные вопросы будут оцениваться рецензентами в соответствии со следующими оценками: выдающиеся, хорошие, справедливые, посредственные или плохие.



\item[10.] Статья может быть: принята (без дальнейшей доработки); направлена на доработку; отклонена; отклонена (по причине очень предварительных результатов); отклонена (не соответствует тематике журнала).



\item[11.] Если статья направлена на доработку, она должна быть снова отправлена на рассмотрение тем же рецензентам.



\item[12.] Решение о~публикации принимается редакционной коллегией журнала на основании экспертных оценок рецензентов с~учётом соответствия представленных материалов тематической направленности журнала, их научной значимости и~актуальности.

\item[13.] Авторы рукописи имеют право рекомендовать рецензентов "--- специалистов в~данной области (не более двух персон).

\item[14.]  Выбор рецензентов редакционная коллегия оставляет за собой.

\item[15.] Редакционная коллегия не привлекают к~работе рецензентов, которые могут иметь конфликт интересов. Авторы могут сообщать имена людей, назначение которых рецензентами нежелательно из-за возможности возникновения конфликта интересов, объяснив свою позицию.



\end{enumerate}



\smallskip



\noindent

{\bf 7. Индексирование}



\begin{enumerate}

\item[1.] Журнал включен в  \href{http://elibrary.ru/rsci_about.asp}{Russian Science Citation Index} на платформе \href{http://wokinfo.com/products_tools/multidisciplinary/rsci/}{Web of Sci\-ence}. 



\item[2.] Журнал индексируется в реферативных базах данных \href{http://www.viniti.ru}{ВИНИТИ РАН}.

\item[3.] Статьи журнала индексируются также в \href{http://scholar.google.com/scholar?q=Vestn.+Samar.+Gos.+Tekhn.+Univ.+Ser.+Fiz.-Mat.+Nauki}{Scholar.Google.com}, \href{https://goo.gl/zrvH0K}{OCLC WorldCat}, Bielefeld Academic Search Engine (BASE), Open Access Infrastructure for Re\-search in Europe (OpenAIRE), Research Papers in Economics (RePEc), Socionet, \href{http://cyberleninka.ru/journal/n/vestnik-samarskogo-gosudarstvennogo-tehnicheskogo-universiteta-seriya-fiziko-matematicheskie-nauki}{СyberLeninka.ru}, \href{http://goo.gl/EeAvfJ}{Math-Net.ru}, \href{https://elibrary.ru/title_about.asp?id=5784}{eLibrary.ru}. 

\item[4.] Журнал интегрирован в поисковые системы 

\href{http://search.crossref.org/}{CrossRef} и \href{http://search.crossref.org/funding}{FundRef}.

\item[5.] Информация о журнале опубликована в ULRICH’S Periodical Directory. 



\end{enumerate}



\smallskip

\smallskip



\noindent

{\bf 8. Этика научных публикаций}\footnote{Настоящие положения в полном объёме взяты с сайта журнала «Вестник ПНИПУ. Механика» (\url{http://vestnik.pstu.ru/mechanics/toauthors/ethics/})}



\smallskip



\smallskip



\smallskip



\noindent

{\bf 8.1. Общие положения}



\begin{enumerate}

\item[\bf 1.1.] Настоящей политикой устанавливаются стандарты этического поведения для сторон, участвующих в процессе публикации: авторов, редакции, рецензента, издателя. Нижеперечисленные стандарты основаны на общепринятой и существующей политике журнала и издателя.

\item[\bf 1.2.]  Федеральное государственное бюджетное образовательное учреждение высшего образования «Самарский государственный технический университет», как учредитель и издатель научного журнала «Вестник Самарского государственного технического университета. Серия: Физико-математические науки», принимает на себя обязательства по контролю за всеми этапами публикации статей и признает свои этические и другие обязательства, связанные с опубликованием статей.

\end{enumerate}





\smallskip



\noindent

{\bf 8.2. Этические стандарты, предъявляемые к авторам публикаций}





\begin{enumerate}

\item[\bf 2.1.] {\bf Стандарт доступа к исходным данным исследования и их хранения.}

Автор обязан представить исходные материалы (данные) исследования по требованию редакции, если в статье не приводятся исходные данные, и должен быть готов предоставить публичный доступ к ним. Автор должен хранить эти данные в течение разумного времени после публикации для возможности их воспроизведения и проверки.



\item[\bf 2.2.] {\bf Стандарт оригинальности (недопустимости плагиата и самоплагиата).}

Автор представляет в редакцию для рассмотрения статью, содержащую результаты оригинального исследования. Если автор в статье использовал работы или включает в свою статью фрагменты из работ (цитаты) других лиц, то такое использование должно быть надлежащим образом оформлено путем указания оригинального источника в библиографическом списке к статье. Плагиат, равно как и автоплагиат, в любой форме является неэтичным и неприемлемым поведением автора.



\item[\bf 2.3.] {\bf Стандарт однократности публикации.}

Автор представляет в редакцию рукопись статьи, ранее не публиковавшейся и не переданной в редакции других журналов. Подача рукописи одновременно в несколько журналов является неэтичной и неприемлема. % Это же касается перевода статьи на иностранный язык.



\item[\bf 2.4.] {\bf  Стандарт подтверждения источников.}

Автор обязуется правильно указать научные и иные источники, которые он использовал в ходе исследования и которые оказали существенное влияние на результаты исследования, в библиографическом списке. Информация, полученная из неофициальных (частных) источников, не должна использоваться при оформлении научной статьи.



\item[\bf 2.5.] {\bf  Стандарт авторства рукописи статьи.}

Все лица, внесшие значительный вклад в получение результатов исследования, должны быть указаны в качестве соавторов статьи. Список авторов должен быть ограничен лишь этими лицами.

Автор, представляющий редакции рукопись, гарантирует, что им указаны все соавторы, что все они видели и одобрили окончательный вариант рукописи и согласны с её представлением в редакцию научного журнала «Вестник Самарского государственного технического университета. Серия: Физико-математические науки» для публикации.



\item[\bf 2.6.] {\bf  Стандарт раскрытия конфликта интересов со стороны автора.}

Авторы должны раскрывать конфликты интересов, которые могут повлиять на оценку и интерпретацию их рукописи. Все источники финансовой поддержки проекта (гранты, госпрограммы, проекты и проч.) должны быть раскрыты и в обязательном порядке указаны в рукописи.



\item[\bf 2.7.] {\bf  Стандарт исправления ошибок в опубликованных работах.}

Если автор обнаружит существенную ошибку или неточность в уже опубликованной статье, то он обязан незамедлительно уведомить об этом редакцию и содействовать ей в исправлении ошибки. Если редакция узнает об ошибке от третьих лиц, то автор обязан незамедлительно устранить ошибку или представить доказательства ее отсутствия.



\end{enumerate}



\smallskip



\noindent

{\bf 8.3. Этические стандарты, предъявляемые к редакции}





\begin{enumerate}

\item[\bf 3.1.] {\bf  Стандарт принятия решения об опубликовании статьи.}

Редакция журнала принимает решение о том, какие из статей, представленных в редакцию, должны быть опубликованы, на основании результатов проверки на предмет выполнения требований к оформлению и результатов рецензирования.

При принятии решения о публикации рукописи редакция руководствуется политикой журнала и не допускает публикацию статей с признаками клеветы, оскорбления, плагиата или нарушения авторских прав. Окончательное решение о публикации статьи или об отказе в таковой принимается главным редактором журнала.



\item[\bf 3.2.] {\bf  Стандарт равенства всех авторов.}

Редакция оценивает рукописи представленных статей независимо от расы, пола, национальности, происхождения, гражданства (подданства), рода занятий, места работы и проживания автора, а также от его политических, философских, религиозных и иных взглядов.

\item[\bf 3.3.] {\bf   Стандарт конфиденциальности.}

Редакция обязуется не раскрывать информацию о представленной рукописи никому, кроме автора, рецензента (при одностороннем слепом рецензировании), а при необходимости "--- издателя.



\item[\bf 3.4.] {\bf  Стандарт раскрытия конфликта интересов со стороны редакции.}\linebreak

Редакция гарантирует, что материалы рукописи, отклоненной от публикации, не будут использоваться в собственных исследованиях членов редколлегии без письменного согласия автора. Редколлегия будет отказываться от рассмотрения рукописи, если она состоит в каких-либо конкурентных отношениях с автором или организацией, связанной с результатами исследования, либо имеется иной конфликт интересов. Редакция обязуется требовать от всех участников процесса опубликования статьи раскрытия конкурирующих интересов.



\item[\bf 3.5.] {\bf  Стандарт рассмотрения претензий на неэтичное поведение автора.}

Редакция оперативно рассматривает каждую претензию на неэтичное поведение авторов рукописей и уже опубликованных статей независимо от времени её получения. Редакция обязуется принимать адекватные разумные меры в отношении таких претензий. Если доводы претензии подтвердятся, то редакция вправе отказаться от публикации статьи, прекратить дальнейшее сотрудничество с автором, опубликовать соответствующее опровержение, а также принять иные необходимые меры для дальнейшего пресечения неэтичного поведения данного автора.





\end{enumerate}



\smallskip



\noindent

{\bf 8.4. Этические стандарты, предъявляемые к рецензентам}



\begin{enumerate}

\item[\bf 4.1.] {\bf Стандарт вклада рецензента в редакционные решения.}

Предоставленная рецензентом экспертная оценка рукописи способствует принятию редакционных решений, а также помогает автору улучшить рукопись. Решение о принятии рукописи к публикации, возвращении ее автору на доработку или отклонении от публикации принимается редколлегией на основании результатов рецензирования.



\item[\bf 4.2.] {\bf Стандарт сроков рецензирования.}

Рецензент обязан предоставить рецензию в сроки, указанные редакцией. Если рассмотрение рукописи и подготовка рецензии в эти сроки невозможны, то рецензент должен уведомить об этом редакцию.



\item[\bf 4.3.] {\bf Стандарт конфиденциальности со стороны рецензента.}

Рукопись статьи, представленная на рецензирование, должна рассматриваться как конфиденциальный документ, независимо от избранной журналом формы рецензирования. Рецензент вправе показывать её или обсуждать с другими лицами только с разрешения главного редактора.

\item[\bf 4.4.] {\bf  Стандарт объективности рецензии.}

Рецензент обязуется проводить экспертную оценку рукописи объективно. Личная критика автора рецензентом недопустима. Все выводы рецензента должны быть строго аргументированы и снабжены ссылками на авторитетные источники.



\item[\bf 4.5.] {\bf  Стандарт подтверждения источников.}

Рецензенты должны указать при наличии работы, которые оказали влияние на результаты исследования, но не были приведены автором. Рецензент обязан обратить внимание редакции на существенное сходство или совпадение между рассматриваемой рукописью и ранее опубликованной другой работой, о которой известно рецензенту.



\item[\bf 4.6.] {\bf  Стандарт раскрытия конфликта интересов.}

Рецензент не может использовать материалы неопубликованной рукописи в своих собственных исследованиях без письменного согласия автора. Рецензент обязан отказаться от рассмотрения рукописи, в связи с которой у него возникает конфликт интересов из-за конкурентных, совместных или иных отношений с автором либо организацией, имеющей отношение к рукописи. 



\end{enumerate}



\smallskip



\pagebreak



\noindent

{\bf 9. Стоимость публикации}



\begin{enumerate}

\item[1.]

Журнал издается на средства издателя.

\item[2.]

Все публикации в журнале бесплатны (нет сборов за публикацию и подготовку статьи).

\end{enumerate}





\noindent

{\bf 10. Политика раскрытия и конфликты интересов}



\begin{enumerate}

\item[1.]

Неопубликованные данные, полученные из представленных к рассмотрению рукописей, не могут быть использованы в личных исследованиях без письменного согласия автора.

Информация или идеи, полученные в ходе рецензирования и связанные с возможными преимуществами, должны сохраняться конфиденциальными и не использоваться с целью получения личной выгоды.

\item[2.]

Все авторы обязаны раскрыть финансовые или другие явные или потенциальные конфликты интересов, которые могут быть восприняты как оказавшие влияние на результаты или выводы, представленные в работе.

\item[3.]

Рецензенты не должны участвовать в рассмотрении рукописей в случае наличия конфликтов интересов вследствие конкурентных, совместных и других взаимодействий и отношений с любым из авторов, компаниями или другими организациями, связанными с представленной работой.

\end{enumerate}



\noindent

{\bf 11. Положение о конфиденциальности}

\begin{enumerate}

\item[1.] Следующая информация об авторах, включая их имена, аффилиацию, контактные сведения, сведения о занимаемой должности и академической степени, ссылки на авторские профили, размещается в статье в разделе «Сведения об авторе».

\item[2.]

Паспортные данные и место жительство авторов, предоставленные в редакционную коллегию журнала, используются исключительно для заключения лицензионных договоров и пересылки корреспонденции и не передаются третьим лицам.

\item[3.]

Имена рецензентов не сообщаются авторам и третьим лицам.

\end{enumerate}





\noindent

{\bf 12. Положения о плагиате, дублировании и избыточности публикации}



\begin{enumerate}

\item[1.]  Каждая статья проверяется на наличие неправомерных заимствований из других работ (плагиат), дублирование и избыточность публикации (включая перевод публикации с других языков).

\item[2.]  

Для выявления плагиата редакционная коллегия использует следующие ресурсы: Антиплагиат (\url{http://samgtu.antiplagiat.ru/}) и Exactus Like (\url{http://like.exactus.ru/index.php/ru/}).

\item[3.]  

Статья, содержащая неправомерные заимствования из других работ (плагиат),  признаки дублирования и избыточности (в том числе признаки автоплагиата), является неприемлемой. Такая статья отклоняется без рецензирования.

\item[4.]  

Расширенные результаты, представленные на конференциях и семинарах, публикуемые в журнале по согласованию с редакционной коллегией и имеющие соответствующие ссылки, не считаются дублированной и избыточной публикацией.

\item[5.]  

Авторам разрешается подавать в журнал статьи, опубликованные на серверах препринтов, кроме опубликованных ранее в других изданиях. Такая статья не считается дублированной и избыточной.

\item[6.]  

Если статья опубликована, но позднее было обнаружено, что она содержит неправомерные заимствования из других работ (плагиат) или является дублированной, то редактор обращается к \href{http://publicationethics.org/resources/flowcharts}{блок-схемам} Комитета по издательской этике (\href{http://publicationethics.org/}{COPE}) для удаления такой статьи.

\end{enumerate}



\noindent

{\bf 13. Повторное использование текста}



\begin{enumerate}

\item[1.] 

Доля повторно использованного текста (текста из других авторских произведений) в представляемой в журнал работе не должна превышать 15\% от общего объёма статьи (не считая списка литературы).

\item[2.]

Повторное использование текста допустимо лишь для включения необходимых определений, формулировки основных принципов и т.п., чтобы читатель мог ознакомиться со статьей без привлечения дополнительных источников. В этом случае повторное использование текста не считается автоплагиатом.

\item[3.]

Если авторы считают необходимым ознакомить читателя с материалами, близкими к тематике статьи или оставшимися за рамками статьи, то они могут приводить достаточно полный список литературы, в которой интересующийся читатель может найти эти материалы.

\end{enumerate}



%\end{document}





\end{document}



\newpage



$\;$





\chead[\fancyplain{}{}]{\fancyplain{}{}}

  \rhead[]{\fancyplain{}{}}

  \lhead[\fancyplain{}{}]{}

  \cfoot[]{}

  \lfoot[]{}

  \rfoot[]{}

\renewcommand{\plainheadrulewidth}{0pt}

\renewcommand{\headrulewidth}{0pt}





  \pagestyle{fancyplain}



  \thispagestyle{plain}





\newpage



$\;$





\chead[\fancyplain{}{}]{\fancyplain{}{}}

  \rhead[]{\fancyplain{}{}}

  \lhead[\fancyplain{}{}]{}

  \cfoot[]{}

  \lfoot[]{}

  \rfoot[]{}

\renewcommand{\plainheadrulewidth}{0pt}

\renewcommand{\headrulewidth}{0pt}





  \pagestyle{fancyplain}



  \thispagestyle{plain}





%\newpage

%

%$\;$

%

%

%\chead[\fancyplain{}{}]{\fancyplain{}{}}

%  \rhead[]{\fancyplain{}{}}

%  \lhead[\fancyplain{}{}]{}

%  \cfoot[]{}

%  \lfoot[]{}

%  \rfoot[]{}

%\renewcommand{\plainheadrulewidth}{0pt}

%\renewcommand{\headrulewidth}{0pt}

%

%

%  \pagestyle{fancyplain}

%

%  \thispagestyle{plain}







\end{document}





%%%%

%%%   Самарский государственный технический университет

%%%   Кафедра прикладной математики и информатики

%%%

%%%   Преамбула для статьи в журнал "Вестник Самарского государственного

%%%   технического университета. Серия: «Физико-математические науки»"

%%%   Ориентирована под MikTEx 2.4 for Win32

%%%

%%%   Хотя сам журнал собирается под GNU/Linux на дистрибутиве TeXLive



\documentclass[11pt,a4paper,twoside]{article}



\usepackage[cp1251]{inputenc}

\usepackage[T2A]{fontenc}

\usepackage[english,russian]{babel}

\usepackage{samgtubib}

\usepackage[dvips]{graphicx}

\usepackage[multiple]{footmisc}



\usepackage{samgtu}

\renewcommand{\VestnikYear}{2009} % Год издания Вестника

\renewcommand{\VestnikNumber}{2\,(19)} % Номер Вестника

\renewcommand{\VestnikMonth}{Ноябрь} % Месяц выхода Вестника

\renewcommand{\VestnikData}{01.11.09} % Дата подписания Вестника в печать

\renewcommand{\VestnikUPL}{26,64} % Усл. печ. л.

\renewcommand{\VestnikUIL}{25,73} % Уч.-изд. л.

\renewcommand{\VestnikRegNumber}{479/09} % Регистрационный номер

\renewcommand{\VestnikOrder}{994} % Номер заказа







%

%

%

%\begin{document}







\renewcommand{\firstpage}{244}

\renewcommand{\lastpage}{254}



%\Partition{Механика}{Mechanics} % Раздел Вестника, в который подаётся статья





\udk{539,3}

\msc{74K20}



\title{Статический изгиб и колебания многослойной прямоугольной пластинки из ортотропного материала при свободном опирании кра\"eв}%

{Статический изгиб и колебания многослойной прямоугольной пластинки\dots}%





\author{П.\,Ф.~Недорезов, А.\,В.~Аристамбекова}%

{Н\,е\,д\,о\,р\,е\,з\,о\,в~П.\,Ф.,\ А\,р\,и\,с\,т\,а\,м\,б\,е\,к\,о\,в\,а~А.\,В.}



\institute{Саратовский государственный университет им.~Н.\,Г.~Чернышевского \\ (национальный исследовательский университет), механико-математический факультет,\\ 

410012, Саратов, ул.~Астраханская, 83.}



\email{E-mails}{p1934n@yandex.ru, aristambekovaav@mail.ru}



\otherinf{\fullauthor{Пётр Феодосьевич Недорезов} \regalia{д.\,т.\,н., проф.} \position{профессор} \dept{каф. математической теории упругости и~биомеханики}

\fullauthor{Асель Валерьевна Аристамбекова} \position{аспирант} \dept{каф. математической теории упругости и~биомеханики}

}



\keywords{ортотропная пластинка, изгиб, вибрация, численное решение, дискретная ортогонализация}



\workabstract{Рассматривается задача изгиба и колебания многослойной прямоугольной пластинки, состоящей из $N$ ортотропных слоёв произвольной толщины. 

Края пластинки предполагаются свободно опёртыми. 

Обсуждается возможность получения аналитического решения. 

Предложена и реализована методика численного решения. Результаты расчётов представлены в~виде таблиц.}



%%%%%%%%%%%%%%%%%%%%%%%%%%%%%%%%%%%%%%%%%%%%%%%%%%%%%%%%%%%%%%%%%%%%%%%%%%%%%%%%%%%%%%%%%%%%%%%%%%%%%



\engtitle{Static bending and vibrations of multilayer orthotropic  rectangular plate with simply supported edges}



\engauthor{P.\,Ph.~Nedorezov, A.\,V.~Aristambekova}

{N\,e\,d\,o\,r\,e\,z\,o\,v~P.\,Ph., A\,r\,i\,s\,t\,a\,m\,b\,e\,k\,o\,v\,a~A.\,V.}



\enginstitute{N.\,G.~Chernyshevsky Saratov State University (National Research University),\\ Faculty of Mathematics and Mechanics, \\

83, Astrakhanskaya st., Saratov, 410012, Russia.}%



\otherenginfm{\fullauthor{Pyotr Ph. Nedorezov} \regalia{Dr. Sci. (Techn.)} \position{Professor} \dept{Dept. of Mathematical Theory of Elasticity and Biomechanics}

\fullauthor{Asel' V. Aristambekova} \position{Postgraduate Student} \dept{Dept. of Mathematical Theory of Elasticity and Biomechanics}}



\engabstract{Bending and vibrations of multilayer rectangular plate are considered. The plate is composed from $N$ orthotropic layers of arbitrary thickness. Edges of the plate are simply supported. Possibility of closed-form solution is discussed. Numerical solution method is developed and implemented. Results are presented in table.}



\engkeywords{orthotropic  plate, bending, vibration, numerical solution, discrete orthogonalization}



\date{30/XII/2010} % Дата поступления статьи в редакцию.

\revisiondate{28/II/2011} % В окончательном виде



%%%%%%%%%%%%%%%%%%%%%%%%%%%%%%%%%%%%%%%%%%%%%%%%%%%%%%%%%%%%%%%%%%%%%%%%%%%%%%%%%%%%%%%%



\maketitle





\Section{Постановка задачи. Основные формулы и уравнения}

Рассматривается толстая прямоугольная пластинка с~размерами в плане $a\times b$ и толщиной $H$, состоящая из жёстко соединённых между собой ортотропных слоёв с толщинами $h_k$ ($k = 1, 2, \dots, N $) (см. рис.~\hyperlink{fig:ned:1}{1}).





Предполагается, что главные направления анизотропии в~каждой точке слоёв параллельны рёбрам пластинки. 

Пластинка изгибается поперечной нагрузкой интенсивности $q(x, y, t)=q_0(x, y)\exp(i\omega t)$, распределённой по плоскости $z=H$; плоскость $z=0$ от нагрузки свободна. 

Края $x=0$, $x=a$ и~$y=0$, $y=b$ считаются свободно опёртыми. 

Деформации пластинки малы и~подчиняются закону Гука~\cite{ned1}.



Для упрощения выкладок в дальнейшем принимается, что

\begin{equation*}

q(x, y)=p_0\sin\frac{ m\pi x }{a}\sin\frac{n\pi y}{b},\quad p_0=\rm const.

%\label{eq:ned:1}

\end{equation*}





Система уравнений для определения проекций вектора смещения $u^{(k)}$, $v^{(k)}$, $w^{(k)}$ и напряжений $\sigma_x^{(k)}$, $\sigma_y^{(k)}$, $\ldots$,$\tau_{xy}^{(k)}$  $k$-того слоя состоит  из уравнений движения сплошной среды и~уравнений обобщённого закона Гука.



\begin{figure}[t!]

\centering \small \hypertarget{fig:ned:1}{}

\includegraphics[scale=.9]{./00Pic/nedorezov1.eps} \\

%\caption{\label{fig:ned:1}}

Рис. 1

\end{figure}





Считая колебательный процесс установившимся, все характеристики напряжённо"=деформированного состояния (НДС)  $k$-того слоя пластинки ($b_{k-1} \hm\le z\le b_k$; $b_0=0$, $b_k=b_{k-1}+h_k$; $k = 1, 2, \dots, N $) будем искать в виде $f^{(k)}(x, y, z, t)=\tilde f^{(k)}(x, y, z)\exp(i\omega t)$.

Тогда для амплитудных значений смещений и напряжений после отделения временной переменной $t$ получим

\begin{gather}

\label{eq:ned:2}

\frac{{\partial\tilde\sigma_x^{(k)}}}{{\partial x}}+\frac{{\partial\tilde\tau_{xy}^{(k)}}}{{\partial y}}+\frac{{\partial\tilde\tau_{xz}^{(k)}}}{{\partial z}}+\rho_k\omega^2\tilde u^{(k)}=0,\\%\quad

%

%\left({x\Leftrightarrow y\Leftrightarrow z,\quad k=\overline{1,N}}\right),

%

\label{eq:ned:3}

\tilde\sigma_x^{(k)}=A_{11}^{(k)}\frac{{\partial\tilde u^{(k)}}}{{\partial x}}+A_{12}^{(k)}\frac{{\partial\tilde v^{(k)}}}{{\partial y}}+A_{13}^{(k)}\frac{{\partial\tilde w^{(k)}}}{{\partial z}},\quad

%%

%%

\tilde\tau_{yz}^{(k)}=A_{44}^{(k)}\Bigl({\frac{{\partial\tilde v^{(k)}}}{{\partial z}}+\frac{{\partial\tilde w^{(k)}}}{{\partial y}}}\Bigr);

\\

\notag

x\Leftrightarrow y\Leftrightarrow z,\quad 1\Leftrightarrow 2\Leftrightarrow 3,\quad 4\Leftrightarrow 5\Leftrightarrow 6,\quad 

\tilde u^{(k)}\Leftrightarrow\tilde v^{(k)} \Leftrightarrow\tilde w^{(k)},\quad  k=1, 2, \dots, N.

\notag

\end{gather}

Здесь  $A_{ij}^{(k)}$ "--- упругие постоянные, $\rho_k$ "--- плотность материала.



Решение уравнений (\ref{eq:ned:2}), (\ref{eq:ned:3}) должно быть подчинено условиям на боковых гранях

\begin{equation}

\label{eq:ned:4}

\begin{array}{lll}

\textrm{при} & x=0,\,\, x=a: &\tilde v^{(k)}=\tilde w^{(k)}=0,\,\,\tilde\sigma_x^{(k)}=0,\\

%

\textrm{при} & y=0,\,\, y=b: &\tilde u^{(k)}=\tilde w^{(k)}=0,\,\,\tilde\sigma_y^{(k)}=0;\\

\end{array}

\end{equation}

на верхней и нижней плоскостях:

\begin{equation}

\label{eq:ned:5}

\begin{array}{lll}

\textrm{при} & z=0: &\tilde\sigma_z^{\left(1\right)}=0,\quad\tilde\tau_{xz}^{\left(1\right)}=0,\quad\tilde\tau_{yz}^{\left(1\right)}=0, \vspace{1mm}\\

%

\textrm{при} & z=H: &\displaystyle\tilde\sigma_z^{\left(1\right)}=-q_0\sin\frac{{m\pi\,x}}{a}\,\sin\frac{{n\pi\,y}}{b},\quad\tilde\tau_{xz}^{\left(1\right)}=0,\quad

\tilde\tau_{yz}^{\left(1\right)}=0;\\

\end{array}

\end{equation}

на поверхностях контакта смежных слоёв $(k=1, 2, \dots, N-1)$ при $z=b_k$:

\begin{equation}

\label{eq:ned:6}

\begin{array}{c}

\tilde u^{(k)}=\tilde u^{\left({k+1}\right)},\quad \tilde v^{(k)}=\tilde v^{\left({k+ 1}\right)},\quad \tilde w^{(k)}=\tilde w^{\left({k+1}\right)},\\

%

\tilde\sigma_z^{(k)}=\tilde\sigma_z^{\left({k+1}\right)},\quad \tilde\tau_{xz}^{(k)}=

\tilde\tau_{xz}^{\left({k+1}\right)},\quad \tilde\tau_{yz}^{(k)}=\tilde\tau_{yz}^{\left({k+1}\right)}.

\end{array}

\end{equation}



Для понижения размерности трехмерной краевой задачи \eqref{eq:ned:2}--\eqref{eq:ned:6} представим функции $\tilde u^{(k)}$, $\tilde v^{(k)}$,

$\tilde w^{(k)}$, $\tilde\sigma_x^{(k)}$, $\tilde\sigma_y^{(k)}$, \ldots, $\tilde\tau_{xy}^{(k)}$ в виде

\begin{equation}

\label{eq:ned:7}

\begin{array}{l}

\tilde u^{(k)}=HU^{(k)}(\varsigma)\cos m\pi\xi\sin n\pi\eta,\\

\tilde v^{(k)}=HV^{(k)}(\varsigma)\sin m\pi\xi\cos n\pi\eta,\\ 

%

\tilde w^{(k)}=HW^{(k)}(\varsigma)\sin m\pi\xi\sin n\pi\eta,\\ 

%%

%%

\tilde\sigma_x^{(k)}=S_x^{(k)}(\varsigma)\sin m\pi\xi\sin n\pi\eta\quad (x\Leftrightarrow y\Leftrightarrow z),\\ 

%

\tilde\tau_{yz}^{(k)}=T_{yz}^{(k)}(\varsigma)\sin m\pi\xi\cos n\pi\eta,\\ 

%%

\tilde\tau_{xz}^{(k)}=T_{xz}^{(k)}(\varsigma)\cos m\pi\xi\sin n\pi\eta,\\

\tilde\tau_{xy}^{(k)}=T_{xy}^{(k)}(\varsigma)\cos m\pi\xi\cos n\pi\eta,

\end{array}

\end{equation}

где  $\xi={x/a}$, $\eta={y/b}$, $\varsigma={z/{H}}$ "--- безразмерные переменные, а $U^{(k)}$, $V^{(k)}$ и~$W^{(k)}$ "--- безразмерные амплитуды составляющих вектора смещения.



Тогда граничные условия \eqref{eq:ned:4} будут выполнены автоматически, а из \eqref{eq:ned:2} и \eqref{eq:ned:3} для новых неизвестных функций после некоторых преобразований получим уравнения

\begin{equation}

\label{eq:ned:8}

\begin{array}{c}

\displaystyle\frac{{dT_{xz}^{(k)}}}{{d\varsigma}}+\tilde h\pi\left({mS_x^{(k)} -ncT_{xy}^{(k)}}\right)+\rho_k H^2\omega ^2 U^{(k)}=0, \vspace{1mm}\\ 

%%

%%

\displaystyle\frac{{dT_{yz}^{(k)}}}{{d\varsigma}}-\tilde h\pi\left({mT_{xy}^{(k)} -ncS_y^{(k)}}\right)+\rho_k H^2\omega ^2 V^{(k)}=0, \vspace{1mm}\\ 

%%

%%

\displaystyle\frac{{dS_z^{(k)}}}{{d\varsigma}}-\tilde h\pi\left({mT_{xz}^{(k)}+ncT_{yz}^{(k)}}\right)+\rho_k H^2\omega ^2 W^{(k)}=0;

\end{array}

\end{equation}

\begin{equation}

\label{eq:ned:9}

\begin{array}{c}

\displaystyle S_x^{(k)}=-\tilde h\pi\left({mA_{11}^{(k)} U^{(k)}+ncA_{12}^{(k)} V^{(k)}}\right)+ A_{13}^{(k)}\frac{{dW^{(k)}}}{{d\varsigma}},\\ [8pt]

%%

%%

\displaystyle S_y^{(k)}=-\tilde h\pi\left({mA_{12}^{(k)} U^{(k)}+ncA_{22}^{(k)} V^{(k)}}\right)+ A_{23}^{(k)}\frac{{dW^{(k)}}}{{d\varsigma}},\\ [8pt]

%%

%%

\displaystyle S_z^{(k)}=-\tilde h\pi\left({mA_{13}^{(k)} U^{(k)}+ncA_{23}^{(k)} V^{(k)}}\right)+ A_{33}^{(k)}\frac{{dW^{(k)}}}{{d\varsigma}},\\ [8pt]

%%

%%

\displaystyle T_{yz}^{(k)}=A_{44}^{(k)}\Bigl({\frac{{dV^{(k)}}}{{d\varsigma}}+\tilde h\pi nc W^{(k)}}\Bigr),\\ [8pt]

%%

%%

\displaystyle T_{xz}^{(k)}=A_{55}^{(k)}\Bigl({\frac{{dU^{(k)}}}{{d\varsigma}}+\tilde hm\pi W^{(k)}}\Bigr),\\ [8pt]

%%

%%

\displaystyle T_{xy}^{(k)}=A_{66}^{(k)}\tilde h\pi\Bigl({ncU^{(k)}+ mV^{(k)}}\Bigr).

\end{array}

\end{equation}

Граничные условия для этих уравнений следуют из \eqref{eq:ned:5} и \eqref{eq:ned:6}.



%%%%%%%%%%%%%%%%%%%%%%%%%%%%%%%%%%%%%%%%%%%%%%%%%%%%%%%%%%%%%%%%%%%%%%%%%%%%%%%%%%%%%%%%%



\smallskip



\Section{Аналитическое решение}

Для системы уравнений \eqref{eq:ned:8}, \eqref{eq:ned:9} нетрудно получить точное аналитическое решение.



\smallskip

%%%%%%%%%%%%%%%%%%%%%%%%%%%%%%%%%%%%%%%%%%%%%%%%%%%%%%%%%%%%%%%%%%%%%%%%%%%%%%%%%%%%%%%%%

{\it а\/\rm) \it  Решение в перемещениях\/}.

Подставляя  из \eqref{eq:ned:9} в \eqref{eq:ned:8} выражения функций $S_x^{(k)}$, $S_y^{(k)}$,~$\ldots$, $T_{xy}^{(k)}$, получаем систему уравнений для $U^{(k)}$, $V^{(k)}$ и $W^{(k)}$:

\begin{multline*}

A_{55}^{(k)}\frac{{d^2U^{(k)}}}{{d\varsigma^2}}=\left[{\tilde h^2\pi^2\left({A_{11}^{(k)} m^2+ A_{66}^{(k)} n^2 c^2}\right)-\rho_k H^2\omega ^2}\right]U^{(k)}(\varsigma)+\\

%

\null\hfill\displaystyle+\left({A_{12}^{(k)}+A_{66}^{(k)}}\right)\tilde h^2\pi ^2 mncV^{(k)}(\varsigma)-\left({A_{13}^{(k)}+A_{55}^{(k)}}\right)\tilde hm\pi\frac{{dW^{(k)}}}{{d\varsigma}},

\end{multline*}

\begin{multline}

\label{eq:ned:10}

A_{44}^{(k)}\frac{{d^2 V^{(k)}}}{{d\,\varsigma ^2}}=\left({A_{12}^{(k)}+A_{66}^{(k)}}\right)\tilde h^2\pi ^2 mncU^{(k)}(\varsigma)+\\

%

+\left[ {\tilde h^2\pi ^2\left({A_{22}^{(k)} n^2 c^2+A_{66}^{(k)} m^2}\right)-\rho_k H^2\omega ^2}\right]V^{(k)}(\varsigma)-\\

%

-\left({A_{23}^{(k)}+A_{44}^{(k)}}\right)\tilde h\pi nc\frac{{dW^{(k)}}}{{d\varsigma}},

\end{multline}

\begin{multline*}

A_{33}^{(k)}\frac{{d^2 W^{(k)}}}{{d\varsigma ^2}}=\tilde h\pi\Bigl[ {\left({A_{13}^{(k)}+A_{55}^{(k)}}\right)\,m\frac{{dU^{(k)}}}{{d\varsigma}}+\left({A_{23}^{(k)}+A_{44}^{(k)}}\right)nc\frac{{dV^{(k)}}}{{d\varsigma}}}\Bigr]+\\

%

+\left[ {\tilde h^2\pi ^2\left({A_{44}^{(k)} n^2 c^2+A_{55}^{(k)} m^2}\right)-\rho_k H^2\omega ^2}\right]W^{(k)}(\varsigma),

\end{multline*}

где $\tilde h={H/{a}}$, $k= 1, 2, \dots, N$.



Для нахождения общего решения системы уравнений \eqref{eq:ned:10} наиболее рациональным, по нашему мнению, является подход, методика которого приводится, например, в \cite[стр. 153--158]{ned2}. Согласно этой методике одна из искомых функций, например $W^{(k)}(\varsigma)$, принимается за основную.

Тогда последовательным дифференцированием третьего уравнения \eqref{eq:ned:10} вычисляются производные функции $W^{(k)}(\varsigma)$ до шестого порядка включительно, причём из результатов дифференцирования исключаются функции $U^{(k)}(\varsigma)$, $V^{(k)}(\varsigma)$ и их производные. 

В результате таких операций для функции $W^{(k)}(\varsigma)$ получается дифференциальное уравнение шестого порядка

\begin{equation}

\label{eq:ned:11}

\frac{{d^6 W^{(k)}}}{{d\,\varsigma ^6}}+d\frac{{d^4 W^{(k)}}}{{d\,\varsigma ^4}}+f\frac{{d^2 W^{(k)}}}{{d\,\varsigma ^2}}+gW^{(k)}(\varsigma)=0

\end{equation}

с известными коэффициентами $d$, $f$ и $g$, а функции $U^{(k)}(\varsigma)$ и $V^{(k)}(\varsigma)$ представляются линейными комбинациями вида

\begin{equation}

\label{eq:ned:12}

\begin{array}{l}

\displaystyle U^{(k)}(\varsigma)=\alpha_1\frac{{d^5 W^{(k)}}}{{d\varsigma ^5}}+\beta_1\frac{{d^3 W^{(k)}}}{{d\varsigma ^3}}+\gamma_1\frac{{dW^{(k)}}}{{d\varsigma}}, \vspace{1mm} \\

%

\displaystyle V^{(k)}(\varsigma)=\alpha_2\frac{{d^5 W^{(k)}}}{{d\varsigma ^5}}+\beta_2\frac{{d^3 W^{(k)}}}{{d\varsigma ^3}}+\gamma_2\frac{{dW^{(k)}}}{{d\varsigma}}.

\end{array}

\end{equation}



Решая уравнение \eqref{eq:ned:11}, получаем выражение $W^{(k)}(\varsigma)$, в котором будет содержаться шесть произвольных постоянных. 

В силу \eqref{eq:ned:12} эти же постоянные будут линейно входить в формулы для $U^{(k)}(\varsigma)$ и $V^{(k)}(\varsigma)$. 

Тогда число неизвестных постоянных в выражениях $U^{(k)}$, $V^{(k)}$ и $W^{(k)}$ для каждого слоя равно шести, и они определяются из уравнений, которые следуют из \eqref{eq:ned:5} и \eqref{eq:ned:6}.



\smallskip

%%%%%%%%%%%%%%%%%%%%%%%%%%%%%%%%%%%%%%%%%%%%%%%%%%%%%%%%%%%%%%%%%%%%%%%%%%%%%%%%%%%%%%%%%

{\it б\/\rm ) \it Решение в смешанной форме\/}.

Решение рассматриваемой задачи можно получить и в смешанной форме, если в качестве основных неизвестных в~$k$-том слое принять $U^{(k)}$, $V^{(k)}$, $W^{(k)}$, $S_z^{(k)}$, $T_{xz}^{(k)}$

и $T_{yz}^{(k)}$.



Тогда для этих функций из \eqref{eq:ned:8} и \eqref{eq:ned:9} получается система уравнений

\begin{gather*}

\frac{{dU^{(k)}}}{{d\varsigma}}=-\tilde hm\pi W^{(k)} +\frac{{T_{xz}^{(k)}}}{{A_{55}^{(k)}}},\quad 

%

\frac{{dV^{(k)}}}{{d\varsigma}}=-\tilde hnc\pi W^{(k)} +\frac{{T_{yz}^{(k)}}}{{A_{44}^{(k)}}},\\ 

%

\frac{{d\,W^{(k)}}}{{d\varsigma}}=\frac{1}{{A_{33}^{(k)}}}\left[ {\tilde h\pi\left({A_{13}^{(k)} mU^{(k)}+A_{23}^{(k)} ncV^{(k)}}\right)+S_z^{(k)}}\right],\\ 

%%

\frac{{dS_z^{(k)}}}{{d\varsigma}}=\tilde h\pi\left[ {mT_{xz}^{(k)}+ncT_{yz}^{(k)}}\right]-\rho_k H^2\omega ^2 W^{(k)},

\end{gather*}

\begin{multline}

\label{eq:ned:13}

\frac{dT_{xz}^{(k)}}{d\varsigma}=

\biggl\{\tilde h^2\pi^2 \biggl[\biggl(A_{11}^{(k)}- \frac{\bigl({A_{13}^{(k)}}\bigr)^2}{A_{33}^{(k)}}\biggr)m^2+ A_{66}^{(k)}n^2c^2\biggr]-\rho_kH^2\omega^2\biggr\}U^{(k)}+\\ 

%

+\tilde h^2\pi ^2 mnc\biggl({A_{12}^{(k)}+A_{66}^{(k)} -\frac{{A_{13}^{(k)} A_{23}^{(k)}}}{{A_{33}^{(k)}}}}\biggr)V^{(k)} -\tilde hm\pi\frac{{A_{13}^{(k)}}}{{A_{33}^{(k)}}}S_z^{(k)},

\end{multline}

%%

\begin{multline*}

\frac{{dT_{yz}^{(k)}}}{{d\varsigma}}=\tilde{h}^2mnc\pi ^2\biggl({A_{12}^{(k)}+A_{66}^{(k)} -\frac{{A_{13}^{(k)} A_{23}^{(k)}}}{{A_{33}^{(k)}}}}\biggr)U^{(k)} +\\ 

%

+\biggl\{ {\tilde h^2\pi ^2\biggl[ {\biggl({A_{22}^{(k)} -\frac{{\bigl({A_{23}^{(k)}}\bigr)^2}}{{A_{33}^{(k)}}}}\biggr)n^2 c^2+A_{66}^{(k)} m^2}\biggr]-\rho_k H^2\omega ^2}\biggr\}V^{(k)}

%

-\tilde hnc\pi\frac{{A_{23}^{(k)}}}{{A_{33}^{(k)}}}S_z^{(k)},

\end{multline*}

где $k=1, 2, \dots, N$, $b_{k-1}\le z\le b_k$.



В результате применения к уравнениям \eqref{eq:ned:13} методики \cite{ned2} для функции $W^{(k)}(\varsigma)$  снова получается дифференциальное уравнение вида \eqref{eq:ned:11}, и его общее решение по"=прежнему будет содержать шесть произвольных постоянных. 

Остальные неизвестные функции будут выражены через производные от  $W^{(k)}(\varsigma)$: $T_{xz}^{(k)}$ и $T_{yz}^{(k)}$ "--- через нулевую, вторую и~четвёртую производные, а~$U^{(k)}$, $V^{(k)}$ и $S_z^{(k)}$  "--- через первую, третью и пятую.



Из-за сложности зависимости корней $\lambda_1$, $\lambda_2$, $\dots$, $\lambda_6$ характеристического уравнения для уравнения \eqref{eq:ned:11} от значений упругих постоянных $A_{ij}^{(k)}$, параметров нагрузки $m$, $n$ и $\omega$ и толщин слоёв $h_k$ выполнить анализ влияния этих параметров на распределение напряжений и смещений по толщине пластинки можно только после выполнения численных расчётов. 

%Поэтому для систем уравнений \eqref{eq:ned:10} или \eqref{eq:ned:13} целесообразно использовать их непосредственное численное решение краевых задач . 

Ниже приводится вариант методики получения численного решения для системы \eqref{eq:ned:13}, так как для  систем такого вида условия контакта слоёв записываются в наиболее простой форме.





\smallskip

%%%%%%%%%%%%%%%%%%%%%%%%%%%%%%%%%%%%%%%%%%%%%%%%%%%%%%%%%%%%%%%%%%%%%%%%%%%%%%%%%%%%%%%%%

\Section{Численное решение}

В настоящее время известен ряд эффективных численных методов решения краевых задач для систем обыкновенных дифференциальных уравнений. 

К~их числу относится метод дискретной ортогонализации С.~К.~Годунова, который предназначен для нахождения решения краевой задачи для системы обыкновенных дифференциальных уравнений первого порядка

\begin{equation}

\frac{{d\overline Y (\varsigma)}}{{d\varsigma}}=D\overline Y(\varsigma).

\label{eq:ned:14}

\end{equation}

Высокая точность метода Годунова проверена разными авторами при решении большого количества тестовых задач.



Чтобы привести совокупность систем \eqref{eq:ned:13} к виду \eqref{eq:ned:14}, введём в~рассмотрение вектор"=функцию

$\overline Y(\varsigma)=\bigl( y_0(\varsigma), y_1(\varsigma), \dots, y_5(\varsigma) \bigr)$, положив при $b_{k-1}\hm\le\varsigma\le b_k$ ($k=1, 2, \dots, N$):

\begin{equation}

\begin{array}{lll}

y_0(\varsigma)=U_k(\varsigma),& y_1(\varsigma)=V_k(\varsigma), & y_2(\varsigma)=W_k(\varsigma),\\

y_3(\varsigma)=S_z^{(k)}(\varsigma), & y_4(\varsigma)=T_{xz}^{(k)}(\varsigma),& y_5(\varsigma)=T_{yz}^{(k)}(\varsigma).

\end{array}

\label{eq:ned:15}

\end{equation}

Тогда после перехода в \eqref{eq:ned:13} к безразмерным переменным и~с~учётом \eqref{eq:ned:14} ненулевые компоненты матрицы 

$ D=\bigl(d_{ij}\bigr)_{i=0,\, j=0}^{5,\; 5}$ при $b_{k-1}\hm\le\varsigma\le b_k$ ($k=1, 2, \dots, N$) определятся по следующим формулам:

$$

\begin{array}{l}

d_{02}=-d_{34}=-\tilde hm\pi,\qquad d_{12}=-d_{35}=-\tilde hnc\pi,\qquad d_{04}=1/A_{55}^{(k)},\\

d_{15}=1/{A_{44}^{(k)}},\qquad d_{23}=1/A_{33}^{(k)},\qquad d_{32}=-\rho_K\omega^2H^2,\\

d_{20}=-d_{43}=\tilde hm\pi A_{13}^{(k)}/A_{33}^{(k)},\qquad d_{21}=-d_{53}=\tilde hnc\pi A_{23}^{(k)}/A_{33}^{(k)}, \vspace{1mm}\\

%%

%%

%d_{23}={1/A_{33}^{(k)}},\quad d_{3,2}=-\rho_K\omega ^2 H^2,\\

%%

%%

d_{40}=\tilde h^2\pi ^2\left[ m^2\Bigl(A_{11}^{(k)}  - \bigl(A_{13}^{(k)}\bigr)^2/A_{33}^{(k)} \Bigr)+n^2 c^2 A_{66}^{(k)} \right]-\rho_k\omega ^2 H^2, \vspace{1mm}\\

%%

%%

d_{51}=\tilde h^2\pi ^2\left[ n^2 c^2\Bigl({A_{22}^{(k)} - \bigl(A_{23}^{(k)}\bigr)^2/A_{33}^{(k)}\Bigr)+m^2 A_{66}^{(k)}}\right]-\rho_k\omega ^2 H^2 , \vspace{1mm}\\

%%

%%

d_{41}=d_{50}=\tilde h^2mnc\pi^2\left[{A_{12}^{(k)}+A_{66}^{(k)}- A_{13}^{(k)} A_{23}^{(k)}}/A_{33}^{(k)} \right] .

\end{array}

$$



Граничные условия %для компонент функции $\overline{Y}(\varsigma)$ 

следуют из \eqref{eq:ned:5}, которые после подстановки выражений \eqref{eq:ned:7} с учётом \eqref{eq:ned:15} записываются в~виде

\begin{equation}

\label{eq:ned:16}

H_1\overline Y(0 )=0,\quad H_2\overline Y(1)=\overline g_2,

\end{equation}

где $H_1=\bigl(h_{i j} \bigr)_{i=0,\, j=0}^{2,\; 5}$ и $H_2=\bigl(h_{i+3\,j} \bigr)_{i=0,\, j=0}^{2,\; 5}$  "--- прямоугольные матрицы, у~которых отличны от нуля компоненты $h_{03}=h_{14}=h_{25}=h_{33}=h_{44}=h_{55}=1$, а~$\overline g_2=(-q_0, 0, 0)$.



При решении краевой задачи \eqref{eq:ned:14}, \eqref{eq:ned:16} методом дискретной ортогонализации параметр $M$  ($M{+}1$ "--- число отрезков,

на которое делится интервал $0\le\varsigma\le 1$) следует выбирать так, чтобы точки контакта слоёв $\varsigma=b_k$ $(k=1, 2, \dots, N-1)$ принадлежали множеству точек ортогонализации $\varsigma=j/M$ $(j=0, 1, \dots, M)$. 

Тогда условия контакта слоёв \eqref{eq:ned:6} будут выполнены автоматически.



После решения задачи \eqref{eq:ned:14}, \eqref{eq:ned:16} максимальные амплитуды напряжений $\sigma_x^{(k)}$, $\sigma_y^{(k)}$, $\tau_{xy}^{(k)}$ в~соответствующих точках вычисляются по формулам

\begin{multline*}

S_x^{(k)}=-\tilde hm\pi\left[ A_{11}^{(k)} -\bigl(A_{13}^{(k)}\bigr)^2 /A_{33}^{(k)} \right]y_0(\varsigma)-\\

%

-\tilde hnc\pi\left[ A_{12}^{(k)} -A_{13}^{(k)} A_{23}^{(k)}/ A_{33}^{(k)} \right]y_1(\varsigma)+

{A_{13}^{(k)} / A_{33}^{(k)} y_3(\varsigma)},

\end{multline*}

%

\begin{multline*}

S_y^{(k)}=-\tilde hm\pi\left[ A_{12}^{(k)} -A_{13}^{(k)} A_{23}^{(k)} / A_{33}^{(k)} \right]y_0(\varsigma)-\\

-\tilde hnc\pi\left[ A_{22}^{(k)} -\bigl(A_{23}^{(k)} \bigr)^2 / A_{33}^{(k)} \right] y_1(\varsigma)+A_{23}^{(k)}/A_{33}^{(k)} y_3(\varsigma),

\end{multline*}

%%

$$

T_{xy}^{(k)}=\tilde h\pi A_{66}^{(k)}\left[ {ncy_0(\varsigma)+my_1(\varsigma)}\right].

$$



Отметим, что все приведённые формулы и уравнения остаются в силе и~при $\omega=0$. 

В~этом случае они определяют НДС многослойной пластинки при статическом изгибе.



\smallskip

%%%%%%%%%%%%%%%%%%%%%%%%%%%%%%%%%%%%%%%%%%%%%%%%%%%%%%%%%%%%%%%%%%%%%%%%%%%%%%%%%%%%%%%%%

\Section{Результаты численных расчётов}

Изложенная методика численного решения была применена для исследования НДС при статическом изгибе \mbox{$(\omega{=}0)$} и~определения резонансных частот при установившихся колебаниях $(\omega{\ne}0)$ квадратных однослойных и~двухслойных пластинок с~размерами в плане $a=b=1$~м и толщиной $H=0,1$~м и $H=0,2$~м. Вычисления выполнялись при $q_0=1,0$~Па, $\rho=1500$~кг/м$^{3}$, $m=n=1$.



В качестве параметров модели принимались параметры, соответствующие ортотропным стеклопластикам (однонаправленный стеклопластик и \linebreak СВАМ~5:1). 

Значения технических упругих постоянных для этих материалов, взятые из \cite{ned3}, приведены в табл.~\ref{ned:tab:1}. 

Для нижнего слоя брались параметры материала СВАМ~5:1, для верхнего "--- однонаправленного стеклопластика; 

$h_*=h_1/H \in \{1,0; 0,95; 0,9; 0,75; 0,5\}$.



% Таблица 1

\begin{table}[b!]

\centering

\small 

\caption{\label{ned:tab:1}}

\vspace{-3mm}

\setlength{\extrarowheight}{1pt}

\setlength{\tabcolsep}{3.5pt}

\begin{tabular}{c|c|c|c|c|c|c|c|c}

\hline

\footnotesize$E_1\cdot 10^{-5}$, &\footnotesize $E_2\cdot 10^{-5}$, & \footnotesize$E_3\cdot 10^{-5}$,  & \footnotesize$G_{12}\cdot 10^{-5}$,  &\footnotesize $G_{23}\cdot 10^{-5}$, & \footnotesize$G_{13}\cdot 10^{-5}$, & \tobot{\footnotesize$\nu_{12}$} & \tobot{\footnotesize$\nu_{23}$} & \tobot{\footnotesize$\nu_{13}$}\\

\footnotesize МПа &\footnotesize  МПа &\footnotesize  МПа & \footnotesize МПа & \footnotesize МПа & \footnotesize МПа &  &  & \\

\hline

\multicolumn{9}{c}{\footnotesize СВАМ 5:1}\\

\hline

\rule{0mm}{14pt}$0,359$ & $0,167$ & $0,889$ & $0,0468$ & $0,0484$ & $0,04$ & $0,18$ & $0,38$ & $0,052$\\

\hline

\multicolumn{9}{c}{\footnotesize Однонаправленный стеклопластик}\\

\hline

\rule{0mm}{14pt}$0,505$ & $0,147$ & $0,147$ & $0,050$ & $0,049$ & $0,055$ & $0,268$ & $0,43$ & $0,07$\\

\hline

\end{tabular}

\end{table}



При статическом изгибе ($\omega=0$) величины прогиба $W(\varsigma)$ и~напряжений $\sigma_x^{(k)}(0,5, 0,5, \varsigma)=S_x^{(k)}(\varsigma)$,

$\sigma_y^{(k)}(0,5, 0,5, \varsigma)=S_y^{(k)}(\varsigma)$ $(k = 1, 2)$ для некоторых характерных значений $\varsigma$ приведены в табл.~\ref{ned:tab:2} ($H=0,1$~м) и в табл.~\ref{ned:tab:3} \mbox{($H=0,2$~м).}

В табл.~\ref{ned:tab:2} и последующих значения при $h_*=0$ и $h_*=1$ соответствуют однослойным пластинкам из СВАМ 5:1 ($h_*=1$) и однонаправленного стеклопластика ($h_*=0$).



% Таблица 2

\begin{table}[b!]

\centering

\small 

\caption{\label{ned:tab:2}}

\vspace{-3mm}

\setlength{\extrarowheight}{1pt}

\setlength{\tabcolsep}{8pt}

\begin{tabular}{l|c|c|c|c|c|c}

\hline

~~~~~~~~\footnotesize $h_{*}$ & \footnotesize 1,0 & \footnotesize 0,95 & \footnotesize 0,90 & \footnotesize 0,75 & \footnotesize 0,50 & \footnotesize 0,0\\

\hline

\rule{0mm}{14pt}%

$W(0,0)\cdot 10^9$ & $-17,048$ & $-14,755$ & $-14,728$ & $-14,643$ & $-14,508$ & $-14,192$\\

%\hline

$W(0,5)\cdot 10^9$ & $-17,193$ & $-14,894$ & $-14,867$ & $-14,782$ & $-14,647$ & $-14,332$\\

%\hline

$W(1,0) \cdot 10^9$ & $-1,7103$ & $-14,781$ & $-14754$ & $-14,670$ & $-14,537$ & $-14,224$\\

%\hline

$S_x(0,0)$, Па & 30,589 & 27,166 & 27,078 & 26,800 & 26,631 & 35,640\\

%\hline

$S_x(0,5)$, Па & $-0,080$ & 0,591 & 0,557 & 0,453 & 0,294 & $-0,054$\\

%\hline

$S_x^{(1)}(h_*)$, Па & $-30,723$ & $-23,058$ & $-20,267$ & $-12,261$ & 0,294 & ---\\

%

$S_x^{(2)}(h_*)$, Па & --- & $-32,384$ & $-28,460$ & $-17,208$ & 0,432 & $-35,792$\\

%\hline

$S_x(1,0)$, Па & $-30,723$ & $-36,388$ & $-36,364$ & $-36,284$ & $-36,147$ & $-35,792$\\

%

$S_y(0,0)$, Па & 16,173 & 13,962 & 13,942 & 13,879 & 13,777 & 12,829\\

%

$S_y(0,5)$, Па & $-0,045$ & $-0,117$ & $-0,112$ & $-0,096$ & $-0,070$ & $-0,059$\\

%

$S_y^{(1)}(h_*)$, Па & $-16,289$ & $-12,761$ & $-11,296$ & $-6,998$ & $-0,070$ & ---\\



$S_x^{(2)}(h_*)$, Па &  ---       & $-12,156$ & $-10,764$ & $-6,680$ & $-0,097$ & $-12,962$\\

%

$S_y(1,0)$, Па & $-16,289$ & $-13,538$ & $-13,510$ & $-13,423$ & $-13,284$ & $-12,962$\\

\hline

\end{tabular}

\end{table}



% Таблица 3

\begin{table}[h!]

\centering

\small 

\caption{\label{ned:tab:3}}

\vspace{-3mm}

\setlength{\extrarowheight}{1pt}

\setlength{\tabcolsep}{10pt}

\begin{tabular}{l|c|c|c|c|c|c}

\hline

~~~~~~~~\footnotesize $h_{*}$ & \footnotesize 1,0 & \footnotesize 0,95 & \footnotesize 0,90 & \footnotesize 0,75 & \footnotesize 0,50 & \footnotesize 0,0\\

\hline

\rule{0mm}{14pt}%

$W(0,0)\cdot 10^9$ & $-1,291$ & $-1,150$ & $-1,148$ & $-1,141$ & $-1,127$ & $-1,092$\\

$W(0,5)\cdot 10^9$ & $-1,336$ & $-1,190$ & $-1,188$ & $-1,181$ & $-1,168$ & $-1,133$\\

$W(1,0)\cdot 10^9$ & $-1,346$ & $-1,179$ & $-1,177$ & $-1,170$ & $-1,157$ & $-1,124$\\

$S_x(0,0)$, Па & 7,577 & 7,101 & 7,074 & 6,989 & 6,814 & $-8,874$\\

$S_x(0,5)$, Па & $-0,017$ & $0,236$ & 0,226 & 0,182 & 0,112 & $-0,03$\\

$S_x^{(1)}(h_*)$,  Па & $-7,808$ & $-5,764$ & $-4,945$ & $-2,810$ & 0,112 & ---\\

$S_x^{(2)}(h_*)$, Па &  --- & $-8,064$ & $-6,914$ & $-3,919$ & 0,174 & $-9,105$\\

$S_x(1,0)$, Па & $-7,808$ & $-9,307$ & $-9,302$ & $-9,284$ & $-9,243$ & $-9,105$\\

$S_y(0,0)$, Па & 4,435 & 3,952 & 3,948 & 3,932 & 3,903 & 3,588\\

$S_y(0,5)$, Па & $-0,032$ & $-0,057$ & $-0,054$ & $-0,048$ & $-0,035$ & $-0,052$\\

$S_y^{(1)}(h_*)$, Па & $-4,603$ & $-3,674$ & $-3,225$ & $-1,962$ & $-0,035$ & ---\\

$S_x^{(2)}(h_*)$, Па & --- & $-3,534$ & $-3,107$ & $-1,905$ & $-0,072$ & $-3,749$\\

$S_y(1,0)$, Па & $-4,603$ & $-3,969$ & $-3,961$ & $-3,933$ & $-3,882$ & $-3,749$\\

\hline

\end{tabular}

\end{table}





Как следует из результатов вычислений, добавление к~однослойной пластинке даже очень тонкого слоя из более жёсткого материала ($E_{1}^{(2)}>E_{1}^{(1)}$) при неизменной общей толщине пакета $H$ существенно влияет на прогибы и~напряжения в~пластинке. 

При этом двухслойная пластинка начинает работать приблизительно так же, как однослойная пластинка из более жёсткого материала. 

Однако следует иметь в~виду, что в двухслойных пластинках напряжения $\sigma_x$, $\sigma_y$ и $\tau_{xy}$ на поверхности контакта слоёв $\varsigma=h_*$ терпят разрыв.

Напряжение $\sigma_z$ при статическом изгибе меньше напряжений $\sigma_x$, $\sigma_y$ и по модулю не превосходит значения $q_0=1,0$~Па.



Установившиеся колебания пластинок исследовались в~диапазоне частот $0<\omega<80000$~с$^{-1}$. 

Исследования показали, что у~рассматриваемых пластинок имеется спектр резонансных частот $\omega_k^{\textrm{рез}}$, при которых все характеристики НДС неограниченно возрастают, а~при переходе через значения $\omega=\omega_k^{\textrm{рез}}$ меняют знак. 

Разным значениям $\omega_k^{\textrm{рез}}$ соответствуют различные законы изменения характеристик НДС по толщине пластинки, которые качественно отличаются друг от друга. 

При колебаниях первого типа (их называем чисто изгибными) амплитуда прогиба $W$ постоянна или очень мало меняется по толщине пакета, а~тангенциальные смещения $U$ и~$V$ и~напряжения $\sigma_x$ и~$\sigma_y$ в~пределах каждого слоя меняются по $\varsigma$ по линейному или близкому к нему закону.



Наряду с чисто изгибными в пластинках реализуются колебания другого типа "--- планарные. 

Для таких колебаний характерно линейное со сменой знака или близкое к нему изменение прогиба $W$ (эффект <<разбухания>> или <<подтягивания>>). Тангенциальные смещения при планарных колебаниях очень мало меняются по толщине, а~нормальные напряжения $\sigma_x$ и $\sigma_y$ в~пределах каждого слоя почти постоянны.



Кроме указанных типов колебаний в~пластинках существуют и~другие, которые являются сложными комбинациями чисто изгибных и~планарных с~преобладанием изгибных или планарных составляющих. 

Такие колебания наблюдаются как при высших резонансных частотах, так и~при отличающихся от них.



Резонансные частоты представляются в виде $\omega_k^{\textrm{рез}}=\omega_k+\delta$ при изгибных или изгибно"=планарных колебаниях и~$\omega_k^{\textrm{рез}}=\Omega_k+\delta$, если колебания планарные. 

Значения $\omega_k$ и $\Omega_k$ при $\delta<1$~с$^{-1}$ для пластинок с~$H=0,1$~м для рассматриваемого диапазона изменения $\omega$ приведены в табл.~\ref{ned:tab:4}, а~в~случае $H=0,2$~м "--- в табл.~\ref{ned:tab:5}.



% Таблица 4

\begin{table}[h!] \centering \small

\caption{\label{ned:tab:4}}

\setlength{\extrarowheight}{4pt}

\setlength{\tabcolsep}{20pt}

\vspace{-3mm}

\begin{tabular}{c|c|c|c|c|c}

\hline

\footnotesize$h_{*}$ &\footnotesize $\omega_1$ &\footnotesize $\omega_2$ &\footnotesize $\omega_3$ &\footnotesize $\Omega_1$ &\footnotesize $\Omega_2$\\

\hline

\rule{0mm}{14pt}%

1,00 & 1952 & 57912 & 79994 & 11159 & 16969\\

%\hline

0,95 & 2097 & 57927 & 60616 & 10885 & 19382\\

%\hline

0,90 & 2099 & 57927 & 60630 & 10883 & 19396\\

%\hline

0,75 & 2105 & 57927 & 60668 & 10877 & 19437\\

%\hline

0,50 & 2115 & 57925 & 60718 & 10868 & 19499\\

%\hline

0,00 & 2138 & 57918 & 60803 & 10848 & 19626\\

\hline

\end{tabular}

\end{table}



% Таблица 5

\begin{table}[h!] \centering \small

\caption{\label{ned:tab:5}}

\setlength{\extrarowheight}{4pt}

\setlength{\tabcolsep}{9pt}

\vspace{-3mm}

\begin{tabular}{c|c|c|c|c|c|c|c|c}

\hline

\footnotesize $h_{*}$ &\footnotesize  $\omega_1$ &\footnotesize  $\omega_2$ &\footnotesize  $\omega_3$ &\footnotesize  $\omega_4$ &\footnotesize  $\omega_5$ &\footnotesize  $\omega_6$ &\footnotesize  $\Omega_1$ &\footnotesize  $\Omega_2$\\

\hline

\rule{0mm}{14pt}%

1,00 & 3446 & 31738 & 39560 & 54081 & 58257 & 76715 & 11147 & 16926\\

%\hline

0,95 & 3649 & 30553 & 34389 & 51961 & 58109 & 61944 & 10900 & 19067\\

%\hline

0,90 & 3652 & 30554 & 34428 & 52079 & 58169 & 62071 & 10896 & 19096\\

%\hline

0,75 & 3663 & 30556 & 34535 & 52356 & 58324 & 62396 & 10885 & 19178\\

%\hline

0,50 & 3685 & 30553 & 34696 & 52632 & 58510 & 62763 & 10867 & 19311\\

%\hline

0,00 & 3742 & 30535 & 34979 & 52777 & 58638 & 62990 & 10830 & 19566\\

\hline

\end{tabular}

\vspace{-3mm}

\end{table}



Как видно из этих таблиц, увеличение толщины пакета существенно уплотняет спектр резонансных частот.

При этом первая резонансная частота $\omega_{1}$ "--- частота чисто изгибных колебаний "--- возрастает,

а частоты изгибно"=планарных и планарно"=изгибных колебаний $\omega_{2}$, $\omega_{3},\ldots$ вследствие сжатия частотного спектра убывают.

Частоты $\Omega_{1}$ и $\Omega_{2}$ чисто планарных колебаний зависят в основном от размеров пластинки в плане и в меньшей степени "--- от материала, с изменением $H$ они меняются мало. Изменение толщины $h_2$ слоя из более жесткого материала сказывается на колебаниях пластинки так же, как и в случае статического изгиба.







Некоторые возможные законы изменения по $\varsigma$ амплитуд прогиба $W$ и тангенциального смещения $U$ при различных частотах иллюстрируют графики, представленные соответственно на рис.~\hyperlink{fig:ned:2}{2} и рис.~\hyperlink{fig:ned:3}{3}. Эти графики построены для пластинки с толщиной пакета $H=0,2$\,м при $h_1=0,15$\,м и $h_2=0,05$\,м. 

Кривая 1 на рис.~\hyperlink{fig:ned:2}{2} соответствует частоте $\omega=36000$\,с$^{-1}$, кривая 2 "--- $\omega=37000$\,с$^{-1}$,

кривая 3 "--- $\omega=54900$\,с$^{-1}$, кривая 4 "--- $\omega=59350$\,с$^{-1}$, кривая 5 "--- $\omega=57500$\,с$^{-1}$,

ординаты кривых 2 и 5 увеличены соответственно в 2~и 10~раз. 

На рис.~\hyperlink{fig:ned:3}{3} кривая 1 построена для частоты $\omega=1000$\,с$^{-1}$,

кривая 2 "--- $\omega=19000$\,с$^{-1}$, кривая 3 "--- $\omega=52000$\,с$^{-1}$.



\begin{figure}[h!]

\centering \small \hypertarget{fig:ned:2}{}

\includegraphics[scale=.8]{./00Pic/nedorezov2.eps} \\

%\caption{\label{fig:ned:2}}

Рис. 2

\end{figure}





\begin{figure}[h!]

\centering \small \hypertarget{fig:ned:3}{}

\includegraphics[scale=.8]{./00Pic/nedorezov3.eps}\\

%\caption{\label{fig:ned:3}}

Рис. 3

\end{figure}





\begin{figure}[h!]

\centering \small \hypertarget{fig:ned:4}{}

\includegraphics[scale=.8]{./00Pic/nedorezov4.eps}\\

%\caption{\label{fig:ned:4}}

Рис. 4

\end{figure}



Отдельный интерес представляют величина и закон изменения по толщине нормального напряжения $\sigma_z$, которое в классической и большинстве других приближенных теорий считаются малыми по сравнению с $\sigma_x$ и $\sigma_y$ и поэтому, как правило, не учитываются.

Графики, представленные на рис.~\hyperlink{fig:ned:4}{4} (сплошные линии "--- $\sigma_z$, штриховые "--- $\sigma_x$), показывают, что в толстой двухслойной ортотропной пластинке напряжения $\sigma_z$ при некоторых значениях $\omega$ меняются по толщине по сложному закону и могут быть меньше

(кривые 1, $\omega=59300$\,с$^{-1}$), одного порядка (кривые 2, $\omega=32300$\,с$^{-1}$) или превосходить (кривые 3, $\omega=53500$\,с$^{-1}$) значения $\sigma_x$.



Сплошные линии на этом рисунке изображают поведение напряжения $\sigma_z$, штриховые "--- $\sigma_x$.

При $\omega=59300$\,с$^{-1}$ в точках контакта слоёв разрыв кривой $\sigma_x$ (штриховая линия 1) при выбранном масштабе незначителен и поэтому на графике практически не отображается.



\begin{thebibliography}{9}



\RBibitem{ned1}

\by Лехницкий~С.\,Г.

\book Анизотропные пластинки

\publaddr М.

\publ ГИТТЛ

\yr 1957

\totalpages 463

\transl 

\by Lekhnitskii~S.\,J.

\book Anisotropic Plates

\publ Gordon and Breach

\publaddr New York 

\yr 1968

\totalpages 534



\RBibitem{ned2}

\by Эльсгольц~Л.\,Э.

\book Дифференциальные уравнения

\publaddr М.

\publ ГИТТЛ

\yr 1957

\totalpages 272

\transl

\by El'sgol'ts~L.\,E.

\book Differential Equations

\publ Gordon and Breach

\publaddr New York

\serial Russian Monographs and Texts on Advanced Mathematics and Physics

\vol IV

\yr 1961

\totalpages 360





\RBibitem{ned3}

\by Джунисбеков~Т.\,М., Андрющенко~О.\,В., Кейкиманова~М.\,Т.

\paper Влияние упругих свойств материала трубы на напряжёно"=деформированное состояние подземного нефтепровода

\jour Вестник ТарГУ

\yr 2008

\pages 155--160

\misctransl 

\by Dzhunisbekov~T.\,M., Andryushchenko~O.\,V.,  Keykimanova~M.\,T.

\paper Vliyanie uprugih svoystv materiala truby na napryazhyono-deformiro\-vannoe sostoyanie podzemnogo nefteprovoda 

\jour Vestnik TarGU

\yr 2008 

\issue 1 

\pages 155--160



\end{thebibliography}





\makeengtitlem





\newpage



%\end{document}

%

